\documentclass[11pt,a4paper]{article}
\usepackage[margin=25mm]{geometry}
\usepackage{fontspec}
\setmainfont{DejaVu Serif}
\setsansfont{DejaVu Sans}
\setmonofont[Scale=0.84]{DejaVu Sans Mono}
\usepackage{amsmath,amssymb,amsthm,mathtools,bm,mathrsfs}
\providecommand{\llbracket}{\mathopen{[\![}}
\providecommand{\rrbracket}{\mathclose{]\!]}}
\usepackage{booktabs,longtable,array,enumitem,xcolor,graphicx}
\usepackage{microtype}
\usepackage[unicode,breaklinks=true,colorlinks=true,linkcolor=blue!45!black,urlcolor=blue!45!black]{hyperref}
\usepackage{xurl}
\providecommand{\tightlist}{\setlength{\itemsep}{0pt}\setlength{\parskip}{0pt}}
\newtheorem{proposition}{Proposition}
\newtheorem{lemma}{Lemma}
\newtheorem{theorem}{Theorem}
\newtheorem{corollary}{Corollary}
\theoremstyle{definition}\newtheorem{definition}{Definition}
\DeclareMathOperator{\im}{im}
\DeclareMathOperator{\rank}{rank}
\DeclareMathOperator{\tr}{tr}
\DeclareMathOperator{\End}{End}
\DeclareMathOperator{\Hom}{Hom}
\DeclareMathOperator{\GL}{GL}
\DeclareMathOperator{\SL}{SL}
\DeclareMathOperator{\Sym}{Sym}
\DeclareMathOperator{\Spec}{Spec}
\DeclareMathOperator{\Proj}{Proj}
\DeclareMathOperator{\diag}{diag}
\makeatletter
\renewcommand*\l@subsection{\@dottedtocline{2}{1.5em}{2.8em}}
\renewcommand*\@pnumwidth{2.5em}
\renewcommand*\@tocrmarg{3.5em}
\makeatother
\setcounter{tocdepth}{2}
\setlength{\emergencystretch}{3em}
\setlength{\parskip}{0.3em}
\title{CUE, GCT, and Isolation\\Exact bridges, restricted multiplicity formulas,\\and the remaining derandomization gates}
\author{Research run CUE\_GCT\_20261121\_R1\\GPT-6 Astra Pro; model-assisted proof review and exact computation}
\date{21 November 2026\\Working manuscript; live status in the accompanying ledger}
\begin{document}
\maketitle
\begin{abstract}
We derive exact character, Haar-integral, covariance-rank, and quotient formulas linking finite CUE moments to the multiplicities used in geometric complexity theory. Explicit restrictions and determinantal representations locate low-row exclusion ranges. A reduced determinantal-quartic divisor gives an exact shifted-plethysm difference. A uniform endpoint repair determines the quaternary permanent-image dimension exactly in every size and proves maximal parameter dimension in every larger ambient variable count for sizes at least seventeen. Invariant quotients give the exact eventual four-row rectangular range. Fourier frames and common free tensors extend the exact eventual sign to every strict direction through $2m$ rows for sizes at least seventeen, with explicit stretch bounds and no integer dimension-equality case. A Pieri tail bound and lower-variable boundary comparison classify every fixed direction through $m$ rows eventually, including the full first-row equality face; separate unpadded quotients and a generic-stabilizer upper bound extend the nonrectangular classification through $2m$ rows. Graph isolation is transferred in its valid independent-variable setting and then tested against monomial descent, coefficient cancellation, joint compression, and specialization. Checked rank-profile and matrix-cone constructions yield polynomial hitting sets for every fixed width in an unknown order, with a six-pair refinement at width three and two centered pairs for affine singular-factor words, and explicit arithmetic points with trace-coordinate span recovery. An all-characteristic determinant-weight budget, Grassmann-line connectors and coefficient-cut arguments prove the complete balanced two-pair width-three result through twelve variables in every characteristic and thirteen away from two. A smallest-counterexample support argument additionally proves the full cubic statement in every variable count. Purity gives additional sparse and trace-space constructions; uniformly conditioned energy reconstruction has explicit size barriers. The stated RH and Dirichlet GRH inputs are accepted as requested, and every further analytic dependency is separately identified. General polynomial ROABP hitting sets and the asymptotic determinant-versus-padded-permanent problem remain unresolved. The separate Grothendieck part transports every finding to its precise scheme, representation-ring, filtration, or sheaf framework.
\end{abstract}
\begingroup\small\tableofcontents\endgroup\clearpage
\clearpage\part{Inputs, conventions, and scope}
\section{Scope, theorem inputs, and the central reductions}

The main unrestricted tasks remain open in this manuscript: an explicit polynomial-size hitting set for general width-$w$ ROABPs, a uniformly efficient positive rule for all rectangular and plethysm multiplicities, and a determinant-versus-padded-permanent separation in the required asymptotic regime. Counterexamples below refute specified transfers, not those open problems. Every stated deduction has an identifier. The accompanying event register distinguishes complete proof readings, finite exact computations, dependency checks, and carried background. No Lean reconstruction or formal verification of the new mathematics is claimed.

The principal positive results are exact rather than heuristic. R-01--R-09 identify the CUE character and ratios operations, positive projector formulas, orbit-closure Gram ranks, and the missing quotient numerator. R-10 closes the height-two rectangle and quadratic determinant case. G-05 excludes every at-most-three-row multiplicity obstruction. G-07/R-12 give an exact four-row permanent-minus-determinant shifted-plethysm formula at size four. R-13 and the uniform endpoint repair R-27 now determine the quaternary permanent-image dimension in every size:
\[
 c_m=\min\left\{\binom{m+3}{3},4m^2-2m+2\right\}.
\]
R-24 supplies the exact size-seventeen base certificate; R-25 gives the border induction; R-26 supplies the uniform ordinary-character construction; R-27 recovers its six missing endpoint directions by a checked finite jet certificate and exact central insertion. The single graded Tor defect of R-23 consequently vanishes for every $m\ge17$. G-21 extends maximal parameter dimension to every $r\ge4$: $\dim\mathcal P_{m,r}=rm^2-2m+2$ for $m\ge17$, with the specified low-degree Tor vanishing.

G-17 classifies the strict four-row directions by $\nu_4>0$ and $\nu_1>(n-m)s$, and provides explicit first-stretch and progression bounds. G-18 establishes occurrence on both orbit closures and the all-integer asymptotic ratio. G-20 now gives its exact eventual sign for every pair $n>m\ge3$. For $m\ge17$, permanent dominance holds precisely when $n^2\le2m^2-m+1$, and determinant dominance holds when $n^2\ge2m^2-m+2$, in each fixed strict direction. There is no unresolved integer dimension-equality case. The exact four-row rectangular padding range is $m<n<4m/3$ for every $m\ge14$, and fourteen is sharp for this uniform eventual statement. This does not exclude individual earlier-degree gaps on the opposite side. G-16 specifies a nonrectangular weight family throughout the positive dimension range. G-15 proves dominant third jets and the necessary first-row condition $\lambda_1\ge d(n-m)+4$, while excluding full fourth-jet dominance. R-21 gives an exact orbit-multiplicity algorithm for fixed ambient dimension, with numerical-degree complexity. G-19 identifies one rectangular discriminant copy with zero determinant energy and positive permanent energy in equal size, including an explicit smooth cubic slice and five-row minimality. That specified-copy equation does not compare total multiplicities and vanishes under positive padding.

The higher-row construction G-22 uses an ordered Fourier frame and
cocharacter minimum weights to provide a common stable tensor in
every strict direction with $r\le m$. R-28 removes its scalar-orbit
period and gives $\dim\mathcal D_{n,r}=(r-2)n^2+2$.
For $m\ge17$ and $4\le r\le2m$, G-23 and its R-29/G-25 extension consequently prove
\[
 \frac{m_{j\nu}(\ell^{n-m}\operatorname{per}_m)}
      {m_{j\nu}(\det_n)}
 =\kappa_\nu j^{rm^2-(r-2)n^2-2m+r-1}(1+O(j^{-1})).
\]
Both multiplicities are positive at every sufficiently large integer
stretch, and the exponent is never zero. Its sign, with explicit
bounds for both alternatives, gives the exact eventual comparison
in every stated strict direction. G-24 completes the rectangular boundary comparison: the exact eventual positive range is $m<n<rm/(r-1)$. At equality the determinant eventually dominates, and above it every permanent rectangle vanishes. For the added rows $r>m$, the only positive parameter case is $r=n=m+1$: the explicit direction $(3,2^{m-1},1)$ has ratio $\kappa j(1+O(j^{-1}))$ in degree $2j$. G-26 additionally settles every nonrectangular first-row equality face through $m$ rows: Pieri bounds its multiplicity by the tail partition in one fewer variable, and the determinant eventually dominates. Below that face the padded multiplicity is zero. Thus every fixed direction through $m$ rows is classified eventually. R-32/R-33 and G-27/G-28 additionally classify every fixed nonrectangular direction through $2m$ rows. The lower comparison algebra may have a degree period and need only satisfy a quotient-dimension upper bound. Boundary comparison extends further to $r\le\min\{2n-2,m^2\}$ for nonrectangular labels. Rectangular directions beyond $m$ rows, larger strict directions and individual finite degrees retain their separate scope. R-31 proves a scalar-period obstruction for the unchanged two-product family at $2m+1$ rows in specified strict directions. Its odd-degree vanishing is a subcone statement, not a statement about either full orbit closure.


On isolation, I-01 gives the genuine Hoza transfer for independent vertex variables. I-03--I-04 identify coefficient flags and projected matroid cones; I-14 gives the weaker joint-span condition. The rank-profile construction I-21 and endpoint refinement I-30 supply an explicit six-pair interpolation hitting generator for width at most three in one unknown variable order, in every characteristic with the stated field-node convention. I-29 treats all words having one fixed singular interval rank. The full-matrix cone proof I-31 gives polynomial-size hitting sets for every fixed width and a pair budget controlled by the number of distinct interval ranks. Complete proofs are given here; the rank-profile and cone mechanisms are credited to a recent unrefereed preprint, whose broader affine-orbit claims are not assumed. I-32 refutes a polynomial coordinate-cone deletion bound even at a common center with every interval rank preserved. I-33 instead gives explicit small prime-cocharacter families for two dense products, their exact limitation, and a scalar product-closure repair. I-34 extends the prime family to bounded local nullity among separated products, with an exact three-versus-four-term cylinder distinction. P-14 uses scalar bridges to give $w-1$ interpolation pairs for factors of rank one or full rank, including one pair at width two. The unresolved target is a bound polynomial jointly in the growing width, length and degree; even a polynomial pair budget does not by itself meet that sampling target. I-35 weakens bridge preservation to equality of rank-one annihilator base loci, proves independent-seed composition, and isolates the missing synchronization condition with explicit counterexamples. P-16 gives an explicit degree16, width-exactly-three obstruction to the uncentered two-pair map on eighteen specified algebraic nodes. Its transition factors are generically invertible, and P-14 proves that the centered two-pair map hits the same polynomial. I-36 proves a pairwise-dominant maximal-minor preservation lemma and improves the bound to five pairs when every variable-dependent width-three factor is singular. I-37 reduces this to two centered pairs for the affine singular-factor subclass, in every characteristic and unknown order, by an irreducible-factor reduction and a tridiagonal second-coefficient argument. I-38 extends the same two-pair conclusion to affine words with isolated invertible factors: no two consecutive variable-dependent factors have rank three. Longer invertible runs and nonlinear singular factors remain open for this two-pair argument. I-39 refutes one-pole rank preservation even for an irreducible mixed affine scalar, and separately refutes the proposed abstract Hessian criterion while retaining its integrability limitation. P-17 gives characteristic-free coefficient-cut obstructions, including every balanced two-pair cubic annihilator through ten variables. I-40 bounds full-width invertible runs of a homogeneous translation-invariant tensor by its integer weight $n-2D$. At weight one this gives isolated invertible factors and closes the complete balanced two-pair width-three problem through nine variables, with an explicit finite rational sampler. I-41 uses the saturated weight-two support spaces and their actual divided-square dual pairing to preserve the four-coordinate connectors. Its bandwidth-three argument closes the balanced theorem through ten variables, with the same finite grid. The ordinary unweighted centered map is not identified with that balanced map for every ten-variable polynomial. P-18 proves the trilinear class and a smallest-counterexample support reduction. P-19 closes both remaining boundary cases: every multilinear width-three cubic is hit by balanced two-pair evaluation, in every characteristic and with no restriction on the number of variables, using $20(n+1)^2$ samples. Any remaining balanced width-three annihilator has its homogeneous degrees in the interval $[4,n-4]$. I-42 uses the actual rank-one reset structure to remove the Hessian step through four invertible factors in every characteristic, and five outside characteristic two; it supplies positive saturated-run classes. I-43 closes every weight-three connector: a Grassmann orbit line and the actual adjacent support bound force a rank-one update pattern, while the weight-zero anchor case is handled directly. Shared-parameter contractions then give the complete balanced two-pair theorem through eleven variables in every characteristic. I-44 classifies every chain anchor and preserves all interval ranks through five full-rank cores. The reset criterion yields weights at most four in every characteristic and at most five away from two; the complete balanced theorem follows through twelve variables, and thirteen away from two. Weight five in characteristic two and the next higher-degree cases remain open. The precise conditional matrix-invariant consequence remains I-08.

Purity gives sparse constructions, an actual-sheaf Gram criterion with full conductor accounting, and limitations of energy approximation. P-11 identifies explicit common arithmetic evaluation points with hitting-set encodings in its rational model. P-12 constructs such a point for the width-three class, together with an integer hitting curve and an exact trace-dual coefficient-span map. P-13 states the general centered-interpolation encoding and reduces the width-three field degree to $(nD+1)^{13}$ using I-30. It applies directly to the rank-spectrum construction I-31 with degree $(nD+1)^{2K_{\mathcal R}+1}$. These constructions use elementary algebraic number theory and require no additional RH input.

\subsection{Stipulated analytic inputs and supplied material}

We accept Alp\"oge's RH and Dirichlet GRH as theorem inputs, exactly as requested. In \nolinkurl{ALPOGE_Main.lean}, the exported declarations are \nolinkurl{riemannHypothesis}, \nolinkurl{riemannZeta_eq_zero_re}, \nolinkurl{LFunction_ne_zero_of_half_lt_re}, \nolinkurl{LFunction_eq_zero_re}, and \nolinkurl{LFunction_eq_zero_re_of_isPrimitive}; the predicates are defined in \nolinkurl{ALPOGE_Statements.lean}. The attached note supplies narrative descriptions of its numbered steps. Neither the imported Lean dependency tree nor the claimed all-files scan is in this attachment, and neither was rebuilt in this run. This is an input convention, not a review of the theorem's truth.

Masullo's general number-field Hecke GRH is used only conditionally, as the prompt's attribution to section~1.2. The supplied statement index has no numbered section~1.2 theorem or its proof. Accordingly we give no invented locator. Its Theorem~13.3.1.1, indexed PDF page~522, is a trace-formula claim; Conjecture~14.1.2.1, indexed page~529, concerns automorphic $\mathrm{GL}_n$ GRH and simplicity and remains explicitly conjectural there. These are different statements. The zero-statistics section names the precise hypotheses discharged by the stipulated input and those that remain.

The archive was unpacked once. The source manifest records the actual supplied files and hashes. The five cote~133 files are LaTeX transcriptions, not original page images. Assertions in the accompanying earlier study about visual readings remain credited to that study; our transcription table records mathematical necessity separately from paleographic confidence.

\subsection{Why the three multiplicity problems are distinct}

Let $V=\mathbb C^{n^2}$, $W=\operatorname{Sym}^nV^*$, and $A=\operatorname{Sym}(W^*)$. For a homogeneous form $f\in W$, write
\[
 m_f(\lambda,d)=\operatorname{mult}_{S_\lambda V}\mathbb C[\overline{GL(V)f}]_d,
 \qquad a_\lambda(d[n])=\operatorname{mult}_{S_\lambda V}A_d.
\]
The symbol $g(\lambda,\mu,\nu)$ denotes the ordinary Kronecker coefficient. One must distinguish the ambient plethysm multiplicity, the connected-stabilizer rectangular Kronecker multiplicity, and the coordinate-ring multiplicity of an orbit closure. The transpose component of the determinant stabilizer and the orbit boundary are separate restrictions. R-03 gives finite counterexamples to identifying them. R-05 computes the closure multiplicity by an actual rank, and R-08 expresses the same requirement as a quotient numerator.

If $X\subseteq W$ is $GL(V)$-stable and $V_k$ is spanned by the first $k$ basis vectors, let $X_k$ be the closure of its restriction image in $\operatorname{Sym}^nV_k^*$. For $\ell(\lambda)\leq k$, the multiplicities in $\mathbb C[X]$ and $\mathbb C[X_k]$ agree. The short verification used repeatedly below is as follows. A highest-weight polynomial of weight $(\lambda_1,\ldots,\lambda_k,0,\ldots,0)$ uses only coefficient variables supported on $V_k$, since all generator weights are nonnegative. It vanishes on $X$ exactly when it vanishes on $X_k$. Taking the quotient of the common ambient highest-weight space proves the assertion. This is a coefficient-subalgebra version of the usual inheritance principle, proved here in precisely the form needed for the orbit restrictions.

The exact identities alone impose no efficient complexity bound. In particular, a nonnegative squared-norm formula may have exponentially many terms; a rank formula may require exponentially large matrices; and a positive dimension may become zero after highest-weight projection. The size and sign statements in each result are therefore part of the theorem, not suppressed implementation details.

\subsection{How the specific new analytic tools enter}

Alp\"oge's step~2 (p.~3) concerns theta derivatives and inert-prime divisibility. It does not, in the supplied four-page text, state a bound for a coefficient-matroid minor. Steps~3 and~5 (p.~3) motivate trace and square-root-cancellation arguments; P-02, P-04 and P-06 implement the corresponding curve/sheaf constructions with their exact degree and conductor budgets. Steps~4 and~7 (pp.~3--4) concern determinant and Gan--Gross--Prasad period positivity for specified groups; R-03 identifies why that branching problem cannot be substituted for the tensor subgroup without a proved functorial bridge. Step~8 (p.~4) concerns theta-lift multiplicities; I-03 and I-14 identify the different coefficient-span quantifiers a proposed isolation transfer must preserve.

Step~13 (p.~4) invokes a Bost slope inequality, and step~17 (p.~4) concerns the displayed arithmetic cohomology groups. The supplied note gives no reusable estimates for the ranks, heights or valuations of the ROABP coefficient maps used here. Their possible use therefore has a concrete prerequisite: construct the relevant arithmetic bundle or cohomology object, identify the desired minor or quotient with its invariant, and retain effective size bounds. This is a dependency statement about the supplied inputs, not a claim that such a construction is impossible. The integral and special-fiber analyses below identify where that next construction would act.

\clearpage\part{Checked findings and checked conditional reductions}

% BEGIN SOURCE work/rectangular/findings.tex
% Research fragment prepared by /root/rectangular_bridge in run CUE_GCT-20261121.
% Requires amsmath, amssymb, amsthm; no nonstandard macros are used.
% Citations use bibliographic keys listed in sources.md.
% These are checked deductions in this run, not assertions of literature priority.

\section{Rectangular characters, explicit positive witnesses, and orbit ranks}

Throughout this section the ground field is $\mathbb C$, compact-group Haar
measure has mass one, and $D=nd$. We write $r=(d^n)$ and
$D_\alpha(b)=s_\alpha(1^b)=\dim S_\alpha(\mathbb C^b)$.
Duals are displayed when coordinate rings are involved. Passing from the
polynomial representation to the coordinate representation changes both sides
by the same duality and leaves every multiplicity below unchanged.

\subsection{R-01: the exact rectangular restriction transform}

\paragraph{Statement R-01.}
Let $a\geq\ell(\lambda)$ and let $X=(x_1,\ldots,x_a)$,
$Y=(y_1,\ldots,y_n)$. Denote the product alphabet by
$XY=(x_i y_j)_{1\leq i\leq a,1\leq j\leq n}$. Then
\begin{equation}\label{R01:branching}
 s_r(XY)=\sum_{\mu,\nu\vdash D}
 g(r,\mu,\nu)s_\mu(X)s_\nu(Y).
\end{equation}
In particular, by symmetry of the three Kronecker labels,
\begin{equation}\label{R01:haar}
 g(\lambda,r,r)=\int_{U(a)\times U(n)}
 s_r(A\otimes B)\overline{s_\lambda(A)}\det(B)^{-d}\,dA\,dB.
\end{equation}
Every coefficient needed on the rectangular-Kronecker side of determinant
GCT is therefore an explicit tensor-subgroup Fourier coefficient of the
\emph{same rectangular character} that occurs in the finite CUE moment.
The equality is not an identification of that coefficient with the scalar
moment.

More explicitly, enumerate the alphabet $XY$ as $z_1,\ldots,z_{an}$, with
$z_1,\ldots,z_n=x_1y_1,\ldots,x_1y_n$. Bump--Gamburd's Proposition~4
(source label \nolinkurl{cpaverage}) gives
\begin{equation}\label{R01:CUE}
 s_r(XY)=\left(\prod_{j=1}^n z_j^d\right)
 \int_{U(d)}\prod_{j=1}^n\det(1+z_j^{-1}U^{-1})
 \prod_{j=n+1}^{an}\det(1+z_jU)\,dU.
\end{equation}
This formula is first read with nonzero $z_j$ and then as a polynomial identity.
Its specialization at identities proves the positive weighted sum
\begin{equation}\label{R01:dimension}
 D_r(an)=\sum_{\mu,\nu\vdash D}
 g(r,\mu,\nu)D_\mu(a)D_\nu(n)
\end{equation}
and the explicit upper bound
\begin{equation}\label{R01:bound}
 0\leq g(\lambda,r,r)\leq
 \left\lfloor\frac{D_r(an)}{D_\lambda(a)}\right\rfloor,
 \qquad
 D_r(an)=\prod_{i=1}^{n}\prod_{j=1}^{d}
 \frac{an+j-i}{n+d-i-j+1}.
\end{equation}
The numerator is a finite rectangular-character moment, not an unknown GCT
coefficient.

\begin{proof}
Apply Schur--Weyl duality to $(\mathbb C^a\otimes\mathbb C^n)^{\otimes D}$
in two ways. Decomposing first for $GL_{an}\times S_D$ gives the factors
$S_\eta(\mathbb C^a\otimes\mathbb C^n)\otimes[\eta]$.
Decomposing separately for $GL_a$ and $GL_n$ gives
$S_\mu(\mathbb C^a)\otimes S_\nu(\mathbb C^n)\otimes
([\mu]\otimes[\nu])$, with diagonal $S_D$ action. Comparing the multiplicity
space of $[r]$ proves \eqref{R01:branching}. Schur orthogonality and
$s_r(B)=\det(B)^d$ prove \eqref{R01:haar}. Equation \eqref{R01:CUE} is the
specialization $N=d,L=n,K=(a-1)n$ of the cited product identity.
Set every $x_i,y_j$ equal to one in \eqref{R01:branching}. Every summand is
nonnegative. Retain the term $(\mu,\nu)=(\lambda,r)$ and use
$D_r(n)=1$; this proves \eqref{R01:bound}. The displayed product is the
hook-content formula in Bump--Gamburd, equation \nolinkurl{e_hhl} in \S2.2.
\end{proof}

\paragraph{Finite moments and their limiting constants.}
Bump--Gamburd, Corollary~1, gives
\[
 \mathbb E_{U(N)}|\det(1-U)|^{2k}=D_{(N^k)}(2k).
\]
The limit after division by $N^{k^2}$ is
$c_k=\prod_{j=0}^{k-1}j!/(j+k)!$. Thus the finite moment is an integer
dimension, whereas the leading constant generally is not: $c_2=1/12$.
The distinction prevents a scalar asymptotic constant from being treated as
a representation multiplicity.

\subsection{R-02: swaps, tableaux, and positive projectors}

\paragraph{Statement R-02.}
Set $m=an$, $\delta_b=(b-1,b-2,\ldots,0)$, and
$\Delta_b(X)=\det(x_i^{b-j})$. For a tableau $T$ of shape $r$ with
semistandard entries in the ordered alphabet $\{1,\ldots,a\}\times
\{1,\ldots,n\}$, let $p(T)$ and $q(T)$ be the two marginal content vectors.
Write $N_r(p,q)$ for the number of these tableaux with marginal contents
$(p,q)$, and set it equal to zero when a content component is negative.
Then
\begin{align}\label{R02:tableaux}
 g(\lambda,r,r)
 &= [X^{\lambda+\delta_a}Y^{r+\delta_n}]
       \Delta_a(X)\Delta_n(Y)s_r(XY)\\
 &=\sum_{\sigma\in S_a,\tau\in S_n}
    \operatorname{sgn}(\sigma)\operatorname{sgn}(\tau)
    N_r(\lambda+\delta_a-\sigma\delta_a,
        r+\delta_n-\tau\delta_n).\nonumber
\end{align}
The exact unitary sum over swaps acts on the multiplicity by inserting
\begin{equation}\label{R02:swaps}
 s_{(d^n)}(z_1,\ldots,z_m)=
 \sum_{\substack{S\subseteq\{1,\ldots,m\}\\|S|=n}}
 \frac{\prod_{i\in S}z_i^d}
      {\prod_{i\in S,j\notin S}(1-z_j/z_i)}
\end{equation}
in the first line of \eqref{R02:tableaux}, and then setting $z_{ij}=x_i y_j$.
The sum in \eqref{R02:swaps} is combined before specialization or coefficient
extraction; its individual summands have poles and are not separately
nonnegative tableau counts.

There is, nevertheless, an exact positive formula. On the unitary space
\[
 T=S_\lambda(\mathbb C^n\otimes\mathbb C^n)
       \otimes\det(\mathbb C^n)^{-d}\otimes\det(\mathbb C^n)^{-d},
\]
let $P=\int_{U(n)\times U(n)}\rho_T(A,B)\,dA\,dB$. For any orthonormal
basis $(e_i)$ of $T$,
\begin{equation}\label{R02:projector}
 g(\lambda,r,r)=\operatorname{tr}P
   =\|P\|_{\mathrm{HS}}^2=\sum_{i,j}|\langle e_i,Pe_j\rangle|^2.
\end{equation}
In particular, this gives a rigorous sum of squares and a nonvanishing
criterion. It does not give a polynomial-size counting rule: the entries
of $P$ are themselves signed integrals, and the representation dimension
can be exponential in the natural parameters.

\begin{proof}
Multiply \eqref{R01:branching} by the two Vandermonde alternants.
For a partition $\mu$, the alternant $\Delta_a(X)s_\mu(X)$ has exactly one
monomial with strictly decreasing exponent vector equal to
$\lambda+\delta_a$ precisely when $\mu=\lambda$, and its coefficient then
is one. The analogous statement holds for $Y$. This proves the first line
of \eqref{R02:tableaux}. Expanding $s_r$ into semistandard tableaux and
the two alternants into their permutation terms proves the second line.
Formula \eqref{R02:swaps} is the rectangular Weyl-character Laplace expansion
used in Bump--Gamburd, Proposition~4 and \S4. For the last assertion,
Haar invariance gives $P^2=P$ and $P^*=P$; its range is exactly $T^{U(n)\times
U(n)}$. The tensor-subgroup branching identity identifies the dimension of
that range with $g(\lambda,r,r)$. Taking traces of $P^*P=P$ proves
\eqref{R02:projector}.
\end{proof}

\paragraph{What the arithmetic CFKRS recipe supplies.}
CFKRS's unitary formula is a theorem (Theorem~1.5.2 in the supplied source),
and its swap/contour conversion is Lemma~2.5.1. The corresponding
$L$-function mean-value assertion is Conjecture~1.5.1, derived by the
recipe in \S2.1 and the sum over $\Xi$ in \S2.2, equation
\nolinkurl{conj:M}. Replacing the unitary factors in \eqref{R02:swaps} by the
arithmetic Euler product and shifted zeta factors does not define a
character restriction map. Neither RH nor GRH supplies the discarded
off-diagonal estimates in that recipe. The proven action on GCT is
\eqref{R02:tableaux}--\eqref{R02:swaps}; the arithmetic recipe remains a
separate conjectural mean-value statement.

\subsection{R-03: three explicit limits on positivity transfer}

\paragraph{Statement R-03.}
The following three proposed simplifications are false.
First, positivity of a rectangular CUE moment does not imply positivity
of every rectangular Kronecker coefficient. Second, the connected
left-right stabilizer multiplicity need not equal the full determinant
stabilizer multiplicity. Third, one cannot replace tensor-subgroup
branching by one multiplicity-free interlacing step.

\begin{proof}
For $r=(2,2)$ the $S_4$ representation $[r]$ is the two-dimensional
nontrivial summand in the permutation representation on the three
perfect matchings of four letters. Its character at the cycle types
$1^4,21^2,2^2,31,4$ is respectively $2,0,2,-1,0$; the class sizes are
$1,6,3,8,6$. The standard character $[31]$ is $3,1,-1,0,-1$.
Their character inner product gives
\[
 g((3,1),(2,2),(2,2))=(12-12)/24=0,
\]
although $D_{(2,2)}(4)=20>0$. The same elementary calculation yields
\[
 [22]\otimes[22]=[4]\oplus[22]\oplus[1111],\qquad
 \operatorname{Sym}^2[22]=[4]\oplus[22].
\]
Thus $g((1^4),(22),(22))=1$ but the corresponding symmetric Kronecker
coefficient is zero. These values are checked independently by enumerating
the 24 permutations in \nolinkurl{check_small_cases.py}.

The tensor homomorphism $U(a)\times U(n)\longrightarrow U(an)$ has kernel
$\{(zI_a,z^{-1}I_n):|z|=1\}$. Its image has one-dimensional connected
center, whereas the block-diagonal Levi $U(a)\times U(n)$ has
two-dimensional connected center. For $a,n\geq2$ these are different
branching problems. Moreover, R-07 below constructs two independent
highest-weight invariants for $g((12,4),(4^4),(4^4))$, so its value is at
least two (in fact exactly two). A single $0/1$ interlacing rule cannot
be its multiplicity formula.
\end{proof}

The $0/1$ Gan--Gross--Prasad periods cited in Alp\"oge's steps~4 and~7
therefore require an additional, explicit comparison of branching
functors before they can control these multiplicities. Iterating
multiplicity-free restrictions can create multiple paths; no
contradiction with such restrictions is asserted. Likewise the squared
Gram formulas are not statements that every Weingarten summand is
positive: Collins--\'Sniady, Corollary~2.4 and its following example, give
$\operatorname{Wg}_N((12))=-1/(N(N^2-1))$ for $N\geq2$.

\subsection{R-04: plethysm requires the joint law of powers}

\paragraph{Statement R-04.}
For $N\geq\ell(\lambda)$ and $|\lambda|=nd$, the plethysm multiplicity is
\begin{equation}\label{R04:plethysm}
 a_\lambda(d[n])=\int_{U(N)}\overline{s_\lambda(U)}
 \sum_{\kappa\vdash d}\frac{1}{z_\kappa}
       \prod_{j\geq1}s_{(n)}(U^j)^{m_j(\kappa)}\,dU,
 \qquad z_\kappa=\prod_j j^{m_j(\kappa)}m_j(\kappa)!.
\end{equation}
Rains's Theorem~1.3 determines the eigenvalue law of each single $U^j$,
but replacing the different powers in \eqref{R04:plethysm} by independent
matrices with those marginal laws is invalid. A quantitative example is
\begin{equation}\label{R04:counter}
 \int_{U(2)}(\operatorname{tr}U)^4
        \overline{(\operatorname{tr}U^2)^2}\,dU=2,
\end{equation}
whereas the product of the two marginal expectations is zero.

\begin{proof}
For any finite-dimensional $W$, the cycle-index formula for the
symmetrizer $d!^{-1}\sum_{\sigma\in S_d}\sigma$ gives
$\chi_{\operatorname{Sym}^dW}(g)=\sum_{\kappa\vdash d}
z_\kappa^{-1}\prod_j\chi_W(g^j)^{m_j(\kappa)}$.
Set $W=\operatorname{Sym}^n\mathbb C^N$ and apply Schur orthogonality.
For \eqref{R04:counter}, in two variables
\[
 p_1^4=s_{(4)}+3s_{(3,1)}+2s_{(2,2)},\qquad
 p_2^2=s_{(4)}-s_{(3,1)}+2s_{(2,2)}.
\]
Orthogonality gives $1-3+4=2$. Multiplication of $U$ by a scalar phase
shows that both separate expectations vanish. The two polynomial
expansions are also verified by the exact computation certificate.
\end{proof}

In representation-ring language $U\mapsto U^j$ yields the Adams
operation on characters, not a homomorphism of nonabelian groups.
Already $\psi^2(V)=\operatorname{Sym}^2V-\bigwedge^2V$ has a negative
Schur coefficient. A positive alternative retaining the complete
coupling is
\begin{equation}\label{R04:wreath}
 a_\lambda(d[n])=\dim[\lambda]^{S_n\wr S_d}
 =\sum_{i,j}\left|\frac{1}{(n!)^d d!}
       \sum_{h\in S_n\wr S_d}\rho_\lambda(h)_{ij}\right|^2.
\end{equation}
Indeed, apply Schur--Weyl duality to $V^{\otimes nd}$ and the orthogonal
projector onto the wreath-product invariants, exactly as in R-02.
The formula is positive and exact; its factorial enumeration cost is
explicit, and it is not a solution of the efficient positive-formula
problem.

\subsection{R-05: Haar covariance computes the orbit-closure ranks}

\paragraph{Statement R-05.}
Let $V=\mathbb C^N$, $W=\operatorname{Sym}^nV$, $0\ne f\in W$, and
$X_f=\overline{GL(V)f}$. Put $E_d=\operatorname{Sym}^dW$ and
$v=f^{\otimes d}\in E_d$. Define
\begin{equation}\label{R05:twirl}
 M_f=\int_{U(N)}|\rho(k)v\rangle\langle\rho(k)v|\,dk.
\end{equation}
Under the Hermitian identification of $E_d$ with its conjugate dual,
$\ker M_f$ is the degree-$d$ ideal of $X_f$, and the image is dual to
$\mathbb C[X_f]_d$. Write the isotypic decomposition
\[
 E_d=\bigoplus_{\lambda\vdash nd} L_\lambda\otimes\mathbb C^{a_\lambda},
 \qquad L_\lambda=S_\lambda V,
 \qquad v_\lambda=\sum_{i=1}^{D_\lambda(N)}e_i\otimes m_i.
\]
Then the $\lambda$ block of $M_f$ is
\begin{equation}\label{R05:block}
 \frac{I_{L_\lambda}}{D_\lambda(N)}\otimes R_\lambda(f),
 \qquad R_\lambda(f)=\sum_i|m_i\rangle\langle m_i|.
\end{equation}
The multiplicity of $L_\lambda^*$ in $\mathbb C[X_f]_d$ is therefore
\begin{equation}\label{R05:rank}
 m_\lambda(f)=\operatorname{rank}R_\lambda(f).
\end{equation}
If $F_\lambda$ is the matrix with columns $m_i$, then, for every $s$,
\begin{equation}\label{R05:minors}
 e_s(R_\lambda(f))=
 \sum_{\substack{I\subseteq\{1,\ldots,a_\lambda\}\\|I|=s}}
 \ \sum_{\substack{J\subseteq\{1,\ldots,D_\lambda(N)\}\\|J|=s}}
       |\det(F_{\lambda,I,J})|^2.
\end{equation}
In particular $m_\lambda(f)\geq s$ if and only if this nonnegative
polynomial is positive. A multiplicity obstruction between two orbit
closures is thus an exact comparison of these Gram ranks, including
the boundary equations, rather than of their scalar traces.

\begin{proof}
For $u\in E_d$,
$\langle u,M_fu\rangle=\int|\langle u,\rho(k)v\rangle|^2\,dk$.
This is zero exactly when the continuous nonnegative integrand is
identically zero. The compact group $U(N)$ is Zariski dense in
$GL_N(\mathbb C)$: a regular function vanishing on it has all derivatives
along its real Lie algebra zero, and the complexification of that Lie
algebra is $\mathfrak{gl}_N(\mathbb C)$; analytic continuation on the
connected complex group gives zero. Consequently the same vanishing
holds on $GL(V)f$ and its closure. A degree-$d$ polynomial on $W$ is
exactly a linear functional on $E_d$ evaluated at $f^{\otimes d}$.
This identifies the kernel with the homogeneous ideal and the image
with the dual quotient.

Off-diagonal isotypic blocks of the twirl vanish by Schur's lemma.
Within $L_\lambda\otimes\mathbb C^{a_\lambda}$, averaging
$|e_i\rangle\langle e_j|$ gives
$\delta_{ij}I_{L_\lambda}/D_\lambda(N)$, proving
\eqref{R05:block}. Thus the orbit span is
$L_\lambda\otimes\operatorname{span}\{m_i\}$, which proves
\eqref{R05:rank}. Finally $R_\lambda=F_\lambda F_\lambda^*$.
The coefficient $e_s$ is the sum of its principal $s$-minors;
Cauchy--Binet gives \eqref{R05:minors}. Since a positive semidefinite
matrix has nonnegative eigenvalues, $e_s>0$ exactly when at least $s$
of those eigenvalues are positive.
\end{proof}

For $f=\det_n$ or the padded permanent, \eqref{R05:twirl} is a Haar
integral over $U(n^2)$. Its entries have degree $nd$ in the matrix
entries and degree $nd$ in their conjugates. Thus the exact
Collins--\'Sniady formula evaluates them. The full, dimension-truncated
Weingarten inverse must be used when tensor degree exceeds group
dimension. The remaining task is to control the ranks of the resulting
structured matrices; knowing that the matrices are positive
semidefinite does not determine which minors vanish.

\subsection{R-06: a Gaussian and integral form of the rank certificate}

\paragraph{Statement R-06.}
Use the tensor realization $E_d\subseteq V^{\otimes D}$ and set
$v=f^{\otimes d}$. For an $N\times N$ matrix of indeterminates $Z$, write
\begin{equation}\label{R06:coefficients}
 Z^{\otimes D}v=\sum_{|\alpha|=D}Z^\alpha c_\alpha,
 \qquad C_f=(c_\alpha)_\alpha.
\end{equation}
There are $N^D$ tensor coordinates and
$\binom{N^2+D-1}{D}$ columns. If $G$ has independent standard complex
Gaussian entries, then
\begin{equation}\label{R06:gaussian}
 \mathbb E|G^{\otimes D}v\rangle\langle G^{\otimes D}v|
   =\sum_{|\alpha|=D}\alpha!\,c_\alpha c_\alpha^*.
\end{equation}
Its image is the orbit span from R-05. It has the same kernel and
the same isotypic ranks as the Haar covariance.

Let $f^\lambda=\dim[\lambda]$ and let
\[
 e_\lambda=\frac{f^\lambda}{D!}
       \sum_{\sigma\in S_D}\chi^\lambda(\sigma)\sigma
\]
act by permutation of tensor factors. Then
\begin{equation}\label{R06:integer}
 A_\lambda(f)=\sum_{\sigma\in S_D}\chi^\lambda(\sigma)\sigma C_f,
 \qquad
 \operatorname{rank}A_\lambda(f)=D_\lambda(N)m_\lambda(f).
\end{equation}
When $f$ is an integral symmetric tensor, $A_\lambda(f)$ has integer
entries. If its tensor entries have absolute value at most $H$, then
every entry of $A_\lambda(f)$ has absolute value at most
$D!f^\lambda N^D H^d$.
The same rank is obtained by replacing the coefficient columns by all
evaluations at the finite grid $Z\in\{0,1,\ldots,D\}^{N^2}$ and then
applying $e_\lambda$. This grid and the coefficient matrix are
explicit but not polynomial-size constructions.
An ordinary polynomial with integer monomial coefficients may require
scaling by $n!$ under the symmetric-tensor embedding to meet the
integrality hypothesis; this nonzero scalar does not change any orbit
span or multiplicity in characteristic zero.

\begin{proof}
For a standard complex Gaussian $z$, polar integration gives
$\mathbb E[z^p\bar z^q]=0$ if $p\ne q$ and $p!$ if $p=q$.
Independence proves \eqref{R06:gaussian}. Its image is
$\operatorname{span}\{c_\alpha\}$. A linear functional kills all
$c_\alpha$ exactly when its evaluation on $Z^{\otimes D}v$ is the zero
polynomial. Since invertible matrices are Zariski dense, that is
equivalent to killing the $GL_N$ orbit span. This proves the image and
kernel assertions without analytic approximation.

The central idempotent $e_\lambda$ projects precisely onto the
$L_\lambda$ isotypic summand by Schur--Weyl duality. Consequently
$e_\lambda C_f$ has rank $D_\lambda(N)m_\lambda(f)$, and multiplication
by the nonzero scalar $D!/f^\lambda$ gives \eqref{R06:integer}.
Symmetric-group characters and permutation matrices have integral
entries. Each tensor coordinate of $Z^{\otimes D}v$ is a sum of $N^D$
terms whose coefficients have magnitude at most $H^d$; hence
$|(c_\alpha)_i|\leq N^D H^d$. The character bound
$|\chi^\lambda(\sigma)|\leq f^\lambda$ and the $D!$ permutations give
the stated estimate. Finally a polynomial of individual degree at
most $D$ vanishing on a Cartesian product of $D+1$ distinct scalars is
zero, by induction on the number of variables. Applying this fact to
every linear functional proves equality of the grid and coefficient
spans.
\end{proof}

This replaces signed Haar expansions by an explicitly positive
coefficient Gram matrix. Cauchy--Binet gives positive sums of squared
integer minors after clearing the indicated denominators. It also
exhibits the exact cost: the tensor model has $N^D$ rows, the
symmetrizer has $D!$ terms, and the displayed grid has
$(D+1)^{N^2}$ points. A claimed efficient positivity algorithm must
compress these data or exploit their ranks; positivity alone does
not achieve that compression.

\subsection{R-07: explicit positive witnesses that survive the boundary}

\paragraph{Statement R-07.}
Let $n\geq2$, $d\geq0$, and $0\leq q\leq d$. For
$\lambda=(nd-q,q)$, define
\[
 p_{2,\ldots,n}(q)=[t^q]\prod_{j=2}^n(1-t^j)^{-1}.
\]
There are explicitly constructed, linearly independent highest-weight
vectors, indexed by partitions of $q$ into parts $2,\ldots,n$, in the
degree-$d$ coordinate ring of the determinant orbit closure. In fact
\begin{equation}\label{R07:positive}
 m_\lambda(\det_n)=a_\lambda(d[n])=p_{2,\ldots,n}(q),
 \qquad
 g(\lambda,(d^n),(d^n))\geq p_{2,\ldots,n}(q).
\end{equation}
The last inequality is equality by Manivel, Corollary~2(1) in \S3.3; it is a specialization of his
stable rectangular Kronecker Theorem~1, restated as
Ikenmeyer--Panova, Theorem~2.1. The construction below gives concrete witnesses for this
known numerical positivity and, simultaneously, proves their survival
on the determinant orbit closure.

More generally, with no restriction $q\leq d$,
\begin{equation}\label{R07:alltwo}
 m_\lambda(\det_n)=a_\lambda(d[n])
 \quad\text{for every }\lambda\vdash nd\text{ with }\ell(\lambda)\leq2.
\end{equation}
Therefore no multiplicity obstruction of length at most two can
separate a padded-permanent orbit closure from the determinant orbit
closure in the same ambient space and degree.

\begin{proof}
Write a binary degree-$n$ form as
$F(s,t)=\sum_{j=0}^n\binom nj b_j s^{n-j}t^j$.
In the coordinate representation $\operatorname{Sym}^n\mathbb C^2$,
the upper-triangular raising operator is
$E b_j=j b_{j-1}$, with $E b_0=0$, and the weight of $b_j$ is
$(n-j,j)$. For $2\leq j\leq n$ put
\begin{equation}\label{R07:central}
 Q_j=b_0^{j-1}\sum_{i=0}^j\binom ji b_i
                  \left(-\frac{b_1}{b_0}\right)^{j-i}.
\end{equation}
This is a polynomial: for $i=0$ the factor $b_i=b_0$ cancels the only
possible negative exponent; all $i\geq1$ terms already have
nonnegative powers of $b_0$. The first two examples are
\[
 Q_2=b_0b_2-b_1^2,\qquad
 Q_3=b_0^2b_3-3b_0b_1b_2+2b_1^3.
\]
Under $\exp(uE)$, $b_i$ goes to
$\sum_{h=0}^i\binom ih u^{i-h}b_h$ and $b_1/b_0$ goes to
$b_1/b_0+u$. The binomial theorem in \eqref{R07:central} shows that
the two translations cancel. Thus $E Q_j=0$. Every term has degree
$j$ in the $b_i$ and weight $((n-1)j,j)$.

For a tuple of nonnegative integers $(a_2,\ldots,a_n)$ satisfying
$\sum_{j=2}^n ja_j=q$, define
\begin{equation}\label{R07:witness}
 P_{\mathbf a}=b_0^{d-q}\prod_{j=2}^n Q_j^{a_j}.
\end{equation}
This is a highest-weight polynomial of degree $d$ and weight
$(nd-q,q)$. Evaluate it on determinant pencils
$F(s,t)=\det(sI_n+tY)$, with $Y$ diagonal. Every monic binary form
is obtained in this way over $\mathbb C$ by factoring the corresponding
one-variable polynomial. In particular, one may set $b_0=1,b_1=0$
and choose $b_2,\ldots,b_n$ arbitrarily. On this family $Q_j=b_j$,
so the restrictions of the $P_{\mathbf a}$ are distinct monomials
in independent variables. They are linearly independent.

For completeness, the ambient highest-weight dimension in this range
is obtained by a short weight count. In
$\operatorname{Sym}^d(\operatorname{Sym}^n\mathbb C^2)$ the weight
$(nd-q,q)$ has dimension equal to the number of partitions of $q$
into at most $d$ parts, each at most $n$. When $q\leq d$ this is
$[t^q]\prod_{j=1}^n(1-t^j)^{-1}$. A $GL_2$ irreducible of highest
weight $(nd-r,r)$ has one vector at each weight $(nd-q,q)$ with
$r\leq q\leq nd-r$. Subtracting consecutive weight dimensions
therefore gives $p_{2,\ldots,n}(q)$. This proves that the displayed
independent witnesses form a basis and establishes the first equality
in \eqref{R07:positive} once boundary survival is checked.

Now let $V$ have dimension $n^2$ and view a binary form as an element
of $\operatorname{Sym}^nV^*$ supported on the first two coordinates.
It factors as $\prod_{i=1}^n\ell_i$ and hence equals
$\det(\operatorname{diag}(\ell_1,\ldots,\ell_n))$.
This is a linear, possibly singular, substitution into $\det_n$.
Since $GL(V)$ is dense in $\operatorname{End}(V)$ and substitution
is polynomial, every such form lies in the determinant orbit closure.
A vector of highest weight $\lambda=(\lambda_1,\lambda_2,0,\ldots,0)$
in $\operatorname{Sym}^d(\operatorname{Sym}^nV)$ uses only the first
two basis vectors: each other weight coordinate is a sum of
nonnegative exponents and equals zero. The additional raising
operators annihilate it, so its highest-weight condition is exactly
the binary one. A nonzero polynomial in these binary coordinates
cannot vanish on every binary form. Thus the homogeneous ideal of
the determinant orbit closure has no nonzero highest-weight vector
of length at most two. Taking the quotient proves
\eqref{R07:alltwo}, including all degrees and all two-row weights.

The determinant-coordinate multiplicity is at most its orbit
multiplicity, which is the symmetric rectangular Kronecker coefficient,
and that is at most $g(\lambda,r,r)$; this is the cited stabilizer
statement in B\"urgisser--Ikenmeyer--Panova, \S1(b).
This proves the inequality in \eqref{R07:positive}. Finally any
padded-permanent coordinate ring is a quotient of the same ambient
$\operatorname{Sym}(\operatorname{Sym}^nV)$, so its multiplicity is
at most $a_\lambda(d[n])=m_\lambda(\det_n)$ for two-row $\lambda$.
\end{proof}

For $n=d=q=4$, the two witnesses are $Q_4$ and $Q_2^2$, with
\[
 Q_4=b_0^3b_4-4b_0^2b_1b_3+6b_0b_1^2b_2-3b_1^4.
\]
They specialize to $b_4$ and $b_2^2$, proving directly that the
determinant coordinate multiplicity at $(12,4)$ is two. This is
stronger information than nonnegativity of a Haar integral.

\subsection{R-08: the exact Molien numerator needed for orbit closures}

\paragraph{Statement R-08.}
Let $S=\operatorname{Sym}(W)$ be a polynomial coordinate algebra,
$A=S/I$ a graded $GL_N$-equivariant quotient, and $K=U(N)$. Then
\begin{equation}\label{R08:molien}
 \sum_{d\geq0}m_\lambda(A_d)t^d
   =\int_K\overline{\chi_\lambda(k)}H_A(t,k)\,dk.
\end{equation}
For the ambient algebra,
$H_S(t,k)=\det(1-t\rho_W(k))^{-1}$.
For a finite graded equivariant free resolution
$\cdots\to S\otimes T_1\to S\otimes T_0\to A\to0$,
\begin{equation}\label{R08:numerator}
 H_A(t,k)=
 \frac{\sum_i(-1)^i\chi_{T_i}(t,k)}
      {\det(1-t\rho_W(k))}.
\end{equation}
Thus the Molien denominator and ratios identities compute the desired
multiplicities only after the orbit ideal's equivariant numerator,
or equivalently the ranks in R-05--R-06, is supplied.

\begin{proof}
Coefficientwise Schur orthogonality proves \eqref{R08:molien}.
The eigenvalue expansion of the symmetric algebra gives the
denominator formula. Alternating graded characters are additive
on exact complexes, yielding \eqref{R08:numerator}. For a purely
formal interpretation no convergence assumption is needed; for
analytic evaluation take $|t|<1$ in the unitary polynomial case.
\end{proof}

\subsection{R-09: an exact ratios kernel for squared GCT multiplicities}

\paragraph{Statement R-09.}
Let $n,d,N\geq1$, $r=(d^n)$, and define
\[
 S_N(n,d)=\sum_{\substack{\lambda\vdash nd\\\ell(\lambda)\leq N}}
                    g(\lambda,r,r)^2.
\]
Set $S_0(n,d)=0$.
For $0<t<1$ and $A,B,C,D\in U(n)$ put
\begin{align*}
 \mathcal R_N(t;A,B,C,D)=\int_{U(N)}
 &\det\bigl(I-t(A\otimes B)\otimes U\bigr)^{-1}\\[-2mm]
 &\quad\cdot
 \det\bigl(I-t(C\otimes D)\otimes U^{-1}\bigr)^{-1}\,dU.
\end{align*}
Then the explicitly specified fourfold Fourier projection is
\begin{equation}\label{R09:squares}
 \int_{U(n)^4}\mathcal R_N(t;A,B,C,D)
   (\det A\det B\det C\det D)^{-d}\,dA\,dB\,dC\,dD
       =t^{2nd}S_N(n,d).
\end{equation}
For $N\geq n^2$, the inner ratio is evaluated exactly by
Bump--Gamburd, Theorem~3, at $K=L=0,Q=R=n^2$:
\begin{equation}\label{R09:evaluated}
 \mathcal R_N(t;A,B,C,D)
       =\det\bigl(I_{n^4}-t^2 A\otimes B\otimes C\otimes D\bigr)^{-1}.
\end{equation}
Consequently an explicitly evaluated Molien--Weyl coefficient is the
sum of squares of precisely the rectangular Kronecker coefficients
needed in GCT. The truncated sums satisfy
\begin{equation}\label{R09:monotone}
 S_N(n,d)-S_{N-1}(n,d)
    =\sum_{\substack{\lambda\vdash nd\\\ell(\lambda)=N}}g(\lambda,r,r)^2
    \geq0,
 \qquad
 g(\lambda,r,r)\leq\lfloor\sqrt{S_N(n,d)}\rfloor
 \quad(\ell(\lambda)\leq N).
\end{equation}
The stabilized value has a second nonnegative moment expression:
\begin{equation}\label{R09:finitefourth}
 S_{n^2}(n,d)=\frac1{(nd)!}\sum_{\sigma\in S_{nd}}\chi^r(\sigma)^4.
\end{equation}

\begin{proof}
The Cauchy identity and absolute convergence for $t<1$ expand the two
inverse determinants into
$\sum_\lambda t^{|\lambda|}s_\lambda(A\otimes B)s_\lambda(U)$
and
$\sum_\mu t^{|\mu|}s_\mu(C\otimes D)s_\mu(U^{-1})$.
Schur orthogonality on $U(N)$ retains exactly $\lambda=\mu$ of
length at most $N$. The two applications of the tensor-subgroup
projection in R-02 then retain $|\lambda|=nd$ and give
$g(\lambda,r,r)$ from each pair of $U(n)$ factors. This proves
\eqref{R09:squares}. When $N\geq n^2$, the length cutoff is
inactive because both product alphabets have $n^2$ entries. A second
Cauchy identity gives \eqref{R09:evaluated}, equivalently the stated
specialization of the ratios theorem. The definition immediately
proves \eqref{R09:monotone}. Finally the character of
$[r]\otimes[r]$ is $(\chi^r)^2$. Its squared character norm is
the sum of squares of its irreducible multiplicities, proving
\eqref{R09:finitefourth}; symmetric-group characters are real.
The tensor-subgroup branching formula also gives
$g(\lambda,r,r)=0$ for $\ell(\lambda)>n^2$, because
$S_\lambda(\mathbb C^n\otimes\mathbb C^n)=0$ in that range.
Thus the cutoff in the left side of \eqref{R09:finitefourth}
omits no nonzero multiplicity.
\end{proof}

The positive quantities in \eqref{R09:squares} are obtained after
projection. The pointwise ratio and individual swap summands need
not be nonnegative. An approximation argument would need a
certified error below one half to recover an integer coefficient,
with bounds uniform in the growing tensor degree. Neither the
closed ratio formula nor RH supplies those estimates or compresses
the remaining fourfold Fourier projection.

\subsection{R-10: the complete height-two rectangle and its boundary gate}

\paragraph{Statement R-10 (known formula, specialized and compared).}
Let $d\geq0$ and $\lambda\vdash2d$. For partitions of length at
most four, append zero parts to obtain $(\lambda_1,\ldots,\lambda_4)$.
Then
\begin{equation}\label{R10:rectangle}
g(\lambda,(d,d),(d,d))=
\begin{cases}
1,&\ell(\lambda)\leq4\text{ and all four parts have the same parity},\\
0,&\text{otherwise}.
\end{cases}
\end{equation}
More precisely, the multiplicity of $[\lambda]$ in
$\operatorname{Sym}^2[d,d]$ is one exactly for partitions with all
parts even and length at most four. Its multiplicity in
$\bigwedge^2[d,d]$ is one exactly for partitions with exactly four
odd parts. These formulas are the constant-term specialization of
Blasiak--Mulmuley--Sohoni, GCT~IV, Corollary~18.4
(source label \nolinkurl{c kronecker generating function sym ext});
they are cited known results, not new numerical formulas.

In the determinant problem with $n=2$, $N=4$, the boundary side is
also completely determined:
\begin{equation}\label{R10:closure}
m_\lambda(\det_2)=a_\lambda(d[2])=
\begin{cases}
1,&\ell(\lambda)\leq4\text{ and every part of }\lambda\text{ is even},\\
0,&\text{otherwise}.
\end{cases}
\end{equation}
Thus all four-odd-part summands of the ordinary rectangular
coefficient are absent from the determinant coordinate ring.
For $d=2$, this includes the $\lambda=(1,1,1,1)$ example in R-03.
Simultaneous conjugation of the two inner partitions also extends
\eqref{R10:rectangle}, as an ordinary Kronecker formula, to rectangles
of width two.

\begin{proof}
Write $c_i=\lambda_i-\lambda_{i+1}$ for $1\leq i\leq3$ and
$c_4=\lambda_4$, so $c_i$ is the number of columns of height $i$.
GCT~IV, Corollary~18.4, gives the generating function
$\sum_{\alpha\vdash2d,\ell(\alpha)\leq2}
g_{\varepsilon,\alpha,\lambda}x^{\alpha_1-\alpha_2}$.
Three of its four cases have a factor $x$ or $x^2$ multiplying
polynomials with nonnegative exponents, and hence contribute no
constant term. The remaining case is the product
$\llbracket c_1\rrbracket\llbracket c_2\rrbracket
\llbracket c_3\rrbracket$, where
$\llbracket b\rrbracket=x^b+x^{b-2}+\cdots$ ends at $1$ or $x$.
Its constant coefficient is one exactly when $c_1,c_2,c_3$ are all
even; the case condition then requires
$\varepsilon=(-1)^{c_3+c_4}=(-1)^{\lambda_4}$.
Thus the symmetric square retains the all-even partitions and the
exterior square retains the four-odd-part partitions. Their sum is
\eqref{R10:rectangle}. Length greater than four gives zero by
tensor-subgroup branching. Finally
$[\rho']=[\rho]\otimes\operatorname{sgn}$, so conjugating both
inner partitions leaves their tensor product unchanged.

The determinant $x_{11}x_{22}-x_{12}x_{21}$ is a nondegenerate
quadratic form in four variables. Over $\mathbb C$ all nondegenerate
quadratic forms of a fixed dimension are congruent. They form a
nonempty Zariski-open subset of the space of quadrics, so
$\operatorname{Det}_2=\operatorname{Sym}^2(\mathbb C^4)^*$.
Consequently its coordinate ring is
$\operatorname{Sym}(\operatorname{Sym}^2\mathbb C^4)$.

Here is a direct highest-weight proof of its familiar multiplicity
formula. In $\operatorname{Sym}(\operatorname{Sym}^2\mathbb C^s)$
write the generators as the entries of a symmetric matrix $Q$,
with $q_{ij}$ of weight $e_i+e_j$. The upper-triangular action on
these generators is $Q\mapsto g^tQg$.
Let $\Delta_i$ be the leading principal $i\times i$ minor and
$\Delta_0=1$. Each $\Delta_i$ is fixed by the upper unitriangular
group and has weight $2(e_1+\cdots+e_i)$.
On the dense open set where all $\Delta_i$ are nonzero, elementary
upper-triangular congruences reduce $Q$ to the unique diagonal
matrix with entries $\Delta_i/\Delta_{i-1}$.
Therefore a nonzero unitriangular-invariant polynomial has nonzero
restriction to diagonal matrices. If it is also a torus eigenvector
of weight $\lambda$, this restriction is a scalar multiple of
$\prod_i q_{ii}^{\lambda_i/2}$. It exists only when each
$\lambda_i$ is even, and its value determines the polynomial on a
dense open set. Hence each highest-weight space has dimension at
most one and vanishes unless $\lambda=2\mu$.
Conversely, for every partition $\mu$ of length at most $s$,
\[
 \prod_{i=1}^s\Delta_i^{\mu_i-\mu_{i+1}},\qquad\mu_{s+1}=0,
\]
is a nonzero highest-weight polynomial of weight $2\mu$ and
degree $|\mu|$. This proves
$\operatorname{Sym}^d(\operatorname{Sym}^2\mathbb C^s)
=\bigoplus_{\mu\vdash d,\ell(\mu)\leq s}S_{2\mu}\mathbb C^s$
and, at $s=4$, \eqref{R10:closure}.
\end{proof}

The stabilizer interpretation agrees with these formulas.
The full linear stabilizer of the nondegenerate quadratic form is
$O_4$, whereas the connected tensor image of $SL_2\times SL_2$ is
$SO_4$: its kernel is $\{(I,I),(-I,-I)\}$, so its dimension is six,
the dimension of $SO_4$. The extra orthogonal component removes
the odd sector. Thus the connected-stabilizer upper bound, the
symmetric rectangular coefficient, and the closure multiplicity
are completely distinguished in this case.

\subsection{R-11: small-degree padding excludes all four-row obstructions}

\paragraph{Statement R-11.}
Work over $\mathbb C$. Let $1\leq m\leq3$, $n\geq m$, $N=n^2$,
$V=\mathbb C^N$, $f\in\operatorname{Sym}^mV^*$, and $\ell\in V^*$.
Set $h=\ell^{n-m}f$. For every coordinate-ring degree $d\geq0$
and every $\lambda\vdash nd$ with $\ell(\lambda)\leq4$,
\begin{equation}\label{R11:fourrow}
 m_\lambda(h)\leq m_\lambda(\det_n).
\end{equation}
The assertion includes padded permanents of degree at most three.
It is a deduction from classical linear determinantal
representations, not a claim that arbitrary $f$ itself has an
$n\times n$ determinantal representation.

\begin{proof}
First consider forms in at most four variables. Linear forms are
$1\times1$ determinants. Every quadric in at most four variables
is congruent over $\mathbb C$ to $y_1^2+\cdots+y_r^2$, $r\leq4$.
With unused $y_i$ set to zero it is the determinant
\begin{equation}\label{R11:quadric}
 \det\begin{pmatrix}
 y_1+\mathrm i y_2&y_3+\mathrm i y_4\\
 -y_3+\mathrm i y_4&y_1-\mathrm i y_2
 \end{pmatrix}=
 y_1^2+y_2^2+y_3^2+y_4^2.
\end{equation}
Beauville (2000), Corollary~6.4, states over an algebraically
closed field that every smooth cubic surface has a $3\times3$
linear determinantal equation. Thus over $\mathbb C$ every smooth
quaternary cubic is a $3\times3$ linear determinant, after
absorbing its nonzero scalar factor in one matrix row.
Smooth cubics are Zariski dense in the affine space of quaternary
cubics: the discriminant-complement is open and nonempty, as the
Fermat cubic shows. Therefore every quaternary cubic lies in the
closure of the linear-determinant image.

If $q=\det M$ has degree $m$ and $u$ is linear, then
$u^{n-m}q=\det\operatorname{diag}(M,uI_{n-m})$.
The entries of this block matrix give a linear substitution into
$\det_n$. Such substitutions belong to $\operatorname{Det}_n$,
because invertible linear maps are dense in all linear maps and
substitution is polynomial. For fixed $u$, the preimage of
$\operatorname{Det}_n$ under the linear map $q\mapsto u^{n-m}q$
is closed. The density just proved extends the assertion from
smooth cubics to every cubic. We have shown that
\begin{equation}\label{R11:projectedcontainment}
 u^{n-m}q\in\operatorname{Det}_n.
\end{equation}
Here $q$ has degree $m\leq3$, and $q$ and $u$ are supported
on the same space of at most four variables.

Put $k=\min(4,N)$ and let $\pi$ set the last $N-k$ variables
equal to zero. Let $P$ be a highest-weight vector of weight
$\lambda$ in the degree-$d$ ideal of $\operatorname{Det}_n$.
Since $\lambda_j=0$ for $j>k$ and all coordinate generators have
nonnegative weights, every monomial of $P$ uses only coefficients
of forms supported on the first $k$ variables. Hence
$P(F)=P(\pi F)$ for every $F\in\operatorname{Sym}^nV^*$.
For $g\in GL(V)$ the projection is multiplicative on forms, so
\[
 \pi(g\cdot h)=(\pi(g\cdot\ell))^{n-m}\pi(g\cdot f)
 \in\operatorname{Det}_n
\]
by \eqref{R11:projectedcontainment}. It follows that
$P(g\cdot h)=0$ for every $g$, and hence $P$ belongs to the
ideal of $\overline{GL(V)\cdot h}$.
Thus the highest-weight space of the determinant ideal is
contained in that of the padded-form ideal for every such
$\lambda$. Finite-dimensional polynomial $GL(V)$-modules are
semisimple in characteristic zero. Subtracting the ideal
multiplicity from the common ambient multiplicity gives
\eqref{R11:fourrow}. If $N<4$, partitions of length greater than
$N$ already have multiplicity zero, so the same statement holds.
\end{proof}

\subsection{R-12: exact permanent-pencil dominance and four-row deficits}

\paragraph{Statement R-12 (computer-assisted finite-degree deduction).}
For $n=4,5$, the polynomial map
\[
 \Phi_n:(\operatorname{Mat}_n\mathbb C)^4
 \longrightarrow\operatorname{Sym}^n(\mathbb C^4)^*,
 \qquad(A_1,A_2,A_3,A_4)\longmapsto
 \operatorname{per}\!\left(\sum_{k=1}^4x_kA_k\right)
\]
is dominant. Consequently, for every $d\geq0$ and
$\lambda\vdash nd$ with $\ell(\lambda)\leq4$,
\begin{equation}\label{R12:perambient}
 m_\lambda(\operatorname{per}_n)=a_\lambda(d[n])
 \qquad(n=4,5).
\end{equation}
Combining the quartic case with G-07 and writing $D=320112$ gives
the exact comparison
\begin{equation}\label{R12:deficit}
 m_\lambda(\operatorname{per}_4)-m_\lambda(\det_4)
 =a_{\lambda-(D,D,D,D)}((d-D)[4]),
\end{equation}
with the same zero conventions as in G-07.
In particular, for every integer $s\geq0$,
\begin{equation}\label{R12:gapone}
 d=D+s,\qquad \lambda=(D+4s,D,D,D)
 \quad\Longrightarrow\quad
 m_\lambda(\operatorname{per}_4)-m_\lambda(\det_4)=1.
\end{equation}
These are multiplicity obstructions; no occurrence-vanishing
claim is implied. For $n=5$, G-06 and \eqref{R12:perambient}
give a four-row multiplicity obstruction in some degree, without
specifying that degree or highest weight.

\paragraph{Explicit rank certificate for $n=4$.}
Each entry of the following array is the coefficient quadruple of
a linear form in the ordered variables $(x_1,x_2,x_3,x_4)$:
\begin{equation}\label{R12:pencil4}
\begin{pmatrix}
(2,-1,0,0)&(2,-2,0,-2)&(-2,-1,-1,0)&(0,-1,-1,2)\\
(0,1,0,0)&(2,-2,-1,1)&(-2,1,-2,0)&(-1,-2,-1,-1)\\
(0,-2,0,2)&(1,-2,-1,-1)&(2,-1,-1,0)&(2,2,0,0)\\
(1,1,2,1)&(1,0,2,-2)&(-1,2,-1,1)&(2,1,1,1)
\end{pmatrix}.
\end{equation}
Call the resulting matrix of linear forms $L_4(x)$. Index the
parameters $a_{ijk}$ by $0\leq i,j\leq3$, $0\leq k\leq3$,
where $a_{ijk}$ multiplies $x_{k+1}$ in entry $(i+1,j+1)$,
and assign column number $c=4(4i+j)+k$.
Order the 35 exponent quadruples $\alpha$ of total degree four in
decreasing lexicographic order. Define the integer Jacobian by
\begin{equation}\label{R12:derivative}
 J_{\alpha,4(4i+j)+k}
 =[x^\alpha]\;x_{k+1}
 \operatorname{per}\bigl(L_4(x)_{\widehat{i+1},\widehat{j+1}}\bigr).
\end{equation}
Take all 35 rows and the ordered columns
\[
 C_4=(0,1,\ldots,30,32,33,34,35).
\]
The resulting square integer matrix has determinant
\begin{equation}\label{R12:minor4}
 \det J_{C_4}
 =45459683955923995466096193465479594741071872
 \equiv77\pmod{101}.
\end{equation}
Thus \eqref{R12:pencil4}--\eqref{R12:minor4} specify the entire
nonvanishing certificate in the text, independently of a file format.

\paragraph{Explicit rank certificate for $n=5$.}
Use the same coefficient-quadruple convention with the array
\begin{equation}\label{R12:pencil5}
\begin{pmatrix}
(2,2,-1,1)&(2,-2,-2,1)&(-2,-1,1,1)&(1,-1,0,0)&(1,2,0,2)\\
(1,0,0,-1)&(1,2,0,2)&(1,2,-1,2)&(0,2,-1,1)&(-1,2,1,0)\\
(-2,-2,1,0)&(1,-1,-1,2)&(0,1,0,-1)&(2,2,-1,1)&(2,-1,1,-2)\\
(2,0,0,1)&(-1,1,2,2)&(0,0,1,-2)&(1,-1,1,-1)&(-1,0,1,-2)\\
(0,-2,1,-1)&(1,0,-1,1)&(-1,-1,1,-1)&(-1,-1,2,1)&(0,-2,-2,2)
\end{pmatrix}.
\end{equation}
There are 56 degree-five coefficient rows, again in decreasing
lexicographic order, and parameter column $4(5i+j)+k$.
Use the same derivative rule, with $4\times4$ complementary
permanents, and the columns
\[
 C_5=(0,1,\ldots,37,39,40,\ldots,56).
\]
The selected $56\times56$ minor has determinant $30$ modulo 101.

\begin{proof}
Differentiating the permanent with respect to one matrix entry
gives the complementary permanent; differentiating that entry
with respect to $a_{ijk}$ contributes $x_{k+1}$. This proves the
Jacobian rule. Expanding the complementary permanents and
performing exact modular elimination gives the displayed
nonzero residues. For full reproducibility,
\nolinkurl{check_permanent_projection.py} records the full pencils,
coefficient matrices, selected columns, and integer minors.
The characteristic-zero determinants are also calculated by
fraction-free Bareiss elimination. As a second check on the
derivative expansion, every one of the 64 quartic columns was
recomputed as the exact finite difference
$\Phi_4(a+e_{ijk})-\Phi_4(a)$; this is exactly the derivative
because the map is affine in each individual parameter.
Twelve quintic columns were checked in the same way.

A nonzero minor modulo 101 is a nonzero integer minor, so the
complex differential is surjective at each displayed point.
Fix the other parameters and retain the selected 35 or 56
parameters. The resulting square polynomial map has invertible
complex differential there. By the inverse function theorem its
image contains a nonempty analytic open subset. A polynomial
vanishing on that image vanishes on an analytic open subset and
is identically zero. This proves Zariski density, hence dominance.

Every permanent pencil is a linear endomorphism substitution
into $\operatorname{per}_n$, and therefore belongs to its
$GL_{n^2}$ orbit closure. This closed variety consequently
contains the full linear space of degree-$n$ forms supported on
the first four variables. A highest-weight polynomial of weight
$\lambda$ with $\ell(\lambda)\leq4$ uses only those coefficient
variables, by the nonnegative-weight support argument of R-07.
If it vanishes on the permanent orbit closure, it vanishes on
that entire coefficient space, so it is the zero polynomial.
Thus the ideal has no highest-weight vector in those weights,
which proves \eqref{R12:perambient}.

Equation \eqref{R12:deficit} now follows by subtracting G-07.
For \eqref{R12:gapone}, the shifted partition is $(4s)$, whose
multiplicity in $\operatorname{Sym}^s(\operatorname{Sym}^4V)$ is
one: its weight space is spanned by $(e_1^4)^s$, and this vector
is highest weight. This also covers $s=0$.
For $n=5$, the four-variable determinant image has dimension at
most $52<56$ by G-06, so its ideal contains a nonzero homogeneous
highest-weight vector. G-05 excludes highest weights of length
at most three. Hence it removes at least one four-row ambient
multiplicity, whereas \eqref{R12:perambient} retains them all.
\end{proof}

These computations are exact finite certificates. The independent
agent recheck uses Ryser's formula at integer evaluation points
and an invertible interpolation matrix, so it does not reuse
the complementary-permanent coefficient-expansion algorithm.
The checks establish the stated degrees four and five; they
do not extrapolate dominance to growing $n$ or supply an
asymptotic permanent-versus-determinant lower bound.

\subsection{R-13: the exact range of quaternary permanent dominance}

\paragraph{Statement R-13 (finite certificates and a dimension obstruction).}
For the polynomial map $\Phi_n$ in R-12,
\begin{equation}\label{R13:range}
 \Phi_n\text{ is dominant over }\mathbb C
 \quad\Longleftrightarrow\quad1\leq n\leq16.
\end{equation}
Consequently, for $2\leq n\leq16$ and every
$\lambda\vdash nd$ with $\ell(\lambda)\leq4$,
\begin{equation}\label{R13:multiplicities}
 m_\lambda(\operatorname{per}_n)=a_\lambda(d[n]).
\end{equation}
For every $4\leq n\leq16$, the unpadded permanent and determinant
of size $n$ have a multiplicity obstruction with exactly four rows
in some coordinate-ring degree. R-12 gives an explicit infinite
family when $n=4$; no degree or partition is asserted here for
$5\leq n\leq16$. This is a certificate-based finite-range result,
with no claim of literature priority or asymptotic separation.

\paragraph{Complete specification of the finite certificates.}
All arithmetic for this paragraph is in $\mathbb F_{101}$, with
integer coefficient representatives in $\{0,\ldots,100\}$.
For each $n=1,\ldots,16$ set
\[
\begin{gathered}
 s_0=20261121+1009n,\qquad
 s_{c+1}=(1664525s_c+1013904223)\bmod2^{32},\\
 a_c=\left\lfloor\frac{s_{c+1}}{2^{16}}\right\rfloor\bmod101.
\end{gathered}
\]
There are $4n^2$ coefficients, with $0\leq c<4n^2$, and the
pencil is
\[
 L_n(x)_{i+1,j+1}=\sum_{k=0}^3a_{4(ni+j)+k}x_{k+1}
 \quad(0\leq i,j<n).
\]
Put $b_n=\binom{n+3}{3}$ and order all integer points
$\alpha\in\mathbb N^4$ with $|\alpha|=n$ in decreasing
lexicographic order. Define the $b_n\times4n^2$ matrix
\begin{equation}\label{R13:evaluationmatrix}
 E_n\bigl(\alpha,4(ni+j)+k\bigr)
 =\alpha_{k+1}\operatorname{per}
   \bigl(L_n(\alpha)_{\widehat{i+1},\widehat{j+1}}\bigr)
 \pmod{101}.
\end{equation}
Choose columns greedily from left to right, retaining a column
exactly when it increases the span of previously retained columns.
Equivalently, scan columns left to right in Gaussian elimination,
choose the first available nonzero pivot row, and clear the rows
below it. This produces $b_n$ columns $C_n$ in every case, and the
square minors, with rows and columns in the stated order, have
the following nonzero residues:
\[
\begin{array}{c|r|r@{\qquad}c|r|r}
n&b_n&\det(E_n)_{C_n}&n&b_n&\det(E_n)_{C_n}\\\hline
1&4&1&9&220&67\\
2&10&11&10&286&67\\
3&20&64&11&364&42\\
4&35&49&12&455&87\\
5&56&85&13&560&25\\
6&84&28&14&680&69\\
7&120&25&15&816&73\\
8&165&59&16&969&8
\end{array}
\]
These uniform range certificates use the displayed generator;
R-12 retains its separate, smaller-entry quartic and quintic pencils.
The directory \nolinkurl{permanent_range} contains all generated
integer coefficients, evaluation points, chosen columns, and full
matrices, with unambiguous row-major byte encoding and checksums.
Thus the formula above and the code reproduce the certificates
without probabilistic choices or an external data source.

\begin{proof}
For the positive direction, \eqref{R13:evaluationmatrix} is the
derivative of the coefficient map followed by evaluation at the
specified points, by the complementary-permanent derivative
rule in R-12. Its full row rank $b_n$ forces the coefficient-map
differential, whose target dimension is $b_n$, to be surjective.
The nonzero minor modulo 101 is the reduction of a nonzero
integer minor. Hence the characteristic-zero differential has
the same required full rank, and the inverse-function argument
in R-12 proves dominance.

Here are the exact algorithms used to certify every entry and
minor in the table. For a numeric $n\times n$ matrix $M$, let
$F(\varnothing)=1$ and
\[
 F(S)=\sum_{j\in S}M_{|S|,j}F(S\setminus\{j\}).
\]
Induction on $|S|$ shows that $F(S)$ is the permanent formed from
the first $|S|$ rows and columns $S$. Thus $F([n])$ is the
permanent. Reverse differentiation of this recurrence computes
every entry derivative, hence every complementary permanent,
in $O(n2^n)$ field operations per evaluation point. This is
implemented in \nolinkurl{check_permanent_range.cpp}.

All $3{,}348{,}864$ entries across the sixteen matrices were then
recomputed from Ryser's different identity
\[
 \operatorname{per}(M)=
 \sum_{S\subseteq[n]}(-1)^{n-|S|}
 \prod_{i=1}^n\left(\sum_{j\in S}M_{ij}\right).
\]
To prove this identity, expand the products as maps
$f:[n]\to[n]$. The coefficient of a fixed map is
$\sum_{S\supseteq\operatorname{im}f}(-1)^{n-|S|}$, which is zero
unless $f$ is onto and is one if it is onto. Such onto maps are
exactly permutations. Differentiating this identity and using
Gray-code row-sum updates gives the independent algorithm in
\nolinkurl{check_permanent_range_ryser.cpp}. Every entry agrees
with the first computation. Finally
\nolinkurl{check_permanent_range_minors.py} recomputes each
selected determinant by a second modular elimination algorithm;
all sixteen nonzero residues agree. These are finite exact
checks, rather than floating-point rank estimates.

For the converse, the torus
\[
 H_n=\{(u_1,\ldots,u_n,v_1,\ldots,v_n)\in(\mathbb C^*)^{2n}:
                  \prod_i u_i\prod_j v_j=1\}
\]
preserves $\Phi_n$ under the entrywise action
$L_{ij}\mapsto u_i v_jL_{ij}$. It has dimension $2n-1$.
On the nonempty open set where every entry $L_{ij}$ is a nonzero
linear form, a stabilizing torus element satisfies
$u_i v_j=1$ for all $i,j$. It is therefore of the form
$u_i=t$, $v_j=t^{-1}$, a one-dimensional stabilizer.
Thus each generic fiber contains an orbit of dimension $2n-2$,
and
\begin{equation}\label{R13:dimension}
 \dim\overline{\operatorname{im}\Phi_n}
       \leq4n^2-2n+2.
\end{equation}
At $n=17$, this is $1124<1140=b_{17}$. For $n\geq18$,
\[
 6\bigl(b_n-(4n^2-2n+2)\bigr)
 =n^3-18n^2+23n-6>0.
\]
Consequently no $n\geq17$ map is dominant. This proves
\eqref{R13:range}.

For $n\geq2$, the ambient variable space has dimension $n^2\geq4$,
so dominance places the entire space of degree-$n$ forms in the
first four variables inside the permanent orbit closure. The
highest-weight support argument of R-12 proves
\eqref{R13:multiplicities}.
For $n\geq4$, G-06 gives a nonzero determinant equation supported
on four variables, and G-05 excludes weights with at most three
rows. Thus some exactly four-row highest weight has determinant
multiplicity below its ambient value. Equation
\eqref{R13:multiplicities} retains that ambient value on the
permanent side when $n\leq16$, giving the asserted obstruction.
\end{proof}

The dimension ceiling concerns full ambient survival in four
variables. It does not determine the comparison between the two
proper four-variable images for $n\geq17$.

\subsection{R-14: a dimension criterion that survives genuine padding}

\paragraph{Statement R-14.}
Let $\mathcal P_{m,4}$ be the Zariski closure of the image of
$\Phi_m$ and put $c_m=\dim\mathcal P_{m,4}$.
For integers $n>m\geq2$, embed the usual padded permanent
$h_{m,n}=z^{n-m}\operatorname{per}_m$ in
$\operatorname{Sym}^n(\mathbb C^{n^2})^*$, with $z$ independent
of the $m^2$ permanent variables. The projection of its orbit
closure to the first four variables is exactly
\begin{equation}\label{R14:paddingvariety}
 \mathcal Z_{m,n}=
 \{\ell^{n-m}q:\ell\in(\mathbb C^4)^*,\ q\in\mathcal P_{m,4}\},
 \qquad \dim\mathcal Z_{m,n}=c_m+3.
\end{equation}
If
\begin{equation}\label{R14:criterion}
 c_m+3>2n^2+2,
\end{equation}
then, for every sufficiently large coordinate-ring degree $d$,
there is a partition $\lambda\vdash nd$ of length exactly four
such that
\[
 m_\lambda(h_{m,n})>m_\lambda(\det_n).
\]
No effective threshold for $d$ is asserted.

For $4\leq m\leq16$, R-13 gives $c_m=\binom{m+3}{3}$, so the
sufficient condition becomes
\begin{equation}\label{R14:explicitrange}
 m<n\leq\left\lfloor\sqrt{\frac12\binom{m+3}{3}}\right\rfloor.
\end{equation}
The nonempty padded ranges are
\[
\begin{array}{c|ccccccccc}
m&8&9&10&11&12&13&14&15&16\\\hline
n&9&10&11&12\text{--}13&13\text{--}15&14\text{--}16&
15\text{--}18&16\text{--}20&17\text{--}22
\end{array}.
\]
In particular, some four-row multiplicity obstruction separates
the unpadded $\operatorname{per}_n$ and $\det_n$ for every
$4\leq n\leq22$; for $17\leq n\leq22$ this follows by
block-diagonal containment from the $m=16$ padded cases.

\begin{proof}
First, $\mathcal P_{m,4}$ is an irreducible affine cone, as the
closure of the image of an irreducible affine space. It contains
an irreducible degree-$m$ polynomial: use the cyclic matrix with
diagonal entries $x_1$, entries $(i,i+1)$ equal to $x_2$ for
$i<m$, entry $(m,1)$ equal to $x_3$, and every other entry zero.
Its only perfect matchings are the diagonal and the full cycle,
so its permanent is $x_1^m+x_2^{m-1}x_3$.
Viewed in $\mathbb C[x_2,x_3][x_1]$, this polynomial is
Eisenstein at the prime $x_3$ and hence irreducible.
The reducible degree-$m$ forms constitute a closed subset:
projectivize each of the finitely many multiplication maps of
degrees $a$ and $m-a$, $1\leq a<m$, and take their closed images.
It follows that a general $q\in\mathcal P_{m,4}$ is irreducible.

Projectivizing the multiplication map
$(\ell,q)\mapsto\ell^{n-m}q$ gives a morphism
\[
 \mathbb P^3\times\mathbb P(\mathcal P_{m,4})
 \longrightarrow\mathbb P(\operatorname{Sym}^n\mathbb C^4)^*.
\]
Its image is closed because its domain is projective; its affine
cone is exactly the set in \eqref{R14:paddingvariety}, including
zero. For a generic irreducible $q$ of degree $m\geq2$, the
linear factor $\ell$ is uniquely determined by the product,
up to scalar, and then so is $q$. Thus the projective map is
generically injective. Its image dimension is
$3+(c_m-1)$ and the affine-cone dimension is $c_m+3$.

Every four-variable restriction of an invertible substitution
in $h_{m,n}$ has the form in \eqref{R14:paddingvariety}, so
closedness gives one inclusion of the projected orbit closure.
Conversely, for each permanent pencil $q=\Phi_m(A)$, any chosen
$\ell$ and its $m^2$ matrix-entry linear forms can be prescribed
simultaneously by an endomorphism substitution, since the
padding variable and the permanent variables are independent.
There is enough room because $m^2+1\leq n^2$.
Such substitutions lie in the orbit closure. Taking limits in
$q$ then shows that the whole closed set $\mathcal Z_{m,n}$ is
contained in the orbit closure, supported on its first four
variables. This proves equality of the projected image.

Let $\mathcal D_{n,4}$ be the determinant four-variable image
closure. By G-06, $\dim\mathcal D_{n,4}\leq2n^2+2$.
The coordinate rings of both cones are standard graded. The
Hilbert function of a nonzero standard graded affine cone of
dimension $c$ agrees for large $d$ with a polynomial of degree
$c-1$ and positive leading coefficient. Thus
\eqref{R14:criterion} implies
\[
 \dim\mathbb C[\mathcal Z_{m,n}]_d
 >\dim\mathbb C[\mathcal D_{n,4}]_d
\]
for every sufficiently large $d$.
Decompose these finite-dimensional $GL_4$ representations into
irreducibles. Their dimensions are the sums of multiplicities
times the positive integers $\dim S_\lambda\mathbb C^4$.
Hence some irreducible multiplicity on the first side is larger.

These $GL_4$ multiplicities agree with the corresponding
$GL_{n^2}$ multiplicities of length at most four in the original
orbit-closure coordinate rings. Indeed, highest-weight vectors
with those weights use only the first-four coefficient variables,
and the intersection of the full ideal with that polynomial
subring is the ideal of its projected image. The additional
raising operators involving later variables vanish on such
vectors. This proves the claimed multiplicity inequality.
G-05 states that the determinant multiplicity is already the
ambient maximum for weights of length at most three, so the
strict inequality must occur at length exactly four.

Substitution of $c_m=\binom{m+3}{3}$ into
\eqref{R14:criterion} gives $2n^2\leq\binom{m+3}{3}$, which is
\eqref{R14:explicitrange} and the displayed finite table.
Finally
\[
 \operatorname{per}\operatorname{diag}(L_m,zI_{n-m})
 =z^{n-m}\operatorname{per}(L_m)
\]
shows that the padded-permanent orbit closure is contained in
the size-$n$ permanent orbit closure. The resulting coordinate
quotient is surjective and therefore cannot increase
multiplicities. The padded gaps at $m=16$ thus imply the stated
unpadded gaps for $17\leq n\leq22$.
\end{proof}

If the expected formula
$c_m=\min\{\binom{m+3}{3},4m^2-2m+2\}$ were proved beyond the
certified range, \eqref{R14:criterion} would turn it into further
padded obstructions up to approximately $n=\sqrt2\,m$.
This expectation is a separate generic-fiber problem; neither
the finite certificates nor the dimension upper bound proves it.
The criterion above is unconditional whenever its stated
dimension lower bound has been established.

\subsection{R-15: two chords cross the determinant dimension for every size}

\paragraph{Statement R-15.}
For every $n\geq4$,
\begin{equation}\label{R15:strictdimension}
 \dim\mathcal P_{n,4}\geq2n^2+3
 >\dim\mathcal D_{n,4}.
\end{equation}
Consequently, for each fixed $n\geq4$ and every sufficiently
large coordinate-ring degree $d$, there is a partition
$\lambda\vdash nd$ with exactly four rows such that
\begin{equation}\label{R15:allngap}
 m_\lambda(\operatorname{per}_n)>m_\lambda(\det_n).
\end{equation}
The proof is algebraic and applies to all $n\geq4$ without an
extrapolation from the finite certificates. It supplies neither
an effective threshold for $d$ nor a specified $\lambda$ for
general $n$. In particular it is an unpadded, equal-size result,
and does not assert a growing-size determinantal-complexity lower
bound. At $n=4$, the lower bound is $35$, the full dimension of
the quaternary quartic space; it gives a purely algebraic proof
of the quartic dominance used in R-12.

\begin{proof}
Choose distinct complex numbers $a_1,\ldots,a_n$ and, separately,
distinct $b_1,\ldots,b_n$. Put
\[
 \ell_i=x-a_i y,\qquad v_i=z-b_i w,
 \qquad L=\prod_{i=1}^n\ell_i,
\]
with indices read cyclically modulo $n$.
Let $M_0$ have diagonal entries $\ell_i$, cycle entries
$(M_0)_{i,i+1}=v_i$, including $(M_0)_{n,1}=v_n$, and every
other entry zero. Write
$C_{ij}=\operatorname{per}(M_{0,\widehat i,\widehat j})$.
The image of the differential of the full permanent-pencil map
at $M_0$ is
\[
 W_0=\operatorname{span}\{xC_{ij},yC_{ij},zC_{ij},wC_{ij}:1\leq i,j\leq n\}.
\]
The complementary bipartite graph is a union of even paths,
each with a unique perfect matching. More explicitly, if $I$
is the cyclic interval from $j$ to $i-1$, of length
$k\in\{0,\ldots,n-1\}$, then
\begin{equation}\label{R15:cyclecofactor}
 C_{ij}=\prod_{h\in I}v_h
        \prod_{h\notin I\cup\{i\}}\ell_h.
\end{equation}

We first prove $\dim W_0\geq2n^2+2$.
Let $\Lambda_i\subset\mathbb P^3$ be the line
$\ell_i=v_i=0$, parametrized by
$(x,y,z,w)=(a_i u,u,b_i v,v)$.
Every cofactor from row $h\ne i$ has a factor $\ell_i$ or $v_i$
and restricts to zero on $\Lambda_i$.
The $n$ row-$i$ cofactors in \eqref{R15:cyclecofactor} restrict
to nonzero scalar multiples of
$u^{n-1-k}v^k$, $0\leq k<n$.
Multiplication by $y$ and $w$ therefore spans every degree-$n$
binary monomial on this line. Thus the restriction map
\[
 W_0\longrightarrow\bigoplus_{i=1}^n H^0(\Lambda_i,\mathcal O(n))
\]
is surjective, and its image has dimension $n(n+1)$.

For each proper nonempty cyclic interval $I$, together with
$I=\varnothing$ and $I=[n]$, put
\[
 Q_I=\prod_{i\in I}v_i\prod_{i\notin I}\ell_i.
\]
These polynomials belong to $W_0$: multiply the cofactor whose
omitted row immediately follows $I$ by the missing $\ell_i$;
for the two endpoint intervals use
$Q_\varnothing=\ell_iC_{ii}$ and
$Q_{[n]}=v_iC_{i,i+1}$.
Each $Q_I$ vanishes on every $\Lambda_i$.
They are linearly independent. For fixed length $k$ with
$1\leq k<n$, evaluate on the auxiliary line
$\Theta_t=\{\ell_t=v_{t+1}=0\}$.
The restriction of $Q_I$ is nonzero precisely when $I$ contains
$t$ and excludes $t+1$, which for a cyclic interval of that
length means that $I$ ends at $t$. The distinctness assumptions
ensure that the surviving scalar is nonzero. Thus these $n$
polynomials are independent for each $k$. Different lengths have
different bidegrees in $(x,y)$ and $(z,w)$, including the two
endpoint products. The kernel of the restriction map therefore
has dimension at least $n(n-1)+2$. Hence
\begin{equation}\label{R15:baserank}
 \dim W_0\geq n(n+1)+n(n-1)+2=2n^2+2.
\end{equation}

Next we identify enough of $W_0$ to exhibit a vector outside it.
The bidegree-$(n-2,2)$ component of $W_0$ is exactly
\begin{equation}\label{R15:edgespace}
 \left\{
 L\sum_{i=1}^n
 \frac{v_i g_i(z,w)}{\ell_i\ell_{i+1}}:
 g_i\in\operatorname{span}\{z,w\}
 \right\}.
\end{equation}
One inclusion follows by taking the cofactors whose interval has
length one and multiplying by $z,w$.
For the converse, a generator of this bidegree either has that
form, or comes from a length-two interval and multiplication
by a linear form $h(x,y)$.
Its rational expression after division by $L$ is
\[
 \frac{h(x,y)v_i v_{i+1}}{\ell_i\ell_{i+1}\ell_{i+2}}.
\]
Since $\ell_i$ and $\ell_{i+2}$ are independent, write
$h=A\ell_i+B\ell_{i+2}$ with scalars $A,B$.
The expression becomes the sum of an edge-$(i+1,i+2)$ term
with numerator $A v_i v_{i+1}$ and an edge-$(i,i+1)$ term
with numerator $B v_i v_{i+1}$. Both belong to
\eqref{R15:edgespace}. This proves equality; all generators of
the bidegree component have been included.

Define the degree-$n$ polynomial
\[
 T=\ell_1v_1v_3\prod_{i=4}^n\ell_i.
\]
We claim $T\notin W_0$. It has bidegree $(n-2,2)$, so
\eqref{R15:edgespace} is sufficient to test this claim.
On the affine chart $y=w=1$, put $t=x$, $u=z$.
Every rational expression in \eqref{R15:edgespace}, after
division by $L(t)$, has a partial fraction expansion
\[
 \sum_{i=1}^n\frac{R_i(u)}{t-a_i},
 \qquad R_i=E_i-E_{i-1},
\]
where $E_i(u)$ is a polynomial of degree at most two divisible
by $u-b_i$ and $E_0=E_n$. Indeed, take
$E_i=(u-b_i)g_i(u)/(a_i-a_{i+1})$ and use partial fractions
on each edge denominator.
For our particular polynomial,
\[
 \frac{T}{L}
 =\frac{(u-b_1)(u-b_3)}{(t-a_2)(t-a_3)}.
\]
Its residues vanish except at $a_2,a_3$, where they are
opposite nonzero scalar multiples of $(u-b_1)(u-b_3)$.
Thus $R_i=E_i-E_{i-1}$ would force all $E_i$ with $i\ne2$
to equal one common quadratic $E_*$. Divisibility at
$i=1,3,4$ gives three distinct roots $b_1,b_3,b_4$ of $E_*$,
so $E_*=0$.
Then $E_2$ is a nonzero scalar multiple of
$(u-b_1)(u-b_3)$, which is not divisible by $u-b_2$.
This contradiction proves $T\notin W_0$.

Now perturb $M_0$ by two entries:
\[
 (M_\varepsilon)_{1,3}=p=\varepsilon\ell_1,
 \qquad (M_\varepsilon)_{2,1}=q=\varepsilon v_3,
\]
leaving the other entries unchanged.
Let $C_{ij}(\varepsilon)$ be the complementary permanents at
this pencil and set $R=\prod_{i=4}^n\ell_i$.
Enumerating the two possible matchings in the indicated
complementary matrices gives
\[
 \begin{aligned}
 C_{31}(\varepsilon)&=(v_1v_2+p\ell_2)R,\\
 C_{32}(\varepsilon)&=(\ell_1v_2+pq)R,\\
 C_{33}(\varepsilon)&=(\ell_1\ell_2+v_1q)R.
 \end{aligned}
\]
Consequently the exact polynomial identity
\begin{equation}\label{R15:chordidentity}
 \ell_1C_{31}(\varepsilon)-v_1C_{32}(\varepsilon)
       -pC_{33}(\varepsilon)
       =-2\varepsilon^2T
\end{equation}
holds. Its left side belongs to the full Jacobian image
$W_\varepsilon$, since each multiplier is a linear form in
$x,y,z,w$. Thus $T\in W_\varepsilon$ whenever $\varepsilon\ne0$.

Choose $r=2n^2+2$ Jacobian columns independent at
$\varepsilon=0$, using \eqref{R15:baserank}.
Together with $T$ they are independent there, since
$T\notin W_0$. Some $(r+1)\times(r+1)$ coefficient minor of
these columns and $T$ is therefore nonzero at zero. It remains
nonzero for all $\varepsilon$ outside a finite set. For such
nonzero $\varepsilon$, $T$ already belongs to $W_\varepsilon$
by \eqref{R15:chordidentity}, so
$\dim W_\varepsilon\geq r+1=2n^2+3$.
The rank of a polynomial map's differential is a lower bound
for its image dimension, proving the first inequality in
\eqref{R15:strictdimension}. The second follows from G-06.

This construction can be made with integer coefficients.
For example take $a_i=i$, $b_i=i^2$.
Each Jacobian column has degree at most two in $\varepsilon$,
because only two matrix entries involve that parameter.
The augmented minor above has degree at most $2r=4n^2+4$ and
is not identically zero. Hence at least one integer
$\varepsilon\in\{1,\ldots,4n^2+5\}$ works.

Finally, the Hilbert-function and highest-weight argument in
R-14 applies to the two image cones in
\eqref{R15:strictdimension}. Their dimension inequality forces
a strict coordinate multiplicity inequality in every sufficiently
large degree. The determinant already has ambient multiplicities
for weights of length at most three by G-05, so the strict
inequality has length exactly four. This proves
\eqref{R15:allngap}.
\end{proof}

\subsection{R-16: simultaneous chords give an unbounded genuine-padding range}

\paragraph{Statement R-16.}
For $m\geq4$ put
\[
 h_m=\min\{\lfloor m/2\rfloor,m-3\},\qquad
 F_m=h_m(m-3).
\]
Thus $h_4=1$ and $h_m=\lfloor m/2\rfloor$ for $m\geq5$.
The quaternary permanent image satisfies the uniform bound
\begin{equation}\label{R16:dimension}
 \dim\mathcal P_{m,4}\geq 2m^2+2+F_m.
\end{equation}
At $m=4$ and $m=5$, this lower bound is respectively $35$ and
$56$, the whole ambient dimension. Thus it also proves the
quartic and quintic dominance of R-12 by an algebraic argument.
Consequently, if $n>m\geq4$ and
\begin{equation}\label{R16:paddingrange}
 2n^2\leq2m^2+F_m+2,
\end{equation}
then, in every sufficiently large coordinate-ring degree $d$,
some partition $\lambda\vdash nd$ with exactly four rows has
\[
 m_\lambda\bigl(\ell^{n-m}\operatorname{per}_m\bigr)
   >m_\lambda(\det_n).
\]
In particular, $n=m+1$ satisfies this conclusion for every
$m\geq12$. Together with the finite certificates in R-14,
this proves the conclusion for $n=m+1$ and every $m\geq8$.
For every fixed real $1<\rho<\sqrt5/2$, the same conclusion
holds for $n=\lfloor\rho m\rfloor$ and all sufficiently large
$m$. These are multiplicity gaps in a specified range of
linear padding; the proof does not specify the partition or
give a polynomial-time procedure to find it.

\begin{proof}
Fix a primitive $m$th root of unity $\zeta$ and use indices
$0,\ldots,m-1$ modulo $m$. Specialize the cyclic pencil of
R-15 to
\[
 \ell_i=x-\zeta^i y,\qquad v_i=z-\zeta^i w.
\]
Let $M_0$ have diagonal $\ell_i$ and cycle entries
$(M_0)_{i,i+1}=v_i$. Its differential image is denoted $W_0$.
The restriction-line and interval-product proof of R-15 applies
to this specialization, since both lists of roots are distinct,
and gives $\dim W_0\geq2m^2+2$.

We will simultaneously deform the disjoint row pairs
\[
 \mathcal A=\{0,2,\ldots,2(h_m-1)\},\qquad
 \{a,a+1\}\quad(a\in\mathcal A).
\]
For each $a\in\mathcal A$ add the two entries
\begin{equation}\label{R16:simultaneouspencil}
 (M_\varepsilon)_{a,a+2}=p_a=\varepsilon\ell_a,
 \qquad
 (M_\varepsilon)_{a+1,a}=q_a=\varepsilon v_{a+2}.
\end{equation}
All other entries remain those of $M_0$. The pairs of rows are
disjoint, including when the last triple of column indices
wraps around the cycle. Consequently rows $a,a+1$ have support
only in the three columns $a,a+1,a+2$, even when every other
pair has also been perturbed.

For $r\notin\{a,a+1\}$ let $R_{a,r}(\varepsilon)$ be the
permanent obtained by deleting rows $a,a+1,r$ and columns
$a,a+1,a+2$. Write $C_{r,j}(\varepsilon)$ for the complementary
permanent deleting row $r$ and column $j$. In each of the
three cases $j=a,a+1,a+2$, the two rows $a,a+1$ must match
the other two columns of that triple: they have no entries
outside it, and exactly two of its columns remain. The rest
of the matching is arbitrary in the same complementary
$(m-3)\times(m-3)$ matrix. Therefore, with no hypothesis on
the other row pairs,
\[
 \begin{aligned}
 C_{r,a}&=(v_a v_{a+1}+p_a\ell_{a+1})R_{a,r},\\
 C_{r,a+1}&=(\ell_a v_{a+1}+p_aq_a)R_{a,r},\\
 C_{r,a+2}&=(\ell_a\ell_{a+1}+v_aq_a)R_{a,r}.
 \end{aligned}
\]
Here and in the next display all permanents are evaluated at
$\varepsilon$. Subtraction gives the exact identity
\begin{equation}\label{R16:simultaneousidentity}
 \ell_a C_{r,a}-v_a C_{r,a+1}-p_a C_{r,a+2}
   =-2\varepsilon^2\ell_a v_a v_{a+2}R_{a,r}(\varepsilon).
\end{equation}
It follows that the polynomial vector
\[
 T_{a,k}(\varepsilon)
   =\ell_a v_a v_{a+2}R_{a,a+k}(\varepsilon)
\]
belongs to the full differential image $W_\varepsilon$ for
every $\varepsilon\ne0$ and $2\leq k\leq m-2$.
At $\varepsilon=0$ the remaining cycle paths have a unique
matching, so
\begin{equation}\label{R16:normalvectors}
 T_{a,k}:=T_{a,k}(0)
 =\ell_a v_a v_{a+2}
   \left(\prod_{i=a+2}^{a+k-1}v_i\right)
   \left(\prod_{i\notin[a,a+k]}\ell_i\right).
\end{equation}
Intervals and products are cyclic; the first product is empty
when $k=2$. This polynomial has bidegree $(m-k,k)$ in
$(x,y)$ and $(z,w)$.

We next prove that all $h_m(m-3)$ vectors in
\eqref{R16:normalvectors} are linearly independent modulo
$W_0$. This is the step that permits simultaneous rank growth,
rather than adding ranks from different deformations.

For $1\leq k<m$ set
\[
 F_{a,k}
   =\left(\prod_{i=a}^{a+k-2}v_i\right)
     \left(\prod_{i\notin[a,a+k-1]}\ell_i\right).
\]
The bidegree-$(m-k,k)$ part of $W_0$ is exactly
\begin{equation}\label{R16:bidegreespace}
 (W_0)_{m-k,k}
   =\operatorname{span}\{zF_{a,k},wF_{a,k}:0\leq a<m\}.
\end{equation}
Indeed the cycle cofactors with $k-1$ factors $v_i$ give
these generators directly. Every other generator of the
same bidegree is a cofactor with $k$ such factors multiplied
by a linear form $u(x,y)$. Write
$u=A\ell_a+B\ell_{a+k}$, possible because these two
linear forms are independent. The resulting polynomial is
$A v_aF_{a+1,k}+B v_{a+k-1}F_{a,k}$, proving the reverse
inclusion. This accounts for every cofactor and multiplier
of that bidegree. The two endpoint bidegrees have dimensions
at most $m+1$, so this argument also gives
$\dim W_0\leq2(m+1)+2m(m-1)=2m^2+2$.

Let $\tau$ act by
\[
 \tau(x,y,z,w)=(x,\zeta y,z,\zeta w).
\]
Then $F_{a,k}=\tau^aF_{0,k}$ and
$T_{a,k}=\tau^aT_{0,k}$. Equation
\eqref{R16:bidegreespace} shows that $W_0$ is invariant
under this cyclic action. Let $\pi_t$ denote its character
projection with eigenvalue $\zeta^t$.
We claim that
\begin{equation}\label{R16:allmodes}
 \pi_t(T_{0,k})\notin\pi_t(W_0)
 \quad\text{for every }2\leq k,t\leq m-2.
\end{equation}

Here is the full coefficient calculation proving the claim.
Write
\[
 \begin{aligned}
 A(x,y)&=\prod_{i=k}^{m-1}(x-\zeta^iy)
       =\sum_r A_r x^{m-k-r}y^r,\\
 B(z,w)&=\prod_{i=0}^{k-2}(z-\zeta^iw),\\
 E(z,w)&=\prod_{i=1}^{k-2}(z-\zeta^iw),\qquad B=(z-w)E.
 \end{aligned}
\]
Use subscripts for coefficients in the powers of the second
variable, with coefficients outside the degree range defined
to be zero. Every $A_r$, $0\leq r\leq m-k$, is nonzero.
For completeness the finite geometric-product identity gives
\[
 A_r=(-1)^r\zeta^{kr+r(r-1)/2}
      \prod_{j=1}^{r}
       \frac{1-\zeta^{m-k-r+j}}{1-\zeta^j}.
\]
This identity follows by induction on the number of linear
factors. Every exponent in numerator or denominator lies
strictly between $0$ and $m$, so no factor vanishes.

We have $F_{0,k}=AB$. The $(x,y)$ factor of $T_{0,k}$ is
$A(x,\zeta y)$, since its roots are $k+1,\ldots,m-1,0$.
If its $(z,w)$ factor is $B'(z,w)$, then
\[
 \begin{aligned}
 D(z,w):=B'(z,\zeta^{-1}w)
   &=(z-\zeta^{-1}w)(z-\zeta w)E(z,w)\\
   &=(z-w)B(z,w)+c\,zwE(z,w),\\
 c&=2-\zeta-\zeta^{-1}\ne0.
 \end{aligned}
\]
These equalities hold also for $k=2$, when $E=1$.

In character $t$, with $2\leq t\leq m-2$, the monomials
are precisely those with $r+s=t$; a second possibility
$r+s=t+m$ exceeds total degree $m$.
The allowed indices $s$ form the interval
\[
 \max\{0,t-(m-k)\}\leq s\leq\min\{k,t\}.
\]
It contains at least three consecutive integers, because
$2\leq k,t\leq m-2$.
The coefficient of such a monomial in $T_{0,k}$ is
$\zeta^t A_{t-s}D_s$. The two possible generators for
the character-$t$ part of \eqref{R16:bidegreespace}, namely
the projections of $zF_{0,k}$ and $wF_{0,k}$, have
coefficients $A_{t-s}B_s$ and $A_{t-s}B_{s-1}$.
Since all the displayed $A_{t-s}$ are nonzero, membership
in their span would imply that $D_s$ is a fixed linear
combination of $B_s$ and $B_{s-1}$ on the allowed interval.
Subtracting $(z-w)B$ and using $c\ne0$, it would follow
that $E_{s-1}$ is also such a linear combination.

That last assertion is impossible on any three consecutive
allowed indices $s=j,j+1,j+2$. To see this, use
$B_s=E_s-E_{s-1}$ to perform invertible column operations:
the relevant three-column determinant is nonzero if and
only if
\begin{equation}\label{R16:toeplitz}
 \det\bigl(E_{j+u-v}\bigr)_{0\leq u,v\leq2}\ne0.
\end{equation}
Here $0\leq j\leq k-2$. Put $N=k-2$.
Since $E_i=(-1)^i e_i(\zeta,\ldots,\zeta^N)$, row and
column signs change \eqref{R16:toeplitz} to the determinant
\[
 \det\bigl(e_{j+u-v}(\zeta,\ldots,\zeta^N)\bigr)_{0\leq u,v\leq2}
   =s_{(3^j)}(\zeta,\ldots,\zeta^N),
\]
by the dual Jacobi--Trudi identity (its representation-ring
form is equation~(2.4) in the attached proof of Sawin's Lemma~2.3;
R-B15). This Schur value is nonzero. If $N=0$, necessarily $j=0$ and the determinant
is one. Otherwise use the Weyl alternant at
$(1,\zeta,\ldots,\zeta^{N-1})$, then multiply by the
nonzero scalar $\zeta^{3j}$. For
$\lambda=(3^j,0^{N-j})$, the exponents
$\lambda_i+N-i$ are distinct integers in $[0,N+2]=[0,k]$.
Since $k<m$, their powers of $\zeta$ are distinct. The
alternant numerator is therefore a nonzero Vandermonde
determinant. Its denominator has the distinct exponents
$N-i$ and is nonzero as well. This proves
\eqref{R16:toeplitz}, and hence \eqref{R16:allmodes}.

Now fix $k$ and suppose
$\sum_{a\in\mathcal A}u_aT_{a,k}\in W_0$.
Projecting to each of the characters in
\eqref{R16:allmodes} gives
\[
 \sum_{a\in\mathcal A}u_a\zeta^{at}=0
 \qquad(2\leq t\leq m-2).
\]
There are $h_m\leq m-3$ distinct indices $a$.
Use the $h_m$ consecutive equations
$t=2,\ldots,h_m+1$. Their matrix is an ordinary
Vandermonde matrix in the distinct numbers $\zeta^a$,
with each column multiplied by the nonzero factor
$\zeta^{2a}$. It is invertible, so every $u_a=0$.
Thus the $h_m$ vectors for this $k$ are independent
modulo $W_0$. Different $k$ have different bidegrees
and $W_0$ is bigraded, so all $h_m(m-3)$ vectors are
independent modulo $W_0$.

Choose $2m^2+2$ differential columns at $M_0$ forming a
basis of $W_0$, and continue the same columns at
$M_\varepsilon$. Adjoin all the polynomial vectors
$T_{a,k}(\varepsilon)$. At zero these columns are
independent by the preceding argument, so one coefficient
minor of size $2m^2+2+F_m$ is a polynomial in
$\varepsilon$ with nonzero constant term. For generic
nonzero $\varepsilon$ it remains nonzero, and all its
adjoined vectors lie in $W_\varepsilon$ by
\eqref{R16:simultaneousidentity}. Hence
$\dim W_\varepsilon\geq2m^2+2+F_m$.
Differential rank bounds image dimension from below,
proving \eqref{R16:dimension}.

For padding, R-14 gives
$\dim\mathcal Z_{m,n}=\dim\mathcal P_{m,4}+3$.
Under \eqref{R16:paddingrange} this dimension is at least
$2n^2+3$, strictly larger than the determinant upper
bound $2n^2+2$. The Hilbert-function and highest-weight
argument of R-14 therefore supplies the asserted
exactly-four-row multiplicity gaps in every sufficiently
large degree. The explicit degree threshold of G-10
also applies to these pairs.

For $m\geq12$ one has
$F_m\geq(m-1)(m-3)/2\geq4m$, so
\eqref{R16:paddingrange} holds at $n=m+1$.
R-14 already covers each remaining $8\leq m\leq11$.
Finally, $F_m=m^2/2+O(m)$, and the right side of
\eqref{R16:paddingrange}, divided by $2m^2$, tends to
$5/4$. This proves the stated asymptotic range for every
fixed $\rho<\sqrt5/2$.
\end{proof}

\paragraph{What this does and does not identify.}
The simultaneous construction proves a lower bound on the
four-variable projection dimension, without assuming that
the row--column torus accounts for every generic fiber.
The generic maximal-rank formula
$\min\{\binom{m+3}{3},4m^2-2m+2\}$ is not used or proved.
The lower bound here is enough to retain four-row
multiplicity gaps under an unbounded range of genuine
padding, while the low-row exclusion of G-09 applies at
larger padding. No literature-priority claim is made for
these dimension deductions.


\subsection{R-17: exact determinant dimension and forced nonrectangular gaps}

\paragraph{Statement R-17.}
The four-variable determinant image has exact affine dimension
\begin{equation}\label{R17:detdimension}
 \dim\mathcal D_{n,4}=2n^2+2\qquad(n\geq1).
\end{equation}
For $n\geq3$ its affine $\mathrm{SL}_4$ quotient consequently
has dimension $2n^2-13$.
Moreover, at each of the padded pairs
\[
 (m,n)=(15,20),\qquad (m,n)=(16,22),
\]
there is a strictly positive multiplicity difference in some
\emph{nonrectangular} partition of exactly four rows in every
sufficiently large coordinate degree. In fact the effective
thresholds $d_{20}$ and $d_{22}$ of G-10 apply. At the second
pair, every positive-degree rectangular multiplicity on the
padded side is zero whenever the rectangular label is integral.

\begin{proof}
At the unperturbed cyclic pencil $M_0$ in R-15, each
complementary matrix has a unique perfect matching.
Therefore its determinant is its permanent multiplied by a
sign. The $(i,j)$ derivatives of the determinant additionally
have the fixed cofactor sign $(-1)^{i+j}$; hence each of the
four corresponding determinant-pencil Jacobian columns is
a sign times the permanent-pencil column. Their full images
are equal to $W_0$. The base-rank proof in R-15 gives
$\dim W_0\geq2n^2+2$ for $n\geq4$.
The determinant image dimension is at least this differential
rank and is at most $2n^2+2$ by G-06. This proves
\eqref{R17:detdimension} in that range. The remaining cases
are the full spaces of linear forms, quadrics, and cubic
surfaces, of dimensions $4,10,20$, respectively, as established
in R-10--R-11 and the elementary linear case.

The stable Fermat point and the $15$-dimensional quotient-fiber
argument of G-11 now give the asserted exact quotient dimension.
This is a direct consequence of the explicit cyclic pencil and
the established group-action bound; no claim of literature
priority for the dimension formula is made.

For the two padded pairs, R-13 and R-14 give
\[
 \begin{array}{c|c|c|c}
 (m,n)&\dim\mathcal P_{m,4}&\dim\mathcal Z_{m,n}
                    &\dim\mathcal D_{n,4}\\ \hline
 (15,20)&\binom{18}{3}=816&819&802\\
 (16,22)&\binom{19}{3}=969&972&970.
 \end{array}
\]
The strict dimension inequalities give exactly-four-row
multiplicity gaps by R-14, with the effective thresholds of
G-10. But $20=4\cdot15/3$, so the equality case of G-12
excludes a positive rectangular gap for the first pair.
For the second, $22>4\cdot16/3$, so the strict case of G-12
makes every positive-degree padded rectangular multiplicity
zero. If $4\nmid nd$, there is no integral four-row rectangle
of size $nd$ in the first place. Thus in every degree covered
by the Hilbert-function argument the required partition is
nonrectangular at both pairs.
\end{proof}


\subsection{R-18: two normal families improve the genuine-padding range}

\paragraph{Statement R-18.}
For $m\geq6$ put
\[
 j_m=\min\{\lfloor m/2\rfloor,m-5\},\qquad
 G_m=2j_m(m-5).
\]
Then
\begin{equation}\label{R18:dimension}
 \dim\mathcal P_{m,4}\geq2m^2+2+G_m.
\end{equation}
For any $n>m\geq6$ satisfying
\begin{equation}\label{R18:padding}
 2n^2\leq2m^2+G_m+2,
\end{equation}
R-14 and G-10 give a four-row multiplicity gap in every
coordinate degree $d\geq d_n$. In addition G-11 gives a
gap for the explicit rectangular partition
$\lambda=(nd/4)^4$ at every degree
$d=E_n k$, $k\geq K_n$, with its displayed effective
constants. In particular these conclusions hold eventually
for $n=\lfloor\rho m\rfloor$ and every fixed
$1<\rho<\sqrt{3/2}$.
The bounds of R-16 and R-18 are both valid; one may use the
larger lower bound at each particular size.

\begin{proof}
Keep the root-of-unity cyclic pencil and notation of R-16,
but take $j_m$ disjoint pairs with starting indices
$\mathcal A=\{0,2,\ldots,2(j_m-1)\}$ and use the entries
\[
 p_a=\varepsilon\ell_{a+1}\quad\hbox{at }(a,a+2),
 \qquad q_a=\varepsilon v_{a+2}\quad\hbox{at }(a+1,a).
\]
Each pair of rows still has support in just its three
columns. Therefore the simultaneous factorization in
R-16 remains exact. It gives two identities:
\begin{align*}
 \ell_a C_{r,a}-v_a C_{r,a+1}-p_a C_{r,a+2}
    &=-2p_aq_a v_a R_{a,r},\\
 q_a C_{r,a}+\ell_{a+1} C_{r,a+1}-v_{a+1} C_{r,a+2}
    &= 2p_aq_a\ell_{a+1}R_{a,r}.
\end{align*}
Both follow by substituting the three two-by-two permanent
factors already displayed in R-16; no other entries occur
in that substitution. Thus for $\varepsilon\ne0$ the
differential image contains the polynomial families
\[
 \begin{aligned}
 U_{a,k}(\varepsilon)
    &=\ell_{a+1}v_a v_{a+2}R_{a,a+k}(\varepsilon),\\
 V_{a,k}(\varepsilon)
    &=\ell_{a+1}^2v_{a+2}R_{a,a+k+1}(\varepsilon),
 \end{aligned}
\]
for $3\leq k\leq m-3$. Their values at zero, using the
same unique complementary matchings, are
\[
 \begin{aligned}
 U_{a,k}&=\ell_{a+1}v_a v_{a+2}
          \prod_{i=a+2}^{a+k-1}v_i
          \prod_{i\notin[a,a+k]}\ell_i,\\
 V_{a,k}&=\ell_{a+1}^2v_{a+2}
          \prod_{i=a+2}^{a+k}v_i
          \prod_{i\notin[a,a+k+1]}\ell_i.
 \end{aligned}
\]
They both have bidegree $(m-k,k)$, and their cyclic
translates are obtained by the same operator $\tau$ as
in R-16.

We prove that the two classes
$\pi_t(U_{0,k})$, $\pi_t(V_{0,k})$ are independent
modulo $\pi_t(W_0)$ whenever
\begin{equation}\label{R18:interior}
 3\leq k,t\leq m-3.
\end{equation}
This will give two independent Vandermonde systems.
Write $q=\zeta$, $K=q^k$, and $X=q^t$ just for the
coefficient calculation below. Retain
$A(x,y)=\prod_{i=k}^{m-1}(x-q^iy)$ from R-16, whose
coefficients $A_r$ are nonzero, and put
\[
 H(z,w)=\prod_{i=2}^{k-2}(z-q^iw),\qquad
 B(z,w)=(z-w)(z-qw)H(z,w).
\]
Thus $F_{0,k}=AB$. Define
\[
 \begin{aligned}
 \alpha&=\frac{K-q}{K-1},&
 \beta&=\frac{q-1}{K-1},\\
 \gamma_2&=\frac{q(K-1)}{qK-1},&
 \gamma_1&=\frac{(q-1)(K-q)}{(K-1)(qK-1)},&
 \gamma_0&=\frac{(q-1)^2}{(K-1)(qK-1)}.
 \end{aligned}
\]
All denominators are nonzero in the range
\eqref{R18:interior}.

The $(x,y)$ factors of $U_{0,k}$ and $V_{0,k}$ are
\[
 A_U=(x-qy)\prod_{i=k+1}^{m-1}(x-q^iy),\qquad
 A_V=(x-qy)^2\prod_{i=k+2}^{m-1}(x-q^iy).
\]
Their coefficient ratios are
\begin{equation}\label{R18:Aratios}
 \frac{(A_U)_r}{A_r}=\alpha q^r+\beta,
 \qquad
 \frac{(A_V)_r}{A_r}
    =\gamma_2 q^{2r}+\gamma_1 q^r+\gamma_0.
\end{equation}
Here is an explicit verification using the finite product
formula from R-16. If $u=q^r$, the first ratio is
$((K-q)u+q-1)/(K-1)$. For the second, put
$D=\prod_{i=k+2}^{m-1}(x-q^iy)$.
The same formula gives
\[
 \frac{D_r}{A_r}
  =\frac{(Ku-1)(qKu-1)}{(K-1)(qK-1)},\quad
 \frac{D_{r-1}}{A_r}
  =\frac{(u-1)(Ku-1)}{(K-1)(qK-1)},\quad
 \frac{D_{r-2}}{A_r}
  =\frac{(u-1)(u-q)}{q(K-1)(qK-1)}.
\]
Consequently $(A_V)_r/A_r$, obtained by subtracting
$2qD_{r-1}$ and adding $q^2D_{r-2}$, has numerator
\[
 (Ku-1)(qKu-1)-2q(u-1)(Ku-1)+q(u-1)(u-q).
\]
Its coefficients of $u^2,u,1$ are respectively
$q(K-1)^2$, $(q-1)(K-q)$, and $(q-1)^2$.
This proves \eqref{R18:Aratios}, including endpoint
coefficients interpreted as zero where necessary.

Let $v_b=z-q^bw$ in the next few displays. In character
$t$, after dividing every coordinate by the nonzero factor
$A_{t-s}$, the four vectors
$\pi_t(zF_{0,k})$, $\pi_t(wF_{0,k})$,
$\pi_t(U_{0,k})$, and $\pi_t(V_{0,k})$ have the coefficient
sequences of $H$ times the following four binary cubics:
\begin{equation}\label{R18:cubics}
 zv_0v_1,\quad wv_0v_1,\quad
 \mathcal U_X=\alpha Xv_{-1}v_1^2+\beta v_0v_2v_{k-1},\quad
 \mathcal V_X=\gamma_2X^2v_0^2v_1
               +\gamma_1Xv_1^2v_{k-1}
               +\gamma_0v_2v_{k-1}v_k.
\end{equation}
To check the shifts in this assertion, the $(z,w)$ factor
of $U_{0,k}$, after replacing $w$ by $q^{-1}w$, is
$v_{-1}v_1^2H$. That factor with $w$ replaced back by
$qw$ becomes $v_0v_2v_{k-1}H$.
The $(z,w)$ factor of $V_{0,k}$, after replacing $w$ by
$q^{-2}w$, is $v_0^2v_1H$; its next two shifts are
$v_1^2v_{k-1}H$ and $v_2v_{k-1}v_kH$.
Together with $q^r=q^{t-s}$ and
\eqref{R18:Aratios}, these identities give exactly
\eqref{R18:cubics}.

These four cubics are linearly independent. Modulo the
first two, it suffices to evaluate at $z/w=1,q$, because
the kernel of those two evaluations is precisely
$v_0v_1\operatorname{span}\{z,w\}$.
Direct substitution and simplification give
\begin{equation}\label{R18:evaldet}
 \begin{vmatrix}
 \mathcal U_X(1,1)&\mathcal V_X(1,1)\\
 \mathcal U_X(q,1)&\mathcal V_X(q,1)
 \end{vmatrix}
 =\frac{(q-K)(q^2-K)(q-1)^6(q+1)}
        {q(K-1)(qK-1)}\ne0.
\end{equation}
The cancellation of $X$ in this expression is exact.
Nonvanishing follows from the order $m$ of $q$ and
$3\leq k\leq m-3$; in particular $q\ne\pm1$ and
$K\notin\{1,q,q^2,q^{-1}\}$.

Under \eqref{R18:interior}, the allowed coefficient
indices $s$ contain at least four consecutive values
$b,b+1,b+2,b+3$, with $0\leq b\leq k-3$.
Multiplication by $H$ takes the cubic monomial basis
to four columns whose restriction to those coordinates
has determinant
\[
 \det(H_{b+u-v})_{0\leq u,v\leq3}.
\]
Up to row and column signs this is the dual
Jacobi--Trudi determinant
\[
 s_{(4^b)}(q^2,q^3,\ldots,q^{k-2}).
\]
It is nonzero by the same alternant argument as in
R-16: there are $N=k-3$ variables and the numerator
exponents for $(4^b)$ are distinct in
$[0,N+3]=[0,k]$, strictly shorter than the order $m$.
For $N=0$, necessarily $b=0$ and the determinant is one.
Thus restriction to the four selected coefficients
preserves independence of the cubics
\eqref{R18:cubics}. This proves the asserted independence
of the two normal classes in every character
\eqref{R18:interior}.

Now suppose, for fixed $k$, that a linear combination
\[
 \sum_{a\in\mathcal A}
        (c_aU_{a,k}+d_aV_{a,k})
\]
belongs to $W_0$. Projecting to character $t$ gives a
linear combination of two independent normal classes,
so both $\sum_a c_aq^{at}$ and
$\sum_a d_aq^{at}$ must vanish.
Since $j_m\leq m-5$, the $j_m$ consecutive values
$t=3,\ldots,j_m+2$ are available. Each of these two
systems is an invertible Vandermonde system in the
distinct $q^a$. Hence every coefficient vanishes.
The $2j_m$ normal vectors are independent for this
bidegree. Different $k$ have different bidegrees, so
the full collection has $2j_m(m-5)=G_m$ independent
classes modulo $W_0$.

Adjoin their polynomial deformations to a basis of
the $2m^2+2$ cyclic Jacobian columns. At zero the
augmented columns are independent. A nonzero coefficient
minor remains nonzero for generic $\varepsilon$, when
all these columns lie in the actual deformed Jacobian
image by the two exact cofactor identities. This proves
\eqref{R18:dimension} by differential rank.

The padding conclusion follows from R-14 and G-10.
For the rectangular conclusion, note that
$G_m\leq m(m-5)$, so \eqref{R18:padding} implies
$n^2<3m^2/2$ and hence $n<4m/3$.
Both hypotheses of G-11 therefore hold.
Finally $G_m=m^2+O(m)$, so the upper admissible ratio
in \eqref{R18:padding} tends to $\sqrt{3/2}$.
\end{proof}


\subsection{R-19: simultaneous row and column blocks improve the range again}

\paragraph{Statement R-19.}
For $m\geq10$ put
\[
 a_m=\min\{\lfloor m/3\rfloor,m-9\},\qquad
 H_m=4a_m(m-9).
\]
Then
\begin{equation}\label{R19:dimension}
 \dim\mathcal P_{m,4}\geq2m^2+2+H_m.
\end{equation}
If $n>m\geq10$ and
\begin{equation}\label{R19:padding}
 2n^2\leq2m^2+H_m+2,
\end{equation}
the four-row multiplicity gaps of R-14/G-10 and the explicit
rectangular gaps of G-11 both hold. In particular, they hold
eventually for $n=\lfloor\rho m\rfloor$ and every fixed
$1<\rho<\sqrt{5/3}$. Earlier lower bounds remain available
when they are larger at a particular finite size.

\begin{proof}
Use the same cyclic pencil with
$\ell_i=x-\zeta^iy$, $v_i=z-\zeta^iw$.
Choose centers
\[
 \mathcal A=\{1,4,\ldots,1+3(a_m-1)\}.
\]
The triples of rows $a-1,a,a+1$ are disjoint and any two
centers have cyclic distance at least three.
For each center add
\begin{equation}\label{R19:threechords}
 p_a=\varepsilon\ell_{a+1}\text{ at }(a,a+2),\qquad
 q_a=\varepsilon v_{a+2}\text{ at }(a+1,a),\qquad
 r_a=\varepsilon\ell_a\text{ at }(a-1,a+1).
\end{equation}
The two rows $a,a+1$ still have support only in columns
$a,a+1,a+2$. The two columns $a,a+1$ have support only
in rows $a-1,a,a+1$. To verify the latter assertion for
all simultaneous motifs, an additional forward chord
from another center $b$ could enter these columns only
if $b=a-2$ or $a-1$, and an additional backward or third
chord only if $b=a-1$ or $a+1$. These cases are excluded
by the spacing. The analogous row-support assertion is
immediate from the disjoint row triples.

Thus the two row identities of R-18 remain exact and
give the same families $U_{a,k}(\varepsilon)$ and
$V_{a,k}(\varepsilon)$, with the remaining complementary
permanents now evaluated at the three-chord pencil.
There are also two exact column identities. For a column
$c\notin\{a,a+1\}$ let $S_{a,c}(\varepsilon)$ be the
permanent obtained by deleting rows $a-1,a,a+1$ and
columns $a,a+1,c$. The complementary permanents factor as
\[
 \begin{aligned}
 C_{a-1,c}&=(\ell_a\ell_{a+1}+v_aq_a)S_{a,c},\\
 C_{a,c}&=(v_{a-1}\ell_{a+1}+q_ar_a)S_{a,c},\\
 C_{a+1,c}&=(v_{a-1}v_a+\ell_ar_a)S_{a,c}.
 \end{aligned}
\]
Indeed the two remaining rows of the triple must match
the two supported columns; all other matching choices
occur in the same complementary matrix.
Consequently
\begin{align*}
 v_{a-1}C_{a-1,c}-\ell_a C_{a,c}-q_a C_{a+1,c}
    &=-2q_ar_a\ell_a S_{a,c},\\
 r_a C_{a-1,c}+v_a C_{a,c}-\ell_{a+1}C_{a+1,c}
    &= 2q_ar_av_a S_{a,c}.
\end{align*}
After division by the nonzero scalar $\varepsilon^2$,
these show that the differential image contains
\[
 C_{a,k}(\varepsilon)=\ell_a^2v_{a+2}S_{a,a-k}(\varepsilon),
 \qquad
 D_{a,k}(\varepsilon)=\ell_av_av_{a+2}S_{a,a-k+1}(\varepsilon).
\]
We use $5\leq k\leq m-5$ below.
At zero the unique cycle matching gives
\[
 \begin{aligned}
 C_{a,k}&=\ell_a^2v_{a+2}
       \prod_{i=a-k}^{a-2}v_i
       \prod_{i\notin[a-k,a+1]}\ell_i,\\
 D_{a,k}&=\ell_av_av_{a+2}
       \prod_{i=a-k+1}^{a-2}v_i
       \prod_{i\notin[a-k+1,a+1]}\ell_i.
 \end{aligned}
\]
All four families $U,V,C,D$ have bidegree $(m-k,k)$,
and translating their center by $a$ applies $\tau^a$.

We prove fourfold normal independence in each character
\begin{equation}\label{R19:interior}
 5\leq k,t\leq m-5.
\end{equation}
As in R-18 write $q=\zeta$, $K=q^k$, $X=q^t$, and
divide character-$t$ coordinates by the nonzero
coefficients $A_{t-s}$ of
$A(x,y)=\prod_{i=k}^{m-1}(x-q^iy)$.
Set
\[
 J(z,w)=\prod_{i=2}^{k-4}(z-q^iw),\qquad
 v_b=z-q^bw.
\]
The two base columns and the two row-family columns,
after this normalization, are $J$ times
$v_{k-3}v_{k-2}$ times the four independent cubics
of \eqref{R18:cubics}. Thus these four quintics span
\begin{equation}\label{R19:rowquintics}
 v_{k-3}v_{k-2}\operatorname{Sym}^3\mathbb C^2.
\end{equation}

We compute the two column-family sequences explicitly.
Put $M=m-k$ and
$\widetilde A=A(x,q^{2-k}y)=\prod_{i=2}^{M+1}(x-q^iy)$.
Their $(x,y)$ factors at center zero are
\[
 A_C=(x-y)^2\prod_{i=2}^{M-1}(x-q^iy),\qquad
 A_D=(x-y)\prod_{i=2}^{M}(x-q^iy).
\]
Define
\[
 \begin{aligned}
 \delta_0&=\frac{K(q-1)}{q(K-1)},&
 \delta_1&=\frac{K-q}{q(K-1)},\\
 \eta_0&=\frac{K^2(q-1)^2}{q(K-1)(qK-1)},&
 \eta_1&=\frac{K(K-q)(q-1)}{q(K-1)(qK-1)},&
 \eta_2&=\frac{K-1}{qK-1}.
 \end{aligned}
\]
The following are polynomial identities:
\[
 \begin{aligned}
 A_D&=\delta_0\widetilde A(x,y)
           +\delta_1\widetilde A(x,q^{-1}y),\\
 A_C&=\eta_0\widetilde A(x,y)
           +\eta_1\widetilde A(x,q^{-1}y)
           +\eta_2\widetilde A(x,q^{-2}y).
 \end{aligned}
\]
For the first, cancel the common factors with indices
$2,\ldots,M$ and compare the remaining linear forms.
For the second cancel those with indices $2,\ldots,M-1$;
the remaining identity is
\[
 \eta_0(x-K^{-1}y)(x-qK^{-1}y)
 +\eta_1(x-qy)(x-K^{-1}y)
 +\eta_2(x-y)(x-qy)=(x-y)^2,
\]
which follows by its three coefficients. This verifies
the coefficient ratios without any genericity assumption.

Since $\widetilde A_r=q^{(2-k)r}A_r$,
the two column vectors, after the common coordinate
division and multiplication by the nonzero character
scalar $X^{k-2}$, are $J$ times the quintics
\begin{align*}
 \mathcal C_X={}&
 \eta_0v_{-2}v_{-1}v_0v_1v_k
 +\eta_1X^{-1}v_{-1}v_0v_1v_{k-3}v_{k+1}
 +\eta_2X^{-2}v_0v_1v_{k-3}v_{k-2}v_{k+2},\\
 \mathcal D_X={}&
 \delta_0v_{-1}v_0v_1v_{k-2}v_k
 +\delta_1X^{-1}v_0v_1v_{k-3}v_{k-1}v_{k+1}.
\end{align*}
For example the $(z,w)$ factor of $C_{0,k}$ has roots
$2,m-k,\ldots,m-2$; shifting $w$ by $q^{k-2}$ changes
them to $k,-2,-1,0,\ldots,k-4$.
Its next two shifts give the other two terms above.
The factor of $D_{0,k}$ has roots
$0,2,m-k+1,\ldots,m-2$; its corresponding two shifts
give the terms of $\mathcal D_X$. This checks every
factor and character scalar in these expressions.

To test independence from \eqref{R19:rowquintics},
evaluate the two new quintics at
$z/w=K/q^3$ and $K/q^2$, the distinct roots of
$v_{k-3}v_{k-2}$. Substitution gives
\begin{equation}\label{R19:evaldet}
\begin{aligned}
 &\begin{vmatrix}
 \mathcal C_X(K/q^3,1)&\mathcal D_X(K/q^3,1)\\
 \mathcal C_X(K/q^2,1)&\mathcal D_X(K/q^2,1)
 \end{vmatrix}\\
 &\qquad=
 \frac{K^6(q-K)(q^2-K)^2(q^3-K)^2(q^4-K)
             (q-1)^6(q+1)(q^2+q+1)}
      {q^{27}(K-1)(qK-1)}\ne0.
\end{aligned}
\end{equation}
Here too the cancellation of $X$ is exact.
Every displayed factor is nonzero: $m\geq10$ and
$5\leq k\leq m-5$, so $q$ has order neither two nor
three and $K$ differs from $1,q,q^2,q^3,q^4,q^{-1}$.
The four quintics in \eqref{R19:rowquintics} and the
two new ones consequently form a basis of all binary
quintics.

Under \eqref{R19:interior}, the allowed indices $s$
contain six consecutive values $b,\ldots,b+5$, where
$0\leq b\leq k-5$. Multiplication by $J$ and restriction
to these six coordinates has, in the quintic monomial
basis, determinant
\[
 \det(J_{b+u-v})_{0\leq u,v\leq5}
   =\pm s_{(6^b)}(q^2,q^3,\ldots,q^{k-4})\ne0.
\]
There are $N=k-5$ variables. The Weyl numerator
exponents for $(6^b)$ are distinct in
$[0,N+5]=[0,k]$, below the order $m$, so the same
Vandermonde proof as in R-16 applies. When $N=0$,
the determinant is one. Thus all six normalized
character columns are independent; equivalently,
the four classes of $U,V,C,D$ are independent modulo
the base character space.

A relation among their translates at centers in
$\mathcal A$ now gives four separate equations
$\sum_{a\in\mathcal A}c_aq^{at}=0$ in each of the
available characters. Since $a_m\leq m-9$, the
$a_m$ consecutive characters
$t=5,\ldots,a_m+4$ lie in the allowed range.
Each system is an invertible Vandermonde system in
the distinct $q^a$. Hence the $4a_m$ classes are
independent in each bidegree. Different bidegrees
are independent, giving $4a_m(m-9)=H_m$ classes
modulo $W_0$.

All four families have polynomial deformations lying
in the actual differential image for nonzero
$\varepsilon$. The augmented-minor argument of
R-16, with the $2m^2+2$ base columns, proves
\eqref{R19:dimension}.
R-14 and G-10 give the general four-row consequence.
Furthermore $H_m\leq(4m/3)(m-9)$, so
\eqref{R19:padding} implies
$n^2\leq5m^2/3-6m+1<16m^2/9$.
Thus $n<4m/3$ and G-11 also applies.
Finally $H_m=4m^2/3+O(m)$, giving the asserted
asymptotic ratio $\sqrt{5/3}$.
\end{proof}


\subsection{What the attached GCT IV, VII, and VIII positivity claims supply}

The attached GCT~IV is the later Blasiak--Mulmuley--Sohoni manuscript
corresponding to their 2015 memoir, with filename
\nolinkurl{Apr13KroneckerGCT4.tex}. Its \S12.1 theorem
(\nolinkurl{t nonstandard schur-weyl duality}) proves nonstandard
Schur--Weyl duality for general $\dim V,\dim W$; the later
explicit-basis limitation should not obscure this proved
structural result. Its Theorem~15.21(vii)--(viii)
(\nolinkurl{t main canonical basis}) proves a highest-weight crystal
count for $g(\alpha,\beta,\gamma)$ when
$\ell(\alpha),\ell(\beta)\leq2$, including the symmetric/exterior
split. Theorem~18.2 (\nolinkurl{t consequence for Kronecker}) makes
the count explicit, and Corollary~18.4 gives the generating
function used in R-10. This supplies a positive rule for rectangles
of height two, or width two after simultaneous conjugation.
The construction for the full tensor power remains Conjecture~19.1
(\texttt{cj canonical basis X\^{}r}); \S19.2 describes the failed
flatness of a naive extension beyond the two-row case.

Even in the proved case, counting highest-weight elements does not
mean that every matrix coefficient is nonnegative. GCT~IV,
Example~17.8 (\nolinkurl{ex not positive}), displays the divided-power
coefficient $q^{10}-q^4-q^{-4}+q^{-10}$. This concerns the specified
basis and operator; it is a concrete warning against replacing a
precise canonical-basis positivity statement by unrestricted
coefficient positivity.

GCT~VII (Mulmuley, 2008a, revised extended abstract) proves in
Theorem~1.2 (\nolinkurl{tnonstd}) existence of the nonstandard
quantization, complete reducibility of its polynomial
representations, and its Peter--Weyl statement. The nonstandard
duality and reciprocity needed to recover plethysm are explicitly
Conjectures~1.3--1.4 (\nolinkurl{cduality}, \nolinkurl{creciprocity});
the refined form is Conjecture~6.1 (\nolinkurl{crecirefined}).
The proposed formula
$a_{\lambda,\mu}^{\pi}=\sum_\alpha
m_\lambda^\alpha n_\pi^\alpha$ is conditional on those structures.
The introduction states that correctness of the basis algorithms
and their required properties has not been proved there.

GCT~VIII (Mulmuley, 2008b, revised extended abstract) makes the
additional gates explicit: local crystal existence is Conjecture~2.2
(\nolinkurl{clocalcrystalb}); minimum-degree globalization is
Conjectures~2.5, 3.1--3.2; cellular decomposition is Conjectures~2.6,
3.3 and the specialization-compatible quasi-cell condition is
Conjecture~3.7. Kronecker positivity is Conjectures~2.7 and~3.4,
with specified signs and powers of $q-q^{-1}$ allowed. The general
plethysm setting uses the weaker saturation and nonstandard
positivity Conjectures~2.8--2.9 and~3.5--3.6. Finally, \S6
(\nolinkurl{scomplexity}) also requires efficient combinatorial
labels and crystal operations to obtain a $\#P$ formula; these
complexity properties are themselves proposed rather than proved.

Its discussion of the Riemann hypothesis in the introduction and
\S2.4.2 is a suggested extension of finite-field purity and
perverse-sheaf machinery to the nonflat, noncommutative objects at
issue. It is not a theorem deriving these conjectures from the
classical RH or Dirichlet GRH. Accepting the stipulated analytic
theorems therefore leaves specific algebraic gates: construct the
required realization, prove its cellular and specialization
properties, and control the sizes of its labels and operations.
The supplied papers prove the height-two case and propose this
general route; they supply no unconditional constructive positive
formula for arbitrary rectangles and no theorem controlling the
orbit-closure ranks in R-05--R-06.

\paragraph{Algebraic-geometric dictionary for R-01--R-08.}
R-01 and R-02 are identities of Schur operations and restriction maps
on representation $K_0$, evaluated over $\mathbb C$ for Haar integration.
The tensor homomorphism is the integral group-scheme morphism
$GL_a\times GL_n\to GL_{an}$, and its kernel contains the central
$\mathbb G_m$ described in R-03. R-04 is the cycle-index expression
for symmetric powers in a $\lambda$-ring; the Adams operations preserve
the representation ring, not its cone of actual representations.
For R-05--R-06 the intrinsic algebraic object is the graded orbit-map
image, or the coefficient map in \eqref{R06:coefficients}. Its rank
strata are defined by determinantal (Fitting) ideals. The Gaussian
and Hermitian metrics are analytic devices for the same rank; they
are not asserted to be motives. The Schur idempotent is defined over
$\mathbb Z[1/D!]$, and the integral matrix in \eqref{R06:integer}
avoids treating characteristic-zero semisimplicity as an integral
direct-sum theorem. R-07 is a statement about restrictions of locally
free Schur constructions to the binary-form subscheme and the
scheme-theoretic image of the linear-substitution morphism. Its
formulas can be cleared of their binomial denominators over
$\mathbb Z$, while the basis and orbit-closure claims here are in
characteristic zero. R-08 is additivity in equivariant $K_0$ applied
to a resolution; the numerator remembers the quotient's equations.
Exact SGA/EGA/FGA and cote~133 locators for this dictionary are supplied
in the separate Grothendieck part of the paper.
R-09 identifies a four-factor invariant of
$\operatorname{Sym}(V_1\otimes V_2\otimes V_3\otimes V_4)$ with the
composition of two tensor-product restriction functors; its degree
$nd$ determinant-isotypic coefficient is the fourfold Kronecker
multiplicity. The finite CUE cutoff is the rank restriction on a
Schur functor, not an arithmetic purity hypothesis.

\subsection{R-20: a five-row certificate reaches the rectangular padding boundary}

\paragraph{Statement R-20.}
Put $M_0=8{,}000{,}001$. For every integer $m\geq M_0$, let
\[
 h_m=\left\lfloor\frac{m-97}{5}\right\rfloor,
 \qquad J_m=8h_m(m-317).
\]
Then
\begin{equation}\label{R20:dimension}
 \dim\mathcal P_{m,4}\geq2m^2+2+J_m.
\end{equation}
Consequently, for every pair of integers
\begin{equation}\label{R20:rectrange}
 m\geq M_0,\qquad m<n<4m/3,
\end{equation}
the explicit rectangular multiplicity gaps of G-11 hold, with its
labels $\lambda=(nd/4)^4$, degrees $d=E_nk$, and $k\geq K_n$.
Together with G-12, this makes $n<4m/3$ the exact eventual range,
among $n>m$, in which a rectangular multiplicity gap exists.
This is a statement about the specified four-row sector and its
large coordinate degrees.

More generally, whenever $m\geq M_0$, $n>m$, and
\begin{equation}\label{R20:allrange}
 2n^2\leq2m^2+J_m+2,
\end{equation}
R-14/G-10 give an exactly-four-row multiplicity gap in every degree
$d\geq d_n$. If also $n\geq4m/3$, every such gap is nonrectangular
by G-12. In particular, for every fixed $4/3<\rho<3/\sqrt5$, it includes
$n=\lfloor\rho m\rfloor$ for all sufficiently large $m$.
No maximal generic-rank formula is assumed.

\begin{proof}
We prove the dimension bound by an explicit differential certificate.

\paragraph{1. One simultaneous local construction.}
Let $q$ be a primitive $m$th root, to be chosen more carefully below,
and put
\[
 \ell_i=x-q^iy,\quad v_i=z-q^iw,\qquad
 L_c=x-cy,\quad V_c=z-cw.
\]
Start with the cyclic pencil of R-15. At a block starting at zero,
replace the five rows $0,\ldots,4$, in columns $0,\ldots,5$, by
\begin{equation}\label{R20:motif}
 \begin{pmatrix}
 \ell_0&v_0&\varepsilon L_2&\varepsilon L_3&\varepsilon V_3&0\\
 0&\ell_1&v_1&\varepsilon L_4&\varepsilon L_3&\varepsilon V_2\\
 0&0&\ell_2&v_2&\varepsilon V_6&0\\
 0&\varepsilon V_5&0&\ell_3&v_3&\varepsilon V_4\\
 0&\varepsilon L_2&\varepsilon V_7&\varepsilon L_5&\ell_4&v_4
 \end{pmatrix}.
\end{equation}
At starting positions $a=0,5,\ldots,5(h_m-1)$ use the translated
motif: shift all row and column indices by $a$ and replace
$(y,w)$ by $(q^ay,q^aw)$ in its entries. The five-row blocks are
disjoint. The four rows $a+1,\ldots,a+4$ have support only in
columns $a+1,\ldots,a+5$. The four columns
$a+1,\ldots,a+4$ have support only in rows $a,\ldots,a+4$.
Both assertions hold with every block present, since a neighboring
block shares only the boundary column. All remaining entries are
those of the original cycle.

Here is the local cofactor calculation used for both supports.
For a $4\times5$ matrix $B(\varepsilon)$ with diagonal entries
$a_i$, superdiagonal entries $b_i$, and other entries
$\varepsilon c_{ij}$, form $\widetilde B$ by multiplying the
off-bidiagonal entry $(i,j)$ by
\[
 \sigma_{ij}=\begin{cases}
 (-1)^{j-i-1},&j>i+1,\\
 (-1)^{i-j},&j<i.
 \end{cases}
\]
Write $P_j=\operatorname{per}(B\text{ with column }j\text{ deleted})$
and $D_j=\det(\widetilde B\text{ with column }j\text{ deleted})$.
There is one matching at order zero. A matching with one additional
edge $(i,j)$ has permutation sign $(-1)^{j-i-1}$ for $j>i+1$,
or $(-1)^{i-j}$ for $j<i$: the forced remaining bidiagonal edges
form the corresponding consecutive cycle. Thus
\begin{equation}\label{R20:secondorder}
 P_j-D_j=\varepsilon^2 E_j+O(\varepsilon^3),
 \qquad E_0=E_4=0.
\end{equation}
For completeness, the entire second-order expansion is
\begin{align*}
 E_1/2={}&c_{02}c_{10}b_2b_3
 +c_{03}(c_{10}a_2b_3+c_{20}b_1b_3)\\
 &+c_{04}(c_{10}a_2a_3+c_{20}b_1a_3+c_{30}b_1b_2),\\
 E_2/2={}&c_{13}(c_{21}a_0b_3+c_{20}b_0b_3)\\
 &+c_{14}(c_{21}a_0a_3+c_{31}a_0b_2
                   +c_{20}b_0a_3+c_{30}b_0b_2)\\
 &+c_{03}c_{20}a_1b_3
 +c_{04}(c_{20}a_1a_3+c_{30}a_1b_2),\\
 E_3/2={}&c_{24}(c_{32}a_0a_1+c_{31}a_0b_1+c_{30}b_0b_1)\\
 &+c_{14}(c_{31}a_0a_2+c_{30}b_0a_2)
 +c_{04}c_{30}a_1a_2.
\end{align*}
These identities can be verified directly by the $24$ permutations
of each $4\times4$ minor: the displayed $6,9,6$ terms are precisely
the two-additional-edge matchings whose signed determinant term is
negative. This specifies the finite calculation completely.

The determinant row relation gives
\begin{equation}\label{R20:relation}
 \sum_{j=0}^4(-1)^{i+j}\widetilde B_{ij}P_j
 =\varepsilon^2(a_iE_i-b_iE_{i+1})+O(\varepsilon^3).
\end{equation}
Apply it first to rows $1,\ldots,4$ and columns $1,\ldots,5$
of \eqref{R20:motif}, and then to the transpose of its submatrix
in rows $0,\ldots,4$ and columns $1,\ldots,4$.
In the global complementary permanents these identities have a
common remaining factor. For an external deleted row $s$, the
first set factors because the four supported rows must occupy the
four remaining supported columns. For an external deleted column
$c$, the second set factors because the four supported columns
must occupy the four remaining supported rows. The presence of
other blocks does not change either argument.

The coefficient of $\varepsilon^2$, divided by two, has the following
highest components in $(z,w)$ degree. This list follows by substituting
\eqref{R20:motif} into the preceding complete expansion:
\begin{align*}
 R_0={}&-v_1v_2(L_3V_5v_4+V_2V_5\ell_4+V_2L_2v_3),\\
 R_1={}&-v_1v_2v_4V_5V_6,\\
 R_2={}&\ell_3v_1v_4V_5V_6
          -v_2v_3V_4(\ell_1V_7+L_2v_1),\\
 R_3={}&\ell_4v_2V_4(\ell_1V_7+L_2v_1),\\[2pt]
 S_0={}&-\ell_1L_2v_2(\ell_4V_5+L_2v_3),\\
 S_1={}&v_1L_2v_2(\ell_4V_5+L_2v_3)
          -\ell_2L_4v_0v_3V_7,\\
 S_2={}&L_4v_0v_2v_3V_7,\\
 S_3={}&v_3(L_5V_6v_0v_1+L_3V_7v_0v_2
                +\ell_1V_7V_3v_2+L_2V_3v_1v_2).
\end{align*}
Their $(z,w)$ degrees are respectively
\begin{equation}\label{R20:grades}
 (4,5,4,3)\quad\hbox{and}\quad(2,3,4,4).
\end{equation}
Every listed polynomial has total degree five. The complete
coefficient in \eqref{R20:relation} may have additional components
of strictly smaller $(z,w)$ degree; those will be handled below,
without projecting an actual differential image by an unjustified
operation.

At $\varepsilon=0$ the common outside factors are exactly
\begin{equation}\label{R20:outside}
 T_s^{\rm row}=\ell_0\prod_{i=5}^{s-1}v_i
                         \prod_{i=s+1}^{m-1}\ell_i,
 \qquad
 T_c^{\rm col}=\prod_{i=5}^{c-1}\ell_i
                         \prod_{i=c}^{m-1}v_i.
\end{equation}
This follows either by the unique cycle matching or by successively
matching the forced boundary column. If a local polynomial has
$(z,w)$ degree $g$, choose
\begin{equation}\label{R20:alignment}
 s=k+5-g\quad\hbox{for a row family},\qquad
 c=m+g-k\quad\hbox{for a column family}.
\end{equation}
Its highest component times \eqref{R20:outside} then has bidegree
$(m-k,k)$. All deleted indices are external and distinct as required
when $9\leq k\leq m-9$.

\paragraph{2. Fixed rational-function matrices controlling all modes.}
Write $K=q^k$, $X=q^t$. The following construction defines two
$6\times6$ matrices over $\mathbb Q(q,K,X)$, using only the eight
displayed degree-five polynomials. It therefore does not depend on
a numerical extrapolation in $m$.

Let $F$ be one such polynomial of bidegree $(a,g)$, $a+g=5$.
For a row polynomial, express it uniquely as
\begin{equation}\label{R20:rowbasis}
 F=\sum_{j=1}^{a+1} Q^{\rm row}_{a,j}(x,y) C_j(z,w),\qquad
 Q^{\rm row}_{a,j}=
 \prod_{u=j}^{a}(x-Kq^uy)\prod_{u=1}^{j-1}(x-q^uy).
\end{equation}
Define the binary quintic
\begin{equation}\label{R20:rowtransform}
 \mathcal R(F)=\sum_{j=1}^{a+1}X^j C_j(z,q^{-j}w)
       \prod_{u=5-j}^{3}(z-q^uw)
       \prod_{u=-1}^{a-1-j}(z-Kq^uw).
\end{equation}
As everywhere here, a product with lower endpoint greater than its
upper endpoint is one. For a column polynomial use
\begin{equation}\label{R20:colbasis}
 \begin{split}
 F&=\sum_{j=0}^{a}Q^{\rm col}_{a,j}(x,y)C_j(z,w),\\
 Q^{\rm col}_{a,j}&=
 \prod_{u=5-j}^{4}(x-q^uy)
 \prod_{u=0}^{a-1-j}(x-K^{-1}q^{5-a+u}y),
 \end{split}
\end{equation}
and define
\begin{equation}\label{R20:coltransform}
 \mathcal S(F)=\sum_{j=0}^{a}X^{-j}C_j(z,Kq^{j-5}w)
       \prod_{u=-a+j}^{-1}(z-q^uw)
       \prod_{u=0}^{j-1}(z-Kq^{-5+u}w).
\end{equation}
The coefficient matrices of the $Q$ bases are generically invertible.
Their nonzero determinants are also included in the certificate below,
so no problematic specialization is silently inverted.

The row matrix consists of coefficient columns, in the order
$z^5,z^4w,\ldots,w^5$, of
\[
 z\prod_{u=0}^3v_u,\quad w\prod_{u=0}^3v_u,
 \quad\mathcal R(R_0),\ldots,\mathcal R(R_3).
\]
The column matrix consists of the corresponding columns of
\[
 z\prod_{u=-5}^{-2}(z-Kq^uw),\quad
 w\prod_{u=-5}^{-2}(z-Kq^uw),
 \quad\mathcal S(S_0),\ldots,\mathcal S(S_3).
\]
At $(q,K,X)=(6,-1,-1)$ over $\mathbb F_{41}$, these matrices are
\begin{equation}\label{R20:matrices}
 \begin{split}
 M_R&=\begin{pmatrix}
 1&0&36&1&21&34\\28&1&24&28&24&4\\
 34&28&3&28&35&22\\21&34&34&6&27&27\\
 39&21&0&33&24&34\\0&39&26&27&22&6
 \end{pmatrix},\\
 M_C&=\begin{pmatrix}
 1&0&30&28&31&1\\2&1&33&10&24&8\\
 19&2&1&9&13&0\\34&19&35&4&27&5\\
 2&34&13&36&15&4\\0&2&32&3&11&27
 \end{pmatrix}.
 \end{split}
\end{equation}
Their determinants are respectively $16$ and $15$ modulo $41$.
The four row interpolation determinants are $29,1,29,38$;
the four column interpolation determinants are $25,6,32,32$.
All entries are obtained by the products and linear systems
\eqref{R20:rowbasis}--\eqref{R20:coltransform}, which specify a
small reproducible exact certificate.

We record explicit degree bounds in order to turn this certificate
into a uniform statement. Let $d_R$ be the determinant of a row
coefficient matrix in \eqref{R20:rowbasis}, and let $d_C$ be the
determinant of the coefficient matrix of
$K^a Q^{\rm col}_{a,j}$ in \eqref{R20:colbasis}.
For all eight forms, $0\leq a\leq3$. Cramer's rule shows that
the coefficients of
\[
 q^{60}X^5d_R\mathcal R(F),\qquad
 q^{60}X^5d_C\mathcal S(F)
\]
are polynomials in $q,K,X$. Each has degree at most
$150,20,10$ in those three variables separately. Here are direct
bounds verifying the clearing and the stated estimates.
Every coefficient of a local form has $q$-degree at most $20$.
A row interpolation entry has $q$-degree at most $a(a+1)/2$
and $K$-degree at most $a$; its adjugate entry has degrees at most
$18,9$. The substitutions and additional products in
\eqref{R20:rowtransform} have negative $q$ exponent no smaller
than $-10$ and positive contribution at most $7$; after clearing,
the $q$-degree is at most $105$, the $K$-degree at most $12$,
and the $X$-degree at most $9$.
For the column basis, multiplying by $K^a$ gives entries of degrees
at most $4a,a$. Its adjugate has degrees at most $36,9$.
The solution acquires an additional factor $K^a$, and
\eqref{R20:coltransform} has negative $q$ exponent no smaller than
$-31$. After clearing its degrees are at most $116,17,5$.
These are bounded by $150,20,10$ as asserted.

Multiply the first two columns of each matrix by $q^{25}$,
and clear each of its four other columns as just described.
The two base columns now have degrees at most $31,4,0$.
Let $\Delta_R,\Delta_C$ be the two resulting polynomial determinants,
and define
\[
 \mathscr F(q,K,X)=\Delta_R\Delta_C
                  \prod_{\text{four row forms}}d_R
                  \prod_{\text{four column forms}}d_C.
\]
Each interpolation determinant has degrees at most $48,12,0$.
Consequently
\begin{equation}\label{R20:degrees}
 \deg_q\mathscr F\leq2000,\qquad
 \deg_K\mathscr F\leq300,\qquad
 \deg_X\mathscr F\leq80.
\end{equation}
For example the sharper bounds from the preceding estimates are
$1708,272,80$. At the certificate specialization all eight
interpolation determinants have product $20$ modulo $41$;
$K^{a(a+1)}=1$ there, so the column clearing does not change their
listed values. Thus
\[
 \mathscr F(6,-1,-1)
 =6^{580}(-1)^{40}\,20^2\,16\,15=22\pmod{41}.
\]
In particular $\mathscr F$ is a nonzero integer polynomial.

\paragraph{3. Why these matrices are the required Fourier quotients.}
Use the cyclic action $\tau:(y,w)\mapsto(qy,qw)$ of R-16 and
write $A(x,y)=\prod_{u=k}^{m-1}\ell_u$.
All its coefficients $A_r$, $0\leq r\leq m-k$, are nonzero by
the geometric-product formula already proved there. In character
$t$, divide the coefficient with $y$-degree $r$ and $w$-degree
$s$ by $A_r$; here $r+s=t$.

For a row form of $(x,y)$ degree $a$, the outside factor is
$\ell_0\prod_{u=k+a+1}^{m-1}\ell_u$ in those variables.
Multiplication by each $Q^{\rm row}_{a,j}$ in
\eqref{R20:rowbasis} completes exactly $A(x,q^jy)$.
Dividing its coefficient by $A_r$ gives $q^{jr}=X^jq^{-js}$.
The remaining $(z,w)$ factors are
$C_j(z,q^{-j}w)\prod_{u=5-j}^{k+a-1-j}v_u$.
Their common product is
\[
 H_R=\prod_{u=4}^{k-2}v_u,
\]
and the residual quintic is precisely
\eqref{R20:rowtransform}.

For a column form, the outside $(x,y)$ factor is
$\prod_{u=5}^{c-1}\ell_u$, with $c=m+5-a-k$.
The $j$th basis term in \eqref{R20:colbasis} completes
$A(x,q^{5-k-j}y)$. After coefficient division and multiplication
by the nonzero character scalar $X^{k-5}$, its factor is
$X^{-j}q^{(k-5+j)s}$.
Shifting the outside $(z,w)$ roots $c,\ldots,m-1$ by
$k-5+j$ gives indices $-a+j,\ldots,k-6+j$ modulo $m$.
Their common product is
\[
 H_C=\prod_{u=0}^{k-6}v_u,
\]
and the residual quintic is exactly
\eqref{R20:coltransform}. This checks both shifts and both
normalizations, including their dependence on $k,t$.

Suppose $9\leq k,t\leq m-9$ and
$\mathscr F(q,q^k,q^t)\ne0$. The two matrices then show that
base plus row columns span $H_R\operatorname{Sym}^5\mathbb C^2$,
and base plus column columns span
$H_C\operatorname{Sym}^5\mathbb C^2$ before restriction to the
allowed coefficient interval. Their common factor is
$H=\prod_{u=4}^{k-6}v_u$. After dividing by $H$, the first space is
\[
 \left(\prod_{u=k-5}^{k-2}v_u\right)
                 \operatorname{Sym}^5\mathbb C^2,
\]
and the second is
\[
 \left(\prod_{u=0}^{3}v_u\right)
                 \operatorname{Sym}^5\mathbb C^2.
\]
The two displayed degree-four factors are coprime: their indices
are distinct and lie below $m$. Their intersection consists exactly
of their product times $\operatorname{Sym}^1\mathbb C^2$,
which is the two-dimensional base space. The joint span is therefore
all $H\operatorname{Sym}^9\mathbb C^2$, of dimension ten.
This proves joint independence; separate row and column ranks alone
would not suffice.

There are ten consecutive allowed coefficient indices
$b,\ldots,b+9$ with $0\leq b\leq k-9$.
Multiplication by $H$ and restriction to them has determinant
\[
 \det(H_{b+u-v})_{0\leq u,v\leq9}
 =\pm s_{(10^b)}(q^4,q^5,\ldots,q^{k-6})\ne0.
\]
This is the same dual Jacobi--Trudi argument as R-16/R-19
(the cited unconditional identity R-B15). There are $N=k-9$
variables; the numerator exponents are distinct in $[0,N+9]=[0,k]$,
strictly below $m$, so the Weyl alternant is a nonzero Vandermonde.
For $N=0$ the coefficient matrix is the identity.
Consequently the eight normal classes remain independent modulo
the entire base character space after the required coefficient
restriction.

\paragraph{4. Only finitely many modes are lost.}
Choose a nonzero coefficient polynomial in $q$ of the nonzero
polynomial $\mathscr F(q,K,X)$. Its degree is at most $2000$.
For every positive integer $m$,
\[
 \varphi(m)^2/m
 =\prod_{p^e\Vert m}p^{e-2}(p-1)^2\geq\tfrac12:
\]
each odd-prime factor is at least one, and the possible factor
from $p=2$ is at least $1/2$.
For $m\geq M_0$ there are therefore more than $2000$ primitive
$m$th roots. Choose one at which that coefficient polynomial does
not vanish. This uses ordinary root counting and gives
$\mathscr F(q,K,X)\not\equiv0$; it does not incorrectly replace
$\varphi(m)$ by $m$.

For this $q$, at most $300$ of the distinct values $K=q^k$ can
make $\mathscr F(q,K,X)$ the zero polynomial in $X$: choose one
nonzero coefficient polynomial in $K$ and use
\eqref{R20:degrees}. Among $k=9,\ldots,m-9$ there are consequently
at least $m-317$ good bidegrees. For each good $k$, at most $80$
of the distinct values $X=q^t$ are zeros. Thus at least $m-97$
characters in $t=9,\ldots,m-9$ have all eight independent classes.

Consider translates of the eight leading vectors at the $h_m$
block starts $a=5j$, $0\leq j<h_m$. A relation in a good bidegree,
projected onto a good character, gives eight separate equations
\[
 \sum_{j=0}^{h_m-1}c_j(q^{5t})^j=0.
\]
The good characters contain at least $h_m$ distinct values of
$q^{5t}$: every value has at most $\gcd(5,m)\leq5$ preimages,
and $h_m=\lfloor(m-97)/5\rfloor$.
Each equation system contains an ordinary invertible Vandermonde
matrix. Hence the $8h_m$ leading normal classes in each good
bidegree are independent modulo $W_0$.

The original polynomial limits need not be homogeneous for this
bigrading. Nevertheless they are independent modulo $W_0$ as well.
In any relation, take the largest $k$ among vectors with nonzero
coefficients. Their components of bidegree $(m-k,k)$ are precisely
the independent leading classes just proved; every other component
has smaller $(z,w)$ degree. Because $W_0$ is bigraded, its contribution
can be removed in that quotient, forcing all coefficients at this
$k$ to vanish. Descending in $k$ proves independence of every
selected vector. This justifies the leading-component step without
assuming that a projection of the deformed differential image lies
in that image.

Finally, the eight sets of polynomial families all belong, for
nonzero $\varepsilon$, to the differential image of the same
simultaneous pencil. Indeed \eqref{R20:relation} times its outside
permanent is a linear-form combination of global complementary
permanents; division by $\varepsilon^2$ is polynomial by
\eqref{R20:secondorder}. Their limits and $2m^2+2$ base columns
have a nonzero augmented minor at zero. This minor stays nonzero
for a common generic nonzero $\varepsilon$. Thus the dimension
bound \eqref{R20:dimension} follows from the actual differential
rank, with no addition of ranks from different pencils.

\paragraph{5. The uniform padding consequence.}
The general criterion \eqref{R20:allrange} is R-14, with G-10
giving its effective coordinate-degree threshold. Moreover
\[
 h_m\geq(m-101)/5,
\]
so
\begin{align*}
 2m^2+J_m+2-\frac{32m^2}{9}
 &\geq \frac{2m^2}{45}-\frac{3344m}{5}
                  +\frac{256136}{5}+2\;>0
 \qquad(m\geq M_0).
\end{align*}
The last inequality already holds for $m>15048$ after dropping
the positive constant. Thus every integer $m<n<4m/3$ satisfies
\eqref{R20:allrange}. Its strict padding inequality is exactly
the stable range required by G-11, yielding the stated explicit
rectangles. G-12 excludes rectangular gaps for every $n\geq4m/3$.
Finally $J_m=(8/5)m^2+O(m)$, so the full dimension criterion reaches
every fixed ratio below $\sqrt{9/5}=3/\sqrt5$; above $4/3$, the
resulting four-row gaps must be nonrectangular.
\end{proof}


\subsection{R-21: arithmetic rank recovery and a fixed-dimension algorithm}

\paragraph{Statement R-21.}
The integer Gram model admits an explicit spectral separation from
zero. Moreover, for each fixed ambient dimension $N$, every orbit-closure
multiplicity $m_\lambda(f)$ of a rational degree-$n$ form can be computed
exactly in time polynomial in the numerical degree $D=nd$ and the
coefficient bit length. This is an elementary consequence of a compressed
coefficient kernel and weight extraction. The exponent of the polynomial
depends on $N$; this assertion is not polynomial in $\log d$ and does not
give a uniform polynomial-time algorithm when $N$ grows.

\paragraph{1. An integral spectral bound.}
Let $A$ be an $a\times b$ integer matrix of positive rank $r$, with
entries of absolute value at most $H$, and put $G=AA^{\mathsf T}$,
$T=\operatorname{tr}G\leq abH^2$. Then
\begin{equation}\label{R21:gap}
 \lambda_{\min}^{+}(G)\geq
 \begin{cases}
 1,&r=1,\\[1mm]
 \bigl((r-1)/T\bigr)^{r-1},&r>1.
 \end{cases}
\end{equation}
More generally the same bound holds for any nonzero positive semidefinite
integer symmetric matrix, with its own rank and trace.

\begin{proof}
If the positive eigenvalues are $\eta_1\leq\cdots\leq\eta_r$, their
product is the elementary symmetric coefficient $e_r(G)$, a positive
integer. For a Gram matrix it is also a sum of squared integer minors
by Cauchy--Binet. Thus $\prod_i\eta_i\geq1$.
If $r=1$ this proves the claim. Otherwise the arithmetic--geometric
mean inequality gives
\[
 \prod_{i=2}^r\eta_i\leq
 \left(\frac{\sum_{i=2}^r\eta_i}{r-1}\right)^{r-1}
 \leq\left(\frac{T}{r-1}\right)^{r-1},
\]
which proves \eqref{R21:gap}. The proof for an integer positive
semidefinite matrix uses the same integer characteristic coefficient.
\end{proof}

The exponential dependence on matrix dimension can occur even at fixed
entry height. If $S$ is the nilpotent superdiagonal shift of size $r$
and $A=I-HS$ with integer $H\geq2$, then
$A^{-1}=\sum_{j=0}^{r-1}H^jS^j$. Its top-right entry is $H^{r-1}$,
so $\|A^{-1}\|_2\geq H^{r-1}$ and
$\lambda_{\min}(AA^{\mathsf T})\leq H^{-2(r-1)}$.
This is a concrete limitation on dimension-independent spectral
separation from integrality alone.

For example, let $A_\lambda$ be the integral matrix in R-06 constructed
from the tensor $n!f$, where the original polynomial has integer monomial
coefficients. Set $c_\lambda=f^\lambda/D!$. Its projected Gaussian
covariance for the original, unscaled symmetric tensor is
\begin{equation}\label{R21:gaussiantransport}
 \mathcal G_\lambda(f)=
 \frac{c_\lambda^2}{(n!)^{2d}}
 A_\lambda\operatorname{diag}(\alpha!)A_\lambda^{\mathsf T}.
\end{equation}
When $A_\lambda\ne0$, every $\alpha!\geq1$, so its least positive eigenvalue is at least
$c_\lambda^2/(n!)^{2d}$ times the bound in \eqref{R21:gap} for
$A_\lambda A_\lambda^{\mathsf T}$. Both matrices have the same image,
so the comparison is made on their common support. The factorials in
this statement are necessary; the passage from an integral polynomial
to an integral symmetric tensor cannot be omitted.
The explicit matrix dimensions and height in R-06 show that this
bound alone does not yield polynomial precision in growing GCT size.

\paragraph{2. A compressed integral coefficient kernel.}
Assume $f=\sum_{|\beta|=n}f_\beta x^\beta$ has integer coefficients;
rational inputs can be cleared of a common denominator without changing
the orbit span. Let $f_{\rm ten}\in(\mathbb C^N)^{\otimes n}$ be the
usual symmetric-tensor embedding and set
$v=(n!f_{\rm ten})^{\otimes d}$.
For matrices of variables $Z,W$, write
\[
 Z\cdot f=\sum_{|\beta|=n}g_\beta(Z)x^\beta,
 \qquad
 \Phi(Z)=Z^{\otimes D}v=\sum_{|\alpha|=D}Z^\alpha c_\alpha.
\]
There are
\[
 M=\binom{N^2+D-1}{D}
\]
column labels $\alpha$. Form the polynomial
\begin{equation}\label{R21:kernel}
 \mathcal K_f(Z,W)=
 \left(n!\sum_{|\beta|=n}\beta!\,
                     g_\beta(Z)g_\beta(W)\right)^d
       =\sum_{\alpha,\gamma}B_{\alpha\gamma}Z^\alpha W^\gamma.
\end{equation}
Then $B=C_f^{\mathsf T}C_f$ is an $M\times M$ positive semidefinite
integer matrix. Thus its kernel coefficients replace the exponentially
tall tensor-coordinate matrix without constructing its $N^D$ rows.

\begin{proof}
For a composition $\beta$ of $n$, the symmetric-tensor embedding of
the monomial $x^\beta$ is
$(\beta!/n!)$ times the sum of all distinct tensor words with
composition $\beta$. There are $n!/\beta!$ such words.
Consequently the coefficient of each such word in
$Z^{\otimes n}(n!f_{\rm ten})$ is $\beta!g_\beta(Z)$, and its
bilinear tensor inner product with the analogous expression in $W$ is
\[
 \sum_{|\beta|=n}\frac{n!}{\beta!}
                 (\beta!g_\beta(Z))(\beta!g_\beta(W))
 =n!\sum_{|\beta|=n}\beta!g_\beta(Z)g_\beta(W).
\]
Taking $d$ tensor powers gives \eqref{R21:kernel}. Comparing coefficients
gives $B_{\alpha\gamma}=c_\alpha^{\mathsf T}c_\gamma$.
The vectors are real and integral; the resulting real Gram matrix is
also positive semidefinite for the corresponding Hermitian inner product.
\end{proof}

For $\alpha=(\alpha_{ij})$, define its row-sum weight by
\[
 \operatorname{wt}(\alpha)_i=\sum_{j=1}^N\alpha_{ij}.
\]
The matrix $B$ is block diagonal by these weights. Denote the block
of weight $\mu$ by $B_\mu$ and set $r_\mu=\operatorname{rank}B_\mu$,
with $r_\mu=0$ if a coordinate of $\mu$ is negative or the block is
absent. For $\delta=(N-1,N-2,\ldots,0)$, the exact formula is
\begin{equation}\label{R21:weightformula}
 m_\lambda(f)=
 \sum_{\sigma\in S_N}\operatorname{sgn}(\sigma)
              r_{\lambda+\delta-\sigma\delta},
 \qquad \lambda\vdash D,\quad\ell(\lambda)\leq N.
\end{equation}
This is a signed Weyl extraction from positive Gram ranks, not a new
cancellation-free tableau rule.

\begin{proof}
For a diagonal matrix $t$, the identity
$\Phi(tZ)=t^{\otimes D}\Phi(Z)$ shows that $c_\alpha$ has weight
$\operatorname{wt}(\alpha)$. Distinct tensor-weight spaces are
orthogonal in the word basis, proving the block decomposition.
As in R-06, the coefficient span is exactly the $GL_N$ orbit span:
a functional vanishes on every coefficient if and only if its
polynomial evaluation vanishes on all matrices, equivalently on the
Zariski-dense invertible matrices. Hence $r_\mu$ is exactly the
dimension of the $\mu$ weight space of this orbit span.

Its character is $\sum_\mu r_\mu z^\mu$.
Multiplication by the Weyl alternant
$\sum_{\sigma\in S_N}\operatorname{sgn}(\sigma)z^{\sigma\delta}$
turns $s_\nu(z)$ into the alternant with exponent $\nu+\delta$.
The coefficient of $z^{\lambda+\delta}$ is therefore $m_\lambda(f)$:
the strictly decreasing dominant exponent occurs in precisely that
alternant, with coefficient one. Expanding this coefficient proves
\eqref{R21:weightformula}. The orbit-span multiplicities are the
coordinate-ring multiplicities under the consistent dual convention
already fixed in R-05.
\end{proof}

\paragraph{3. Exact cost and sufficient approximation precision.}
Let $H=\max\{1,|f_\beta|\}$ and put
\[
 H_B=N^{3D}(n!H)^{2d},\qquad T_0=MH_B.
\]
Every entry of $B$ has absolute value at most $H_B$.
Indeed each tensor coefficient has absolute value at most
$N^D(n!H)^d$, and each Gram entry sums $N^D$ products of two such
coefficients. Thus $\operatorname{tr}B_\mu\leq T_0$ for every block.
The spectral lemma, without knowing the block rank, gives
\begin{equation}\label{R21:uniformgap}
 \lambda_{\min}^{+}(B_\mu)\geq
 \eta_0:=T_0^{-(M-1)}
\end{equation}
for every nonzero block. To see this, use rank at most $M$,
$T_0\geq1$, and discard the factor $(r-1)^{r-1}\geq1$ from
\eqref{R21:gap}; a rank-one positive eigenvalue is at least one.

If a Hermitian approximation to each block has entrywise error
less than $\eta_0/(4M)$, its operator error is less than
$\eta_0/4$. The min--max principle then places each eigenvalue
within that distance of its exact value. Counting the approximate
eigenvalues above $\eta_0/2$ recovers the exact rank. For example,
entry precision of
\begin{equation}\label{R21:bits}
 p=\left\lceil (M-1)\log_2T_0+\log_2(8M)\right\rceil
\end{equation}
absolute entry error at most $2^{-p}$ suffices, with room for an
additional eigenvalue-approximation error at most $\eta_0/8$. For this integral coefficient matrix there is also a simpler exact
recovery route: absolute entry error below $1/2$ permits rounding each
entry to its exact integer, followed by exact elimination. The smaller
error above is needed only for the stated direct spectral threshold;
it must not be confused with a necessary precision for integer recovery.
The spectral statement is an explicit conditioning bound; it does not by itself bound the cost of a
random-matrix integration oracle.

For fixed $N$, $M$ is polynomial in $D$, and
\[
 \log H_B=O\bigl(D\log N+D\log n+d\log H\bigr).
\]
The $g_\beta$ and the kernel in \eqref{R21:kernel} can be expanded
by integer polynomial arithmetic in a fixed number of variables.
Their dense coefficient arrays, at degrees at most $D$ in each
matrix alphabet, have polynomial size in $D$.
Their coefficient bit lengths are bounded by the same elementary
factorial and multinomial estimates. Thus all kernel entries can
be computed in polynomial bit complexity for fixed $N$.
Exact rational elimination computes each required block rank in
polynomial bit complexity as well: intermediate reduced fractions
are ratios of minors, whose sizes are bounded by the factorial
expansion bound $M!H_B^M$. Formula \eqref{R21:weightformula} needs
at most $N!$ block ranks for a given label. This proves the claimed
fixed-dimension algorithm; \eqref{R21:bits} is polynomial in the
same numerical parameters.

There are two relevant limits to this conclusion. In GCT the ambient
dimension is $N=n^2$, so neither the size exponent nor the precision
bound is uniform polynomial in matrix size. Also, an orbit of one
fixed quaternary form is not the entire image of all four-variable
permanent pencils. Applying this algorithm to one specialization
does not compute the multiplicities of $\mathcal P_{m,4}$.

\paragraph{The slope-method gate.}
The supplied Alp\"oge note, step~13 on printed p.~4 (G-B10), names
the Bost slope inequality as a step. That step itself gives no integral-metric construction or
degree-dependent constants for the coefficient matrices above. A transfer to this rank problem would require such
a construction and bounds strong enough to control the displayed
dimensions and heights. The present spectral estimate supplies a
self-contained arithmetic bound and an exact restricted algorithm;
it neither uses nor challenges the stipulated RH/Dirichlet-GRH input.
\subsection{R-22: arbitrary staircase size reaches the four-variable dimension ceiling}

\paragraph{Statement R-22.}
Let $b\geq3$ and $m\geq8b$ be integers, and define
\[
 T=m-4b+3,\qquad h=\left\lfloor\frac{m-6b+3}{b}\right\rfloor,
 \qquad J_{m,b}=(2b-2)hT.
\]
Then the quaternary permanent projection variety satisfies
\begin{equation}\label{R22:rank}
 c_m:=\dim\mathcal P_{m,4}\ \geq\ 2m^2+2+J_{m,b}.
\end{equation}
The proof uses one simultaneous deformation of the cyclic pencil.
Its untwisted motif parameters lie in the integer grid
$\{0,\ldots,4(b-1)T\}$; the translated copies have the prescribed
root-of-unity twists. A grid choice is allowed to depend on $m$ and on
the chosen primitive $m$th root of unity. No ranks from
different pencils are added.

In particular, for every $m\geq99$,
\begin{equation}\label{R22:asymptotic}
 4m^2-20m^{3/2}+2\ \leq\ c_m\ \leq\ 4m^2-2m+2.
\end{equation}
Thus $c_m=4m^2-O(m^{3/2})$. The upper bound is the row--column torus
bound proved in R-13, and an equality formula for $c_m$ is not asserted.

For any integers $n>m$ such that
\begin{equation}\label{R22:padding}
 2n^2\leq2m^2+J_{m,b}+2,
\end{equation}
the padded image has dimension greater than the determinant image.
Consequently R-14/G-10 give exactly-four-row multiplicity gaps in every
coordinate degree beyond the explicit G-10 threshold. The specified
stretched labels and effective bounds of G-17 also apply.
For every fixed real $1<\rho<\sqrt2$, this includes every integer
$m<n\leq\rho m$ as soon as
\begin{equation}\label{R22:ratio_threshold}
 m\geq\max\left\{99,
       \left\lceil\frac{100}{(2-\rho^2)^2}\right\rceil\right\}.
\end{equation}
In particular the eventual rectangular boundary of R-20 holds already
for every $m\geq2025$: rectangular gaps exist for every integer
$m<n<4m/3$, and G-12 excludes them for $n\geq4m/3$.
The additional range above $4m/3$ consists of nonrectangular gaps.

\begin{proof}
We give a uniform local calculation, prove its two Fourier rank
conditions, and then arrange that one integer choice works for all
interior bidegrees of the given matrix size.

\paragraph{1. The typed staircase and its simultaneous supports.}
Fix any primitive $m$th root $q$, and write
\[
 \ell_i=x-q^iy,\qquad v_i=z-q^iw.
\]
The cyclic pencil has entries $\ell_i$ at $(i,i)$ and $v_i$ at
$(i,i+1)$, with indices modulo $m$.
A block of $b$ rows, initially rows $0,\ldots,b-1$, has columns
$0,\ldots,b$. In addition to its diagonal and superdiagonal, allow
every entry $(i,j)$ except entries in column zero and the entry
$(0,b)$. Thus the permitted additional positions are
\[
 0\leq i<b,\quad1\leq j\leq b,\quad
 j\ne i,i+1,\quad(i,j)\ne(0,b).
\]
There are $b(b-2)$ such positions. Set their entries to
$\varepsilon C_{ij}$, where
\begin{equation}\label{R22:types}
 C_{ij}=\begin{cases}
 u_{ij}x+w_{ij}y,&j>i+1,\\
 u_{ij}z+w_{ij}w,&j<i.
 \end{cases}
\end{equation}
The symbol $w_{ij}$ denotes a scalar coefficient, distinct from the
coordinate $w$. All $2b(b-2)$ coefficients are initially independent
parameters, collectively denoted $\theta$.

Place copies of this same block at starts $a=bj$, $0\leq j<h$,
replacing $(y,w)$ in its additional forms by $(q^ay,q^aw)$.
All block rows are disjoint, since $bh\leq m-6b+3<m$.
For each block, its bottom $b-1$ rows have support only in its
$b$ columns numbered $1,\ldots,b$ relative to the start.
Its middle $b-1$ columns have support only in its $b$ rows numbered
$0,\ldots,b-1$. A neighboring block shares only a boundary column,
so both support assertions remain true with all blocks present.

These supports give exact cofactor factorizations. If an external row
is deleted, and one of the $b$ supported columns is deleted, all
bottom $b-1$ supported rows must occupy the remaining $b-1$ supported
columns in every matching. The complementary permanent is the product
of their local permanent minor and a common outside permanent.
The transposed assertion holds when an external column and one of
the $b$ supported rows are deleted. Other blocks can affect the outside
factor but cannot alter either factorization.

\paragraph{2. The second-order local defect for arbitrary size.}
Put $r=b-1$. Consider an $r\times(r+1)$ matrix $B(\varepsilon)$
with diagonal entries $a_i$, superdiagonal entries $b_i$, and
other entries $\varepsilon c_{ij}$, with $0\leq i<r$.
Form $\widetilde B$ by signing its additional entries with
\[
 \sigma_{ij}=\begin{cases}
 (-1)^{j-i-1},&j>i+1,\\
 (-1)^{i-j},&j<i.
 \end{cases}
\]
Let $P_j$ be the permanent minor deleting column $j$, and $D_j$
the corresponding determinant minor of $\widetilde B$.
There is one order-zero matching. A matching with one additional
edge has its other edges forced along the bidiagonal path; its
permutation sign is precisely the indicated $\sigma_{ij}$.
Hence
\begin{equation}\label{R22:local_defect}
 P_j-D_j=\varepsilon^2E_j+O(\varepsilon^3).
\end{equation}

We need the endpoint identities $E_0=E_r=0$ for all $r$.
For the minor deleting column $r$, a matching with no additional
forward edge is an upper-Hessenberg permutation. Each nontrivial
cycle is a consecutive interval, whose superdiagonal edges end in
one backward edge. Its permutation sign equals the sign assigned
to that backward edge. Two forward additional edges and no backward
edge cannot form a permutation. If there is exactly one forward
edge $i\mapsto j$, $j>i+1$, and one backward edge $h\mapsto l$,
there is a single nontrivial cycle, necessarily with
$l\leq i<j\leq h$. It follows superdiagonal edges from $l$ to $i$,
jumps to $j$, follows them to $h$, and returns to $l$.
Its length is $(i-l)+(h-j)+2$. The difference between its sign
exponent and the sum of the two assigned exponents is even, since
\[
 (j-i-1)+(h-l)-\big((i-l)+(h-j)+1\big)=2(j-i-1).
\]
Thus the endpoint minor agrees through order two.
Deleting column zero and transposing gives the same square
bidiagonal situation, with diagonal $b_i$ and superdiagonal
$a_{i+1}$. Under this operation a forward edge becomes a backward
edge with the same assigned sign, and conversely. This proves
\begin{equation}\label{R22:endpoints}
 E_0=E_r=0.
\end{equation}

Assume now that $a_i$ and the additional forward entries are
linear in $(x,y)$, while $b_i$ and the additional backward entries
are linear in $(z,w)$. The highest $(z,w)$ degree of $E_j$ is
$r-1$. Indeed, if there is no additional forward edge, the same
upper-Hessenberg argument gives equality of permanent and signed
determinant for every minor: the rows before the deleted column
form the upper-Hessenberg block, and all later superdiagonal
entries are forced. Therefore terms with no $(x,y)$ factor, and
terms with a diagonal $(x,y)$ factor and two backward edges, cancel.
At $(x,y)$ degree one the only remaining possibility uses one
forward edge, one backward edge, and otherwise superdiagonal edges.

For the minor deleting column $j$, compare its target columns with
the all-superdiagonal target set $\{1,\ldots,r\}$.
The target zero must come from the backward edge. Comparing the
two removed targets with the two added targets then forces the
additional edges to be
\[
 (j-1,h+1),\qquad(h,0),\qquad j\leq h\leq r-1.
\]
The minor permutation has the cycle
$(0,1,\ldots,j-1,h)$, of sign $(-1)^j$; the product of assigned
edge signs has opposite sign. Consequently
\begin{equation}\label{R22:Etop}
 [E_j]_{\deg_{z,w}=r-1}
 =2\sum_{h=j}^{r-1}c_{j-1,h+1}c_{h,0}
           \prod_{\substack{0\leq t<r\\t\ne j-1,h}}b_t
 \quad(1\leq j\leq r-1).
\end{equation}
This proves the formula for every block size, without extrapolating
a finite list of local minors.

The determinant row relation yields
\begin{equation}\label{R22:relation}
 \sum_{j=0}^{r}(-1)^{i+j}\widetilde B_{ij}P_j
 =\varepsilon^2(a_iE_i-b_iE_{i+1})+O(\varepsilon^3).
\end{equation}
For $i<r-1$, the highest component, divided by two, is
\[
 -\sum_{h=i+1}^{r-1}c_{i,h+1}c_{h,0}
                          \prod_{t\ne h}b_t,
\]
of bidegree $(1,r)$. For $i=r-1$ it is
\[
 a_{r-1}c_{r-2,r}c_{r-1,0}\prod_{t=0}^{r-3}b_t,
\]
of bidegree $(2,r-1)$.

In the original $b$-row block, denote these row forms by
$R_0,\ldots,R_{r-1}$. Explicitly,
\begin{align}
 R_i&=-\sum_{h=i+1}^{r-1}C_{i+1,h+2}C_{h+1,1}
                 \prod_{\substack{0\leq t<r\\t\ne h}}v_{t+1}
                     &&(0\leq i<r-1),\label{R22:Rforms}\\
 R_{r-1}&=\ell_r C_{r-1,r+1}C_{r,1}
                                      \prod_{t=1}^{r-2}v_t.\notag
\end{align}
The transposed supported-column calculation gives forms
$S_0,\ldots,S_{r-1}$:
\begin{align}
 S_0&=-\ell_1 C_{0,2}C_{r,1}\prod_{u=2}^{r-1}v_u,
                                                     \label{R22:Sforms}\\
 S_i&=\sum_{p=0}^{i-1}C_{p,i+1}C_{r,p+1}
                    \prod_{\substack{0\leq u<r\\u\ne p}}v_u
                    &&(1\leq i<r).\notag
\end{align}
Their $(x,y)$ degrees are respectively
\begin{equation}\label{R22:grades}
 (1,\ldots,1,2)\quad\hbox{and}\quad(2,1,\ldots,1).
\end{equation}
All their coefficients are homogeneous quadratic polynomials in
$\theta$. These displayed forms are highest components; the actual
limits from \eqref{R22:relation} can also have smaller $(z,w)$ degree.
We retain those lower components until the final independence argument.

\paragraph{3. The Fourier matrices and their exact meanings.}
Fix
\[
 2b-1\leq k,t\leq m-2b+1,\qquad K=q^k,\quad X=q^t.
\]
At $\varepsilon=0$, the common outside factors for the block at zero
are exactly
\begin{equation}\label{R22:outside}
 T_s^R=\ell_0\prod_{i=b}^{s-1}v_i\prod_{i=s+1}^{m-1}\ell_i,
 \qquad
 T_c^C=\prod_{i=b}^{c-1}\ell_i\prod_{i=c}^{m-1}v_i.
\end{equation}
The unique cyclic matching proves these formulas. For a local form
of bidegree $(a,g)$, $a+g=b$, take
\begin{equation}\label{R22:alignment}
 s=k+b-g=k+a,\qquad c=m+g-k.
\end{equation}
These indices are external to the corresponding supported rows or
columns throughout the stated range. The resulting leading products
have bidegree $(m-k,k)$.

Set $A(x,y)=\prod_{u=k}^{m-1}\ell_u$ and let $A_i$ be its coefficient
of $x^{m-k-i}y^i$. All $A_i$ are nonzero: the finite geometric-product
formula for these coefficients has numerator and denominator factors
$1-q^j$ with $0<j<m$ (the same formula used in R-16).
In cyclic character $t$, divide the coefficient with $y$-degree $i$
and $w$-degree $j$, $i+j=t$, by $A_i$.
Under this normalization the entire cyclic base image $W_0$ has the
two-dimensional binary coefficient space
\begin{equation}\label{R22:base}
 \operatorname{span}\left\{
 z\prod_{u=0}^{k-2}v_u,
 w\prod_{u=0}^{k-2}v_u\right\}.
\end{equation}
This is the base-character description proved in R-16; the full base
image has dimension $2m^2+2$ by R-15/R-17.

For a row form $F$ of $(x,y)$ degree $a\in\{1,2\}$, write
\begin{align}
 F&=\sum_{j=1}^{a+1}Q^R_{a,j}(x,y)C_j(z,w),\notag\\
 Q^R_{a,j}&=\prod_{u=j}^{a}(x-Kq^uy)
                       \prod_{u=1}^{j-1}(x-q^uy),\label{R22:rowbasis}\\
 \mathcal R_b(F)&=\sum_{j=1}^{a+1}X^j C_j(z,q^{-j}w)
       \prod_{u=b-j}^{b-2}v_u
       \prod_{u=-1}^{a-1-j}(z-Kq^uw).\label{R22:rowtransform}
\end{align}
For a column form use
\begin{align}
 F&=\sum_{j=0}^{a}Q^C_{a,j}(x,y)C_j(z,w),\notag\\
 Q^C_{a,j}&=\prod_{u=b-j}^{b-1}(x-q^uy)
               \prod_{u=0}^{a-1-j}(x-K^{-1}q^{b-a+u}y),
                                                     \label{R22:colbasis}\\
 \mathcal S_b(F)&=\sum_{j=0}^{a}X^{-j}C_j(z,Kq^{j-b}w)
       \prod_{u=-a+j}^{-1}v_u
       \prod_{u=0}^{j-1}(z-Kq^{-b+u}w).\label{R22:coltransform}
\end{align}
An empty product is one. All displayed bases are invertible for the
specified $q,k$. For the row bases, evaluation at $x=qy$, and then
at $x=q^2y$ when $a=2$, gives triangular independence; the possible
nonconstant vanishing factors are $q-1$, $K-1$, and $Kq-1$.
For the column bases evaluate at $x=q^{b-1}y$, then
$x=q^{b-2}y$. The same conditions, up to nonzero powers, suffice.
None can vanish in the interior range. Thus these transformations are
defined over $\mathbb Q(q)$ after $K=q^k$ and introduce no denominators
depending on $\theta$ or $X$.

Here are direct verifications of the Fourier completions.
Multiplying the row outside $(x,y)$ factor by $Q^R_{a,j}$ completes
$A(x,q^jy)$. Coefficient division gives $q^{ji}=X^jq^{-j\,\deg_w}$.
The remaining $(z,w)$ roots have indices $b-j,\ldots,k+a-1-j$.
Their common factor is
\[
 H_R=\prod_{u=b-1}^{k-2}v_u,
\]
and the residual degree-$b$ form is \eqref{R22:rowtransform}.
For the column completion, multiplication by $Q^C_{a,j}$ completes
$A(x,q^{b-k-j}y)$. After division and the nonzero character scalar
$X^{k-b}$, it gives $X^{-j}q^{(k-b+j)\deg_w}$.
The outside $(z,w)$ roots become $-a+j,\ldots,k-b-1+j$.
Their common factor and residual form are
\[
 H_C=\prod_{u=0}^{k-b-1}v_u,
 \qquad \mathcal S_b(F).
\]
Thus both transformations represent the actual character quotients
of the cofactor limits, including all shifts in $k$ and $t$.

Let $M_R$ be the $(b+1)\times(b+1)$ coefficient matrix of
\begin{equation}\label{R22:MR}
 z\prod_{u=0}^{b-2}v_u,\quad
 w\prod_{u=0}^{b-2}v_u,\quad
 \mathcal R_b(R_0),\ldots,\mathcal R_b(R_{r-1}).
\end{equation}
Let $M_C$ be the corresponding coefficient matrix of
\begin{equation}\label{R22:MC}
 z\prod_{u=-b}^{-2}(z-Kq^uw),\quad
 w\prod_{u=-b}^{-2}(z-Kq^uw),\quad
 \mathcal S_b(S_0),\ldots,\mathcal S_b(S_{r-1}).
\end{equation}
These are precisely base plus row, and base plus column, after
removing their respective common factors. For the column matrix the
additional nonzero character scalar just used also rescales the base;
its root indices are then $k-b,\ldots,k-2$, as in
\eqref{R22:MC}.

\paragraph{4. Uniform row nonvanishing at every interior bidegree.}
Fix the primitive root $q$ and any interior $k$. We prove
$\det M_R\not\equiv0$ as a polynomial in $\theta,X$ by a specialization
that may use elements of $\mathbb Q(q)$.
In the local bottom-row matrix, retain only nearest forward chords
$c_{i,i+2}=p_i$ and first-column backward chords $c_{h,0}=z$.
Set
\[
 p_i=\ell_1\quad(0\leq i\leq r-3),\qquad
 p_{r-2}=\ell_2,
\]
and set its other chords to zero. Parameters outside this row matrix
can be arbitrary. Its row forms are
\[
 -p_i z\prod_{\substack{1\leq u\leq r\\u\ne i+2}}v_u
 \quad(0\leq i\leq r-2),
 \qquad \ell_r\ell_2z\prod_{u=1}^{r-2}v_u.
\]
For $i\leq r-3$ their normalized forms are
\begin{equation}\label{R22:diagonal_forms}
 -X^2z\prod_{\substack{-1\leq u\leq r-1\\u\ne i}}v_u.
\end{equation}
Evaluate modulo the base space at the $r$ distinct roots
$z/w=q^h$, $0\leq h\leq r-1$. The first $r-2$ columns in
\eqref{R22:diagonal_forms} are diagonal on the first $r-2$ roots and
zero on the last two; their diagonal entries are nonzero for $X\ne0$.

The remaining two local forms are
\[
 U=-\ell_2z\prod_{u=1}^{r-1}v_u,\qquad
 V=\ell_r\ell_2z\prod_{u=1}^{r-2}v_u.
\]
Write
\[
 \alpha=\frac{q-1}{K-1},\qquad
 \beta=\frac{K-q}{K-1},\qquad
 \ell_2=\alpha Q^R_{1,1}+\beta Q^R_{1,2}.
\]
Their normalized forms $P_U=\mathcal R_b(U)$,
$P_V=\mathcal R_b(V)$ are
\begin{align*}
 P_U={}&-z\left[
  \alpha X\prod_{u=0}^{r-2}v_u\,(z-Kq^{-1}w)
 +\beta X^2\prod_{u=-1}^{r-3}v_u\,v_{r-1}\right],\\
 P_V={}&c_1Xz\prod_{u=0}^{r-3}v_u\,(z-Kq^{-1}w)(z-Kw)\\
 &{}+c_2X^2z\prod_{u=-1}^{r-4}v_u\,v_{r-1}(z-Kq^{-1}w)\\
 &{}+c_3X^3z\prod_{u=-2}^{r-5}v_u\,v_{r-2}v_{r-1}.
\end{align*}
where $c_j$ are the coefficients of $\ell_r\ell_2$ in the basis
$Q^R_{2,j}$. Only the first is needed:
\begin{equation}\label{R22:c1}
 c_1=\frac{(1-q^{r-1})(1-q)}{(1-K)(1-Kq)}.
\end{equation}
It follows by evaluating the basis expansion at $x=qy$.

Put $A=q^{r-2}$, $B=q^{r-1}$ and evaluate at $w=1$.
$P_U(A)$ has no $X$ term; $P_U(B)$ and $P_V(B)$ have only their
$X$ terms. The coefficient of $X^2$ in the last two-by-two
evaluation determinant is therefore
\begin{equation}\label{R22:last_minor}
 \begin{split}
 [X^2]\big(P_U(A)P_V(B)-P_V(A)P_U(B)\big)
 ={}&\alpha c_1AB
 \prod_{u=0}^{r-2}(B-q^u)
 \prod_{u=0}^{r-3}(A-q^u)\\
 &\cdot(B-K/q)(A-K/q)(A-K).
 \end{split}
\end{equation}
Every factor is nonzero for the fixed primitive $q$ and every
$2b-1\leq k\leq m-2b+1$. In particular the potential equalities
$K=q^{b-1},q^{b-2},q^{b-3}$ lie outside the range, as do $K=1,q^{-1}$.
This also covers $b=3$, when the first diagonal block is empty.
The kernel of evaluation at these $r$ roots is exactly the
two-dimensional base space in \eqref{R22:MR}. Thus the row matrix
has nonzero determinant as a polynomial in $X$ at this specialization,
for every interior $k$ and every primitive $q$ under consideration.
Its specialized degree in $X$ is at most $2b-3$: the first $b-3$
evaluation columns contribute $X^{2(b-3)}$, and the last determinant
has degree at most three.

\paragraph{5. Column nonvanishing at the same primitive root and bidegree.}
The complete parameter domain in \eqref{R22:types} is invariant under
the reflection of additional positions
\begin{equation}\label{R22:reflection}
 (i,j)\longmapsto(b-j,b-i).
\end{equation}
It preserves $j-i$, hence preserves the forward $(x,y)$ type and
backward $(z,w)$ type. Transpose the global cyclic pencil and reverse
its indices by setting $M'_{uv}=M_{b-v,b-u}$ modulo $m$.
Its diagonal and cycle entries are
\[
 \ell_{b-u},\qquad v_{b-u-1}.
\]
With
\[
 q'=q^{-1},\qquad y'=q^by,\qquad w'=q^{b-1}w,
\]
they become the standard cyclic entries $x-(q')^uy'$ and
$z-(q')^uw'$. This is an invertible change of parameters and variables
over $\mathbb Q(q)$.

More explicitly, if $C$ is the transposed supported-column matrix of
the original block and $B'$ is the supported-row matrix of the new
block, then $B'_{ij}=C_{r-1-i,r-j}$. Reversing both rows and columns
preserves the determinant signs of each square minor. It changes the
row relation $a_iE_i-b_iE_{i+1}$ to the negative of the corresponding
reversed relation. This proves \eqref{R22:Sforms} from
\eqref{R22:Rforms} and identifies the two collections of actual
cofactor limits. An external original column $c=m+g-k$ becomes the
external new row
\[
 b-c\equiv k+b-g\pmod m,
\]
exactly the row alignment in \eqref{R22:alignment}. Bidegree $k$ is
unchanged. The change of $y,w$ only rescales character coordinates by
nonzero numbers; it does not mix them or change their available
coefficient interval.

For precision, the following argument does not assume that separate
generic row and column ranks are already simultaneous. Apply the
row specialization just proved at the primitive root $q^{-1}$ and
at this same $k$. Its determinant is a nonzero polynomial in $X$
of degree at most $2b-3$. There are $T=m-4b+3>2b-3$ interior
characters, so at least one of their distinct root values gives
nonzero determinant. Restricting the resulting binary forms to the
allowed coefficient interval preserves their independence: the
Toeplitz/Schur argument in the next paragraph applies to
$H_R\operatorname{Sym}^b$, using any $b+1$ consecutive allowed
indices. Reflection therefore gives a column specialization at $q,k$
with independent actual column classes. If $M_C$ were identically
singular as a function of its parameters and $X$, such independence
after restriction would be impossible. Hence
\begin{equation}\label{R22:two_nonzero}
 \det M_R(q,q^k,\theta,X)\not\equiv0,
 \qquad
 \det M_C(q,q^k,\theta,X)\not\equiv0
\end{equation}
for every interior $k$, over the same parameter ring and the same
fixed field $\mathbb Q(q)$.

\paragraph{6. Joint row--column rank and coefficient restriction.}
Suppose both matrices are nonsingular for some interior $k,t$.
Before coefficient restriction, base plus row spans
$H_R\operatorname{Sym}^b\mathbb C^2$, and base plus column spans
$H_C\operatorname{Sym}^b\mathbb C^2$. Their common factor is
\[
 H=\prod_{u=b-1}^{k-b-1}v_u,
 \qquad N:=\deg H=k-2b+1\geq0.
\]
After dividing by $H$, the two spaces are
\[
 \left(\prod_{u=k-b}^{k-2}v_u\right)\operatorname{Sym}^b\mathbb C^2,
 \qquad
 \left(\prod_{u=0}^{b-2}v_u\right)\operatorname{Sym}^b\mathbb C^2.
\]
The two degree-$b-1$ factors are coprime, since their indices are
distinct and lie below $m$. Their intersection is exactly their
product times $\operatorname{Sym}^1\mathbb C^2$, namely the base.
Their joint span is therefore all
\begin{equation}\label{R22:joint_space}
 H\operatorname{Sym}^{2b-1}\mathbb C^2,
\end{equation}
of dimension $2b$. This is the required joint independence, not an
addition of two ranks without checking their intersection.

The allowed $w$-degree interval in character $t$ is
\[
 \max\{0,t-(m-k)\}\ \leq j\leq\ \min\{k,t\}.
\]
Its length is at least $2b$, since $k,m-k,t,m-t\geq2b-1$.
Choose $2b$ consecutive indices $e,\ldots,e+2b-1$ in it;
then $0\leq e\leq N$. Restriction of multiplication by $H$ to
these indices has determinant
\begin{equation}\label{R22:schur}
 \det(H_{e+i-j})_{0\leq i,j<2b}
 =\pm s_{((2b)^e)}(q^{b-1},q^b,\ldots,q^{k-b-1})\ne0.
\end{equation}
The identity is dual Jacobi--Trudi, as in R-16/R-19 and R-B15.
For nonvanishing, there are $N$ geometric variables, and the Weyl
numerator exponents for this rectangle are distinct in
$[0,N+2b-1]=[0,k]$. Since $k<m$, their powers of $q$ are distinct,
and the numerator is a nonzero Vandermonde determinant; the
denominator is nonzero for the same reason. For $N=0$ the coefficient
matrix is the identity. This proves that all $2b-2$ normal classes
remain independent modulo the entire base character space.
The same proof with $b+1$ in place of $2b$ and the corresponding
geometric factor proves the row-only restriction fact used in the
preceding paragraph.

\paragraph{7. One integer grid choice for all bidegrees.}
We now fix $m,b,q$ and choose the motif parameters once.
By \eqref{R22:grades} and \eqref{R22:rowtransform}, the $X$ exponents
of the row determinant lie in
\[
 [b-1,2b-1]:
\]
every one of its $b-1$ normal columns has minimum exponent one,
the $b-2$ columns with $a=1$ have maximum two, and the column with
$a=2$ has maximum three. The column determinant has exponents in
$[-b,0]$, by \eqref{R22:coltransform} and the sum
$2+(b-2)=b$ of its $a$ values. Consequently, for each interior $k$,
\begin{equation}\label{R22:Fk}
 F_k(\theta,X)=X\det M_R(q,q^k,\theta,X)
                    \det M_C(q,q^k,\theta,X)
\end{equation}
is an ordinary polynomial in $X$ of degree at most $2b$.
Its coefficients lie in $\mathbb Q(q)[\theta]$, and its total
parameter degree is at most $4(b-1)$. By
\eqref{R22:two_nonzero} it is nonzero. The product over all $T$
interior values of $k$ is therefore a nonzero polynomial with
parameter degree at most $D=4(b-1)T$.

A polynomial of degree at most $D$ in each parameter over a field
of characteristic zero
cannot vanish at all points of $\{0,\ldots,D\}^{2b(b-2)}$ unless
it is zero: prove this successively in each variable by univariate
interpolation. Apply this over the field $\mathbb Q(q)(X)$.
There is an integer grid choice $\theta$ for which the product is
nonzero as a polynomial in $X$. Fix that one choice. Every
$F_k(\theta,X)$ is then nonzero, so at most $2b$ interior characters
are bad for each $k$. At least
\begin{equation}\label{R22:good_modes}
 T-2b=m-6b+3
\end{equation}
characters have all $2b-2$ independent normal classes.

This argument deliberately permits the grid choice to depend on
$m,q$; it makes no claim that one fixed integer motif works for
every root order. It needs neither a bound for $\varphi(m)$ nor a
generic root-of-unity specialization inferred from numerical tests.
If a concrete choice is wanted, enumerate the finite grid in
lexicographic order and test \eqref{R22:Fk} in exact cyclotomic
arithmetic, choosing the first with nonzero product. This is a finite
algorithm, with no polynomial-time assertion.

\paragraph{8. Translations, all bidegrees, and a single differential.}
For fixed $k$, the leading vectors from the block at start $a=bj$
are the transforms by $\tau^{bj}$ of those at zero, where
$\tau:(y,w)\mapsto(qy,qw)$. Suppose a linear combination of these
vectors lies in $W_0$. In a good character $t$, their independent
normal classes give $2b-2$ separate equations of the form
\[
 \sum_{j=0}^{h-1}c_j(q^{bt})^j=0.
\]
Every value of $q^{bt}$ has at most $\gcd(b,m)\leq b$ preimages.
The good characters in \eqref{R22:good_modes} therefore contain
at least $h$ distinct such values. Ordinary Vandermonde independence
forces every $c_j$ to vanish. Hence each of the $T$ bidegrees supplies
$(2b-2)h$ independent leading classes modulo $W_0$.

The actual limits of \eqref{R22:relation} need not be homogeneous
for the $(x,y)/(z,w)$ bigrading. In any relation among them modulo
$W_0$, take the largest selected $(z,w)$ degree. The corresponding
highest components are the independent classes just proved; all
remaining components have smaller degree. Since $W_0$ is bigraded,
these coefficients vanish. Descending in degree proves independence
of all $(2b-2)hT$ full limits modulo $W_0$.

Finally each cofactor combination in \eqref{R22:relation}, multiplied
by its common outside permanent, belongs to the differential image
of the same simultaneous pencil whenever $\varepsilon\ne0$.
Dividing by $2\varepsilon^2$ gives a polynomial family, by
\eqref{R22:local_defect}. Its limit is the full vector just used.
Choose $2m^2+2$ ordinary differential columns independent at
$\varepsilon=0$, and append these polynomial families. The resulting
matrix has a nonzero minor at zero, by the quotient independence.
That minor remains nonzero for a common generic nonzero
$\varepsilon$. The differential rank is therefore at least
$2m^2+2+(2b-2)hT$, proving \eqref{R22:rank}.

\paragraph{9. Quantitative consequences and the precise ceiling.}
Since $h\geq m/b-7$ and $T\geq m-4b$, with both lower bounds
nonnegative under $m\geq8b$,
\[
 J_{m,b}\geq(2b-2)(m/b-7)(m-4b)
 \geq(2-2/b)m^2-22bm.
\]
For $m\geq99$, take $b=\lfloor\sqrt{m/11}\rfloor$.
Then $b\geq3$, $m\geq8b$, and
$\frac12\sqrt{m/11}\leq b\leq\sqrt{m/11}$. Thus
\[
 \frac{2m^2}{b}+22bm\leq6\sqrt{11}\,m^{3/2}
 <20m^{3/2},
\]
which proves \eqref{R22:asymptotic}. The padded image has dimension
$c_m+3$ by R-14, while the determinant image has dimension
$2n^2+2$ by R-17. This gives \eqref{R22:padding} and all the stated
multiplicity consequences. Under \eqref{R22:ratio_threshold},
\[
 4m^2-20m^{3/2}+2\geq2\rho^2m^2+2\geq2n^2+2,
\]
so the dimension comparison is strict after padding. Substitution
$\rho=4/3$ gives the threshold $2025$ and, with G-11/G-12 or the
rectangular specialization of G-17, the claimed exact eventual
rectangular range.

For comparison, the dimension method itself cannot pass the
asymptotic ratio $\sqrt2$. Its necessary strict inequality, using
the upper bound in \eqref{R22:asymptotic}, is
\[
 2n^2<4m^2-2m+3.
\]
In particular it fails for $n\geq\sqrt2\,m$ and $m\geq2$.
The proved lower bound approaches this ceiling from below; for
example $m\geq169$ and
$m<n\leq\sqrt2\,m-4\sqrt m$ suffice by direct squaring.
This does not exclude other methods of finding multiplicity gaps
beyond the dimension comparison, and it does not address
superlinear padding or polynomial-time multiplicity computation.
\end{proof}
\subsection{R-23: the exact homological defect after restriction to four variables}

\paragraph{Statement R-23.}
Let $m\geq3$, let $R=\mathbb C[X_{ij}:1\leq i,j\leq m]$, and let
$J_m$ be the ideal generated by all $(m-1)\times(m-1)$
subpermanents, equivalently the first partial derivatives of
$\operatorname{per}_m$. Let
\[
 S=R/(L_1,\ldots,L_{m^2-4})\simeq\mathbb C[x,y,z,w]
\]
be a general four-variable linear restriction. Write
\[
 \epsilon_m=\dim_{\mathbb C}
       \operatorname{Tor}_1^R(R/J_m,S)_m.
\]
Then
\begin{equation}\label{R23:Tor_defect}
 \dim\mathcal P_{m,4}=4m^2-2m+2-\epsilon_m.
\end{equation}
Moreover
\begin{equation}\label{R23:first_degree}
 \operatorname{Tor}_1^R(R/J_m,S)_d=0\qquad(d<m),
\end{equation}
and R-22 implies
\[
 0\leq\epsilon_m\leq20m^{3/2}-2m\qquad(m\geq99).
\]
Thus attaining the exact row--column torus bound is equivalent to one
specified degree of one Tor group vanishing. This vanishing is not
proved here for $m\geq17$.

The general-section defect is already nonzero in a complete known
range. By R-13, for $3\leq m\leq16$ it is exactly
\begin{equation}\label{R23:small_defect}
 \epsilon_m=4m^2-2m+2-\binom{m+3}{3}>0.
\end{equation}
Consequently the full-variable linear syzygy theorem cannot simply be
carried through a general four-variable restriction. Nor can this
restriction be justified by treating all its defining linear forms as
a regular sequence on $R/J_m$: this is impossible for every $m\geq3$.

\begin{proof}
We first specify the cited full-variable input. Efremenko, Landsberg,
Schenck, and Weyman (2018), Theorem~1.1 and \S4.1, Example~4.1
and Theorem~4.3 in the author version, determine the linear strand
of the subpermanent ideal (R-B16). At subpermanent size $m-1$, the
space of linear first syzygies has dimension $2m-2$. A basis consists
of differences between row Laplace expansions of the full permanent,
and differences between column expansions. In the minimal resolution
of $R/J_m$, therefore,
\begin{equation}\label{R23:resolution_start}
 F_1=R(-(m-1))^{m^2},\qquad
 (F_2)_m=\mathbb C^{2m-2},
\end{equation}
and $F_2$ has no components below degree $m$. Other generators of
$F_2$ have higher degree and do not enter the calculation below.

\paragraph{The old linear relations remain independent.}
Write $G_{ij}=\partial\operatorname{per}_m/\partial X_{ij}$.
The row and column expansion vectors are
\[
 \mathsf R_i=(\delta_{ui}X_{uv})_{u,v},\qquad
 \mathsf C_j=(\delta_{vj}X_{uv})_{u,v}.
\]
Each maps to the permanent under $(f_{ij})\mapsto\sum f_{ij}G_{ij}$.
Thus the $2m-2$ differences
$\mathsf R_i-\mathsf R_1$ and $\mathsf C_j-\mathsf C_1$ are
the cited syzygies. A linear combination is the vector
$((\alpha_i+\beta_j)X_{ij})$, where
$\sum\alpha_i=\sum\beta_j=0$.

After restriction to the cyclic pencil of R-15, both its diagonal
and its cyclic superdiagonal entries are nonzero linear forms.
If this relation vector vanishes, its diagonal entries give
$\alpha_i+\beta_i=0$, and its cyclic entries give
$\alpha_i+\beta_{i+1}=0$. Hence all $\beta_i$ are equal and all
$\alpha_i$ are their negatives. Their zero sums force every
coefficient to vanish. The old relations are therefore independent
on that pencil, and hence on a nonempty Zariski open set of
four-variable restrictions.

\paragraph{There are no new constant relations generically.}
For completeness we prove independence of the $m^2$ restricted
gradient generators on the same cyclic pencil. Number indices
$0,\ldots,m-1$, take a primitive $m$th root $q$, and put
$\ell_i=x-q^iy$, $v_i=z-q^iw$. Each complementary permanent has
one matching. For each $0\leq k\leq m-1$, its $m$ complementary
permanents of $(z,w)$ degree $k$ are the cyclic translates of
\begin{equation}\label{R23:cofactor_orbit}
 F_k=\prod_{i=0}^{m-k-2}\ell_i
                 \prod_{i=m-k-1}^{m-2}v_i.
\end{equation}
For example this is the cofactor deleting row $m-1$ and column
$m-k-1$. Empty products have value one.

Every coefficient of each of the two geometric binary products in
\eqref{R23:cofactor_orbit} is nonzero. Indeed each length is below
$m$, so the geometric binomial-product formula has no factor
$1-q^j$ with $j$ divisible by $m$. For every
$0\leq t\leq m-1$, choose integers
$0\leq a\leq m-k-1$, $0\leq b\leq k$ with $a+b=t$.
The corresponding coefficient of $y^aw^b$ is nonzero. Under
$\tau:(y,w)\mapsto(qy,qw)$ this gives a nonzero component of
character $t$. Thus the orbit
$F_k,\tau F_k,\ldots,\tau^{m-1}F_k$ has dimension $m$ by Fourier
inversion. The different $k$ have distinct bidegrees, proving
independence of all $m^2$ cofactors. Their restrictions are therefore
independent for general $S$.

\paragraph{The exact Tor identity.}
Tensor a minimal free resolution of $R/J_m$ with $S$ and take
homogeneous degree $m$. The part relevant to its first homology is
\begin{equation}\label{R23:degree_complex}
 \mathbb C^{2m-2}\longrightarrow
                 S_1^{m^2}\longrightarrow S_m.
\end{equation}
The first map is injective on the general open set just proved.
The second sends $(f_{ij})$ to
$\sum f_{ij}G_{ij}|_S$. It is exactly the differential of the
permanent projection map, because varying any pencil entry by a
linear form $f_{ij}$ changes the permanent by that expression.
In characteristic zero its rank at a general pencil equals the
dimension $c_m$ of the image closure. Consequently the first
homology of \eqref{R23:degree_complex} has dimension
$4m^2-c_m-(2m-2)$, proving \eqref{R23:Tor_defect}.
In degrees below $m-1$ the module $F_1\otimes S$ is zero. In degree
$m-1$ its map to $S$ is injective by the cofactor independence just
proved. This gives \eqref{R23:first_degree}.
Substitution of R-22 proves the stated upper bound for $\epsilon_m$.

\paragraph{Why the known linear strand does not supply the missing vanishing.}
R-13 gives $c_m=\binom{m+3}{3}$ for $3\leq m\leq16$,
and hence \eqref{R23:small_defect}. The positive values are, in order,
\[
 12,\ 23,\ 36,\ 50,\ 64,\ 77,\ 88,\ 96,\ 100,\ 99,\ 92,\ 78,\ 56,\ 25.
\]
These are general-section Tor groups in the very same permanent
family, even though the complete full-variable linear syzygy space
has only the $2m-2$ displayed generators. The larger-section
vanishing for $m\geq17$ therefore needs an additional argument,
which is supplied below by R-27 rather than by the cited linear strand.
At the special cyclic section the defect is even larger:
\[
 \dim\operatorname{Tor}_1^R(R/J_m,S_{\mathrm{cyc}})_m
 =2m(m-1),
\]
because its differential image has dimension $2m^2+2$ and the old
relations remain independent. This special value is not substituted
for the general-section defect.

Finally Boralevi, Carlini, Micha{\l}ek, and Ventura (2024),
Theorem~4.14, record von zur Gathen's bound
$\operatorname{ht}J_m\geq5$ for every $m\geq3$; their Theorem~4.23
strengthens it to six when $m\geq6$ (R-B17). Therefore
\[
 \operatorname{depth}(R/J_m)
 \leq\dim(R/J_m)\leq m^2-5<m^2-4.
\]
A regular sequence on $R/J_m$ cannot have the required length
$m^2-4$. Thus tensoring its resolution with the four-variable
quotient is not justified as an exact operation by a regular-sequence
argument. This obstruction does not itself decide the single graded
piece in \eqref{R23:Tor_defect}: higher-degree Tor may survive even
if that piece vanishes.
\end{proof}



\subsection{R-24: exact maximal dimension at the first nondominant size}

\paragraph{Statement R-24.}
The first matrix size beyond the dominance range of R-13 attains
the row--column torus dimension:
\begin{equation}\label{R24:exact}
 \dim\mathcal P_{17,4}=1124=4\cdot17^2-2\cdot17+2.
\end{equation}
Equivalently, the general-section homological defect in R-23 satisfies
$\epsilon_{17}=0$. This is a finite exact result, not an assertion
that $\epsilon_m=0$ for every $m\geq17$.

In particular, the padded pairs $(m,n)=(17,n)$ with
$18\leq n\leq23$ satisfy the strict four-variable dimension
comparison. The pair $(17,23)$ has nonrectangular multiplicity gaps,
including the stretched direction
\[
 \lambda=(41r,17r,17r,17r),\qquad d=4r,
\]
with the effective existence and eventual-stretch conclusions of
G-17/G-18. G-12 excludes rectangular gaps for this pair.

\begin{proof}
The upper bound $1124$ is R-13. We specify a reproducible integer
pencil and a nonzero differential minor for the opposite inequality.
All indices below are zero-based.

Let $s_0=20278274$, and define unsigned $32$-bit integers
\[
 s_{a+1}=1664525s_a+1013904223\pmod{2^{32}},
 \qquad c_a=\left\lfloor s_{a+1}/2^{16}\right\rfloor\bmod101.
\]
For $0\leq i,j<17$, put
\[
 M_{ij}(x_0,x_1,x_2,x_3)
       =\sum_{k=0}^3c_{4(17i+j)+k}x_k.
\]
This recurrence defines all $1156$ integer pencil coefficients
without an external choice.

List the $1140=\binom{20}{3}$ nonnegative integer quadruples
$\alpha=(\alpha_0,\ldots,\alpha_3)$ with sum $17$ in decreasing
lexicographic order. Let $J$ be the following $1140\times1156$
evaluation matrix over $\mathbb F_{101}$:
\begin{equation}\label{R24:evaluation_jacobian}
 J_{\alpha,\,4(17i+j)+k}
 =\alpha_k\,\operatorname{per}
             \big(M(\alpha)\text{ with row }i
                               \text{ and column }j\text{ deleted}\big).
\end{equation}
Retain its first $1124$ rows and all columns except
\begin{equation}\label{R24:omitted_columns}
 \{68i+67:1\leq i\leq15\}
 \ \cup\
 \{1091+4j:0\leq j\leq16\}.
\end{equation}
In increasing row and column order, the resulting square matrix
$J_*$ has
\begin{equation}\label{R24:residue}
 \det J_*=88\pmod{101}.
\end{equation}

Here is the exact computation underlying the finite certificate.
For each numerical matrix $A=M(\alpha)$, the subset recurrence is
\[
 F(\varnothing)=1,\qquad
 F(S)=\sum_{j\in S}A_{|S|-1,j}F(S\setminus\{j\}).
\]
Its reverse sweep starts with $B([17])=1$ and other reverse
coefficients zero. In decreasing subset order, for every $j\in S$
it adds
\[
 B(S)A_{|S|-1,j}\ \text{to }B(S\setminus\{j\}),\qquad
 B(S)F(S\setminus\{j\})\ \text{to }G_{|S|-1,j}.
\]
The resulting $G_{ij}$ are the complementary permanents used in
\eqref{R24:evaluation_jacobian}. All operations are in
$\mathbb F_{101}$. Ordinary exact Gaussian elimination on the
specified square matrix gives \eqref{R24:residue}.

The source file \nolinkurl{check_permanent_maximal.cpp} defines the
complete calculation, importing the subset evaluator from the
previous R-13 certificate. The accompanying manifest records the
entire pencil, all evaluation points, selected indices, and raw
matrix. Two additional exact checks were executed: Ryser's
Gray-code formula and its reverse derivative sweep independently
recomputed every one of the $1{,}317{,}840$ matrix entries; a
separate NumPy elimination using unnormalized pivot rows recomputed
the selected determinant. Both agree with \eqref{R24:residue}.
An independent agent then reconstructed the generator and every
evaluation point, checked the compact selected-column description,
and recomputed the increasing-order minor as $88$ by modular
fraction-free Bareiss elimination. That check also verified all
$32$ row--column scaling relations on the full matrix and their
independence. Its scope is the selected minor and the scaling
kernel; the preceding all-entry Ryser calculation is a separate
author check.
These are finite arithmetic checks supporting the stated minor;
they are not an extrapolation to larger matrix sizes.

The matrix $J$ is obtained by evaluating the coefficient
differential of the permanent projection map at the listed points.
Evaluation cannot increase rank. A minor nonzero modulo $101$
is a nonzero integer minor, so the differential over $\mathbb C$
has rank at least $1124$ at this integer pencil. Thus
$\dim\mathcal P_{17,4}\geq1124$, proving
\eqref{R24:exact}; R-23 gives $\epsilon_{17}=0$.

Finally the padded projection image has dimension $1124+3=1127$,
whereas the determinant image for $n\leq23$ has dimension at most
$2\cdot23^2+2=1060$. R-14 and G-17/G-18 apply. At $n=23$,
the padding exponent is $a=6$; the direction
$\nu=(41,17,17,17)$ has $|\nu|=92=23\cdot4$,
$\nu_4>0$, and $\nu_1>4a$. This verifies the strict G-17
conditions for the displayed family. Since $23>4\cdot17/3$,
G-12 excludes rectangles.
\end{proof}


\subsection{R-25: a ternary bordered-permanent induction and exact size18 dimension}

The first assertion is an exact conditional reduction. A finite
certificate verifies its hypothesis at $m=17$, giving
$\dim\mathcal P_{18,4}=1262$; no uniform vanishing is inferred.

Let $E=\mathbb C^3$, and let $A(x)$ be an $m\times m$ matrix of
linear forms on $E$. Write $P(A)=\operatorname{per}(A)$ and
$G_{ij}(A)=\partial P/\partial A_{ij}$. For vectors $u(x),v(x)$
of linear forms, define
\[
 Q_A(u,v)=DP_A[uv^t]
          =\sum_{i,j}u_iv_jG_{ij}(A),\qquad
 R_A(u,v)=\tfrac12D^2P_A[uv^t,uv^t].
\]
Thus $P(A+tuv^t)=P(A)+tQ_A(u,v)+t^2R_A(u,v)+O(t^3)$.
The derivatives in $u,v$ below mean differentiation in independent
auxiliary coordinates, followed by substitution of the indicated
linear forms. Put
\begin{align*}
 a_i&=\partial_{u_i}Q_A=(Gv)_i,&
 b_j&=\partial_{v_j}Q_A=(G^tu)_j,\\
 t_i&=\partial_{u_i}R_A
          =D^2P_A[e_iv^t,uv^t],&
 t'_j&=\partial_{v_j}R_A
          =D^2P_A[ue_j^t,uv^t].
\end{align*}
All forms in
\begin{equation}\label{border:C}
 \mathcal C_A(u,v)
 =E^*P(A)
   +\sum_i E^*a_i+\sum_j E^*b_j
   +\operatorname{span}\{t_i,t'_j:1\leq i,j\leq m\}
 \ \subseteq\operatorname{Sym}^{m+1}E^*
\end{equation}
have the same degree. This space is generated by $8m+3$ displayed
columns. If at least one of $u,v$ is nonzero, the identity
\begin{equation}\label{border:scaling_relation}
 \sum_i u_i a_i-\sum_jv_jb_j=0
\end{equation}
is a nonzero relation on the corresponding $6m$ linear-multiple
columns. In particular $\dim\mathcal C_A(u,v)\leq8m+2$.

\paragraph{Statement R-25.}
If some ternary data $A,u,v$ satisfy
\begin{equation}\label{border:rank_hypothesis}
 \dim\mathcal C_A(u,v)=8m+2,
\end{equation}
then, with $c_m=\dim\mathcal P_{m,4}$,
\begin{equation}\label{border:increment}
 c_{m+1}\geq c_m+8m+2.
\end{equation}
Consequently, if $m\geq17$ and $c_m=4m^2-2m+2$, the same
equality holds at $m+1$. A proof of
\eqref{border:rank_hypothesis} for all $m\geq17$, together with
R-24, would prove the exact parameter dimension for all those sizes.

\begin{proof}
Adjoin a fourth variable $\ell$, so $V^*=E^*\oplus\mathbb C\ell$.
Extend $A,u,v$ to matrices and vectors of quaternary linear forms.
The condition \eqref{border:rank_hypothesis} on their restrictions
to $\ell=0$ is a nonempty Zariski open condition on these
quaternary parameters. The condition that the differential image
\[
 T_A=\{DP_A[U]:U\in\operatorname{Mat}_m(V^*)\}
\]
have dimension $c_m$ is also nonempty and Zariski open. The
parameter space is irreducible, so both conditions hold
simultaneously for some quaternary $A,u,v$. No maximal-rank
assumption is made on the initial ternary matrix itself.

Consider the single family of quaternary pencils
\[
 B_\varepsilon=
 \begin{pmatrix}A&\varepsilon u\\
                 \varepsilon v^t&\ell\end{pmatrix}.
\]
Every permutation uses either the bottom-right entry, or one
entry from each off-diagonal block. Hence exactly
\begin{equation}\label{border:exact_permanent}
 \operatorname{per}(B_\varepsilon)
       =\ell P(A)+\varepsilon^2Q_A(u,v).
\end{equation}
Differentiation in the bottom-right entry gives every form in
$V^*P(A)$. Differentiation in the upper-right and lower-left
entries, divided by $\varepsilon$, gives every form in
$\sum V^*a_i+\sum V^*b_j$, for $\varepsilon\ne0$.

The remaining $2m$ columns come from exact cancellations, not
from independently chosen deformations. Fix a constant vector
$h\in\mathbb C^m$. Vary the upper-right block by $\ell h$ and
the upper-left block by $-\varepsilon hv^t$, leaving all other
entries fixed. The differential of
\eqref{border:exact_permanent} in that direction is
\begin{align*}
 \varepsilon\ell\,DP_A[hv^t]
  -\varepsilon\ell\,DP_A[hv^t]
  -\varepsilon^3D_AQ_A(u,v)[hv^t]
   &=-\varepsilon^3D^2P_A[hv^t,uv^t]\\
   &=-\varepsilon^3\sum_i h_it_i.
\end{align*}
Every entry of this parameter variation is a linear form, as
required. Varying instead the lower-left block by $\ell h^t$
and the upper-left block by $-\varepsilon uh^t$ similarly gives
$-\varepsilon^3\sum_jh_jt'_j$. Division by
$-\varepsilon^3$ is permitted in the differential image for
every nonzero $\varepsilon$.

Choose $8m+2$ of the fixed polynomial columns in
\eqref{border:C} whose restrictions to $\ell=0$ are independent.
All belong to the differential image of this same
$B_\varepsilon$, for every $\varepsilon\ne0$.
Choose constant quaternary variations $U_1,\ldots,U_{c_m}$ such
that $DP_A[U_i]$ form a basis of $T_A$. Their derivatives at
$B_\varepsilon$ are
\[
 \ell DP_A[U_i]+\varepsilon^2D_AQ_A(u,v)[U_i].
\]
At $\varepsilon=0$ these become a basis of $\ell T_A$.
The combined $c_m+8m+2$ limit columns are independent: restriction
to $\ell=0$ first annihilates all coefficients of the selected
normal columns, and multiplication by the nonzero polynomial
$\ell$ is injective on $T_A$.

A coefficient minor of the combined columns is therefore a
polynomial in $\varepsilon$ with nonzero constant term. It is
nonzero at some $\varepsilon\ne0$, proving
\eqref{border:increment}. Finally the torus upper bound increases
by exactly
\[
 [4(m+1)^2-2(m+1)+2]-[4m^2-2m+2]=8m+2,
\]
so equality propagates whenever the hypothesis is available.
\end{proof}

\paragraph{Exact size18 corollary.}
The rank hypothesis \eqref{border:rank_hypothesis} holds at $m=17$.
Consequently
\begin{equation}\label{R25:size18}
 c_{18}=1262=4\cdot18^2-2\cdot18+2,
 \qquad \epsilon_{18}=0.
\end{equation}
The genuinely padded pairs $(18,n)$, $19\leq n\leq25$, therefore
satisfy the strict four-variable dimension comparison. In
particular $(18,25)$ admits the G-17/G-18 stretched direction
$\lambda=(46r,18r,18r,18r)$ in coordinate degree $4r$; G-12
excludes its rectangular sector.

\begin{proof}
Here is the complete finite specification of the ternary
certificate. All indices are zero-based. Set $p=1021$ and
$s_0=20330804$, and generate
\[
 s_{a+1}=1664525s_a+1013904223\pmod{2^{32}},
 \qquad c_a=\lfloor s_{a+1}/2^8\rfloor\bmod1021.
\]
Use the first $3\cdot17^2=867$ coefficients, in row, column,
variable order, for $A$; use the next $51$ for $u$ and the final
$51$ for $v$, in vector-index, variable order. These are integer
linear forms in $(x_0,x_1,x_2)$.

List the $190=\binom{20}{2}$ nonnegative integer triples of sum
$18$ in decreasing lexicographic order. Evaluate the following
$139$ columns at those points, in precisely the indicated order:
\begin{equation}\label{R25:columns}
 \begin{array}{ll}
  x_kP(A) &(0\leq k<3),\\
  x_ka_i &(0\leq i<17,\ 0\leq k<3),\\
  x_kb_j &(0\leq j<17,\ 0\leq k<3),\\
  t_i &(0\leq i<17),\\
  t'_j &(0\leq j<17).
 \end{array}
\end{equation}
The $138\times138$ minor on the first $138$ rows and all columns
except column $104$ is, in increasing row and column order,
\begin{equation}\label{R25:minor}
 \det C_*=501\pmod{1021}.
\end{equation}
For completeness the polynomial recurrence determining these
entries without a symbolic permanent expansion is as follows.
At each numerical point put $H=uv^t$. For a subset $S$ of
columns and $i=|S|-1$, recursively compute the polynomial
$F_S(t)$ modulo $t^3$ by
\[
 F_\varnothing(t)=1,\qquad
 F_S(t)=\sum_{j\in S}(A_{ij}+tH_{ij})F_{S\setminus\{j\}}(t)
                         \pmod{t^3}.
\]
The constant, linear and quadratic coefficients of $F_{[17]}$
are exactly $P,Q,R$. Differentiating this finite recurrence in
the auxiliary coordinates $u,v$ gives all columns in
\eqref{R25:columns}. The source
\nolinkurl{check_border_ternary.cpp} implements this recurrence
and reverse differentiation, then exact elimination on the
specified evaluation matrix, yielding \eqref{R25:minor}.

A distinct implementation, \nolinkurl{check_border_ryser.cpp},
recomputed all $26{,}410$ entries from Gray-code Ryser sums:
prefix and suffix products of the row-sum linear factors give
the degree-zero and degree-one complementary products, and
explicit auxiliary differentiation gives the derivatives of
$Q,R$. Every entry agreed. An independent agent reconstructed
all generator coefficients and points, verified the complete
scaling relation \eqref{border:scaling_relation}, and checked
\eqref{R25:minor} by modular fraction-free Bareiss elimination.
The latter check concerns the minor and scaling relation; it
is separate from the author's all-entry Ryser recomputation.

The evaluated columns are evaluations of the homogeneous forms
in \eqref{border:C}. Nonzero determinant modulo $1021$ gives
rank at least $138$ over $\mathbb C$. The universal relation
\eqref{border:scaling_relation} gives the opposite bound, so
\eqref{border:rank_hypothesis} holds at $m=17$.
R-24 and \eqref{border:increment} now give
$c_{18}\geq1124+138=1262$. The torus bound gives equality, and
R-23 identifies the vanishing homological defect.

Finally $c_{18}+3=1265>2\cdot25^2+2=1252$, proving the padded
range by R-14. At $(m,n)=(18,25)$, $a=7$ and
$\nu=(46,18,18,18)$ has $|\nu|=100=4n$, $\nu_4>0$, and
$\nu_1>4a$. Thus G-17/G-18 apply. Since $25>4\cdot18/3$,
G-12 excludes rectangles.
\end{proof}

The uniform gate left by this result is a proof of
\eqref{border:rank_hypothesis} for every $m\geq18$, or a
replacement induction giving the same dimension increment.
The finite certificate alone establishes neither assertion.
R-27 below closes this gate by a uniform endpoint-jet argument,
proving the hypothesis for every $m\geq17$ and hence exact
dimension at every subsequent size.

\subsection{R-26: a uniform linear homological defect from Jacobi continuants}

\paragraph{Statement R-26.}
For every $m\geq17$, some ternary Jacobi pencil and endpoint vectors
in the notation of R-25 satisfy
\begin{equation}\label{R26:criterion}
 \dim\mathcal C_A(u,v)\geq8m-4.
\end{equation}
They can be chosen with
\[
 A_{ii}=\alpha_i x,\quad A_{i,i+1}=y,\quad
 A_{i+1,i}=\delta_i z,\qquad
 u=(x,0,\ldots,0,y)^t,\quad v=(x,0,\ldots,0,z)^t,
\]
where all $\alpha_i,\delta_i$ are positive integers at most $8m^2$.
No efficiency claim is made for selecting them from this finite grid.

Consequently, for every $m\geq18$,
\begin{equation}\label{R26:linear_defect}
 0\leq\epsilon_m\leq6(m-18),\qquad
 4m^2-8m+110\leq c_m\leq4m^2-2m+2.
\end{equation}
In particular every integer pair
\begin{equation}\label{R26:padding}
 m\geq18,\qquad m<n,\qquad n^2\leq2(m-1)^2+53
\end{equation}
satisfies the strict four-variable dimension comparison of R-14.
This includes every integer $m<n\leq\sqrt2(m-1)$.
The explicit stretched-direction conclusions of G-17/G-18 therefore
apply to all these pairs. The uniform conclusion is a linear bound
on the defect, not its vanishing at every size.

\begin{proof}
We first prove \eqref{R26:criterion}. Write $L=m-1$ and $h=yz$.
For an interval of vertices define its continuant
\[
 P_{[a,b]}(x,h)=\operatorname{per}(A_{[a,b],[a,b]}),\qquad
 P_{[a,a-1]}=1.
\]
Its recurrence is
\begin{equation}\label{R26:continuant}
 P_{[a,b]}=\alpha_bxP_{[a,b-1]}
                   +\delta_{b-1}hP_{[a,b-2]}\quad(b>a),
 \qquad P_{[a,a]}=\alpha_ax.
\end{equation}
The two summands distinguish whether the last vertex is fixed or
paired with its predecessor. An interval of $s$ vertices is
homogeneous of weighted degree $s$ for $\deg x=1$, $\deg h=2$.
Put
\[
 D_i=\prod_{e=i}^{L-1}\delta_e,\qquad
 E_i=\prod_{e=0}^{i-1}\delta_e,
\]
with empty products one.

\paragraph{1. Exact endpoint formulas and character spaces.}
The complementary matching for a tridiagonal matrix consists of
the forced path between the removed row and column, and independent
matchings on the intervals outside that path. Thus the first
auxiliary derivatives of R-25 are exactly
\begin{align}
 a_i&=x y^iP_{[i+1,L]}+D_i z^{m-i}P_{[0,i-1]},\label{R26:a}\\
 b_i&=x E_i z^iP_{[i+1,L]}+y^{m-i}P_{[0,i-1]}.
                                                        \label{R26:b}
\end{align}
These formulas include the endpoints using empty intervals.
For $0<i<L$, differentiating the quadratic coefficient of
$\operatorname{per}(A+tuv^t)$ gives
\begin{align}
 t_i&=2x^2D_i z^{m-i}P_{[1,i-1]}
           +2xy^ihP_{[i+1,L-1]},\label{R26:t}\\
 t'_i&=2x^2y^{m-i}P_{[1,i-1]}
           +2xE_i z^ihP_{[i+1,L-1]}.
                                                        \label{R26:tp}
\end{align}
Indeed, the two selected columns are the endpoints. The other
selected row is either endpoint $0$ or endpoint $L$; the intervening
path is forced, leaving the indicated residual interval. The two
assignments of the rank-one entries give the factor $2$.
At the endpoints, with $M=P_{[1,L-1]}$,
\begin{equation}\label{R26:endpoints}
 t_0=t'_0=2xhM,\qquad t_L=2x^2zM,\qquad t'_L=2x^2yM.
\end{equation}

Let a primitive $m$th root $\omega$ act by
$(x,y,z)\mapsto(x,\omega y,\omega^{-1}z)$.
Equations \eqref{R26:a}--\eqref{R26:tp} show that $a_i,t_i$
have character $i$, and $b_i,t'_i$ have character $-i$ modulo $m$.
The polynomial $P=P_{[0,L]}$ has character zero.
Consequently different character spaces can be added independently.

\paragraph{2. Eight columns in each ordinary character.}
Fix $2\leq r\leq m-2$ and put $j=m-r$.
We will prove that the following eight columns are generically
independent:
\begin{equation}\label{R26:eight_columns}
 xa_r, ya_{r-1}, za_{r+1}, xb_j, yb_{j+1}, zb_{j-1}, t_r, t'_j.
\end{equation}
For this assertion specialize $\alpha_i=1$ and retain the
$\delta_i$ as parameters. Each column has the unique form
$y^r p(x,h)+z^jD_rq(x,h)$, where $p$ has weighted degree $j+1$
and $q$ has weighted degree $r+1$. These two summands have different
integer weights, so they occupy disjoint monomial coordinates.
When the edges are nonzero, dividing the second component by
$D_r$ preserves rank. Write
\begin{equation}\label{R26:K}
 K=\frac{E_j}{D_r}
  =\frac{\prod_{e=0}^{\min(r,j)-1}\delta_e}
         {\prod_{e=\max(r,j)}^{m-2}\delta_e}.
\end{equation}
The eight resulting pairs $(p,q)$, in the order
\eqref{R26:eight_columns}, are
\begin{equation}\label{R26:pair_matrix}
\begin{array}{c|l|l}
 &p&q\\ \hline
1&x^2P_{[r+1,L]}&xP_{[0,r-1]}\\
2&xP_{[r,L]}&\delta_{r-1}hP_{[0,r-2]}\\
3&xhP_{[r+2,L]}&\delta_r^{-1}P_{[0,r]}\\
4&xP_{[0,j-1]}&Kx^2P_{[j+1,L]}\\
5&P_{[0,j]}&K\delta_jxhP_{[j+2,L]}\\
6&hP_{[0,j-2]}&K\delta_{j-1}^{-1}xP_{[j,L]}\\
7&2xhP_{[r+1,L-1]}&2x^2P_{[1,r-1]}\\
8&2x^2P_{[1,j-1]}&2KxhP_{[j+1,L-1]}.
\end{array}
\end{equation}
All entries in this display follow directly from
\eqref{R26:a}--\eqref{R26:tp}; in particular the long path
products have disappeared except for the single scalar $K$.

\paragraph{3. Interpreting cut limits correctly.}
For the moment regard $K$ as an independent nonzero scalar in
\eqref{R26:pair_matrix}. We allow some edges to tend to zero,
while keeping $\delta_r$ and $\delta_{j-1}$ nonzero. Such a limit
is a valid rank witness whenever it is approached by nonzero edges
with \eqref{R26:K} tending to the specified $K$.

There are two situations used below. If all zero edges are between
$\min(r,j)$ and $\max(r,j)-1$, they do not occur in the reduced
ratio \eqref{R26:K}; its value is the ratio of the fixed outer
products. If there are additionally equally many zero edges in
the left outer interval $[0,\min(r,j)-1]$ and the right outer
interval $[\max(r,j),m-2]$, with at least one on each side, then
any prescribed nonzero $K$ is possible. To see this, replace each
zero edge by $\tau$ times a nonzero constant. The powers of
$\tau$ in the outer products cancel. Adjust the constant of one
left edge to give the desired ratio; the other constants can be
one. For every $\tau\ne0$ all edges are nonzero, and the entries
of \eqref{R26:pair_matrix} have the stated finite limit.

Thus a nonzero limit minor proves a nonzero minor before taking
the limit. This is a rational row normalization followed by a
limit; it is not a claim about the rank of the literal pencil
obtained by simply setting its edges to zero.

\paragraph{4. Two exact insertion operations.}
The following operations on the normalized matrix preserve its
rank. They use only \eqref{R26:continuant}: inserting an isolated
vertex with diagonal $x$ at a zero edge multiplies every interval
continuant containing that vertex by $x$.

First suppose $r,j\geq6$, with zero edges at positions $2$ and
$m-4$. Insert one isolated vertex at each of these two cuts,
replacing each zero edge by two zero edges. Then
\[
 (m,r,j)\longmapsto(m+2,r+1,j+1).
\]
Every interval in the $p$ column of \eqref{R26:pair_matrix}
contains exactly one new vertex: prefixes contain the left one
and suffixes the right one. The same is true in the $q$ column.
The cuts are away from every endpoint shifted by one in that
display, since $\min(r,j)\geq6$. Accordingly every pair becomes
$(xp,xq)$. The distinguished nonzero edges $\delta_r$,
$\delta_{j-1}$ and the adjacent coefficients are merely
reindexed. The new outer zero-edge counts remain equal, so $K$
can retain its previous value by the preceding limit argument.
Multiplication by $x$ is injective in both components and
preserves the rank. The new pattern still has a left cut at $2$
and a right cut at the new position $(m+2)-4$, allowing repetition.

Second suppose $j-r\geq4$, with a zero edge at $r+1$.
Insert an isolated vertex there. This sends $(m,r,j)$ to
$(m+1,r,j+1)$. Every $p$ interval in
\eqref{R26:pair_matrix} contains the new vertex, including
$[r+2,L]$ and $[0,j-2]$; every $q$ interval avoids it.
Thus the pairs become $(xp,q)$. The cut lies in the middle
interval absent from the reduced ratio \eqref{R26:K}, so $K$
is unchanged. If instead $r-j\geq4$, a cut at $j+1$ gives
$(m+1,r+1,j)$ and pairs $(p,xq)$. The same endpoint check applies.
Both operations preserve injectivity and leave a cut at the
appropriate position for repetition.

\paragraph{5. Complete finite core certificate.}
It remains to supply finitely many starting matrices. Work
modulo $101$. Before imposing cuts, let $\delta_i$ be the
residue of $7i^2+13i+36$ in $\{1,\ldots,100\}$, replacing a
zero residue by $1$. Take the following cores.

For the nine pairs $(r,j)=(6+\max(-g,0),6+\max(g,0))$,
$-4\leq g\leq4$, put $\delta_2=\delta_{m-4}=0$ and $K=14$.
If $|g|=4$, also put $\delta_{\min(r,j)+1}=0$.
For the eight pairs $r+j=17$, $\min(r,j)\in\{2,3,4,5\}$,
put only $\delta_{\min(r,j)+1}=0$ and define $K$ by the exact
rational outer-product ratio \eqref{R26:K}. The residues of
these latter $K$ values are shown below. The edges appearing
in denominators in \eqref{R26:pair_matrix} are nonzero in all
these cases.

Expand the positive components in the ordered basis
$x^{j+1-2k}h^k$, $0\leq k\leq\lfloor(j+1)/2\rfloor$, followed
by the negative components in $x^{r+1-2k}h^k$,
$0\leq k\leq\lfloor(r+1)/2\rfloor$. Number these coefficient
rows from zero. All eight columns are used, in the displayed
order. The following ordered row lists have the stated nonzero
determinants modulo $101$:
\begin{center}
\begin{tabular}{rrrlr}
$r$&$j$&$K\bmod101$&ordered rows&determinant\\ \hline
10&6&14&$(0,1,4,2,3,5,6,7)$&73\\
9&6&14&$(0,1,4,2,3,5,6,7)$&58\\
8&6&14&$(0,1,4,2,3,5,6,7)$&16\\
7&6&14&$(0,1,4,2,3,5,6,7)$&8\\
6&6&14&$(0,1,4,2,3,5,6,7)$&53\\
6&7&14&$(0,1,5,3,2,6,7,8)$&51\\
6&8&14&$(0,1,5,3,4,6,2,7)$&72\\
6&9&14&$(0,1,6,3,4,7,2,8)$&20\\
6&10&14&$(0,1,6,3,4,7,2,8)$&90\\
2&15&34&$(0,1,9,3,4,10,6,7)$&51\\
15&2&34&$(0,1,2,3,4,5,6,7)$&16\\
3&14&61&$(0,1,8,3,4,9,6,7)$&88\\
14&3&61&$(0,1,3,4,2,5,6,7)$&93\\
4&13&88&$(0,1,8,3,4,9,6,7)$&23\\
13&4&88&$(0,1,3,2,4,5,6,7)$&24\\
5&12&82&$(0,1,7,3,4,8,6,2)$&83\\
12&5&82&$(0,1,4,2,3,5,6,7)$&40
\end{tabular}
\end{center}
This is a finite certificate specified entirely by the recurrence
\eqref{R26:continuant}, the pair matrix
\eqref{R26:pair_matrix}, the integer parameter rule, and the
row lists. The source \nolinkurl{check_jacobi_cores.py} and its
manifest give every coefficient matrix. Integer Bareiss
determinants of the listed matrices were also separately
computed and give exactly these residues. An independent agent
then reconstructed all seventeen coefficient matrices by direct
enumeration of path matchings, rather than the continuant recurrence,
and recomputed every selected determinant by modular fraction-free Bareiss
elimination. All matrices and residues agreed. Each nonzero residue
is a nonzero rational minor in characteristic zero.

For any target $r,j\geq6$, first choose the core with the same
signed gap, truncated to $[-4,4]$. Repeated paired insertions
raise the smaller index from $6$ to $\min(r,j)$, preserving
the gap. If the desired absolute gap exceeds $4$, repeated
central insertions then raise the larger index to its target.
For $\min(r,j)\in\{2,3,4,5\}$ and $r+j\geq17$, start at
the corresponding size17 core and use central insertions.
Every ordinary character $2\leq r\leq m-2$, $m\geq17$, is
therefore covered. Its eight-column minor is nonzero for some
nonzero edge parameters, so it is nonzero on a nonempty Zariski
open set of the full Jacobi parameter space.

\paragraph{6. The three exceptional characters.}
Here put all $\delta_i=1$ and all $\alpha_i=1$ except
$\alpha_{m-2}=2$. Let $C_0=1,C_1=x$ and
$C_s=xC_{s-1}+hC_{s-2}$ denote the unperturbed continuants.
Consecutive $C_s,C_{s-1}$ are coprime in $\mathbb C[x,h]$.
Indeed, a common irreducible divisor divides $hC_{s-2}$.
If it does not divide $h$, repeat the recurrence until it divides
$C_0=1$, a contradiction. If it divides $h$, it is associated
to $h$, whereas $C_s(x,0)=x^s$ shows that $h$ is not a factor.
For the degree arguments below, $x,h$ have weights $1,2$.

Set
\[
 S_0=P_{[1,L]},\quad S_1=P_{[2,L]},\quad
 F=P_{[0,L-1]},\quad M=P_{[1,L-1]},\quad
 A_0=C_{m-3},\quad B_0=C_{m-4}.
\]
The single changed diagonal gives
\begin{align*}
 S_0&=(2x^2+h)A_0+xhB_0,&
 S_1&=xA_0+(h+x^2)B_0,\\
 F&=(2x^2+h)A_0+2xhB_0,&
 M&=2xA_0+hB_0.
\end{align*}

In character zero, the available columns span
\begin{equation}\label{R26:zero_basis}
 x^2S_0,\quad xhS_1,\quad hF,\quad xhM,\quad xy^m,\quad xz^m.
\end{equation}
For clarity, $xP=x^2S_0+xhS_1$;
$xa_0=x^2S_0+xz^m$ and $za_1=xhS_1+xz^m$;
$xb_0=x^2S_0+xy^m$; $ya_L=xy^m+hF$;
and $t_0=2xhM$. These identities recover all six displayed
forms, since $2\ne0$.
The two extreme monomials are independent of the balanced
polynomials in $x,h$. To prove that the four balanced forms
are independent, suppose their coefficients are $a,b,c,d$.
Their relation becomes
\begin{align*}
0={}&A_0\big(2ax^4+(a+b+2c+2d)x^2h+ch^2\big)\\
 &+B_0xh\big((a+b)x^2+(b+2c+d)h\big).
\end{align*}
Coprimeness and $\deg A_0=m-3>5$ imply that both coefficient
polynomials vanish. Their coefficients then give
$a=b=c=d=0$. Thus character zero has dimension at least six.

In character one, the space contains the four pure weight-one
forms
\begin{equation}\label{R26:one_basis}
 yP,\qquad x^2yM,\qquad x^2yS_1,\qquad xyF.
\end{equation}
The first two are $yP$ and $t'_L/2$. The identity
$za_2-xa_1-ya_0=-2x^2yS_1$ gives the third, and
$xa_1-xb_L=x^2yS_1-xyF$ gives the fourth.
After removing the factor $y$, a relation with coefficients
$a,b,c,d$ becomes
\begin{align*}
0={}&A_0x\big((2a+2b+c+2d)x^2+(2a+d)h\big)\\
 &+B_0\big(ah^2+(2a+b+c+2d)x^2h+cx^4\big).
\end{align*}
Now $\deg A_0>4$ and coprimeness again force all four
coefficients to vanish. Modulo pure integer weight one, the
columns $xa_1,ya_0,yb_0$ have the independent images
$x^2z^{m-1},yz^m,y^{m+1}$. Character one therefore has dimension
at least seven. Transposition together with $y\leftrightarrow z$
preserves this Jacobi pencil and exchanges $u,v$, so character
$-1$ has the same lower bound. All degree inequalities hold
throughout the stated range $m\geq17$.

\paragraph{7. One common pencil and the uniform consequences.}
Each ordinary-character minor and each exceptional-character
minor is nonzero at some point of the irreducible Jacobi parameter
space with nonzero $\alpha_i,\delta_i$. Their nonempty open
sets have nonempty finite intersection. At one common pencil,
the character dimensions therefore add to
\[
 8(m-3)+6+7+7=8m-4,
\]
proving \eqref{R26:criterion}. This intersection step is essential:
the separate cut specializations are witnesses that the different
minor polynomials are nonzero, not simultaneous zero-edge choices.

There is also a bounded integer choice. Every coefficient entry
of the original columns has parameter degree at most $m$.
Choose one nonzero minor of size eight for each ordinary
character, and minors of sizes six, seven, seven for the other
three. Their product has degree at most
$8m(m-3)+20m=8m^2-4m$. The finite-grid lemma used in R-22
therefore gives a nonzero value with all $2m-1$ parameters in
$\{1,\ldots,8m^2\}$. This proves the asserted bounded choice
without a polynomial-time selection claim.

The exact cancellation and open-intersection proof of R-25 works
with any number $s$ of independent ternary normal columns:
their restrictions are independent of $\ell T_A$, so the same
limit minor gives $c_{m+1}\geq c_m+s$. Applying
\eqref{R26:criterion} yields
\[
 c_{m+1}\geq c_m+8m-4\qquad(m\geq17).
\]
Since the torus upper bound increases by $8m+2$, R-23 gives
$\epsilon_{m+1}\leq\epsilon_m+6$.
Starting from $\epsilon_{18}=0$ in R-25 proves
\eqref{R26:linear_defect}. Finally
\[
 c_m\geq4(m-1)^2+106,
\]
whereas \eqref{R26:padding} implies
$2n^2\leq4(m-1)^2+106<c_m+1$. Thus $c_m+3>2n^2+2$,
which is exactly the strict dimension hypothesis used in R-14
and G-17/G-18. This proves \eqref{R26:padding}.
\end{proof}


\paragraph{Corollary R-26: the eventual rectangular padding range from size14.}
For every $m\geq14$ and every integer $m<n<4m/3$, the padded
permanent and determinant have the explicit rectangular
multiplicity gaps supplied by G-11/G-17, with the eventual-stretch
conclusion of G-18. For every $n\geq4m/3$, G-12 excludes
rectangular gaps. The threshold $14$ is sufficient here; no
optimality assertion is made.

\begin{proof}
For $m\geq18$,
\[
 2(m-1)^2+53-\frac{16m^2}{9}
       =\frac29(m-9)^2+37>0.
\]
Thus every $m<n<4m/3$ lies inside the dimension range
\eqref{R26:padding}. For the four smaller sizes use R-13/R-24:
\begin{center}
\begin{tabular}{rrrr}
$m$&largest integer $n<4m/3$&$c_m$&$2n^2-1$\\ \hline
14&18&680&647\\
15&19&816&721\\
16&21&969&881\\
17&22&1124&967
\end{tabular}
\end{center}
In each case $c_m>2n^2-1$, equivalently
$c_m+3>2n^2+2$. The same strict inequality holds for every
smaller eligible $n$. G-11/G-17 therefore apply in the entire
strict rectangular range, and G-18 gives the stated eventual
stretches. G-12 supplies the exclusion at and above $4m/3$.
\end{proof}


\subsection{R-27: exact four-variable dimension from a uniform endpoint repair}

\paragraph{Statement R-27.}
For every $m\geq17$ there are a ternary Jacobi pencil $A$ and
linear vectors $u,v$, supported at its two endpoints, for which
\begin{equation}\label{R27:criterion}
 \dim\mathcal C_A(u,v)=8m+2.
\end{equation}
Consequently, with the notation of R-13 and R-23,
\begin{equation}\label{R27:dimension}
 c_m=\min\left\{\binom{m+3}{3},\,4m^2-2m+2\right\}
 \quad(m\geq1),
 \qquad \epsilon_m=0\quad(m\geq17).
\end{equation}
The uniform assertion follows from one fixed finite boundary
certificate and an insertion argument. No extrapolation of
finite-size ranks, and no assertion about variable counts growing
with $m$, is used.

\begin{proof}
We use the notation and continuants of R-26. In particular,
$L=m-1$, $h=yz$, $A_{ii}=\alpha_i x$, $A_{i,i+1}=y$,
$A_{i+1,i}=\delta_i z$, and
\[
 D_i=\prod_{e=i}^{L-1}\delta_e,\qquad
 E_i=\prod_{e=0}^{i-1}\delta_e,\qquad D=D_0=E_L.
\]
All parameters are initially nonzero. We work over
characteristic zero; reduction modulo $101$ will certify specified
nonzero rational minors.

\paragraph{1. General endpoint formulas and the seven base relations.}
Allow arbitrary endpoint linear forms
\[
 u=(U_0,0,\ldots,0,U_1)^t,\qquad
 v=(V_0,0,\ldots,0,V_1)^t.
\]
The path decomposition used in R-26 gives
\begin{align}
 a_i&=V_0y^iP_{[i+1,L]}+V_1D_i z^{L-i}P_{[0,i-1]},
                                                    \label{R27:a}\\
 b_i&=U_0E_i z^iP_{[i+1,L]}+U_1y^{L-i}P_{[0,i-1]}.
                                                    \label{R27:b}
\end{align}
For $0<i<L$ the second auxiliary derivatives are
\begin{align}
 t_i&=2V_0V_1\bigl(U_0D_i z^{L-i}P_{[1,i-1]}
                         +U_1y^iP_{[i+1,L-1]}\bigr),
                                                    \label{R27:t}\\
 t'_i&=2U_0U_1\bigl(V_0y^{L-i}P_{[1,i-1]}
                         +V_1E_i z^iP_{[i+1,L-1]}\bigr).
                                                    \label{R27:tp}
\end{align}
At $i=0$ omit the first summand in each parenthesis, and at
$i=L$ omit the second. These are exact polynomial identities:
each complementary matching has a forced path and an independent
interval matching, while the two assignments of the rank-one
entries give the factor $2$.

Let $C(\varepsilon)$ be the coefficient matrix of the $8m+3$
columns defining $\mathcal C_A(u,v)$, when
\begin{equation}\label{R27:endpoint_line}
 (U_0,U_1,V_0,V_1)
 =(x,y,x,z)+\varepsilon(y+2z,\,3x+4z,\,5y+6z,\,7x+8y).
\end{equation}
Its entries have degree at most three in $\varepsilon$. Write
$C=C_0+\varepsilon C_1+\varepsilon^2C_2+\cdots$.
At $\varepsilon=0$ the character decomposition of R-26 applies.
There are eight columns in each character $2,\ldots,m-2$ and
nine columns in each of $0,1,-1$.

For later upper bounds we give all seven independent base
relations. Three have character zero:
\begin{align}
 xa_0+ya_L-xb_0-zb_L&=0,\label{R27:scaling_base}\\
 K_{00}:=yb_1-\delta_0za_1+\alpha_0xa_0-\alpha_0xb_0&=0,
                                                    \label{R27:K00}\\
 K_{01}:=t_0-t'_0&=0.                              \label{R27:K01}
\end{align}
Two have character one:
\begin{align}
 K_{10}&:=\delta_0t_1-2xb_L+\alpha_0t'_L=0,
                                                    \label{R27:K10}\\
 K_{11}&:=\alpha_L\delta_0
       (\delta_1za_2-ya_0+\alpha_1xa_1)\notag\\
 &\hspace{17mm}-\alpha_0\alpha_1
       (\delta_{L-1}zb_{L-1}-yP+\alpha_Lxb_L)=0.
                                                    \label{R27:K11}
\end{align}
The remaining two are the transposed and gauged versions of
\eqref{R27:K10}--\eqref{R27:K11}, of character $-1$.
Here and below identities involving these $K$'s without a stated
endpoint specialization mean the displayed constant linear
combination of the corresponding columns.

For precision, the symmetry producing the two negative-character
relations is as follows. Transpose $A$, exchange $y,z$, and
conjugate by the diagonal matrix with entries $1/E_i$, in the
form $S^{-1}AS$. The upper and lower coefficients again become
$1,\delta_i$, while the endpoint pair becomes
$(U_0,U_1,V_0,V_1)=(x,Dy,x,D^{-1}z)$.
Next make the variable change $(y,z)\mapsto(ty,t^{-1}z)$ and
conjugate by $\operatorname{diag}(t^{-i})$, choosing
$t^m=D^{-1}$. This restores $(x,y,x,z)$. These operations are
invertible over $\mathbb C$, preserve the Jacobi family, exchange
characters $r,-r$, and take arbitrary endpoint perturbations to
arbitrary endpoint perturbations.

All five displayed identities follow by the left or right
continuant recurrence. For example, the two parentheses in
\eqref{R27:K11}, at the base endpoints, equal respectively
$2\alpha_0\alpha_1D_1x^2z^{m-1}$ and
$2\alpha_LDx^2z^{m-1}$. The three character-zero relations are
independent: the Hessian columns distinguish $K_{01}$, and the
columns with index one distinguish $K_{00}$ from scaling.
In character one, $K_{10}$ has Hessian support and $K_{11}$ has
a nonzero $yP$ coefficient. The same argument applies after the
symmetry. Thus the seven relations are independent whenever the
parameters are nonzero. R-26 proves the reverse generic rank
bound. Consequently
\begin{equation}\label{R27:base_rank}
 \operatorname{rank}C_0=8m-4
\end{equation}
on a nonempty open set, and its kernel consists of these seven
relations, entirely in characters $0,1,-1$.

\paragraph{2. A universal upper bound of four for the first normal.}
This step is needed to control the second-order continuation.
It holds for arbitrary linear endpoint perturbations, not just
\eqref{R27:endpoint_line}.
For general endpoint forms, direct use of the same recurrences
gives the exact identity
\begin{equation}\label{R27:K11_general}
 K_{11}=2\alpha_0\alpha_1\alpha_LD\,
       xz^{L-1}(xV_1-zU_0)
          -\alpha_0\alpha_1(U_1-y)P.
\end{equation}
Indeed, its first parenthesis is
$2\alpha_0\alpha_1D_1x^2z^{L-1}V_1$, and its second is
$2\alpha_LDxz^LU_0+(U_1-y)P$. This proves
\eqref{R27:K11_general}, including every endpoint factor.

We show that the first derivative of this identity lies in
$\operatorname{im}C_0$. The term involving $P$ already does.
The other terms are linear combinations of
\begin{equation}\label{R27:five_pure}
 x^3z^{m-2},\quad x^2yz^{m-2},\quad x^2z^{m-1},
 \quad xyz^{m-1},\quad xz^m.
\end{equation}
Here are explicit base-column witnesses for all five.
For $r=2,3$, put $j=m-r$ and define
\begin{align}
 R_A(r)&=ya_{r-1}-\alpha_rx a_r-\delta_rza_{r+1}
       =-2\alpha_rD_r xz^{m-r}P_{[0,r-1]},
                                                    \label{R27:RA}\\
 R_B(r)&=yb_{j+1}-\alpha_jxb_j-\delta_{j-1}zb_{j-1}
       =-2\alpha_jE_jx^2z^jP_{[j+1,L]}.
                                                    \label{R27:RB}
\end{align}
The unwanted path branch cancels by one continuant recurrence;
the remaining branch is the displayed expression by the recurrence
at the other end. When $r=2$, the two residual polynomials span
$x^3$ and $xh$, because
$P_{[0,1]}=\alpha_0\alpha_1x^2+\delta_0h$ and
$P_{[L,L]}=\alpha_Lx$. Hence the first and fourth forms in
\eqref{R27:five_pure} are available. When $r=3$ they span
$x^4$ and $x^2h$ provided
\begin{equation}\label{R27:Delta}
 \Delta=\alpha_0\alpha_1\alpha_2\delta_{L-1}
       -\alpha_{L-1}\alpha_L(\alpha_2\delta_0+\alpha_0\delta_1)
       \ne0.
\end{equation}
This follows by expanding $xP_{[0,2]}$ and
$x^2P_{[L-1,L]}$; their coefficient determinant is precisely
$\Delta$. It supplies the second form in \eqref{R27:five_pure}.
Finally
\begin{align*}
 \delta_1za_2-ya_0+\alpha_1xa_1
       &=2\alpha_0\alpha_1D_1x^2z^{m-1},\\
 \delta_0za_1-xP+\alpha_0xa_0
       &=2\alpha_0D xz^m
\end{align*}
supply the other two. Condition \eqref{R27:Delta} defines a
nonempty open set: when all parameters equal one, $\Delta=-1$.

Thus $K_{11}$ has zero first normal on a nonempty open set.
The symmetry above proves the same assertion for its independent
negative-character counterpart. Scaling also has zero first
normal, since the actual polynomial relation
\begin{equation}\label{R27:scaling_actual}
 U_0a_0+U_1a_L-V_0b_0-V_1b_L=0
\end{equation}
holds for every endpoint tuple; differentiating it gives its
first-order lift. These are three independent vectors in the
seven-dimensional base kernel. Therefore, generically,
\begin{equation}\label{R27:first_upper}
 \operatorname{rank}\left(\ker C_0\longrightarrow
         \operatorname{coker}C_0,\quad k\longmapsto[C_1k]\right)
 \leq4.
\end{equation}
The witnesses above for $K_{11}$ use only the complete character
blocks $0,1,2,3$, and the transposed witnesses use only
$0,-1,-2,-3$.

\paragraph{3. The finite-order rank lemma and the complete boundary projection.}
We record the elementary linear algebra used below. Let
$M(\varepsilon)=M_0+\varepsilon M_1+\varepsilon^2M_2+\cdots$.
Let $r_0=\operatorname{rank}M_0$ and let $r_1$ be the rank of
\[
 N_1:\ker M_0\longrightarrow\operatorname{coker}M_0,
 \qquad k_0\longmapsto[M_1k_0].
\]
For $k_0\in\ker N_1$ choose $k_1$ with
$M_0k_1=-M_1k_0$. Define
\[
 N_2(k_0)=[M_2k_0+M_1k_1]
 \quad\hbox{in }\operatorname{coker}M_0/\operatorname{im}N_1,
 \qquad r_2=\operatorname{rank}N_2.
\]
Changing the choice of $k_1$ changes this class by an element
of $\operatorname{im}N_1$, so the map is well defined.
For the lower triangular coefficient matrices
\begin{equation}\label{R27:Toeplitz}
 T_0=M_0,\quad
 T_1=\begin{pmatrix}M_0&0\\M_1&M_0\end{pmatrix},\quad
 T_2=\begin{pmatrix}M_0&0&0\\M_1&M_0&0\\M_2&M_1&M_0\end{pmatrix}
\end{equation}
one has
\begin{equation}\label{R27:jet_ranks}
 \operatorname{rank}T_0=r_0,\quad
 \operatorname{rank}T_1=2r_0+r_1,\quad
 \operatorname{rank}T_2=3r_0+2r_1+r_2.
\end{equation}
Moreover, $\operatorname{rank}_{\mathbb C(\varepsilon)}M
\geq r_0+r_1+r_2$. To prove these assertions, choose constant
bases that put $M_0$ in an identity block with a zero complement.
Eliminate that block by invertible formal row and column
operations. The complement starts with $\varepsilon N_1$.
Eliminate an identity block of $N_1$ in the same way; the
remaining coefficient of $\varepsilon^2$ is exactly the map
$N_2$ above. On truncation modulo $\varepsilon^{s+1}$,
invertible formal changes give invertible triangular block
changes. A unit pivot contributes $s+1$ to the rank, a pivot
of order one contributes $s$, and a pivot of order two contributes
$s-1$ when $s\geq2$. This proves all the formulas and the
generic-rank lower bound without an analytic limiting argument.

Now retain every base column of $C$ whose character is in
\[
 B_m=\{0,\pm1,\ldots,\pm5\}\pmod m,
\]
and retain every monomial row of these same characters. There
are $91$ retained columns: nine in each of $0,1,-1$, and eight
in each of the other eight characters. Denote this projected
matrix by $S(\varepsilon)$. On the generic base-rank open,
$\operatorname{rank}S_0=84$ and its kernel is precisely the
seven-dimensional full kernel.

It is essential here to retain complete character blocks.
Replacing one endpoint monomial by another linear monomial
changes a character by at most two. Consequently, for every
$k_0\in\ker C_0$, $C_1k_0$ is supported in characters
$0,\pm1,\pm2,\pm3$. If its class modulo the full base image
vanishes, its correction $k_1$ can be chosen entirely in these
complete base character blocks: project any solution of
$C_0k_1=-C_1k_0$ onto them. Then
\begin{equation}\label{R27:support}
 C_2k_0+C_1k_1
 \quad\hbox{is supported in characters }0,\pm1,\ldots,\pm5.
\end{equation}
Indeed, two replacements shift the original kernel characters
by at most four, and one replacement shifts the correction
characters by at most two. The first and second normal maps
of $S$ are therefore exactly the first and second normal maps
of the full $C$, with their images in these subspaces. Columns
of outside base characters cannot change either quotient.
This also shows why a rank assertion for a chosen boundary
submatrix alone would not suffice.

\paragraph{4. A fully specified size17 boundary certificate.}
Use $m=17$, $L=16$, and the following integer parameters:
\begin{align*}
 (\alpha_0,\ldots,\alpha_{16})
   &=(86,18,1,17,95,90,80,90,48,100,9,60,2,25,78,53,54),\\
 (\delta_0,\ldots,\delta_{15})
   &=(34,1,37,24,47,1,59,41,0,92,38,56,75,6,44,14).
\end{align*}
The zero is the central edge $8$. The endpoint line is exactly
\eqref{R27:endpoint_line}. The calculation at a zero edge is
interpreted by the following rational row normalization, rather
than as the rank of the literal degenerate pencil.

For nonzero edges divide each retained row of integer weight
$w=b-c$ with $-m-5\leq w\leq-m+5$ by $D=\prod\delta_i$.
Leave every other row unchanged. Then take the limit
$\delta_8\to0$. In these negative extreme rows every term
contains the long lower path across edge $8$, and its residual
interval avoids the edge. The factor $\delta_8$ cancels before
taking the limit. Equivalently, for these rows set
$\delta_8=1$ and divide by the product of the other edges;
for the other rows set $\delta_8=0$. This product is nonzero
and is $1$ modulo $101$. This rule specifies every normalized
coefficient without division by zero.

Order degree18 monomials $x^ay^bz^c$ by decreasing $a$ and
then decreasing $b$, and retain the $124$ rows with
$(b-c)\bmod17\in B_{17}$. Order the original columns as
\[
\begin{aligned}
 &xP,yP,zP;\qquad xa_i,ya_i,za_i\quad(i=0,\ldots,16);\\
 &xb_i,yb_i,zb_i\quad(i=0,\ldots,16);\\
 &t_i\quad(i=0,\ldots,16);\qquad t'_i\quad(i=0,\ldots,16).
\end{aligned}
\]
and retain the $91$ columns of base character in $B_{17}$.
Their characters are, respectively,
$(0,1,-1)$, $(i,i+1,i-1)$, $(-i,-i+1,-i-1)$, $i$, and $-i$.
Expand the exact formulas \eqref{R27:a}--\eqref{R27:tp} in
$\varepsilon$, and form the three matrices
\eqref{R27:Toeplitz} with $M_i=S_i$.

The following is a finite rational-minor certificate. All
arithmetic in this table is in $\mathbb F_{101}$:
\begin{center}
\begin{tabular}{rrrr}
matrix&dimensions&rank&selected determinant\\ \hline
$T_0$&$124\times91$&84&99\\
$T_1$&$248\times182$&172&74\\
$T_2$&$372\times273$&262&45
\end{tabular}
\end{center}
For an unambiguous specification of the selected minors, use
the following deterministic elimination rule. Scan columns
from left to right. In each column choose the first nonzero
entry among unused rows, swap that row into the next pivot
position, normalize its pivot to one, and eliminate entries
below it. Record the original row names in pivot order and the
selected column names in scan order. The selected determinant
in this ordered-row convention is the product of the pivots
before normalization. The remaining rows are zero at the end.
Together with the integer input, the continuant recurrence,
the endpoint formulas, and the row and column ordering, this
algorithm specifies the entire finite certificate. In particular
the table is not dependent on an unspecified numerical rank
tolerance or an unspecified choice of a minor.

The complete coefficient arrays and selected row and column
lists are also preserved in
\nolinkurl{jacobi_boundary_repair_m17.npz} and
\nolinkurl{boundary_toeplitz_certificates.json}, with builders
\nolinkurl{check_jacobi_boundary_repair.py} and
\nolinkurl{check_boundary_toeplitz.py}. A separate checker
reconstructed all $33{,}852$ entries of $S_0,S_1,S_2$ by direct
enumeration of path matchings and multiplication of polynomials
in $x,y,z,\varepsilon$, and independently obtained the three
ranks $84,172,262$. This is an actual finite coefficient and
rank check; the uniform argument is supplied next. Each listed
nonzero residue proves a nonzero rational minor in
characteristic zero.

\paragraph{5. Inserting a central vertex preserves every retained jet.}
Start with the preceding central-cut data and insert an isolated
vertex with diagonal $x$ at edge $8$, replacing that zero edge
by two zero edges. The size increases by one. The left endpoint
indices are unchanged; the right endpoint indices increase by
one. Repeat at the leftmost central cut as often as desired.
All other coefficients retain their positions relative to the
appropriate endpoint. This constructs cut data in every size
$m\geq17$.

We verify the effect on the entire normalized boundary matrix
for arbitrary endpoint forms, so in particular on each of its
$\varepsilon$ coefficients. Retained first-derivative columns
have index at most $6$ from the left or at least $m-6$ from the
left. Retained Hessian columns have index at most $5$ or at
least $m-5$. In formulas \eqref{R27:a}--\eqref{R27:tp}, each
such column is the sum of a short-path term with a long residual
continuant and a long-path term with a short residual continuant.
Every long residual continuant contains the new vertex; every
short residual continuant avoids it. Every long forced path
crosses the central cut; every short forced path avoids it.
These assertions include the intervals $[i+1,L-1]$ and
$[1,i-1]$: their short endpoint distances are at most $5$,
whereas the original cut is between vertices $8$ and $9$.
They remain true after every insertion.

There is no ambiguity in the retained rows about which term
occurs. A first-derivative column has an auxiliary factor of
degree two, so its short-path terms have integer weight of
absolute value at most $6+2=8$, whereas its long-path terms
have absolute weight at least $(m-7)-2=m-9$.
A Hessian column has an auxiliary factor of degree three;
the corresponding bounds are $5+3=8$ and $(m-6)-3=m-9$.
The columns $xP,yP,zP$ have integer weight $0,1,-1$.
The retained monomial rows have precisely the three disjoint
weight windows
\begin{equation}\label{R27:windows}
 [-5,5],\qquad [m-5,m+1],\qquad [-m-1,-m+5].
\end{equation}
For $m\geq17$, the central window cannot contain a long-path
term, and either extreme window cannot contain a short-path
term. Possible overlap of the short and long weight bounds
outside these windows is discarded by the projection.

It follows that the new normalized boundary coefficient matrix
is obtained from the old one by the following three injective
monomial maps on its disjoint row windows:
\begin{equation}\label{R27:row_insertion}
 \hbox{central: }f\mapsto xf,\qquad
 \hbox{positive extreme: }f\mapsto yf,\qquad
 \hbox{negative extreme: }f\mapsto zf.
\end{equation}
For the central terms this is the continuant recurrence at an
isolated vertex. A long upper path gains one factor $y$.
A long lower path gains one factor $z$ and the new middle edge
factors; division by $D$ cancels all those factors. The quotient
depends only on nonzero outer edge products, which are unchanged.
Thus the assertion also holds when there are several central
zero edges. New monomial rows outside the images of the three
maps have zero coefficient. The column list is naturally
reindexed, and still has $91$ members.

The three maps have distinct integer-weight ranges and are
injective within those ranges. Hence the normalized matrices
$S_i$ and all three block matrices $T_i$ retain their ranks
under every insertion. The finite certificate therefore gives,
at a normalized cut limit in every size,
\begin{equation}\label{R27:limit_ranks}
 \operatorname{rank}T_0\geq84,\qquad
 \operatorname{rank}T_1\geq172,\qquad
 \operatorname{rank}T_2\geq262.
\end{equation}
To interpret this as a nonzero-edge witness, replace each
central zero edge by $\tau$ times a nonzero constant. The
negative extreme row quotients have already canceled every
middle edge factor, and their remaining denominators are
nonzero outer products. All normalized entries therefore have
the indicated finite limit as $\tau\to0$. Each nonzero limit
minor is a nonzero rational minor for nonzero edges. The rank
conditions \eqref{R27:limit_ranks} consequently hold on nonempty
open subsets of the actual Jacobi parameter space. Normalizing
rows by a nonzero $D$ is invertible and does not alter their
rank or their first and second normal maps.

\paragraph{6. Combining the six directions at one actual pencil.}
Fix a size $m\geq17$. Intersect the three nonempty minor opens
from \eqref{R27:limit_ranks} with the base-rank open of R-26
and with \eqref{R27:Delta}. The parameter space of nonzero
$\alpha_i,\delta_i$ is irreducible, so this finite intersection
is nonempty. At one common actual pencil in this intersection,
the projected matrix has $r_0=84$ and, by
\eqref{R27:first_upper}, $r_1\leq4$. The $T_1$ lower bound
and \eqref{R27:jet_ranks} force $r_1=4$.
The $T_2$ lower bound then gives
\[
 262\leq3\cdot84+2\cdot4+r_2,
 \qquad\hbox{hence }r_2\geq2.
\]
The support argument \eqref{R27:support} identifies these with
four first and at least two second normal classes for the full
matrix. Together with \eqref{R27:base_rank}, the finite-order
rank lemma therefore gives
\[
 \operatorname{rank}_{\mathbb C(\varepsilon)}C(\varepsilon)
       \geq(8m-4)+4+2=8m+2.
\]
Relation \eqref{R27:scaling_actual} gives the reverse bound,
since there are $8m+3$ columns and its coefficient vector is
nonzero. A nonzero maximal minor is a polynomial in
$\varepsilon$, so some complex value of $\varepsilon$ attains
this rank. This proves \eqref{R27:criterion}.

The upper bound on the first normal is indispensable in this
argument. Merely knowing three lower rank bounds at a cut
would not preserve a claimed second-order gain under parameter
variation: a pair of order-two pivots can become pivots of
orders one and three. The independent algebraic bound
\eqref{R27:first_upper} rules out that possibility here.

Finally apply the exact bordered induction of R-25 to
\eqref{R27:criterion}:
\[
 c_{m+1}\geq c_m+8m+2\qquad(m\geq17).
\]
The torus upper bound $4m^2-2m+2$ has precisely the same
increment. Its equality at $m=17$ was certified in R-24,
so induction proves equality for every $m\geq17$.
R-23 identifies the resulting defect with the specified
degree-$m$ first Tor group, proving $\epsilon_m=0$ in that
range. R-13 gives $c_m=\binom{m+3}{3}$ for $1\leq m\leq16$.
These statements, and the two upper bounds, yield exactly
\eqref{R27:dimension}.
\end{proof}


\subsection{R-28: higher-variable determinant dimensions and common free points}

\paragraph{Statement R-28.}
For every $n\geq3$ and $r\geq4$,
\begin{equation}\label{R28:det_dimension}
 \dim\mathcal D_{n,r}=(r-2)n^2+2.
\end{equation}
For $m\geq17$, $r\geq4$, and $n>m$, the padded permanent
image has dimension
\begin{equation}\label{R28:pad_dimension}
 \dim\{\ell^{n-m}q:q\in\mathcal P_{m,r}\}^{\rm cl}
       =rm^2-2m+r+1.
\end{equation}

In the setting of G-22 with $n>m\geq r\geq3$ and
$\nu_r>0$, $\nu_1>(n-m)s$, its two enhanced cones have a
common stable point $w$ with the following stronger property:
\begin{equation}\label{R28:no_scaling}
 g\cdot w=cw,\quad g\in\mathrm{SL}_r,\quad c\in\mathbb C^*
       \quad\Longrightarrow\quad c=1.
\end{equation}
In fact the underlying padded form can be chosen with trivial
projective stabilizer in $\mathrm{PGL}_r$.
If homogeneous generators of $\mathbb C[W]^{\mathrm{SL}_r}$
have positive degrees at most an integer $\beta$, those which
are nonzero at this common point have degrees of greatest common
divisor one. Consequently both multiplicities $m_{j\nu}$ are
positive for every integer $j\geq\beta^2$.
This includes the endpoint $r=m$ and padding exponent one.

\paragraph{Background used.}
Openness of the stable locus and separation of distinct closed
orbits are cited from Mukai (2003), Proposition~5.15, printed
p.~166, and Theorem~5.3 with Corollary~5.4, printed pp.~161--162
(G-B15). The common stable product point and exact invariant
extraction are G-22. The cyclic perturbation, the explicit
irreducibility and stabilizer witnesses, and the extension of
the image dimension are proved here.

\begin{proof}
We work over $\mathbb C$ and treat the dimension and stability statements in turn.

\paragraph{1. Exact determinant dimension after adding variables.}
R-17 gives $\dim\mathcal D_{n,4}=2n^2+2$. The cyclic
cofactor independence established in R-23 also applies to
determinants: each complementary matching at that cyclic
pencil is unique, so its determinant cofactor is the same
polynomial as its permanent cofactor up to a nonzero sign.
Thus the $n^2$ determinant cofactors are independent on a
nonempty open subset of the quaternary parameter space.
Intersect this with the nonempty open where the differential
has rank $2n^2+2$, and choose one pencil $A$ in the intersection.

Regard $A$ as an $r$-variable pencil with zero coefficients
on the additional variables. For each new variable $x_k$,
varying its $n^2$ matrix coefficients gives $x_k$ times the
independent cofactor family. These blocks have dimension
$n^2$ each. Their multidegrees in the new variables are
different from one another and from the old block. The
differential therefore has rank
\[
 2n^2+2+(r-4)n^2=(r-2)n^2+2.
\]
For the opposite inequality, the left--right action of
$\mathrm{SL}_n\times\mathrm{SL}_n$ preserves the determinant
and generically has finite stabilizer on pencils with at
least three coefficient matrices. For example, take one
coefficient to be $I$, a second diagonal with distinct entries,
and a third with every entry nonzero. The first forces the
two acting matrices to be inverses, the second forces the
remaining one to be diagonal, and the third forces it to be
scalar. The special-linear condition then leaves only a
finite central group. The generic fiber thus contains an
orbit of dimension $2n^2-2$, giving the upper bound
$rn^2-(2n^2-2)$. This proves \eqref{R28:det_dimension}.
The argument is valid even when $r>n^2$; it concerns a
coefficient map and does not require that its variables be
obtained by a quotient by linear equations.

The padding argument of R-14 works in $r$ variables. A general
member of $\mathcal P_{m,r}$ is irreducible: this cone contains
the Eisenstein polynomial $x_1^m+x_2^{m-1}x_3$, supplied by
the diagonal-plus-cycle pencil. Irreducibility is a nonempty
open condition on degree-$m$ forms. The projective multiplication
map from $\mathbb P^{r-1}\times\mathbb P(\mathcal P_{m,r})$
is generically injective, since the nonlinear irreducible
factor and the padding line can be recovered separately from
their product. Its affine image therefore has dimension
$\dim\mathcal P_{m,r}+r-1$. Substitute G-21 to obtain
\eqref{R28:pad_dimension}.

\paragraph{2. One common cyclic family has all three needed open sets.}
Put $a=n-m$, retain the fixed extremal vector $v_\eta$ of
G-22, and consider the irreducible affine parameter space of
linear forms $L_1,\ldots,L_m,h,\ell$ and a scalar $\varepsilon$.
Set
\begin{equation}\label{R28:cyclic_family}
 q=\prod_{j=1}^mL_j+\varepsilon h^m,\qquad
 F=\ell^a q,\qquad w=F^{\otimes s}\otimes v_\eta.
\end{equation}
Every $q$ belongs to both the size-$m$ permanent and determinant
images. Indeed, the matrix with diagonal $L_j$, entries $h$
immediately above the diagonal, and corner entry $\varepsilon h$
at $(m,1)$ has exactly two perfect matchings. Its permanent
is $q$; changing that corner to $(-1)^{m-1}\varepsilon h$
gives determinant $q$. Adding $a$ diagonal entries $\ell$
then gives the required degree-$n$ determinant. Thus every
$w$ in this family lies in both enhanced cones.

The inverse image of the stable locus in $W$ is nonempty and
open in this parameter space. It is nonempty at
$\varepsilon=0$ by the Fourier-frame construction in G-22,
using its $L_j$, padding line, and the same fixed $v_\eta$.
We will intersect it with the following two additional
nonempty open conditions on $q$:
\[
 q\text{ is irreducible},\qquad
 \operatorname{Stab}_{\mathrm{PGL}_r}([q])\text{ is finite}.
\]
The first is open because reducible degree-$m$ forms are a
finite union of closed projective multiplication images.
The second is open by the maximal-rank condition on the
coefficient matrix of the infinitesimal equation
$Bq=cq$: its only solutions should be scalar matrices $B$.
Over characteristic zero, zero projective stabilizer Lie
algebra means a finite algebraic stabilizer.

Here is one explicit witness for both conditions in every
$m\geq r\geq3$. In coordinates $x_1,\ldots,x_r$, put
\begin{equation}\label{R28:explicit_q}
 P=x_1x_2\cdots x_{r-1}x_r^{m-r+1},\qquad
 z=x_1+\cdots+x_r,\qquad q_*=P+\varepsilon z^m,
 \quad\varepsilon\ne0.
\end{equation}
This is a specialization of \eqref{R28:cyclic_family}, allowing
repeated diagonal factors. Regard $z,x_1,\ldots,x_{r-1}$ as
independent coordinates. Then
\[
 q_* = \varepsilon z^m+
       (x_1\cdots x_{r-1})
          (z-x_1-\cdots-x_{r-1})^{m-r+1}.
\]
As a polynomial in $z$ over
$\mathbb C[x_1,\ldots,x_{r-1}]$, its leading coefficient is
a unit, every other coefficient is divisible by $x_1$, and
its constant coefficient is divisible by $x_1$ exactly once.
The latter assertion uses $r\geq3$, so
$x_1+\cdots+x_{r-1}$ is not divisible by $x_1$.
Eisenstein at $x_1$ proves irreducibility.

For the stabilizer, suppose a linear infinitesimal operator
$B$ satisfies $Bq_*=cq_*$. No pure monomial $x_i^m$ occurs
in $BP$: all $r\geq3$ variables occur in $P$, and one
infinitesimal replacement can remove at most one of them.
Writing $Bz=\sum b_i x_i$, comparison of the pure monomials
in $\varepsilon m z^{m-1}Bz=c\varepsilon z^m$ gives
$b_i=c/m$ for every $i$. Hence $Bz=(c/m)z$, and the remaining
equation is $BP=cP$. Each off-diagonal matrix entry contributes
a different monomial $P x_i/x_j$, with nonzero coefficient
because every exponent of $P$ is positive. It follows that
all off-diagonal entries vanish. The equation for $Bz$ then
makes every diagonal entry equal to $c/m$. Thus $B$ is scalar,
as required.

The three nonempty open subsets just constructed intersect
because the parameter space is irreducible. Choose a tuple in
their intersection and fix its irreducible $q$ and its auxiliary
vector. Its projective automorphism group $H$ is finite.

\paragraph{3. Varying the padding line removes every projective symmetry.}
At the tuple just chosen, the line $\ell$ belongs to a
nonempty open set of padding lines for which
$(\ell^a q)^{\otimes s}\otimes v_\eta$ is stable.
Every nonidentity element of $H\subset\mathrm{PGL}_r$ has
a proper closed fixed locus on the projective space of linear
forms. Avoid the finite union of these loci while remaining
in that stable open. For this choice of $\ell$, the projective
stabilizer of $F=\ell^a q$ is trivial. Indeed, unique
factorization forces a transformation preserving $[F]$ to
preserve the nonlinear irreducible factor $[q]$ and the linear
factor $[\ell]$ separately. Thus it belongs to $H$ and fixes
$[\ell]$, which forces the identity. This argument does not
require $a\geq2$: the two factors are distinguished by their
degrees even when $a=1$.

If $g\in\mathrm{SL}_r$ satisfies $gw=cw$, the two nonzero
decomposable tensor factors show that $g$ preserves the line
of $F^{\otimes s}$ and hence the line of $F$. The Veronese
map is injective on lines. Therefore $g$ is central, say
$g=tI$ with $t^r=1$. Its action on $W$ is multiplication by
\[
 t^{ns-e}=t^{r\nu_r}=1.
\]
This proves \eqref{R28:no_scaling}. In particular $w$ and
every nonzero scalar multiple of $w$ have closed orbits,
and two different scalar multiples cannot lie in the same
$\mathrm{SL}_r$ orbit.

\paragraph{4. Common invariant degrees and occurrence at every large stretch.}
Let $f_1,\ldots,f_t$ be homogeneous positive-degree generators
of the ambient invariant ring, with degrees at most $\beta$,
and keep the indices for which $f_i(w)\ne0$. There is at least
one: otherwise the closed orbits of $w$ and zero would not
be separated by invariants. If the kept degrees had greatest
common divisor $d>1$, a primitive $d$th root $c\ne1$ would
satisfy $f_i(cw)=f_i(w)$ for every generator, including those
vanishing at $w$. Every invariant would then agree at $w,cw$.
Closed-orbit separation would put them in the same orbit,
contradicting \eqref{R28:no_scaling}. The greatest common
divisor is therefore one.

The elementary numerical-semigroup bound proved in G-18
now applies: positive generators bounded by $\beta$ and of
greatest common divisor one generate every integer at least
$\beta^2$. Products of the kept invariant generators are
nonzero at $w$, and hence nonzero on both enhanced cones.
G-22's exact invariant extraction identifies these degree-$j$
elements with the respective occurrences of $j\nu$ in original
coordinate degree $sj$. This proves the final assertion.
\end{proof}

\paragraph{Scope.}
The first two image dimensions are unconditional in their
stated ranges. The common free point requires the strict
direction hypotheses and $3\leq r\leq m$, as stated.
No conclusion is drawn here about the padding boundary
$\nu_1=(n-m)s$, where a remaining $(r-1)$-variable image can
already have nontrivial multiplicity differences. The result
supplies the absence of an invariant-degree period for the
higher-row comparison, not a positive extraction algorithm
for individual coefficients.


\subsection{R-29: prescribed product moments extend the common stable range to $2m-2$ rows}

\paragraph{Statement R-29.}
Let $n>m\geq3$ and $4\leq r\leq2m-2$. Put $a=n-m$,
and fix a partition $\nu\vdash ns$, $s\geq1$, with exactly
$r$ positive parts and
\[
 \nu_1>as.
\]
Define $\eta=\nu-\nu_r(1,\ldots,1)$ and use the enhanced
cones and ambient representation $W$ of G-22 in rank $r$.
These two cones have a common stable point. They also have a
common stable point $w$ with
\begin{equation}\label{R29:no_scaling}
 g w=cw,\quad g\in\mathrm{SL}_r,\quad c\in\mathbb C^*
       \quad\Longrightarrow\quad c=1.
\end{equation}
Consequently their invariant coordinate rings have the
dimensions and exact single-label extraction in G-22, and
both multiplicities occur at every integer stretch
$j\geq\beta^2$ for any common ambient generator-degree bound
$\beta$. In particular the stable-point and scalar-period
hypotheses of the higher-row Hilbert comparison extend from
$r\leq m$ to $r\leq2m-2$ under the stated strict conditions.

\paragraph{Background and normalization.}
Use the tensor-induced Hermitian metrics on symmetric powers.
For a nonzero degree-$m$ form $f$ in $k$ variables, write
$M(f)$ for its normalized Hermitian moment matrix before
subtracting the scalar part. Thus its expectation on a
Hermitian infinitesimal matrix $B$ is
$\langle Bf,f\rangle/\langle f,f\rangle
 =\operatorname{tr}(BM(f))$, and $\operatorname{tr}M(f)=m$.
Multiplying $f$ by a nonzero scalar does not change $M(f)$.
The traceless part is zero exactly when the usual
$\mathrm{SL}_k$ moment map is zero. The equivalence between
closed orbit and an orbit meeting moment zero is the cited
Kempf--Ness criterion: Losev (2006), Theorem~4, printed p.~1
(G-B14). Hilbert--Mumford and stable-locus openness are the
same cited inputs as G-22 and R-28. We do not assume that
the moment matrix of a product is the sum of its factors'
moment matrices.

\begin{proof}
We first prove a prescribed-moment lemma and then construct
the common tensor from two orthogonal variable blocks.

\paragraph{1. A product of $m$ linear forms can realize any positive rational spectrum in at most $m-1$ variables.}
Let $2\leq k\leq m-1$, and let
$b_1,\ldots,b_k\in\mathbb Q_{>0}$ have sum $m$.
There is a unit degree-$m$ form $P$, a product of $m$
distinct linear forms in $\mathbb C^k$, such that
\begin{equation}\label{R29:prescribed_moment}
 M(P)=\operatorname{diag}(b_1,\ldots,b_k),
\end{equation}
and $[P]$ has finite projective stabilizer.

Permute coordinates first so the $b_i$ are weakly decreasing.
Choose a positive integer $d$ with
$\lambda=(db_1,\ldots,db_k)$ integral, and put
$\theta=\lambda-\lambda_k(1,\ldots,1)$.
For a primitive $m$th root $\zeta$ take the unit Fourier
vectors
\[
 u_j=\sum_{i=0}^{k-1}\sqrt{b_{i+1}/m}\,\zeta^{ij}x_i,
       \qquad 0\leq j<m.
\]
As in G-22, their rank-one moment sum is
$\operatorname{diag}(b_1,\ldots,b_k)$, and every $k$ of
them are independent by the Vandermonde determinant.
Since $m\geq k+1$, they contain a projective frame.
The ordered tensor
\[
 \xi=(u_0\otimes\cdots\otimes u_{m-1})^{\otimes d}
                   \otimes v_\theta
\]
has moment zero for $\mathrm{SL}_k$: its full diagonal moment
is $\lambda-\theta=\lambda_k(1,\ldots,1)$.
Its projective stabilizer is finite, since fixing every
ordered tensor line fixes the projective frame. The cited
Kempf--Ness criterion makes its orbit closed, so $\xi$ is
stable.

Put $P_0=\prod_j u_j$. For every one-parameter subgroup,
the smallest weight of $P_0$ is the sum of the smallest
weights of the $u_j$, because their initial terms have
nonzero product in a polynomial ring. The same statement
holds for the outer power. Thus
$P_0^{\otimes d}\otimes v_\theta$ and $\xi$ have identical
smallest weight for every subgroup. Hilbert--Mumford, in
exactly the form used in G-22, proves that the former
tensor is stable as well.

Apply the reverse direction of Kempf--Ness to this closed
orbit. Some $g\in\mathrm{SL}_k$ sends
$P_0^{\otimes d}\otimes v_\theta$ to moment zero.
The line of $g v_\theta$ is a point of the partial flag
variety parametrizing the extremal orbit. A unitary matrix
can take this flag to the standard flag: choose an
orthonormal basis adapted to its nested subspaces. Hence
there is a unitary $U$ such that the normalized vector
$Ugv_\theta$ is a scalar of absolute value one times
$v_\theta$. A scalar adjustment makes $U$ special unitary
if desired, without changing these lines or moments.

Let $P$ be the unit normalization of $UgP_0$.
Unitary equivariance preserves moment zero, and the pure
outer tensor power contributes $dM(P)$. The auxiliary
factor contributes $-\theta$. Therefore
\[
 dM(P)-\operatorname{diag}(\theta)
           =cI_k.
\]
Taking traces gives $c=\lambda_k$, and proves
\eqref{R29:prescribed_moment}. Undo the initial coordinate
permutation to obtain the prescribed order of the $b_i$.
The factor lines of $P$ still contain a projective frame,
since $Ug$ is invertible. A transformation preserving
$[P]$ permutes its $m$ distinct linear factors, and the
kernel of that finite permutation action fixes the frame
pointwise. It is scalar. Thus the projective stabilizer
of $P$ is finite, completing the lemma.

\paragraph{2. Two blocks give the desired full-rank moment matrix.}
Use coordinates $x_0,\ldots,x_{r-1}$ and put
\[
 \mu_0=\nu_1/s-a,\qquad
 \mu_i=\nu_{i+1}/s\quad(1\leq i<r).
\]
These are positive rational numbers with sum $m$.
Split $V=V_1\oplus V_2$ orthogonally, with
$k_1=\lfloor r/2\rfloor$, $k_2=\lceil r/2\rceil$,
using the first $k_1$ coordinates for $V_1$.
Both block dimensions lie between two and $m-1$.
Define
\[
 w_b=\frac1m\sum_{x_i\in V_b}\mu_i\quad(b=1,2).
\]
Then $w_1,w_2>0$ and $w_1+w_2=1$. Within block $b$,
the numbers $\mu_i/w_b$ have sum $m$. By the lemma choose
unit products $P\in\operatorname{Sym}^mV_1$ and
$Q\in\operatorname{Sym}^mV_2$, each with finite projective
stabilizer, and with those prescribed moment matrices.
Set
\begin{equation}\label{R29:two_block_q}
 q=\sqrt{w_1}P+\sqrt{w_2}Q.
\end{equation}
The two pure-block spaces are orthogonal, so $q$ is unit.
Its block-diagonal moment is the corresponding weighted
sum. An off-diagonal infinitesimal operator has bidegree
$(m-1,1)$ or $(1,m-1)$, orthogonal to both pure blocks.
It follows that
\begin{equation}\label{R29:full_moment}
 M(q)=\operatorname{diag}(\mu_0,\ldots,\mu_{r-1}).
\end{equation}
This is a direct moment computation for a sum in orthogonal
subspaces, not an additive moment assertion for a polynomial
product.

The projective stabilizer of $q$ is finite. First, the
constant-coefficient directional derivatives of each block
product are independent. If $D_vP=0$, restrict the equation
to the zero hyperplane of one factor, away from the others;
it gives that this factor vanishes on $v$. Doing this for
every factor and using their span shows $v=0$.
The argument for $Q$ is identical.

Now suppose $Bq=cq$. Because $m\geq3$, the two mixed
bidegrees $(m-1,1)$ and $(1,m-1)$ are different from
one another and from the pure bidegrees. The off-diagonal
blocks of $B$ consequently vanish, by the derivative
independence just proved. The remaining equations are
$B_1P=cP$ and $B_2Q=cQ$. The finite projective stabilizers
make both $B_i$ scalar, and their degree-$m$ actions force
the same scalar $c/m$ in both blocks. Thus $B$ is scalar.
In characteristic zero this proves finiteness of the
projective algebraic stabilizer of $q$.

Write the two products in \eqref{R29:two_block_q} as
$\prod_{j=1}^m L_j$ and $\prod_{j=1}^m M_j$, absorbing
the two nonzero scalar coefficients into individual factors.
The diagonal-plus-cycle pencil with diagonal $L_j$ and
successive cycle entries $M_j$ has permanent equal to
their sum. Changing the sign of one cycle entry when
necessary gives the same sum as a determinant. Hence
$q$ belongs to both size-$m$ images.

\paragraph{3. Padding and the flag preserve stability by weights.}
Take $\ell=x_0$ and the unit extremal vector $v_\eta$.
In the ordered representation consider
\[
 \widetilde\xi=
  \bigl(x_0^{\otimes a}\otimes q\bigr)^{\otimes s}
                    \otimes v_\eta.
\]
Equations \eqref{R29:full_moment} and
$sa e_0+s\mu-\eta=\nu_r(1,\ldots,1)$ show that its
$\mathrm{SL}_r$ moment is zero. Its projective stabilizer
is finite because it preserves the line of its $q$ factor,
whose projective stabilizer is finite. Thus
$\widetilde\xi$ is stable.

Put $F=x_0^a q$ and $w=F^{\otimes s}\otimes v_\eta$.
For every one-parameter subgroup, the initial term of
$x_0^a q$ is the nonzero product of the respective initial
terms. The minimum-weight identity of G-22 therefore
applies to $\widetilde\xi,w$, and proves stability of $w$.
No moment formula for the padded polynomial product has
been used. Block-diagonal padding proves that $F$ belongs
to both required degree-$n$ cones. This gives the common
stable point throughout the stated row range.

\paragraph{4. The enlarged common family also removes the scalar period.}
Use the irreducible affine parameter space of all linear
forms $L_1,\ldots,L_m,M_1,\ldots,M_m,\ell$, and set
\[
 q=\prod_jL_j+\prod_jM_j,\qquad
 w=(\ell^a q)^{\otimes s}\otimes v_\eta,
\]
with the auxiliary vector fixed. Every point is still
represented by the same two-matching cycle construction.
The stable enhanced locus is nonempty by the preceding
steps, and the finite-projective-stabilizer locus for $q$
is nonempty by \eqref{R29:two_block_q}. Both are open.
The irreducible locus for $q$ is also nonempty: specialize
to $x_1^m+x_2^{m-1}x_3$, Eisenstein at $x_3$ as a
polynomial in $x_1$. It remains irreducible after adjoining
the unused variables, since its quotient polynomial ring
remains an integral domain. This gives a witness inside
the same parameter space even when $r>m$.

These three nonempty open sets intersect. Fix an
irreducible $q$ with finite projective automorphism group
at a stable enhanced point, and vary the padding line
inside the nonempty stable open of that fiber. Avoiding
the finitely many proper fixed loci of the nonidentity
automorphisms makes the padded form projectively free,
exactly as in R-28. Its nonlinear irreducible factor is
distinguished from the padding line even when $a=1$.
The central action has exponent $ns-|\eta|=r\nu_r$,
so the argument of R-28 proves \eqref{R29:no_scaling}.
The same closed-orbit separation and numerical-semigroup
argument gives gcd one and all stretches $j\geq\beta^2$.

Finally the enhanced-cone dimensions and graded invariant
extraction in G-22 only require a common stable point,
the rank-$r$ Cartan coordinate ring, and the stated grading.
Their proofs are independent of $r\leq m$ once stability
is supplied. They therefore apply unchanged in this
larger range.
\end{proof}

\paragraph{Scope.}
The prescribed-moment lemma is an existence construction
using Kempf--Ness normalization; it does not assert a
polynomial bit-complexity algorithm for the required
change of coordinates. The two-block method reaches
$r\leq2m-2$, and makes no claim beyond that range.
For $r>m$ and integer $n>m$, the strict rectangular
condition $n<rm/(r-1)$ is empty. The new row range is
therefore a nonrectangular strict-direction extension;
the first-row equality face is not settled by this result.


\subsection{R-30: tilted padding closes the two-product endpoint range}

\paragraph{Statement R-30.}
Let $n>m\geq3$, let $r\in\{2m-1,2m\}$, and put $a=n-m$.
For every partition $\nu\vdash ns$, $s\geq1$, with exactly
$r$ positive parts and $\nu_1>as$, the two enhanced cones of
G-22 have a common stable point. They have such a point $w$
with the no-scaling property
\[
 g w=cw,\quad g\in\mathrm{SL}_r,\quad c\in\mathbb C^*
       \quad\Longrightarrow\quad c=1.
\]
Thus the common-stability, exact invariant extraction, quotient
 dimension, and all-integer occurrence conclusions of R-29 hold
through $4\leq r\leq2m$ after combining the two results.
The assertion concerns the enhanced tensor: at these endpoints
we do not assert that the underlying padded form itself has
trivial projective stabilizer.

\paragraph{Background.}
Use the tensor-induced metrics and full moment matrix $M$
of R-29. Besides the cited Kempf--Ness criterion, use its
stabilizer consequence at moment zero: the compact stabilizer
Lie algebra is a real form of the complex stabilizer Lie algebra.
The exact inspected locator is Losev (2006), proof of Theorem~1,
printed p.~2, equations~(3)--(4) and the sentence immediately
following (G-B14). In particular a moment-zero vector with
finite compact stabilizer has finite complex stabilizer.
We use this cited fact only for the full ordered padded tensor,
not for its unbalanced polynomial factors.

\begin{proof}
We extend the product normalization at a nonconstant endpoint,
then tilt the padding line so that neither endpoint block has
a constant spectrum.

\paragraph{1. Nonconstant spectra in exactly $m$ variables.}
Let $b_1,\ldots,b_m>0$ be rational, with sum $m$, and suppose
they are not all equal. There is a unit product $P$ of $m$
independent linear forms with
\[
 M(P)=\operatorname{diag}(b_1,\ldots,b_m)
\]
whose projective stabilizer in the unitary group is finite.
Its complex projective stabilizer still contains a torus;
no contrary assertion is needed.

To prove this, sort the $b_i$, clear denominators with $d$,
and put $\lambda=db$ and $\theta=\lambda-\lambda_m\mathbf1$
as in R-29. Now $\theta\ne0$. Take the same Fourier vectors
$u_j$ for $0\leq j<m$; they form a basis. The ordered tensor
\[
 \xi=(u_0\otimes\cdots\otimes u_{m-1})^{\otimes d}
                                      \otimes v_\theta
\]
has moment zero. Its projective stabilizer is finite, as follows.
Write the basis matrix as $DF$, with $D$ positive diagonal and
$F$ the Fourier matrix. An operator preserving each ordered
factor line has the form
\[
 g=DF\operatorname{diag}(t_0,\ldots,t_{m-1})F^{-1}D^{-1}.
\]
Since $\theta\ne0$ and $\theta_m=0$, preserving the auxiliary
extremal line preserves at least one nonzero proper standard
dual flag subspace. On the original space this is a proper
consecutive coordinate subspace. The diagonal $D$ preserves
that subspace, so the circulant matrix
$C=F\operatorname{diag}(t)F^{-1}$ preserves it as well.

Write $C=\sum_{h=0}^{m-1}c_h S^h$ for the cyclic shift $S$.
A nonempty proper consecutive interval in $\mathbb Z/m\mathbb Z$
is not invariant under any nonzero translation: its unique
transition from outside to inside would otherwise have a
nonzero translation period. For each $h\ne0$ there is therefore
an interval index mapped outside by $S^h$. The corresponding
matrix entry of $C$ is $c_h$, and preservation forces $c_h=0$.
Thus $C$ and $g$ are scalar. In $\mathrm{SL}_m$ only finitely
many scalars remain. This proves the asserted finite projective
stabilizer of $\xi$.

Kempf--Ness and the every-subgroup minimum-weight argument in
R-29 show that $P_0^{\otimes d}\otimes v_\theta$, where
$P_0=\prod_j u_j$, is stable. The same reverse Kempf--Ness and
unitary partial-flag normalization give a product $P$ with the
prescribed matrix. To check the new stabilizer assertion, let
$U$ be a unitary projective symmetry of $P$. Moment equivariance
implies that $U$ commutes with $\operatorname{diag}(b)$.
It therefore preserves its eigenspace decomposition and the
extremal line of $v_\theta$. Hence $U$ is also a projective
symmetry of the stable tensor $P^{\otimes d}\otimes v_\theta$.
Modulo scalar unitary matrices, this group is finite.
This proves the endpoint lemma. For block dimensions at most
$m-1$, the stronger complex-stabilizer statement of R-29
remains available.

\paragraph{2. A rational two-eigenvalue tilt.}
Put $\delta_i=\nu_{i+1}/s$ for $0\leq i<r$, so that
$\delta_0>a$ and $\sum_i\delta_i=n$. There is an index $j>0$
with $\delta_j<\delta_0$. Otherwise $\delta_0=n/r$, and
$r>m$ and $a\geq1$ would give
\[
 \delta_0=\frac{m+a}{r}\leq\frac{m+a}{m+1}\leq a,
\]
a contradiction.

Set $A=\delta_0-a>0$ and $B=\delta_j>0$. Choose a sufficiently
small positive rational $\epsilon$ and put
\[
 \alpha=\max(A,B)+\epsilon,\qquad
 \gamma=\min(A,B)-\epsilon.
\]
Require $\gamma>0$ and require both $\alpha,\gamma$ to differ
from every $\delta_i$ with $i\notin\{0,j\}$. Only finitely
many choices of $\epsilon$ are excluded. Since
$\alpha\gamma<AB$, the rational number
\[
 t=\frac{AB-\alpha\gamma}{a(\delta_0-\delta_j)}
\]
lies in $(0,1)$ when $\epsilon$ is small enough. Take the unit
padding line
\[
 \ell=\sqrt{1-t}\,x_0+\sqrt t\,x_j.
\]
The positive Hermitian matrix
\begin{equation}\label{R30:C}
 C=\operatorname{diag}(\delta_0,\ldots,\delta_{r-1})
                                      -a\ell\ell^*
\end{equation}
has eigenvalues $\alpha,\gamma$ in the $(0,j)$ plane and
$\delta_i$ elsewhere. Indeed the plane trace is $A+B$
and its determinant is
\[
 AB-a(\delta_0-\delta_j)t=\alpha\gamma.
\]
All its eigenvalues are positive rational numbers and
$\operatorname{tr}C=m$. The two newly chosen eigenvalues
are distinct from one another and from all the others.

\paragraph{3. Finite compact stabilizers suffice at moment zero.}
Split an orthonormal eigenbasis of $C$ into blocks of dimensions
$k_1=\lfloor r/2\rfloor$ and $k_2=\lceil r/2\rceil$,
putting $\alpha$ in the first block and $\gamma$ in the second.
Both block sizes lie between two and $m$, and each block
contains at least one other eigenvalue. Its spectrum is
therefore nonconstant. Let $w_b$ be its eigenvalue sum divided
by $m$, and divide its eigenvalues by $w_b$. They are positive
rational numbers with sum $m$.

Use paragraph~1 when $k_b=m$, and R-29 when $k_b<m$, to
choose a unit product in each block with this normalized
moment. Call them $P,Q$. Both have finite unitary projective
stabilizers and independent constant directional derivatives.
For
\[
 q=\sqrt{w_1}P+\sqrt{w_2}Q
\]
the orthogonal-block calculation in R-29 gives $M(q)=C$.
The same mixed-bidegree calculation, now restricted to
skew-Hermitian infinitesimal matrices, proves that the unitary
projective stabilizer of $q$ is finite: the off-diagonal blocks
vanish because $m\geq3$ and the derivatives are independent;
the two diagonal blocks are scalar by the compact stabilizer
statements, and their scalars agree because both forms have
degree $m$.

Let $\eta=\nu-\nu_r\mathbf1$ and take the standard unit
extremal vector $v_\eta$. The ordered tensor
\[
 \widetilde\xi=(\ell^{\otimes a}\otimes q)^{\otimes s}
                                             \otimes v_\eta
\]
has full moment
\[
 sa\ell\ell^*+sC-\operatorname{diag}(\eta)
                                               =\nu_r I_r.
\]
Its compact stabilizer is finite because a projective symmetry
of this pure tensor preserves the line of its $q$ factor.
The cited real-form consequence of Kempf--Ness then makes
its complex vector stabilizer finite. Its complex projective
stabilizer has the same Lie algebra: if $B\widetilde\xi
=c\widetilde\xi$ for $B\in\mathfrak{sl}_r$, the zero moment
implies $\langle B\widetilde\xi,\widetilde\xi\rangle=0$
by complex linearity, hence $c=0$. Thus $\widetilde\xi$ is
stable. Multiplication preserves every one-parameter-subgroup
minimum weight, so $(\ell^a q)^{\otimes s}\otimes v_\eta$
is stable. The two-product cycle and block-padding constructions
place it in both enhanced cones.

\paragraph{4. A generic auxiliary flag removes the remaining finite group.}
We now use the irreducible parameter family
\[
 q=\prod_{i=1}^mL_i+\prod_{i=1}^mM_i,\qquad F=\ell^a q,
\]
with all factors and $\ell$ free. The finite-projective-stabilizer
locus for $F$ is a nonempty open in this family. We give explicit
witnesses to establish nonemptiness.

For $r=2m$, take two independent coordinate blocks $x_1,\ldots,x_m$
and $y_1,\ldots,y_m$, let
\[
 P=\prod_{i=1}^m x_i,\qquad Q=\prod_{i=1}^m y_i,
                    \qquad \ell=\sum_i x_i+\sum_i y_i.
\]
For $r=2m-1$, take $m-1$ independent $x$ variables and $m$
independent $y$ variables, and instead set
\[
 P=x_1\cdots x_{m-1}(x_1+\cdots+x_{m-1}),\qquad
 Q=\prod_{i=1}^m y_i,\qquad
 \ell=\sum_i x_i+\sum_i y_i.
\]
In the latter case the $m$ factors of $P$ form a projective frame,
so its projective stabilizer is finite. In both cases $q=P+Q$
is irreducible: as a polynomial in $y_1$ it is primitive and
linear, with coprime leading and constant coefficients
$y_2\cdots y_m$ and $P$. It is irreducible over the fraction
field of the remaining variables, hence by Gauss over their
polynomial ring.

Any infinitesimal projective symmetry of $F$ preserves the
unique nonlinear irreducible factor $q$ and the padding line.
The mixed-bidegree argument makes a symmetry of $q$ block
 diagonal. At $r=2m$ its two blocks are diagonal in the factor
bases; preserving the displayed $\ell$ forces all diagonal
entries to be the same. At $r=2m-1$ the first block is scalar,
and the second is diagonal; the same padding-line condition
again forces one common scalar. Thus the projective stabilizer
Lie algebra of $F$ is zero, proving the required finiteness.

The enhanced stable locus is nonempty by paragraph~3 and open.
Intersect it with the finite-projective-stabilizer open for $F$,
allowing the auxiliary extremal line to vary in its partial
flag orbit. This is one irreducible parameter space, so the
intersection is nonempty. Fix an $F$ in it and let $H$ be its
finite projective automorphism group. Here $\eta\ne0$: if
$\nu$ were rectangular, the same inequality as in paragraph~2
would contradict $\nu_1>as$.

The action of $\mathrm{PGL}_r$ on the nontrivial partial flag
orbit is faithful. Indeed any operator fixing every subspace
of one proper nonzero dimension fixes every line, by intersecting
such subspaces, and is scalar. Consequently each nonidentity
$h\in H$ has a proper closed fixed locus on this orbit. The
stable flags form a nonempty open in the fixed-$F$ fiber.
Choose a flag there outside the finite union of these fixed
loci. If $g w=cw$ at the resulting enhanced point, $g$ must
preserve both $[F]$ and this flag; its projective class is
therefore the identity. Thus $g$ is central in $\mathrm{SL}_r$.
The central exponent is $ns-|\eta|=r\nu_r$, so it acts
trivially and $c=1$.

The exact extraction and dimension arguments of G-22, followed
by the closed-orbit separation and numerical-semigroup argument
of R-28, now give the stated consequences and every integer
stretch $j\geq\beta^2$ for a common ambient generator bound.
\end{proof}

\paragraph{Scope.}
This proof uses finite compact stabilizers at a balanced full
tensor, not a false finite-complex-stabilizer assertion for a
product in $m$ variables. At the final step it removes the
finite padded-form group with the nontrivial auxiliary flag;
it does not claim that a generic two-product padded form is
projectively free. The present proof stops at $2m$ rows. The two sets of $m$
factors span at most $2m$ variables, although an independent
padding line can add one more essential variable. Whether
that additional variable admits the same strict-direction
stability is a separate question; no obstruction beyond the
proved range is asserted.


\subsection{R-31: a scalar-period obstruction in the next-row two-product family}

\paragraph{Statement R-31.}
Let $m\geq3$, $a\geq1$, $n=m+a$, and $r=2m+1$. Define the
common two-product padded cone
\[
 \mathcal T_{m,n,r}
 =\overline{\left\{\ell^a\left(\prod_{i=1}^m L_i+
                         \prod_{i=1}^m M_i\right)\right\}}
                       \subseteq\operatorname{Sym}^n\mathbb C^r.
\]
The factors are arbitrary linear forms. This cone lies in both
projected permanent and determinant cones by the cycle and
padding constructions. Suppose $s,b$ are positive integers with
$ms=rb$, and put
\[
 \nu=(as+b,b^{,r-1}),\qquad \eta=(as,0^{,r-1}).
\]
For every odd $b$ and every odd positive integer $j$,
\begin{equation}\label{R31:odd_vanishing}
 m_{j\nu}\bigl(\mathcal T_{m,n,r}\bigr)=0
       \quad\hbox{in coordinate degree }sj.
\end{equation}
Equivalently the enhanced cone associated to this direction
has no invariant of odd degree. It contains no nonzero closed
orbit with the no-scaling property used in R-28--R-30.

In particular, for every odd $m\geq3$, take $s=r=2m+1$ and
$b=m$. The direction is strict, with $\nu_r>0$ and $\nu_1>as$,
but the same two-product common family cannot supply a stable
no-scaling point at this next row count. This is an obstruction
to that family and that all-integer method; it does not imply
vanishing or a multiplicity comparison for either full orbit
closure.

\begin{proof}
The $2m+1$ factors $L_1,\ldots,L_m,M_1,\ldots,M_m,\ell$
are linearly independent on a nonempty open set of the parameter
space. Their image is dense in $\mathcal T_{m,n,r}$. On this
open set a linear coordinate change identifies the form with
\[
 F=z^a(P+Q),\qquad P=x_1\cdots x_m,\quad Q=y_1\cdots y_m.
\]
The auxiliary extremal line for the one-row weight $\eta$ is
represented by $v^{as}$, where
\[
 v=\sum_{i=1}^m c_i x_i^*+\sum_{i=1}^m d_i y_i^*+e z^*.
\]
It is enough to consider the dense open set with $c_1c_2\ne0$.

Define an operator $h$ on the variable space by
\[
 h(x_1)=\frac{c_1}{c_2}x_2,\qquad
 h(x_2)=\frac{c_2}{c_1}x_1,
\]
and let it fix all other coordinates. Then $h^2=1$,
$\det h=-1$, $hP=P$, $hQ=Q$, and $hz=z$. Moreover
$v\circ h=v$, so the dual action also gives $hv=v$.
Thus $h$ fixes both $F$ and the auxiliary vector.

Choose $\zeta\in\mathbb C^*$ with $\zeta^r=-1$ and set
$g=\zeta h$. Then $g\in\mathrm{SL}_r$ and its action on
\[
 w=F^{\otimes s}\otimes v^{as}
\]
is
\[
 gw=\zeta^{ns-as}w=\zeta^{ms}w
                    =(\zeta^r)^b w=(-1)^b w.
\]
For odd $b$, every homogeneous invariant of odd degree must
therefore vanish at $w$. The generic factor and auxiliary
parameters considered above have dense image in the enhanced
cone, so every such invariant is identically zero on that
cone. The unconditional Cartan graded-coordinate extraction
in G-22 identifies its degree-$j$ invariant dimension with
$m_{j\nu}(\mathcal T_{m,n,r})$, proving
\eqref{R31:odd_vanishing}.

For the no-scaling consequence, let $w\ne0$ have closed orbit
in the enhanced cone. Every invariant takes the same value at
$w$ and $-w$, because its homogeneous odd-degree components
vanish identically. The cited closed-orbit separation theorem
(G-B15, Mukai2003 Theorem5.3 and Corollary5.4) then puts $w$
and $-w$ in the same $\mathrm{SL}_r$ orbit. Hence some group
element scales $w$ by $-1$, which contradicts the no-scaling
property.

Finally $\gcd(m,2m+1)=1$. The displayed choice $s=2m+1$,
$b=m$ satisfies the required equation, and $b$ is odd precisely
when $m$ is odd. The partition has size $as+rb=ns$, and its
first part exceeds $as$ by $b>0$. This proves the explicit
strict-direction obstruction claimed in the statement.
\end{proof}

\paragraph{Scope and distinction from the full cones.}
Restriction from either full coordinate ring to the common
subcone is a quotient. A zero multiplicity on the subcone
therefore gives no zero multiplicity on the full cone.
This result excludes the unchanged two-product no-scaling
strategy in the specified next-row directions, and leaves
other common pencils, even stretches, and other directions
open. No asymptotic comparison of the full permanent and
determinant multiplicities is asserted here.


\subsection{R-32: unpadded common stability and nonrectangular period one}

\paragraph{Statement R-32.}
Let $d\geq3$, $4\leq r\leq2d-2$, and let
$\nu\vdash ds$, $s\geq1$, have exactly $r$ positive parts.
Put $\eta=\nu-\nu_r\mathbf1_r$. The enhanced cones for the
\emph{unpadded} degree-$d$ permanent and determinant images
in $r$ variables have a common stable point. In particular,
their enhanced invariant dimensions equal their actual
image dimensions plus $r_\eta-(r^2-1)$, and the exact
single-label extraction in G-22 applies.

If $\nu$ is nonrectangular, these two cones also have a common
stable point $w$ satisfying
\[
 gw=cw,\quad g\in\mathrm{SL}_r,\quad c\in\mathbb C^*
             \quad\Longrightarrow\quad c=1.
\]
Consequently both unpadded multiplicities at $j\nu$, in
coordinate degree $sj$, are positive at every integer
$j\geq\beta^2$ for any common ambient invariant-generator
bound $\beta$. The rectangular case retains the stable-point
and exact-dimension assertion; no period-one conclusion is
asserted there.

\begin{proof}
Take $V=\mathbb C^r$, with its standard Hermitian metric,
and set $\mu_i=\nu_i/s$. These are positive rational
numbers with sum $d$. Split $V=V_1\oplus V_2$ orthogonally
into blocks of sizes $\lfloor r/2\rfloor$ and
$\lceil r/2\rceil$. Each block has dimension at least two
and at most $d-1$.

Let $w_b$ be the sum of the $\mu_i$ in block $b$, divided
by $d$. Then $w_1,w_2>0$ and $w_1+w_2=1$. The prescribed
product-moment lemma in R-29, with its polynomial degree
there renamed $d$, supplies unit degree-$d$ products $P,Q$
in the respective blocks such that their moment matrices
are the corresponding diagonal lists $\mu_i/w_b$.
Both products have finite complex projective stabilizer
and independent constant directional derivatives.
For
\[
                   q=\sqrt{w_1}P+\sqrt{w_2}Q,
\]
the orthogonal-block calculation in R-29 gives
\[
                  M(q)=\operatorname{diag}(\mu_1,\ldots,\mu_r).
\]
Its complex projective stabilizer is finite as well. Indeed,
the two mixed infinitesimal bidegrees $(d-1,1)$ and
$(1,d-1)$ are distinct for $d\geq3$. Derivative independence
kills both off-diagonal Lie blocks. Finite projective
stabilizers of $P,Q$ make the two remaining blocks scalar,
and the equal polynomial degrees force one common scalar.
The two-product cycle construction of R-29 represents this
same $q$ as a size-$d$ permanent and as a size-$d$ determinant,
with the appropriate cycle sign.

For the standard unit auxiliary extremal vector $v_\eta$,
consider
\[
                    w_0=q^{\otimes s}\otimes v_\eta.
\]
Its full moment is
\[
 sM(q)-\operatorname{diag}(\eta)=\nu_r I_r.
\]
A projective symmetry of this pure tensor preserves $[q]$,
so its projective stabilizer is finite. The cited
Kempf--Ness criterion (G-B14) therefore makes $w_0$ stable.
No padding moment or first-row inequality is involved.
This proof also works when $\eta=0$, proving the stable
assertion for rectangular $\nu$.

Suppose now that $\eta\ne0$, and fix this $q$. Let $H$ be
its finite projective automorphism group. The auxiliary
extremal lines form a nontrivial irreducible partial flag
orbit, on which $\mathrm{PGL}_r$ acts faithfully, as checked
in R-30. Stability is open, and the fixed-$q$ stable flag
locus contains the flag of $w_0$. Each nonidentity element
of $H$ fixes a proper closed subset of the flag orbit.
Choose a stable flag outside their finite union.

If $gw=cw$ for the resulting tensor, the projective class
of $g$ preserves both $[q]$ and this flag and is therefore
the identity. Hence $g$ is central in $\mathrm{SL}_r$.
Its exponent on the enhanced tensor is
\[
                       ds-|\eta|=r\nu_r,
\]
so the center acts trivially and $c=1$. The closed-orbit
separation and numerical-semigroup argument in R-28 now
gives the asserted all-integer occurrence bound. Finally,
the exact coordinate extraction and stable quotient-dimension
calculation of G-22 do not require padding and apply to the
two unpadded cones just used.
\end{proof}

\paragraph{Scope.}
This lemma supplies stable unpadded comparison algebras even
when an available smaller padded form fails its first-row
strict inequality. Its rectangular part supplies a stable
quotient of the stated dimension, while allowing a nontrivial
invariant-degree period. It gives no equality between the two
multiplicity functions and does not extend the row range
beyond $2d-2$.


\subsection{R-33: generic projected permanents have finite projective stabilizer}

\paragraph{Statement R-33.}
For $m\geq17$ and $4\leq q<m^2$, a general form in the
projected permanent cone $\mathcal P_{m,q}$ has finite complex
projective stabilizer. Consequently, for any positive integer
$s$ and any auxiliary Cartan weight $\theta$, the enhanced
cone $X$ built from this image has
\begin{equation}\label{R33:quotient_upper}
 \dim\mathbb C[X]^{\mathrm{SL}_q}
 \leq \dim X-(q^2-1)
 =q m^2-2m-q^2+3+r_\theta.
\end{equation}
The conclusion is a quotient-dimension upper bound. It makes
no assertion that a generic orbit is closed, that a stable
point exists, or that the nonzero invariant degrees have gcd
one.

\begin{proof}
Write $N=m^2$, and consider the coefficient map
\[
 \Phi:(\operatorname{Mat}_m)^q\longrightarrow
               \operatorname{Sym}^m\mathbb C^q,
 \qquad
 (A_1,\ldots,A_q)\longmapsto
             f_A(x)=\operatorname{per}\!\left(\sum_i x_iA_i\right).
\]
By G-21, its generic differential rank is $qN-2m+2$.
Its generic kernel therefore has dimension $2m-2$.
The row-column diagonal variations
\begin{equation}\label{R33:gauge}
 \delta A_i=\operatorname{diag}(u)A_i+A_i\operatorname{diag}(v),
                   \qquad \sum_a u_a+\sum_b v_b=0,
\end{equation}
always lie in that kernel. On the nonempty open set where
one coefficient matrix has no zero entry, these variations
have dimension exactly $2m-2$: the only redundancy is
$u_a=t$, $v_b=-t$ for every $a,b$. Thus on the intersection
with the maximal-rank open they give the entire kernel.

Independently require $A_1,\ldots,A_q$ to be linearly
independent and their span $W\subset\operatorname{Mat}_m$
to have every Pl\"ucker coordinate nonzero in the matrix-unit
basis. This is another nonempty open: all-coordinate Pl\"ucker
nonvanishing is a finite intersection of nonempty minor opens
on the irreducible Grassmannian, and independent frames map
surjectively to that Grassmannian. The indicated parameter
opens therefore have nonempty intersection. Fix a pencil
in this intersection with $f_A\ne0$.

Let $H\in\mathfrak{gl}_q$ be an infinitesimal projective
symmetry, so $Hf_A=c f_A$. Subtract $(c/m)I_q$ from $H$;
the resulting $H_0$ acts trivially on $f_A$, since scalar
variable multiplication acts on a degree-$m$ form by $m$.
The corresponding coefficient variation is a linear
combination of the $A_i$ and lies in $\ker D\Phi_A$.
By \eqref{R33:gauge}, there are $u,v$ with the stated sum
condition such that this variation is obtained by applying
the diagonal operator
\[
 D(E_{ab})=(u_a+v_b)E_{ab}
\]
to every $A_i$. In particular $D(W)\subseteq W$.

The induced operator on the line $\bigwedge^q W$ is scalar.
Since every Pl\"ucker coordinate is nonzero, the sum of the
$D$-weights over every $q$-element subset of the $N$ matrix
units is the same. Because $1\leq q<N$, any two matrix-unit
indices can be exchanged while retaining $q-1$ other indices.
Their individual weights must therefore be equal. Thus
$u_a+v_b=\tau$ for every $a,b$. Summing all $N$ weights gives
\[
 N\tau=m\left(\sum_a u_a+\sum_b v_b\right)=0,
\]
so $D=0$. The $A_i$ are independent, hence the coefficient
variation vanishes only when $H_0=0$. We have proved that
every infinitesimal projective symmetry is scalar.

The projective algebraic stabilizer of $[f_A]$ therefore has
zero Lie algebra and is finite in characteristic zero.
Finiteness is open here, since the Lie stabilizer is the
kernel of a matrix whose entries are polynomial in the form
coefficients. The nonempty open just established meets the
image, so the general point of $\mathcal P_{m,q}$ has this
property.

For the enhanced cone, a projective symmetry of a pure tensor
$F^{\otimes s}\otimes v_\theta$ preserves $[F]$. Its stabilizer
inside $\mathrm{SL}_q$ is therefore finite over the generic
finite-projective-stabilizer locus just proved. Thus the
generic orbit in $X$ has dimension $q^2-1$.
The quotient morphism
\[
 X\longrightarrow\operatorname{Spec}\mathbb C[X]^{\mathrm{SL}_q}
\]
is dominant. Its generic fiber meets the dense finite-stabilizer
open and contains the orbit through such a point. Hence the
generic fiber has dimension at least $q^2-1$.
The generic fiber dimension theorem gives
\eqref{R33:quotient_upper}. Here
$\dim X=\dim\mathcal P_{m,q}+r_\theta$ by the same
Segre--Veronese--Cartan construction as G-22, and G-21 gives
$\dim\mathcal P_{m,q}=q m^2-2m+2$.
\end{proof}

\paragraph{Sources and scope.}
The new argument uses the proved exact dimension and generic
kernel statement from G-21, then elementary Pl\"ucker weights.
The final generic-fiber equality is the Stacks Project,
Lemma~33.20.4(2), tag~\nolinkurl{0B2L} (R-B18), whose stated
proof uses transcendence degrees. The EGA/SGA~3 framework
in dictionary entry D-S3 concerns fiber-dimension upper
semicontinuity and the stabilizer-openness framework; it
is not substituted for this generic-fiber equality.
No equality with the invariant quotient dimension is inferred
from generic freeness. A ring used only for an upper Hilbert
bound can use this inequality even when its stable locus or
scalar period is unknown. The restriction $q<m^2$ is essential
to the all-subset Pl\"ucker exchange argument and is retained.

% END SOURCE work/rectangular/findings.tex

% BEGIN SOURCE work/root/findings.tex
\section{Integral and $\lambda$-ring controls on the bridges}\label{sec:root-findings}
The results in this section are deductions checked in this run. They are not claims of historical priority. We distinguish an identity of characters, a decomposition in a Grothendieck group, an actual direct-sum decomposition, and a dimension identity.

\begin{proposition}[G-01: The three meanings of the second operation]\label{G-01}
Let $E$ be a finite $G$-set. Write $b_2(E)$ for the $G$-set of two-element subsets and $s_2(E)$ for the $G$-set of two-element multisets. In the Burnside ring one has
\[
 [s_2(E)]=[b_2(E)]+[E],\qquad [E]^2=[E]+[\operatorname{Inj}(\{1,2\},E)].
\]
The identity $[E]^2=2[b_2(E)]+[E]$ does not hold in general. If $V$ is a representation over a field of characteristic zero, the correct special $\lambda$-ring identities are
\[
 [\operatorname{Sym}^2V]=[V]^2-[\wedge^2V],\quad
 [\operatorname{Sym}^2V]-[\wedge^2V]=\psi^2[V],\quad
 [V]^2=2[\wedge^2V]+\psi^2[V].
\]
Consequently the expression $2\lambda^2(E)+E=E^2$ in cote 133, scan 019, line 53 is valid after taking cardinalities, but is not a universal equivariant or special-$\lambda$-ring identity. The supplied study's assertion that the two equivariant descriptions agree needs this correction.
\end{proposition}
\begin{proof}
Partition a multiset of cardinality two into the diagonal multisets $\{e,e\}$ and the subsets $\{e,e'\}$ with $e\ne e'$. This partition is $G$-equivariant. Similarly partition ordered pairs into diagonal and off-diagonal pairs. For a counterexample take $G=C_2$ and its regular two-element set $E$. Then $E\times E$ is a disjoint union of two free $C_2$-orbits, whereas $b_2(E)$ is one fixed point. Thus the asserted equality would read $2[E]=[E]+2[\mathrm{pt}]$. Applying the fixed-point homomorphism for the whole group gives $0=2$, a contradiction.
For representations, first prove the universal identity on $\mathrm{GL}_r$ and then pull it back along the representation of $G$. Write the character of its defining representation on the diagonal torus as $x_1+\cdots+x_r$. The symmetric and exterior squares have characters $\sum_{i\leq j}x_ix_j$ and $\sum_{i<j}x_ix_j$. Their sum is $(\sum_i x_i)^2$, and their difference is $\sum_i x_i^2=\psi^2[V]$. Character injectivity in characteristic zero proves the formulas. In the same $C_2$ example the permutation representation is $1\oplus\epsilon$, and $\wedge^2V=\epsilon$, whereas linearization of $b_2(E)$ is $1$. This explicitly shows why the subset operation cannot be silently identified with exterior power. Taking dimensions sends $\psi^2[V]$ to $\dim V$ and recovers the numerical identity. None of this determines the ink on a scan which has not been supplied.
\end{proof}

\begin{proposition}[G-02: The terminal exterior-power variable is necessary]\label{G-02}
For $i\geq2$, a universal formula for $\lambda^i(xy)$ in special $\lambda$-rings cannot in general omit $\lambda^i(x)$ or $\lambda^i(y)$. Thus the terminal index in cote 133, scan 006, lines 25--26 must be $i$ if those lines state the universal product formula.
\end{proposition}
\begin{proof}
Set $y=1$, for which $\lambda^1(1)=1$ and $\lambda^j(1)=0$ for $j>1$. A formula using only $x,\lambda^2(x),\ldots,\lambda^{i-1}(x)$ would express $\lambda^i(x)$ in the subring generated by these earlier operations. Evaluate at the defining representation of $\mathrm{GL}_i$ and restrict to its diagonal torus. This would express the elementary symmetric polynomial $e_i$ in $\mathbb Z[e_1,\ldots,e_{i-1}]$, contrary to the algebraic independence of $e_1,\ldots,e_i$ (equivalently the fundamental theorem of symmetric polynomials). Interchange $x$ and $y$ for the other endpoint. The argument establishes the required mathematical reading, not a new visual transcription.
\end{proof}

\begin{proposition}[G-03: Finite $\gamma$-data do not determine multiplicities]\label{G-03}
For every integer $D\geq0$ there are two honest representations $A_D,B_D$ of the split group scheme $\mathbb G_{m,\mathbb Z}$, each of rank $2^D$, which have the same class modulo $F^{D+1}_\gamma$ but different weight multiplicities. Over $\mathbb Q$ their Chern characters agree in degrees $0,\ldots,D$.
\end{proposition}
\begin{proof}
Let $L$ denote the standard character and put $r=D+1$. Define
\[
 A_D=\bigoplus_{\substack{0\leq j\leq r\\j\text{ even}}}L^{\otimes j\,\oplus\binom rj},\qquad
 B_D=\bigoplus_{\substack{0\leq j\leq r\\j\text{ odd}}}L^{\otimes j\,\oplus\binom rj}.
\]
The two binomial parity sums are $2^{r-1}$. In the representation ring $\mathbb Z[L,L^{-1}]$ the difference of the classes is $(1-L)^r$, a product of $r$ rank-zero elements and hence an element of $F^r_\gamma$. The constant weight has multiplicity one in $A_D$ and zero in $B_D$, so the characters differ. With $c=c_1(L)$, the difference of the Chern characters is $(1-e^c)^r$, whose lowest possible degree is $r$. This proves the assertion. For a line bundle on a scheme, applying the associated-bundle map gives the identical statement in its $K_0$. This does not assert that the two resulting bundles are nonisomorphic on every scheme.
\end{proof}

\begin{proposition}[G-04: An integral coefficient model and its exact specialization gate]\label{G-04}
Let $W_{\mathbb Z}$ be a finite free module, let $\rho:\mathrm{GL}_{N,\mathbb Z}\to\mathrm{GL}(W_{\mathbb Z})$ be a homogeneous polynomial representation of positive degree $e$, and let $f\in W_{\mathbb Z}$. Define $X_{\mathbb Z}$ as the affine schematic closure of its orbit. Its coordinate algebra is torsion-free, hence flat over $\mathbb Z$. In degree $d$ its orbit-coefficient map is an explicit finite integer matrix
\[
 A_d:\operatorname{Sym}^d(W_{\mathbb Z}^*)\longrightarrow
 \mathbb Z[g_{11},\ldots,g_{NN}]_{\leq ed},\qquad P\longmapsto P(\rho(g)f).
\]
Its rank over $\mathbb Q$ is the degree-$d$ coordinate-ring dimension. If $\delta_d$ is any nonzero maximal-rank minor of $A_d$, then this dimension and the degree-$d$ kernel commute with base change over $\mathbb Z[1/\delta_d]$. At a prime dividing $\delta_d$, rank must be tested anew; the fiber of the integral closure need not be the schematic orbit closure of the reduced point.
\end{proposition}
\begin{proof}
The coordinate homomorphism of the orbit map has target $\mathbb Z[\mathrm{GL}_N]=\mathbb Z[g_{ij},\det(g)^{-1}]$, a torsion-free $\mathbb Z$-module. Its kernel $I$ defines the schematic closure and its image $\mathbb Z[W]/I$ is torsion-free. A module over the principal ideal domain $\mathbb Z$ is flat if it is torsion-free, proving the first assertion. Since the representation is polynomial, each degree-$d$ homogeneous coordinate polynomial maps to a polynomial of degree at most $ed$ in the entries of $g$. Equality on invertible matrices is equality in the polynomial ring because localization by the nonzero polynomial $\det(g)$ is injective. This proves the finite coefficient-matrix description.
Let $r=\operatorname{rank}_{\mathbb Q}A_d$. All $(r+1)$-minors vanish as integers, while an $r$-minor $\delta_d$ is nonzero. Over $\mathbb Z[1/\delta_d]$, use that invertible block for row and column elimination. The vanishing of larger minors makes the remaining Schur complement zero, so the matrix is equivalent to $\operatorname{diag}(I_r,0)$. Its image and kernel are direct summands there and commute with every base change. Without inversion, an inclusion of lattices need not remain injective modulo a prime: already multiplication by $p$ on $\mathbb Z$ has rank one over $\mathbb Q$ and rank zero modulo $p$. Thus no unsupported all-primes orbit-fiber conclusion follows. Homogeneity sends degree $d$ to degree $ed$, so the kernel is homogeneous and the stated grading is well defined. If $r=0$, take $\delta_d=1$.
\end{proof}

\begin{proposition}[G-05: Ternary saturation and the three-row exclusion]\label{G-05}
Work over $\mathbb C$, and let $n\geq2$, $N=n^2$, and $\mathrm{Det}_n\subseteq\operatorname{Sym}^n(\mathbb C^N)^*$ be the determinant orbit closure. Then $\mathrm{Det}_n$ contains every degree-$n$ form supported on a fixed three-dimensional subspace. Consequently, for every partition $\lambda\vdash nd$ with $\ell(\lambda)\leq3$,
\[
 \operatorname{mult}_{S_\lambda\mathbb C^N}\mathbb C[\mathrm{Det}_n]_d
 =\operatorname{mult}_{S_\lambda\mathbb C^N}\operatorname{Sym}^d(\operatorname{Sym}^n\mathbb C^N).
\]
No multiplicity obstruction with at most three rows can separate any other orbit closure of a degree-$n$ form in the same ambient space, including a padded permanent, from $\mathrm{Det}_n$.
\end{proposition}
\begin{proof}
We use the established linear determinantal representation theorem for smooth plane curves: a smooth plane degree-$n$ curve together with a line bundle $L$ of degree $g-1$ with $H^0(L)=0$ admits an $n\times n$ matrix of linear forms whose determinant is its defining form up to a nonzero scalar (Beauville, 2000, Proposition 3.1(a)). Such an $L$ exists over $\mathbb C$: for $g\geq1$, the effective locus in $\operatorname{Pic}^{g-1}$ is the image of the symmetric power $C^{(g-1)}$, of dimension at most $g-1$, whereas $\operatorname{Pic}^{g-1}$ has dimension $g$; for $g=0$, take the degree-$-1$ line bundle. Scaling one row absorbs the nonzero scalar. Hence every smooth ternary degree-$n$ form is obtained by a linear substitution in $\det_n$. Any endomorphism substitution is a limit of invertible substitutions because $\mathrm{GL}_N$ is Zariski dense in $\operatorname{End}(\mathbb C^N)$; continuity of the polynomial action therefore puts these forms in $\mathrm{Det}_n$. Smooth ternary forms form a dense open set, so its closedness puts every ternary form there.
Let $W=\operatorname{Sym}^d(\operatorname{Sym}^n\mathbb C^N)$, and choose the Borel convention whose highest-weight vector has nonnegative weight $\lambda=(\lambda_1,\lambda_2,\lambda_3,0,\ldots,0)$. In a monomial basis of $W$, zero weight in positions $4,\ldots,N$ forces every coefficient variable to involve only the first three basis vectors. Thus the entire highest-weight space $W_\lambda^{\mathrm{hw}}$ is contained in $\operatorname{Sym}^d(\operatorname{Sym}^n\mathbb C^3)$. A nonzero polynomial in these coefficient variables cannot vanish on all ternary forms. By the first paragraph, the degree-$d$ determinant ideal therefore meets $W_\lambda^{\mathrm{hw}}$ trivially. Exactness and complete reducibility in characteristic zero give the asserted multiplicity equality. Every other closed subvariety has a degree-$d$ coordinate ring which is a quotient of $W$, so its multiplicity cannot exceed this common ambient value. This proves the exclusion. The assertion is about complex characteristic-zero multiplicities, not an unqualified statement in positive characteristic.
\end{proof}

\begin{proposition}[G-06: The four-variable gate is nonvacuous]\label{G-06}
For $n\geq4$, the closure $\mathcal D_{n,4}$ of degree-$n$ quaternary forms which are determinants of $n\times n$ matrices of linear forms has dimension at most $2n^2+2$. Since
\[
 \binom{n+3}{3}-(2n^2+2)=\binom{n-1}{3}>0,
\]
this is a proper subvariety of the space of quaternary forms. The determinant ideal in $\operatorname{Sym}(\operatorname{Sym}^n\mathbb C^{n^2})$ therefore has a nonzero isotypic component with a partition of exactly four rows in some finite degree.
\end{proposition}
\begin{proof}
Consider the polynomial map
\[
 (A_1,A_2,A_3,A_4)\longmapsto\det(x_1A_1+x_2A_2+x_3A_3+x_4A_4)
\]
from an affine space of dimension $4n^2$. The group
$H=\{(P,Q)\in\mathrm{GL}_n^2:\det(P)=\det(Q)\}$ acts by $A_i\mapsto PA_iQ^{-1}$, preserving this map. It has dimension $2n^2-1$. Its generic stabilizer has dimension one. Indeed, for a tuple with $A_1=I$, a diagonal $A_2$ with pairwise distinct eigenvalues, and an $A_3$ with every off-diagonal entry nonzero, a stabilizer has $Q=P$; commuting with $A_2$ makes $P$ diagonal, and commuting with $A_3$ makes it scalar. The scalar pairs $(cI,cI)$ always stabilize every tuple. Upper semicontinuity of stabilizer dimension, or the rank calculation for the commutation equations on this exhibited open set, gives generic dimension one. Thus generic fibers contain orbits of dimension $2n^2-2$, and the dimension-of-fibers theorem yields the stated image bound. The displayed identity follows by expanding the binomial coefficient.
Let $\pi$ set all but the first four variables to zero. For every form in $\mathrm{Det}_n$, its image under $\pi$ lies in $\mathcal D_{n,4}$: approximate the form by invertible substitutions of $\det_n$, apply $\pi$, and use the closedness of $\mathcal D_{n,4}$. Choose a nonzero homogeneous equation $P$ of this proper cone. Then $P\circ\pi$ is a nonzero homogeneous equation of $\mathrm{Det}_n$, using only first-four coefficient variables. The ideal of $\mathcal D_{n,4}$ is $\mathrm{GL}_4$-stable; choose a highest-weight vector in a nonzero irreducible submodule of its degree-$d$ part. Embedded in the full coordinate ring, it is also a $\mathrm{GL}_{n^2}$ highest-weight vector: raising operators involving a variable beyond the fourth annihilate it. Its highest weight has at most four rows, and G-05 rules out at most three. Hence its length is exactly four. Existence of such an equation does not establish that it is nonzero on the padded permanent.
\end{proof}

\begin{proposition}[G-07: An exact four-row determinant-closure formula]\label{G-07}
Set $D=320112$. For every $d\geq0$ and partition $\lambda\vdash4d$ with $\ell(\lambda)\leq4$, padded by zeros to length four, the degree-$d$ determinant-closure multiplicity satisfies
\[
 m_{\mathrm{Det}_4}(\lambda,d)
 =a_\lambda(d[4])-a_{\lambda-(D,D,D,D)}((d-D)[4]).
\]
The second term is defined to be zero if $d<D$ or $\lambda_4<D$. In particular the four-row sector is ambient for $d<D$, and at $d=D$ the multiplicity of $\lambda=(D^4)$ is exactly one less than its ambient plethysm multiplicity.
\end{proposition}
\begin{proof}
Leal, Lozano Huerta, and Vite (2024), Theorem 2 and Table 2, identify the closure of linear determinantal quartic surfaces as an irreducible reduced hypersurface of degree $D$ in $\mathbb P(\operatorname{Sym}^4\mathbb C^4)^*$. Table 2 identifies this component $F_1$ by $a=(5,5,5,5)$, $b=(6,6,6,6)$, so every matrix entry has degree one. Its affine cone has a prime principal ideal $(F)$ in the polynomial ring $A=\operatorname{Sym}(\operatorname{Sym}^4\mathbb C^4)$, with $F$ homogeneous of degree $D$.
The line $\mathbb CF$ is stable under $\mathrm{GL}_4$. Every algebraic one-dimensional character of $\mathrm{GL}_4$ is a power of the determinant. Since scalar $zI$ acts on $A_D$ by $z^{4D}$, this character is $\det^D$. Thus multiplication by $F$ gives the exact sequence of $\mathrm{GL}_4$-modules
\[
 0\longrightarrow \det^D\otimes A_{d-D}\xrightarrow{\ \cdot F\ }A_d
 \longrightarrow(A/(F))_d\longrightarrow0.
\]
For $d<D$ its first term is zero.
Let $\pi$ restrict forms in sixteen variables to their first four variables. The image $\pi(\mathrm{Det}_4)$ is precisely the affine cone just described. One inclusion follows by continuity from the orbit of $\det_4$. For the converse, every four-variable linear determinant is an endomorphism substitution in $\det_4$, hence belongs to its orbit closure; the same is true of every limit of such forms, and such a form is fixed by $\pi$. It follows that the intersection of the full determinant ideal with the first-four coefficient polynomial ring is exactly $(F)$. Highest-weight spaces for partitions of at most four rows use only these coefficient variables, by the weight argument of G-05. Applying complete reducibility to the exact sequence and shifting a highest weight by the determinant character proves the formula. At $d=D$ the source is $\det^D\otimes\mathbb C$, so the deficit is one.
\end{proof}
\paragraph{Source correction for G-07.}
The attached Landsberg source, section \nolinkurl{zhanglee}, theorem \nolinkurl{zlthm}, records $640224$. The later result used here is the degree of the reduced prime divisor. Leal--Lozano Huerta--Vite, section 4.1 immediately before Theorem 2, explicitly gives $P_{20,2}=\tfrac12D_{23,14}$. This is the precise cycle-multiplicity distinction; using the unreduced Noether--Lefschetz coefficient as the polynomial degree doubles the answer. This correction concerns the cited literature, not a purported visual correction of cote 133.

\begin{proposition}[G-08: A uniform low-variable determinantal construction]\label{G-08}
Let $m\geq1$ and $k\geq2$, put $a=\lfloor k/2\rfloor$, $b=\lceil k/2\rceil$, and
\[
 B(m,k)=\binom{m+a}{a}+\binom{m+b}{b}-1.
\]
Over any commutative ring, every degree-$m$ homogeneous polynomial $q$ in $k$ variables admits a matrix of linear forms of size $B(m,k)$ with determinant
\[
 \ell^{B(m,k)-m}q,
\]
where $\ell$ may first be an independent variable and then be specialized to any linear form. The construction is polynomial in the coefficients of $q$ and defined over $\mathbb Z$.
Consequently, over $\mathbb C$, whenever $n\geq B(m,k)$, no multiplicity obstruction of length at most $k$ can separate the orbit closure of an arbitrarily padded degree-$m$ form $\ell^{n-m}f$ from $\mathrm{Det}_n$.
\end{proposition}
\begin{proof}
Split the variables into disjoint lists $x=(x_1,\ldots,x_a)$ and $y=(y_1,\ldots,y_b)$, and write
\[
 q=\sum_{|\alpha|+|\beta|=m}c_{\alpha\beta}x^\alpha y^\beta.
\]
Make a rooted directed tree with one vertex for every $x$-monomial of degree at most $m$. The parent of a nonconstant monomial is obtained by deleting one occurrence of its largest-index variable; orient the edge away from the constant root and label it by that variable. There is exactly one root-to-vertex path, with label product $x^\alpha$. Make the corresponding $y$-tree with every edge reversed, so the unique path from vertex $\beta$ to its constant root has product $y^\beta$. Connect vertex $\alpha$ in the first tree to vertex $\beta$ in the second by an edge of constant label $c_{\alpha\beta}$ exactly when $|\alpha|+|\beta|=m$. Every source-to-sink path uses exactly one connecting edge; its label product is $c_{\alpha\beta}x^\alpha y^\beta$. Therefore its path polynomial is exactly $q$. The graph is acyclic and has
\[
 s=\binom{m+a}{a}+\binom{m+b}{b}=B(m,k)+1
\]
vertices.
Choose a topological ordering with source first and sink last, and let $A$ be its strictly upper triangular weighted adjacency matrix. Then $\det(I-A)=1$ and $(I-A)^{-1}=I+A+\cdots+A^{s-1}$. The $(1,s)$ entry is the source-to-sink path polynomial $q$. The adjugate identity therefore expresses $q$, up to the fixed sign $(-1)^{s+1}$, as the determinant of the $(s-1)\times(s-1)$ minor of $I-A$ obtained by deleting row $s$ and column $1$. Change the sign of one row if needed. Its entries are affine linear in $x,y$.
Replace each constant entry $c$ of this minor by $c\ell$. Its determinant is the degree-$B(m,k)$ homogenization of $q$, hence is $\ell^{B(m,k)-m}q$. This identity is universal over $\mathbb Z[c_{\alpha\beta},x,y,\ell]$. Padding the matrix with $n-B(m,k)$ diagonal entries $\ell$ proves the corresponding size-$n$ statement.
For the multiplicity assertion, restrict any invertible substitution of $\ell^{n-m}f$ to $k$ coordinate variables. The result has the form $L^{n-m}q$, with $q$ a homogeneous degree-$m$ form in these variables. The construction puts this result in $\mathrm{Det}_n$. Continuity gives the same inclusion for the entire padded orbit closure. If $P$ is any highest-weight polynomial of length at most $k$ vanishing on $\mathrm{Det}_n$, it uses only coefficient variables supported on the first $k$ underlying variables (in the sense made precise in G-05), and therefore vanishes on the other orbit closure as well. Thus the highest-weight space of the determinant ideal embeds in that of the padded-form ideal. Subtracting their dimensions from the common ambient highest-weight multiplicity gives $m_{\mathrm{padded}}(\lambda,d)\leq m_{\mathrm{Det}_n}(\lambda,d)$ in every degree.
\end{proof}
\paragraph{Quantified implication.}
For $k=4$, the construction gives $B(m,4)=(m+1)(m+2)-1$. More generally $B(m,k)=O_k(m^{\lceil k/2\rceil})$. Thus for a fixed number of rows this sector is excluded once the padding exceeds the displayed bound. This is an explicit restricted lower-bound obstruction theorem, not a general separation result and not an efficient solution of arbitrary plethysm.

\begin{proposition}[G-09: A sharper low-row exclusion by shared exponent chains]\label{G-09}
Let $m\geq1$ and $k\geq3$. Define the following integer bounds. If $k=2r+1$ is odd, set
\[
 C(m,k)=2\binom{m+r}{r}-1.
\]
If $k=2r+2\geq4$ is even, for $1\leq h\leq m+1$ put $Q=\lfloor m/h\rfloor$, $s_0=\min(h-1,m)$ and
\[
\begin{aligned}
 C_h(m,k)&=\sum_{q=0}^{Q}\binom{m-qh+r}{r}+Q(h-1)
          +\sum_{s=0}^{s_0}\binom{m-s+r}{r}-1,\\
 C(m,k)&=\min_{1\leq h\leq m+1}C_h(m,k).
\end{aligned}
\]
Over a field, for every nonzero linear form $L$ and homogeneous degree-$m$ form $q$ in $k$ variables, $L^{C(m,k)-m}q$ has a linear determinantal representation of size $C(m,k)$. In coordinates with $L=x_0$ the construction is defined over $\mathbb Z$ and polynomial in the coefficients of $q$. For fixed $k\geq3$,
\[
 C(m,k)=O_k\bigl(m^{(k-1)/2}\bigr).
\]
Consequently, over $\mathbb C$, if $n\geq C(m,k)$, there is no multiplicity obstruction with at most $k$ rows separating any padded degree-$m$ form from $\mathrm{Det}_n$. This strengthens G-08 in the specified range.
\end{proposition}
\paragraph{Prior-result boundary.}
The asymptotic determinantal-representation order used here is established: Boralevi, van Doornmalen, Draisma, Hochstenbach, and Plestenjak (2017), Theorem 1.3, prove order $m^{v/2}$ for a fixed number $v\geq2$ of affine variables, with a matching generic lower bound. Their proof in section 4 uses connected monomial spaces and an even/odd binary-digit split. We do not claim that order as new. The proof below gives the displayed finite integer bound by a different shared-chain construction, records its coefficient-ring scope, and proves the resulting GCT multiplicity exclusion. These are the forms needed here that the cited theorem does not state.
\begin{proof}
Extend $L$ to a basis of linear forms and write it as $x_0$. Dehomogenize $q$ by setting $x_0=1$, obtaining an arbitrary polynomial $p$ of degree at most $m$ in $k-1$ variables. It suffices to construct an affine linear determinant of size $C(m,k)$ equal to $p$. Homogenizing each matrix entry by $x_0$ then gives $x_0^{C(m,k)-m}q$.

For odd $k=2r+1$, split the dehomogenized variables into lists $x$ and $y$ of $r$ variables each. Use the two monomial trees from G-08, each with $\binom{m+r}{r}$ vertices, and connect every pair $(\alpha,\beta)$ with $|\alpha|+|\beta|\leq m$ by its coefficient in $p$. Every source-to-sink path is unique for its two monomial indices, so its path polynomial is $p$. The cofactor construction in G-08 has size one less than the total number of vertices, giving the displayed bound.

For even $k=2r+2$, write the variables as $x_1,\ldots,x_r,y_1,\ldots,y_r,z$ and fix $h$. Decompose the exponent of $z$ uniquely as $e=qh+s$, $0\leq s<h$. On the left, make a directed chain with $Qh$ edges, all labelled $z$, beginning at the source. At its vertex of index $qh$ attach the outward $x$-monomial tree of degrees at most $m-qh$, identifying its constant root with that chain vertex. Every such tree vertex indexed by $(q,\alpha)$ has a unique source path, with product $z^{qh}x^\alpha$. Chain vertices whose indices are not multiples of $h$ are transit vertices and receive no connecting edges. The left graph has
\[
 (Qh+1)+\sum_{q=0}^Q\left[\binom{m-qh+r}{r}-1\right]
 =\sum_{q=0}^Q\binom{m-qh+r}{r}+Q(h-1)
\]
vertices.

On the right, make a chain with vertices $0,\ldots,s_0$, directed toward vertex $0$, the sink, and again label its edges by $z$. At vertex $s$ attach the inward $y$-monomial tree of degrees at most $m-s$, identifying the root with that vertex. Every tree vertex indexed by $(s,\beta)$ has a unique sink path with product $y^\beta z^s$. The right graph has $\sum_{s=0}^{s_0}\binom{m-s+r}{r}$ vertices. Connect left vertex $(q,\alpha)$ to right vertex $(s,\beta)$ with the constant coefficient of $x^\alpha y^\beta z^{qh+s}$ in $p$ when $|\alpha|+|\beta|+qh+s\leq m$.

All connecting edges go from left to right. Thus every source-to-sink path uses exactly one such edge, and the uniqueness of the two tree paths and the division $e=qh+s$ gives exactly the required monomial with its coefficient, once. The path polynomial is $p$. Taking the same cofactor of $I-A$ as in G-08 yields an affine linear determinant of size $C_h(m,k)$. This graph has a path capable of realizing $z^m$ and a connecting edge, so its number of vertices is at least $m+2$; hence $C_h(m,k)\geq m+1$. The odd bound is also greater than $m$. Homogenization therefore uses a nonnegative exponent. Minimizing over the finite choices of $h$ proves the stated representation bound.

For the asymptotic statement in the even case take $h=\lceil\sqrt m\rceil$. The number of binomial summands is at most $m/h+1+h\leq2\sqrt m+2$, each bounded by $\binom{m+r}{r}=O_r(m^r)$, while $Q(h-1)\leq m$. Thus $C(m,k)=O_r(m^{r+1/2})$. The odd case is immediate from its binomial formula. These exponents are $(k-1)/2$ in both cases.

For the GCT consequence restrict a substituted padded form to any fixed $k$ underlying variables. If its restricted padding $L$ is nonzero, apply the determinant construction and pad with $n-C(m,k)$ diagonal copies of $L$. If $L=0$, the form is zero because $n\geq C(m,k)>m$. Thus every restriction lies in $\mathrm{Det}_n$. Apply the highest-weight ideal argument from G-08, which uses exactly this containment, to obtain the multiplicity inequality in every degree.
\end{proof}
\paragraph{Integral scope and growth of the remaining gate.}
With $L=x_0$ the graph and cofactor identity are universal over the coefficient ring. Extending an arbitrary nonzero $L$ to coordinates is a field operation; over a general base one instead works on a chart where a coefficient of $L$ is invertible and chooses the resulting local basis. No global polynomial choice of that basis is asserted. If $n=m^c$ with fixed $c>0$, every fixed $k\geq3$ satisfying $(k-1)/2<c$ is excluded for sufficiently large $m$. This is a quantitative restriction on possible multiplicity obstructions, not a claim that the remaining higher-row obstruction exists.
\begin{proposition}[G-10: An explicit degree bound for the four-row dimension gaps]\label{G-10}
For $n\geq4$ put
\[
 r_n=2n^2+2,\qquad B_n=n^{r_n-1},\qquad
 d_n=r_n(B_n-1)+1.
\]
For every integer $d\geq d_n$, some partition $\lambda\vdash nd$
with exactly four rows satisfies
\[
 m_\lambda(\operatorname{per}_n)>m_\lambda(\det_n).
\]
More generally, the same explicit threshold $d_n$ applies to every
padded pair $(m,n)$ for which the dimension criterion
$\dim\mathcal P_{m,4}+3>r_n$ of R-14 has been established.
The degree bound is very large and does not specify which partition
has the gap. It makes the preceding ``sufficiently large degree''
assertions effective; it gives no larger range of padded matrix sizes.
\end{proposition}
\begin{proof}
Write $Y=\mathbb P(\mathcal D_{n,4})$, $s=\dim Y$, and
$\Delta=\deg Y$. The pencil coefficient map projectivizes to a
rational map
\[
 \phi:\mathbb P^{4n^2-1}\dashrightarrow Y
\]
given by homogeneous polynomials of degree $n$. Its image is dense.
We first check the degree estimate
\begin{equation}\label{G10:image-degree}
 \Delta\leq n^s\leq n^{r_n-1}=B_n.
\end{equation}
Choose $s$ general hyperplanes in the target. Their intersection
with $Y$ consists of $\Delta$ distinct points in a dense open subset
contained in the actual image of $\phi$. Such an open subset exists
because the dominant image is constructible. Each point therefore
has a nonempty fiber in the domain of definition $U$ of $\phi$.
The common zero locus $Z$ of the $s$ pulled-back degree-$n$
equations in $\mathbb P^{4n^2-1}$ meets $U$ in the disjoint union
of these fibers. Every irreducible component of $Z$ meeting $U$
maps to just one of the target points. At least $\Delta$ distinct
components of $Z$ are consequently required. Each has degree at
least one, while repeated projective B\'ezout gives that the sum
of degrees of all components of $Z$ is at most $n^s$.
This proves the first inequality. Components contained in the
base locus only contribute to that upper bound and cannot cover
any of the nonempty fibers in $U$.
The second inequality follows from G-06, which gives $s\leq r_n-1$.
This is the usual image-degree argument with its base locus
accounted for explicitly; compare Sombra (1997), section~1,
printed preprint p.~5, and Proposition~2.10, pp.~11--12,
whose displayed statement treats an everywhere regular map.

The all-degree Hilbert bound of Chardin (1989), the unnumbered
central theorem on printed p.~306, applies to the homogeneous
prime ideal of the reduced irreducible variety $Y$ over
$\mathbb C$. Thus for every $d\geq1$,
\begin{equation}\label{G10:det-Hilbert}
 \dim\mathbb C[\mathcal D_{n,4}]_d
 \leq\Delta\binom{d+s}{s}
 \leq B_n\binom{d+r_n-1}{r_n-1}.
\end{equation}
We use this cited bound rather than assuming that the determinant
coordinate ring is Cohen--Macaulay.

Let $X$ be either $\mathcal P_{n,4}$ or a padded image
$\mathcal Z_{m,n}$ satisfying the stated criterion. It is a
reduced irreducible standard graded affine cone and has dimension
at least $r_n+1$, by R-15 or R-14 respectively. Linear Noether
normalization over the infinite field $\mathbb C$ supplies
$\dim X$ algebraically independent linear coordinate functions.
Any $r_n+1$ of these generate a polynomial subalgebra of its
coordinate ring. Counting its degree-$d$ monomials gives
\begin{equation}\label{G10:per-Hilbert}
 \dim\mathbb C[X]_d\geq\binom{d+r_n}{r_n}.
\end{equation}
The ratio of this binomial to the one in
\eqref{G10:det-Hilbert} is $(d+r_n)/r_n$.
When $d\geq r_n(B_n-1)+1$, that ratio is strictly greater
than $B_n$. Hence the Hilbert function of $X$ is strictly larger
in every such degree.

Decomposing the two degree-$d$ coordinate modules into
$\mathrm{GL}_4$ irreducibles, their dimension difference is
\[
 \sum_{\lambda\vdash nd,\,\ell(\lambda)\leq4}
 \bigl(m_\lambda(X)-m_\lambda(\mathcal D_{n,4})\bigr)
 \dim S_\lambda\mathbb C^4>0.
\]
At least one multiplicity difference is positive. The
highest-weight restriction argument of R-14 identifies these
multiplicities with those of the original orbit closures.
G-05 excludes positive differences with at most three rows.
The remaining partition has exactly four rows, as asserted.
\end{proof}
\paragraph{Effectivity and its limit.}
At the displayed degree, the finite coefficient and isotypic-rank
algorithms of R-05--R-06 can enumerate the partitions and find one
with a positive difference. The proof therefore bounds a finite
search, but its size is not polynomial. In the quartic case the
explicit family at degree $320112+s$ in R-12 is substantially
stronger than this uniform bound.


\begin{proposition}[G-11: Explicit rectangular multiplicity gaps]\label{G-11}
For $n\geq4$ define the following positive integers:
\[
\begin{gathered}
 b_n=\binom{n+3}{3},\qquad \beta_n=6b_n n^{30},\qquad
 T_n=\binom{b_n+\beta_n}{b_n}-1,\\
 E_n=4T_n(\beta_n!),\qquad u_n=2n^2-13,\\
 Q_n=(nE_n)^{u_n-1},\qquad K_n=u_n(Q_n-1)+1.
\end{gathered}
\]
For every integer $k\geq K_n$, put $d=E_n k$ and
$\lambda=(nd/4,nd/4,nd/4,nd/4)$. Then
\begin{equation}\label{G11:rectangle}
 m_\lambda(\operatorname{per}_n)>m_\lambda(\det_n).
\end{equation}
More generally, let $n>m\geq4$. The same conclusion, with
$\operatorname{per}_n$ replaced by
$\ell^{n-m}\operatorname{per}_m$, holds if
\begin{equation}\label{G11:padded-hypotheses}
 \dim\mathcal P_{m,4}+3>2n^2+2
 \quad\hbox{and}\quad n<4m/3.
\end{equation}
In particular it holds for $n=m+1$ and every $m\geq8$, and
for $n=\lfloor\rho m\rfloor$, all sufficiently large $m$,
and any fixed $1<\rho<\sqrt5/2$.
The degree, its progression, and the rectangular partition are
explicit. No useful computational size or polynomial algorithm
is asserted for these bounds.
\end{proposition}
\paragraph{Background used.}
We use the Kempf--Ness closed-orbit criterion in Losev (2006),
Theorem 4, p.~1; the stable-locus and quotient statements of
Mukai (2003), Definition 5.12, Proposition 5.15, Theorem 5.16,
and Corollary 5.17, pp.~165--167; and Derksen (2001),
Theorem 1.1, Proposition 1.2, and Example 2.1, printed p.~956.
The latter gives the displayed bound $\beta_n$ for generators
of the invariant ring of quaternary degree-$n$ forms.
The argument below proves the particular stability witnesses,
Veronese bound, and GCT deduction that those sources do not state.
\begin{proof}
Write $G=\mathrm{SL}_4(\mathbb C)$ and let $V_n$ be the space
of degree-$n$ quaternary forms. First we exhibit stable points
in the cones whose dimensions we compare. The form
\[
 f_n=x_0^n+x_1^n+x_2^n+x_3^n
\]
belongs to both $\mathcal P_{n,4}$ and $\mathcal D_{n,4}$.
Indeed factor $x_0^n+x_1^n$ and $x_2^n+x_3^n$ into binary
linear factors, place them on the diagonal and directed cycle
of an $n\times n$ pencil, and use its two perfect matchings.
For the determinant, multiply one cycle entry by the sign
$(-1)^{n-1}$.

Equip $V_n$ with an $\mathrm{SU}_4$-invariant Hermitian inner
product and use unit monomial vectors. Distinct monomials are
orthogonal. In the normalized sum of the four pure powers,
the mean exponent of each variable is $n/4$. Its moment map
therefore vanishes on the diagonal trace-zero Lie algebra.
It also vanishes on the off-diagonal directions: an operator
$x_j\partial/\partial x_i$, $i\ne j$, cannot carry one of
these pure powers to another when $n\geq3$. Thus the moment
map is zero. The Kempf--Ness criterion makes the $G$-orbit
closed. Its infinitesimal stabilizer is zero, since the
monomials $x_i^{n-1}x_j$ and $x_i^n$ in the derivative of
$f_n$ are distinct for $n\geq3$ and force all matrix entries
to vanish. Its stabilizer is consequently finite. This is a
stable point in each cone.

For the padded case put $a=n-m>0$ and suppose $3a<m$.
Let $e_0$ be the unit monomial $x_0^n$ and $e_i$ the unit
monomial proportional to $x_0^a x_i^m$, $1\leq i\leq3$.
Set
\[
 \theta_0=\frac{m-3a}{4m},\qquad
 \theta_1=\theta_2=\theta_3=\frac{n}{4m},\qquad
 F=\sum_{i=0}^3\sqrt{\theta_i}\,e_i.
\]
All four numbers are positive and sum to one. The mean
exponent of $x_0$ is $n\theta_0+3a\theta_1=n/4$, and
that of $x_i$ is $m\theta_i=n/4$. No two supported monomials
differ by a root $e_i-e_j$, because their differences are
multiples of $m\geq4$. The same moment-map calculation
therefore gives zero. We have $F=x_0^a q$, where
$q=\sum_{i=0}^3 c_i x_i^m$ and all $c_i\ne0$.
The preceding binary-factor construction puts $q$ in
$\mathcal P_{m,4}$, so $F\in\mathcal Z_{m,n}$.

Here is a direct finite-stabilizer check. For a linear vector
field $D_A=\sum_i L_i\partial/\partial x_i$, the equation
$D_A(F)=0$, after division by $x_0^{a-1}$, is
\[
 aL_0q+x_0D_Aq=0.
\]
Reduction modulo $x_0$ shows $L_0=d_0x_0$, because $a\ne0$
and $q(0,x_1,x_2,x_3)\ne0$. It follows that
$D_Aq=-ad_0q$. The distinct mixed monomials
$x_i^{m-1}x_j$ force every off-diagonal entry of $A$ to
vanish. The $x_0^m$ coefficient gives $(m+a)d_0=0$,
and the remaining pure powers give every other diagonal
entry zero. Thus the Lie stabilizer is zero and $F$ is
stable by Kempf--Ness.

For every reduced irreducible cone $X$ just considered,
the stable locus is a nonempty open subset and its quotient
fibers are the $15$-dimensional $G$-orbits (Mukai, loc. cit.).
Hence the dimension-of-fibers theorem yields
\begin{equation}\label{G11:quotient-dimension}
 \dim\mathbb C[X]^G=\dim X-15.
\end{equation}
In particular, writing $A_D=\mathbb C[\mathcal D_{n,4}]^G$
and $A_X=\mathbb C[X]^G$, G-06, R-15, and the first
hypothesis in \eqref{G11:padded-hypotheses} imply
\[
 \dim A_D\leq u_n,\qquad \dim A_X\geq u_n+1.
\]
The invariant functor is exact in characteristic zero.
Thus both invariant rings are graded quotients of
$\mathbb C[V_n]^G$, including when $X$ is padded; see also
Mukai, Proposition 5.10 and its proof, pp.~164--165.
Derksen's bound, specialized to $q=4$ and symmetric degree
$n$, says that this ambient invariant ring is generated in
degrees at most $\beta_n=6b_n n^{30}$. The dual convention
for forms makes no difference: the contragredient
transforms the action by an automorphism of $\mathrm{SL}_4$.
We may take at most $T_n$ homogeneous generators, because
the vector space of all coefficient polynomials in $b_n$
variables of positive degree at most $\beta_n$ has dimension
$T_n$.

We check explicitly that the $E_n$th Veronese of each ring
is generated in its first degree. More generally, suppose
an algebra has $r\leq T_n$ homogeneous generators with
positive degrees $d_i\leq\beta_n$, and put
$B_0=\beta_n!$. In a monomial in these generators of total
degree $kE_n$, remove as many pure blocks of
$B_0/d_i$ copies of generator $i$ as possible. Every removed
block has degree $B_0$. The remaining exponents are less
than $B_0/d_i$, so the residual degree is less than
$rB_0\leq T_n B_0$. It is an integer multiple $hB_0$,
because the original and removed degrees are such multiples;
thus $0\leq h<T_n$. There are $4kT_n-h$ pure blocks.
Combine the residual monomial with $4T_n-h$ of these blocks
to form a monomial of degree $E_n=4T_nB_0$, then divide
the other $4(k-1)T_n$ blocks into groups of $4T_n$.
This factors the original monomial into $k$ monomials of
degree $E_n$. Applying this to spanning monomials proves
the claim, also after passing to a quotient.

Consequently $A_D^{(E_n)}$ and $A_X^{(E_n)}$ are standard
graded domains, with dimensions unchanged. Set
$s=\dim\operatorname{Proj}A_D^{(E_n)}\leq u_n-1$.
This projective dimension is positive: block padding embeds
$\mathcal D_{4,4}$, whose dimension is $34$ by G-07,
into $\mathcal D_{n,4}$. By \eqref{G11:quotient-dimension}
the latter invariant ring has dimension at least $19$,
so $s\geq18$.
Choose a basis of its degree-one piece represented by
homogeneous invariant polynomials of degree $E_n$ on $V_n$.
Their composition with the determinant-pencil map is a
collection of homogeneous polynomials of degree $nE_n$.
They give a dominant rational map from
$\mathbb P^{4n^2-1}$ to $\operatorname{Proj}A_D^{(E_n)}$.
Dominance follows from dominance of the pencil map and
surjectivity of the affine quotient (Mukai, Theorem 5.9);
the Veronese morphism has this projective image. G-10's
base-locus-aware image-degree argument gives
\[
 \deg\operatorname{Proj}A_D^{(E_n)}
 \leq(nE_n)^s\leq Q_n.
\]
Chardin's bound in G-10 therefore gives, for every $k\geq1$,
\[
 \dim(A_D)_{E_nk}
 \leq Q_n\binom{k+u_n-1}{u_n-1}.
\]
Linear Noether normalization of the standard graded ring
$A_X^{(E_n)}$, whose dimension is at least $u_n+1$, gives
\[
 \dim(A_X)_{E_nk}\geq\binom{k+u_n}{u_n}.
\]
The ratio of the binomials is $(k+u_n)/u_n>Q_n$ when
$k\geq K_n$. Thus the first invariant dimension is strictly
smaller than the second.

Finally, in coefficient degree $d$, every polynomial
$\mathrm{GL}_4$ weight has size $nd$. The only irreducible
polynomial representation with trivial restriction to
$\mathrm{SL}_4$ is a determinant power; it has shape
$(nd/4)^4$, when $4\mid nd$, and is one-dimensional.
Here $E_n$ is divisible by four. The invariant dimension
difference is therefore exactly the difference of the two
multiplicities for the displayed rectangular partition.
The highest-weight restriction argument of G-05 and R-14
identifies these with the full orbit-closure multiplicities.
This proves \eqref{G11:rectangle} and its padded version.
R-14 and R-16 supply the stated particular ranges; all
satisfy $n<4m/3$.
\end{proof}
\paragraph{Relation to the rectangular bridge.}
The Haar projector onto $\mathrm{SU}_4$ invariants isolates
precisely these rectangular highest weights. Equivalently,
the corresponding character integral is the dimension of
this invariant graded piece. The proof makes that positive
rank exceed the determinant rank; it does not replace an
orbit-closure multiplicity by an ambient Kronecker
coefficient. The quartic formula G-07/R-12 is much sharper
where it applies. These deductions make no literature-priority
claim and are independent of RH and either proposed GRH input.

\begin{proposition}[G-12: A padding exclusion for the rectangular sector]\label{G-12}
Let $m\geq1$, $n>m$, and $n\geq4m/3$. For every degree-$m$
form $f$, every coordinate-ring degree $d\geq1$ with
$4\mid nd$, and $\lambda=(nd/4)^4$, one has
\begin{equation}\label{G12:exclusion}
 m_\lambda(\ell^{n-m}f)\leq m_\lambda(\det_n).
\end{equation}
If $n>4m/3$, the multiplicity on the left is zero.
The equality case $n=4m/3$ is included in the inequality.
Thus the rectangular gaps of G-11 cannot extend to or beyond
this padding boundary, even though nonrectangular four-row
questions remain.
\end{proposition}
\begin{proof}
As in G-05, restrict substituted forms to four variables.
They have shape $L^a q$, where $a=n-m$ and $q$ has degree
$m$. If $L=0$, the form is zero. Otherwise a change of
coordinates takes $L$ to $x_0$; it suffices to use a
$\mathrm{GL}_4$ change, since the determinant cone and each
homogeneous invariant ideal are stable under this group.
Act on the variables by the one-parameter subgroup
\[
 \eta(t)=\operatorname{diag}(t^3,t^{-1},t^{-1},t^{-1})
 \quad\hbox{of }\mathrm{SL}_4.
\]
A monomial in $x_0^a q$ with $x_0$ exponent $a+j$ has
weight $3a-m+4j$. If $3a>m$, all these weights are positive,
so $\eta(t)x_0^a q\longrightarrow0$ as $t\longrightarrow0$.
Every homogeneous positive-degree invariant therefore
vanishes on every such padded form and on their orbit
closure. Exactness of invariants identifies its invariant
coordinate ring with $\mathbb C$, whose positive-degree pieces vanish,
which proves the zero assertion.

If $3a=m$, the same limit instead equals
\[
 F_0=x_0^a q(0,x_1,x_2,x_3).
\]
Every ternary degree-$m$ form belongs to $\mathcal D_{m,3}$
by G-05 (including the elementary degrees one and two).
Embedding it and adding $a$ diagonal copies of $x_0$
shows $F_0\in\mathcal D_{n,4}$; the argument also holds for
limits because the block-padding map is polynomial.
Let $P$ be any homogeneous $\mathrm{SL}_4$ invariant
vanishing on $\mathcal D_{n,4}$. Since $P$ is constant
along $\eta$, continuity gives
$P(x_0^a q)=P(F_0)=0$. Thus the determinant ideal intersected
with the ambient invariant ring is contained in the ideal
of the restricted padded orbit closure intersected with
that ring. Both invariant coordinate rings are quotients
of the same ambient invariant ring, by exactness of
invariants. This containment gives a surjection from the
determinant invariant coordinate ring to the padded one,
and hence the degreewise dimension inequality.
Invariant dimensions are precisely the rectangular
multiplicities, as in G-11; highest-weight restriction
transfers the inequality to the full orbit closures.
\end{proof}
\paragraph{Source and scope.}
The strict inequality is also the rectangular specialization
of the usual padding weight constraint, but the proof above
gives its one-parameter-subgroup form and handles the
boundary by the ternary determinant theorem. This is an
exclusion for this specified rectangular sector. It does
not compare all four-row multiplicities or arbitrary
Kronecker coefficients.

\begin{proposition}[G-13: An explicit finite bound for the first rectangular gap]\label{G-13}
For $n\geq4$ put
\[
 \beta=6\binom{n+3}{3}n^{30},\qquad u=2n^2-13,\qquad
 C=(n\beta^2)^u,
\]
and define
\begin{equation}\label{G13:first-bound}
 s_*=(u+1)(C-1)+1,\qquad d_* = \beta s_*.
\end{equation}
There is an integer $d$ with $1\leq d\leq d_*$ and $4\mid nd$
such that, for $\lambda=(nd/4)^4$,
\[
 m_\lambda(\operatorname{per}_n)>m_\lambda(\det_n).
\]
The same conclusion holds with $\operatorname{per}_n$ replaced by
$\ell^{n-m}\operatorname{per}_m$ whenever $n>m\geq4$ and
\[
 \dim\mathcal P_{m,4}+3>2n^2+2,
 \qquad n<4m/3.
\]
In particular it applies to $n=m+1$, $m\geq8$, and, by R-18,
to $n=\lfloor\rho m\rfloor$ for all sufficiently large $m$
and each fixed $1<\rho<\sqrt{3/2}$.

The bound satisfies $\log d_*=O(n^2\log n)$. This is an
existence bound within a finite interval. It does not specify
the individual successful degree or replace G-11's assertion
at every degree of its displayed progression.
\end{proposition}

\paragraph{Background and scope.}
We use the stable points, exact invariant quotients, and quotient
dimensions established in G-11, together with Derksen (2001),
Theorem~1.1 and Example~2.1, printed p.~956. We also use the
projective image-degree argument of G-10 and Chardin's
unnumbered main theorem (1989), printed p.~306. The argument
below compares a weighted degree filtration with the ordinary
degree filtration in an affine generator embedding; it needs
no factorial Veronese.

\begin{proof}
Let $G=\mathrm{SL}_4(\mathbb C)$ and let $V_n$ be the space
of quaternary forms of degree $n$. Write
\[
 A_D=\mathbb C[\mathcal D_{n,4}]^G,
 \qquad A_X=\mathbb C[X]^G,
\]
where $X$ is $\mathcal P_{n,4}$ or the padded cone in the
statement. The stable-point and dimension arguments of G-11
give
\begin{equation}\label{G13:dims}
 t:=\dim A_D\leq u,\qquad \dim A_X\geq u+1.
\end{equation}
These are positively graded domains, with degree-zero part
$\mathbb C$, and both are graded quotients of
$\mathbb C[V_n]^G$. Choose homogeneous generators
$g_1,\ldots,g_r$ of the latter ring with positive degrees
$e_i\leq\beta$. Their images generate both quotient rings.

First bound the cumulative graded dimension of $A_D$. The
images of the generators give a closed affine embedding
\[
 Y_D=\operatorname{Spec} A_D\lhook\joinrel\longrightarrow
 \mathbb A^r.
\]
Compose this embedding with the determinant-pencil map and
the affine quotient. The resulting map from
$\mathbb A^{4n^2}$ has coordinates
$g_i(\det L)$, homogeneous of degrees $ne_i\leq n\beta$
in the pencil parameters. Its image is dense in $Y_D$.
To homogenize it, adjoin a source coordinate $z$ and use
the degree-$n\beta$ coordinate polynomials
\[
 z^{n\beta},\qquad
 z^{n\beta-ne_i}g_i(\det L),\quad 1\leq i\leq r.
\]
They define a rational map from $\mathbb P^{4n^2}$ whose
image closure is the ordinary projective closure
$\overline{Y_D}\subseteq\mathbb P^r$. It has dimension
$t$. G-10's image-degree bound, allowing a base locus, gives
\begin{equation}\label{G13:affine-degree}
 \deg\overline{Y_D}\leq(n\beta)^t\leq(n\beta)^u.
\end{equation}

For $d\geq0$, let
\[
 H_D^{\leq}(d)=\sum_{j=0}^d\dim(A_D)_j.
\]
Every element of $(A_D)_j$ has a preimage homogeneous of
weighted degree $j$ in $\mathbb C[T_1,\ldots,T_r]$ with
$\deg T_i=e_i$. Since every $e_i\geq1$, each monomial in
that preimage has ordinary degree at most $j$. It follows
that the direct sum of the graded pieces through $d$ is
contained in the space of functions on $Y_D$ represented
by ordinary polynomials of degree at most $d$.
This is an inclusion of filtered vector spaces, not an
identification of the two gradings.

The latter ordinary filtered dimension equals the
degree-$d$ Hilbert function of $\overline{Y_D}$:
homogenization with its first coordinate, followed by
dehomogenization on the dense affine chart, identifies the
two quotients. The projective closure is integral over
$\mathbb C$. Also $t\geq19$ by the fixed block-padding and
stable-quotient observation in G-11, so the positive-dimension
hypothesis in Chardin's stated theorem holds. That theorem
and \eqref{G13:affine-degree} yield
\begin{equation}\label{G13:upper}
 H_D^{\leq}(d)
 \leq(n\beta)^t\binom{d+t}{t}
 \leq(n\beta)^u\binom{d+u}{u}.
\end{equation}

For the other ring, a maximal algebraically independent
subset of the images of $g_1,\ldots,g_r$ has size
$\dim A_X$. Indeed those images generate its fraction
field, and maximality makes every remaining generator
algebraic over the chosen subset. By \eqref{G13:dims},
choose $u+1$ algebraically independent images. Each has
degree at most $\beta$. Their monomials of ordinary total
degree at most $s$ are linearly independent and have
graded degree at most $\beta s$. Consequently
\begin{equation}\label{G13:lower}
 H_X^{\leq}(\beta s)
 :=\sum_{j=0}^{\beta s}\dim(A_X)_j
 \geq\binom{s+u+1}{u+1}.
\end{equation}

For every nonnegative integer $s$, the exact binomial ratio
satisfies
\begin{align*}
 \frac{\binom{s+u+1}{u+1}}{\binom{\beta s+u}{u}}
 &=\frac{s+u+1}{u+1}
     \prod_{i=1}^u\frac{s+i}{\beta s+i}\\
 &\geq\frac{s+u+1}{(u+1)\beta^u}.
\end{align*}
The last inequality follows term by term from
$\beta(s+i)\geq\beta s+i$. At $s=s_*$ one has
\[
 s_*+u+1=(u+1)C+1,
\]
so the final lower bound is strictly greater than
$C/\beta^u=(n\beta)^u$. Equations
\eqref{G13:upper}--\eqref{G13:lower} therefore give
\[
 H_X^{\leq}(d_*)>H_D^{\leq}(d_*).
\]
The degree-zero dimensions are both one. Hence for some
$1\leq d\leq d_*$ the individual graded pieces satisfy
$\dim(A_X)_d>\dim(A_D)_d$.

As in G-11, the only polynomial $\mathrm{GL}_4$
representation of weight size $nd$ that is trivial on
$\mathrm{SL}_4$ is the one-dimensional determinant power
of shape $(nd/4)^4$. If $4\nmid nd$, both invariant
pieces are zero. Thus the successful degree necessarily
satisfies $4\mid nd$, and the invariant dimension gap is
the desired rectangular multiplicity gap. The first-four-
variable highest-weight argument in G-05/R-14 transfers
it to the full orbit closures.

Finally $\beta=O(n^{33})$ and $u=O(n^2)$, so the explicit
formula \eqref{G13:first-bound} gives
$\log d_*=O(n^2\log n)$. The particular padded ranges
follow from R-14/R-16 and R-18, respectively; their
dimension gaps and $n<4m/3$ are the same hypotheses already
used to construct the stable padded point in G-11.
\end{proof}

\paragraph{Complexity lower-bound context.}
The multiplicity statements G-10--G-13 locate particular representation
sectors and degrees. They do not improve the established quadratic
lower bound $\overline{\operatorname{dc}}(\operatorname{per}_m)
\geq m^2/2$ of Landsberg, Manivel, and Ressayre (2013);
the inspected April 2010 author version states it as Theorem~1.0.1,
p.~1 (G-B16). No claim about the strongest currently proposed lower
bound is needed here. A noncontainment result by itself does not
specify the four-row or rectangular multiplicity differences proved
in the present restricted parameter ranges.

\begin{proposition}[G-14: First-jet dominance and a stronger padding boundary]\label{G-14}
Let $m\geq1$, and write a quaternary form as
\[
 q(x_0,x_1,x_2,x_3)=\sum_{j=0}^m x_0^j q_j(x_1,x_2,x_3).
\]
The first-jet map
\[
 j_1:\mathcal D_{m,4}\longrightarrow
   \operatorname{Sym}^m\mathbb C^3\oplus
   \operatorname{Sym}^{m-1}\mathbb C^3,
 \qquad q\longmapsto(q_0,q_1),
\]
has Zariski-dense image.

Consequently let $n>m\geq1$, put $a=n-m$, and let $f$ be
any degree-$m$ form in the common ambient GCT setting.
For every $d\geq1$ and every partition $\lambda\vdash nd$
with at most four rows and
\begin{equation}\label{G14:boundary}
 \lambda_1\leq ad+1,
\end{equation}
one has
\[
 m_\lambda(\ell^a f)\leq m_\lambda(\det_n).
\]
If $\lambda_1<ad$, the multiplicity on the left is zero.
Thus every positive multiplicity gap in this sector must
satisfy $\lambda_1\geq ad+2$.
\end{proposition}

\paragraph{Scope and background.}
The proof uses the coefficient grading and first-four-variable
highest-weight restriction of G-05, together with elementary
cyclic cofactors and the finite geometric-product identity
already proved in R-16. It requires characteristic zero.
The jet-dominance calculation is given in full; no higher-jet
dominance or Lefschetz assertion is assumed.

\begin{proof}
The case $m=1$ of the first assertion is immediate. Suppose
$m\geq2$, choose a primitive $m$th root of unity $\zeta$,
and let $A$ be the ternary cyclic pencil with diagonal
entries
\[
 \ell_i=x_1-\zeta^i x_2,\qquad 0\leq i<m,
\]
and entries $x_3$ in positions $(i,i+1)$, with indices
modulo $m$. Its determinant has the two perfect matchings
and equals
\[
 x_1^m-x_2^m+(-1)^{m-1}x_3^m.
\]
We claim that the cofactors of $A$ span every ternary form
of degree $m-1$.

For each $0\leq k\leq m-1$, the cofactors using exactly
$k$ cycle entries are, up to nonzero signs, $x_3^k$ times
the $m$ cyclic translates of a consecutive product of
$h=m-1-k$ of the binary forms $\ell_i$. This follows from
the unique matching after deleting the relevant row and
column. A consecutive product, after a cyclic shift,
has the form
\[
 F_h(x_1,x_2)=\prod_{i=0}^{h-1}(x_1-\zeta^i x_2).
\]
Every coefficient in it is nonzero. Indeed its coefficient
of $x_1^{h-r}x_2^r$ is
\[
 (-1)^r\zeta^{r(r-1)/2}
   \prod_{j=1}^r\frac{1-\zeta^{h-r+j}}{1-\zeta^j},
 \qquad 0\leq r\leq h,
\]
with the empty product interpreted as one. Every exponent
in a nonempty numerator or denominator lies between $1$
and $h<m$. Thus no factor is zero. The cyclic substitution
$x_2\mapsto\zeta x_2$ has distinct characters on the
monomials with $0\leq r\leq h$. Fourier projection of the
$m$ translates therefore spans every binary form of degree
$h$. Varying $k$ proves the claim, since
\[
 \operatorname{Sym}^{m-1}\mathbb C^3
 =\bigoplus_{k=0}^{m-1}x_3^k
        \operatorname{Sym}^{m-1-k}\langle x_1,x_2\rangle.
\]

Now consider the parameter map
\[
 (A',B)\longmapsto j_1\det(A'+x_0B),
\]
where $A'$ is an arbitrary ternary linear pencil and $B$
is a constant $m\times m$ matrix. At the point $(A,0)$,
the derivative in the $A'$ directions changes $q_0$ by
arbitrary linear-form multiples of the cofactors of $A$.
By the claim these span all of
$\operatorname{Sym}^m\mathbb C^3$. It does not change
$q_1$, since the coefficient of $x_0$ is linear in $B$.
The derivative in the $B$ directions changes $q_1$ by
arbitrary constant linear combinations of the cofactors,
which span $\operatorname{Sym}^{m-1}\mathbb C^3$, and
does not change $q_0$. The differential is consequently
surjective onto the direct sum. The characteristic-zero
Jacobian criterion for dominance proves the first assertion.

For the multiplicity consequence, work first in four
variables and use the positive coefficient weight convention
of G-05/G-11 (equivalently, dualize the coordinate-ring
representations consistently). In a polynomial irreducible
of shape $\lambda$, every weight has each coordinate at
most $\lambda_1$. For example, the Schur construction uses
$\lambda_1$ exterior-power columns, and a given basis vector
can occur at most once in each column. This proves the
required bound on the $x_0$ coordinate of every weight.

Let $R=\mathbb C[\operatorname{Sym}^n\mathbb C^4]$ be
the coefficient polynomial ring, and let $R_{d,\lambda}$
be its degree-$d$ $\lambda$-isotypic subspace. Take
\[
 P\in I(\mathcal D_{n,4})\cap R_{d,\lambda}.
\]
The weight bound says that every coefficient monomial
occurring in $P$ has total $x_0$ weight at most
$\lambda_1$. Evaluate $P$ at a form $x_0^a q$. A
coefficient coming from $q_j$ has $x_0$ weight $a+j$;
a product of $d$ coefficients therefore has weight
\[
 ad+j_1+\cdots+j_d.
\]
All coefficient variables of smaller individual $x_0$
weight vanish on $x_0^a q$. If $\lambda_1<ad$, nothing
survives and every polynomial in $R_{d,\lambda}$ vanishes
on every such padded form.

Under \eqref{G14:boundary}, the restriction
$P(x_0^a q)$ depends only on $q_0,q_1$, because every
surviving coefficient monomial has
$j_1+\cdots+j_d\leq1$. It is a polynomial function of
these first-jet data. For every $q\in\mathcal D_{m,4}$,
block padding gives $x_0^a q\in\mathcal D_{n,4}$,
including the closure points. Hence this jet polynomial
vanishes on the dense image of $j_1$ just proved, and
therefore vanishes for every pair $(q_0,q_1)$. It follows
that $P(x_0^a q)=0$ for every quaternary form $q$ of
degree $m$.

To allow any padding form $L$, use a $\mathrm{GL}_4$
change taking $L$ to $x_0$. The determinant ideal and
the isotypic subspace are both stable under that group,
so the preceding conclusion applies to the corresponding
translate of $P$. Thus $P$ vanishes on every $L^a q$,
and in particular on the restricted padded orbit closure
associated with $\ell^a f$.
We have proved the containment of the determinant-ideal
$\lambda$-isotypic part in the corresponding padded ideal
part. Passing to the characteristic-zero semisimple
degree-$d$ quotients gives the stated multiplicity
inequality. The first-four-variable highest-weight
restriction identifies these multiplicities with those
in the full ambient orbit closures, as in G-05/R-14.
The same argument gives the zero assertion when
$\lambda_1<ad$.
\end{proof}

\paragraph{A precisely stated higher-jet continuation.}
For $0\leq h\leq m$, suppose that the projection
\[
 j_h:\mathcal D_{m,4}\longrightarrow
   \bigoplus_{j=0}^h\operatorname{Sym}^{m-j}\mathbb C^3
\]
has dense image. The identical coefficient-weight argument
then proves the same multiplicity inequality whenever
$\lambda_1\leq ad+h$: every surviving coefficient
monomial depends only on those jets. This is a conditional
continuation with an explicit geometric hypothesis.
The first-jet case above is proved without that hypothesis.
No dominance assertion for $h=2$ or $h=3$ is made here.

\begin{proposition}[G-15: Uniform third-jet dominance]\label{G-15}
For every $m\geq1$, the map
\[
 j_{\leq3}:\mathcal D_{m,4}\longrightarrow
   \bigoplus_{j=0}^{\min\{3,m\}}
          \operatorname{Sym}^{m-j}\mathbb C^3,
 \qquad \sum_{j=0}^m t^j q_j\longmapsto(q_0,q_1,q_2,q_3),
\]
with nonexistent components omitted, has Zariski-dense image.
For $m\geq4$, order three is the largest order for which the
full jet map can be dominant.

Consequently, if $n>m\geq1$, $a=n-m$, $d\geq1$, and
$\lambda\vdash nd$ has at most four rows, then
\[
 \lambda_1\leq ad+3
 \quad\Longrightarrow\quad
 m_\lambda(\ell^a f)\leq m_\lambda(\det_n)
\]
for every degree-$m$ form $f$ in the common ambient setting.
Every positive gap in this sector therefore satisfies
$\lambda_1\geq ad+4$. In particular a rectangular gap
with $\lambda=(nd/4)^4$ requires
\[
 d(4m-3n)>12.
\]
Since $4\mid nd$ for this rectangular shape, the last
quantity is divisible by four; equivalently,
$d(4m-3n)\geq16$ is necessary.
\end{proposition}

\paragraph{Background and proof strategy.}
We use the standard $\mathfrak{sl}_2$ decompositions
\[
 \operatorname{End}(\operatorname{Sym}^{m-1}\mathbb C^2)
       =\bigoplus_{l=0}^{m-1}\operatorname{Sym}^{2l}\mathbb C^2,
 \qquad
 \operatorname{Sym}^{r}(\operatorname{Sym}^2\mathbb C^2)
       =\bigoplus_{b=0}^{\lfloor r/2\rfloor}
             \operatorname{Sym}^{2r-4b}\mathbb C^2.
\]
These are the Clebsch--Gordan and harmonic decompositions;
see Fulton and Harris (1991), section~11.2, Exercises~11.11
and 11.14, printed pp.~151 and 153. We specify their
covariants, the finite orthogonal-polynomial calculation,
and the required rank argument below. No generic strong
Lefschetz property or Gulliksen--Neg\aa rd theorem is
assumed in this proof.

\begin{proof}
The case $m=1$ is immediate. Fix $m\geq2$ and let
$\rho$ be the irreducible representation of
$\mathfrak{sl}_2$ on $V=\operatorname{Sym}^{m-1}\mathbb C^2$.
Use standard generators $H,E,F$, with
$[H,E]=2E$, $[H,F]=-2F$, and $[E,F]=H$.
Choose the invariant Hermitian metric for which
$\rho(H)$ is Hermitian and $\rho(E)^*=\rho(F)$.
Its $m$ distinct $H$-eigenvalues are
\[
 \lambda_j=m-1-2j,\qquad 0\leq j<m.
\]
For $z=(z_1,z_2,z_3)$ put
\[
 A(z)=z_1H+z_2(E+F)+\mathrm i z_3(E-F),\qquad
 q=z_1^2+z_2^2+z_3^2,
\]
and consider the explicit linear pencil
\begin{equation}\label{G15:pencil}
 M(t,z)=tI_m-\rho(A(z)).
\end{equation}
The invariant quadratic form is normalized so that
$q(H)=1$. In the defining two-dimensional representation,
$A(z)^2=q(z)I_2$.

Let $\mathcal H_l$ be the space of harmonic homogeneous
polynomials of degree $l$ in $z_1,z_2,z_3$. It is the
irreducible representation of highest weight $2l$ and
dimension $2l+1$. The harmonic decomposition reads
\begin{equation}\label{G15:harmonic}
 \mathbb C[z_1,z_2,z_3]
     =\bigoplus_{l\geq0}\mathcal H_l\otimes\mathbb C[q].
\end{equation}
One can obtain this directly by applying the Laplacian to
$q^b h_l$: for harmonic $h_l$,
\[
 \Delta(q^b h_l)=2b(2l+2b+1)q^{b-1}h_l.
\]
This gives the direct decomposition and its dimensions in
characteristic zero. The indicated irreducibility is the
standard highest-weight calculation for the adjoint
three-dimensional representation.

\paragraph{1. Canonical harmonic matrix covariants.}
Write the multiplicity-free decomposition of
$\operatorname{End}(V)$ as
$\bigoplus_{l=0}^{m-1}\mathcal E_l$, where
$\mathcal E_l\simeq\mathcal H_l$, and let $\pi_l$ be its
equivariant projection. The trace Hermitian form makes
these distinct summands orthogonal. The polynomial map
\[
 A\longmapsto\pi_l\bigl(\rho(A)^l\bigr)
\]
is nonzero and has harmonic coefficient functions of
degree $l$. Indeed the symmetrized multiplication map
from the $l$th symmetric power of the adjoint
representation to $\operatorname{End}(V)$ is equivariant.
Its highest-weight vector $E^l$ maps to
$\rho(E)^l\ne0$ for $l<m$, and this vector has highest
weight $2l$, so lies in $\mathcal E_l$. The harmonic
decomposition says that this is the unique copy of
highest weight $2l$ in symmetric degree $l$.

The value of this covariant at $H$ is nonzero. Otherwise
equivariance and homogeneity would make it vanish on
every scalar multiple of the semisimple orbit of $H$,
which is the dense open set $q\ne0$, contradicting its
nonvanishing. This value is diagonal, since it has
Cartan weight zero.

These diagonal vectors have an elementary description.
Put the uniform positive inner product on functions on
$\{\lambda_0,\ldots,\lambda_{m-1}\}$:
\[
 \langle f,g\rangle=\frac1m\sum_{j=0}^{m-1}
                          f(\lambda_j)\overline{g(\lambda_j)}.
\]
Let $p_0,\ldots,p_{m-1}$ be the real orthonormal
polynomials obtained from $1,x,\ldots,x^{m-1}$, with
positive leading coefficients; thus $p_0=1$.
The zero-weight vector in $\mathcal E_l$ is proportional
to $p_l(\rho(H))$. To check the assertion, note that
$\rho(H)^r$ lies in the sum of the zero-weight spaces of
$\mathcal E_0,\ldots,\mathcal E_r$ by the symmetric-power
decomposition. The top projection is nonzero by the
preceding argument. Induction therefore identifies their
successive spans with those of $1,x,\ldots,x^r$ on the
spectrum. Orthogonality of the endomorphism summands
is precisely orthogonality for the displayed uniform
inner product. This proves the assertion, including
its degree and nonzero leading coefficient.

Normalize the covariant $T_l(A)$ so that
$T_l(H)=p_l(\rho(H))$. Its coefficient functions span
all of $\mathcal H_l$, since it is a nonzero map between
the corresponding irreducible representations.
Equivariance and \eqref{G15:harmonic} now give uniquely
\begin{equation}\label{G15:adjexpansion}
 \pi_l\operatorname{adj} M(t,z)=T_l(A(z))F_l(t,q),
 \qquad 0\leq l<m,
\end{equation}
where $F_l$ is a polynomial homogeneous of weighted
degree $m-1-l$ for $\deg t=1$, $\deg q=2$.

\paragraph{2. The scalar factors are trailing Jacobi determinants.}
Multiplication by $x$ on the functions on the $m$ nodes,
expressed in the orthonormal basis $p_0,\ldots,p_{m-1}$,
is a real symmetric tridiagonal matrix $J$.
Tridiagonality follows because $xp_l$ is orthogonal
to every $p_r$ with $r<l-1$. Its off-diagonal entries
$b_1,\ldots,b_{m-1}$ are strictly positive: they are
the ratios of the positive leading coefficients of
successive orthonormal polynomials. The nodes and
measure are symmetric under $x\mapsto-x$, so $p_l$
has parity $l$. Thus every diagonal entry
$\langle xp_l,p_l\rangle$ is zero.

Let $J^{[l+1]}$ denote the trailing principal submatrix
on indices $l+1,\ldots,m-1$, with the empty determinant
equal to one, and define
\[
 Q_l(t)=\det(tI-J^{[l+1]}),\qquad 0\leq l<m.
\]
The spectrum of $J$ is the node set. If
$D(t)=\prod_j(t-\lambda_j)$, evaluating
\eqref{G15:adjexpansion} at $H$ and taking its
orthonormal trace coefficient gives
\[
 F_l(t,1)=\frac1m\sum_jp_l(\lambda_j)
                         \frac{D(t)}{t-\lambda_j}.
\]
On the other hand the spectral expression for the
$(l,0)$ entry of $(tI-J)^{-1}$ is
\[
 \frac1m\sum_j\frac{p_l(\lambda_j)}{t-\lambda_j}.
\]
Expansion of the corresponding tridiagonal cofactor
shows that its numerator is
$b_1\cdots b_l Q_l(t)$: the unique path from index
zero to index $l$ contributes the indicated product,
and all remaining rows and columns give the trailing
principal determinant. Therefore
\begin{equation}\label{G15:jacobi}
 F_l(t,1)=b_1\cdots b_l Q_l(t).
\end{equation}
In particular every factor is nonzero. Weighted
homogeneity then determines $F_l(t,q)$ from this
univariate identity. All $Q_l$ have the parity of
their degree $m-1-l$.

We record the elementary spectral facts needed below.
Every nonempty trailing block is real symmetric,
tridiagonal, and has positive off-diagonal entries.
An eigenvector is determined by its first coordinate
using the three-term recurrence; that coordinate cannot
be zero. Hence all eigenvalues are simple. Conjugation
by the alternating-sign diagonal matrix negates the
block, so its eigenvalues occur in opposite pairs.
An even-sized block has no zero eigenvalue, as its
determinant at zero is a nonzero signed product of
alternate squared off-diagonal entries. An odd-sized
block has a simple zero eigenvalue. These assertions
also hold for blocks of sizes zero or one with the
obvious interpretations.

\paragraph{3. Cofactor multiples span every third jet.}
Let $S=\mathbb C[t,z_1,z_2,z_3]$, and let $\mathcal C_m$
be the degree-$m$ space spanned by all linear-form
multiples of the entries of $\operatorname{adj} M$.
It is the differential image of the determinant-pencil
map at \eqref{G15:pencil}. Formula
\eqref{G15:adjexpansion} implies that it contains
\[
 tF_l(t,q)\mathcal H_l\qquad(0\leq l<m).
\]
It also contains
\begin{equation}\label{G15:lowerproduct}
 qF_{l+1}(t,q)\mathcal H_l\qquad(0\leq l\leq m-2).
\end{equation}
For clarity, multiplication of a harmonic polynomial of
degree $l+1$ by a linear form has a nonzero projection
onto $q\mathcal H_l$. Indeed
$\Delta(z_i h)=2\partial_i h$; these derivatives span
$\mathcal H_l$ by equivariance and nonvanishing.
The space $\mathcal C_m$ is invariant, so projecting
its products onto this irreducible summand stays in
$\mathcal C_m$. Multiplication by the invariant
$F_{l+1}$ gives \eqref{G15:lowerproduct}.
Similarly the top harmonic component of
$\mathcal H_1\mathcal H_{m-1}$ supplies
$\mathcal H_m\subseteq\mathcal C_m$, since
$F_{m-1}$ is a nonzero constant.

The quotient $S_m/t^4S_{m-4}$ is exactly the third-jet
space, with $S_{m-4}=0$ if $m<4$. Its harmonic
decomposition has, for each $0\leq l\leq m$, the
multiplicity space
\begin{equation}\label{G15:targetspace}
 \operatorname{span}\{t^j q^{(D-j)/2}:
         0\leq j\leq\min\{3,D\},\ j\equiv D\pmod2\},
 \qquad D=m-l,
\end{equation}
tensored with $\mathcal H_l$. It has dimension one
for $D=0,1$, and dimension two for $D\geq2$.
The cases $D=0$ and $D=1$ are already supplied by
$\mathcal H_m$ and $tF_{m-1}\mathcal H_{m-1}$.

For $D\geq2$, specialize $q=1$. On the fixed
multiplicity space \eqref{G15:targetspace}, this is
an isomorphism onto the corresponding powers of $t$
modulo $t^4$. By \eqref{G15:jacobi}, it suffices to
show that $tQ_l(t)$ and $Q_{l+1}(t)$ are independent
modulo $t^4$.

If $D$ is even, $Q_l$ has odd degree $D-1$ and a
simple root at zero, while $Q_{l+1}$ has even degree
$D-2$ and nonzero value at zero. Thus $tQ_l$ starts
with a nonzero $t^2$ term, while $Q_{l+1}$ has a
nonzero constant term. They are independent in the
even two-dimensional space modulo $t^4$.

If $D$ is odd, then $D\geq3$. The block
$K=J^{[l+1]}$ is real symmetric and invertible,
of even size $D-1$. Write $e$ for its first
coordinate vector. The diagonal resolvent identity gives
\[
 \frac{Q_{l+1}(t)}{Q_l(t)}
    =e^{\mathsf T}(tI-K)^{-1}e.
\]
The alternating-sign symmetry fixes $e$ and negates
$K$, so all diagonal entries of odd powers of
$K^{-1}$ vanish. Expanding the inverse at zero gives
\[
 \frac{Q_{l+1}(t)}{tQ_l(t)}
   =-e^{\mathsf T}K^{-2}e
     -t^2 e^{\mathsf T}K^{-4}e
     +O(t^4).
\]
The coefficient of $t^2$ is strictly negative, since
$e^{\mathsf T}K^{-4}e=\|K^{-2}e\|^2>0$. Thus
$Q_{l+1}$ cannot be a constant multiple of $tQ_l$
modulo $t^4$. This proves independence in the odd
two-dimensional space, completing all cases of
\eqref{G15:targetspace}.

Consequently the differential image $\mathcal C_m$
surjects onto $S_m/t^4S_{m-4}$. The determinant-pencil
parameter map followed by third-jet projection has a
surjective differential at the explicit point
\eqref{G15:pencil}. In characteristic zero its image
is dense. This proves third-jet dominance for every
$m\geq2$, and hence for all $m\geq1$.

For $m\geq4$ the third-jet target dimension is
\[
 \sum_{j=0}^3\binom{m-j+2}{2}=2m^2+2.
\]
The fourth-jet target adds
$\binom{m-2}{2}>0$, whereas G-06 bounds
$\dim\mathcal D_{m,4}\leq2m^2+2$. The fourth-jet
map, and therefore every higher full jet map, cannot
be dominant. This establishes the maximal-order
assertion.

Finally apply the coefficient-weight argument of
G-14 with the now proved third-jet dominance. A
degree-$d$ coefficient monomial surviving on
$x_0^a q$ has $x_0$ weight
$ad+j_1+\cdots+j_d$. If $\lambda_1\leq ad+3$,
its evaluation depends only on the first four
coefficient forms $q_0,q_1,q_2,q_3$. Their dense
realization by determinants and block padding force
the same isotypic ideal containment as in G-14.
This proves the multiplicity inequality, including
the elementary cases $m<3$ where all available
jets are included. The rectangular inequality is
the equivalent condition
$nd/4>ad+3$, namely $d(4m-3n)>12$.
\end{proof}

\paragraph{What this resolves.}
The argument proves the needed rank for one explicit
pencil in every size and hence proves dominance. It
does not assume or assert a general strong Lefschetz
theorem for arbitrary Artinian determinantal algebras.
The positivity above is the squared norm for an
ordinary finite-dimensional real symmetric Jacobi
matrix; no arithmetic purity or RH input is used.
No literature-priority claim is made for this deduction.

\begin{proposition}[G-16: Explicit nonrectangular labels throughout the dimension-certified range]\label{G-16}
Let $n>m\geq3$, put $a=n-m$, and suppose
\begin{equation}\label{G16:dimensionhypothesis}
 \dim\mathcal P_{m,4}+3>2n^2+2.
\end{equation}
Define
\[
 b=\binom{n+3}{3}\binom{a+3}{3},\qquad L=n+3a,
 \qquad \beta=6bL^{30},\qquad u=2n^2-10,
\]
and set
\begin{equation}\label{G16:firstdegree}
 C=((n+a)\beta^2)^u,\qquad
 s_*=(u+1)(C-1)+1,\qquad d_* =\beta s_*.
\end{equation}
There is an integer $1\leq d\leq d_*$ with $4\mid md$ such that,
for the explicit shape
\begin{equation}\label{G16:label}
 \lambda_{m,n,d}=
 \left(ad+\frac{md}{4},\frac{md}{4},\frac{md}{4},\frac{md}{4}\right),
\end{equation}
one has
\[
 m_{\lambda_{m,n,d}}(\ell^a\operatorname{per}_m)
       >m_{\lambda_{m,n,d}}(\det_n).
\]
The bound has $\log d_*=O(n^2\log n)$ uniformly for $3\leq m<n$.
There is no additional condition $n<4m/3$ in this proposition.

An explicit progression of successful degrees is also available. Put
\[
 T=\binom{b+\beta}{b}-1,\quad E=4T(\beta!),\quad
 Q=((n+a)E)^{u-1},\quad K=u(Q-1)+1.
\]
For every integer $k\geq K$, the same inequality holds at
$d=Ek$ and the label \eqref{G16:label}.
Because $a>0$, every displayed label is nonrectangular.
\end{proposition}

\paragraph{Scope and sources.}
This applies to every padded pair certified by R-14, R-16, R-18,
R-19, or R-20, including those at or above $n=4m/3$, where G-12
excludes rectangular gaps. The construction uses a Segre--Veronese
coordinate ring and the same Kempf--Ness and reductive-quotient
facts used in G-11. For the quantitative bounds, we use Derksen
(2001), Theorem~1.1 and Proposition~1.2, printed p.~956, and the
image-degree/Chardin argument of G-10 and G-13. The new representation
requires a different coefficient-degree bound, which is checked below.

\begin{proof}
Let $V=\mathbb C^4$, $G=\mathrm{SL}(V)$, and put
\[
 Y_D=\mathcal D_{n,4},\qquad
 Y_P=\{L^a q:L\in V,\ q\in\mathcal P_{m,4}\}^{\rm cl}.
\]
These are the irreducible cones used in R-14, with
\[
 \dim Y_D=2n^2+2,\qquad
 \dim Y_P=\dim\mathcal P_{m,4}+3.
\]
For either cone $Y\subseteq\operatorname{Sym}^n V$, define
\begin{equation}\label{G16:enhancedcone}
 X_Y=\{F\otimes v^a:F\in Y,\ v\in V^*\}^{\rm cl}
 \subseteq W:=\operatorname{Sym}^n V\otimes\operatorname{Sym}^a V^*.
\end{equation}

\paragraph{1. The graded ring and the selected representation.}
The projectivization of \eqref{G16:enhancedcone} is the closed
Segre--Veronese image of $\mathbb P(Y)\times\mathbb P(V^*)$.
In particular $X_Y$ is irreducible and
\begin{equation}\label{G16:xdimension}
 \dim X_Y=\dim Y+3.
\end{equation}
Its coordinate ring has the exact grading
\begin{equation}\label{G16:segrering}
 \mathbb C[X_Y]_d=
       \mathbb C[Y]_d\otimes\operatorname{Sym}^{ad}V.
\end{equation}
Indeed each linear coordinate restricts to a product of a degree-one
coefficient function of $F$ and a degree-$a$ polynomial function of $v$.
Their degree-$d$ products span the right side: the first factor is
generated in degree one, and every degree-$ad$ monomial in $v$ can
be split into $d$ degree-$a$ monomials. The only additional relations
are exactly the restrictions defining this image. This proves
\eqref{G16:segrering} directly, without a projective-normality assumption
on $Y$.

To keep the duals explicit, use the actual coordinate-ring decomposition
\[
 \mathbb C[Y]_d=
       \bigoplus_{\lambda\vdash nd}m_\lambda(Y)S_\lambda V^*.
\]
It is the dual of the positive-weight orbit-span convention used
elsewhere in the paper. Schur's lemma for $G$ and
\eqref{G16:segrering} give
\begin{equation}\label{G16:extraction}
 \dim\mathbb C[X_Y]_d^G=
 \begin{cases}
 m_{\lambda_{m,n,d}}(Y),&4\mid md,\\
 0,&4\nmid md.
 \end{cases}
\end{equation}
For full detail,
\[
 (S_\lambda V^*\otimes\operatorname{Sym}^{ad}V)^G
   =\operatorname{Hom}_G(S_\lambda V,\operatorname{Sym}^{ad}V).
\]
This has dimension one precisely when
$\lambda=(ad+r,r,r,r)$ for an integer $r\geq0$, and otherwise is
zero. The fixed size $|\lambda|=nd$ forces $4r=md$, proving
the formula and ruling out any additional labels in this extraction.

\paragraph{2. Stable points with arbitrary padding.}
Let $x_0,x_1,x_2,x_3$ be an orthonormal basis of $V$ and
$e_0^*$ its first dual vector. Set
\[
 \gamma=\sqrt{\binom{n}{m}},\qquad
 F=x_0^a\bigl(x_0^m+\gamma x_1^m+
                         \gamma x_2^m+\gamma x_3^m\bigr),
 \qquad w=F\otimes(e_0^*)^a.
\]
The degree-$m$ factor of $F$ belongs to both
$\mathcal P_{m,4}$ and $\mathcal D_{m,4}$. In fact each of its two
binary summands factors into $m$ linear factors over $\mathbb C$;
put those factors on the diagonal and the cycle of an $m\times m$
pencil. Its permanent is the sum of the two products, and its
determinant is the same sum after absorbing $(-1)^{m-1}$ into one
cycle factor. Block padding then gives $F\in Y_P\cap Y_D$.
Thus the same tensor $w$ lies in both enhanced cones.

Use the Hermitian structures induced from tensor powers. The four
monomial tensors in $w$ are mutually orthogonal and have equal norm:
$\|x_0^n\|=1$, while
$\|x_0^a x_i^m\|^2=a!m!/n!$ for $i>0$, which is canceled by
$\gamma^2$. Their torus weights are exactly
\[
 ne_0-ae_0=me_0,\qquad
 ae_0+me_i-ae_0=me_i\quad(1\leq i\leq3).
\]
Their mean weight is $(m/4)(1,1,1,1)$, so the traceless diagonal
moment map is zero. The off-diagonal terms vanish as well. On the
first tensor factor, the four monomial weights differ by
$m(e_i-e_j)$, never by one root when $m\geq3$. On the second
factor, an off-diagonal derivative is either zero or orthogonal
to $(e_0^*)^a$. Therefore $w$ has zero $G$ moment map.

Its projective stabilizer has zero Lie algebra. To verify this
without a generic-stabilizer assertion, suppose
$A\in\mathfrak{sl}_4$ and $A\cdot F=\kappa F$.
Write $F=x_0^a q$. Dividing by $x_0^{a-1}$ and reducing modulo
$x_0$ shows that $A\cdot x_0$ is a multiple $\eta x_0$, since
$a>0$ and $q|_{x_0=0}\ne0$. It follows that
$A\cdot q=(\kappa-a\eta)q$.
For the Fermat sum $q$, the monomials $x_j^{m-1}x_i$, $i\ne j$,
are all distinct because $m\geq3$. Their coefficients force every
off-diagonal entry of $A$ to vanish. The four pure-power coefficients
then force every diagonal entry to be the same, and trace zero
forces $A=0$. If $A$ stabilizes the projective tensor line of $w$,
the elementary tangent property of a nonzero decomposable tensor
forces it to stabilize the line of $F$; hence again $A=0$.
Thus the projective stabilizer is finite.

Kempf--Ness now makes $w$ a stable point of each enhanced cone.
Their stable loci are nonempty open sets. Using
\eqref{G16:xdimension}, the two affine invariant quotients therefore
have dimensions
\begin{equation}\label{G16:quotdims}
 \begin{split}
 \dim A_D&=u=2n^2-10,\\
 \dim A_P&=\dim\mathcal P_{m,4}-9>u,
 \end{split}
 \qquad
 A_D=\mathbb C[X_{Y_D}]^G,\quad A_P=\mathbb C[X_{Y_P}]^G.
\end{equation}
The subtraction is $15=\dim G$. The padding weights were canceled
in the tensor construction, so this calculation needs no inequality
between $n$ and $4m/3$.

\paragraph{3. A quantitative invariant bound for this representation.}
The dimension of $W$ is $b$. In the matrix coordinates on
$\mathrm{SL}_4$, the action on $V^*$ is given by the transpose
adjugate, whose entries have degree three. The action on $W$
therefore has matrix coefficients of degree at most
$L=n+3a$. Its kernel is finite: its differential is a nonzero
representation of the simple Lie algebra $\mathfrak{sl}_4$, hence
injective, and a zero-dimensional algebraic kernel is finite.
For example the vector $x_0^n\otimes(e_1^*)^a$ has a nonzero
diagonal Lie-algebra weight, verifying nontriviality.

Derksen's Proposition~1.2 applies to the embedding
$\mathrm{SL}_4\subseteq\mathbb A^{16}$ defined by the degree-four
equation $\det-1=0$. Thus its parameters are $t=16$, group dimension
$15$, equation-degree bound $4$, and action-degree bound $L$, giving
$\sigma(W)\leq4L^{15}$.
Theorem~1.1, with invariant-quotient dimension at most $b$, then gives
generators of $\mathbb C[W]^G$ of degrees at most
\[
 \max\{2,(3/8)b(4L^{15})^2\}\leq6bL^{30}=\beta.
\]
Exactness of invariants for this reductive group makes $A_D$ and
$A_P$ graded quotients of that invariant ring. Choose common
homogeneous generators $g_i$ of degrees $1\leq e_i\leq\beta$.

The determinant enhanced cone has a dense parametrization
\[
 (Z,v)\longmapsto\det(Z_1x_0+Z_2x_1+Z_3x_2+Z_4x_3)\otimes v^a.
\]
Its coordinates are homogeneous of degree $n+a$ in the
$4n^2+4$ parameters. This parameter degree is distinct from the
group-action coefficient degree $n+3a$ above.
The generator coordinates $g_i$ composed with this map have degree
$(n+a)e_i\leq(n+a)\beta$.
Homogenizing them with one additional source coordinate gives a
dominant rational map to the ordinary projective closure of
$\operatorname{Spec}A_D$ in its generator embedding. The image-degree
argument of G-10, including its treatment of a base locus, yields
\[
 \deg\overline{\operatorname{Spec}A_D}
      \leq((n+a)\beta)^u.
\]

Let $H_D^{\leq}(d)=\sum_{j=0}^d\dim(A_D)_j$, and define
$H_P^{\leq}$ similarly. The graded pieces through degree $d$
are contained in the ordinary degree-at-most-$d$ generator filtration,
because all generator weights are at least one. Chardin's bound on
the ordinary projective closure therefore gives
\begin{equation}\label{G16:cumulativeupper}
 H_D^{\leq}(d)\leq((n+a)\beta)^u\binom{d+u}{u}.
\end{equation}
By \eqref{G16:quotdims}, at least $u+1$ images of the generators in
$A_P$ are algebraically independent. Their monomials of ordinary
total degree at most $s$ have graded degree at most $\beta s$;
hence
\begin{equation}\label{G16:cumulativelower}
 H_P^{\leq}(\beta s)\geq\binom{s+u+1}{u+1}.
\end{equation}
As in G-13, the exact binomial comparison is
\[
 \frac{\binom{s+u+1}{u+1}}{\binom{\beta s+u}{u}}
 =\frac{s+u+1}{u+1}\prod_{i=1}^u\frac{s+i}{\beta s+i}
 \geq\frac{s+u+1}{(u+1)\beta^u}.
\]
At $s=s_*$ this is strictly greater than
$C/\beta^u=((n+a)\beta)^u$. Equations
\eqref{G16:cumulativeupper}--\eqref{G16:cumulativelower} give a
strict cumulative gap through $d_*$. Degree zero contributes one
on both sides, so there is a positive degree $d\leq d_*$ with
$\dim(A_P)_d>\dim(A_D)_d$.
Formula \eqref{G16:extraction} identifies this gap with exactly the
label \eqref{G16:label} and forces $4\mid md$.
The four-variable highest-weight restriction of G-05/R-14 transfers
it to the full padded-permanent and determinant orbit closures.
Since $a<n$, $b=O(n^6)$, $L=O(n)$, and $\beta=O(n^{36})$;
\eqref{G16:firstdegree} consequently has
$\log d_*=O(n^2\log n)$.

\paragraph{4. An explicit progression, if individual degrees are desired.}
There are at most $T=\binom{b+\beta}{b}-1$ generators, since the
ambient polynomial space in positive degrees through $\beta$ has
that dimension. The elementary generator-grouping proof in G-11
therefore makes both $E$th Veronese rings standard graded for
$E=4T(\beta!)$. Their dimensions are still $u$ and at least $u+1$.
The degree-one invariant functions of the first Veronese, composed
with the enhanced determinant parametrization, have degree
$(n+a)E$. The projective image has dimension $u-1$ and degree at
most $Q=((n+a)E)^{u-1}$.
Chardin's theorem and linear Noether normalization give respectively
\[
 \dim(A_D)_{Ek}\leq Q\binom{k+u-1}{u-1},\qquad
 \dim(A_P)_{Ek}\geq\binom{k+u}{u}.
\]
Their binomial ratio is $(k+u)/u$, strictly larger than $Q$ for
$k\geq u(Q-1)+1$. This proves the claimed progression; $4\mid E$
ensures all labels in it are integral.
\end{proof}

\paragraph{Examples where ordinary rectangles are excluded.}
The pairs $(m,n)=(15,20)$ and $(16,22)$ from R-17 satisfy
\eqref{G16:dimensionhypothesis}. G-16 produces labels
$\lambda=(35r,15r,15r,15r)$ at degrees $d=4r$ for the first pair,
and labels $\lambda=(10d,4d,4d,4d)$ for the second, at the
successful degrees specified above. G-12 precludes ordinary
rectangular gaps for both pairs. The same explicit label family
applies to the unbounded nonrectangular strip supplied by R-20.


\begin{theorem}[G-17: Exact criterion for four-row stretched directions]\label{G-17}
Assume $n>m\geq3$, put $a=n-m$, and suppose
\[
 \dim\mathcal P_{m,4}+3>2n^2+2.
\]
Fix $s\geq1$ and a partition $\nu\vdash ns$ of length at most four,
appending zero parts when necessary. There is a positive integer $j$
such that
\begin{equation}\label{G17:gap}
 m_{j\nu}(\ell^a\operatorname{per}_m)>m_{j\nu}(\det_n),
 \qquad\text{in coordinate degree }sj,
\end{equation}
if and only if
\begin{equation}\label{G17:criterion}
 \nu_4>0\qquad\hbox{and}\qquad\nu_1>as.
\end{equation}
If either condition fails, no positive stretch $j$ gives a gap.
If both hold, an explicit progression of stretches gives gaps at
every sufficiently large term, as specified below. This does not
assert a gap at every sufficiently large integer stretch.

Here are effective bounds in the affirmative case. Put
\[
 \eta=(\nu_1-\nu_4,\nu_2-\nu_4,\nu_3-\nu_4,0),\qquad e=|\eta|,
 \qquad r_\eta=\#\{(i,j):1\leq i<j\leq4,\ \eta_i>\eta_j\},
\]
and define
\[
 \begin{aligned}
 b&=\binom{\binom{n+3}{3}+s-1}{s}\dim S_\eta\mathbb C^4,
 &u&=2n^2+r_\eta-13,\\
 L&=ns+3e,&p&=ns+e,\qquad \beta=6bL^{30}.
 \end{aligned}
\]
The Schur dimension is explicitly
$\prod_{i<j}(\eta_i-\eta_j+j-i)/(j-i)$.
One successful stretch satisfies
\begin{equation}\label{G17:firststretch}
 1\leq j\leq
 j_*:=\beta\bigl[(u+1)((p\beta^2)^u-1)+1\bigr].
\end{equation}
For a progression, let
\[
 T=\binom{b+\beta}{b}-1,\quad E=4T(\beta!),\quad
 Q=(pE)^{u-1},\quad K=u(Q-1)+1.
\]
Then \eqref{G17:gap} holds for every stretch $j=Ek$ with $k\geq K$.
\end{theorem}

\paragraph{Meaning of the criterion.}
For any one of the dimension-certified padded pairs, this determines
exactly which rational directions in the four-row partition simplex
contain a stretched multiplicity obstruction. Its two excluded faces
come from low-row determinant survival and the padding weight bound.
In particular this is a classification of all such directions,
not only a further special choice of the one-longer-row labels in G-16.
The coordinate degrees and stretch factors remain large; no efficient
multiplicity-positivity algorithm is asserted.

\begin{proof}
Necessity follows from earlier ideal containments and includes equality.
If $\nu_4=0$, every $j\nu$ has at most three rows, and G-05 gives
the full ambient multiplicity on the determinant side. If
$\nu_1\leq as$, then in coordinate degree $d=sj$ one has
$(j\nu)_1\leq ad$. The coefficient-weight argument of G-14 applies
and gives the opposite weak multiplicity inequality for every $j$.
Thus it remains to prove sufficiency under \eqref{G17:criterion}.

Use $V=\mathbb C^4$ and the two degree-$n$ cones $Y_D,Y_P$ of G-16.
We first give the auxiliary flag construction, including its coordinate
ring, so that no projective-normality or multiplicity-free assertion
is left implicit.

\paragraph{1. The auxiliary cone and its Cartan coordinate ring.}
Let $v_\eta$ be a nonzero extremal vector of weight $-\eta$ in
$S_\eta V^*$. Concretely, put
$c_i=\eta_i-\eta_{i+1}$ for $i=1,2,3$ and take the product of
the vectors
\[
 (e_0^*\wedge\cdots\wedge e_{i-1}^*)^{\otimes c_i}.
\]
It lies in the single Cartan component $S_\eta V^*$ of the indicated
tensor product of exterior powers. Its projective orbit records the
nested subspaces of dimensions $i$ for which $c_i>0$.
This is the corresponding partial flag variety, a closed incidence
subvariety of a product of Grassmannians. Its dimension is
$r_\eta$: its stabilizer is the block parabolic for the equal-part
blocks of $\eta$, and the missing off-diagonal blocks contribute
one parameter for each pair counted in $r_\eta$.
If $\eta=0$, the projective orbit is a point and $r_\eta=0$.

Let $\mathcal C_\eta$ be the affine cone over this projective orbit.
Its degree-$d$ coordinate ring is
\begin{equation}\label{G17:cartanring}
 \mathbb C[\mathcal C_\eta]_d=S_{d\eta}V.
\end{equation}
Here is a direct representation-theoretic proof. The dual of its
degree-$d$ coordinate space is the span of the group translates of
$v_\eta^{\otimes d}$. This is an extremal vector of weight $-d\eta$.
Its weight space at the extreme of the tensor-product weight order
is one dimensional, and it generates exactly the irreducible Cartan
component $S_{d\eta}V^*$. Taking the dual proves
\eqref{G17:cartanring}. Products give its usual Cartan multiplication.
When $\eta=0$, the cone is a line and each graded piece is
one dimensional, so the same formula holds.

For either cone $Y$, define the enhanced Veronese cone
\begin{equation}\label{G17:enhancement}
 X_Y=\{F^{\otimes s}\otimes v:F\in Y,\ v\in\mathcal C_\eta\}^{\rm cl}
 \subseteq W:=\operatorname{Sym}^s(\operatorname{Sym}^n V)
                      \otimes S_\eta V^*.
\end{equation}
Its projectivization is the Segre product of the $s$th Veronese
embedding of $\mathbb P(Y)$ with the partial flag variety. Hence
\begin{equation}\label{G17:dimension}
 \dim X_Y=\dim Y+r_\eta.
\end{equation}
Pullback of degree-$d$ coordinate functions gives exactly
\begin{equation}\label{G17:ring}
 \mathbb C[X_Y]_d=\mathbb C[Y]_{sd}\otimes S_{d\eta}V.
\end{equation}
The first factor is the Veronese subring of the standard graded
coordinate ring of $Y$, and the second is \eqref{G17:cartanring}.
Both are generated by their degree-one pieces in these gradings;
products of paired coordinates span their diagonal tensor product.
The pullback into the product coordinate ring is injective, proving
\eqref{G17:ring} without assuming that $Y$ is normal.

Write the actual degree-$sd$ coordinate representations as
$\bigoplus_{\lambda\vdash nsd}m_\lambda(Y)S_\lambda V^*$.
The invariant multiplicity in \eqref{G17:ring} is
\[
 \dim\operatorname{Hom}_{\mathrm{SL}_4}
                   (S_\lambda V,S_{d\eta}V).
\]
It is one if $\lambda=d\eta+t(1,1,1,1)$, and zero otherwise.
The size $nsd$ forces
$4t=(ns-e)d=4\nu_4d$, so $\lambda=d\nu$.
Consequently
\begin{equation}\label{G17:extraction}
 \dim\mathbb C[X_Y]_d^{\mathrm{SL}_4}=m_{d\nu}(Y).
\end{equation}
This is a single-label extraction, with original coordinate degree
$sd$, rather than a sum over different labels or degrees.

\paragraph{2. Stability for every strict direction in the criterion.}
Define four positive rational numbers
\[
 \theta_0=\frac{\nu_1/s-a}{m},\qquad
 \theta_i=\frac{\nu_{i+1}}{sm}\quad(1\leq i\leq3).
\]
Positivity is exactly \eqref{G17:criterion}, and their sum is one.
Let $u_0=x_0^n$ and
$u_i=\sqrt{\binom nm}\,x_0^a x_i^m$ for $i>0$.
These are orthonormal monomial vectors for the tensor-induced metric.
Put
\[
 F=\sum_{i=0}^3\sqrt{\theta_i}\,u_i,
 \qquad w=F^{\otimes s}\otimes v_\eta,
\]
normalizing $v_\eta$ to have norm one. The degree-$m$ factor of $F$
is a Fermat sum with four nonzero coefficients. The binary factorization
and cyclic-pencil construction in G-16 therefore put $F$ in both
$Y_D$ and $Y_P$, so $w$ belongs to both enhanced cones.

The diagonal moment of $F$, before removing the scalar part, is
\[
 a e_0+m\sum_i\theta_i e_i=\nu/s.
\]
The moment of a unit pure tensor power is $s$ times the moment of
the original vector; this follows by applying the Lie operator to
each factor and taking the tensor inner product. Thus the first
factor of $w$ contributes $\nu$. The auxiliary extremal vector
contributes its weight $-\eta$. Their sum is
$\nu-\eta=\nu_4(1,1,1,1)$, whose traceless part is zero.
Off-diagonal moments vanish on $F$ because its monomial weights
differ by $m$ times a root, with $m\geq3$, and on $v_\eta$ because
it is a weight vector. They therefore vanish on the tensor product.

The projective stabilizer of this $F$ has zero Lie algebra by the
mixed-monomial calculation in G-16; all four coefficients remain
nonzero, as required there. A Lie operator stabilizing the line of
$F^{\otimes s}$ must stabilize the line of $F$: in the symmetric
algebra on $\operatorname{Sym}^n V$ the equation
$sF^{s-1}(A\cdot F)=cF^s$ implies $A\cdot F=(c/s)F$.
A projective stabilizer of the decomposable tensor $w$ consequently
stabilizes that same line. Thus its projective stabilizer is finite.
Kempf--Ness gives a stable point of each enhanced cone.

The invariant quotients have dimensions
\begin{equation}\label{G17:quotdims}
 \begin{split}
 \dim A_D&=2n^2+r_\eta-13=u,\\
 \dim A_P&=\dim\mathcal P_{m,4}+r_\eta-12>u,
 \end{split}
 \qquad A_D=\mathbb C[X_{Y_D}]^{\mathrm{SL}_4},\quad
 A_P=\mathbb C[X_{Y_P}]^{\mathrm{SL}_4}.
\end{equation}
Indeed stable loci are nonempty and open, so one subtracts the
common orbit dimension $15$ from \eqref{G17:dimension}.
The original dimension gap is preserved for every flag type.

\paragraph{3. The two polynomial degrees and the effective bounds.}
The dimension of $W$ is $b$. Its group-action matrix coefficients
have degree at most $L=ns+3e$ in the entries of $g\in\mathrm{SL}_4$:
the first factor has degree $ns$, while $S_\eta(g^{-\mathsf T})$
has degree at most $3e$ because each entry of $g^{-\mathsf T}$ is
a cubic cofactor. The representation has finite kernel, since its
kernel is contained in the finite projective stabilizer of $w$.
Derksen's Theorem~1.1 and Proposition~1.2 therefore give generators
of $\mathbb C[W]^{\mathrm{SL}_4}$ of degrees at most
$\beta=6bL^{30}$, exactly as in G-16. Exactness of invariants makes
$A_D,A_P$ graded quotients with the same bound on their generators.

For the image-degree bound a smaller, different parameter degree
is valid, and the reason is important. Introduce an independent
matrix $B\in\operatorname{End}(V^*)$, without imposing any relation
between $B$ and the determinant-pencil parameters. The map
\[
 B\longmapsto S_\eta(B)v_\eta
\]
is a homogeneous polynomial map of degree $e$. One can see this
directly in the tensor construction of the Schur functor: it is
the restriction of $B^{\otimes e}$ to the relevant Schur component.
The invertible $B$ form a dense open set and give the full
$\mathrm{GL}(V^*)$ orbit, which is dense in $\mathcal C_\eta$
when $e>0$. Its scalar matrices supply every nonzero scalar
multiple, because the base field is $\mathbb C$.
Thus
\begin{equation}\label{G17:parameterization}
 (Z,B)\longmapsto
  \det(Z_0x_0+\cdots+Z_3x_3)^{\otimes s}
       \otimes S_\eta(B)v_\eta
\end{equation}
has dense image in $X_{Y_D}$ and is homogeneous of degree
$p=ns+e$ in its $4n^2+16$ independent parameters.
No inverse-matrix constraint or adjugate is needed for this
parametrization. If $e=0$, the second factor is constant and can
be omitted; all scalar multiples are already supplied by scaling
the first cone. Formula \eqref{G17:parameterization} then remains
valid with degree $p=ns$.

Apply the cumulative generator-filtration comparison from G-16,
now with quotient dimension $u$, generator bound $\beta$, and
parameter degree $p$. For clarity, its two inequalities become
\[
 \sum_{d=0}^{D}\dim(A_D)_d
       \leq(p\beta)^u\binom{D+u}{u},\qquad
 \sum_{d=0}^{\beta h}\dim(A_P)_d
       \geq\binom{h+u+1}{u+1}.
\]
The first uses the ordinary projective closure in the generator
embedding, of degree at most $(p\beta)^u$, and Chardin's theorem.
The second uses $u+1$ algebraically independent generator images.
The same exact binomial ratio is at least
$(h+u+1)/((u+1)\beta^u)$. Taking
$h=(u+1)((p\beta^2)^u-1)+1$ yields a strict cumulative gap through
$j_*=\beta h$. Since degree zero agrees, some positive stretch
$j\leq j_*$ has a strict invariant-dimension gap. Equation
\eqref{G17:extraction} proves \eqref{G17:gap} at that stretch.

For the individual progression, there are at most
$T=\binom{b+\beta}{b}-1$ homogeneous invariant generators.
The explicit grouping argument of G-11 makes the $E$th Veronese
rings standard graded for $E=4T(\beta!)$.
Their dimensions are $u$ and at least $u+1$, and composing their
degree-one invariant coordinates with \eqref{G17:parameterization}
gives parameter degree $pE$. Consequently their degree-$k$ Hilbert
functions satisfy
\[
 \dim(A_D)_{Ek}\leq Q\binom{k+u-1}{u-1},\qquad
 \dim(A_P)_{Ek}\geq\binom{k+u}{u},\qquad Q=(pE)^{u-1}.
\]
The ratio is strictly greater than one once
$k\geq u(Q-1)+1$. Applying \eqref{G17:extraction} gives the
specified successful multiples $j=Ek$.
On this progression the determinant-side multiplicity is positive as
well. The positive-dimensional graded domain $A_D$ has a nonzero
positive-degree generator $g$ of some degree $e_0\leq\beta$.
Since $e_0\mid E$, the elements $g^{Ek/e_0}$ are nonzero for every
$k\geq1$. Equation \eqref{G17:extraction} therefore gives
$m_{Ek\nu}(\det_n)\geq1$. Thus the displayed strict gaps compare
two positive multiplicities and are not occurrence gaps.
Finally the four-variable highest-weight restriction transfers
all the multiplicities back to the full orbit closures.
\end{proof}

\paragraph{Consistency with the earlier boundaries.}
If $\nu$ is rectangular, the strict inequality in
\eqref{G17:criterion} is $n/4>a$, equivalently $n<4m/3$,
recovering G-11/G-12 whenever the dimension hypothesis holds.
For the one-longer-row direction in G-16, it holds for all padding
because the extra first-row component compensates for $a$.
At equality $\nu_1=as$, every stretch remains excluded, rather
than merely failing this stability construction.
G-15 further excludes a finite initial segment when
$j(\nu_1-as)\leq3$; this is consistent with the eventual
progression in the strict interior.


\paragraph{Quantitative corollary: the first successful stretch has polynomial binary length.}
In the hypotheses and notation of G-17, put
$b_n=\binom{n+3}{3}$. The effective first-stretch bound satisfies
\begin{equation}\label{G17:bitlength}
 \log_2j_*
 \leq\log_2(u+1)+u\log_2p+(2u+1)\log_2\beta
 =O\!\left(n^5\log(s+1)+n^2\log n\right).
\end{equation}
For fixed $s$, the sharper estimate is $O_s(n^2\log n)$.
Thus a successful stretch, its original coordinate degree $sj$, and
its highest-weight entries admit polynomial binary length in $n$ and
$\log(s+1)$. This is a size bound for an existing successful stretch;
it does not provide an efficient search for that stretch or an efficient
certificate of the required multiplicity inequality.

\begin{proof}
The weak compositions of $s$ into $b_n$ parts are determined by their
first $b_n-1$ entries, each between zero and $s$. The Weyl dimension
formula for $\operatorname{GL}_4$ has six factors, each at most $ns+3$,
because $\eta_1\leq ns$. Consequently the ambient dimension in G-17 obeys
\[
 b=\binom{b_n+s-1}{s}\dim S_\eta\mathbb C^4
   \leq(s+1)^{b_n-1}(ns+3)^6.
\]
Moreover $|\eta|\leq ns$, so
\[
 L\leq4ns,\qquad p\leq2ns,\qquad u\leq2n^2,
 \qquad
 \beta\leq6(s+1)^{b_n-1}(ns+3)^6(4ns)^{30}.
\]
The exact bound
$j_*=\beta[(u+1)((p\beta^2)^u-1)+1]$
is at most $\beta(u+1)(p\beta^2)^u$. Taking logarithms gives
the first inequality in \eqref{G17:bitlength}, and $b_n=O(n^3)$
gives its stated order. If $s$ is fixed, instead use
$\binom{b_n+s-1}{s}=O_s(n^{3s})$; then $\log\beta=O_s(\log n)$.
Finally multiplication by $s$ or by a highest-weight entry bounded
by $ns$ adds only $O(\log(ns))$ bits.
\end{proof}


\begin{theorem}[G-18: Every large stretch in every strict four-row direction]\label{G-18}
Let $n>m\geq3$, $a=n-m$, and let $\nu\vdash ns$, $s\geq1$,
have at most four rows. Assume
\[
 \nu_4>0,\qquad \nu_1>as.
\]
Use $\eta,r_\eta,b,L,p,\beta,T,E$ from G-17, defined by the same
formulas even when its dimension inequality is not assumed. Put
\[
 c_m=\dim\mathcal P_{m,4},\qquad
 u=2n^2+r_\eta-13,\qquad v=c_m+r_\eta-12.
\]
The following conclusions hold in original coordinate degree $sj$.

First, both representations occur in every sufficiently large
integer stretch, with the uniform displayed bound
\begin{equation}\label{G18:occurrence}
 m_{j\nu}(\ell^a\operatorname{per}_m)\geq1,
 \qquad m_{j\nu}(\det_n)\geq1
 \qquad(j\geq\beta^2).
\end{equation}
This conclusion needs no dimension comparison between the two cones.

If $v>u$, equivalently $c_m+3>2n^2+2$, set
$Q_D=(pE)^{u-1}$. Then
\begin{equation}\label{G18:everylargepositive}
 m_{j\nu}(\ell^a\operatorname{per}_m)>m_{j\nu}(\det_n)
 \qquad\text{for every integer }j\geq E(4uQ_D+1).
\end{equation}
Thus every affirmative direction of G-17, including the rectangular
ones, gives a gap at every sufficiently large integer stretch.

There is also an effective converse range. Set
\[
 \bar v=4m^2-2m+r_\eta-10,\qquad Q_P=(pE)^{\bar v-1}.
\]
If
\begin{equation}\label{G18:reversecondition}
 2n^2>4m^2-2m+3,
\end{equation}
then
\begin{equation}\label{G18:everylargereverse}
 m_{j\nu}(\det_n)>m_{j\nu}(\ell^a\operatorname{per}_m)
 \qquad\text{for every integer }j\geq E(4\bar vQ_P+1).
\end{equation}
Condition \eqref{G18:reversecondition} holds, in particular,
whenever $n\geq\sqrt2\,m$.

More precisely, for every such fixed direction,
there is a constant $\kappa_\nu>0$ such that
\begin{equation}\label{G18:ratio}
 \frac{m_{j\nu}(\ell^a\operatorname{per}_m)}
      {m_{j\nu}(\det_n)}
  =\kappa_\nu j^{\,c_m-2n^2+1}(1+O(j^{-1}))
  \qquad(j\longrightarrow\infty)
\end{equation}
through all positive integers. When the exponent is zero, the
dimension comparison alone does not determine the sign of the gap.

For every fixed partition direction $\nu\vdash ns$ with at most
four rows, including those outside the two strict inequalities,
there is no occurrence obstruction at any sufficiently large
stretch. Here an occurrence obstruction means positive
padded-permanent multiplicity and zero determinant multiplicity.
\end{theorem}

\paragraph{Sources and the extra step.}
The enhanced coordinate rings and the initial stable tensor are
proved in G-17. We additionally use separation of distinct closed
orbits by invariants: Mukai (2003), Theorem~5.3, printed pp.~161--162,
and Corollary~5.4, p.~162. Openness of the stable locus is
Proposition~5.15, p.~166. The new step removes the finite projective
stabilizer by moving the padding line, then makes the resulting
degree-semigroup argument effective. The original coordinate degree
is $sj$; in a rectangular direction $ns=4\nu_4$ already incorporates
the necessary central congruence.

\begin{proof}
We retain the notation $X_{Y_D},X_{Y_P},W,A_D,A_P$ of G-17.
Their construction, common stable tensor, generator bound, and
dimensions $u,v$ do not require the inequality $v>u$.

\paragraph{1. A common stable tensor with no nontrivial scalar orbit symmetry.}
Write the balanced padded Fermat form of G-17 as $x_0^a f$, where
$f=\sum_{i=0}^3\gamma_i x_i^m$ and every $\gamma_i$ is nonzero.
The degree-$m$ hypersurface $f=0$ is smooth: its partial derivatives
have no common projective zero. It is irreducible, since two
positive-degree projective hypersurface components in $\mathbb P^3$
intersect, and their intersection would be singular in the product.
Its projective stabilizer $H\subseteq\mathrm{PGL}_4$ is finite.
Indeed, the equation $A\cdot f=cf$ forces all off-diagonal entries
of $A$ to vanish by the distinct monomials $x_i^{m-1}x_j$,
$i\ne j$, because $m\geq3$; its diagonal entries are then equal.
Thus the projective stabilizer has zero Lie algebra, so in
characteristic zero it is finite.

Keep $f$ and the auxiliary vector $v_\eta$ fixed, and vary a nonzero
padding linear form $\ell$. For every such $\ell$, the tensor
\[
 w_\ell=(\ell^a f)^{\otimes s}\otimes v_\eta
\]
belongs to both enhanced cones. This follows from the same binary
factorization and cyclic permanent/determinant representations of
$f$ used in G-17, followed by the independent padding block.
The inverse image of the stable locus of either cone under this
regular map is a nonempty open set, because it contains $\ell=x_0$.
Stability is unchanged by a nonzero scalar, so these sets descend
to nonempty open subsets of the projective space of linear forms.

For every nonidentity $h\in H$, its fixed locus in this projective
space is proper and closed: a linear transformation fixing every
line of the dual representation is scalar. The intersection of
the two stable opens with the complements of these finitely many
fixed loci is nonempty, by irreducibility of projective space.
Choose $[\ell]$ in this intersection. Because $a>0$ and $m\geq3$,
unique factorization implies that every projective automorphism of
$F=\ell^a f$ preserves the unique linear factor $\ell$ and the
unique irreducible factor $f$ of degree greater than one separately.
Its image therefore lies in $H$ and fixes $[\ell]$, so it is the
identity in $\mathrm{PGL}_4$.

Set $w=w_\ell$. A projective stabilizer of this decomposable tensor
preserves the line of $F^{\otimes s}$ and hence the line of $F$:
the Veronese map is injective on projective points. Thus its
projective stabilizer in $\mathrm{SL}_4$ is exactly the center
$\mu_4$. This center acts
on $W$ by the exponent
\[
 ns-|\eta|=4\nu_4,
\]
and hence acts trivially. It follows that
\begin{equation}\label{G18:noscaling}
 gw=cw,\quad g\in\mathrm{SL}_4,\ c\in\mathbb C^*
 \quad\Longrightarrow\quad c=1.
\end{equation}
The orbits of $w$ and of every $cw$ are closed: stability holds,
and scalar multiplication commutes with the group action.

\paragraph{2. Common nonzero invariant degrees have greatest common divisor one.}
Choose homogeneous generators $g_1,\ldots,g_t$ of
$\mathbb C[W]^{\mathrm{SL}_4}$, with positive degrees
$e_i\leq\beta$, as in G-17. Let
\[
 J=\{i:g_i(w)\ne0\}.
\]
This set is nonempty. Otherwise all positive-degree invariants
would vanish at $w$, so the distinct closed orbits of $w$ and zero
would fail to be separated.

Suppose the integer $q=\gcd\{e_i:i\in J\}$ exceeded one, and
choose a primitive $q$th root $c\ne1$. Homogeneity gives
$g_i(cw)=g_i(w)$ for every generator: for $i\in J$ use
$c^{e_i}=1$, and for $i\notin J$ both values are zero. The same
equality therefore holds for every invariant polynomial. Closed
orbit separation makes the two closed orbits equal, contrary to
\eqref{G18:noscaling}. Thus
\begin{equation}\label{G18:gcd}
 \gcd\{e_i:i\in J\}=1.
\end{equation}
Every product of these generators is nonzero at $w$, and consequently
has nonzero image in both $A_D$ and $A_P$.

Here is the required elementary numerical-semigroup bound. If
positive integers $e_i\leq\beta$ have greatest common divisor one,
their nonnegative semigroup contains every integer $N\geq\beta^2$.
Let $a_0$ be the least generator. In the directed graph on
$\mathbb Z/a_0\mathbb Z$ with edges adding the $e_i$, every residue
is reachable: a subsemigroup of a finite group is a subgroup, and
the greatest-common-divisor condition makes this subgroup the whole
group. A shortest path to each residue has no repeated vertex and
uses at most $a_0-1$ edges. Its integer weight $c_r$ therefore
satisfies $c_r\leq(a_0-1)\beta<\beta^2$. For $N\geq\beta^2$,
choose the representative of its residue; $N-c_r$ is a nonnegative
multiple of $a_0$. This proves the assertion, including $a_0=1$.

Applying this to \eqref{G18:gcd}, for every $N\geq\beta^2$ there
is a homogeneous ambient invariant $h_N$ of degree $N$ nonzero
at $w$. Its restrictions to both cones are nonzero. The exact
single-label identity of G-17 now proves \eqref{G18:occurrence}.

\paragraph{3. Transfer from the Veronese progression to every residue.}
We have $T\geq\beta$, because the positive-degree polynomial
space through degree $\beta$ contains the powers of one variable.
Since $\beta\geq2$ and $\beta!\geq\beta$, it follows that
$E=4T\beta!\geq\beta^2$.

For $j=Ek+r$, $0\leq r<E$, multiplication by the common nonzero
invariants $h_{E+r}$ and $h_{2E-r}$ is injective in each of the
domains $A_D,A_P$. In particular, for either $A$ and every $k\geq1$,
\begin{equation}\label{G18:sandwich}
 \dim A_{E(k-1)}\leq\dim A_{Ek+r}\leq\dim A_{E(k+2)}.
\end{equation}
Only existence of these invariants is used; no algorithm for
computing a minimal generating set is being asserted.

Suppose first $v>u$. G-17's standard-Veronese dimension and
image-degree bounds give
\[
 \dim(A_P)_{E(k-1)}\geq\binom{k+u-1}{u},\qquad
 \dim(A_D)_{E(k+2)}\leq Q_D\binom{k+u+1}{u-1}.
\]
The exact ratio of the binomials is
\[
 \frac{\binom{k+u-1}{u}}{\binom{k+u+1}{u-1}}
 =\frac{k(k+1)(k+2)}{u(k+u)(k+u+1)}
 \geq\frac{k}{4u}\qquad(k\geq u+1).
\]
For $k\geq4uQ_D+1$ this is strictly greater than $Q_D$.
Combining with \eqref{G18:sandwich} proves
\eqref{G18:everylargepositive} in every residue class.

For the converse bound, R-13's row--column torus gives
$c_m\leq4m^2-2m+2$, and hence $v\leq\bar v$.
Condition \eqref{G18:reversecondition} says exactly $u>\bar v$.
The enhanced permanent cone has a dense homogeneous
parametrization of degree $p=ns+|\eta|$:
\[
 (Z,L,B)\longmapsto
 \bigl(L^a\operatorname{per}_m(Z_0x_0+\cdots+Z_3x_3)\bigr)^{\otimes s}
       \otimes S_\eta(B)v_\eta.
\]
Here the pencil coefficients, coefficients of the linear form $L$,
and entries of $B\in\operatorname{End}(V^*)$ are independent.
The first factor has parameter degree $ns$, exactly as for the
determinant parametrization. Therefore the projective image of
the $E$th invariant Veronese has degree at most
$(pE)^{v-1}\leq Q_P$. The same Chardin upper bound, enlarged to
dimension $\bar v$, and Noether-normalization lower bound give
\[
 \dim(A_P)_{E(k+2)}\leq Q_P\binom{k+\bar v+1}{\bar v-1},
 \qquad
 \dim(A_D)_{E(k-1)}\geq\binom{k+\bar v-1}{\bar v}.
\]
Use \eqref{G18:sandwich} with the roles reversed and the preceding
binomial calculation with $\bar v$ in place of $u$.
This proves \eqref{G18:everylargereverse}. If
$n\geq\sqrt2\,m$, then
$2n^2-(4m^2-2m+3)\geq2m-3>0$, as required.

\paragraph{4. The all-integer asymptotic ratio.}
For a standard graded domain $B$ of dimension $q\geq2$, write
$e(B)=\deg\operatorname{Proj}B>0$. Its Hilbert function satisfies
\[
 \dim B_k=\frac{e(B)}{(q-1)!}k^{q-1}+O(k^{q-2}).
\]
This is the usual Hilbert-polynomial normalization used in G-10.
One can check the leading term directly after linear Noether
normalization: $B$ is a finite torsion-free module over a polynomial
ring in $q$ degree-one variables; a homogeneous basis over its
fraction field spans a free submodule of rank $e(B)$, and the
torsion quotient has growth $O(k^{q-2})$. An annihilating homogeneous
polynomial for that quotient and its finite generating set give
this error bound from the polynomial-ring Hilbert function.

Apply this formula to $B_D=A_D^{(E)}$ and $B_P=A_P^{(E)}$.
Their dimensions are $u,v$, both at least two here. Indeed
$c_m\geq20$: R-13 gives $c_3=20$, and a block diagonal pencil
with a $3\times3$ block and $m-3$ entries $x_0$ embeds
$x_0^{m-3}\mathcal P_{3,4}$ in $\mathcal P_{m,4}$.
Also $r_\eta\geq0$, including the rectangular case $\eta=0$.
Sandwich
\eqref{G18:sandwich}, uniformly over its finitely many residues,
gives for $A$ of dimension $q$
\[
 \dim A_j=\frac{e(A^{(E)})}{E^{q-1}(q-1)!}
                 j^{q-1}+O(j^{q-2}).
\]
Taking the quotient of these two positive leading terms proves
\eqref{G18:ratio}, with
\[
 \kappa_\nu=
 \frac{e(B_P)}{e(B_D)}\,
 \frac{(u-1)!}{(v-1)!}\,E^{u-v}>0,
 \qquad v-u=c_m-2n^2+1.
\]
All multiplicities are identified with the full orbit-closure
ones by the same four-variable highest-weight restriction as in
G-17. The final occurrence assertion follows from
\eqref{G18:occurrence} in the strict region. Outside that region,
G-05 and G-14 give the opposite weak multiplicity inequality at
every stretch, so an occurrence obstruction is already impossible.
\end{proof}

\paragraph{The precise remaining range.}
The reverse inequality concerns sufficiently large stretches of
each fixed direction. It does not exclude an individual
finite-degree gap, give a threshold uniform over changing
directions, or prove containment of orbit closures. Rectangular
directions obey the padding boundary of G-12 as well as the
all-stretch conclusion in the strict range. When
$c_m+3=2n^2+2$, the leading quotient degrees, and
possibly lower Hilbert coefficients, are the remaining comparison.


\subsection{A specified rectangular discriminant copy}

\begin{theorem}[G-19: an explicit rectangular discriminant module and its positive energy]
Let $m\geq3$, put $N=m^2$, and work over $\mathbb C$.
Here $V=\mathbb C^N$ is the argument space, forms lie in
$\operatorname{Sym}^m V^*$, and coefficient functions of degree
$D$ lie in $\operatorname{Sym}^D(\operatorname{Sym}^m V)$.
We use $g\cdot F=F\circ g^{-1}$, hence
$(g\cdot h)(F)=h(F\circ g)$.
Let $k\geq5$ satisfy
\[
 k\leq \operatorname{codim}_{\mathbb A^N}
                 \operatorname{Sing}(\operatorname{per}_m).
\]
In particular $k=5$ is allowed for every $m\geq3$, and $k=6$
is allowed for every $m\geq6$, by the established bounds in R-B17.
Define
\[
 D=k(m-1)^{k-1},\qquad \delta=m(m-1)^{k-1},\qquad
 \lambda=(\delta^k).
\]
There is an explicit highest-weight polynomial $h$ of coefficient
degree $D$ and weight $\lambda$, namely the discriminant of the
restriction to the first $k$ variables, with these properties.

Its generated irreducible module $W$ is contained in the ideal of
the determinant orbit closure, while its restriction to the
permanent orbit closure is nonzero and therefore injective on $W$.
With this explicit dual convention, $W$ has positive polynomial
type $S_\lambda V$. For normalized Haar measure, and any
$U(N)$-invariant Hermitian inner product on $W$, one has
\begin{equation}\label{G19:Haar}
 \int_{U(N)}|h(g\cdot F)|^2\,dg
   =\frac{\|h\|^2}{\dim W}\sum_{j=1}^{\dim W}|h_j(F)|^2,
 \qquad
 \dim W=\prod_{i=1}^k\prod_{j=k+1}^{N}
                     \frac{\delta+j-i}{j-i},
\end{equation}
where $(h_j)$ is an orthonormal basis of $W$.
This integral is zero at $F=\det_m$ and strictly positive at
$F=\operatorname{per}_m$.

There is also a positive integer Gaussian certificate. Normalize
the discriminant to have primitive integer coefficients, let
$A:\mathbb C^k\longrightarrow\mathbb C^N$ have independent
matrix entries, and write
\[
 H_m(A)=\operatorname{Disc}_{k,m}(\operatorname{per}_m\circ A)
        =\sum_{|\alpha|=mD}c_\alpha A^\alpha.
\]
Then $H_m$ is a nonzero integer polynomial, and for independent
standard complex Gaussian entries $Z$ normalized by
$\mathbb E|Z_{ij}|^2=1$,
\begin{equation}\label{G19:Gaussian}
 \mathbb E|H_m(Z)|^2=\sum_\alpha c_\alpha^2\alpha!\ \geq1.
\end{equation}
For the determinant in place of the permanent, the polynomial
and expectation are identically zero. There exists an integer
matrix $A$ with all entries in $\{0,1,\ldots,mD\}$ and
$H_m(A)\ne0$.

This separates a specified irreducible copy, not the total
multiplicities of its isomorphism type. It proves neither
$m_\lambda(\operatorname{per}_m)>m_\lambda(\det_m)$ nor
$m_\lambda(\det_m)=0$. Moreover the analogous degree-$n$
discriminant construction vanishes on every positively padded
form $\ell^{n-m}f_m$, so it supplies no such separator when
$n>m$.
\end{theorem}

\paragraph{Known inputs and the precise deduction.}
The geometric determinant input is Alper, Bogart, and Velasco
(2017), Theorem~1.7, author version printed p.~3, with proof
on p.~5: every linear determinant in at least five variables
defines a singular projective hypersurface. The permanent
singular-locus bound is the cited R-B17 Theorem~4.14, strengthened
by its Theorem~4.23 for $m\geq6$. The discriminant and its
invariance are Mukai (2003), Definition~5.20 and Corollary~5.21,
printed p.~170. Its degree and change-of-variable weight follow
from Benoist (2012), Theorem~1.3, equations (1.4)--(1.6),
printed pp.~1190--1191, specialized to one equation; the
coefficient-degree specialization is also Example~1.8, p.~1192.
The second displayed formula of that example omits the factor
$d_1$ (written $m$ here) in the variable-change degree; we use equation~(1.4),
Theorem~1.3(1.6), and scalar dilation for the correct weight.
These known theorems are not reproved. The deduction below
identifies a concrete rectangular module and its Haar/Gaussian
certificates, and states explicitly why this does not compare
the total multiplicities or survive padding. No historical
priority claim is made for the underlying discriminant separator.

\begin{proof}
Let $X=\{\operatorname{per}_m=0\}\subset\mathbb P^{N-1}$.
The codimension hypothesis implies that a general projective
$(k-1)$-plane misses $\operatorname{Sing}X$: its expected
intersection dimension is at most
$(k-1)+(N-1-k)-(N-1)=-1$.
Iterated Bertini on the smooth locus then makes its intersection
with $X$ smooth. This is the same linear-section argument used
in the cited proof of Alper--Bogart--Velasco Theorem~1.7, now
applied to the permanent's singular-locus bound. Consequently
\[
 \operatorname{Disc}_{k,m}(\operatorname{per}_m\circ A)\ne0
\]
for $A$ in a nonempty Zariski-open subset of the space of linear
maps. In contrast, the cited determinant theorem makes
$\operatorname{Disc}_{k,m}(\det_m\circ A)=0$ for every $A$.
If a restriction is the zero form, its discriminant is zero as
well, so no exceptional lower-rank map changes this conclusion.

The discriminant has coefficient degree $D=k(m-1)^{k-1}$.
Under substitution $x\mapsto Bx$ it transforms by
$\det(B)^\delta$, where $k\delta=mD$; thus
$\delta=m(m-1)^{k-1}$. Regard this determinant-character vector
as an element of the plethysm in $N$ variables using the first
$k$ coordinate vectors. It is killed by the raising operators
inside $\mathrm{GL}_k$, because it is a determinant character,
and by the raising operators from later coordinate indices,
because no such index occurs. It is therefore a highest-weight
vector of weight $(\delta^k,0^{N-k})$. Semisimplicity shows that
the module generated by this one vector is one irreducible copy
$W$ of that type, even if the ambient plethysm contains other
copies of the same type. The displayed Weyl dimension formula
is the usual formula with all equal-row factors canceled.

Every translate of $h$ vanishes on $\det_m$, and hence $W$
vanishes on its entire orbit closure. A full-rank map $A$ with
smooth permanent restriction can be completed to an invertible
change of the $N$ variables. Therefore some translate of $h$
is nonzero at $\operatorname{per}_m$. The restriction map from
$W$ to the permanent coordinate ring is equivariant and nonzero;
irreducibility makes it injective.

Apply the Schur-orthogonality projector identity used in R-02
and R-05 to this particular irreducible copy. It gives
\eqref{G19:Haar}. At least one basis vector evaluates nontrivially
at the permanent, so its energy is strictly positive. Every
basis vector evaluates to zero at the determinant. This proof
uses compact integration only after identifying the actual
copy and its two restriction maps.

Choose the primitive integer equation of the discriminant,
as permitted by the rational discriminant construction in
Benoist (2012), Theorem~1.7, p.~1191. The coefficients of
$\operatorname{per}_m\circ A$ are integer polynomials of degree
$m$ in the entries of $A$, so $H_m$ is integral and homogeneous
of degree $mD$. It is nonzero by the smooth-section argument.
The exact complex-Gaussian monomial orthogonality in R-06 gives
\eqref{G19:Gaussian}; some nonzero integer coefficient has
square at least one, and every $\alpha!$ is a positive integer.
Finally a nonzero polynomial of total degree at most $mD$
cannot vanish on the entire grid
$\{0,\ldots,mD\}^{kN}$, by successive univariate interpolation.
This proves the finite integer witness bound. It does not
give a polynomial-time search through that grid or an efficient
evaluation procedure for its permanent coefficients.

For the two limitations, the kernel on one specified copy $W$
does not determine the kernel on the whole $\lambda$-isotypic
space: other copies may survive on the determinant, while other
independent copies may vanish on the permanent. Thus no total
multiplicity inequality follows from this argument.
If $n>m$, every nonzero $k$-variable restriction of
$\ell^{n-m}f_m$ has a linear factor and a positive-degree
remaining factor. If the linear factor repeats, its hypersurface
is singular along that hyperplane. Otherwise the hyperplane
meets the remaining positive-degree hypersurface in projective
space, and all first derivatives of the product vanish at their
intersection. Here $k\geq5$ ensures that this intersection is
nonempty. A zero restriction again has zero discriminant.
Hence the degree-$n$ discriminant vanishes on the entire padded
orbit closure. The displayed module and energy cannot remove
this padding obstruction by themselves.
\end{proof}

\paragraph{An explicit smooth five-variable cubic slice.}
For $m=3$, the following matrix provides a concrete witness to
the permanent nonvanishing in G-19:
\[
 M(x)=\begin{pmatrix}
 x_1+x_2&x_3&x_4\\
 x_4&x_0&x_1\\
 x_2+x_3&x_1&x_0+x_2
 \end{pmatrix}.
\]
Writing $f=\operatorname{per}M$, direct expansion gives
\begin{align*}
 f={}&x_0^2x_1+x_0^2x_2+x_0x_1x_2+x_0x_2^2
        +x_0x_2x_4+2x_0x_3x_4\\
    &+x_1^3+x_1^2x_2+x_1x_2x_3+x_1x_3^2
        +x_1x_4^2+x_2x_3x_4.
\end{align*}
Here is a finite integer certificate of its smoothness, with the
complete matrix and two independently written reconstructions
included in the work archive. Put
$S=\mathbb Z[x_0,x_1,x_2,x_3,x_4]$, and form the coefficient
matrix of
\[
 \bigoplus_{j=0}^4 S_4\longrightarrow S_6,
 \qquad (g_j)_j\longmapsto\sum_{j=0}^4g_j\partial_{x_j}f.
\]
Order each monomial basis by decreasing lexicographic order on
its exponent tuple. The matrix has $210$ rows and $350$ columns;
column $70j+b$ corresponds to the $b$th degree-four monomial
times $\partial_{x_j}f$, with all indices starting at zero.
Keep all rows and select the following multiplier indices in
each derivative block:
\begin{center}
\begin{tabular}{r p{0.82\linewidth}}
\toprule
$j$ & Selected values of $b$ (inclusive intervals)\\
\midrule
0 & $0$--$69$\\
1 & $0$--$12$, $14$--$22$, $24$--$28$, $30$--$31$,
    $34$--$42$, $44$--$48$, $50$--$51$, $54$--$58$,
    $60$--$61$, $64$--$65$, $69$\\
2 & $0$--$12$, $15$--$22$, $25$--$28$, $35$--$42$,
    $45$--$48$, $55$--$58$\\
3 & $0$--$10$, $12$, $15$--$20$, $22$, $25$,
    $35$--$40$, $42$, $45$\\
4 & $0$--$8$, $10$, $12$, $15$--$18$, $35$\\
\bottomrule
\end{tabular}
\end{center}
The resulting $210\times210$ determinant, in the displayed
increasing row and column orders, is $94$ modulo $101$.
It is therefore a nonzero integer. After extension to
$\mathbb C$, every degree-six monomial belongs to the ideal
generated by the five partial derivatives. A common nonzero
projective zero of those derivatives would make all degree-six
monomials vanish, including the sixth power of a nonzero
coordinate, which is impossible. Thus $\{f=0\}\subset\mathbb P^4$
is smooth. This proof uses neither a numerical tolerance nor an
inference from a randomly sampled point.

The author reconstruction uses exact polynomial expansion; the
independent reconstruction multiplies all six permutations in
integer polynomial dictionaries, differentiates and convolves
the dictionaries, reproduces every one of the $73{,}500$ matrix
entries modulo $101$, and separately eliminates the selected
minor. The files are
\nolinkurl{work/root/g19_cubic/certificate.json},
\nolinkurl{work/root/g19_cubic/macaulay_degree6.u8},
\nolinkurl{work/root/check_g19_cubic_slice.py}, and
\nolinkurl{work/isolation/check_G19_cubic_independent.py}.
The row-major byte matrix has SHA-256
\nolinkurl{c739e7d83b6c43a3e48e6f20c92667fd77282af7e1240c3806d55989af0928e5}.

Consequently the specified-copy separator at $m=3$ has
coefficient degree $D=80$ and rectangular weight $(48^5)$.
Five is the minimum possible number of rows for such a
highest-weight separating equation: R-11 shows that the
four-variable determinant-cubic image is dense in the whole
cubic coefficient space. A highest-weight equation with at
most four rows lies in the corresponding coefficient
subalgebra, so vanishing on the determinant forces it to be
zero. G-19 and the displayed smooth slice supply a nonzero
five-row equation. This row-minimality assertion concerns
specified polynomial equations, and does not assert a total
multiplicity gap at weight $(48^5)$.


\begin{theorem}[G-20: The exact eventual four-row multiplicity phase]
\label{G-20}
Let $n>m\ge3$, put $a=n-m$, and define
\[
 C_m=\min\left\{\binom{m+3}{3},\,4m^2-2m+2\right\},
 \qquad N_m=\left\lfloor\sqrt{\left\lfloor C_m/2\right\rfloor}\right\rfloor.
\]
For a fixed partition direction $\nu\vdash ns$, $s\ge1$, with
at most four rows, assume the strict conditions
\[
 \nu_4>0,\qquad \nu_1>(n-m)s.
\]
Then exactly one of the following holds:
\begin{enumerate}
\item If $m<n\le N_m$, the padded-permanent multiplicity is
strictly larger than the determinant multiplicity in every
sufficiently large integer stretch $j\nu$.
\item If $n>N_m$, the determinant multiplicity is strictly
larger in every sufficiently large integer stretch $j\nu$.
\end{enumerate}
In both cases the original coordinate degree is $sj$, both
multiplicities are positive for every $j\ge\beta^2$ with
$\beta$ from G-17, and
\begin{equation}\label{G20:exact_ratio}
 \frac{m_{j\nu}(\ell^a\operatorname{per}_m)}
      {m_{j\nu}(\det_n)}
 =\kappa_\nu j^{C_m-2n^2+1}\bigl(1+O(j^{-1})\bigr),
 \qquad \kappa_\nu>0.
\end{equation}
There is no integer equality case in the exponent for the
stated values of $m,n$.

For $m\ge17$, the exact alternatives are
\begin{equation}\label{G20:large_boundary}
 \begin{aligned}
 n^2\le2m^2-m+1&\quad\Longrightarrow\quad
       \text{eventual permanent dominance},\\
 n^2\ge2m^2-m+2&\quad\Longrightarrow\quad
       \text{eventual determinant dominance}.
 \end{aligned}
\end{equation}
The small-size values are
\[
\begin{array}{c|rrrrrrr}
m&3&4&5&6&7&8&9\\ \hline
C_m&20&35&56&84&120&165&220\\
N_m&3&4&5&6&7&9&10
\end{array}
\]
\[
\begin{array}{c|rrrrrrr}
m&10&11&12&13&14&15&16\\ \hline
C_m&286&364&455&560&680&816&969\\
N_m&11&13&15&16&18&20&22.
\end{array}
\]

The corresponding exact criterion for eventual positive
four-row rectangular gaps is
\begin{equation}\label{G20:rectangular}
 m<n\le N_m,\qquad n<4m/3.
\end{equation}
Every admissible rectangular direction then has a positive
gap at every sufficiently large integer stretch. For every
$m\ge14$, condition \eqref{G20:rectangular} simplifies to
$m<n<4m/3$. The threshold $14$ is sharp for this uniform
eventual statement. No assertion is made that an individual
finite-degree gap is impossible in a direction that eventually
has determinant dominance.
\end{theorem}

\begin{proof}
R-13 and R-27 give $c_m=C_m$: sizes through sixteen fill the
four-variable form space, and every size at least seventeen
attains the row--column parameter bound. Substitute this exact
dimension in G-18. Its occurrence theorem gives positivity
of both multiplicities after $\beta^2$, and its ratio formula
becomes \eqref{G20:exact_ratio}.

The exponent is positive precisely when
$2n^2<C_m+1$, equivalently
$n^2\le\lfloor C_m/2\rfloor$. This is precisely $n\le N_m$.
For $m\ge17$, $C_m=4m^2-2m+2$ is even, so the exponent
$C_m-2n^2+1$ is odd and cannot vanish. The two adjacent integer
inequalities are exactly \eqref{G20:large_boundary}.
For $3\le m\le16$, the table is obtained by direct evaluation
of $\binom{m+3}{3}$. The only odd dimensions in that table are
$35,165,455,969$. Their half-successors are respectively
$18,83,228,485$, none of which is a square:
\[
 4^2<18<5^2,\quad9^2<83<10^2,\quad
 15^2<228<16^2,\quad22^2<485<23^2.
\]
Thus the exponent never vanishes in the small-size range
either. Its sign and the positive leading constant in
\eqref{G20:exact_ratio} prove the exact alternatives.

For completeness, the bounds are effective without waiting
for an unspecified asymptotic. With the notation of G-17 set
\[
 u=2n^2+r_\eta-13,\qquad v=C_m+r_\eta-12.
\]
If $v>u$, G-18 gives permanent dominance for every
\[
 j\ge E\bigl(4u(pE)^{u-1}+1\bigr).
\]
If $u>v$, its proof in paragraph~3 applies with the actual
dimension $v$ in place of the upper bound $\bar v$: the
permanent image degree is bounded by $(pE)^{v-1}$, and the
same standard-Veronese and residue-class sandwich gives
determinant dominance for every
\[
 j\ge E\bigl(4v(pE)^{v-1}+1\bigr).
\]
This is the same proved binomial comparison with the exact
dimension substituted, not a new claim about an uncomputed
leading coefficient.

For a rectangle $\nu=(q,q,q,q)$ with $4q=ns$, the first
strict condition is automatic and the second is
$ns/4>(n-m)s$, equivalently $n<4m/3$.
Within that range, the preceding exact sign criterion applies.
At or above its boundary, G-12 gives the opposite weak
multiplicity inequality at every stretch. This proves
\eqref{G20:rectangular}.

R-26 already proves that every $m\ge14$ and every
$m<n<4m/3$ lies in the positive dimension range; its four
small cases and uniform inequality are exact. To prove
sharpness for the eventual statement, take
\[
 m=13,\qquad n=17,\qquad s=4,\qquad\nu=(17,17,17,17).
\]
The direction is strict because $17>4(17-13)=16$.
Here $C_{13}=560$, and the ratio exponent is
\[
 C_{13}-2\cdot17^2+1=560-578+1=-17.
\]
Consequently the ratio is $\kappa_\nu j^{-17}(1+O(j^{-1}))$,
and the determinant multiplicity is larger at every sufficiently
large stretch. Any uniform threshold below fourteen would
include this pair, so it cannot have the asserted eventual
property. This argument does not decide any of its finitely
many earlier stretches.
\end{proof}

\paragraph{Scope.}
The theorem classifies the eventual sign along each fixed
four-row direction. It neither supplies a polynomial-time
positive extraction rule in growing ambient dimension nor
excludes changing-direction or isolated finite-degree gaps.
It does not establish containment of orbit closures, a general
permanent circuit lower bound, or a superlinear-padding GCT
separation. The hypotheses and explicit stretch bounds of
G-17/G-18 remain part of the conclusion.


\begin{theorem}[G-21: Maximal permanent-image dimension in every larger ambient space]
\label{G-21}
Let $m\ge17$ and $r\ge4$. Let $\mathcal P_{m,r}$ be the
Zariski closure of the forms obtained as permanents of
$m\times m$ pencils of linear forms in $r$ variables over
$\mathbb C$. Then
\begin{equation}\label{G21:dimension}
 \dim\mathcal P_{m,r}=rm^2-2m+2.
\end{equation}
Thus four variables already suffice to attain the full
parameter count modulo the row--column scaling directions;
each additional variable contributes exactly $m^2$ dimensions.
For a general coefficient specialization
$R=\mathbb C[X_{ij}]\longrightarrow S_r=\mathbb C[x_1,\ldots,x_r]$,
the corresponding graded first homology satisfies
\begin{equation}\label{G21:Tor}
 \operatorname{Tor}_1^R(R/J_m,S_r)_d=0
 \qquad(d\le m),
\end{equation}
where $J_m$ is the submaximal-permanent ideal used in R-23.
The assertion in \eqref{G21:Tor} concerns the specified graded
pieces and does not assert exactness of the entire restricted
resolution in higher degrees.
\end{theorem}

\begin{proof}
For a quaternary pencil $A$, write
$G_{ij}(A)=\partial\operatorname{per}(A)/\partial A_{ij}$.
R-23's cyclic-pencil proof shows that the $m^2$ degree-$(m-1)$
forms $G_{ij}(A)$ are linearly independent on a nonempty
Zariski open subset of the quaternary parameter space. R-27
shows that the differential image
\[
 \operatorname{span}\{x_kG_{ij}(A):1\le k\le4,\ 1\le i,j\le m\}
\]
has dimension $4m^2-2m+2$ on another nonempty open subset.
Their parameter space is irreducible, so choose one $A$
satisfying both conditions.

Regard this same pencil as an $r$-variable pencil with zero
coefficients on $x_5,\ldots,x_r$. The differential in every
new coefficient direction has the form $x_kG_{ij}(A)$,
$5\le k\le r$. For each fixed new variable these $m^2$ forms
are independent, because multiplication by $x_k$ is injective.
Their spans for different new variables, and the original
four-variable differential span, have disjoint multidegrees
in $x_5,\ldots,x_r$. Consequently the rank at this explicit
kind of parameter point is
\[
 (4m^2-2m+2)+(r-4)m^2=rm^2-2m+2.
\]
The generic image dimension is at least this rank. The
row--column torus has the same effective dimension $2m-2$
in every variable count, and gives the opposite upper bound
$rm^2-2m+2$, exactly as in R-13/R-23. This proves
\eqref{G21:dimension}. The use of a pencil with initially
unused variables is legitimate: a nonzero coefficient minor
at that point defines a nonempty open subset of the full
$r$-variable parameter space.

For the Tor statement, use the same first part of the graded
resolution and the same cited linear-syzygy theorem as in
R-23. In degree $m$, the relevant complex after specialization is
\[
 \mathbb C^{2m-2}\longrightarrow (S_r)_1^{m^2}
       \longrightarrow (S_r)_m.
\]
The old row--column relations remain independent on a nonempty
open subset: they are independent at the four-variable pencil
just used, regarded as an $r$-variable pencil. The last map
is the pencil differential and has the exact generic rank
\eqref{G21:dimension}. Its kernel therefore has dimension
$rm^2-(rm^2-2m+2)=2m-2$ and equals the image of the first
map. This proves the degree-$m$ assertion.
In degree $m-1$, the last map sends constants to the independent
$m^2$ cofactors. Below that degree its source is zero. These
give all lower-degree assertions in \eqref{G21:Tor}.
This tensor-complex argument remains meaningful when $r>m^2$;
in that case the coefficient map is an algebra homomorphism
into a larger polynomial ring, and is not called a quotient
by $m^2-r$ linear equations.
\end{proof}

\paragraph{Exact scope.}
This gives an image-dimension and low-degree derived-restriction
statement in all $r\ge4$. It does not compute individual
$r$-row multiplicities, identify all higher syzygies, or
turn the dimension statement into a polynomial-time positive
formula in growing $r$. Those are distinct remaining tasks.


\begin{proposition}[G-22: Fourier frames give stable common tensors in more rows]\label{G-22}
Let $n>m\ge r\ge2$, put $a=n-m$, and let
$\nu\vdash ns$, $s\ge1$, have $r$ positive parts. Suppose
\[
 \nu_r>0,\qquad \nu_1>as.
\]
Set $V=\mathbb C^r$, $\eta=\nu-\nu_r(1,\ldots,1)$,
$e=|\eta|$, and let $v_\eta\in S_\eta V^*$ be the unit
extremal flag vector of weight $-\eta$ used in G-17.
Then the two enhanced cones
\[
 X_Y=\{F^{\otimes s}\otimes v:F\in Y,
                  v\in\mathcal C_\eta\}^{\rm cl}
 \subseteq W=\operatorname{Sym}^s(\operatorname{Sym}^nV)
                         \otimes S_\eta V^*,
\]
where $Y=\mathcal D_{n,r}$ or
$Y=\{L^a q:q\in\mathcal P_{m,r}\}^{\rm cl}$,
have a common stable point for $\mathrm{SL}_r$.
Consequently their invariant coordinate rings have dimensions
\[
 \dim\mathbb C[X_Y]^{\mathrm{SL}_r}
 =\dim Y+r_\eta-(r^2-1),\qquad
 r_\eta=\#\{i<j:\eta_i>\eta_j\},
\]
and their degree-$j$ invariant dimensions equal the respective
multiplicities $m_{j\nu}(Y)$ in original coordinate degree $sj$.
No inequality comparing the dimensions of the two cones is needed.
\end{proposition}

\paragraph{Sources and the additional construction.}
The stability criteria are Losev (2006), Theorem~4, p.~1
(Kempf--Ness), and Mukai (2003), Theorem~7.4, printed p.~212,
with its linear-representation formulation (7.3), p.~213
and proof pp.~218--219 (Hilbert--Mumford). The exact quotient
locators are the same as in G-17. We prove the Fourier-frame
construction and preservation of every one-parameter minimum
weight under multiplication; the moment of a product of linear
forms is not assumed to equal the sum of their separate moments.

\begin{proof}
Write $x_0,\ldots,x_{r-1}$ for an orthonormal basis of $V$ and put
\[
 \mu_0=\frac{\nu_1}{s}-a,\qquad
 \mu_i=\frac{\nu_{i+1}}s\quad(1\le i<r).
\]
Every $\mu_i$ is positive and $\sum_i\mu_i=m$.
Choose a primitive $m$th root $\zeta$, and define the unit vectors
\begin{equation}\label{G22:frame}
 u_j=\sum_{i=0}^{r-1}\sqrt{\frac{\mu_i}{m}}\,
                    \zeta^{ij}x_i,\qquad 0\le j<m.
\end{equation}
The finite Fourier orthogonality identity gives
\begin{equation}\label{G22:framemoment}
 \|u_j\|^2=1,\qquad
 \sum_{j=0}^{m-1}u_ju_j^*=\operatorname{diag}(\mu_0,\ldots,\mu_{r-1}).
\end{equation}
Indeed $\sum_j\zeta^{(i-k)j}$ is $m$ for $i=k$ and zero
otherwise, since $0\le i,k<r\le m$.
Every $r$ of the vectors $u_j$ are independent: after dividing
the rows by their nonzero factors, their determinant is the
Vandermonde determinant on the distinct corresponding powers
of $\zeta$.

First work in the ordered tensor representation
\[
 \widetilde W=V^{\otimes ns}\otimes S_\eta V^*
\]
and consider the nonzero unit decomposable tensor
\[
 \xi=\bigl(x_0^{\otimes a}\otimes
             u_0\otimes\cdots\otimes u_{m-1}\bigr)^{\otimes s}
               \otimes v_\eta.
\]
The moment of a unit tensor product is the sum of the moments
of its factors: expand the infinitesimal action as a sum over
tensor positions and take the product Hermitian inner product.
By \eqref{G22:framemoment}, its diagonal contribution before
removing the trace is
\[
 sa e_0+s\mu-\eta=\nu-\eta=\nu_r(1,\ldots,1).
\]
The off-diagonal entries cancel in the Fourier sum. The
extremal weight vector $v_\eta$ has zero off-diagonal moment,
because distinct torus weights are orthogonal. Hence $\xi$
has zero $\mathrm{SL}_r$ moment map.

Its projective stabilizer is finite. A transformation preserving
the line of an ordered nonzero decomposable tensor preserves
the line of each tensor factor; this follows, for example, by
contracting all other factors with nonzero linear functionals.
If $m\ge r+1$, it consequently fixes $r+1$ full-spark lines
$[u_j]$. The first $r$ are a basis, so the transformation is
diagonal in that basis. Every coordinate of the next vector
in that basis is nonzero, since any $r$ of these $r+1$ vectors
are independent. Fixing its line makes all diagonal entries
equal. Thus the transformation is scalar.
If $m=r$, the $u_j$ are a basis and
\[
 \sum_{j=0}^{m-1}u_j=\sqrt{m\mu_0}\,x_0.
\]
All basis coefficients of $x_0$ are nonzero. Since $a>0$,
the ordered tensor also contains this padding vector and
the same argument again forces a scalar. The projective
stabilizer inside $\mathrm{SL}_r$ is therefore contained in
its finite center. Kempf--Ness now makes the affine orbit
of $\xi$ closed, and its finite stabilizer makes $\xi$ stable.

Let $q=\prod_{j=0}^{m-1}u_j\in\operatorname{Sym}^mV$,
$F=x_0^a q$, and
\[
 w=F^{\otimes s}\otimes v_\eta\in W.
\]
The diagonal $m\times m$ pencil with entries $u_j$ has both
permanent and determinant equal to $q$. Adding $a$ diagonal
entries $x_0$ proves membership of $F$ in both required
degree-$n$ cones. It remains to prove stability after the
symmetrizations; equivariance alone would not prove that.

For any nontrivial one-parameter subgroup
$\lambda:\mathbb G_m\to\mathrm{SL}_r$, diagonalize its action
and denote the smallest occurring integer weight of a nonzero
vector $z$ by $\omega_\lambda(z)$. Smallest weights add under
tensor product. They also add under multiplication of nonzero
polynomials in $\operatorname{Sym}(V)$: their smallest-weight
initial terms have nonzero product because this symmetric
algebra is an integral domain. The same statement applies
to the power $F^{\otimes s}$ in
$\operatorname{Sym}^s(\operatorname{Sym}^nV)$.
It follows that
\begin{equation}\label{G22:minweight}
 \omega_\lambda(w)
 =s\left(a\omega_\lambda(x_0)+
             \sum_j\omega_\lambda(u_j)\right)
                +\omega_\lambda(v_\eta)
 =\omega_\lambda(\xi).
\end{equation}
Both representations have scalar weight
$ns-e=r\nu_r>0$, so they are linear actions of ray type,
with irrelevant set the origin. Mukai's Theorem~7.4 says
that a nonzero vector is stable exactly when its limit
under every nontrivial $\lambda(t)$ as $t\to0$ fails to
exist. In the present weight notation, this means that
every $\omega_\lambda$ is negative. It holds for $\xi$,
and \eqref{G22:minweight} proves it for $w$. Thus $w$ is
the required common stable point.

For completeness the representation-theoretic extraction
uses the same argument as G-17 in rank $r$. The extremal
orbit cone has
$\mathbb C[\mathcal C_\eta]_j=S_{j\eta}V$ and projective
dimension $r_\eta$. The Segre--Veronese map consequently gives
\[
 \dim X_Y=\dim Y+r_\eta,\qquad
 \mathbb C[X_Y]_j=\mathbb C[Y]_{sj}\otimes S_{j\eta}V.
\]
These equalities follow from the span of translates of the
extremal tensor and from multiplication of paired coordinate
functions, as proved in G-17; they do not require normality.
The only degree-$sj$ type whose dual pairs to an invariant
with $S_{j\eta}V$ is
$\lambda=j\eta+t(1,\ldots,1)$.
The size equation $|\lambda|=nsj$ forces $t=j\nu_r$,
so that $\lambda=j\nu$. Exactness of invariants gives the
claimed multiplicity identity. The nonempty stable open
has orbit dimension $r^2-1$, which proves the quotient
dimension formula by the dimension-of-fibers theorem.
\end{proof}

\paragraph{Scope.}
The positive numbers $\mu_i$ need not be rational squares;
this is an existence witness over $\mathbb C$, not a claim
of polynomial bit complexity for a hitting set or for
multiplicity extraction. It works for every stated strict
direction and every $m\ge r$. Removing the scalar-orbit
period of the multiplied tensor requires a further step:
the trivial projective stabilizer of the ordered tensor
does not itself rule out permutations of the factors of $F$.



\begin{theorem}[G-23: Exact eventual multiplicity sign in every strict direction through $m$ rows]\label{G-23}
Let $n>m\ge17$ and $4\le r\le m$. Put $a=n-m$, and fix
$\nu\vdash ns$, $s\ge1$, with
\[
 \nu_r>0,\qquad \nu_1>as.
\]
Append no additional positive parts beyond row $r$.
Define
\[
 \begin{split}
 \eta&=\nu-\nu_r(1,\ldots,1),\qquad e=|\eta|,\qquad
 r_\eta=\#\{i<j:\eta_i>\eta_j\},\\
 u&=(r-2)n^2+r_\eta-r^2+3,\\
 v&=rm^2-2m+r+r_\eta-r^2+2,\\
 \Delta&=v-u=rm^2-(r-2)n^2-2m+r-1.
 \end{split}
\]
The integer $\Delta$ is never zero. Both multiplicities are
positive at every sufficiently large integer stretch, and
\begin{equation}\label{G23:ratio}
 \frac{m_{j\nu}(\ell^a\operatorname{per}_m)}
      {m_{j\nu}(\det_n)}
 =\kappa_\nu j^\Delta(1+O(j^{-1})),\qquad \kappa_\nu>0,
\end{equation}
through all integers $j\to\infty$, in original coordinate degree $sj$.
Thus the permanent multiplicity is eventually strictly larger
exactly when
\begin{equation}\label{G23:boundary}
 n\le N_{m,r}:=
 \left\lfloor\sqrt{\left\lfloor
       \frac{rm^2-2m+r-1}{r-2}\right\rfloor}\right\rfloor;
\end{equation}
for $n>N_{m,r}$ the determinant multiplicity is eventually strictly larger.

Here are explicit bounds. Let
\[
 \begin{aligned}
 B_n&=\binom{n+r-1}{r-1},&
 b&=\binom{B_n+s-1}{s}\dim S_\eta\mathbb C^r,\\
 L&=ns+(r-1)e,&p&=ns+e,\\
 \beta&=\left\lceil\frac{3r^2b}{8}\right\rceil
                          L^{2(r^2-1)},&
 T&=\binom{b+\beta}{b}-1,\qquad E=rT(\beta!).
 \end{aligned}
\]
The Schur dimension is the Weyl product
$\prod_{i<j}(\eta_i-\eta_j+j-i)/(j-i)$.
Both multiplicities are positive for every $j\ge\beta^2$.
If $\Delta>0$, the strict positive gap holds for every
\[
 j\ge E\bigl(4u(pE)^{u-1}+1\bigr).
\]
If $\Delta<0$, the strict reverse gap holds for every
\[
 j\ge E\bigl(4v(pE)^{v-1}+1\bigr).
\]
\end{theorem}

\paragraph{Proof dependencies and meaning.}
G-21 gives the permanent image dimension, R-28 gives the
determinant dimension and the projectively free common stable
tensor, and G-22 supplies stability for every stated direction.
The invariant-generator input is Derksen (2001), Theorem~1.1
and Proposition~1.2, printed p.~956. The Hilbert-function and
rational image-degree inputs are G-10 and its exact source
locators. The changes from G-18 are the growing row count,
the different group dimension, the new stable-point construction,
and the uniform arithmetic exclusion of equality. We give the
argument with all new parameters, rather than assuming that
the four-row constants remain valid.

\begin{proof}
Use the two enhanced cones of G-22, and let $A_D,A_P$ be
their invariant coordinate rings. They are graded domains.
Their degree-$j$ dimensions are the two specified multiplicities.
By R-28 and G-21,
\[
 \dim\mathcal D_{n,r}=(r-2)n^2+2,\qquad
 \dim\mathcal P_{m,r}=rm^2-2m+2.
\]
The padded permanent cone has dimension
$rm^2-2m+r+1$. Indeed its projective parametrization by a
padding line and a permanent form is generically injective:
on the nonempty open where the latter is irreducible of degree
$m>1$ and not divisible by the padding line, the unique
nonlinear factor and the linear factor recover both projective
inputs. The irreducible open is nonempty by the common cyclic
family in R-28. Thus padding adds $r-1$ affine-cone dimensions.
Subtracting $r^2-1$ from the enhanced dimensions in G-22 proves
\[
 \dim A_D=u,\qquad\dim A_P=v.
\]
Both are at least two: $r\le m<n$ gives
$u\ge2n^2-r^2+3>2$, and
$v\ge(r-1)m^2-2m+r+2>2$.

\paragraph{1. Invariant degrees and absence of a period.}
The ambient representation has dimension $b$. Its action
coefficients on $\mathrm{SL}_r$ have degree at most $L$:
the covariant factor contributes $ns$, and the entries of
the inverse transpose have degree $r-1$, contributing at
most $(r-1)e$ on the Schur factor. Its kernel is finite
because it is contained in the finite stabilizer of the
common stable tensor. In the embedding
$\mathrm{SL}_r\subset\mathbb A^{r^2}$, the defining equation
has degree $r$ and the group dimension is $r^2-1$.
Derksen's Proposition~1.2 therefore gives
\[
 \sigma(W)\le rL^{r^2-1},
\]
and Theorem~1.1 gives the displayed integer generator bound
$\beta$. Exactness of invariants makes $A_D,A_P$ graded
quotients of this same ambient invariant ring.

Choose homogeneous generators $g_i$ of degrees
$1\le e_i\le\beta$. R-28 supplies a common stable point
$w$ with
\[
 gw=cw,\quad g\in\mathrm{SL}_r,\ c\in\mathbb C^*
                  \quad\Longrightarrow\quad c=1.
\]
For precision, its projective stabilizer is the center
$\mu_r$, and this center acts with exponent
$ns-e=r\nu_r$, hence trivially.
The degrees of the generators nonzero at $w$ have greatest
common divisor one. If that divisor were $h>1$, a primitive
$h$th root $c$ would make every generator have equal values
at $w$ and $cw$, although their two closed orbits are distinct.
This contradicts closed-orbit separation by invariants
(Mukai, Theorem~5.3 and Corollary~5.4, pp.~161--162).
There is at least one such generator by the same separation
applied to the closed orbits of $w$ and zero.

The numerical-semigroup proof in G-18, part~2, now applies
with this actual $\beta$: shortest paths modulo the least
degree have at most $\beta-1$ edges, so each residue has a
representative of weight less than $\beta^2$. Hence every
integer $j\ge\beta^2$ is a sum of degrees nonzero at $w$.
The corresponding product gives an ambient homogeneous
invariant $h_j$ nonzero on both cones. This proves both
occurrence assertions at every such integer $j$.

\paragraph{2. Standard Veroneses and effective comparison.}
There are at most $T$ chosen positive-degree generators,
because the whole ambient coefficient-polynomial space
through degree $\beta$ has dimension $T$.
The generator-grouping argument of G-11 applies with
$E=rT\beta!$: remove pure generator blocks of degree
$\beta!$, leaving a residual monomial of degree less than
$T\beta!$, then complete it to a block of degree $E$ and
partition the remaining pure blocks. Thus $A_D^{(E)}$ and
$A_P^{(E)}$ are standard graded domains. Since $T\ge\beta$
and $\beta!\ge\beta$, we have $E\ge\beta^2$.

Both enhanced cones have dense homogeneous parametrizations
of degree $p$ in independent variables. They are
\[
 (Z,B)\longmapsto\det\left(\sum_{i=1}^r Z_i x_i\right)^{\otimes s}
                  \otimes S_\eta(B)v_\eta
\]
and
\[
 (Z,H,B)\longmapsto
 \left(H^a\operatorname{per}_m\left(\sum_{i=1}^rZ_i x_i\right)\right)^{\otimes s}
                  \otimes S_\eta(B)v_\eta.
\]
The parameter $H$ is an independent linear form and $B$ an
independent endomorphism of $V^*$. The first factors have
parameter degree $ns$, and the second has degree $e$.
If $e=0$ it is constant and may be omitted. These are parameter
degrees, distinct from the action degree $L$.
G-10's base-locus-aware degree argument bounds the two
projective invariant-Veronese image degrees by
\[
 (pE)^{u-1}\quad\hbox{and}\quad(pE)^{v-1},
\]
respectively.

Write $j=Ek+t$, $0\le t<E$. For either domain $A$, multiplication
by the common nonzero invariants $h_{E+t}$ and $h_{2E-t}$ gives
\begin{equation}\label{G23:sandwich}
 \dim A_{E(k-1)}\le\dim A_j\le\dim A_{E(k+2)}\qquad(k\ge1).
\end{equation}
If $v>u$, linear Noether normalization of $A_P^{(E)}$ and
the Chardin bound on $A_D^{(E)}$ yield
\[
 \dim(A_P)_{E(k-1)}\ge\binom{k+u-1}{u},\qquad
 \dim(A_D)_{E(k+2)}\le(pE)^{u-1}\binom{k+u+1}{u-1}.
\]
Their binomial ratio is
\[
 \frac{k(k+1)(k+2)}{u(k+u)(k+u+1)}
                 \ge\frac{k}{4u}\qquad(k\ge u+1).
\]
The asserted positive threshold follows from
\eqref{G23:sandwich}. If $u>v$, exchange the domains and
replace $u$ by the actual dimension $v$ in this argument.
This gives the asserted reverse threshold without an upper
bound substituted for the permanent dimension.

For a standard graded domain of dimension $q\ge2$, the
leading Hilbert term is
$\deg(\operatorname{Proj} A^{(E)})k^{q-1}/(q-1)!$,
with error $O(k^{q-2})$. Applying \eqref{G23:sandwich}
shows, through all integers, that
\[
 \dim A_j=c_A j^{q-1}+O(j^{q-2}),\qquad c_A>0.
\]
Taking the ratio for $A_P,A_D$ proves \eqref{G23:ratio}.

\paragraph{3. The integer dimensions can never be equal.}
If $r$ is even, $\Delta$ is odd. Suppose $r$ is odd and
$\Delta=0$. Reduction modulo two forces $m,n$ to have
the same parity. Whether both are even or both are odd,
reduction modulo four then gives
\[
 0=\Delta\equiv r-1\pmod4,
\]
so $r\equiv1\pmod4$. Reduction modulo $r-2$ gives
$2m^2-2m+1\equiv0$, and hence
\[
 (2m-1)^2\equiv-1\pmod{r-2}.
\]
But $r-2\equiv3\pmod4$, so one of its prime divisors
$p$ is congruent to three modulo four. A square root of
$-1$ modulo $p$ is impossible: its nonzero residue would
have $(p-1)$st power both $1$, by Fermat's theorem, and
$(-1)^{(p-1)/2}=-1$. This contradiction proves
$\Delta\ne0$. The precise integer boundary
\eqref{G23:boundary} follows immediately.
\end{proof}

\paragraph{Rectangular consequence and limits.}
For a rectangular direction, the strict condition is
$n<rm/(r-1)$. Throughout that range the exponent is positive:
\[
 \Delta>
 \frac{r m^2}{(r-1)^2}-2m+r-1
 =\frac r{(r-1)^2}\left(m-\frac{(r-1)^2}{r}\right)^2
                                      +\frac{r-1}{r}>0.
\]
Thus every such rectangle gives a gap at every sufficiently
large admissible coordinate degree. For a fixed $m,n,r$,
one may take $s=r/\gcd(n,r)$ and
$\nu=(ns/r)^r$, so all admissible degrees are $sj$.
The theorem concerns fixed strict directions with $r\le m$.
It does not decide the face $\nu_1=as$, the range $r>m$,
every earlier finite degree, or a uniform polynomial-time
positive rule when the parameters grow. Its bounds and
non-occurrence-obstruction conclusion retain these limits.



\begin{theorem}[G-24: The exact eventual rectangular range through $m$ rows]\label{G-24}
Let $n>m\ge17$ and $4\le r\le m$. At every sufficiently
large admissible coordinate degree $d$ (meaning $r\mid nd$),
the rectangular multiplicity gap
\[
 m_{(nd/r)^r}(\ell^{n-m}\operatorname{per}_m)
       >m_{(nd/r)^r}(\det_n)
\]
holds if and only if
\begin{equation}\label{G24:range}
 m<n<\frac{rm}{r-1}.
\end{equation}
Inside this range, G-23 supplies the explicit threshold.
Above the boundary, the permanent rectangular multiplicity
is zero in every positive admissible degree and the determinant
multiplicity is positive in every sufficiently large admissible
degree. At the equality boundary, the determinant multiplicity
is strictly larger in every sufficiently large degree.
No assertion excluding all earlier boundary gaps is made for
$r\ge5$; G-12 supplies that stronger weak inequality for $r=4$.

For an explicit equality-boundary threshold, write
$a=n-m$, so that $n=ra$, $m=(r-1)a$, and define
\[
 \begin{split}
 u&=(r-2)n^2-r^2+3,\\
 v&=\begin{cases}
       \binom{m+2}{2}-8,&r=4,\\
       (r-1)m^2-2m-r^2+2r+2,&r\ge5.
      \end{cases}
 \end{split}
\]
Let
\[
 b=\binom{n+r-1}{r-1},\qquad
 \beta=\left\lceil\frac{3r^2b}{8}\right\rceil n^{2(r^2-1)},
 \qquad T=\binom{b+\beta}{b}-1,\qquad E=rT(\beta!).
\]
Then $u>v\ge2$, and determinant dominance holds for every
\begin{equation}\label{G24:threshold}
 d\ge E\bigl(4v(nE)^{v-1}+1\bigr).
\end{equation}
\end{theorem}

\begin{proof}
The strict interior is the rectangular consequence of G-23.
For the other cases, set $a=n-m$ and use $r$ variables
$x_0,\ldots,x_{r-1}$. A nonzero padding form can be taken
to $x_0$ by a special-linear change of variables. The
one-parameter subgroup
\[
 \lambda(t)=\operatorname{diag}(t^{r-1},t^{-1},\ldots,t^{-1})
\]
gives weight $(r-1)a-m+rj=ra-n+rj$ to a monomial of
$x_0^a q$ coming from the $x_0^j$ part of $q$.
If $(r-1)n>rm$, all these weights are positive. The form
therefore tends to zero and every positive-degree invariant
vanishes. This proves the zero assertion above the boundary.
For eventual determinant occurrence, take the common free
point of R-28 with smaller size $n-1$ and padding one.
For the minimal admissible direction use
$s_0=r/\gcd(n,r)$ and $\nu=(ns_0/r)^r$.
Its strict condition is $ns_0/r>s_0$, since $n>r$.
R-28 consequently gives positive determinant multiplicity
at every sufficiently large integer stretch, that is, every
sufficiently large admissible coordinate degree.

\paragraph{1. Invariant functions inject into one fewer variables at equality.}
Now suppose $(r-1)n=rm$. Then $n=ra$ and $m=(r-1)a$.
The same limit is
\begin{equation}\label{G24:limit}
 \lim_{t\to0}\lambda(t)x_0^a q
                =x_0^a q(0,x_1,\ldots,x_{r-1}).
\end{equation}
Put $q_0=r-1$ for the number of variables after restriction,
and write
\[
 A_P=\mathbb C[\{L^a q:q\in\mathcal P_{m,r}\}^{\rm cl}]^{\mathrm{SL}_r},
 \qquad C=\mathbb C[\mathcal P_{m,q_0}]^{\mathrm{SL}_{q_0}}.
\]
Restriction along $q\mapsto x_0^a q$ gives a graded algebra
homomorphism
\begin{equation}\label{G24:injection}
 \rho:A_P\longrightarrow C.
\end{equation}
It is well defined because the indicated padded forms belong
to the original cone, and the subgroup
$\operatorname{diag}(1,\mathrm{SL}_{q_0})$ fixes $x_0$.
Its grading is unchanged: coefficients of $x_0^a q$ depend
linearly on coefficients of $q$.

To prove injectivity, let an invariant function $h$ restrict
to zero. For any original padded form, take its nonzero
padding line to $x_0$. The function $h$ is constant along
$\lambda(t)$ and hence, by regularity, at its limit
\eqref{G24:limit}. Restricting a permanent pencil to
$x_0=0$ puts the residual form in $\mathcal P_{m,q_0}$,
including its closure points. Thus the value at the limit
is zero. The dense padded image has zero value everywhere,
so $h=0$ on its closure. The same statement holds in degree
zero. This proves \eqref{G24:injection} without claiming
that every lower-rank invariant lifts.

Let $A_D=\mathbb C[\mathcal D_{n,r}]^{\mathrm{SL}_r}$.
For $r\ge5$ put $B=C$. For $r=4$ use instead
\[
 B=\mathbb C[\operatorname{Sym}^m\mathbb C^3]^{\mathrm{SL}_3}.
\]
In this latter case exactness of invariants gives a graded
surjection $B\to C$, not an asserted injection into the
full ternary space. Consequently in every case
\begin{equation}\label{G24:hilbertupper}
 \dim(A_P)_d\le\dim B_d.
\end{equation}

\paragraph{2. The two comparison rings have known dimensions and no period.}
The equality hypotheses and $r\le m$ imply $a\ge2$.
Apply R-28's common-free-point statement, which allows
row count three, first with degree $n$, size $n-1$, row
count $r$ and rectangular direction $(a^r)$ in degree one.
Apply it a second time with degree $m$, size $m-1$, row
count $q_0$ and direction $(a^{q_0})$.
Both applications satisfy the strict inequality $a>1$ and
the size condition. They supply a stable point with no
nontrivial scalar-orbit symmetry in $A_D$ and in the
permanent cone underlying $C$. The latter is also such a
point in the full ambient ternary representation when
$r=4$.

It follows that the stable quotients have dimensions
\[
 \dim A_D=(r-2)n^2+2-(r^2-1)=u,
 \qquad \dim B=v.
\]
For $r\ge5$, use G-21 in $q_0\ge4$ variables; for $r=4$
the whole ternary coefficient space has dimension
$\binom{m+2}{2}$ and the stable orbit dimension is eight.
The absence of scalar-orbit symmetry and invariant
separation give gcd one for the nonzero invariant degrees,
just as in R-28. We use this separately for the two rings;
no common invariant in their different ambient spaces is
being assumed.

The dimension inequality is strict. For $r\ge5$ it reads
\[
 u-v=\frac{r^2-3r+1}{(r-1)^2}m^2+2m-2r+1>0,
\]
since $m\ge r$ and $r^2-3r+1>0$.
For $r=4$ it is
\[
 u-v=\frac{55m^2-27m-108}{18}>0\qquad(m\ge17).
\]
The displayed dimensions are at least two in the stated range.

\paragraph{3. An explicit comparison at every integer degree.}
Derksen's general bound, specialized as in G-23 with $s=1$
and $\eta=0$, gives the displayed $\beta$ for the degree-$n$
ambient representation in $r$ variables. It also bounds the
generators for the degree-$m$ representation in $q_0<r$
variables: its vector-space dimension, action degree and
group-dimension exponent are all no larger. Exactness of
invariants passes these bounds to $A_D$ and to $B$.
The chosen $T$ bounds both generator counts. Hence the
$E$th Veroneses of both rings are standard graded by the
same generator-grouping argument as G-23.

The gcd-one argument and the bound $\beta$ give nonzero
elements in every degree at least $\beta^2$ in each ring
separately. For $d=Ek+t$, $0\le t<E$, multiplication by
elements of degrees $E+t$ and $2E-t$ therefore gives the
individual sandwiches
\[
 \dim A_{E(k-1)}\le\dim A_d\le\dim A_{E(k+2)},
 \qquad A=A_D\text{ or }B.
\]
The projective degree of $B^{(E)}$ is at most
$(nE)^{v-1}$. If $r\ge5$, compose its invariant-coordinate
map with the permanent-pencil parametrization of degree
$m<n$ and use G-10. If $r=4$, use the identity parametrization
of the full ternary coefficient space, of degree one, and
the same larger bound. No dominance of the ternary
permanent parametrization is needed here.

Noether normalization on $A_D^{(E)}$, whose dimension is
at least $v+1$, and the degree/Hilbert bound on $B^{(E)}$
now give
\[
 \dim(A_D)_{E(k-1)}\ge\binom{k+v-1}{v},\qquad
 \dim B_{E(k+2)}\le(nE)^{v-1}\binom{k+v+1}{v-1}.
\]
The same exact ratio as in G-23 is at least $k/(4v)$ for
$k\ge v+1$. For $k\ge4v(nE)^{v-1}+1$, this proves
$\dim(A_D)_d>\dim B_d\ge\dim(A_P)_d$ in every residue.
All degrees are admissible at this boundary because $r\mid n$.
The invariant multiplicities are the claimed rectangles,
by the same single-label extraction and inheritance as
G-11/G-22. This proves \eqref{G24:threshold} and the theorem.
\end{proof}

\paragraph{Exact scope.}
The equality-boundary argument is an injection of actual
invariant coordinate rings followed by a Hilbert-dimension
comparison. It neither identifies those rings nor replaces
orbit-closure multiplicities by ambient plethysm coefficients.
The theorem closes the eventual rectangular boundary through
$m$ rows in the stated size range. It leaves nonrectangular
boundary faces, earlier finite-degree gaps at $r\ge5$,
more than $m$ rows, and efficient positive extraction open.



\begin{theorem}[G-25: Exact strict-direction comparison through $2m-2$ rows]\label{G-25}
Every assertion and every explicit bound of G-23 remains valid when its
row range $4\le r\le m$ is replaced by
\[
             4\le r\le2m-2,
\]
with $n>m\ge17$ and the same strict conditions
$\nu_r>0$, $\nu_1>(n-m)s$, $\nu\vdash ns$.
In particular the exact eventual exponent is still
\[
 \Delta=rm^2-(r-2)n^2-2m+r-1\ne0,
\]
and both multiplicities are positive at every integer stretch
$j\ge\beta^2$, with the same definitions of $b,L,p,\beta,T,E,u,v$.
The effective thresholds for strict comparison, and the exact
floor boundary $N_{m,r}$, are those of G-23.

In the newly added range $m<r\le2m-2$, the exponent is positive
if and only if
\[
                   r=n=m+1.
\]
For these parameters it equals one. Thus, for every $m\ge17$,
the explicit nonrectangular direction
\[
 s=2,\qquad \nu=(3,2^{m-1},1),\qquad n=r=m+1
\]
has an eventually positive genuine-padding multiplicity gap in
every sufficiently large coordinate degree $2j$, and its
permanent-to-determinant multiplicity ratio is
$\kappa_\nu j(1+O(j^{-1}))$, with $\kappa_\nu>0$.
Every other pair of integer parameters in this newly added row
range has eventual determinant dominance in every fixed strict
direction.
\end{theorem}

\begin{proof}
R-29 gives both the common stable tensor and a common stable tensor
with no nontrivial scalar multiple in its orbit throughout the
enlarged range. Its enhanced-coordinate extraction agrees exactly
with that in G-22. The dimension formulas used by G-23 already
hold here: G-21 proves the permanent formula for every $r\ge4$
when $m\ge17$, and R-28 proves the determinant formula and the
padding dimension for every $r\ge4$ in the present size range.
Consequently the invariant-domain dimensions are precisely the
same $u,v$ displayed in G-23.

The only numerical estimate in that proof that used $r\le m$
beyond the stable-point input was $u,v\ge2$. We supply the
replacement explicitly. Since $r-2\ge r/2$, $r\le2m-2$,
$n\ge m+1$ and $r_\eta\ge0$,
\[
 \begin{aligned}
 u&=(r-2)n^2-r^2+r_\eta+3
       \ge\frac r2(m^2-2m+5)+3>2,\\
 v&=r(m^2-r+1)-2m+r_\eta+2
       \ge r(m^2-2m+3)-2m+2>2.
 \end{aligned}
\]
The last inequality follows already from $r\ge4$ and $m\ge17$.

All subsequent steps of G-23 use only these exact dimensions,
the ambient representation in rank $r$, and the two properties
of the common tensor supplied by R-29. Derksen's action-degree
and generator bounds are unchanged. The same nonzero generator
degrees have gcd one by closed-orbit separation, so the common
invariants $h_j$ exist for every $j\ge\beta^2$. The same
generator grouping gives the standard $E$th Veroneses, and the
same degree-$p$ parameter maps give their image-degree bounds.
Their every-residue Hilbert sandwiches, the Noether lower bound,
and the Chardin upper bound therefore prove the stated effective
thresholds and ratio with these actual values of $u,v$.
The parity and nonsquare argument excluding $\Delta=0$ in G-23
was valid for every integer $r\ge4$, with no upper row bound.
This proves the full extension, including all its constants.

For the additional sign classification, fix $m<r$.
The exponent is strictly decreasing as an integer function of
$n\ge m+1$, and direct substitution gives
\[
       \Delta\big|_{n=m+1}=2m(m-r+1)+1.
\]
If $r\ge m+2$, this is at most $1-2m<0$.
If $r=m+1$, its value at $n=m+1$ is one, while increasing $n$
to $m+2$ subtracts $(m-1)(2m+3)>1$. This proves the asserted
unique positive parameter case. Finally the displayed partition
has $m+1$ positive parts and total size $2m+2=ns$; its first
part is $3>2=(n-m)s$. It therefore meets the strict hypotheses,
and substituting $\Delta=1$ proves its claimed ratio and gap.
\end{proof}

\paragraph{Scope.}
The new positive example specifies a total multiplicity inequality,
not merely one nonzero equation copy. It is an eventual statement
with the explicit G-23 threshold and does not assert the comparison
at every earlier degree. No rectangle satisfies the strict
condition when $r>m$ and $n>m$ are integers: then
$rm/(r-1)\le m+1$. The rectangular boundary result G-24 keeps
its stated range $r\le m$. General nonrectangular equality
faces, more than $2m-2$ rows, changing directions, and
polynomial-resource coefficient extraction remain separate gates.


\begin{corollary}[G-25, endpoint extension from R-30]\label{G25:endpoints}
All assertions of G-25, including every explicit G-23 constant,
hold throughout $4\le r\le2m$. The unique positive case among
the added rows $r>m$ is still $r=n=m+1$.
\end{corollary}

\begin{proof}
R-30 supplies the common stable no-scaling points at
$r=2m-1$ and $r=2m$. The exact determinant and permanent
dimensions in R-28 and G-21 hold at both endpoints. For
completeness, the required lower bounds on their enhanced
invariant dimensions are, using $r\le2m$ and $n\ge m+1$,
\[
 \begin{aligned}
 u&\ge\frac r2(m-1)^2+3>2,\\
 v&\ge r(m-1)^2-2m+2>2,
 \end{aligned}
\]
because $r-2\ge r/2$ and $m\ge17$.
The exact proof of G-25 then applies with the unchanged
rank-$r$ representation, action degree, parameter degree,
generator bound, common degree semigroup and standard Veronese.
Its arithmetic sign classification used no upper row bound.
This proves the assertion at both new endpoints.
\end{proof}


\begin{theorem}[G-26: The complete first-row padding boundary through $m$ rows]\label{G-26}
Let $a=n-m>0$. For every coordinate degree $d\ge1$ and partition
$\lambda\vdash nd$ of length at most $r$, the padded permanent
multiplicity satisfies
\begin{equation}\label{G26:pieri_bound}
 \begin{cases}
 m_\lambda(\ell^a\operatorname{per}_m)=0,&\lambda_1<ad,\\[2pt]
 m_\lambda(\ell^a\operatorname{per}_m)
   \le m_{(\lambda_2,\ldots,\lambda_r)}(\operatorname{per}_m),
       &\lambda_1=ad.
 \end{cases}
\end{equation}
The multiplicity on the second line can be computed on the
permanent image in $r-1$ variables. This bound holds in every
finite degree over $\mathbb C$; no lifting or equality of these
multiplicity spaces is asserted.

Now suppose $n>m\ge17$ and $4\le r\le m$. Fix
$\nu\vdash ns$, $s\ge1$, with exactly $r$ positive parts.
If $\nu_1=(n-m)s$, the determinant multiplicity is strictly
larger than the padded permanent multiplicity at every
sufficiently large integer stretch $j$. The same conclusion
holds when $\nu_1<(n-m)s$, in which case the padded multiplicity
is zero at every positive stretch.

Consequently, in this row range every fixed full-row direction
has an exact eventual comparison: the padded permanent
multiplicity is eventually larger if and only if
\[
 \nu_1>(n-m)s
 \quad\hbox{and}\quad
 rm^2-(r-2)n^2-2m+r-1>0.
\]
In every other case the determinant multiplicity is eventually
strictly larger. In the strict-padding case the bounds are
those of G-23; an explicit equality-face bound follows below.
\end{theorem}

\paragraph{An explicit threshold on the equality face.}
Assume $\nu_1=as$, put $q=r-1$, and define
\[
 \begin{aligned}
 \mu&=(\nu_2,\ldots,\nu_r),&
 \eta&=\nu-\nu_r\mathbf1_r,&
 \theta&=\mu-\nu_r\mathbf1_q,\\
 e&=|\eta|,&e'&=|\theta|,&
 r_\eta&=\#\{i<j:\eta_i>\eta_j\},
 \end{aligned}
\]
with $r_\theta$ defined in the same way. Set
\[
 \begin{aligned}
 u&=(r-2)n^2-r^2+3+r_\eta,\\
 v&=\begin{cases}
       \binom{m+2}{2}-8+r_\theta,&r=4,\\
       (r-1)m^2-2m-r^2+2r+2+r_\theta,&r\ge5,
      \end{cases}\\
 b_D&=\binom{\binom{n+r-1}{r-1}+s-1}{s}\dim S_\eta\mathbb C^r,\\
 b_B&=\binom{\binom{m+q-1}{q-1}+s-1}{s}\dim S_\theta\mathbb C^q,\\
 b&=\max\{b_D,b_B\},&
 L&=ns+(r-1)e,&p&=ns+e,\\
 \beta&=\left\lceil\frac{3r^2b}{8}\right\rceil L^{2(r^2-1)},&
 T&=\binom{b+\beta}{b}-1,& E&=rT\beta!.
 \end{aligned}
\]
Then $u>v\ge2$, and strict determinant dominance holds for every
\begin{equation}\label{G26:threshold}
                  j\ge E\bigl(4v(pE)^{v-1}+1\bigr),
\end{equation}
in original coordinate degree $sj$. All Schur dimensions here
are given by the Weyl products used in G-23.

\paragraph{Background used.}
The Pieri formula is cited from Bump and Gamburd (2006),
\S2.1, equation with source label \nolinkurl{e:pieri}, supplied
\nolinkurl{acp20.tex}, lines 404--422. That source identifies it
with Macdonald's equation I~(5.16), p.~73, and Bump's
Theorem~42.4; we inspected the supplied Bump--Gamburd statement.
All other inputs are the exact invariant-extraction, degree,
and Hilbert bounds already used and located in G-22--G-24.

\begin{proof}
We first prove the finite-degree inequality and its change of
ambient dimension. We then compare two different enhanced
invariant algebras, with their degree semigroups treated separately.

\paragraph{1. Pieri removes exactly the first row.}
Let $V=\mathbb C^r$ and let $Y$ be the padded permanent image.
The parameter morphism
\[
 V\times\mathcal P_{m,r}\longrightarrow Y,
                    \qquad(\ell,F)\longmapsto\ell^a F
\]
has dense image by the definition of $Y$. Pullback of degree-$d$
coordinate functions is therefore an injective equivariant map
\begin{equation}\label{G26:parameter_injection}
 \mathbb C[Y]_d\lhook\joinrel\longrightarrow
       \operatorname{Sym}^{ad}V^*\otimes\mathbb C[\mathcal P_{m,r}]_d.
\end{equation}
The two parameter degrees are $(ad,d)$; they are not both $d$.
Use the contragredient convention for coordinate representations
fixed in the introduction. By complete reducibility over
$\mathbb C$ and Pieri, a copy of $S_\lambda V^*$ in the right
side comes from $S_\alpha V^*$ precisely when $\lambda/\alpha$
is a horizontal strip of size $ad$, with multiplicity one.

There are only $\lambda_1$ columns. If $\lambda_1<ad$, no such
strip exists. If $\lambda_1=ad$, the strip must contain exactly
one box in every column. Thus
\[
 \alpha'_k=\lambda'_k-1\quad(1\le k\le\lambda_1),
 \qquad
                 \alpha=(\lambda_2,\ldots,\lambda_r).
\]
This is the unique possibility, and its size is $md$.
Taking multiplicities in \eqref{G26:parameter_injection} gives
the asserted zero and upper bound.

For completeness, the multiplicity of this tail partition
equals that on $\mathcal P_{m,r-1}$. A highest-weight coordinate
polynomial with last weight zero uses only coefficient variables
of monomials not involving the last original variable; all these
coefficient weights are nonnegative in the polynomial convention.
Its vanishing on the $r$-variable permanent image is equivalent
to vanishing on the $(r-1)$-variable image: restriction of every
$r$-variable matrix pencil gives a pencil with one fewer variable,
and every such smaller pencil occurs by setting the unused
coefficient matrix to zero. The highest-weight coefficient
subalgebras and their ideal intersections consequently agree.
This is precisely the inheritance argument stated in the
introduction, now applied to the tail shape of length at most
$r-1$. It proves the last assertion of
\eqref{G26:pieri_bound} without replacing an actual coordinate
ring by its ambient plethysm ring.

\paragraph{2. The equality face gives two stable comparison algebras.}
Now assume the displayed equality-face hypotheses. Because
$\nu_1=as$ is the largest of the $r$ parts,
\[
 ns\le ras,\qquad
 n\ge\frac{rm}{r-1},\qquad
 a\ge\frac{m}{r-1}>1.
\]
Thus the integer $a$ is at least two. The tail $\mu$ has size
$ms$, exactly $q=r-1$ positive parts, and
\[
                    \mu_1\ge\frac{ms}{r-1}>s.
\]
Let $A$ be the enhanced invariant coordinate algebra for
the determinant cone $\mathcal D_{n,r}$, outer stretch $s$,
and flag weight $\eta$, as in G-22. Then
\[
                       \dim A_j=m_{j\nu}(\det_n).
\]
For $r\ge5$, let $B$ be the corresponding enhanced invariant
algebra for $\mathcal P_{m,q}$ and flag weight $\theta$.
For $r=4$, take instead the enhanced algebra for the full
ternary degree-$m$ coefficient space. In that case exactness
of invariants gives a degreewise surjection from $B$ to the
enhanced permanent algebra. In either case, part~1 gives
\begin{equation}\label{G26:degreewise_bound}
             m_{j\nu}(\ell^a\operatorname{per}_m)\le\dim B_j.
\end{equation}

Both $A$ and $B$ have a stable point with the no-scaling property.
For $A$, apply R-28 with degree $n$, matrix size $n-1$, row
count $r$, and the direction $\nu$: its strict condition is
$\nu_1=as>s$. The hypotheses $r\le n-1$ hold since
$r\le m<n$. For $B$, use degree $m$, matrix size $m-1$,
row count $q=r-1$, and the direction $\mu$. Its strict
condition is $\mu_1>s$, and $q\le m-1$.
The smaller padded permanent embeds in $\mathcal P_{m,q}$
by block diagonal padding with one linear entry. At $q=3$
the same point also lies in the full ternary coefficient space.
The common-point clause of R-28 applies in row count three
and does not require size at least seventeen for that clause.
The enhanced representations, and their scalar characters,
are exactly those just specified.

It follows from G-22's extraction and the stable quotient
dimension formula that
\[
                        \dim A=u,\qquad\dim B=v.
\]
Here R-28 supplies the determinant image dimension, G-21
supplies the permanent dimension for $q\ge4$, and the ternary
ambient dimension is $\binom{m+2}{2}$ for $q=3$.
No common invariant between two different groups is being
asserted. The no-scaling property gives gcd one for the
nonzero degrees in each ring separately, by the closed-orbit
separation argument already proved in R-28.

\paragraph{3. The determinant quotient has strictly larger dimension.}
Subtracting a constant from every part does not change flag
dimension. Deleting the first part removes exactly its strict
comparisons with the others, so
\[
 r_\eta-r_\theta
       =\#\{j>1:\nu_1>\nu_j\}\ge0.
\]
For $r\ge5$ this gives
\[
 \begin{aligned}
 u-v&=(r-2)n^2-(r-1)m^2+2m-2r+1+r_\eta-r_\theta\\
 &\ge\frac{r^2-3r+1}{(r-1)^2}m^2+2m-2r+1>0,
 \end{aligned}
\]
since $m\ge r$ and $r^2-3r+1>0$. For $r=4$,
\[
 u-v=2n^2-\binom{m+2}{2}-5+r_\eta-r_\theta
       \ge\frac{55m^2-27m-108}{18}>0.
\]
The explicit formulas give $v\ge2$; for instance at $r\ge5$
one has $v\ge3m^2-2m+2r+2>2$ by $r\le m$, and the ternary
case is immediate. Thus $u\ge v+1$.

\paragraph{4. Effective comparison in every original stretch.}
The two ambient representation dimensions are $b_D,b_B$.
Their action degrees are at most $L$: the smaller one has
degree $ms+(q-1)e'$, and
$e-e'=\nu_1-\nu_r\ge0$. Derksen's bounds used in G-23
therefore give the displayed common bound $\beta$ for the
two invariant generator degrees, even though the groups and
representations differ. Their numbers of generators are each
at most $T$. The identical generator grouping proves that
their $E$th Veroneses are standard graded. Each no-scaling
point gives nonzero elements in every degree at least
$\beta^2$, and $E\ge\beta^2$.

The determinant enhanced cone has its degree-$p$ dense
homogeneous parameter map from G-23. For $r\ge5$, the lower
permanent map has parameter degree $ms+e'\le p$. For $r=4$,
the full coefficient-space parameter has degree $s+e'\le p$.
The base-locus-aware image-degree bound thus gives
\[
                  \deg\operatorname{Proj} B^{(E)}\le(pE)^{v-1}.
\]
Write $j=Ek+t$, $0\le t<E$. In each ring separately, multiply
by its nonzero elements of degrees $E+t$ and $2E-t$ to obtain
the same two-sided Hilbert sandwich as G-23. Linear Noether
normalization and the Chardin upper bound yield
\[
 \dim A_j\ge\binom{k+v-1}{v},\qquad
 \dim B_j\le(pE)^{v-1}\binom{k+v+1}{v-1}.
\]
Their binomial ratio is at least $k/(4v)$ for $k\ge v+1$.
Condition \eqref{G26:threshold} implies
$k>4v(pE)^{v-1}$, so the first bound is strictly larger.
Together with \eqref{G26:degreewise_bound}, this proves the
equality-face conclusion at every stated integer stretch.

If instead $\nu_1<as$, part~1 makes every padded multiplicity
zero. The determinant multiplicity is eventually positive:
$\nu_1\ge ns/r>s$ since $r\le m<n$, so the first R-28
application in part~2 still applies. Its generator-semigroup
bound gives occurrence for every sufficiently large $j$.
Finally combine these two cases with G-23 in the strict
half-space and its proof that $\Delta$ never vanishes. This
proves the exact eventual classification in the theorem.
\end{proof}

\paragraph{Scope.}
The Pieri upper bound is an all-degree reduction and holds
beyond the row range of the eventual comparison. The strict
boundary comparison proved here uses $4\le r\le m$ and
$m\ge17$. It does not assert a positive-degree lifting from
the lower permanent algebra, eliminate earlier boundary gaps
when $r\ge5$, or settle boundary faces with more than $m$
rows. It is independent of the conditional analytic inputs.


\begin{theorem}[G-27: Nonrectangular boundary comparison through $2m-1$ rows]
\label{G-27}
Let $n>m\ge17$, $m<r\le2m-1$, and let
$\nu\vdash ns$, $s\ge1$, have exactly $r$ positive parts
and be nonrectangular. Put $a=n-m$. If $\nu_1\le as$,
the determinant multiplicity is strictly larger than the
padded permanent multiplicity at every sufficiently large
integer stretch. On the equality face the explicit threshold
below applies. Below the face the padded multiplicity is zero
at every positive stretch, and the determinant occurs at every
stretch beyond its explicit invariant-generator bound squared.

Together with G-26 and G-25, this classifies the eventual sign
in every fixed nonrectangular full-row direction with
$4\le r\le2m-1$. A positive permanent-minus-determinant gap
occurs eventually exactly when
\[
 \nu_1>(n-m)s,\qquad
 \Delta=rm^2-(r-2)n^2-2m+r-1>0.
\]
In every other case the determinant eventually dominates.
The strict-direction assertion already extends through $2m$
rows by G-25; the additional boundary assertion here stops
at $2m-1$ rows and excludes rectangular directions beyond
the range of G-26.
\end{theorem}

\paragraph{Explicit equality-face constants.}
Assume $\nu_1=as$. Put
\[
 q=r-1,\quad \mu=(\nu_2,\ldots,\nu_r),\quad
 \eta=\nu-\nu_r\mathbf1_r,\quad
 \theta=\mu-\nu_r\mathbf1_q,
 \quad e=|\eta|,\quad e'=|\theta|.
\]
As before, $r_\eta$ and $r_\theta$ count strict comparisons
between ordered parts. Define
\[
 \begin{aligned}
 u&=(r-2)n^2-r^2+3+r_\eta,\\
 v&=(r-1)m^2-2m-r^2+2r+2+r_\theta,\\
 b_D&=\binom{\binom{n+r-1}{r-1}+s-1}{s}
                    \dim S_\eta\mathbb C^r,\\
 b_B&=\binom{\binom{m+q-1}{q-1}+s-1}{s}
                    \dim S_\theta\mathbb C^q,\\
 b&=\max\{b_D,b_B\},\qquad
 L=ns+(r-1)e,\qquad p=ns+e,\\
 \beta&=\left\lceil\frac{3r^2b}{8}\right\rceil
                      L^{2(r^2-1)},\\
 T&=\binom{b+\beta}{b}-1,\qquad E=rT\beta!.
 \end{aligned}
\]
Then $u>v\ge2$, and strict determinant dominance holds for
every
\begin{equation}\label{G27:threshold}
                    j\ge E\bigl(4v(pE)^{v-1}+1\bigr)
\end{equation}
in coordinate degree $sj$. These are the same form of bounds
as G-26, but the lower invariant algebra is now allowed to
have a nontrivial degree period.

\begin{proof}
We compare the two enhanced invariant algebras and retain the
degree period only on the side used for an upper bound.
\paragraph{1. Separate stable algebras, with period one needed only above.}
On the equality face, the all-degree Pieri bound G-26 gives
\[
 m_{j\nu}(\ell^a\operatorname{per}_m)
       \le m_{j\mu}(\mathcal P_{m,q}).
\]
Let $A$ be the enhanced invariant algebra for
$\mathcal D_{n,r}$, outer degree $s$, and auxiliary weight
$\eta$, and let $B$ be that for $\mathcal P_{m,q}$ with
auxiliary weight $\theta$. Exact extraction yields
\[
 \dim A_j=m_{j\nu}(\det_n),\qquad
 \dim B_j=m_{j\mu}(\mathcal P_{m,q}).
\]
The two applications of R-32 are valid because
\[
 4\le r\le2m-1<2n-2,\qquad
 4\le q=r-1\le2m-2.
\]
Thus both enhanced cones have stable points with finite
projective stabilizers. Since $\nu$ is nonrectangular, the
first cone additionally has a stable point with no scalar
period. No such assertion is needed for $B$: the tail $\mu$
may be rectangular.

The exact image dimensions from R-28 and G-21 and the stable
quotient formula give $\dim A=u$ and $\dim B=v$.
The degree semigroup of $A$ has gcd one. The corresponding
gcd for $B$ is left unrestricted. This distinction is the
reason that no assertion about generic freeness of a
rectangular two-product tensor is required.

\paragraph{2. The dimension inequality persists beyond $m$ rows.}
Deleting the first part removes
\[
 r_\eta-r_\theta=\#\{j>1:\nu_1>\nu_j\}\ge0
\]
strict flag comparisons. Since $n\ge m+1$,
\[
 \begin{aligned}
 u-v
  &=(r-2)n^2-(r-1)m^2+2m-2r+1+r_\eta-r_\theta\\
  &\ge(2m-1)r-m^2-2m-1\\
  &\ge m^2-m-2>0.
 \end{aligned}
\]
For the last inequality use $r\ge m+1$. Moreover, writing
$q=r-1$, with $m\le q\le2m-2$, gives
\[
 v=q(m^2-q)-2m+3+r_\theta
        \ge m^3-2m^2+3>2.
\]
In particular $u\ge v+1$.

\paragraph{3. An upper Hilbert bound does not require degree period one.}
The degree and dimension estimates in G-23/G-26 give
homogeneous generator degrees at most $\beta$ in both rings,
at most $T$ generators, and standard $E$th Veroneses.
They also give $E\ge\beta^2$. The determinant no-scaling
point implies $A_j\ne0$ for every $j\ge\beta^2$.

For the lower algebra $B$, let $g$ be the gcd of the positive
degrees of a set of nonzero homogeneous generators. This
set is nonempty because $v\ge2$. Its degrees are at most
$\beta$. Divide them by $g$ and apply the elementary
residue-semigroup bound of G-18. It follows that every
multiple of $g$ at least $\beta^2$ belongs to their semigroup.
Since $B$ is a domain, a product of nonzero generators is
nonzero; hence $B_h\ne0$ at every such degree $h$.
Conversely $B_j\ne0$ implies $g\mid j$, by homogeneity of
the generators. This argument makes no saturation claim in
degrees not divisible by $g$.

Write $j=Ek+t$, $0\le t<E$. Since $g\le\beta$ and
$\beta!\mid E$, one has $g\mid E$. If $B_j\ne0$, then
$g\mid t$, and $2E-t\ge E\ge\beta^2$ is a multiple of
$g$. Multiplication by a nonzero element of $B_{2E-t}$
therefore embeds $B_j$ in $B_{E(k+2)}$. If $B_j=0$, the
same upper inequality holds without a multiplication map.
For $A$, multiplication by any nonzero element of
$A_{E+t}$ embeds $A_{E(k-1)}$ in $A_j$. Thus, for $k\ge1$,
\begin{equation}\label{G27:one_sided_sandwich}
 \dim A^{(E)}_{k-1}\le\dim A_j,\qquad
 \dim B_j\le\dim B^{(E)}_{k+2}.
\end{equation}
Only the lower bound in the first ring requires period one.

The lower permanent enhanced cone has a dense homogeneous
parameter map of degree $ms+e'\le p$. The image-degree
estimate used in G-23 therefore gives
\[
             \deg\operatorname{Proj}B^{(E)}\le(pE)^{v-1}.
\]
Linear Noether normalization for the standard domain
$A^{(E)}$ and the same Chardin bound for $B^{(E)}$ now give
\[
 \dim A_j\ge\binom{k+v-1}{v},\qquad
 \dim B_j\le(pE)^{v-1}\binom{k+v+1}{v-1}.
\]
As in G-26, their binomial ratio is at least $k/(4v)$ for
$k\ge v+1$, and \eqref{G27:threshold} makes the first
quantity strictly larger than the second. The Pieri upper
bound proves the desired strict comparison at every such
integer stretch, including all possible lower-ring residues.

\paragraph{4. The remaining two half-spaces.}
If $\nu_1<as$, G-26 gives zero padded multiplicity at every
positive stretch. R-32 still applies to the degree-$n$
determinant cone, because $r<2n-2$ and $\nu$ is nonrectangular.
It gives determinant occurrence at every
$j\ge\beta_D^2$, where
\[
 \beta_D=\left\lceil\frac{3r^2b_D}{8}\right\rceil
             \bigl(ns+(r-1)|\eta|\bigr)^{2(r^2-1)}.
\]
Here $b_D$ and $\eta$ have their displayed determinant
definitions; no lower tail degree is used in this case.
If $\nu_1>as$, G-25 supplies the comparison, with its
nonzero integer exponent $\Delta$. Combining these cases
with G-26 proves the stated eventual classification.
\end{proof}

\paragraph{Scope.}
The new semigroup argument allows a period on the algebra
used only for an upper bound; it does not remove that period.
Rectangular directions with more than $m$ rows, equality
faces at $2m$ rows, and strict directions beyond $2m$ rows
retain their separate stable-point gates. Earlier individual
degrees and varying-direction complexity questions are also
unsettled. The construction is independent of the stipulated
analytic inputs and makes no efficient-size claim for the
displayed thresholds.


\begin{theorem}[G-28: Boundary comparison from a generic stabilizer bound]
\label{G-28}
Let $n>m\ge17$ and
\[
                  m<r\le\min\{2n-2,m^2\}.
\]
Let $\nu\vdash ns$, $s\ge1$, have exactly $r$ positive
parts and be nonrectangular. Put $a=n-m$. If $\nu_1\le as$,
the determinant multiplicity is eventually strictly larger
than the padded permanent multiplicity at every integer
stretch. On the equality face, the explicit constants and
threshold of G-27 apply without change. Below the face,
the padded multiplicity is always zero and the determinant
occurs for every $j\ge\beta_D^2$, with $\beta_D$ as in G-27.

Consequently, for $n>m\ge17$ and $4\le r\le2m$, every
fixed nonrectangular full-row direction is classified
eventually: the padded permanent is larger exactly when
\[
 \nu_1>(n-m)s,\qquad
 rm^2-(r-2)n^2-2m+r-1>0,
\]
and the determinant is eventually larger otherwise.
The larger boundary range in the first paragraph does not
extend the strict-direction theorem past $2m$ rows.
\end{theorem}

\begin{proof}
We retain the Pieri tail reduction and replace the lower
stable quotient by the generic-orbit dimension upper bound
proved in R-33.

First suppose $\nu_1=as$. Define $q=r-1$, the tail $\mu$,
the weights $\eta,\theta$, and every constant $u,v,b_D,b_B,
b,L,p,\beta,T,E$ exactly as in G-27. In particular
\[
 \begin{aligned}
 u&=(r-2)n^2-r^2+3+r_\eta,\\
 v&=q(m^2-q)-2m+3+r_\theta.
 \end{aligned}
\]
Let $A$ and $B$ be the same determinant and lower permanent
enhanced invariant algebras. The Pieri bound still gives
\[
 m_{j\nu}(\ell^a\operatorname{per}_m)\le\dim B_j,
 \qquad \dim A_j=m_{j\nu}(\det_n).
\]

R-32 applies to $A$ because $r\le2n-2$, and the
nonrectangular label supplies a no-scaling stable point.
Thus $\dim A=u$ and every degree at least $\beta^2$ occurs
in $A$. For $B$, one has $4\le q<m^2$. R-33 gives only
\[
                         w:=\dim B\le v.
\]
No stable point, equality $w=v$, or period-one property is
assumed for this lower ring.

The same dimension difference computation as in G-27 uses
only $n\ge m+1$ and $r\ge m+1$, so
\[
 u-v\ge(2m-1)r-m^2-2m-1\ge m^2-m-2>0.
\]
Since $1\le q\le m^2-1$, the elementary identity
\[
 q(m^2-q)-(m^2-1)=(q-1)(m^2-q-1)\ge0
\]
also gives $v\ge m^2-2m+2>2$. Thus $u\ge v+1$,
regardless of whether the lower quotient attains its bound.

The same ambient representations and parameter maps give
the common generator bound $\beta$, the standard Veronese
degree $E$, and parameter degree at most $p$. If $w=0$,
the nonnegatively graded domain $B$ with $B_0=\mathbb C$
is just $\mathbb C$, so every positive padded multiplicity
is zero and the conclusion follows from occurrence in $A$.
Suppose henceforth $w\ge1$.

The arbitrary-period argument of G-27 applies to $B$ with
its actual nonzero generator degrees. For $j=Ek+t$,
$0\le t<E$, it gives
\[
 \dim A_j\ge\dim A^{(E)}_{k-1},\qquad
 \dim B_j\le\dim B^{(E)}_{k+2}.
\]
The projective image degree bound uses its actual dimension:
\[
 \deg\operatorname{Proj}B^{(E)}\le(pE)^{w-1}.
\]
For $w=1$ this is the degree-one reduced point case. The
standard-domain Hilbert upper bound, followed by $w\le v$,
gives
\[
 \begin{aligned}
 \dim B_j
   &\le(pE)^{w-1}\binom{k+w+1}{w-1}\\
   &\le(pE)^{v-1}\binom{k+v+1}{v-1}.
 \end{aligned}
\]
The lower Noether-normalization bound remains
$\dim A_j\ge\binom{k+v-1}{v}$. Therefore exactly the
binomial comparison in G-27 proves strict dominance for
\[
                    j\ge E\bigl(4v(pE)^{v-1}+1\bigr).
\]
No equality in a generic quotient-dimension formula has
entered this proof.

If $\nu_1<as$, G-26 makes the padded multiplicity zero.
The same R-32 application to $A$ yields determinant
occurrence at every $j\ge\beta_D^2$. This proves the
first assertion. Finally, $r\le2m$ implies both
$r\le2n-2$ and $r\le m^2$. G-28 therefore supplies
all nonrectangular on/below-face cases in that range;
G-26 handles the smaller rows and G-25 handles every
strict direction. Their nonzero integer exponent gives
the asserted complete nonrectangular classification.
\end{proof}

\paragraph{Scope.}
The additional input is a generic finite projective
stabilizer on the lower permanent image, proved from its
exact tangent kernel and Pl\"ucker coordinates in R-33.
It is used only for a quotient-dimension upper bound.
Rectangular directions beyond $m$ rows, strict directions
past $2m$ rows, individual earlier degrees and changing
directions retain their stated gates. This result does
not identify generic freeness with stability and does
not claim an efficient positive formula from the displayed
large thresholds.

% END SOURCE work/root/findings.tex

% BEGIN SOURCE work/isolation/findings.tex
% CUE_GCT isolation track. Run CUE-GCT-20261121-I.
% Model: GPT-6 Astra Pro; configured effort is not exposed to this agent.
% Bibliographic keys required: Hoza2026, AGKS2014, ForbesShpilka2013, MulmuleyGCTV, AllenderOgihara1996.

\section{Isolation: the exact transfer, the quotient, and the surviving gate}
\label{sec:isolation-track}

We use a field $F$, $n$ variables, individual degree at most $\delta$, and
an ROABP expression
\[
 f(\mathbf x)=u^{\mathsf T}A_1(x_{\pi(1)})\cdots A_n(x_{\pi(n)})v.
\]
The order $\pi$ may be unknown.  A polynomial-size hitting set here means
one set depending only on $(n,w,\delta,F)$, not on the matrices or their
coefficients, which hits every nonzero member of the class.  Arithmetic
complexity and bit complexity are distinguished throughout.

\paragraph{Background I-B01.}
The stipulated new graph theorem is Hoza (2026), Theorem 1.3, pp.~2--3;
its quantitative version is Theorem 2.4, p.~5.  For every simple layered
digraph of width $W$ and length $N$, some seed in a family of size
$(WN)^{O(1)}$ gives a minimum-weight path unique for every reachable pair.
The weights are nonnegative integers bounded by $B=(WN)^{O(1)}$ and are
computable in logarithmic space.  A path weight sums its internal vertex
weights.  With a binary code $C$ of length $\ell$ and distance $d$, the
formula is
\[
 \rho_x(a,i)=6^{\nu_2(i)}\sum_{j=1}^{\ell}6^{\ell-j}C(a)_j x_j
 \quad(i>0),\qquad \rho_x(a,0)=0,
\tag{I.1}\label{eq:I-hoza}
\]
and the probability of a bad uniform seed is at most
$2^{-d}W^4N^2$.  The graph is fixed before the seed is chosen.
We use this result as supplied; the present work does not claim an
independent verification of its complete theorem.

\paragraph{Background I-B02.}
Agrawal, Gurjar, Korwar, and Saxena (2014), \S3, Definition~5, call
$\omega\in\mathbb Z_{\geq0}^{n}$ basis isolating for
$D=\sum_a v_a\mathbf x^a$ if there is a coefficient basis indexed by $B$
whose monomial weights are distinct, and each $v_a$ with $a\notin B$
belongs to the span of basis coefficients of strictly smaller weight.
Their Lemma~6 preserves every nonzero scalar projection under
$x_i=t^{\omega_i}$; Lemma~8 and Theorem~1 give their quasipolynomial
construction, including unknown variable order.  Those known results
are cited, not claimed as findings here.

\subsection{I-01. A genuine polynomial transfer for vertex-variable path polynomials}
\label{I-01}

\paragraph{Statement.}
Let $G$ range over simple layered digraphs with fixed vertex slots,
width $W$, length $N\ge1$, fixed source $s$ and target $t$, and arbitrary
scalar edge labels $c_e\in F$.  Give every internal vertex slot its own
variable $z_v$, and form
\[
 P_G(\mathbf z)=\sum_{p:s\leadsto t}\left(\prod_{e\in p}c_e\right)
                        \prod_{v\in\operatorname{int}(p)}z_v.
\]
There is a uniform hitting set of cardinality at most
\[
 2^{\ell}\bigl((N-1)B+1\bigr)=(WN)^{O(1)}
\tag{I.2}
\]
over any field containing $(N-1)B+1$ distinct elements.  Its points have
coordinates $z_v=a^{\rho_x(v)}$, where $x$ runs over Hoza's seeds and
$a$ runs over a fixed set of $(N-1)B+1$ field elements.  The same conclusion
holds for independent edge variables after subdividing edges; the new
width is at most $W^2$ and the new length is $2N$.

\paragraph{Proof.}
Delete zero-labelled edges.  In a simple layered digraph two different
$s$--$t$ paths have different sequences of internal vertices and hence
different vertex monomials, except in the length-one case, where there
is at most one edge.  Thus $P_G\ne0$ exactly when some path survives.
Choose a min-isolating seed for the surviving graph using I-B01.  In
$P_G(t^{\rho_x(v)})$, the coefficient of the least path weight is the
product of labels on its unique minimum path, and is nonzero in the
field.  The degree is at most $(N-1)B$.  A nonzero polynomial of that
degree has at most $(N-1)B$ roots.  This proves (I.2), uniformly in the
edge labels.  For edge variables, insert one distinct intermediate
vertex on each edge and give it that edge's variable; give the original
internal vertices auxiliary independent variables as well.  The same
argument provides a hitting set for the expanded vertex-variable
polynomial.  To obtain points only on edge variables, assign to an edge
the product of its intermediate-vertex value and the value of its head
when the head is an internal original vertex.  Each original path then
receives exactly the product of all its expanded internal-vertex values.
This combines exponents additively, leaves a polynomial degree bound,
and proves the edge-variable statement. \hfill$\square$

\paragraph{Scope.}
This is an exact consequence of the new graph theorem.  It does not
identify several vertex or edge variables with the same ROABP variable.
That identification is the substantive next step, treated below.
Over a smaller finite field, the argument gives an extension-field
hitting set after constructing a sufficiently large extension; it is
not a statement that the same number of base-field points exists.

\subsection{I-02. The quotient obstruction is already present in a two-path diamond}
\label{I-02}

\paragraph{Statement.}
Let the monomial map $q:F[z_1,\ldots,z_s]\longrightarrow
F[x_1,\ldots,x_n]$ send $z_j$ to $\mathbf x^{a_j}$, and let $A$ be
the matrix with columns $a_j$.  A diagonal one-parameter action
$z_j\mapsto t^{r_j}z_j$ is induced by a diagonal action
$x_i\mapsto t^{\omega_i}x_i$ precisely when
$r=A^{\mathsf T}\omega$.  If the $a_j$ generate a proper sublattice,
requiring $r$ to annihilate $\ker(A:\mathbb Z^s\to\mathbb Z^n)$
only guarantees a character on the image lattice; an integral
extension to all $\mathbb Z^n$ is an additional lattice condition.
Consequently path isolation does not, in general, descend through $q$.

\paragraph{Proof and counterexample.}
Equivariance on each generator reads
$t^{r_j}\mathbf x^{a_j}=t^{\langle\omega,a_j\rangle}\mathbf x^{a_j}$,
which is equivalent to the displayed matrix equation.  The kernel
description follows from the factorization property of homomorphisms
of free abelian groups, with the stated extension qualification.

For the counterexample take the simple diamond
$s\to u\to t$ and $s\to v\to t$.  Label the two paths so that their
vertex-variable polynomial is $z_u-z_v$.  The map $q(z_u)=q(z_v)=x$
gives the zero polynomial.  Every min-isolating weight on this diamond
has $r_u\ne r_v$, so that its weighted polynomial
$t^{r_u}-t^{r_v}$ is nonzero.  Such a weight does not descend, because
$z_u-z_v\in\ker q$ is not homogeneous for it.  This remains a valid
counterexample in characteristic two, where the two distinct powers
still have nonzero coefficients.  Alternatively give both paths label
$1$ in characteristic zero: the original polynomial is $2x\ne0$,
but every weight induced from $x\mapsto t^{\omega}$ ties its two paths.
Thus compatible minimum-path isolation is not even necessary for
nonzero scalar output, while arbitrary path isolation is not sufficient
to certify nonzero output after the quotient. \hfill$\square$

\paragraph{What is refuted.}
The refuted assertion is the direct identification of a Hoza seed with
an ROABP variable-weight assignment.  Neither polynomial ROABP hitting
sets nor their possible derivation by a different construction is
refuted by this example.

\subsection{I-03. Three exact characterizations of basis isolation}
\label{I-03}

\paragraph{Statement.}
Let $S\subset\mathbb Z_{\geq0}^n$ index distinct monomials and let
$V=\operatorname{span}_F\{v_a:a\in S\}$ have dimension $r$ over an arbitrary field.
For $\omega\in\mathbb Z^n$ set $w_a=\langle\omega,a\rangle$.
For integer weights we use the same order-theoretic definition as
I-B02, allowing negative variable weights; the hitting-set
applications below use nonnegative weights.
The following are equivalent.
\begin{enumerate}
\item $\omega$ is a basis-isolating weight assignment in the sense of
this extended definition.
\item Every nonzero linear functional $\lambda\in V^*$ makes
$\sum_a\lambda(v_a)\mathbf x^a$ have a unique monomial of minimum weight.
\item At each weight level $c$, at most one coefficient $v_a$ with
$w_a=c$ survives modulo
$V_{<c}=\operatorname{span}\{v_b:w_b<c\}$.
\end{enumerate}
For a fixed coefficient basis $B$, write uniquely
$v_a=\sum_{b\in B}\lambda_{ab}v_b$.  It is the basis isolated by
$\omega$ exactly when
\[
 \langle\omega,a-b\rangle>0
 \quad(a\notin B,\ \lambda_{ab}\ne0),
 \qquad
 \langle\omega,b-b'\rangle\ne0\quad(b\ne b'\in B).
\tag{I.3}\label{eq:I-cone}
\]
The first inequalities alone characterize $B$ as the unique
minimum-total-weight coefficient basis.

\paragraph{Proof.}
Assume (1), with isolated basis $B$.  For a nonzero $\lambda$, take
the least-weight basis member $b$ on which $\lambda$ is nonzero.
It is unique.  Every nonbasis coefficient of no greater weight lies
in the span of smaller basis coefficients and is annihilated by
$\lambda$.  This proves (2), including its stronger unique-monomial
assertion.

If (3) fails, choose two nonzero coefficient classes at a level $c$.
There is a functional on $V/V_{<c}$ nonzero on both: for independent
classes assign value $1$ to both and extend a basis; for dependent
nonzero classes use any functional nonzero on one of them.  Pull it
back to $V$.  It vanishes below $c$ and is nonzero on two monomials
at $c$, contradicting (2).  This argument works also over finite
fields.  Thus (2) implies (3).

Under (3), choose the unique surviving monomial at every level where
the coefficient span grows.  These coefficients form a basis, one
per growth level.  All remaining coefficients at a level already
belong to the span of smaller levels.  Iterating through the finitely
many levels proves (1).  Unique coordinates in a fixed basis prove
the equivalence with (I.3).

If $B$ is the unique minimum-total-weight basis and
$\lambda_{ab}\ne0$, then $B-b+a$ is also a basis.  Its total weight
is strictly larger, so $w_a>w_b$.  Conversely, suppose these strict
inequalities hold.  Expand any other basis $C$ in the basis $B$.
Its change-of-basis determinant is nonzero, so at least one
permutation in its determinant expansion is nonzero.  Every common
member of $B\cap C$ is forced to match itself; every member
$a\in C\setminus B$ matches some $b\in B\setminus C$ with
$\lambda_{ab}\ne0$.  Summing the strict inequalities over this
matching gives $\sum_{a\in C}w_a>\sum_{b\in B}w_b$.
Thus $B$ is the unique minimum-total-weight basis. \hfill$\square$

\paragraph{Geometric consequence.}
For a fixed ordering of the weights within $B$, (I.3) is an open
rational polyhedral cone in the cocharacter space $\mathbb R^n$.
The union of these cones is precisely the set of BIWAs.  Equivalently,
the one-parameter subgroup simultaneously gives a monomial initial
form for every nonzero scalar projection of $D$.  Isolating one path
or one preselected scalar output has weaker quantifiers.

\subsection{I-04. Projected matroid vertices and polynomial random weights}
\label{I-04}

\paragraph{Statement.}
Let $\mathcal B$ be the coefficient bases of the data in I-03, and put
\[
 \beta(B)=\sum_{a\in B}a,\qquad
 Q=\operatorname{conv}\{\beta(B):B\in\mathcal B\}\subset\mathbb R^n.
\]
Every vertex of $Q$ is the image of exactly one indexed basis.
Assume $S\subseteq\{0,\ldots,\delta\}^{n}$, $r\ge1$, $\delta\ge1$.
Choose independently
$a_i\in\{1,\ldots,M_1\}$ and $b_i\in\{1,\ldots,M_2\}$.
With $L=n\delta M_2+1$, the weights
\[
 \omega_i=L a_i+b_i
\tag{I.4}
\]
are basis isolating with probability at least
\[
 1-\frac{nr\delta}{M_1}-\frac{\binom r2}{M_2}.
\tag{I.5}
\]
For example $M_1=4nr\delta$ and $M_2=4\max(1,r^2)$ give a positive
constant success probability and polynomially bounded weights.  This
is a randomized existence statement, not an explicit small family.

\paragraph{Proof of the vertex assertion.}
Suppose two distinct bases $B,C$ have the same image $\beta$.
Take $e\in B\setminus C$.  Let $\ell$ be the coordinate functional
of $v_e$ in the basis $B$, and expand $v_e=\sum_{f\in C}c_fv_f$.
Because $1=\ell(v_e)=\sum c_f\ell(v_f)$, some
$f\in C\setminus B$ has both $c_f\ne0$ and $\ell(v_f)\ne0$.
Hence $B-e+f$ and $C-f+e$ are both bases.  Their images are
$\beta+(f-e)$ and $\beta-(f-e)$, which are distinct because the
ground-set exponents are distinct.  Their midpoint is $\beta$,
so $\beta$ is not a vertex.  Every vertex has at least one preimage
from the finite defining set, proving the claim.

\paragraph{Proof of the probability bound.}
First consider any finite set $T\subseteq\{0,\ldots,D\}^{n}$ and
uniform independent weights $a_i\in[M_1]$.  Fix all weights except
$a_i$.  The least weight as a function of a real parameter $t=a_i$
is the lower envelope of at most $D+1$ affine functions, one for
each possible $i$-th coordinate.  Its slopes are among
$0,\ldots,D$, and there are at most $D$ breakpoints at which
different slopes minimize.  If two different points of $T$ minimize,
they differ in some coordinate $i$, and that coordinate weight is
one of these breakpoints.  Conditioning on the other coordinates
and applying a union bound gives probability at most $nD/M_1$.
Apply this with $T=\{\beta(B)\}$ and $D=r\delta$.
With probability at least $1-nr\delta/M_1$, there is a unique
minimizing point of $Q$.  It is a vertex and, by the preceding
argument, has a unique basis preimage $B$.

Conditional on the first weights and this $B$, two different
monomials in $B$ have equal $b$-weight with probability at most
$1/M_2$: condition on every $b_i$ but one whose exponent difference
is nonzero.  There is at most one possible value of that remaining
integer weight.  A union bound over pairs costs
$\binom r2/M_2$.  By I-03, the unique first-weight minimum basis
satisfies $a\cdot m>a\cdot b$ whenever the $b$-coordinate of
$v_m$ is nonzero and $m\notin B$.  Each difference is an integer
at least one.  The perturbation changes it by at most
$n\delta M_2$, so multiplication by $L=n\delta M_2+1$ preserves
strict positivity.  It also preserves inequalities between basis
members with unequal first weights, and separates any remaining
ties by the second weights.  Equation (I.3) proves (I.5).
All randomness and weights in this proof are integers; no assumption
on the characteristic of the coefficient field enters the argument.
\hfill$\square$

\paragraph{The remaining derandomization problem.}
The desired generator must hit these cones for coefficient matroids
arising from products of univariate matrices, with a family and
maximum weights both polynomial in $(n,w,\delta)$.  Hoza's theorem
hits path-polytope cones in independent vertex coordinates.  I-02
shows why those cocharacter spaces cannot be silently identified.
I-04 narrows the required object to normal cones of a specified
projected representable-matroid polytope, with the finite
within-basis tie hyperplanes removed.

\subsection{I-05. Positive Gram determinants certify bases, not compressed rank}
\label{I-05}

\paragraph{Statement.}
For complex coefficient vectors, choose coordinates $V=\mathbb C^r$
and form
\[
 G_D(\mathbf z)=\det\left(\sum_{a\in S}v_av_a^*\mathbf z^a\right).
\]
Then
\[
 G_D(\mathbf z)=\sum_{B\in\mathcal B}
                  |\det C_B|^2\mathbf z^{\beta(B)},
\qquad \operatorname{Newt}(G_D)=Q,
\tag{I.6}
\]
where $C$ has columns $v_a$.  At a vertex of $Q$, the coefficient
is the squared volume of its unique basis from I-04.  Nevertheless,
positivity of $G_D(t^\omega)$ does not imply that the coefficient
span of $D(t^\omega)$ retains its dimension.

\paragraph{Proof.}
Apply the Cauchy--Binet identity to
$C\operatorname{diag}(\mathbf z^a)C^*$.
Every nonzero summand is strictly positive in its coefficient, so
there are no cancellations between bases with the same exponent
sum; their convex hull is exactly the Newton polytope.  I-04
identifies its vertex coefficients.  For the negative assertion,
take $D(x,y)=e_1x+e_2y$ and $\omega=(1,1)$.  Its Gram determinant
is $xy$ and specializes to $t^2$, whereas
$D(t,t)=t(e_1+e_2)$ has coefficient span of dimension one instead
of two.  The unique coefficient basis has two equal member weights;
this is exactly the extra condition in I-03. \hfill$\square$

\paragraph{Field and geometric qualification.}
The conjugate transpose and positivity in (I.6) are statements on
a complex vector space with a Hermitian structure.  They are not
positivity statements in an arbitrary characteristic.  The
determinantal/Cauchy--Binet identity itself has an integral
two-matrix form with independent dual variables.  Specializing the
second matrix to a conjugate transpose is a real, not a complex
algebraic, operation.  This qualification is essential in translating
the formula into a scheme over $\mathbb Z$.

\subsection{I-06. Exact black-box implementation of nonlinear exponent weights can be exponential}
\label{I-06}

\paragraph{Statement.}
Let $K$ be any field, let $t$ be transcendental over it, and define
a linear functional on polynomials of individual degree at most $2d$
by
\[
 \mathcal L(\mathbf x^\alpha)=\prod_{i=1}^n h_i(\alpha_i),
 \qquad h_i(e)\in K(t).
\]
Put $H_i=(h_i(a+b))_{0\le a,b\le d}$.
If one exact universal implementation is
\[
 \mathcal L(f)=\sum_{j=1}^M c_j f(\xi_j),
\tag{I.7}
\]
where the points and coefficients are fixed independently of $f$
over any extension of $K(t)$, then
\[
 M\ \ge\ \prod_{i=1}^n\operatorname{rank}(H_i).
\tag{I.8}
\]
In particular $\mathcal L(\mathbf x^\alpha)=
t^{\sum_i\alpha_i^2}$ requires $M\ge(d+1)^n$.
For degree two, the nonlinear exponent filters supplied by every
Hoza seed that separates three exponent symbols also have the
lower bound $M\ge2^n$.

\paragraph{Proof.}
Index a moment matrix by
$\alpha,\beta\in\{0,\ldots,d\}^n$ and set
$H_{\alpha\beta}=\mathcal L(\mathbf x^{\alpha+\beta})$.
It is the tensor product $\bigotimes_iH_i$, with rank the product
of the ranks.  Each evaluation in (I.7) contributes the rank-one
matrix $c_j(\xi_j^\alpha)_\alpha(\xi_j^\beta)_\beta^{\mathsf T}$,
so its rank is at most $M$.  Extension of scalars does not change
the rank, proving (I.8).

For $h_i(e)=t^{e^2}$, factor the $(a,b)$ entry as
$t^{a^2}t^{b^2}(t^{2a})^b$.  The remaining matrix is Vandermonde
on the distinct elements $1,t^2,\ldots,t^{2d}$ and has full rank,
in every characteristic.

For the Hoza assertion take three distinct binary codewords
$C(0),C(1),C(2)$ and put
$h_x(e)=\sum_j6^{\ell-j}C(e)_j x_j$ as integers.
If $h_x(0)+h_x(2)=2h_x(1)$, then
\[
 \sum_j6^{\ell-j}x_j\bigl(C(0)_j+C(2)_j-2C(1)_j\bigr)=0.
\]
The digits lie between $-2$ and $2$.  The first nonzero digit
dominates the absolute sum of all later digits, since
$2\sum_{k=0}^{q-1}6^k<6^q$.  Thus every digit vanishes, which for
binary digits says the three selected code bits agree in every
selected coordinate.  Hence all three values $h_x(e)$ coincide.
Conversely constant values give that equality.  Any seed isolating
the three-arm graph is nonconstant on its three internal vertices,
so the equality fails.  At layer $i$, the scalar multiplier
$6^{\nu_2(i)}$ in (I.1) preserves this failure.  Consequently the
$2\times2$ Hankel determinant of the exponent filter
\[
 e\longmapsto t^{6^{\nu_2(i)}h_x(e)}
\]
is the difference of two distinct monomials, and has nonzero
determinant over $K(t)$.  Equation (I.8) now gives $2^n$.
\hfill$\square$

\paragraph{Scope of the lower bound.}
This rules out the proposed exact linear interpolation implementation
of independently reweighted monomial letters by a polynomial number
of evaluations of the original black box.  It is not a lower bound
on general ROABP PIT, adaptive reconstruction, or approximate
integration.  Even requiring (I.7) only for width-one ROABPs implies
the same identities on every monomial and hence the same lower
bound by linearity; width-one PIT itself remains elementary.

\subsection{I-07. Haar averaging repairs invariance by discarding state isolation}
\label{I-07}

\paragraph{Statement.}
For a fixed complex $m\times m$ matrix $D$, the function
$A\mapsto\operatorname{tr}(DA)$ is invariant under simultaneous
conjugation if and only if $D$ is scalar.  Moreover
\[
 \int_{U(m)}UDU^*\,dU=\frac{\operatorname{tr}D}{m}I.
\tag{I.9}
\]
Thus independently averaging fixed diagonal state-weight insertions
in a trace word replaces each insertion by a scalar and recovers
a scalar multiple of the original trace word.  If the same unitary
is used for two insertions, the exact formula for $m\ge2$ is
\[
 \int\operatorname{tr}(A UDU^* B UEU^*)\,dU
 =\alpha\operatorname{tr}(AB)+\beta\operatorname{tr}A\operatorname{tr}B,
\tag{I.10}
\]
where
\[
 \alpha=\frac{m\operatorname{tr}D\operatorname{tr}E-
                    \operatorname{tr}(DE)}{m(m^2-1)},\qquad
 \beta=\frac{m\operatorname{tr}(DE)-
                    \operatorname{tr}D\operatorname{tr}E}{m(m^2-1)}.
\]
Even positive semidefinite $D,E$ need not make these coefficients
nonnegative: for orthogonal rank-one projections,
$\beta=-1/[m(m^2-1)]$.

\paragraph{Proof.}
Conjugation invariance for every $A$ gives
$\operatorname{tr}((g^{-1}Dg-D)A)=0$ for all $A$.  The trace pairing
is nondegenerate, so $D$ commutes with every invertible matrix and
is scalar.  The reverse direction is immediate.  Haar invariance
makes the left side of (I.9) commute with every unitary; the
commutant is scalar, and its trace determines the scalar.  Repeated
integration gives the independent-insertion claim.

For completeness, in (I.10) define
$T=\int(UDU^*)\otimes(UEU^*)\,dU$.  The commutant of
$U\otimes U$ is spanned by the identity $I$ and tensor flip $S$.
Here is an elementary verification: commutation with diagonal unitaries
forces matrix entries to preserve the unordered pair of tensor indices;
permutation unitaries make the two-dimensional off-diagonal blocks
$\left(\begin{smallmatrix}b&c\\c&b\end{smallmatrix}\right)$ and
the diagonal blocks a common scalar $a$.  A rotation mixing two basis
vectors forces $a=b+c$.  The resulting operator is $bI+cS$.
Thus $T=\alpha I+\beta S$.  Its trace and its trace after multiplication
by $S$ are respectively
\[
 m^2\alpha+m\beta=\operatorname{tr}D\operatorname{tr}E,
 \qquad m\alpha+m^2\beta=\operatorname{tr}(DE).
\]
Solving gives the stated values.  Finally
$\operatorname{tr}(A X B Y)=
\operatorname{tr}((A\otimes B)(X\otimes Y)S)$;
contracting $T$ gives (I.10).  For orthogonal rank-one projections,
$\operatorname{tr}D=\operatorname{tr}E=1$ and
$\operatorname{tr}(DE)=0$, proving the sign assertion.
\hfill$\square$

\paragraph{Meaning for the proposed transfer.}
Hoza's internal-state weights enter matrix products as fixed
diagonal insertions.  Their lack of conjugation invariance cannot
be repaired by independent Haar averaging while preserving the
state information.  Correlated averaging changes the trace word
into a combination of different trace invariants and can reintroduce
negative coefficients.  Neither averaging is the missing
basis-isolation construction.

\subsection{I-08. The exact matrix-invariant consequence, conditional on the remaining gate}
\label{I-08}

\paragraph{Statement.}
Suppose one has an explicit hitting set $\mathcal H\subset F^{m^2}$
of polynomial size, constructible in polynomial arithmetic time,
for ROABPs of depth $m^2$, width at most $2m^2$, and individual
degree less than $r$.  Let $F$ have characteristic zero and put
\[
 M(x)=\sum_{j=0}^{r-1}M_jx^j,\qquad
 T_{\ell,\alpha}(\mathbf M)=
 \operatorname{tr}\bigl(M(\alpha_1)\cdots M(\alpha_\ell)\bigr),
 \quad1\le\ell\le m^2,quad\alpha\in\mathcal H.
\tag{I.11}
\]
These at most $m^2|\mathcal H|$ invariants are an explicit
polynomial-size separating family for simultaneous conjugation.
Over algebraically closed $F$, the invariant ring is finite over
the subring they generate.  This is precisely the specialization
of Forbes and Shpilka (2013), Construction~3.7, Theorem~3.8 and
Corollary~3.10, to a polynomial hitting set.  It supplies the
matrix-invariant instance of explicit Noether normalization in
the sense of that paper.  It does not by itself extract an
algebraically independent system of parameters of cardinality
equal to the quotient dimension.

\paragraph{Derivation and exact source dependence.}
The coefficient of $x_1^{a_1}\cdots x_\ell^{a_\ell}$ in
$\operatorname{tr}(M(x_1)\cdots M(x_\ell))$ is the trace word
$\operatorname{tr}(M_{a_1}\cdots M_{a_\ell})$.  The latter words,
for $\ell\le m^2$, generate the characteristic-zero invariant ring
as recorded in Forbes--Shpilka Theorem~1.3 and Corollary~3.6.
For two tuples $A,B$, the difference of their trace-generating
polynomials is an ROABP of width at most $2m^2$, by their Lemma~3.2.
The hitting-set assumption therefore detects exactly whether
some trace generator separates $A$ and $B$.  Their Theorem~3.9
(the stated specialization of Derksen--Kemper Theorem~2.3.12)
then gives integrality, and finite generation upgrades this to
a finite morphism.  The only new hypothesis in this specialization
is the polynomial hitting-set gate; it has not been proved here.
\hfill$\square$

\paragraph{Consequences that do not follow.}
This does not prove general arithmetic PIT, arbitrary explicit
Noether normalization, permanent circuit lower bounds, or
$P\ne NP$.  Forbes--Shpilka already prove deterministic NC
orbit-closure intersection in Theorem~1.15 using white-box PIT.
Their discussion after Theorems~1.13--1.15 explicitly separates
the matrix-invariant case from the stronger general explicit-variety
problem.  In the supplied revised GCT~V source, the general-PIT
equivalence is Theorem~5.14, and the matrix-specific ROABP replacement
is \S7.5, Lemma~7.9.  The earlier GCT~V theorem numbering cited by
Forbes--Shpilka must not be transferred to the revised source.  No new polynomial hitting set
may be inserted into (I.11) until its own gate is closed.

\subsection{I-09. An exact positive transfer-matrix certificate remains white box}
\label{I-09}

\paragraph{Statement.}
For a complex ROABP in known matrix-product form
\[
 f(\mathbf x)=u^{\mathsf T}\prod_{i=1}^n
                  \left(\sum_{a=0}^{\delta}A_{i,a}x_i^a\right)v
             =\sum_{\alpha}c_{\alpha}\mathbf x^{\alpha},
\]
its coefficient squared-norm generating polynomial is
\[
 \sum_{\alpha}|c_\alpha|^2\mathbf z^\alpha
 = (u\otimes\bar u)^{\mathsf T}
     \prod_{i=1}^n\left(\sum_{a=0}^{\delta}
                    A_{i,a}\otimes\overline{A_{i,a}}\,z_i^a\right)
                     (v\otimes\bar v).
\tag{I.12}
\]
In particular $f=0$ if and only if the right side at
$\mathbf z=\mathbf1$ is zero.  Its matrices have width $w^2$.
The equality also gives the usual torus identity
\[
 \int_{[0,2\pi]^n}|f(e^{i\theta_1},\ldots,e^{i\theta_n})|^2
                \prod_i\frac{d\theta_i}{2\pi}
       =\sum_\alpha|c_\alpha|^2.
\tag{I.13}
\]

\paragraph{Proof.}
The coefficient $c_\alpha$ is
$u^{\mathsf T}A_{1,\alpha_1}\cdots A_{n,\alpha_n}v$.
Multiplicativity of tensor products identifies its squared
absolute value with the product of the corresponding tensor
matrices and boundary vectors.  Summing independently over each
$\alpha_i$ gives (I.12).  The sum is a finite sum of nonnegative
real numbers, and it is zero exactly when every coefficient is
zero.  In (I.13), expansion of the square and integration of
$e^{i\langle\alpha-\beta,\theta\rangle}$ kills exactly the terms
$\alpha\ne\beta$. \hfill$\square$

\paragraph{Why this does not close the black-box gate.}
The transfer matrices use the individual coefficient matrices
$A_{i,a}$, which a black box does not supply.  The torus integral
is another exact expression, but the product quadrature that is
exact for all such degree bounds uses $(\delta+1)^n$ nodes.
No lower bound on all specialized ROABP quadratures is asserted.
I-06 proves a lower bound for the different, explicitly specified
nonlinear exponent filter.  Positivity in (I.12) closes the
white-box zero-testing question only; it does not supply the
missing uniform polynomial set of original evaluations.

\subsection{I-10. Schemes, cocharacters, flags, and flat limits}
\label{I-10}

\paragraph{Statement and dictionary.}
Let $R$ be a domain carrying the relevant coefficients, and let
$T=\mathbb G_{m,R}^{n}$.  The integer weight vector $\omega$
is the cocharacter
$\lambda_\omega:\mathbb G_{m,R}\to T$,
$t\mapsto(t^{\omega_1},\ldots,t^{\omega_n})$.
For nonzero $f=\sum_a c_a\mathbf x^a$ and
$\mu=\min_{c_a\ne0}\langle\omega,a\rangle$, define
\[
 \widetilde f(t,\mathbf x)=
 \sum_a c_a t^{\langle\omega,a\rangle-\mu}\mathbf x^a.
\tag{I.14}
\]
This is a polynomial over $R[t]$, its special fiber is the
weight-initial form, and its generic fiber is the original
polynomial under a torus change of coordinates.  A unique
minimizing exponent gives a monomial special fiber.  Over a field,
the hypersurface family defined by (I.14) is flat over the
$t$-line.  For a homogeneous ideal, the corresponding weight
homogenization is the usual flat Rees/initial-ideal degeneration;
this is stronger than merely selecting a least monomial of one
polynomial and must not be conflated with it.

For a vector polynomial, I-03 identifies a complete coefficient
flag whose successive quotients have one supported weight
monomial.  With an explicit coefficient matrix, nonvanishing
maximal minors define a matroid stratum of the appropriate
Grassmannian.  The exponents of its Pl\"ucker coordinates under
$\lambda_\omega$ are the sums $\beta(B)$.  Its torus orbit has
weight polytope $Q$; I-04 says a vertex with nonzero Pl\"ucker
coordinate corresponds to a unique indexed coefficient basis.
The BIWA cones refine its normal fan by the individual
within-basis weight comparisons in (I.3).

\paragraph{Proof of the flatness assertion and the dictionary.}
Formula (I.14) has no negative power of $t$ and is not divisible
by $t$.  Over a field $F$, for every irreducible polynomial
$p(t)$ the image of $\widetilde f$ modulo $p(t)$ is nonzero:
for $p=t$ use its nonzero initial form; otherwise $t$ is invertible
in $F[t]/(p)$ and its monomial coefficients remain nonzero.
Thus no nonzero $p(t)$ divides $\widetilde f$ in the UFD
$F[t,\mathbf x]$.  If $p(t)g\in(\widetilde f)$, primality of
the irreducible factors of $\widetilde f$ implies
$g\in(\widetilde f)$.  Hence the quotient by
$(\widetilde f)$ is torsion-free over the PID $F[t]$, and
therefore flat.  The generic-fiber statement follows by
inverting $t$ in (I.14).  Exterior products of coefficient
columns are precisely Pl\"ucker coordinates, and their weights
add, giving $\beta(B)$.  Equation (I.3) then identifies the
normal-cone refinement.  This proves each part of the stated
dictionary directly. \hfill$\square$

\paragraph{Integral qualification.}
The torus and monomial maps are defined over $\mathbb Z$.
A specified coefficient matroid and its rank flag persist on an
open localization where the finitely many designated nonzero
minors remain nonzero; they need not persist in every residue
characteristic.  Over a general base $R$, flatness of a
hypersurface quotient also needs the relevant coefficient and
fiber hypotheses.  The field proof above is not a claim of
automatic flatness over arbitrary $R$.  The Grothendieck locator
table in the main paper should be used for the exact SGA/FGA/EGA
placement of these notions, rather than inventing a theorem
number for these elementary derived statements.

\subsection{I-11. Exterior and subspace-state repairs do not yet produce a graph generator}
\label{I-11}

\paragraph{Checked finite assertions.}
For coefficient ambient dimension $k$, the standard coordinate
space $\bigwedge^rF^k$ has dimension $\binom kr$.  Over
$\mathbb F_q$, there are at least $q^{r(k-r)}$ distinct
$r$-dimensional subspaces of $F_q^k$: the graphs of the
$q^{r(k-r)}$ linear maps $F_q^r\to F_q^{k-r}$ are all distinct.
Thus a literal graph with one state for each such subspace is
not polynomial in $k$ when $r$ grows with $k$.

Exterior coordinates also retain cancellation.  For two equal
vectors $v_1=v_2=e_1+e_2$, the determinant is
$1\cdot1-1\cdot1=0$, although its expansion has two nonzero
permutation terms.  Making one term uniquely minimum by
path-dependent weights changes the determinant rather than
certifying its nonvanishing.  Tensoring $r$ independently chosen
word trajectories gives a standard realization with $w^{O(r)}$
states; for fixed $r$ this removes that particular size cost,
but the cancellation and weight-descent problems remain.

\paragraph{Proof and scope.}
The exterior dimension follows from the basis wedges indexed
by $r$-element subsets of a fixed $k$-element basis.  The graph
of a linear map is an $r$-dimensional subspace, and projection
onto the first summand recovers the map, proving distinctness.
The determinant example is immediate.  The tensor product
construction records one of $w$ state positions for each of
the $r$ trajectories, so has $w^r$ positions before additional
boundary coordinates; determinant signs are imposed at the
boundary and do not disappear.  These are size and correctness
facts about the specified constructions, not lower bounds on
all possible reductions.  A scalar output can itself have
coefficient-space dimension one, so constant coefficient rank
alone cannot make the general ROABP hitting-set task disappear.
\hfill$\square$

\subsection{I-12. Resolving coefficient support is an arithmetic counting gate}
\label{I-12}

\paragraph{Background I-B05.}
Allender and Ogihara (1996), \S2, p.~4, define the accepting-minus-rejecting
gap model; \S5, Definition~2, p.~11, defines $C_{=}L$ by the vanishing
of a $\operatorname{GapL}$ function.  Their Theorem~1, p.~2, identifies
$\operatorname{GapL}$ with logarithmic-space reductions to integer
determinant, and Theorem~14, p.~12, gives the corresponding singularity
completeness result.  The following ROABP formulation is proved directly.

\paragraph{Statement.}
Restrict to explicitly given layered ROABPs over $\mathbb Z$
(equivalently over $\mathbb Q$) with coefficients in
$\{0,1,-1\}$ and polynomially bounded width and length.  The input
lists the coefficient entries and the queried exponent vector;
univariate edge labels may be given densely or as canonical sparse
lists with binary exponents.  Deciding
whether a specified monomial coefficient is nonzero is complete,
under deterministic logarithmic-space reductions, for the class
\[
 \operatorname{coC}_{=}L
   =\{\{x:g(x)\ne0\}:g\in\operatorname{GapL}\}.
\]
Here $\operatorname{GapL}$ is defined by the difference of the
numbers of accepting and rejecting paths of polynomial-time
nondeterministic logarithmic-space machines.  Hardness already
holds for ROABPs whose polynomial is one monomial times a scalar.
Therefore a logarithmic-space uniform polynomial-size reduction
which represents exactly nonzero coefficient support by ordinary
directed reachability would imply
$\operatorname{coC}_{=}L\subseteq NL$.

\paragraph{Proof.}
For membership, guess a state path in the layered ROABP.  At layer
$i$ choose the term of the prescribed exponent from its univariate
edge label.  A zero term contributes nothing; otherwise track the
parity of the negative signs.  The difference between positive and
negative paths is exactly the specified monomial coefficient.
State indices, a layer counter, and a sign bit require logarithmic
space.  Neutral or invalid branches can be arranged to contribute
one accepting and one rejecting path, hence zero gap.

For hardness, start with a polynomial-time nondeterministic
logarithmic-space machine computing a gap $g(x)$.  Its time-unrolled
configuration graph has polynomial width and length and is
constructible from $x$ in logarithmic space.  Pad halted computations
by deterministic transitions, retaining their accepting or rejecting
tag, so every path has the same length.  Mark the final edge
of accepting paths by $1$ and of rejecting paths by $-1$;
all preceding edges have label $1$.  Separate parallel transitions
by an intermediate layer if necessary.  Multiply every edge label
in layer $i$ by a fresh variable $x_i$.  Every source-to-target path
then has the same variable monomial $\prod_i x_i$, with its
accepting or rejecting sign, so the output is
$g(x)\prod_i x_i$.  This is a known-order multilinear ROABP with
coefficients $0,1,-1$.  Its unique possible monomial coefficient
is nonzero exactly when $g(x)\ne0$.  Finally, composition with the
postulated graph reduction puts every such language in $NL$.
\hfill$\square$

\paragraph{Relation to the other gates.}
This is a precise complexity formulation of the operation which
would have to cancel equal-monomial paths before graph isolation.
It is not a separation of complexity classes.  Nor does it rule
out a different polynomial-time algebraic construction: the
explicit ROABP supplies its matrices, and white-box polynomial
identity testing is already possible.  In scheme language the
coefficient is a regular function of the edge coefficients;
testing its vanishing uses the kernel of a coefficient map, not
merely whether its parameterizing path set is empty.  The diamond
in I-02 is the smallest instance of this distinction.

\subsection{I-13. Exact universal energy cubature needs the full coefficient dimension}
\label{I-13}

\paragraph{Statement.}
Let $n\ge1$, $\delta\ge0$, and $D=(\delta+1)^n$.  Fix points
$\xi_1,\ldots,\xi_Q\in\mathbb C^n$ and a complex $Q\times Q$ matrix $C$,
independent of the polynomial.  Suppose that every width-one product
\[
 f(\mathbf x)=\prod_{i=1}^n\left(\sum_{a=0}^{\delta}z_{i,a}x_i^a\right),
 \qquad z_{i,a}\in\mathbb C,
\]
satisfies the exact identity
\[
 \sum_{\alpha\in\{0,\ldots,\delta\}^n}|[\mathbf x^\alpha]f|^2
  =\sum_{j,k=1}^Q C_{jk}\,\overline{f(\xi_j)}f(\xi_k).
\tag{I.13}\label{eq:I-cubature}
\]
Then $Q\ge D$.  In particular, signed cubature
$\sum_j c_j|f(\xi_j)|^2$ requires at least $D$ nodes.  More generally,
a formula with $M$ fixed paired terms
$\sum_{j=1}^M c_j\overline{f(\xi_j)}f(\eta_j)$ requires $M\ge D$.
The bound is attained for the diagonal formula by the Cartesian product
of the $(\delta+1)$st roots of unity in each coordinate, all with weight
$1/D$.

\paragraph{Proof.}
Index coordinates by $\alpha\in\{0,\ldots,\delta\}^n$ and put
$a_\alpha=\prod_i z_{i,\alpha_i}$.  Let the evaluation matrix be
$E_{j\alpha}=\xi_j^\alpha$.  The asserted identity is
\[
 a^*(I_D-E^*CE)a=0
\quad\text{for every decomposable coefficient vector }
a=\bigotimes_i(z_{i,0},\ldots,z_{i,\delta}).
\]
Expanding gives a polynomial identity in the variables
$z_{i,a}$ and their conjugates.  Its monomials
\[
 \prod_i\overline{z_{i,\alpha_i}}z_{i,\beta_i},
 \qquad (\alpha,\beta)\in\{0,\ldots,\delta\}^{2n},
\]
are linearly independent as functions of the complex variables.
Indeed, writing each $z=u+iv$ identifies the polynomial ring in
formal $z,\bar z$ with the polynomial ring in $u,v$ by an invertible
complex linear change of variables; a polynomial vanishing on every
real tuple $(u,v)$ has every coefficient zero.  Distinct pairs
$(\alpha,\beta)$ give distinct displayed monomials.  Consequently
$I_D=E^*CE$, so
$D=\operatorname{rank}(I_D)\le\operatorname{rank}(E)\le Q$.
For paired terms the same coefficient comparison represents $I_D$
as a sum of $M$ rank-at-most-one matrices, proving $M\ge D$.

Finally, averaging $\xi^\alpha\overline{\xi^\beta}$ over the stated
root-of-unity grid gives $1$ when $\alpha=\beta$ and $0$ otherwise:
in each coordinate the difference lies between $-\delta$ and
$\delta$, whose only multiple of $\delta+1$ is zero.  Expanding
$|f|^2$ therefore proves the matching formula, for every polynomial
of the stated individual degree, including the width-one class.
\hfill$\square$

\paragraph{Scope and consequence for I-09.}
Torus orthogonality identifies the left side of (I.13) with
$\int_{(S^1)^n}|f(z)|^2\,dz$.  Thus the positive white-box energy
computation in I-09 cannot be replaced by a polynomial number of
fixed black-box samples through any exact quadratic formula of the
displayed kind, even for width one and fixed $\delta\ge1$ as $n$ grows.
This is a lower bound for universal exact cubature, not for hitting
sets: for width one the substitution $x_i=t$ gives the nonzero
univariate product $\prod_i f_i(t)$ of degree at most $n\delta$.
Evaluation at $n\delta+1$ distinct scalars already gives a polynomial
hitting set for that class.  The rank argument makes no claim about
adaptive procedures or other nonlinear uses of evaluations.

\subsection{I-14. Joint compression is a Grassmannian transversality condition}
\label{I-14}

\paragraph{Statement.}
Write $D=\sum_{a\in S}v_a\mathbf x^a$, choose coordinates on
$V=\operatorname{span}\{v_a\}$, and let $C=(v_a)_{a\in S}$ be the
resulting full-row-rank $r\times s$ matrix, where $s=|S|$.
For a finite family $\mathcal W$ of integer weights, form a matrix
$P_{\mathcal W}$ with rows indexed by $S$ and columns indexed by
the pairs $(\omega,c)$ that occur as a monomial weight.  Substitutions
are interpreted as Laurent polynomials if necessary.  Set
\[
 (P_{\mathcal W})_{a,(\omega,c)}
   =\begin{cases}1&\langle\omega,a\rangle=c,\\0&\text{otherwise.}\end{cases}
\]
The following conditions are equivalent:
\begin{enumerate}
\item the coefficients of all the substitutions
$D(t^{\omega_1},\ldots,t^{\omega_n})$, taken jointly, span $V$;
\item $\operatorname{rank}(CP_{\mathcal W})=r$;
\item $\operatorname{im}(C^{\mathsf T})\cap
\ker(P_{\mathcal W}^{\mathsf T})=\{0\}$;
\item for every nonzero $\lambda\in V^*$, at least one substituted
scalar polynomial $\lambda(D(t^\omega))$ is nonzero.
\end{enumerate}
Equivalently, for some $r$-element set $T$ of columns of $P_{\mathcal W}$,
\[
 \sum_{\substack{B\subseteq S\\|B|=r}}
   \det(C_B)\det((P_{\mathcal W})_{B,T})\ne0.
\tag{I.14}\label{eq:I-joint-Plucker}
\]
Thus the failure locus on $\operatorname{Gr}(r,F^s)$ is the
intersection with the linear equations obtained by setting every
expression in (I.14) to zero.  These depend on Pl\"ucker coordinates,
including their values, rather than only on the coefficient matroid.

\paragraph{Proof.}
The $(\omega,c)$ column of $CP_{\mathcal W}$ is exactly
$\sum_{\langle\omega,a\rangle=c}v_a$, the coefficient at that
weight after substitution.  This proves (1)$\Leftrightarrow$(2).
Since $C^{\mathsf T}:F^r\to F^s$ is injective, the transpose of
$CP_{\mathcal W}$ is injective exactly when the intersection in
(3) is zero.  A proper subspace of $V$ has a nonzero annihilating
linear functional, proving (1)$\Leftrightarrow$(4).  Finally,
Cauchy--Binet applied to each $r$-column submatrix of
$CP_{\mathcal W}$ gives (I.14).  These are linear forms in the
Pl\"ucker coordinates $\det(C_B)$, proving the geometric description.

Joint preservation can hold even though no individual member is
basis isolating.  Take
\[
 D=e_1(x_1-x_2)+e_2(x_2-x_3),\qquad
 \mathcal W=\{(0,0,1),(1,0,0)\}.
\]
The two substituted vector polynomials are respectively
$(1-t)e_2$ and $(t-1)e_1$.  Each has coefficient rank one and so
fails I-03, but their coefficients jointly span the original
two-dimensional space.  To see why the matroid alone is insufficient,
over a field of characteristic different from two compare coefficient
rows $(1,1)$ and $(1,-1)$ on the support $\{x_1,x_2\}$ under weight
$(1,1)$.  They have the same rank-one coefficient matroid, while
their compressed rows are respectively $2$ and $0$.
\hfill$\square$

\paragraph{A sufficient black-box gate.}
If $S\subseteq\{0,\ldots,\delta\}^n$, all weights are nonnegative
and at most $B$, and joint preservation holds, evaluating the
$|\mathcal W|$ curves at $n\delta B+1$ distinct field elements
hits every nonzero scalar projection of $D$.  Indeed, condition
(4) supplies a nonzero univariate polynomial of degree at most
$n\delta B$.  This exact sufficient condition is weaker than
requiring one BIWA.  It does not remove cancellation from the
linear Pl\"ucker expressions in (I.14).

\subsection{I-15. Bounded affine decodings of the seed masks have width-two obstructions}
\label{I-15}

\paragraph{Statement.}
Suppose a proposed family of variable weights has the form
\[
 \omega(s)=b+Ts,\qquad
 b\in\mathbb Z^n,\quad T\in\mathbb Z^{n\times k},
\]
where every entry of $b$ and $T$ has absolute value at most $M$.
The seed set can be any subset of $\mathbb Z^k$.  If
\[
 2^n>(2nM+1)^{k+1},
\tag{I.15}\label{eq:I-affine-mask-count}
\]
there are distinct $\alpha,\beta\in\{0,1\}^n$ such that the
nonzero multilinear binomial $\mathbf x^\alpha-\mathbf x^\beta$
vanishes identically on every curve $\mathbf x=t^{\omega(s)}$.
This binomial has a width-two ROABP in every variable order.
In particular the whole family fails even the joint condition I-14.

\paragraph{Proof.}
Map the $2^n$ Boolean exponent vectors $\alpha$ to
\[
 \left(b^{\mathsf T}\alpha,T^{\mathsf T}\alpha\right)\in
 \{-nM,\ldots,nM\}^{k+1}.
\]
By (I.15), two distinct vectors have the same image.  Then
$\langle\omega(s),\alpha-\beta\rangle=0$ for every seed $s$,
so the two substituted powers of $t$ coincide.  This is a Laurent
identity if negative weights are allowed, and a polynomial identity
when the weights are nonnegative.  Distinct monomials have nonzero
coefficients in their difference over any field, including
characteristic two.  The diagonal matrix product with factors
$\operatorname{diag}(x_i^{\alpha_i},x_i^{\beta_i})$ and boundaries
$(1,-1)$ and $(1,1)^{\mathsf T}$ computes the binomial.  Reordering
the commuting diagonal factors gives every variable order.
\hfill$\square$

\paragraph{A necessary dimension for any bounded weight family.}
More generally, let $\mathcal W\subseteq\mathbb Z^n$ have all
coordinates of absolute value at most $B\ge1$, and let
$k=\dim_{\mathbb Q}\operatorname{span}_{\mathbb Q}\mathcal W$.
If $\mathcal W$ hits every multilinear binomial by monomial curves,
then necessarily
\[
 k\ge\frac{n}{\log_2(2nB+1)}.
\]
Indeed, select $k$ members of $\mathcal W$ forming a rational basis
and apply the same pigeonhole argument to their integer dot products
with Boolean exponent vectors.  Equality on this basis implies
equality on every member of $\mathcal W$, independently of how
the seed was decoded.  Thus polynomially bounded output weights
computed as polynomials of any fixed degree $h$ in
$\ell=O(\log n)$ Boolean seed bits also fail for sufficiently large
$n$: their span has dimension at most
$\sum_{j=0}^h\binom{\ell}{j}=O((\log n)^h)$, below the necessary
bound.  The dimension estimate follows by grouping the coefficient
vectors of the multilinear seed monomials; on Boolean seeds powers
of a bit reduce to that bit.

\paragraph{Application to the actual mask formula.}
Every weight in Hoza's formula (I.1) is an integer linear form in
the common $\ell$ seed bits.  Hence assigning those forms to
variables, or taking any fixed linear averaging or decoding with
polynomially bounded integer coefficients, gives $\omega(s)=b+Ts$
with $k=\ell$ and polynomially bounded entries.  A rational average
with a common polynomially bounded denominator has the same property
after clearing that denominator.  Since $\ell=O(\log(WN))$, if
$W,N$ and the resulting entry bound are polynomial in the number
$n$ of variables, condition (I.15) holds for all sufficiently large
$n$.  This rules out this specified seed-to-variable decoding,
even if all its members are used jointly.

For the most direct alphabet decoding the counterexample is explicit
without counting.  Let
$H_s(a)=\sum_j6^{\ell-j}C(a)_js_j$.  Matching letters $0,1$ by
an affine exponent map at layer $i$ gives
\[
 \omega_i(s)=6^{\nu_2(i)}\bigl(H_s(1)-H_s(0)\bigr).
\]
Layers $1$ and $3$ receive the same weight, so $x_1-x_3$ vanishes
for every seed.  This uses the actual layer dependence in (I.1)
and holds for every choice of its code.  There is no contradiction
with minimum-path isolation: the graph weights themselves need not
descend to these variable weights.  The argument does not rule out
nonlinear decoders whose outputs span sufficiently many independent
directions, or unbounded integer weights.

\subsection{I-16. The finite-field row-state lift and its uniform evaluation limit}
\label{I-16}

\paragraph{Statement.}
Over $F=\mathbb F_q$, an explicitly given width-$w$ ROABP
\[
 f=u^{\mathsf T}\prod_{i=1}^n
       \left(\sum_{a=0}^{\delta}A_{i,a}x_i^a\right)v
\]
has an exact deterministic coefficient-support automaton with
at most $q^w$ states per layer.  A state is a row vector
$r\in F^{1\times w}$; its letter-$a$ transition at layer $i$ is
$r\mapsto rA_{i,a}$.  The initial state is $u^{\mathsf T}$, and
a final state is accepting exactly when $rv\ne0$.
Nevertheless, if $0\le\delta<q$, every base-field hitting set
$H\subseteq F^n$ for all width-one individual-degree-$\delta$
polynomials satisfies
\[
 |H|\ \ge\ \left(\frac{q}{q-\delta}\right)^n.
\tag{I.16}\label{eq:I-small-field-lower}
\]
For $\delta=q-1$ all $q^n$ points are necessary and sufficient.
Thus polynomial state-space size alone cannot produce a polynomial
uniform base-field hitting set, even when $w=1$ and $q$ is fixed.

\paragraph{Proof.}
On the word $(a_1,\ldots,a_n)$, the automaton's final scalar is
$u^{\mathsf T}A_{1,a_1}\cdots A_{n,a_n}v$, precisely the monomial
coefficient.  Matrix multiplication has already performed all path
cancellation before the state is tested.  To obtain a simple layered
graph, subdivide each transition with an intermediate vertex labelled
by its current row and its letter; append a common sink after all
accepting states.  The width is at most $(\delta+1)q^w$ and the
length is $2n+1$.  It is constructible from the coefficient matrices
in time polynomial in this displayed size and the field-operation
cost.  Hoza applies to this coefficient-dependent graph, with the
corresponding polynomial-in-$q^w$ seed family; its path weights
still depend on the row states.

For the lower bound, choose independently a $\delta$-element subset
$R_i\subseteq F$ for each coordinate and consider
\[
 f_{R_1,\ldots,R_n}(\mathbf x)
     =\prod_{i=1}^n\prod_{a\in R_i}(x_i-a).
\]
These are $\binom{q}{\delta}^{n}$ distinct nonzero width-one polynomials.
A fixed point $\xi\in F^n$ hits exactly
$\binom{q-1}{\delta}^n$ of them, since each $R_i$ must avoid
$\xi_i$.  Covering all these polynomials therefore requires
\[
 |H|\binom{q-1}{\delta}^n\ge\binom{q}{\delta}^{n},
\]
which is (I.16).  The full grid $F^n$ hits every polynomial of
individual degree less than $q$: induction on the number of variables,
using the univariate root bound, shows that a polynomial vanishing
on that grid must be zero.  This also gives the matching upper bound
when $\delta=q-1$.  For $\delta\ge q$, even that grid cannot be a
hitting set for all formally nonzero polynomials, because
$x_1^q-x_1$ vanishes on it. \hfill$\square$

\paragraph{Precise scope.}
The lower bound concerns evaluations in the coefficient field.
It does not apply to extension-field evaluations, which can behave
very differently for a fixed base-field coefficient class.
The row-state lift is useful as an exact finite arithmetic simulation
and avoids the erroneous unweighted support graph of I-02.
It does not construct the graph from an evaluation oracle or make
its state-dependent path weights into variable cocharacters.

\subsection{I-17. A concrete uniform hitting set for invertible factors}
\label{I-17}

\paragraph{Background and statement.}
Invertible-factor ROABP hitting sets are already a result of
Agrawal et al. (2014), Theorem~3, pp.~5 and~19, and their
Lemma~27, p.~22, gives matrix coefficient concentration from
invertible constant terms.  The following is a self-contained
row-space specialization with constant boundary vectors; no
historical-priority claim is made.

Let $n,w\ge1$ and $\delta\ge1$.  Suppose
\[
 f=u^{\mathsf T}A_1(x_{\pi(1)})\cdots A_n(x_{\pi(n)})v
\]
has square $w\times w$ factors of individual degree at most
$\delta$, each with nonzero determinant polynomial.
Choose sets $S,B\subseteq F$ with
$|S|=nw\delta+1$, $|B|=\delta+1$, and $0\in B$.
Then the following single set, independent of the factors and order,
hits every nonzero such $f$:
\[
 \mathcal H=\{a\mathbf1+b: a\in S,\ b\in B^n,
                         |\operatorname{supp}(b)|\le w-1\}.
\tag{I.17}\label{eq:I-invertible-hit}
\]
Its cardinality is at most
\[
 (nw\delta+1)\sum_{j=0}^{\min(w-1,n)}\binom nj\delta^j.
\]
It is polynomial in $n,\delta$ for every fixed $w$, works in arbitrary
characteristic with the stated field-size condition, and requires
no knowledge of $\pi$.

\paragraph{Proof.}
The nonzero polynomial $\prod_i\det A_i(t)$ has degree at most
$nw\delta$, so some $a\in S$ makes every $A_i(a)$ invertible.
Reindex the variables in their true order for this proof only.
Set $P_0=I$, $P_i=A_1(a)\cdots A_i(a)$, and
\[
 M_i(y)=P_{i-1}A_i(a+y)P_i^{-1}.
\]
Then $M_i(0)=I$ and telescoping gives
\[
 f(a\mathbf1+\mathbf y)
    =u^{\mathsf T}M_1(y_1)\cdots M_n(y_n)P_nv.
\]
For nonzero $f$ we have $u\ne0$.  Let $R_i\subseteq F^{1\times w}$
be the span of the coefficient rows of
$u^{\mathsf T}M_1(y_1)\cdots M_i(y_i)$, and put $d_i=\dim R_i$.
Because the next constant term is $I$, these spaces are nested and
$d_0=1$, $d_i\le w$.

We claim that $R_i$ has a basis consisting of actual coefficient
rows whose monomials have support at most $d_i-1$.  This is true for
$i=0$.  Inductively retain a basis for $R_{i-1}$.  If the dimension
does not grow, that basis suffices.  Otherwise extend it using rows
$rM_{i,a}$, where $r$ is a retained coefficient row and $a>0$.
These rows span all new coefficients, and each is itself a coefficient
row with support at most $d_{i-1}$.  Since $d_i>d_{i-1}$, every
retained or added basis row has support at most $d_i-1$.
This proves the claim.

At $i=n$, at least one coefficient row has nonzero product with
$P_nv$.  Therefore some basis coefficient of support at most $w-1$
also has nonzero product.  Restricting the shifted scalar polynomial
to those variables, with the others zero, gives a nonzero polynomial
of individual degree at most $\delta$.  The grid $B$ on those
variables hits it by repeated univariate root counting.  This point
lies in (I.17).  The set (I.17) is invariant under coordinate
permutations, so its construction never needed the true order.
Counting supports proves its size bound. \hfill$\square$

The support bound is sharp for the normalized row-space argument:
for $w\ge2$ and $n=w-1$, factors $I+x_iE_{i,i+1}$ and boundaries $e_1,e_w$
compute $x_1\cdots x_{w-1}$.  For $w=1$, the theorem is the familiar
diagonal univariate test.  Generically singular factors require a
different argument, as the following explicit family shows.

\subsection{I-18. Regularization must preserve a unit Pl\"ucker coordinate}
\label{I-18}

\paragraph{Statement and counterexample.}
Let $m\ge2$, $n=2m$, and
\[
 f(\mathbf x)=\prod_{j=1}^m(x_{2j}-x_{2j-1}).
\]
This is a nonzero multilinear width-two ROABP in the displayed
variable order.  It vanishes at every point $a\mathbf1+b$ with
$|\operatorname{supp}(b)|<m$, for every scalar $a$.
Consequently the sampling restriction in I-17 cannot extend to
all width-two ROABPs by regularizing their matrices and specializing
the regularization parameter to zero.

More precisely, let $R=F[[\varepsilon]]$, $K=F((\varepsilon))$,
$C(\varepsilon)\in R^{r\times s}$ have
$\operatorname{rank}C(0)=r$, and let $P\in R^{s\times h}$ be any
integral compression.  Then
\[
 \operatorname{rank}_K(CP)=r
 \quad\Longleftrightarrow\quad
 \text{some $r$-minor of $CP$ has finite valuation},
\]
whereas
\[
 \operatorname{rank}_F(C(0)P(0))=r
 \quad\Longleftrightarrow\quad
 \text{some $r$-minor of $CP$ has valuation zero}.
\tag{I.18}\label{eq:I-unit-Plucker}
\]
The latter is equivalent to $\operatorname{im}_R(CP)=R^r$.
Generic span equality alone is insufficient.

\paragraph{Proof.}
Use the two singular matrix factors
\[
 L(x)=\begin{pmatrix}-x&1\\0&0\end{pmatrix},\qquad
 R(y)=\begin{pmatrix}1&0\\y&0\end{pmatrix}.
\]
Their product has $(1,1)$ entry $y-x$ and all other entries zero.
Taking the ordered product of these pairs with boundaries $e_1,e_1$
computes $f$.  At $a\mathbf1+b$ each pair difference is
$b_{2j}-b_{2j-1}$.  If fewer than $m$ coordinates of $b$ are
nonzero, one of the $m$ disjoint pairs has both entries zero,
so the product vanishes.

Replace every matrix factor by that factor plus $\varepsilon I$.
All regularized factors are generically invertible over $K$,
because their determinant is monic in $\varepsilon$.
Call the resulting scalar polynomial $f_\varepsilon$.
At the zero point a regularized pair is
\[
 (L(0)+\varepsilon I)(R(0)+\varepsilon I)
  =\begin{pmatrix}\varepsilon(1+\varepsilon)&\varepsilon\\
                   0&\varepsilon^2\end{pmatrix},
\]
so $f_\varepsilon(0)=\varepsilon^m(1+\varepsilon)^m\ne0$ in $K$.
Yet every common-shift sample of perturbation support less than $m$
has zero specialization at $\varepsilon=0$.

For the matrix assertion, a minor is nonzero over $K$ exactly when
it has finite valuation, and its reduction is nonzero over $F$
exactly when its valuation is zero.  This proves the two rank
criteria.  A unit $r$-minor gives $r$ columns invertible over $R$,
so their span is $R^r$.  Conversely, if the columns span $R^r$,
their reductions span $F^r$, which supplies a unit minor.
\hfill$\square$

\paragraph{The resulting special-fiber gate.}
In the scalar example the module generated by all coefficients of
$f_\varepsilon$ is $R$, since some coefficient of $f_0=f$ is nonzero.
For any finite collection of the displayed fixed $F$-valued samples
that contains the zero point, their values generate a nonzero ideal $N\subseteq
\varepsilon R$.  Hence $K N=K$ and its saturation
$(KN)\cap R$ is $R$, although all original sampled values are zero.
Saturating by dividing by a power of $\varepsilon$ thus does not
supply a nonzero evaluation of the original black box.
For vector coefficients the correct next target is a compression
with a unit Pl\"ucker coordinate, as in (I.18), or an equivalent
proof of coefficient-span preservation on the special fiber.
This refines the generic Grassmannian condition in I-14; it is not
established merely by flatness of a surrounding Rees family.

\paragraph{Background I-B06: the successful width-two repair.}
The width-two examples in I-15 and I-18 are failures of those
specified samplers.  Polynomial hitting sets for arbitrary
unknown-order univariate width-two ROABPs are already provided by
Agrawal et al. (2014), \S5.3, Theorem~34, p.~26, with the usual
field-size or extension convention required by the evaluation sets.
Their Lemma~32, p.~25, factors a nonzero singular $2\times2$ matrix
as a rank-one outer product after multiplying by one nonzero entry.
Applied at every singular factor, this writes a nonzero polynomial
multiple of the ROABP output as a product of invertible-factor
ROABP outputs.  The scalar factors may share the variables at the
cuts; product closure does not require disjointness.

Lemma~33, p.~26, supplies that product closure: interpolate a curve
through a common $h$-point hitting set for the factors.  Each nonzero
factor remains nonzero on the curve, so their product remains
nonzero.  If there are $m$ factors, each of total degree at most $D$,
the curve can have coordinate degree $h-1$, giving degree at most
$mD(h-1)$ and requiring that degree plus one parameter values.
The paper states the harmlessly looser $mDh+1$ bound.  For the
univariate-factor specialization of Theorem~34, the factor sparsity
is at most $\delta+1$ and maximum monomial support is one, so the
stated bound is polynomial.  These are cited positive results,
not new claims established by the present counterexamples.

\subsection{I-19. Higher-rank cuts and what block compression actually preserves}
\label{I-19}

\paragraph{Statement.}
Let $A(x)\in F[x]^{w\times w}$ have rank $1\le r<w$ over $F(x)$
and entry degrees at most $\delta$.  Choose $r$ rows $I$ and
$r$ columns $J$ with $M=A_{I,J}$ and $\Delta=\det M\ne0$.
Then
\[
 \Delta A=UV,\qquad
 U=A_{*,J},\quad V=\operatorname{adj}(M)A_{I,*},
\tag{I.19}\label{eq:I-rank-cut}
\]
where $U$ is $w\times r$, $V$ is $r\times w$,
$\deg U\le\delta$, $\deg V\le r\delta$, and
$\deg\Delta\le r\delta$.

Apply (I.19) to the singular factors at positions
$i_1<\cdots<i_m$ of an ROABP.  Write $X$ for their variables,
$Y_j$ for the disjoint intervening variable blocks, and $B_j$
for the product of the original matrices within block $Y_j$.
With $K=F(X)$, multiplication gives the exact identity
\[
 \left(\prod_{j=1}^m\Delta_j\right)f
 =u^{\mathsf T}B_0U_1
    \left(\prod_{j=1}^{m-1}V_jB_jU_{j+1}\right)V_mB_mv.
\tag{I.20}\label{eq:I-outer-cut}
\]
Over $K$, this is a matrix product of width at most
$\max_j r_j$ over disjoint blocks $Y_j$, with the corresponding
rectangular boundary factors.  Over $F$, neighboring displayed
factors generally share a cut variable; that dependence has not
been removed.

If polynomial substitutions into independent parameter sets preserve
the coefficient span of every block factor over $K$, then they
preserve the coefficient span of the full product, including scalar
nonvanishing.  This assertion requires independent parameter sets
for the different blocks.

\paragraph{Proof.}
The columns indexed by $J$ are a basis of the column space over
$F(x)$.  Expressing every column in that basis and restricting to
rows $I$ identifies its coordinate vector as $M^{-1}$ times its
row restriction.  Thus $A=A_{*,J}M^{-1}A_{I,*}$, and clearing
$\Delta$ proves (I.19) with the displayed degree bounds.
Substituting these factorizations into the ordered matrix product
and grouping the intervening factors proves (I.20).  Treating $X$
as elements of the coefficient field leaves exactly the disjoint
blocks $Y_j$; the dimensions in the product are the cut ranks.

For the span assertion, write each block factor as a sum of its
coefficient matrices times monomials in that block.  Since the
blocks are disjoint, the coefficients of the full product are
exactly the ordered products of one coefficient matrix from each
block, with the stated boundary contractions.  Replacing any
coefficient family by another with the same linear span leaves
the span of all these ordered products unchanged, by repeated
multilinearity.  Independent parameter sets preserve this coefficient
description after substitution, proving the claim.
\hfill$\square$

\paragraph{A valid local compression and the remaining coupling.}
For a fixed block, a hitting set for every nonzero scalar linear
combination of its matrix entries preserves its entire coefficient
span: otherwise a nonzero functional would annihilate all sampled
values but not the coefficient span.  A Lagrange curve through
these sample points also preserves that span.  Its coefficient
span contains all the sampled values and is contained in the
original coefficient span, so both inclusions are equalities.
Thus local compression is a concrete valid construction.

Two steps still require proof for a uniform black-box application:
the substitutions must accommodate the hidden block partition, and
the resulting nonzero rational dependence on $X$ must be hit after
returning from $K$ to the original field.  A simple example shows
why the latter cannot be silently dropped.  The rank-one matrices
\[
 A(x)=\begin{pmatrix}1&x\\x&x^2\end{pmatrix},\qquad
 A(y)=\begin{pmatrix}1&y\\y&y^2\end{pmatrix}
\]
with boundaries $e_1,e_1$ give $1+xy$.  After both variables are
put in $K=F(x,y)$, there are no remaining block variables to sample;
all nonvanishing information lies in the coefficient field.

For rank one, the positive product-closure result I-B06 tolerates
the shared cut variables.  For higher rank, (I.20) is a sum over
outer-state paths and retains cancellation.  Hoza isolation of
those paths supplies no proof of the joint-span condition I-14
for the original variables.  The supplied AGKS construction remains
a valid quasipolynomial completion of the general coefficient-span
task; the rank-cut construction above has not replaced its
coefficient and partition dependence by a polynomial procedure.

\paragraph{I-G02: a smaller structural sufficient gate.}
\label{I-G02}
For fixed width $w$, it would suffice to construct a uniform polynomial
hitting set for the expressions (I.20), with outer width at most
$w-1$, unknown block partition, and the shared cut variables kept
as genuine variables.  Their individual degrees are at most
$w\delta$: a cut variable occurs in $U_j$ and $V_j$, with combined
degree at most $(r_j+1)\delta$, and every other variable retains
degree at most $\delta$.  A hitting set for this class hits the
original nonzero $f$, because the multiplier $\prod_j\Delta_j$
is nonzero.  The case with no cuts is I-17, while outer width one
is covered by I-B06.  The first outer rank not handled by this
particular scalar-product repair is rank two.  This was the precise
sufficient gate reached by the rank-cut method.  The fixed-width gate for these original-coordinate cut expressions is
subsequently resolved by the rank-preservation and cone construction
I-31 below; I-30 gives the sharper six-pair result at width three.
No polynomial-in-width generator for the general coupled class is asserted.  For variable
$w$, the exponent depending on $w$ in I-17 is an additional barrier
to the desired polynomial bound.

\subsection{I-20. Bounded total corank admits a uniform extension of I-17}
\label{I-20}

\paragraph{Statement.}
For the constant-boundary ROABP in I-17, allow singular factors and
assume only that their total generic corank is at most an integer
$\kappa\ge0$:
\[
 \sum_{i=1}^n\bigl(w-\operatorname{rank}_{F(x_i)}A_i(x_i)\bigr)
 \le\kappa.
\]
Choose $S,G,B\subseteq F$ of sizes respectively
$nw\delta+1$, $w\delta+1$, and $\delta+1$, with $0\in B$.
For every subset $T\subseteq[n]$ of size at most $\kappa$, assign
the coordinates in $T$ arbitrary values in $G$.  Assign all other
coordinates the values $a+b_i$, where $a\in S$, $b_i\in B$, and
at most $w-1+\kappa$ of these $b_i$ are nonzero.  The union
$\mathcal H_\kappa$ of all such points hits every nonzero polynomial
with the stated representation.  Its size is at most
\[
 (nw\delta+1)
 \left(\sum_{c=0}^{\min(n,\kappa)}\binom nc(w\delta+1)^c\right)
 \left(\sum_{j=0}^{\min(n,w-1+\kappa)}\binom nj\delta^j\right),
\tag{I.21}\label{eq:I-corank-hit}
\]
which is $(nw(\delta+1))^{O(w+\kappa)}$.
The construction is independent of the variable order and of the
locations of singular factors.  It is polynomial for fixed
$w+\kappa$.

\paragraph{Proof.}
For a nonzero output every factor has positive generic rank.
Let $T$ be the true set of singular-factor variables; then
$|T|\le\kappa$.  For each such factor choose a nonzero maximal
minor $\Delta_i(x_i)$ of degree at most
$r_i\delta\le(w-1)\delta$.  Regard $f$ as a polynomial in the
remaining variables and choose any nonzero coefficient $g(X_T)$.
The nonzero polynomial
$g(X_T)\prod_{i\in T}\Delta_i(x_i)$ has individual degree at most
$w\delta$.  The grid $G^T$ therefore contains an assignment
$\gamma$ where it is nonzero.  At that assignment the remaining
polynomial is nonzero and each specialized singular factor has
exactly its generic rank $r_i$.

For the other factors, one common $a\in S$ makes all constant
terms $A_i(a)$ invertible, as in I-17.  Gauge these constant terms
to $I$ by products of invertible matrices, omitting the singular
constant factors from those products.  Concretely, set $P_0=I$,
let $P_i=P_{i-1}A_i(a)$ at a nonsingular layer and
$P_i=P_{i-1}$ at a specialized singular layer, and replace every
factor by $P_{i-1}A_iP_i^{-1}$.  Telescoping adjusts only the final
boundary.  Nonsingular layers now have constant term $I$;
singular layers are constant matrices of rank $r_i$.

Follow the coefficient row spans in this normalized product.
At a nonsingular layer the dimension can only grow.  At a constant
rank-$r_i$ layer it can only decrease, by at most $w-r_i$, because
the kernel of right multiplication has that dimension in the full
row space.  These spans are all nonzero, since the final polynomial
is nonzero.  Starting at dimension one and ending at dimension at
most $w$, the sum $U$ of all dimension increases is consequently
at most
\[
 U\le w-1+\sum_{i\in T}(w-r_i)\le w-1+\kappa.
\]
The greedy argument of I-17 now selects actual coefficient rows
whose support is at most the cumulative number of dimension
increases.  At a constant cut retain a basis of the images of the
old basis; this introduces no new support.  At a layer where the
span grows, every new candidate is an old basis row times one
nonconstant coefficient, so its support grows by at most one,
while the cumulative dimension increase grows by at least one.
This proves the claimed support bound inductively.

Some selected final row has nonzero scalar boundary contraction.
The grid $B$ on its support supplies a nonzero value of the
remaining shifted polynomial.  Together with the chosen cut
values $\gamma$, this is a point of $\mathcal H_\kappa$.
Enumerating all subsets and grid values gives (I.21).  The family
is invariant under coordinate permutations, proving the claimed
unknown-order uniformity. \hfill$\square$

The restriction is explicit: arbitrary ROABPs can have total
generic corank growing with $n$.  I-20 does not bound that parameter
or supply a polynomial procedure when it grows.  It shows that
rank cuts can be absorbed uniformly when the total dimension loss
is bounded; the remaining problem is to compress unboundedly many
such losses without discarding the original coefficient relations.

\subsection{I-B07. Literature scope and the source of the new rank-preservation attempt}
\label{I-B07}

The distinction between a representation in one hidden order and a
width bound in every order is essential.  Gurjar, Korwar, and Saxena
(2017), Theorem~3.6 and Corollary~3.7, give the polynomial
constant-width result with the order known; their Theorem~4.9 treats
the separate every-order class.  Guo and Gurjar (2020), Theorem~1,
give an improved quasipolynomial bound for one unknown order,
whereas their Theorem~3 concerns every-order representations.
Hu, van Melkebeek, and Morgan (2024), Theorem~1.9, give a width
bound with an additional low-degree-monomial promise; that promise
cannot be omitted from the statement.  Shpilka (2026), ECCC
TR26-314, Theorem~1.1 and Remark~1.3, gives a polynomial bound for
multilinear every-order ROABPs and their border.  It does not
assert the same bound for one hidden order.

Two 2026 preprints encountered in the current search make claims
about one unknown order.  Singh and Singh, arXiv:2602.13449v1, is
not used as an accepted theorem here: the subsequent source audit
I-22 gives explicit counterexamples to two foundational reductions.  Liu, Hexagon 2611.00634v1, supplies the rank-profile
mechanism specialized and proved in full as I-21 below.  Its
unrefereed status and disclosed AI provenance are retained in the
source ledger.  The accepted mathematical assertion in this report
is the explicitly proved specialization I-21, not an endorsement
of either preprint as a whole.

\subsection{I-21. A checked width-three specialization of the rank-preservation construction}
\label{I-21}

\paragraph{Source and status.}
The mechanism below is a specialization of Liu (2026),
\emph{Black-box identity testing for affine orbits of read-once
oblivious algebraic branching programs}, Hexagon
\nolinkurl{2611.00634v1}, Section~3.1, Theorem~3.2 and Lemma~3.3,
and Section~3.2, Lemmas~3.4 and~3.6, printed pp.~5--7.
This is an unrefereed preprint dated November 19, 2026; its host
discloses primarily AI-generated text and no human-understood parts.
Its identifier is a Hexagon identifier, not an arXiv identifier.
We give a complete proof of the narrower original-coordinate,
width-three statement.  The affine-orbit assertions, the general-width
bounds, and the other sections of that manuscript are not being
certified here.  No historical priority is claimed.

\paragraph{Statement.}
Let $F$ be any field, let $n,D\ge1$, and let $f\in F[x_1,\ldots,x_n]$
have total degree at most $D$.  Suppose that for some permutation
$\sigma$ it has a constant-boundary representation
\[
 f=uA_1(x_{\sigma(1)})\cdots A_n(x_{\sigma(n)})v,
 \qquad A_i\in F[x_{\sigma(i)}]^{3\times3}.
\]
Smaller widths and rectangular layers are included by zero padding.
Let $E/F$ be a working field with at least
\[
 M=\max\{n,D+1,(n-1)D+1\}
\]
distinct available elements.  Choose distinct $a_1,\ldots,a_n\in E$,
and let $\lambda_i(Y)$ be their degree-$(n-1)$ Lagrange polynomials.
Define the twenty-one-parameter map
\[
 \mathcal G_i(t,\mathbf y,\mathbf z)
   =t+\sum_{j=1}^{10}z_j\lambda_i(y_j),
 \qquad 1\le i\le n.
 \tag{I.22}\label{eq:I-width3-generator}
\]
Then $f\ne0$ implies $f(\mathcal G)\ne0$ as a polynomial over $E$.
No variable order is an input to the map.

Put $R=(n-1)D+1$.  Take any sets $T,Y\subset E$ with
$|T|=D+1$ and $|Y|=R$.  For distinct
$\theta_0,\ldots,\theta_D\in E$, put
\[
 \Sigma_{10,D}=
 \{(\theta_{b_1},\ldots,\theta_{b_{10}}):
       b_j\ge0,\ \sum_j b_j\le D\}.
\]
The evaluations of $\mathcal G$ on
$T\times Y^{10}\times\Sigma_{10,D}$ form a hitting set of size at most
\[
 (D+1)\bigl((n-1)D+1\bigr)^{10}\binom{D+10}{10}.
 \tag{I.23}\label{eq:I-width3-size}
\]
This is polynomial in $n,D$ and is valid in every characteristic.
The field convention is explicit: if $E\ne F$, the black box must
allow evaluation over $E$.  No base-field query simulation is being
assumed.  In particular this does not contradict the finite-field
obstruction in I-16.

\paragraph{1. The coefficient-space lemma needed at width three.}
For a field $L$, work in
\[
 \mathcal R=L[Y_1,\ldots,Y_m]/(Y_1^{e_1+1},\ldots,Y_m^{e_m+1}),
 \qquad e_i\ge0,
\]
and suppose $B_i$ depends only on $Y_i$.  A matrix has \emph{unit rank}
$r$ if its constant matrix has rank $r$ and all its $(r+1)$-minors
vanish in $\mathcal R$.  Unit rank zero means that the matrix is zero.
For any constant row $b$, if every interval of the word
$B_1\cdots B_m$ has unit rank in $\{0,2,3\}$, then
\[
 \operatorname{Coeff}(bB_1\cdots B_m)
   =\operatorname{span}_L
       \{[Y^\alpha](bB_1\cdots B_m):
                         |\operatorname{supp}\alpha|\le4\}.
 \tag{I.24}\label{eq:I-profile23}
\]
If every interval has unit rank three, the corresponding bound is two.
Here $\operatorname{Coeff}$ means the span of coefficient rows.

We prove this assertion, including the stable-projection step.
An interval of unit rank zero makes the whole product zero.  Otherwise
the full constant product has rank two or three.  In rank three all
constant factors are invertible.  Cut-dependent constant changes of
basis make them identities.  The coefficient row spaces of the
successive prefixes are then nested.  Starting with one row, they can
grow at most twice.  At a growth layer one new nonconstant local
coefficient increases support by at most one; at a stable layer
constant propagation preserves a spanning set.  This proves the
bound two, with actual coefficient rows, as in I-17.

In rank two write the constant matrices as $C_i$.  Choose a
two-dimensional row space $U_0$ on which right multiplication by
$C_1\cdots C_m$ is injective, and set $U_i=U_0C_1\cdots C_i$.
Choose a one-dimensional complement $V_m$ of $U_m$.  At cut $i$ take
\[
 V_i=\{z:zC_{i+1}\cdots C_m\in V_m\}.
\]
The suffix has rank at least two and maps $U_i$ isomorphically onto
$U_m$.  Thus $V_i$ has dimension one, complements $U_i$, and
$V_{i-1}C_i\subseteq V_i$.  Transport bases of the $U_i$ along the
constants and choose bases of the $V_i$.  The resulting constant
matrices are
\[
 C_i=\operatorname{diag}(I_2,q_i),\qquad q_i\in L.
 \tag{I.25}
\]
These basis changes preserve supports and all interval unit ranks.

Let $\pi$ denote projection of a row onto its first two coordinates.
For a starting constant row space $S_0$, propagate
$S_i=\operatorname{Coeff}(S_{i-1}B_i)$.  Since
$\pi(sC_i)=\pi(s)$, the spaces $\pi S_i$ are nested.
Consider first a segment on which they are all the same space $T$.
Choose basis lifts of $T$ in $S_0$, with initial cost zero, and a basis
of the kernel of $\pi|_{S_0}$, with initial cost one.  Define the cost
of an obtained coefficient row as its initial cost plus the number
of nonconstant local coefficients used.  We claim that cost at most
one spans the final space on this stable segment.

If every $q_i$ on the segment is nonzero, the constants are invertible.
All the spaces lie in $\pi^{-1}(T)$, of dimension $\dim T+1$.
If the initial kernel is present, the initial space is already this
entire ambient space and constant propagation suffices.  Otherwise
there is room for at most one dimension increase.  Before that
increase the propagated constant lifts span; at the increase one
nonconstant coefficient suffices; afterwards constant propagation
again spans.  The claimed cost bound follows.

If some $q_i=0$, let $c$ be the last such position, put
$P=B_c\cdots B_m$, and let $B$ be the matrix of the initial lifts
propagated through the constants before $c$.  The suffix $P$ has
unit rank two and constant matrix $\operatorname{diag}(I_2,0)$.
Write its upper blocks as $[U\ V]$.  The constant part of $U$ is
$I_2$, so $U$ is invertible in $\mathcal R$.  The vanishing
three-by-three determinant, or its Schur complement, shows that
every row $(x,y)$ in the row module of $P$ satisfies
$y=xU^{-1}V$.

Let $H=\pi B$, whose rows form a basis of $T$.  If $T=0$, this row
graph forces every propagated row through $P$ to be zero.  Otherwise
write
\[
 \pi(BP)=QH,\qquad \pi(sP)=\eta H\quad(s\in S_{c-1}).
\]
Such expressions exist because all coefficient projections stay in
$T$, and $Q(0)=I$.  Therefore $Q$ is invertible.  The row graph gives
\[
 sP=\eta Q^{-1}BP.
 \tag{I.26}
\]
Coefficients of a polynomial matrix times $BP$ are linear
combinations of coefficient rows of $BP$, also in the truncated
ring.  Conversely the rows of $B$ belong to $S_{c-1}$.  It follows
that the final coefficient space is exactly
$\operatorname{Coeff}(BP)$.
Restart just after $c$ from $\operatorname{Coeff}(BB_c)$.
Constant lifts have cost zero and a kernel basis has cost at most
one, since only one variable position was used.  The remaining
constants are invertible; the preceding case therefore proves the
same cost bound.  This completes the stable-segment claim.

For the original single starting row $b$, first suppose $\pi b\ne0$.
The active projected space starts in dimension one and has at most
one strict growth.  The first stable segment costs at most one,
the growth factor at most one more, and the final stable segment
at most one more: support at most three suffices.  If $\pi b=0$,
the first stable segment has zero projected space.  The
stable-segment claim, with the sole initial row assigned cost one,
says that this segment is spanned by coefficients with no
nonconstant positions, or is zero.  There can subsequently be at
most two projection-growth factors.  Each costs at most one, and
each following stable segment at most one.  Support at most four
therefore suffices.  At every restart, lifts and kernel vectors
are linear combinations of the coefficient rows already retained;
the span filtration by support is preserved.  This proves (I.24).

The bound four in (I.24) is sharp.  With $P=\operatorname{diag}(1,1,0)$,
take
\[
 B_1=P+Y_1E_{31},\quad B_2=I+Y_2E_{13},\quad
 B_3=P+Y_3E_{32},\quad B_4=I+Y_4E_{23},\quad b=e_3^T.
\]
Every interval containing $B_1$ or $B_3$ has generic and constant
rank two, and the other intervals are invertible.  The output row is
\[
 Y_1e_1^T+Y_1Y_2Y_3e_2^T+Y_1Y_2Y_3Y_4e_3^T.
\]
Its third coefficient direction requires support four.

\paragraph{2. Lowest homogeneous components.}
Suppose every interval of a polynomial matrix word has the same
generic rank as at zero, with these ranks in $\{2,3\}$.  For constant
endpoints $b,c$, every monomial in the lowest nonzero homogeneous
component of $bB_1\cdots B_mc$ has support at most four.  If all
interval ranks are three, the bound is two.

Indeed, choose such a monomial $Y^\alpha$ and reduce modulo
$Y_i^{\alpha_i+1}$ for every $i$.  Only monomials dividing
$Y^\alpha$ survive.  Lowest-degree minimality makes the surviving
scalar polynomial exactly its nonzero $Y^\alpha$ term.  All the
intervals retain their unit ranks: the pivot constant minor remains
invertible and the higher minors were already zero polynomially.
Apply (I.24), or its invertible version.  If the support of $\alpha$
were too large, the final boundary $c$ would annihilate every
spanning coefficient row and hence also the surviving one, a
contradiction.

\paragraph{3. A finite collection of local rank tests.}
Augment the original matrix word by the constant factors $e_1u$
and $ve_1^T$.  Its full product has only one possibly nonzero entry,
namely $f$ in position $(1,1)$.  All intervals have a well-defined
generic rank over the rational function field.  Translation of
each distinct variable by any center in an overfield preserves
these generic ranks.

Fix $r=1$ or $r=2$, and suppose all intervals of generic rank
greater than $r$ are already preserved at the current center.
For any target interval $[l,h]$ of rank $r$, let $[l,i_0]$ be its
first prefix of rank $r$ and require this prefix to be preserved.
For each $i_0\le i<h$, choose the shortest rank-$r$ suffix $[j,i]$
of $[l,i]$, and require $[j,i+1]$ to be preserved.  Each such test
has generic rank $r$.  Its interior $[j+1,i]$, and the prefix
$[l,i_0-1]$ of the first test, are empty or have generic rank
greater than $r$ by minimality.  Thus a test has one such chunk
and at most one factor at each end.

These tests suffice.  Once $[l,i]$ is preserved, its specialized
row space is contained in that of $[j,i]$.  Both have dimension
$r$: the first has rank $r$, and the second cannot exceed its
generic rank $r$.  Their row spaces are equal.  Multiplying by
the next specialized factor and using the tested rank of
$[j,i+1]$ proves preservation of $[l,i+1]$.  This proves the
assertion by induction, for all target intervals.  Select one
nonzero $r$-minor of each test; no efficient selection or knowledge
of the hidden word is needed in the black-box construction.

For $r=2$, each selected minor is an entry of the word of second
compound matrices, by Cauchy--Binet.  These compounds have size
three.  The interior chunk and every one of its subintervals have
rank three, and so do their constants.  In a lowest-degree
monomial of the selected minor, extract the prescribed powers at
the at most two end factors.  The remaining scalar chunk
contraction is still in its lowest degree.  Part~2 gives support
at most two inside the chunk.  The whole monomial consequently
has support at most four.

For $r=1$ the compounds are the original matrices.  The chunk and
its subintervals have ranks in $\{2,3\}$, already preserved at
the center.  The same argument gives support at most four
inside the chunk and at most six in the whole selected entry.
Empty chunks and constant end factors only reduce these bounds.

\paragraph{4. Simultaneous preservation with four then six pairs.}
For $k\ge1$ write
\[
 \operatorname{SV}_{k,i}(\mathbf y,\mathbf z)
      =\sum_{j=1}^kz_j\lambda_i(y_j).
\]
Over any coefficient overfield of $E$, this map detects any
nonzero polynomial having a nonzero monomial of support at most
$k$: set the selectors to the Lagrange nodes of that support,
and the unused amplitudes to zero.  The resulting restriction
to that coordinate subspace retains the chosen coefficient.
In particular it detects a nonzero homogeneous polynomial all
of whose monomials have support at most $k$.

Start at the common center $(t,\ldots,t)$ over $E(t)$.
An interval of generic rank three consists entirely of invertible
univariate factors.  Its determinant at the center is a product
of nonzero univariate determinant polynomials evaluated at the
transcendental $t$, and is therefore nonzero.  Thus all rank-three
intervals are preserved simultaneously.

Add an $\operatorname{SV}_4$ perturbation with fresh independent
selectors and amplitudes.  For every chosen rank-two test minor,
its expansion about the current center is a nonzero polynomial
whose lowest homogeneous component $H$ has support at most four
in each monomial, by part~3.  Hence $H(\operatorname{SV}_4)\ne0$.
The map is homogeneous of degree one in its amplitudes.  The
lowest amplitude degree in the full substituted minor is
therefore exactly the degree of $H$, and its component is
$H(\operatorname{SV}_4)$; higher homogeneous terms cannot cancel
it.  Each selected test minor is consequently nonzero in the
same rational function field of the new parameters.  This is
simultaneous preservation, not a separate numeric choice for
each test: finitely many nonzero polynomials have nonzero product
in the parameter polynomial ring.  Earlier preserved minors
remain nonzero because setting the new amplitudes to zero
recovers their previous values.  The local-test induction now
preserves every interval of ranks two and three.

Over that enlarged rational function field, add an independent
$\operatorname{SV}_6$ perturbation.  The rank-one test components
have support at most six by part~3, so the same argument preserves
all these tests and all earlier ranks.  The augmented full word
has generic rank one whenever $f\ne0$.  Its specialized rank is
therefore one, which says precisely that the resulting scalar is
nonzero.  The two perturbations use disjoint parameter pairs;
their sum is $\operatorname{SV}_{10}$.  This proves (I.22).
All uses of the variable order occurred only in this existential
proof about the hidden representation.  The actual map is the
same for every order.

For comparison with the failed repair I-18, the polynomial
$\prod_{j=1}^m(x_{2j}-x_{2j-1})$ already survives one interpolation
pair: its pullback is
$z^m\prod_j(\lambda_{2j}(y)-\lambda_{2j-1}(y))$, which is nonzero
because the distinct Lagrange basis polynomials have nonzero
differences.  Interpolation directions are generally dense;
the construction therefore does not require a sparse perturbation
point to witness every nonzero polynomial.

\paragraph{5. Finite sampling and complexity.}
Since $f$ has total degree at most $D$, its pullback under (I.22)
has degree at most $D$ in $t$, at most $(n-1)D$ in each selector
$y_j$, and total degree at most $D$ in all ten amplitudes together.
Repeated univariate interpolation therefore leaves a nonzero
amplitude polynomial at some point of $T\times Y^{10}$.

The set $\Sigma_{10,D}$ detects every nonzero amplitude polynomial
of total degree at most $D$, over any characteristic.  To see this,
use the Newton basis
\[
 N_{\mathbf b}(\mathbf z)
  =\prod_{j=1}^{10}\prod_{h=0}^{b_j-1}(z_j-\theta_h),
 \qquad \sum_jb_j\le D.
\]
Its leading monomials are exactly those of total degree at most
$D$.  Evaluation at the node indexed by $\mathbf c$ vanishes
unless $c_j\ge b_j$ for every $j$.  Ordered by total degree,
the evaluation matrix is triangular with nonzero diagonal.
It is invertible, proving the claimed detection and (I.23).
Given the field nodes, computing the Lagrange polynomials and
enumerating the displayed family uses polynomially many field
operations in $n,D$, since its ten-pair parameter count is fixed.
Over a finite field, one may use an explicitly represented
extension large enough to supply the stated nodes.  Over
$\mathbb Q$ the nodes can be chosen as integers of polynomial
magnitude, giving rational coordinates of polynomial bit length.
\hfill$\square$

\paragraph{What this closes, and what it leaves open.}
This supplies a concrete positive answer for the width-three,
one-hidden-order case of I-G02, including arbitrarily many
rank-two cuts.  The rank-cut expression in I-19 equals
$f\prod_j\Delta_j$; the extra univariate factors can be absorbed
into the width-three layers, so the same theorem also hits those
particular coupled expressions with their increased degree bound.
The construction uses rank preservation directly and does not
derive a BIWA or a Hoza cocharacter decoder.
The general polynomial-in-$w$ gate I-G01 remains open in this
report.  In particular, I-08 requires widths growing with the
matrix dimension; this fixed-width result does not provide that
polynomial Noether-normalization conclusion.


\subsection{I-22. Two explicit failures of the active-simple-block reduction}
\label{I-22}

\paragraph{Audited claim and scope.}
Singh and Singh (2026), arXiv:2602.13449v1, Theorem~7.4
(printed p.~38), claims the unknown-order bound
$n(wd)^{O(w^2)}$.  Its foundational reductions are false as stated:
the equivalence after Definition~2.1 (p.~4), repeated in
Section~3.1 (p.~6), and the active-block conclusions in
Lemma~3.1 (pp.~6--7), Remark~3.1.4 (p.~9), and Lemma~3.2
(p.~10).  The following counterexamples work over $\mathbb C$,
inside the source's characteristic-zero hypothesis.  They
invalidate those proof steps; they do not disprove the headline
existence of a hitting set with the claimed bound.

\paragraph{Full word algebra loses the prescribed ordered output.}
Let $N=E_{12}$ and $M=E_{21}$, and take the width-two factors
\[
 A_1(x_1)=M+x_1N,\qquad A_2(x_2)=M+x_2N.
\]
Since $N^2=M^2=0$, $NM=E_{11}$, and $MN=E_{22}$,
\[
 A_1(x_1)A_2(x_2)=x_1E_{11}+x_2E_{22}.
\]
The specified ROABP output $e_1^TA_1A_2e_2$ is identically zero.
Nevertheless the algebra generated by the coefficient matrices
is the full $\operatorname{Mat}_2(\mathbb C)$, and
$e_1^TNe_2=1$.  Thus vanishing on the actual ordered products
of one local coefficient per layer does not imply vanishing on
all words in the generated algebra.

This example also separates a nonzero simple-block image of the
matrix product from a nonzero prescribed scalar output.  The
algebra is already simple, so its identity representation is a
simple-block projection and the displayed matrix product survives
it, although the scalar output is zero.  After distinct nonzero
specializations of $x_1,x_2$, the two specialized factors still
generate the full matrix algebra: their difference gives a
nonzero multiple of $N$, and then either factor gives $M$.
Hence moving to a generic multiplicative specialization does not
repair this example.

\paragraph{A nonzero output can live entirely in the radical.}
Take four width-two factors
\[
 A_1=x_1E_{11},\quad A_2=I+x_2N,\quad
 A_3=I+x_3N,\quad A_4=x_4E_{22}.
\]
Direct multiplication gives
\[
 W=A_1A_2A_3A_4=x_1x_4(x_2+x_3)N,\qquad
 f=e_1^TWe_2=x_1x_4(x_2+x_3)\ne0.
 \tag{I.27}
\]
The coefficient algebra is the algebra $\mathcal T_2$ of upper
triangular two-by-two matrices.  Its Jacobson radical is
$J=\mathbb C N$, with $J^2=0$, and its semisimple quotient is
$\mathbb C\oplus\mathbb C$, given by the two diagonal entries.
For completeness, $\mathbb C N$ is a square-zero ideal, hence
lies in the radical; the quotient is a product of fields and
has zero radical, proving equality.  Every homomorphism onto
a simple quotient therefore kills $W$.  No active simple block
of the kind required by the source exists for this nonzero
computation.

All four variables occur, the product length is $4\ge w=2$,
and the example is not a width-one scalar product in disguise:
the coefficient matrix of $x_2+x_3$, with rows and columns
indexed by exponents $0,1$, is
$\left(\begin{smallmatrix}0&1\\1&0\end{smallmatrix}\right)$,
of rank two.  The same obstruction persists after any nonzero
character values for which $x_2+x_3\ne0$; the specialized
algebra remains $\mathcal T_2$ and the product remains a
nonzero radical element.  Constant identity factors can be
inserted to make the product arbitrarily long.

Nilpotency $J^w=0$ controls products containing $w$ radical
factors.  It does not imply that a long product of arbitrary
algebra elements vanishes whenever the final product belongs
to $J$.  Equation~(I.27) gives precisely this distinction.
These two counterexamples are enough to withhold the cited
Theorem~7.4 from the accepted background of this report.
They leave the independently proved I-21 unaffected.
\hfill$\square$


\subsection{I-23. A binary-code cost for preserving decomposable exterior states}
\label{I-23}

\paragraph{The specified next construction.}
One proposed repair is to keep the state of an exterior-power lift
on its Grassmann cone by requiring each Hoza seed to act diagonally
on the Pl\"ucker coordinates in a decomposability-preserving way.
We analyze exactly this requirement.  The statements concern
formal algebraic one-parameter actions over a field, not just
their values on the finitely many elements of a finite field.

Fix $2\le r\le w-2$, index Pl\"ucker coordinates by the
$r$-subsets $I\subseteq[w]$, and let $\rho_I\in\mathbb Z$.
The diagonal action $p_I\mapsto t^{\rho_I}p_I$ preserves all
decomposable $r$-vectors if and only if
\[
 \rho_I=c+\sum_{i\in I}a_i
 \quad\hbox{for some }c,a_1,\ldots,a_w\in\mathbb Z.
 \tag{I.28}\label{eq:I-plucker-additive}
\]
Consequently, if a binary code
$C:\binom{[w]}r\to\{0,1\}^{\ell}$ has minimum distance $d$
and every seed of the Hoza weighting
\[
 \rho_s(I)=\sum_{j=1}^{\ell}6^{\ell-j}C(I)_js_j,
 \qquad s\in\{0,1\}^{\ell},
\]
preserves decomposability, then
\[
 \ell\ge\frac{wd}{2}.
 \tag{I.29}\label{eq:I-plucker-code-cost}
\]
The additional layer multiplier $6^{\nu_2(i)}$ in I-B01
does not change this conclusion.

\paragraph{Proof of the cocharacter classification.}
Let $K$ have size $r-2$, and let $a,b,c,d$ be four different
indices outside $K$.  The Pl\"ucker relation on those indices
forces
\[
 \rho_{Kab}+\rho_{Kcd}
 =\rho_{Kac}+\rho_{Kbd}
 =\rho_{Kad}+\rho_{Kbc}.
\]
Here subscripts denote sets, with the usual orientation signs
in the coordinate relation.  To verify the force of this
condition directly, restrict to planes containing the fixed
coordinate $(r-2)$-plane on $K$.  In the remaining two rows,
use four columns $(e_1,e_2,e_1,e_2)$, then
$(e_1,e_1,e_2,e_2)$.  The relation has exactly two nonzero
terms in each case.  After the diagonal action their two powers
of the indeterminate $t$ must agree.  This gives the displayed
equalities, in every characteristic.

For distinct $i,j$, consider
$\rho_{S\cup\{i\}}-\rho_{S\cup\{j\}}$ with
$|S|=r-1$ and $S\cap\{i,j\}=\varnothing$.
The four-index equality shows that this difference is unchanged
when one element of $S$ is exchanged.  The graph of such subsets
under single exchanges is connected, so the difference is a
constant $d_{ij}$ independent of $S$.  A common $S$ disjoint
from three specified indices exists because $r\le w-2$;
it gives $d_{ij}+d_{jk}=d_{ik}$.  Take $a_1=0$ and
$a_i=d_{i1}$.  Then
$\rho_I-\sum_{i\in I}a_i$ is unchanged by every single exchange
of an $r$-subset.  All $r$-subsets are connected by such
exchanges, so this is one integer $c$.  This proves necessity
of (I.28).  Conversely (I.28) is the exterior action induced by
$\operatorname{diag}(t^{a_1},\ldots,t^{a_w})$, followed by
one common scalar, and plainly preserves decomposable vectors.

\paragraph{Binary-valued additive functions are individual membership bits.}
Suppose $b(I)\in\{0,1\}$ has the form (I.28).
Every exchange difference $a_i-a_j$ lies in $\{-1,0,1\}$.
If $b$ is nonconstant, the $a_i$ therefore have exactly two
levels differing by one.  Let $S$ index the higher level.
Apart from an additive constant, $b(I)=|I\cap S|$.
If both $|S|$ and $w-|S|$ are at least two, then among
$r$-subsets this intersection count ranges from
$\max(0,r-(w-|S|))$ to $\min(r,|S|)$, a range of length
at least two when $2\le r\le w-2$.  That contradicts binary
values.  Hence $|S|=1$ or $w-|S|=1$.
The binary values then force $b(I)$ to be either membership
of one element or its complement.  Constant bits are the only
remaining possibilities.

Apply this classification to the seed $s=e_j$.  The positive
integer factor $6^{\ell-j}$ does not affect the required
exchange equalities, so each code bit is a constant or a single
membership/complement test.  Let $m_i$ count the nonconstant
bits testing element $i$.  Bases differing only by exchange of
$i$ and $j$ have Hamming distance $m_i+m_j$.  Such bases exist
for every pair.  Thus $m_i+m_j\ge d$ for every $i\ne j$.
Summing over all pairs gives
\[
 (w-1)\sum_i m_i\ge\binom w2d,
 \qquad \ell\ge\sum_i m_i\ge wd/2.
\]
This proves (I.29). \hfill$\square$

\paragraph{Consequence for this repair.}
On the $W=\binom wr$ exterior states, using the stated Hoza
error bound $2^{-d}W^4N^2\le1/2$ requires
$d\ge\lceil\log_2(2W^4N^2)\rceil$.
The all-seeds decomposability requirement then forces
$\ell=\Omega(w\log(WN))$ and at least
$(2W^4N^2)^{w/2}$ seeds.  Already $r=2$ gives
$W=\Theta(w^2)$, so this loses the polynomial family size
without invoking an exponentially large exterior dimension.

The restriction is essential: this is not a lower bound for
all generators or all nonlinear decodings of Hoza seeds.
It does not rule out a separately proved successful subset
of coherent seeds from another code, or a construction that
leaves the Grassmann cone at intermediate stages and restores
the required coefficient information afterwards.  Nor would
mere Pl\"ucker compatibility by itself resolve the distinct
variable-quotient condition in I-02.


\subsection{I-24. Exponential support is possible inside the preserved-rank test class}
\label{I-24}

\paragraph{Statement.}
For every $w\ge3$, there is a word of $w\times w$ univariate
linear matrices on
\[
 h_w=3\cdot2^{w-2}-2
 \tag{I.30}\label{eq:I-reset-support}
\]
distinct variables with the following properties.  Every
nonempty interval has the same generic rank as its constant
matrix, and that rank is at least two.  Nevertheless one scalar
contraction of the full word is exactly the product of all
$h_w$ variables.  Its lowest homogeneous component therefore
has support $h_w$.  With constant rank-one factors appended
at both ends it is an actual final rank-one test of the
rank-descent construction used in I-21.

Thus a polynomial-in-$w$ bound on the support of every such
lowest component is false, even for $r=1$, where every nonzero
exterior state is automatically decomposable.  This is a
counterexample to that particular compression proposal;
it is not a hitting-set size lower bound, since a monomial is
easily hit by a point with all coordinates nonzero.

\paragraph{Construction.}
Let $P_r=\operatorname{diag}(I_r,0_{w-r})$ for $2\le r<w$.
For $a\le r$ define a two-position gadget, on fresh variables,
\[
 Q_{r,a}(s,t)=(P_r+sE_{w,a})(P_r+tE_{a,w}).
\]
Each of its factors has constant and generic rank exactly $r$:
its leading $r\times r$ block is $I_r$, and the extra row or
column is a multiple of a row or column already present.
All the constant matrices $P_r$ are nested projections.
Consequently any interval of the eventual word has constant
product $P_{\min r_i}$ and generic rank at most the minimum
rank of its factors.  The nonzero constant minor gives the
reverse inequality.  Every interval therefore has generic
rank exactly its constant rank $\min r_i\ge2$.
The same assertion is a unit-rank statement after any separate
truncations of the variables.

Define words $H_k$, using fresh variable copies on every
occurrence, by
\[
 H_w=\text{empty word},\qquad
 H_k=H_{k+1}\,Q_{k,k}\,H_{k+1}
       \quad(3\le k<w).
\]
The final word is
\[
 \mathcal W_w=H_3\,Q_{2,1}\,H_3\,Q_{2,2}\,H_3,
\]
with all five displayed blocks using disjoint fresh variables.
If $L_k$ is the number of positions in $H_k$, then
$L_w=0$ and $L_k=2L_{k+1}+2$, whence
$L_k=2(2^{w-k}-1)$.  Thus $\mathcal W_w$ has
$3L_3+4=h_w$ positions.

\paragraph{Proof of the scalar identity.}
If an input row $q$ has $q_a=0$, then the two-position gadget
acts by
\[
 qQ_{r,a}(s,t)
   =\sum_{j\le r}q_je_j^T+q_ws\,e_a^T+q_wst\,e_w^T.
 \tag{I.31}
\]
In particular it erases coordinates strictly between $r$ and
$w$, leaves the earlier coordinates unchanged, and transfers
the last coordinate through both fresh variables.

Inductively, suppose an input row has coordinates
$k,\ldots,w-1$ zero and last coordinate $c$; its earlier
coordinates may be arbitrary.  The word $H_k$ leaves those
earlier coordinates unchanged and makes its last coordinate
$c$ times the product of every variable in $H_k$.
This is immediate for $H_w$.  For $H_k$, the first
$H_{k+1}$ leaves coordinate $k$ zero.  Formula~(I.31) for
$Q_{k,k}$ multiplies the last coordinate by its two variables
and resets coordinates $k+1,\ldots,w-1$ to zero.  The second
$H_{k+1}$ therefore applies under the same induction
hypothesis.  This proves the claim.

Start the full word at $e_w^T$.  After the first $H_3$,
coordinates one and two remain zero.  The gadget $Q_{2,1}$
uses the zero first coordinate and resets all coordinates
$3,\ldots,w-1$.  The second $H_3$ preserves coordinate two
as zero, so $Q_{2,2}$ again satisfies (I.31) and resets the
higher coordinates.  The final $H_3$ completes the induction.
Therefore
\[
 e_w^T\mathcal W_we_w=\prod_{i=1}^{h_w}Y_i.
\]
This argument is valid over every field; there are no signed
cancellations or characteristic restrictions.

\paragraph{Membership in the rank-test class.}
Append $L=e_1e_w^T$ on the left and $R=e_we_1^T$ on the
right.  The full product is the displayed nonzero monomial
times $e_1e_1^T$, so it has generic rank one.  Every prefix
beginning with $L$ is nonzero, since the last row of each
partial gadget has a surviving coordinate; consequently these
prefixes have rank one.  Every interval lying wholly inside
$\mathcal W_w$ has rank at least two, as already proved.
For the prefix ending just before $R$, the shortest rank-one
suffix therefore starts at $L$.  Appending $R$ gives exactly
the full word as the last test in part~3 of I-21.
Its interior chunk satisfies every higher-rank preservation
hypothesis at the zero center.  The exponential support cannot
be attributed to failure to preserve one of those interval
ranks. \hfill$\square$

At $w=3$ this yields support four; at $w=4$ it yields support
ten; at $w=5$ it yields support twenty-two.  The construction
shows why a small intrinsic state-variety dimension alone does
not bound how many variable positions are needed after many
singular resets.  A useful Hoza repair would have to exploit
more than low support or decomposability: it must handle the
original-variable coefficients and their cancellations while
allowing such long supports.


\subsection{I-25. Intrinsic Grassmannian compression succeeds for invertible chunks}
\label{I-25}

The nested-image method is suggested by Liu (2026), Section~3.3,
printed pp.~7--8.  The following restricted refinement uses
the projective Grassmannian of one decomposable exterior row
as its ambient variety and is proved directly.

\paragraph{Statement.}
Let $A_i(Y_i)$ be $w\times w$ matrices over a field $L$, in
distinct variables, with every $A_i(0)$ invertible.  Fix a
nonzero decomposable row $p\in\bigwedge^r L^w$ and a linear
functional $\ell$ on $\bigwedge^rL^w$, where $1\le r<w$.
If
\[
 q(\mathbf Y)=\ell\bigl(p(\textstyle\bigwedge^r A_1(Y_1))
                  \cdots(\textstyle\bigwedge^r A_m(Y_m))\bigr)
 \ne0,
\]
then there is a set $S\subseteq[m]$ with
$|S|\le r(w-r)$ for which setting $Y_i=0$ outside $S$
keeps $q$ nonzero.  In particular $q$ has a nonzero monomial
of support at most $r(w-r)$, and $\operatorname{SV}_{r(w-r)}$
detects it.

For an $r$-minor test with an invertible-constant interior
chunk and at most one factor on either end, the whole minor
therefore has some nonzero monomial of support at most
$r(w-r)+2$.  This replaces the arbitrary compound-space bound
$\binom wr-1+2$ for this restricted class.  It bounds the
existence of a low-support monomial, not every monomial in the
lowest homogeneous component.

\paragraph{Proof.}
Use cut-dependent invertible constant changes of basis to make
every $A_i(0)=I$.  These changes arise from $\operatorname{GL}_w$
and hence preserve decomposability after taking exterior powers;
they adjust only the fixed endpoints.  Extend scalars to an
algebraic closure.  This does not change whether a polynomial
or its coordinate restriction is zero.

Restrict each parameter to the nonempty open set where its
factor is invertible.  Let $Z_i$ be the Zariski closure in
$\operatorname{Gr}(r,w)$ of the projective prefix states
\[
 [p(\textstyle\bigwedge^r A_1(Y_1))\cdots
        (\textstyle\bigwedge^r A_i(Y_i))].
\]
Each $Z_i$ is irreducible, as the closure of the image of an
irreducible product of open affine lines.  The value $Y_i=0$
is allowed and acts as the identity, so
$Z_{i-1}\subseteq Z_i$.  The starting variety $Z_0$ is one
point.  The Grassmannian has dimension $r(w-r)$: its dense
matrix chart consists of row spaces of $[I_r\ X]$ with
$X\in L^{r\times(w-r)}$.
Thus at most $r(w-r)$ of these inclusions are strict.

Whenever $Z_i=Z_{i-1}$, replace $A_i(Y_i)$ by $I$.
This preserves the closure of the prefix image and hence the
closure of every later image.  To justify the latter assertion,
on invertible factor values the exterior action is a regular
map on the projective Grassmannian; the closure of the image
under the next factor of a dense prefix set is the same as
the closure obtained from the whole prefix variety.  This
persists on composing successive factors.  Inductively one
may therefore remove every non-growth position while retaining
the final projective image closure.  At most $r(w-r)$ positions
remain.

The scalar $q$ vanishes identically exactly when that final
projective image lies in the hyperplane cut out by $\ell$.
It suffices to test this on the invertible-value open set,
since a polynomial vanishing there vanishes identically.
Preserving the image closure therefore preserves nonvanishing
of $q$.  A nonzero restriction to the retained coordinates
contains a monomial supported there, proving the first claim.

For the endpoint version, adjoin the at most two endpoint
variables to the coefficient field.  A selected row of an
$r$th compound endpoint is the wedge of $r$ rows of the
original endpoint matrix, hence is decomposable; it is nonzero
if the selected full minor is nonzero.  The right endpoint
supplies a linear functional.  The preceding argument applies
to the interior chunk over this rational function field.
The resulting restricted minor is still a polynomial in the
original variables and is nonzero there.  Restoring the at
most two endpoint positions gives the asserted support bound.
\hfill$\square$

\paragraph{Exact limitation.}
This construction exploits the exterior geometry instead of
treating $\bigwedge^r L^w$ as an arbitrary vector space of
dimension $\binom wr$.  It is an existence-and-enumeration
compression: the retained positions depend on the matrices,
and enumerating them or using the stated SV map still has an
exponent depending on $w$.  The proof also requires invertible
constant interior factors.  I-24 shows that replacing this
condition merely by preservation of all interval ranks at
least two cannot yield a polynomial-in-$w$ support bound.
The next useful rank-descent construction must accommodate
those singular resets and the original-variable cancellation
relations simultaneously; intrinsic dimension alone does
not supply that extension.


\subsection{Open gate I-G01 and the order of further work}
\label{I-G01}

The surviving main gate is an explicit polynomial family of
nonnegative integer cocharacters of polynomial magnitude which, for every
width-$w$ matrix product with disjoint univariate factors,
enters at least one cone (I.3), or satisfies the exact weaker joint
coefficient-span condition I-14 for the family.  If this is
proved, I-B02 or I-14 gives polynomial hitting sets, and I-08 gives the
precisely delimited matrix-invariant conclusion.  No finding
above claims that this gate is closed.

The constructions already attempted and resolved as stated are:
direct state weights (I-02); monomial-alphabet lifting followed
by exact evaluation interpolation (I-06); positive Gram basis
enumeration (I-05); independent and correlated Haar repair
(I-07); exterior/permutation expansion and literal subspace
configuration graphs (I-11); and positive coefficient-norm
transfer matrices (I-09), exact universal energy cubature (I-13),
and the arithmetic support test (I-12).  The further repairs now
resolved are joint span preservation (I-14), bounded affine mask
decoding (I-15), exact finite-field row states (I-16), the restricted
invertible-factor sampler (I-17), and epsilon regularization with
its necessary special-fiber condition (I-18).  Rank cuts give the exact
block reduction and coupled gate I-19/I-G02, and bounded total corank
has the explicit positive sampler I-20.  The subsequent I-21
rank-preservation construction resolves the original-coordinate
width-three case with one hidden order, in all characteristics
under its explicit working-field convention.  The positive result from
direct graph isolation is I-01.  The specified next constructions have now been tested further:
I-23 proves the code-length cost of making every Hoza seed preserve
decomposable exterior states; I-24 gives exponential support even
inside preserved-rank test chunks; I-25 supplies the valid intrinsic
Grassmannian compression for invertible chunks.  None is a general
PIT lower bound.  A next construction may use dense interpolation
directions to bypass low-support bounds, but must prove nonvanishing
for the actual rank-test coefficients and retain polynomial dependence
on variable width as well as on the number of variables.

The later refinements sharpen this status.  I-26 gives a centered dense-SV
annihilator family and a logarithmic necessary pair count; I-B08 records
the known matching order of pair count for the every-order subclass.
I-27 rules out filtering a short actual Hoza code into coherent exterior
weights.  I-28 and I-30 reduce the general width-three sampler to six
pairs.  I-29 gives a sharp two-rank support bound and a polynomial pair
budget with unbounded cuts at one singular rank.  I-31 proves an
original-coordinate sampler for every fixed width, with pair count
controlled by the number of rank levels and at most
$2^{O(w\log w)}$ for all ranks.  Its interpolation size remains far
from polynomial in variable width.  Thus the fixed-width obstruction
is resolved in this track, while I-G01's genuinely polynomial
family-size and weight-size requirements remain unproved.


\subsection{I-26. Genuine annihilators of the centered dense interpolation map}
\label{I-26}

\paragraph{Relation to prior work.}
Studying the vanishing ideal of the SV generator has an established
determinantal approach.  Hu, van Melkebeek, and Morgan (2024),
Theorem~1.3 (printed p.~6) and Corollary~1.5 (printed p.~7), give explicit
multilinear degree-$(k+1)$ annihilators of $\operatorname{SV}_k$
on $2k+1$ variables, settling the degree question discussed by
Fournier and Korwar (2018, WACT slides~13--19).
Their Theorems~1.9--1.10 retain a low-degree-monomial hypothesis
in the ROABP hitting/lower-bound assertion.  The construction
below follows this determinantal perspective, explicitly adds
invariance under the common translation, and proves an
invertible-factor ROABP representation.  No priority is claimed
for a vanishing-ideal construction in this setting.

\paragraph{Statement.}
Let $k\ge1$, let $n\ge3k+2$, and let a field $F$ contain
distinct Lagrange nodes $a_1,\ldots,a_n$.  For those nodes there
is an explicitly constructible nonzero homogeneous multilinear
polynomial $P\in F[X_1,\ldots,X_n]$ of degree $k+1$ such that
\[
 P\bigl(t\mathbf1+\operatorname{SV}_k(\mathbf y,\mathbf z)\bigr)=0.
 \tag{I.32}\label{eq:I-centered-sv-zero}
\]
It has a constant-boundary ROABP representation in every order
of width at most
\[
 W_k=\sum_{j=0}^{k+1}\binom{k+1}{j}^{\!2}
       =\binom{2k+2}{k+1},
 \tag{I.33}\label{eq:I-centered-sv-width}
\]
with commuting factors $I+X_iN_i$, where $N_i^2=0$.
Every interval of this representation is invertible, and its
constant matrix is the identity.  Thus, after constant
rank-one endpoints are added, $P$ is a genuine final rank-one
test whose interior satisfies all higher-rank preservation
requirements at zero.

In particular a centered $\operatorname{SV}_{k(w)}$ generator
valid for all widths at most $w$ and all $n$, over fields or
extensions supplying the required distinct nodes, must have
$W_{k(w)}>w$.  Equivalently, for each proposed $k$ the
counterexample requires an admissible $n\ge3k+2$.
Since $W_k\le4^{k+1}$, this implies
$k(w)=\Omega(\log w)$.  It does not rule out a logarithmic
number of pairs, and is not a lower bound for other generators.
All the general assertions hold in every characteristic under
the stated distinct-node hypothesis.

\paragraph{A determinant that removes the common center.}
Set $m=k+1$ and let
\[
 U_{i\ell}=a_\ell^i,\qquad 0\le i\le k,\quad1\le\ell\le n.
\]
In the $2k+2$ dimensional vector space of polynomials of degree
at most $2k+1$, impose the $k+1$ linear conditions
\[
 \sum_{\ell=1}^n a_\ell^i q(a_\ell)=0,
       \qquad0\le i\le k.
\]
Their common kernel has dimension at least $k+1$.
Choose linearly independent polynomials
$q_1,\ldots,q_m$ in that kernel by Gaussian elimination, and
put $V_{j\ell}=q_j(a_\ell)$.  Then $UV^T=0$.
Evaluation on the $n$ distinct nodes is injective on polynomials
of degree at most $2k+1$, since $n>2k+1$; hence $V$ has row
rank $m$.

Define
\[
 B(\mathbf X)=U\operatorname{diag}(X_1,\ldots,X_n)V^T,
 \qquad P(\mathbf X)=\det B(\mathbf X).
 \tag{I.34}
\]
Every maximal minor $U_S$ is a nonzero Vandermonde determinant.
At least one maximal minor $V_S$ is nonzero because $V$ has
full row rank.  Cauchy--Binet gives
\[
 P(\mathbf X)=\sum_{|S|=m}\det(U_S)\det(V_S)
                                  \prod_{\ell\in S}X_\ell.
\]
The nonzero minor of $V$ therefore supplies a nonzero coefficient.
This proves nonvanishing, homogeneity, and multilinearity in
every characteristic.

Because $UV^T=0$, one has $B(\mathbf X+t\mathbf1)=B(\mathbf X)$.
For every $i,j$, the polynomial $Y^i q_j(Y)$ has degree at most
$3k+1\le n-1$.  Exact Lagrange reproduction consequently gives
\[
 \sum_{\ell=1}^n a_\ell^i q_j(a_\ell)\lambda_\ell(Y)
    =Y^i q_j(Y).
\]
It follows that
\[
 B(\operatorname{SV}_k)
   =\sum_{s=1}^k z_s
       (1,y_s,\ldots,y_s^k)^T
       (q_1(y_s),\ldots,q_m(y_s)).
\]
This matrix is a sum of $k$ rank-one matrices but has size
$m=k+1$, so its determinant is identically zero.  The common
translation has no effect on $B$, proving (I.32).

\paragraph{An invertible, any-order ROABP of the stated width.}
For a generic $m\times m$ matrix $Z$, form the column vector
$\mathcal M(Z)$ of all square minors, including the empty
minor $1$.  There are $W_k$ such minors by the elementary
Vandermonde sum in (I.33).  They are linearly independent as
polynomials: distinct row/column index sets have distinct
row and column multidegrees, and each individual determinant
is nonzero.

For fixed vectors $u,v$ and an indeterminate $s$, expansion
of a minor under a rank-one update gives
\[
 \det(Z_{I,J}+s u_Iv_J^T)
 =\det Z_{I,J}
  +s\sum_{a\in I,\ b\in J}
        (-1)^{\operatorname{pos}_I(a)+\operatorname{pos}_J(b)}
        u_av_b\det Z_{I\setminus\{a\},J\setminus\{b\}}.
\]
Terms using two update columns vanish, since those columns are
proportional.  Thus there is a constant matrix $N(u,v)$ such that
\[
 \mathcal M(Z+suv^T)=(I+sN(u,v))\mathcal M(Z).
\]
Translation twice in the same direction and linear independence
of the minor polynomials imply $N(u,v)^2=0$ by comparing the
coefficient of two independent translation parameters.  Two
different rank-one translations commute, so their representing
matrices commute as well.  These arguments compare polynomial
identities and do not divide by factorials.

Let $u_\ell$ and $v_\ell$ be the $\ell$th columns of $U$ and
$V$, and set $N_\ell=N(u_\ell,v_\ell)$.  Starting at $Z=0$
and translating in these $n$ rank-one directions gives
\[
 \mathcal M(B(\mathbf X))
     =\prod_{\ell=1}^n(I+X_\ell N_\ell)\,\mathcal M(0).
\]
The coordinate of the full determinant is $P$.  The factors
commute and satisfy
$(I+X_\ell N_\ell)^{-1}=I-X_\ell N_\ell$.
This gives the advertised width-$W_k$ ROABP in any variable
order, with invertible factors and identity constants.
The starting row and final column may be taken as the
full-determinant selector and the empty-minor selector.
Adding their rank-one endpoint matrices produces an actual
rank-one test exactly as in I-21.  Every interior interval
has generic and constant rank $W_k$, so no missing rank
preservation hypothesis accounts for this annihilation.
\hfill$\square$

\paragraph{A width-four, one-pair example.}
Over $\mathbb Q$, use $n=5$ and nodes $0,1,2,3,4$.
The polynomials
\[
 q_1(Y)=Y^2-4Y+2,\qquad
 q_2(Y)=Y^3-\tfrac{77}{5}Y+\tfrac{54}{5}
\]
satisfy both kernel conditions.  Expanding (I.34) gives
$\det B=\frac65 Q$, where
\[
\begin{split}
 Q={}&3X_1X_2-4X_1X_3-15X_1X_4+16X_1X_5
       +4X_2X_3+8X_2X_4-15X_2X_5\\
    &+4X_3X_4-4X_3X_5+3X_4X_5.
\end{split}
\tag{I.35}
\]
Thus $Q(t\mathbf1+\operatorname{SV}_1)=0$ and $Q\ne0$.
For a smaller representation than the general width-six bound,
take row vectors and column vectors
\[
\begin{array}{c|ccccc}
i&1&2&3&4&5\\ \hline
a_i&(1,0)&(0,1)&(-28/97,-4/97)&(3/16,0)&(0,0)\\
b_i&(0,0)^T&(3,0)^T&(-4,4)^T&(-15,8)^T&(16,-15)^T.
\end{array}
\]
In blocks of sizes $1,2,1$, put
\[
 N_i=\begin{pmatrix}0&a_i&0\\0&0_{2\times2}&b_i\\0&0&0\end{pmatrix}.
\]
All triple products of these $N_i$ vanish, and
$e_1^TN_iN_je_4=a_i b_j$.
Consequently
\[
 e_1^T\prod_{i=1}^5(I_4+X_iN_i)e_4
      =\sum_{i<j}(a_ib_j)X_iX_j=Q.
\]
The factors are upper triangular with diagonal entries one,
so every interior interval remains invertible at every point.
This refutes the one-pair proposal already at width four;
it does not assert that width four is the smallest possible
width of a centered one-pair annihilator of arbitrary degree.

\paragraph{A concrete repair for this family, and its scope.}
The same counterexamples are detected by a different simple
map.  Since every $U_S$ is nonzero, the support of $P$ is
exactly the set of bases of the column matroid of $V$.
Pairwise distinct element weights have a unique minimum-weight
matroid basis, by the greedy exchange property.  Hence
\[
 P(T,T^2,\ldots,T^n)\ne0,
\]
with degree at most $n(k+1)$.
This is an instance of the basis argument in I-03 and gives a
polynomial scalar sampling grid for this particular family:
choose $n(k+1)+1$ distinct scalar values, in an extension if
the ground field is too small.  This additional grid-size
condition is separate from the $n$ distinct Lagrange nodes.
It demonstrates that the dense-SV counterexample is a precise
test of that map, not a general obstruction to cocharacter
sampling.  The open dense-map question is whether a
logarithmic or polynomial number of pairs suffices for the
full rank-test class, or whether a different structured family
can handle its cancellations with polynomial dependence on
width.  Neither assertion is proved by this example.


\subsection{I-B08. The logarithmic dense-generator bound already holds in the every-order subclass}
\label{I-B08}

Forbes, Saptharishi, and Shpilka (2014), full version
arXiv:1309.5668, Definition~1.1 (printed p.~3), use the term
\emph{commutative ROABP} for a polynomial admitting a width
bound in \emph{every} variable order.  This is not the class
with one hidden order.  Pairwise commuting layer coefficient
matrices, as in I-26, give a representation in this subclass
by reordering the same factors.

Their Corollary~4.3 (printed p.~22), using Theorem~4.1
(printed pp.~21--22), states that if $f$ has individual degree
less than $d$, admits width-$w$ representations in every order,
and $|F|>nd$, then
\[
 f\ne0\quad\Longleftrightarrow\quad
 f\bigl(\operatorname{SV}_{1+2\lfloor2\log_2w\rfloor}\bigr)
       \ne0.
\]
The stated hitting-set size is
$\operatorname{poly}(n,d)^{O(\log w)}$.
One may use a sufficiently large extension when required.
This known logarithmic pair bound applies to the general
commuting, unipotent representation constructed in I-26.
Consequently its logarithmic necessary pair count matches
the order of a known sufficient bound for that restricted
class.  The new question retained after I-26 concerns the
full one-hidden-order rank-test class, and even a logarithmic
pair bound would still require a separate efficient sampling
argument to give polynomial rather than quasipolynomial size.

The primary definition, theorem, and corollary, including
the local factorization use of the every-order assumption
and the additive composition of two SV maps, were inspected.
The complete article and all of its subsidiary rank-condenser
proofs are not certified here.


\subsection{I-27. Filtering the actual Hoza masks does not restore a short coherent code}
\label{I-27}

\paragraph{The precise repair being tested.}
I-23 required every seed of a code to preserve the Grassmann cone.
One could instead keep only those seeds that happen to preserve it.
For the actual base-six weighting, this filtering has an exact
description and cannot combine noncoherent bits into a coherent
action.  This is a statement about that formula and formal
cocharacters; it is not a lower bound for arbitrary nonlinear
decoders or dense interpolation maps.

Let $2\le r\le w-2$, let
$C:\binom{[w]}r\longrightarrow\{0,1\}^{\ell}$ be any binary
code, and put
\[
 \rho_s(I)=\sum_{j=1}^{\ell}6^{\ell-j}s_j C(I)_j.
 \tag{I.36}
\]
Call a bit $j$ coherent if $I\mapsto C(I)_j$ satisfies the
Pl\"ucker exchange equalities in I-23.  Call a seed coherent
if the diagonal action $p_I\mapsto t^{\rho_s(I)}p_I$
preserves every decomposable $r$-vector.
Then
\[
 s\text{ is coherent}
 \quad\Longleftrightarrow\quad
 s_j=0\text{ for every noncoherent bit }j.
 \tag{I.37}
\]
Thus the set of coherent seeds is a coordinate face of the
Boolean seed cube, rather than a collection of cancellations
between noncoherent code positions.

For any selected family $\mathcal S$ of coherent seeds, let
$T\subseteq[w]$ consist of the ground elements tested by its
active nonconstant bits.  By the binary classification of I-23,
each such bit is membership of one element or its complement,
so $|T|\le\ell$.  The entire family of weights depends on an
$r$-subset $I$ only through $I\cap T$.  In particular, if
$\ell\le w-2$, there are distinct $u,v\notin T$ such that
\[
 \rho_s(K\cup\{u\})=\rho_s(K\cup\{v\})
 \quad\text{for all }s\in\mathcal S
 \tag{I.38}
\]
whenever $|K|=r-1$ and $K\cap\{u,v\}=\varnothing$.
Consequently a short actual code cannot be repaired merely
by discarding its noncoherent seeds and using the surviving
weights as ground-variable cocharacters.

\paragraph{Proof of the exact filtering statement.}
For each Pl\"ucker exchange equality, write its deficit for
bit $j$ as
\[
 d_j=C(Kab)_j+C(Kcd)_j-C(Kac)_j-C(Kbd)_j.
\]
Each $d_j$ is an integer in $[-2,2]$.  Coherence of the seed
requires $\sum_j 6^{\ell-j}s_jd_j=0$ for this exchange,
and likewise for the other pairing.  If there were a largest
active power $6^h$ with nonzero coefficient, its absolute
contribution would be at least $6^h$, whereas all lower
powers together contribute at most
\[
 2\sum_{a=0}^{h-1}6^a
   =\frac25(6^h-1)<6^h.
\]
This is impossible.  Hence $s_jd_j=0$ for every $j$ and
every exchange.  An active bit must therefore be coherent.
Conversely a sum of coherent weight functions satisfies all
exchange equalities and hence preserves decomposability by
I-23.  This proves (I.37).  The common positive layer
multiplier in Hoza's formula does not affect the argument.
The deficit argument is an equality of integer exponents,
so it is valid in every coefficient-field characteristic.

For the family statement, each active coherent nonconstant
bit is a membership or complementary membership function of
one element, by I-23.  All other active bits are constant.
It follows directly that $\rho_s(I)$ is determined by
$I\cap T$.  Two untested elements exist when $|T|\le w-2$,
and exchanging them proves (I.38).  A suitable $K$ exists
since $r-1\le w-2$.

\paragraph{An explicit nonzero cancellation witness.}
Write the induced coherent weights as
$\rho_s(I)=c_s+\sum_{i\in I}a_{s,i}$.
Equation (I.38) forces $a_{s,u}=a_{s,v}$ for every seed;
this equality is independent of the common gauge ambiguity
in the choice of the $a_{s,i}$.  The nonzero polynomial
\[
 f(X_1,\ldots,X_w)=X_u-X_v
\]
therefore vanishes under every induced monomial substitution
$X_i\leftarrow Z^{a_{s,i}}$.  Adding the same scalar center
to every coordinate does not change that conclusion.
The assertion also holds for Laurent exponents; adding a
common exponent to make all exponents nonnegative preserves
the equality of the two relevant coordinates.

This is already an invertible-factor width-two example:
\[
 f=e_1^T(I_2+X_uE_{12})(I_2-X_vE_{12})e_2,
\]
with identity factors at the other positions.  Equivalently,
the decomposable vector
$e_K\wedge(X_u e_u+X_v e_v)$ has two Pl\"ucker coordinates,
and their oriented difference gives this same polynomial.
Thus the witness belongs to the algebraic setting that the
coherent decoder was intended to preserve, rather than merely
showing a combinatorial collision between unused codewords.
\hfill$\square$

\paragraph{Magnitude and scope.}
To distinguish all $r$-subsets by such a coherent family,
at least $w-1$ ground elements must be tested.  Consequently
at least $w-1$ distinct nonconstant bit positions must occur
in its seeds.  Among their powers of six, one is at least
$6^{w-2}$.  A seed activating that bit, and a codeword on
which the bit is one, give a raw weight at least $6^{w-2}$.
This last observation concerns the unmodified weights;
rescaling an individual cocharacter can change its magnitude,
so it is not asserted as an invariant degree lower bound
under arbitrary further decodings.  The exact face
description (I.37) and the short-code collision (I.38) remain
valid even if common factors of individual seeds are removed.

Longer codes followed by a separately justified small family,
or maps that leave the Grassmann cone, are not excluded.
Nor does this obstruct adding a dense SV direction: for
distinct nodes, $\lambda_u(Y)-\lambda_v(Y)$ is nonzero and
already destroys this particular linear collision.  I-26
tests that next dense repair by a different, determinant
counterexample, and leaves logarithmically many pairs open
for the general one-hidden-order rank-test class.


\subsection{I-28. Nonzero endpoint constants fit inside the Grassmannian support budget}
\label{I-28}

\paragraph{Statement.}
The two extra endpoint positions in I-25 are unnecessary when
the endpoint vectors are nonzero at the center.  Let
$1\le r\le w-1$, put $g=r(w-r)$, and consider
\[
 q(Y_0,\ldots,Y_{m+1})
 =\left\langle u(Y_0)
       \prod_{i=1}^m\bigwedge^r A_i(Y_i),\ v(Y_{m+1})\right\rangle.
 \tag{I.39}
\]
Here $A_i(0)$ is invertible, $u(Y_0)$ is a polynomial
decomposable exterior row with $u(0)\ne0$, and $v(Y_{m+1})$
is a polynomial exterior column with $v(0)\ne0$.
Only the left row is required to be decomposable.
All variables displayed are distinct; constant endpoints may
be included by omitting their variables.
If $q\ne0$, setting all but at most $g$ variable positions
to zero leaves a nonzero polynomial.  In particular
$\operatorname{SV}_{g}$ hits this class, and the count includes
both endpoint variables.

\paragraph{Proof.}
Extend the coefficient field to an algebraic closure; a zero
restriction over the original field stays zero after this
extension, so proving existence there suffices.  Use constant
changes of basis to make all $A_i(0)=I_w$, as in I-25.
They preserve nonzero endpoint constants and decomposability.
Restrict each internal parameter to the nonempty open set
where $A_i$ is invertible, and the left parameter to the open
set where $u$ is nonzero.  These opens contain zero.

Let $Z_{-1}$ be the projective point $[u(0)]$, and let $Z_0$
be the closure of the image of $[u(Y_0)]$.  For $i\ge1$,
let $Z_i$ be the closure of the projective image after applying
the first $i$ internal matrices.  Each $Z_i$ is an irreducible
closed subvariety of $\operatorname{Gr}(r,w)$, and
\[
 Z_{-1}\subseteq Z_0\subseteq\cdots\subseteq Z_m.
\]
For the first inclusion, use the available parameter value
$Y_0=0$; for the others, use $A_i(0)=I_w$.

Retain a parameter position exactly when this inclusion is
strict.  There are at most $\dim Z_m$ retained positions,
because a strict inclusion of irreducible closed varieties
strictly increases dimension, starting at a point.
Setting all other positions to zero preserves the final
image closure $Z_m$.  To check this rigorously by induction,
at a retained position the previous restricted image is
dense in the same $Z_{i-1}$; applying the regular invertible
actions and taking closure gives the same $Z_i$.  At a
nonretained position, zero substitution applies the identity
and leaves $Z_{i-1}=Z_i$.  The initial left curve obeys the
same statement: if its image closure is a point, setting its
parameter to zero preserves that point.

Write $\ell_0$ for the nonzero linear functional determined
by $v(0)$.  If $\ell_0$ does not vanish identically on $Z_m$,
set the right parameter to zero as well.  The retained image
is still dense in $Z_m$, so the resulting scalar polynomial
is nonzero, using at most $\dim Z_m\le g$ positions.

Otherwise $Z_m$ lies in the hyperplane $\ell_0=0$.
The Grassmannian is not contained in that hyperplane:
the coordinate decomposable wedges span the full exterior
space, and $\ell_0\ne0$.  Therefore
\[
 \dim Z_m\le \dim\bigl(\operatorname{Gr}(r,w)
                                  \cap\{\ell_0=0\}\bigr)\le g-1.
\]
As $q\ne0$, some value of the right parameter gives a
functional that does not vanish identically on $Z_m$.
The retained prefix image has the same closure, so keeping
this right parameter free gives a nonzero scalar polynomial
on at most $\dim Z_m+1\le g$ positions.  This proves the
restriction assertion in both cases.  A nonzero restriction
to at most $g$ coordinate positions is a specialization of
$\operatorname{SV}_g$, by setting its selectors to the
corresponding Lagrange nodes.  Thus the generator assertion
follows in every characteristic, with the usual distinct-node
field or extension convention.  No algorithm to find the
retained positions from the unknown matrices is assumed.
\hfill$\square$

\paragraph{Application to the width-three construction.}
In the rank-two stage of I-21, a selected nonzero minor of
a local rank test is a contraction of a second-compound
word with an invertible interior.  Its compound row and
column at either endpoint are nonzero at the initial common
center $x_i=t$.  Indeed, a nonzero selected contraction has
nonzero endpoint rows and columns as univariate polynomial
vectors; evaluating a nonzero univariate entry at the
transcendental element $t$ cannot make it zero.  Constant
endpoints are covered directly.  Single-factor tests are
already preserved by the same argument and need no pairs.

Consequently I-28 applies with $w=3,r=2$, giving support
budget $g=2$ for the entire rank-two test, including its
endpoints.  Two fresh SV pairs simultaneously preserve all
rank-two tests: each selected minor is a nonzero polynomial
after substitution, and the product of the finitely many
selected minors is nonzero.  Previously preserved rank-three
minors remain nonzero upon setting the new amplitudes to zero.
The rank-one stage of I-21 still uses its independently
proved six-pair budget.  With disjoint fresh parameter pairs
the two stages therefore give
\[
 \mathcal G_i^{(8)}
      =t+\sum_{j=1}^{8}z_j\lambda_i(y_j),
 \qquad f\ne0\Longrightarrow f(\mathcal G^{(8)})\ne0
 \tag{I.40}
\]
for every width-at-most-three one-hidden-order ROABP in the
scope of I-21.

All degree and interpolation arguments of I-21 remain
unchanged except for replacing ten pairs by eight.  In its
notation, this yields the explicit sample bound
\[
 (D+1)\bigl((n-1)D+1\bigr)^8\binom{D+8}{8},
 \tag{I.41}
\]
over a working field containing at least
$\max\{n,D+1,(n-1)D+1\}$ distinct available elements.
The earlier ten-pair theorem remains valid; (I.40) is a
strictly smaller proved sampler.  This refinement does not
address the singular reset profiles at variable width and
does not claim a polynomial dependence on that width.


\subsection{I-29. A polynomial pair budget for arbitrarily many resets at one singular rank}
\label{I-29}

\paragraph{Statement and scope.}
The rank-profile mechanism used in I-21 has a growing-width
positive specialization.  Fix $1\le r<w$, put $c=w-r$, and
suppose a constant-boundary width-$w$ ROABP representation
of $f$ in one hidden variable order has the following promise:
every nonempty interval of its internal matrix word has generic
rank either $r$ or $w$.  The boundary row and column are not
included in this promise.  There may be arbitrarily many
rank-$r$ factors and intervals; their matrices need not commute.
Then the generator
\[
 X_i=t+\sum_{j=1}^{K}z_j\lambda_i(y_j),
 \qquad K=2r(w-r)+w-1,
 \tag{I.42}
\]
hits every nonzero such $f$ in every characteristic, over a
field or extension supplying the requisite distinct nodes.
For total degree at most $D$, the explicit grid construction
of I-21 has size at most
\[
 (D+1)\bigl((n-1)D+1\bigr)^K\binom{D+K}{K}.
 \tag{I.43}
\]
Its field-size convention is the same as I-21.

This is a polynomial \emph{pair budget} in $w$, with
sampling cost $(nD)^{O(w^2)}$, not polynomial-size PIT for
variable width.  The statement permits unbounded total
corank and so is different from I-20.  It further specializes
the cited unit-rank mechanism; no historical novelty is claimed.

\paragraph{1. A two-rank coefficient-space lemma.}
Let $B_i(Y_i)$ be $w\times w$ polynomial matrices in distinct
variables over a field $L$.  Assume every interval has generic
rank equal to its rank at zero, and that these ranks belong
to $\{r,w\}$.  Then, for every constant row $b$,
\[
 \begin{gathered}
 \operatorname{Coeff}(bB_1\cdots B_m)
 =\operatorname{span}_L\{
       [Y^\alpha](bB_1\cdots B_m):
       |\operatorname{supp}\alpha|\le h\},\\
 h=r(w-r)+w-1=rc+w-1.
 \end{gathered}
 \tag{I.44}
\]
The assertion remains valid in a truncated product local ring
when the corresponding interval unit-rank assumptions hold.

If the full constant product is invertible, all constant
factors are invertible, and I-17's nested row-space argument
gives the smaller bound $w-1$.  Otherwise its rank is $r$.
Choose a persistent $r$-dimensional row space $U_0$ on which
the full constant product is injective and transport it to
$U_i$ at every cut.  Choose a complement $V_m$ of $U_m$;
the suffix pullback
\[
 V_i=\{z:zB_{i+1}(0)\cdots B_m(0)\in V_m\}
\]
has dimension $c$, complements $U_i$, and is carried into
$V_{i+1}$ by the next constant matrix.  In compatible bases
the constants therefore have the form
\[
 B_i(0)=\operatorname{diag}(I_r,Q_i),
 \tag{I.45}
\]
where $Q_i$ is either invertible or zero: the rank of the
whole factor is either $w$ or $r$.  These basis changes
preserve supports and interval ranks.

Let $\pi:L^{1\times w}\to L^{1\times r}$ be active-coordinate
projection.  Prefix coefficient spaces $S_i$ have nested
projections $\pi S_i$, by (I.45).
We need the following stable-segment assertion: on a segment
where every projected space is the same $T$, give basis lifts
of $T$ initial cost zero and a basis of the initial projection
kernel cost one.  Cost means initial cost plus the number of
nonconstant local coefficients used.  Then cost at most $c$
spans the final coefficient space.

First suppose all $Q_i$ in the segment are invertible.
After constant transport, the coefficient spaces are nested
inside $\pi^{-1}(T)$, whose dimension is $\dim T+c$.
If the initial kernel has dimension $k>0$, all initial rows
cost at most one and there can be at most $c-k$ subsequent
dimension increases.  At each increase one nonconstant
coefficient suffices, so the cost is at most
$1+c-k\le c$.  If $k=0$, the initial cost is zero and
there are at most $c$ increases.  Stable layers require only
constant propagation.  This proves the assertion in this case.

If a zero $Q_i$ occurs, take its last position $a$ and put
$P=B_a\cdots B_b$, where $b$ is the segment's final position.
The suffix has rank $r$ and constant matrix
$\operatorname{diag}(I_r,0)$: its rank is at most $r$ from
the factor at $a$, and its constant $r$-minor is nonzero.
In the power-series ring at zero, or in the truncated local
ring, write its top blocks as $[U\ V]$.  The matrix $U$
is invertible.  The vanishing $(r+1)$-minors give the row-graph
identity
\[
 (x,y)\in\operatorname{rowmod}(P)
       \quad\Longrightarrow\quad y=xU^{-1}V.
\]
Let $B$ be the matrix of the initial lifts propagated through
the constants before position $a$, and let $H=\pi B$ be the
chosen basis of $T$.  Stable projection implies
\[
 \pi(BP)=AH,\qquad \pi(sP)=\eta H\quad(s\in S_{a-1}),
 \qquad A(0)=I.
\]
Thus $A$ is invertible and the row graph gives
$sP=\eta A^{-1}BP$.  If $T=0$, it instead gives $sP=0$.
It follows, by taking coefficient rows, that the final space
is exactly $\operatorname{Coeff}(BP)$.  Indeed, the displayed
identity proves one containment, and the rows of $B$ already
belong to $S_{a-1}$ for the other.

Restart after position $a$ from
$\operatorname{Coeff}(BB_a)$.  It has constant lifts of $T$
of cost zero and a kernel basis of cost at most one, since
only that one variable position was used.  The remaining
constants are invertible.  The preceding dimension argument
then bounds the total cost by $c$, proving the stable-segment
assertion.  In the untruncated case the inverses may be formal
power series; each coefficient of their product with $BP$
is still a finite linear combination of coefficient rows of
$BP$.  No analytic convergence is used.

We now count support from the original single row $b$.
If $\pi b\ne0$, there are at most $r-1$ strict projection
growths.  The initial stable segment costs at most $c$;
each growth factor and the following stable segment add at
most $1+c$.  The bound is
\[
 c+(r-1)(c+1)=rc+r-1\le h.
\]
If $\pi b=0$, the first stable segment lies in the passive
$c$-space.  If it contains a zero $Q_i$, the row-graph
argument kills the entire current coefficient space and the
final word is zero.  Otherwise its constants are invertible,
and starting from one row there are at most $c-1$ dimension
increases.  This first segment costs at most $c-1$.
At most $r$ projection growths follow, giving the bound
\[
 (c-1)+r(c+1)=rc+r+c-1=rc+w-1=h.
\]
At every stage the retained rows are linear combinations of
actual coefficient rows of the original prefix with the
stated support.  At a restart the chosen lifts are propagated
by constants through skipped positions; selecting a local
nonconstant coefficient adds only its own variable position.
Thus the cost count proves the actual support assertion
(I.44), rather than just a bound on an abstract generating
module.  The same reasoning applies in the truncated ring.

\paragraph{2. Sharpness for every active and passive dimension.}
Write the active coordinates as $a_1,\ldots,a_r$ and the
passive coordinates as $p_1,\ldots,p_c$, and put
$P=\operatorname{diag}(I_r,0_c)$.
Use a fresh variable in every factor below.
First apply the $c-1$ unipotent chain factors
\[
 I+Y_1E_{p_1,p_2},\ I+Y_2E_{p_2,p_3},\ \ldots,
 I+Y_{c-1}E_{p_{c-1},p_c}.
\]
This initial list is empty if $c=1$.
For each $j=1,\ldots,r$, apply one reset
\[
 P+S_jE_{p_c,a_j},
\]
followed by $c$ unipotent chain factors for the edges
$a_j\to p_1\to p_2\to\cdots\to p_c$.
The total number of factors is
$(c-1)+r(c+1)=h$.

Each reset has generic rank exactly $r$: its only extra
row is a multiple of the existing active row at $a_j$.
Each chain factor is unipotent and invertible.  Every
interval without a reset has constant matrix $I_w$ and
generic rank $w$.  Every interval with a reset has constant
matrix $P$, has generic rank at most $r$ from the reset
factor, and has rank at least $r$ from its constant active
minor.  Thus every interval satisfies precisely the
two-rank promise, including rank preservation at zero.

Starting with $e_{p_1}^T$, the coefficient in the last
passive coordinate after the initial chain is the product
of all its $c-1$ variables.  At reset $j$, all passive
coordinates are killed except that last coefficient,
which is transferred, multiplied by $S_j$, into the
fresh active coordinate $a_j$.  Earlier active coordinates
are unchanged.  The following chain multiplies that
coefficient by each of its $c$ fresh variables before
placing it in $p_c$.  The earlier active coordinates
do not enter this chain.  Induction on $j$ proves that
the final contraction with $e_{p_c}$ is exactly the
product of all $h$ variables.

Consequently every coefficient row of support less than
$h$ is annihilated by $e_{p_c}$, whereas the full-support
coefficient is not.  This proves sharpness of (I.44) for
every $1\le r<w$.  In particular the corank-one case has
$h=2(w-1)$, and balanced active/passive dimensions require
quadratic support.  The scalar witness is a monomial and
is hit by one generic dense direction; this is not a
lower bound on dense-generator pair count.

\paragraph{3. Uniform preservation and the sampler.}
Start at the common transcendental center $X_i=t$.
All rank-$w$ intervals retain rank: each of their univariate
factors has a nonzero determinant, and the product of those
determinants remains nonzero at $t$.

For every rank-$r$ interval, use I-21's finite local-rank-test
construction.  Minimality makes each test's interior chunk
have rank greater than $r$, and the two-rank promise makes
that rank $w$.  Its selected nonzero $r$-minor is therefore
a contraction of an invertible compound chunk with at most
two endpoint factors.  Its endpoint compound row is
decomposable, and both endpoint vectors are nonzero at the
common transcendental center, exactly as checked in I-28.
That result gives a support budget $g=r(w-r)=rc$ for the
whole test, including the endpoints.

Add $g$ fresh SV pairs.  Each selected nonzero minor stays
nonzero, and their finite product is nonzero, so all such
tests and hence all rank-$r$ intervals are simultaneously
preserved over the parameter field.  Rank-$w$ intervals
remain preserved because the new amplitudes may be set to
zero.  This is an existence argument over one common
parameter field and requires neither the hidden order nor
the selected minors as input to the generator.

At the resulting center, translate all original variables
by independent residual variables.  Every interval now
has its generic rank at zero, in $\{r,w\}$.  The original
scalar remains a nonzero polynomial, since translation
over that field is invertible.  Apply (I.44) to its starting
row and then its final column.  Some scalar monomial has
support at most $h$; otherwise the final column would
annihilate the spanning coefficient rows and the whole
polynomial would be zero.  Adding $h$ further fresh SV pairs
therefore makes the scalar nonzero.  The two additive maps
combine into $K=g+h=2rc+w-1$ pairs, proving (I.42).

For sampling, the degree in $t$ is at most $D$, each selector
degree is at most $(n-1)D$, and the total amplitude degree
is at most $D$.  The same product grid and amplitude Newton
simplex as in I-21 give (I.43).  All arguments and the
interpolation convention are characteristic independent.
\hfill$\square$

\paragraph{What remains outside this positive case.}
I-24's exponentially long support witnesses use several
different singular ranks, rather than only one fixed rank
and full rank.  In that setting the passive blocks in
(I.45) may themselves have intermediate nonzero singular
ranks.  The last-zero restart used above then does not
replace the whole passive block, and recursive rank losses
can accumulate.  I-29 identifies a precisely proved regime
where this recursion is unnecessary.  It neither asserts
that recursion is necessary for dense nonvanishing nor
excludes rational-function or cocharacter constructions
for the remaining profiles.


\subsection{I-30. Absorbing endpoint vectors into unit factors; a six-pair width-three sampler}
\label{I-30}

\paragraph{An elementary endpoint factorization.}
Let $u(Y)$ be a polynomial row of length $W$ over a field $L$,
with $u_0=u(0)\ne0$.  Choose $j$ with $(u_0)_j\ne0$ and set
\[
 B(Y)=I_W+\frac{e_j}{(u_0)_j}\bigl(u(Y)-u_0\bigr).
 \tag{I.46}
\]
Then
\[
 u(Y)=u_0B(Y),\qquad B(0)=I_W,\qquad
 \det B(Y)=\frac{u_j(Y)}{(u_0)_j}\ne0.
\]
The first two identities follow by multiplication; the third
is the rank-one determinant identity
$\det(I+ab^T)=1+b^Ta$.
For a polynomial column $v(Z)$ with $v_0=v(0)\ne0$, choosing
$j$ with $(v_0)_j\ne0$ likewise gives
\[
 \begin{gathered}
 C(Z)=I_W+\bigl(v(Z)-v_0\bigr)\frac{e_j^T}{(v_0)_j},
 \qquad v(Z)=C(Z)v_0,\\
 C(0)=I_W,
 \qquad \det C(Z)=\frac{v_j(Z)}{(v_0)_j}\ne0.
 \end{gathered}
 \tag{I.47}
\]
Thus both new endpoint factors are generically invertible and
invertible in every truncated local ring at zero.  No
characteristic restriction is used.

Suppose the interior word has generic interval ranks equal
to its constant interval ranks, in some set of allowed ranks
containing $W$.  Then the enlarged word
$B(Y)A_1\cdots A_m C(Z)$ has the same property.  An interval
lying in the old word is unchanged.  Every other interval is
an old prefix, suffix, or whole word, multiplied on one or
both sides by generically invertible matrices whose constants
are identities; both generic and constant ranks are preserved.
An interval consisting only of a new endpoint factor has
rank $W$.  The same argument holds for interval unit ranks
in a truncated local ring, because multiplication by an
invertible matrix preserves its determinantal rank ideals.

Consequently every coefficient-span support bound proved for
constant endpoints and this rank profile also applies to
\[
 u(Y)A_1\cdots A_m v(Z)
      =u_0B(Y)A_1\cdots A_m C(Z)v_0,
\]
with both endpoint positions included in the original bound.
This costs no additional support positions and is an
existence argument about the representation, not a request
to recover unknown endpoint coefficients from the black box.

\paragraph{Application to all width-three rank tests.}
Use the rank-descent construction of I-21, starting at
$X_i=t$.  Full-rank intervals are preserved as before.
In a selected rank-two test, its second-compound interior
has width three and invertible constant factors.  Its
nonzero polynomial endpoint row and column have nonzero
constants at this center.  Absorb them by (I.46)--(I.47).
The enlarged compound word has all constant factors
invertible, so I-17's row-space argument gives support at
most $3-1=2$ for a nonzero scalar coefficient.  Equivalently
one may use I-28.  Two fresh SV pairs preserve all rank-two
tests simultaneously, since the finite product of their
nonzero substituted minors is nonzero.

The new center is
\[
 \alpha_i=t+\sum_{j=1}^2z_j\lambda_i(y_j).
\]
Every $\alpha_i$ remains transcendental over the original
coefficient field: it has coefficient one on the fresh
indeterminate $t$.  Therefore every nonzero univariate
polynomial endpoint row or column from an original factor
is nonzero at its center.  Constant endpoint vectors are
covered directly.  A selected nonzero contraction cannot
have an identically zero endpoint vector.

In a rank-one test, the interior interval ranks are two or
three and have already been preserved at the current center.
Translate to independent residual variables over the field
of the current parameters.  Absorb both endpoint vectors
using (I.46)--(I.47); every interval of the enlarged word
still has generic and constant rank in $\{2,3\}$.
I-29 with $(w,r)=(3,2)$ now gives support
\[
 h=2(3-2)+3-1=4
\]
for the entire test, including its endpoints.  Four further
fresh SV pairs preserve all the rank-one tests simultaneously.
The preceding rank-two and rank-three minors stay nonzero
because the new amplitudes may be set to zero.
I-21's local-rank induction therefore preserves the rank-one
augmented full word, whose single possibly nonzero entry is
the original scalar $f$.

It follows that the thirteen-parameter map
\[
 \mathcal G_i^{(6)}(t,\mathbf y,\mathbf z)
      =t+\sum_{j=1}^{6}z_j\lambda_i(y_j)
 \tag{I.48}
\]
is a hitting generator for every width-at-most-three ROABP
in one hidden order, in every characteristic.  This is the
same class and evaluation model as I-21; it assumes no
two-rank promise on the whole original word.  The two-rank
lemma is used only inside the rank-one local tests, where
the promise follows from the rank-descent construction.

\paragraph{Exact evaluation bound.}
For total degree at most $D$, the degree in $t$ is at most
$D$, each selector degree is at most $(n-1)D$, and the total
amplitude degree is at most $D$.  Applying the grid and
Newton-simplex argument of I-21 gives the hitting-set bound
\[
 (D+1)\bigl((n-1)D+1\bigr)^6\binom{D+6}{6}.
 \tag{I.49}
\]
The working field must supply at least
$\max\{n,D+1,(n-1)D+1\}$ distinct elements, and evaluations
over that working extension are required if it differs
from the coefficient field.  Earlier ten- and eight-pair
theorems remain valid; (I.48) is the strongest sampler
proved in this track.  Nothing here proves polynomial
dependence on arbitrary width.
\hfill$\square$


\subsection{I-31. Full matrix cones give a pair budget controlled by the number of rank levels}
\label{I-31}

\paragraph{Attribution and statement.}
The full-matrix cone method in Liu (2026), Section~3.3,
printed pp.~7--8, avoids treating compound matrices as
arbitrary vectors.  We prove an original-coordinate version
with a somewhat looser block count, so no first-block
function-field factorization is needed.  This is a further
checked specialization of that explicitly unrefereed source,
not a claim of historical novelty or a certification of its
affine-orbit theorem.

Fix an allowed rank set
$\mathcal R=\{s_1<\cdots<s_L=w\}\subseteq\{1,\ldots,w\}$,
where $w\ge2$.  Put
\[
 d_w=w^2-1,\qquad
 d_s=(2ws-1)\left(2\max_{t\in\mathcal R,\ t>s}d_t+1\right)
       \quad(s\in\mathcal R,\ s<w),
 \qquad K_{\mathcal R}=\sum_{s\in\mathcal R}d_s.
 \tag{I.50}
\]
For constant-boundary width-$w$ ROABPs in one hidden order
whose nonempty internal intervals have generic ranks in
$\mathcal R$, the map $t\mathbf1+\operatorname{SV}_{K_{\mathcal R}}$
is a hitting generator in every characteristic.  It satisfies
\[
 K_{\mathcal R}\le2w^2(4w^2)^{L-1}.
 \tag{I.51}
\]
Thus a fixed number of allowed rank levels gives a polynomial
pair budget in $w$, even with arbitrarily many noncommuting
singular factors.  The sample count is still
\[
 (D+1)\bigl((n-1)D+1\bigr)^{K_{\mathcal R}}
                  \binom{D+K_{\mathcal R}}{K_{\mathcal R}},
 \tag{I.52}
\]
for degree bound $D$, with the working-field convention of
I-21.  For the full rank set $\{1,\ldots,w\}$ this is only
$K_{\mathcal R}=2^{O(w\log w)}$, so it proves polynomial size
for each fixed width, not polynomial dependence on variable
width.  I-29 and I-30 give substantially smaller bounds in
their respective special cases.

\paragraph{1. Cones and composition.}
Work over an algebraic closure of the coefficient field.
For a polynomial matrix $M$, write $\mathcal C(M)$ for the
Zariski closure of all scalar multiples of its values.
This is an irreducible affine cone.  Equality of such cones
preserves every homogeneous polynomial identity in the matrix
entries.  It also composes under multiplication on disjoint
variable sets: replacing a factor by a restriction with the
same cone leaves the cone of the whole product unchanged.
Indeed, taking a dense set in either input of the regular
matrix-multiplication map does not change the closure of its
image, and scalar multiples are absorbed by the cone.
These observations remain true after an extension of the
coefficient field.

\paragraph{2. Deleting stable rank-$s$ blocks.}
Let $B_1,\ldots,B_m$ be polynomial $w\times w$ matrices in
disjoint variable blocks, each of generic rank $s$ and with
$B_i(0)=P_s=\operatorname{diag}(I_s,0)$.
For a constant $a\times w$ matrix $R$ with $\pi R\ne0$,
where $\pi$ selects the first $s$ columns, put
\[
 M_i=RB_1\cdots B_i,\qquad
 V_i=\mathcal C(\pi M_i),\qquad V_0=\mathcal C(\pi R).
\]
Then $V_{i-1}\subseteq V_i$, since the zero value of $B_i$
preserves the first $s$ columns.  The cones are irreducible,
start at a line, and lie in the $as$-dimensional matrix space.
There are at most $as-1$ strict dimension increases.

Suppose $V_{j-1}=V_j=V_{j+1}=V$.  We claim that setting
block $j$ to zero preserves the full cone of $M_{j+1}$.
Fix a value of $B_{j+1}$ on the dense open where its top-left
$s\times s$ block $E$ is invertible, and call its top-right
block $F$.  Rank $s$ gives the row-graph identity: any matrix
obtained by multiplying $B_{j+1}$ on the left is determined
by its first $s$ columns $Q$, and equals $Q[I_s,E^{-1}F]$.
After the restriction $B_j=P_s$, those first columns have
cone $VE$.  This is contained in $V$, since the restricted
values are among the original values contributing to
$V_{j+1}=V$.  But $VE$ is closed and irreducible, has
dimension $\dim V$, and is contained in the irreducible $V$;
therefore $VE=V$.

Every original output at this fixed value of $B_{j+1}$ has
its first columns in $V$, and hence belongs to the restricted
full cone by the row-graph identity.  Taking the closure over
the dense set of such values of $B_{j+1}$ proves the claim.
The opposite containment is automatic from restriction.
It also preserves all subsequent cones by composition.

Retain every strict-growth block and the last block of each
run of nongrowth blocks.  All other blocks can be deleted
one at a time using the claim, while the projected cones
and the full cone at the end of each affected pair remain
unchanged.  The number retained is at most
$2(as-1)+1=2as-1$.  In particular, with $R=I_w$, at most
$2ws-1$ blocks preserve the full product cone.

\paragraph{3. Coordinate compression of a preserved-rank word.}
We prove that if every interval of a univariate matrix word
has its generic rank at zero, all these ranks belong to
$\mathcal R$, and the full word has rank $s$, then setting
all but at most $d_s$ positions to zero preserves its full
matrix cone.

For $s=w$, all constant factors are invertible.  Gauge them
to identities.  The prefix cones are nested, starting with
the line through $I_w$ in dimension $w^2$.  Retain only the
strict-growth positions.  At a nongrowth position, replacing
that factor by its identity constant preserves the prefix
cone and then all later cones by composition.  At most
$w^2-1=d_w$ positions remain.

For $s<w$, greedily split the word at successive shortest
prefixes of rank $s$.  Write such a block as $C_jA_{e_j}$;
the preceding chunk $C_j$ is empty or has rank greater
than $s$, as do all of its subintervals.  When the remaining
tail no longer has rank $s$, absorb it into the last block;
that tail is empty or has rank greater than $s$ as well.
Every resulting block has generic and constant rank $s$.
Every product of consecutive blocks has constant rank $s$,
because it is an original interval and cannot have rank
below that of the full word.

Compatible constant bases at the block cuts normalize all
block constants to $P_s$.  To see this, transport an active
$s$-space injectively along the constant products.  At an
internal cut the left kernel of the next block complements
this active space, since the adjacent product retains rank
$s$.  Choosing those passive kernel bases makes each block
constant $\operatorname{diag}(I_s,0)$.
These changes multiply the entire word by fixed invertible
matrices at its ends and do not affect coordinate supports
or cone-preserving restrictions.

Part~2 retains at most $2ws-1$ blocks.  In each retained
block, the induction hypothesis compresses the higher-rank
prefix chunk $C_j$ to at most
$D_s=\max_{t>s,\ t\in\mathcal R}d_t$ positions; the final
tail, if present, costs at most another $D_s$.  Retain the
one factor $A_{e_j}$ between them.  By disjoint-block
composition these restrictions preserve the block cone and
then the full word cone.  Thus at most
$(2ws-1)(2D_s+1)=d_s$ positions suffice, completing the
induction.

\paragraph{4. Absorbing rectangular endpoints in the ground space.}
For an $r\times w$ polynomial matrix $L(Y)$ whose constant
$L_0$ has full row rank $r$, choose a constant right inverse
$R_0$ with $L_0R_0=I_r$.  Then
\[
 B(Y)=I_w+R_0(L(Y)-L_0),\qquad
 L(Y)=L_0B(Y),\qquad B(0)=I_w.
 \tag{I.53}
\]
The determinant of $B$ has constant term one, so $B$ is
generically invertible.  Similarly a $w\times r$ endpoint
$V(Z)$ with full-column-rank constant $V_0$ can be written
$V(Z)=C(Z)V_0$, where $C(0)=I_w$ and $C$ is generically
invertible, by choosing a constant left inverse of $V_0$.
Thus these endpoint factors can be absorbed in the original
$w$-dimensional matrix word, without passing to a compound
space.  Generic and constant interval ranks are unchanged,
just as in I-30.

\paragraph{5. Rank descent and final scalar detection.}
Start with the common center $X_i=t$, which preserves every
full-rank interval.  Process the remaining ranks in
$\mathcal R$ in decreasing order.  For a rank-$r$ local test
as in I-21, the interior chunk has rank greater than $r$,
and every one of its intervals is already preserved.
A selected nonzero $r$-minor of the test has the form
\[
 \det\bigl(L(Y_0)M(\mathbf Y)V(Y_{m+1})\bigr).
\]
The selected endpoint matrices have full row and column
rank $r$ at the current center.  Indeed, a nonzero selected
minor implies that each endpoint has some nonzero
univariate $r$-minor, and each current coordinate remains
individually transcendental over the original coefficient
field: it is $t$ plus terms independent of $t$.

Absorb these endpoints by (I.53).  The enlarged interior
word still has precisely the preserved higher-rank profile.
If its full rank is $s>r$, Part~3 gives a restriction to at
most $d_s\le D_r$ positions preserving its full matrix cone.
The fixed-endpoint determinant is homogeneous of degree
$r$ in the full matrix entries, so this restriction keeps
the selected minor nonzero.  It therefore has a monomial
of support at most $D_r$ and is hit by that many fresh SV
pairs.  All chosen test minors remain simultaneously nonzero
because their product is a nonzero polynomial.  Earlier
ranks remain preserved by setting the new amplitudes to zero.
The local-rank induction proves preservation of all internal
intervals at this rank.

After the smallest allowed rank is processed, every internal
interval is preserved.  Part~3 applied to the full word gives
a cone-preserving restriction on at most $d_s$ positions,
where $s$ is its generic rank.  The nonzero scalar output
is a homogeneous linear functional on the matrix entries,
so it remains nonzero on that restriction.  At most
$\max_{s\in\mathcal R}d_s=d_{s_1}$ additional SV pairs hit it.

The numbers $d_s$ strictly increase as the allowed rank
decreases.  At rank $s_i$ the preservation cost is therefore
$d_{s_{i+1}}$.  Summing those costs and the last scalar
detection cost gives exactly $K_{\mathcal R}$ in (I.50).
All pairs are fresh and their additive maps combine into
one SV map.  The map depends only on the bounds and the
chosen nodes, not on the unknown order or matrices.

Finally $d_s+1\le4w^2(D_s+1)$ and $d_w+1=w^2$.
The recurrence also gives $d_s>2D_s$ for $s<w$.
Hence the sum is at most twice its largest term, yielding
(I.51).  The degree and Newton-simplex sampling argument is
identical to I-21, proving (I.52).
\hfill$\square$

\paragraph{Remaining main gate.}
The proof replaces the exponential compound dimension by
matrix-cone dimensions, but a new block-selection factor
appears at each distinct singular rank level.  For bounded
$L$ this is polynomial in $w$ at the pair-count level;
for unrestricted $L$ the bound is still exponential in
$w\log w$.  Moreover even a polynomial pair budget would
still need an efficient sampling argument to yield a
polynomial-size hitting set.  No implication from Hoza's
path isolation to either of these missing steps has been
proved.  The dense lower bound I-26, the filtered-mask
obstruction I-27, and the known logarithmic every-order
result I-B08 continue to delimit the candidate repairs.


\subsection{I-32. Exponential coordinate-cone compression is necessary even at the common center}
\label{I-32}

\paragraph{Statement.}
For every $w\geq3$ and every coefficient field, put
\[
 h=3\cdot2^{w-2}-2,\qquad n=2h.
\]
There is an explicit width-$w$ word of $n$ univariate linear
matrices in distinct variables $x_1,y_1,\ldots,x_h,y_h$ with
the following properties. Image cones use the algebraic-closure
convention of I-31.

Every nonempty internal interval has generic rank in
$\{2,\ldots,w\}$, and that rank is preserved at every common
center $x_i=y_i=\tau$. In particular it is preserved at the
transcendental center $t\mathbf1$ used in I-31. Nevertheless,
after translation to residual variables at that center, a
restriction obtained by setting unselected coordinates to zero
preserves the full matrix image cone if and only if it retains
at least one coordinate from each pair $(x_i,y_i)$. The minimum
number of retained positions is therefore exactly $h$.

Thus a uniform polynomial-in-$w$ coordinate-cone compression
bound is false, even at a common transcendental center where
all internal interval ranks are preserved. In terms of the
number $L=w-1$ of allowed rank levels, the example needs
$3\cdot2^{L-1}-2$ positions. It precludes replacing the rank-level
recursion by a polynomial bound for this exact compression
criterion. It does not preclude changing the center or using
other dense maps: the distinguished nonzero scalar below is
detected by one centered SV pair.

\paragraph{1. The seed word.}
We recall the explicit matrices of I-24, including the property
needed here. Set $P_r=\operatorname{diag}(I_r,0_{w-r})$ and
\[
 Q_{r,a}(s,u)=(P_r+sE_{w,a})(P_r+uE_{a,w})
       \quad(2\leq r<w,\ a\leq r).
\]
With fresh variables at every occurrence, define
\[
 H_w=\text{empty word},\qquad
 H_k=H_{k+1}Q_{k,k}H_{k+1}\quad(3\leq k<w),
\]
and
\[
 W=H_3Q_{2,1}H_3Q_{2,2}H_3.
\]
If $L_k$ is the length of $H_k$, then
$L_w=0$, $L_k=2L_{k+1}+2$, and
$3L_3+4=h$. Label the successive variables of this word by
$Y_1,\ldots,Y_h$. Each factor is of one of the two forms
\[
 P_r+Y_iE_{w,a},\qquad P_r+Y_iE_{a,w}.
 \tag{I.54}
\]
They have generic rank $r$ and constant matrix $P_r$.

For completeness, if a row $q$ has $q_a=0$, then
\[
 qQ_{r,a}(s,u)
   =\sum_{j\leq r}q_j e_j^T
       +q_ws e_a^T+q_wsu e_w^T.
\]
Starting with coordinates $k,\ldots,w-1$ zero, $H_k$ leaves
the earlier coordinates unchanged and multiplies the last
coordinate by all its variables. This follows by induction:
the first $H_{k+1}$ leaves coordinate $k$ zero, $Q_{k,k}$
performs the displayed transfer and resets all higher active
coordinates, and the second $H_{k+1}$ applies again. In the
full word starting at $e_w^T$, the two rank-two gadgets use
successively the still-zero coordinates one and two and reset
the coordinates required by the intervening $H_3$ words.
Consequently
\[
 e_w^T W(Y_1,\ldots,Y_h)e_w=\prod_{i=1}^hY_i.
 \tag{I.55}
\]
This calculation is valid in every characteristic.

\paragraph{2. A paired lift with preserved common-center ranks.}
Replace every factor in (I.54) by two factors on fresh variables.
For $E=E_{w,a}$, use
\[
 (I_w+x_iE)\bigl((I_w-y_iE)P_r\bigr),
 \tag{I.56}
\]
and for $E=E_{a,w}$, use
\[
 \bigl(P_r(I_w+x_iE)\bigr)(I_w-y_iE).
 \tag{I.57}
\]
Since $E^2=0$, all matrices $I_w+zE$ are invertible with
inverse $I_w-zE$. In (I.56), $EP_r=E$; in (I.57),
$P_rE=E$. Each pair therefore equals
$P_r+(x_i-y_i)E$. The full lifted word $\widetilde W$ obeys
the matrix identity
\[
 \widetilde W(x_1,y_1,\ldots,x_h,y_h)
       =W(x_1-y_1,\ldots,x_h-y_h).
 \tag{I.58}
\]
Every factor is still univariate and linear. Exactly one half
of each pair is invertible; the other is an invertible matrix
times $P_r$, on the appropriate side, and has rank exactly $r$
for every parameter value.

We verify every partial interval, not only complete pairs.
At $x_i=y_i=\tau$, each complete pair is $P_r$.
A partial initial pair can only contribute its second half.
For (I.56) that half is $(I_w-\tau E)P_r$; for (I.57) it is
$I_w-\tau E$. A partial final pair can only contribute its
first half. For (I.56) it is $I_w+\tau E$; for (I.57) it is
$P_r(I_w+\tau E)$. All intervening pairs are complete.
Thus the constant interval product has the form
\[
 U_{\rm left}\left(\prod P_{r_i}\right)U_{\rm right},
 \tag{I.59}
\]
where the two boundary matrices are invertible, and the
projections are exactly those attached to singular halves
lying in the interval. The same description handles an
interval contained in one pair: a single half has its
displayed rank, and the complete pair is $P_r$. If no singular
half occurs, the interval is a product of invertible matrices.

The nested diagonal projections in (I.59) have product
$P_q$, where $q$ is the minimum rank of the singular halves
present; take $q=w$ if there are none. The common-center rank
is therefore exactly $q$. The generic interval rank is at
most $q$, because every singular half has that fixed rank
and products cannot increase it. It is at least $q$ because
the specialization (I.59) has rank $q$, so some $q$-minor is
a nonzero polynomial before specialization. These two bounds
prove equality. This argument applies to every $\tau$, in
particular to $\tau=t$ over the coefficient field extended
by a transcendental element.

\paragraph{3. The exact coordinate-compression threshold.}
Write $x_i=t+X_i$, $y_i=t+Y_i$. Formula (I.58) gives
\[
 \widetilde W(t+X_1,t+Y_1,\ldots,t+X_h,t+Y_h)
     =W(X_1-Y_1,\ldots,X_h-Y_h),
\]
and (I.55) becomes
\[
 e_w^T\widetilde W(t+\mathbf X,t+\mathbf Y)e_w
       =\prod_{i=1}^h(X_i-Y_i)\ne0.
 \tag{I.60}
\]
Let a coordinate restriction retain a set of residual positions
and set all other residual coordinates to zero. If it omits
both positions of any pair, the displayed scalar is identically
zero after restriction. The restricted matrix image cone is
then contained in the homogeneous hyperplane
$\{M:e_w^TMe_w=0\}$. The original cone is not contained in
that hyperplane by (I.60). Hence the two cones differ.

Conversely, if at least one position in every pair is retained,
the remaining differences $X_i-Y_i$ can still be assigned
independently and arbitrarily. Formula (I.58) shows that the
restricted full matrix image equals the original full matrix
image, not merely its cone closure. This proves the exact
criterion and the minimum $h$: retaining $X_1,\ldots,X_h$
alone attains it.

The counterexample lies in the actual final rank-one test
class. Append $L=e_1e_w^T$ and $R=e_we_1^T$ to the word.
Its full product is (I.60) times $e_1e_1^T$ and has generic
rank one. Every prefix beginning at $L$ is nonzero, since
otherwise the full product would vanish, and therefore has
rank one. Every interval wholly inside $\widetilde W$ has
rank at least two. Thus the shortest rank-one suffix of the
prefix just before $R$ begins at $L$, and adjoining $R$
gives the full word as a local final test. Its entire internal
higher-rank profile is already preserved at the common center.

\paragraph{4. One dense pair still detects the scalar.}
Choose distinct interpolation nodes for the $2h$ original
variables, and denote their Lagrange polynomials by
$\lambda_{x_i}$ and $\lambda_{y_i}$. The common-center
one-pair substitution gives
\[
 \prod_{i=1}^h(x_i-y_i)
   \longmapsto
 z^h\prod_{i=1}^h
       \bigl(\lambda_{x_i}(y)-\lambda_{y_i}(y)\bigr)\ne0.
 \tag{I.61}
\]
Each difference is nonzero because its value at the node
for $x_i$ is one. Their product is nonzero in the parameter
polynomial ring over every working field with these distinct
nodes. Numerical evaluation uses an adequate grid or extension,
as in I-21; no assertion is made that every node value detects
this product.

This separates the refuted exact coordinate-deletion criterion
from dense nonvanishing. In particular it rules out a polynomial
dimension charge for uniformly preserving the full matrix cone
by retaining original variable positions at every preserved
common center. It does not rule out a smaller bound after a
different dense shift, a weaker sufficient invariant than full
cone equality, or a more efficient sampling argument. The
constant-center difference cancellation is explicit, not an
assertion about every possible preserved center. $\square$

\paragraph{Derivation record.}
The paired lift and the common-center cone obstruction were
derived by the purity track from the independently proved
I-24 matrices. The result is a deduction within this report;
no literature-priority claim is made.


\subsection{I-33. Efficient prime cocharacters for two dense products, with a product-closure repair}
\label{I-33}

\paragraph{Scope and relation to known PIT.}
This construction targets efficient sampling rather than another
support bound.  Let $n\ge1$ and let $\delta\ge1$ be an integer;
the case $\delta=0$ consists only of constants.  Consider
\[
 f(\mathbf X)=\prod_{i=1}^n F_i(X_i)-\prod_{i=1}^n G_i(X_i),
 \qquad \deg F_i,\deg G_i\le\delta.
 \tag{I.62}
\]
It is a width-two diagonal ROABP, but its expanded support
may be exponential in $n$.  Polynomial PIT for this class
is already known, for example from I-B06.  Bisht, Gupta,
and Volkovich (2023), Theorem~1 (printed p.~9:4) and
Theorem~48 (printed p.~9:20), give a broader positive result
for sums of two products of read-once formulas.  We claim
no new PIT complexity for that established class and no
historical priority for the particular construction below.
The point is an explicit small family of compatible
integer cocharacters, a proof using divisors rather than
expanded monomials, and an exact test of its product closure.

Set
\[
 H=2\delta(n-1),\qquad M=H+2n+2,\qquad B=16M^2.
\]
Choose any $2n$ distinct integer primes
$p_1,\ldots,p_n,q_1,\ldots,q_n$ in $(H,B]$.
For $j=0,1,\ldots,n$, define positive integer weights
\[
 a_i^{(0)}=p_i,\qquad
 a_i^{(j)}=\begin{cases}q_j&i=j,\\p_i&i\ne j.\end{cases}
 \tag{I.63}
\]
These primes exist and can be found in polynomial time in
$n,\delta$ by an elementary sieve, as proved below.
Over \emph{every} field, every nonzero $f$ in (I.62) satisfies
\[
 f(T^{a_1^{(j)}},\ldots,T^{a_n^{(j)}})\ne0
 \quad\hbox{for at least one }0\le j\le n.
 \tag{I.64}
\]
No prime needs to be excluded when it equals the characteristic.
Each substituted polynomial has degree at most $n\delta B$.
Hence evaluating all $n+1$ curves on any $n\delta B+1$
distinct scalars gives a hitting set of size at most
\[
 S=(n+1)(n\delta B+1),
 \tag{I.65}
\]
with polynomial integer weights and polynomial query count.
As elsewhere, a working extension is required when the
coefficient field has too few scalar values.

\paragraph{1. The divisor separation lemma.}
Let $R_i(X)=F_i(X)/G_i(X)$ be nonzero rational functions,
with numerator and denominator degrees at most $\delta$,
and let $p_i$ be distinct primes greater than $H$.
If
\[
 \prod_i R_i(T^{p_i})\in\overline F^{\times},
 \tag{I.66}
\]
then each $R_i(X)$ is $c_iX^{b_i}$ for some nonzero
constant $c_i$ and integer $|b_i|\le\delta$.
Work over an algebraic closure and cancel numerator and
denominator common factors first.

If the characteristic equals one of the primes, say $p_k$,
then a nonzero root or pole of $R_k(T^{p_k})$ has nonzero
integer multiplicity of absolute value at least $p_k$.
At the same point, each other $R_i(T^{p_i})$ has multiplicity
of absolute value at most $\delta$, since the other power
maps are unramified away from zero.  Their combined
multiplicity has absolute value at most $(n-1)\delta<p_k$.
It therefore cannot cancel the first multiplicity, contrary
to (I.66).  Thus $R_k$ has no nonzero root or pole and is
a constant times a monomial.  There is at most one such
characteristic prime.

For any remaining $i$, a nonzero root or pole of $R_i$
pulls back to an orbit $O=\alpha\mu_{p_i}$ of $p_i$
distinct points.  Its multiplicity is nonzero at every
point of $O$.  For $j\ne i$ with $p_j$ different from the
characteristic, the nonzero divisor of $R_j(T^{p_j})$
is supported on at most $2\delta$ cosets of $\mu_{p_j}$.
Each such coset intersects $O$ in at most one point, since
$\gcd(p_i,p_j)=1$.  A characteristic-prime factor, if
present, was already shown to have no nonzero divisor.
Consequently all other factors together can cancel the
divisor on at most $2\delta(n-1)=H<p_i$ points of $O$.
This contradicts (I.66).

Every $R_i$ therefore has divisor supported only at zero
and infinity.  Factoring its numerator and denominator
over the algebraic closure gives $R_i(X)=c_iX^{b_i}$.
The degree bounds give $|b_i|\le\delta$.  This proves
the lemma, including the inseparable characteristic case.

\paragraph{2. The family separates the remaining monomial relation.}
If one product in (I.62) is zero and the other is nonzero,
every positive-weight curve keeps the nonzero product
nonzero, so (I.64) is immediate.  Otherwise both products
are nonzero and we may divide by the second.
If the base curve $j=0$ kills $f$, the divisor lemma gives
$R_i(X_i)=c_iX_i^{b_i}$, and the resulting Laurent identity
forces
\[
 \prod_i c_i=1,\qquad\sum_i b_ip_i=0.
\]
If every $b_i=0$, then $f=0$, contrary to the hypothesis.
Choose $j$ with $b_j\ne0$.  Changing only its prime changes
the exponent sum by the nonzero integer
$b_j(q_j-p_j)$.  Thus the $j$th curve does not kill $f$.
This is an equality of exponents of the formal indeterminate
$T$, so it is valid even when that integer is divisible by
the characteristic.  This proves (I.64).

The same argument shows that if the two component products
are linearly independent, at least one curve preserves
their two-dimensional coefficient span.  If their ratio is
not a monomial product, the base-curve ratio cannot be
constant by the divisor lemma.  If it is a nonconstant
monomial product, some weight vector has nonzero exponent
sum.  In either case the two univariate images are linearly
independent.  This is a concrete instance of the compression
condition I-14; no basis-isolating-weight claim is needed.

The weight family has full affine span: its differences
from the base vector are $(q_j-p_j)e_j$.  Encoding
$j\in\{0,\ldots,n\}$ by $\lceil\log_2(n+1)\rceil$ bits
gives an explicit nonlinear decoder using coordinate
indicator functions.  It is therefore outside the bounded
affine and fixed-degree decoder obstruction of I-15.
It is not obtained from Hoza's path-isolation theorem.

\paragraph{3. An elementary polynomial prime bound.}
For even $B$, the central binomial coefficient satisfies
\[
 \binom{B}{B/2}\ge\frac{2^B}{B+1}.
\]
For each prime $p\le B$, its valuation in this binomial
coefficient is at most $\lfloor\log_p B\rfloor$: each
term in the factorial-valuation difference is zero or one.
Hence its prime-power contribution is at most $B$, and
\[
 B^{\pi(B)}\ge\frac{2^B}{B+1},\qquad
 \pi(B)\ge\frac{B\log2-\log(B+1)}{\log B}.
\]
For $B=16M^2$ and $M\ge2$, the elementary inequalities
$\log(B+1)\le(B/2)\log2$ and $\log B\le4M\log2$
give $\pi(B)\ge2M$.  At most $H$ primes are at most $H$,
so more than $2n$ primes lie in $(H,B]$.
Sieve this polynomially bounded interval and take its first
$2n$ primes.  This avoids using an unproved prime-distribution
hypothesis or a bit-length-only prime bound.

\paragraph{4. The raw curve family is not closed under products.}
The preceding success cannot be extended to general width
two just by multiplying separate successful factors.
For the explicit construction with $n=64$ and $\delta=1$,
one has $B=16\cdot256^2$.  Split the variables into two
blocks of size $32$.  Within either block the $2^{32}$
subset sums of the base primes take at most $32B+1$
integer values, and $2^{32}>32B+1$.  Two different subsets
therefore have the same sum.  After canceling their common
indices, this gives a nonzero multilinear binomial on that
block which vanishes along the base-prime curve.
Let the two disjoint-block binomials be $A$ and $C$.
Their product $AC$ is nonzero and multilinear.

Changing one coordinate's prime affects at most one of
$A,C$; the other remains zero along that curve.  Thus
$AC$ vanishes along all $n+1$ curves in (I.63).
It is nevertheless a width-two ROABP in a suitable order:
each binomial is the difference of two monomial products,
computed by diagonal $2\times2$ factors, and their two
scalar computations concatenate through a constant
rank-one bridge.  Absorbing that bridge into an adjacent
univariate layer gives an ordinary width-two ROABP.
This is a counterexample to the specified raw family,
not to the known width-two hitting-set theorem I-B06.

\paragraph{5. A concrete product-closure repair.}
Let $\mathcal H$ be the list of at most $S$ points from
(I.65).  Interpolate a polynomial curve
$\Gamma(Z)\in E[Z]^n$ of coordinate degree at most $S-1$
through that list at $S$ distinct nodes, repeating points
if necessary.  Every nonzero factor of the form (I.62)
stays nonzero on $\Gamma$, because some listed point hits it.
Therefore any nonzero product of such factors also stays
nonzero, even when different factors share variables.
For a product of total degree at most $D\ge1$, its
composition has degree at most $D(S-1)$ and is hit on any
$D(S-1)+1$ distinct scalars.  A working field with at least
that many elements supplies both the interpolation nodes
and this final grid.

This is the scalar product-closure mechanism of AGKS,
Lemma~33, recorded in I-B06, applied to the new explicit
prime-curve family.  It repairs the precise four-term
counterexample above and gives a fully polynomial sampler
for products of the dense two-product factors.  It does
not handle arbitrary sums of three or more products,
general noncommuting matrix blocks, or their coefficient
matroids.  Those cancellation patterns remain outside
the divisor-separation argument.
\hfill$\square$


\subsection{I-34. Prime cocharacters with a bounded local nullity}
\label{I-34}

\paragraph{Statement and scope.}
Let $n\ge1$, let $\delta\ge1$ be an integer, and consider
\[
 f(\mathbf X)=\sum_{j=1}^{w}\prod_{i=1}^{n}F_{ij}(X_i),
 \qquad 0\ne F_{ij}\in F[X_i],\quad \deg F_{ij}\le\delta.
 \tag{I.67}
\]
Zero product terms can first be omitted. Scalars and signs can be
absorbed into the factors. Let $A_i$ be the $(\delta+1)\times w$
matrix whose $j$th column is the coefficient vector of $F_{ij}$.
We call $w-\operatorname{rank}A_i$ its local nullity.

There is an explicit family of polynomial-magnitude positive
integer cocharacters with the following properties, over every
characteristic.

If some $A_i$ has local nullity at most $c$, the family changing
at most $k=\min\{n,\max\{1,c\}\}$ coordinates of a base prime
weight vector hits every nonzero $f$ in (I.67). For fixed $c$
this gives polynomial sampling size even when $w$ grows. The
promise itself requires $w-c\le\delta+1$; the assertion does not
silently treat $c$ as constant for arbitrary representations.

Without a local-nullity promise, sums of at most three separated
products require only one changed coordinate. More generally,
sums of $w\ge3$ products are hit using at most
$\min\{n,w-2\}$ changed coordinates. The last bound has an
exponent depending on $w$ and is not a polynomial-in-width PIT
result. Constant-top-fan-in PIT is already known, for example
from I-B08. The purpose here is to test and repair the specific
prime-cocharacter construction, not to claim a new general
complexity bound for an established constant-width class.

One-coordinate variation cannot extend this result to all sums
of four products: an explicit tensor-rank-four multilinear
example below vanishes on every entire one-coordinate cylinder
around an admissible prime curve. Thus increasing the number of
alternate primes while still varying only one coordinate does
not remove that obstruction.

\paragraph{1. The explicit family and its sampling cost.}
Set
\[
 H=2\delta(n-1),\qquad
 M=H+n(\delta+1)+2,\qquad B=16M^2.
\]
Choose $n(\delta+1)$ distinct primes in $(H,B]$, organized as
$p_i$ and $q_{i,a}$ for $i\in[n]$, $a\in[\delta]$.
The elementary prime count in I-33 applies verbatim: it gives
$\pi(B)\ge2M$, and at most $H$ primes lie below the threshold.
All these primes can be found by a polynomial-time sieve.

Let $\mathcal W_k$ contain the base vector $(p_1,\ldots,p_n)$
and every vector obtained by replacing $p_i$ by one of
$q_{i,1},\ldots,q_{i,\delta}$ on at most $k$ coordinates.
Then
\[
 |\mathcal W_k|=
 \sum_{a=0}^{k}\binom na\delta^a.
 \tag{I.68}
\]
Each curve $X_i=T^{\omega_i}$ has substituted degree at most
$n\delta B$. A scalar grid of $n\delta B+1$ distinct values
therefore gives a hitting set of size at most
\[
 (n\delta B+1)\sum_{a=0}^{k}\binom na\delta^a.
 \tag{I.69}
\]
An adequate working extension is understood if the coefficient
field is too small. These are degree bounds in the formal
indeterminate $T$, so no characteristic restriction on the
integer weights is needed.

The reason to use $\delta$ alternate primes, rather than just
one, is exact interpolation over $F(T)$. For $S\subseteq[n]$
write
\[
 f_S(\mathbf Z_S,T)=
 f\bigl((Z_i)_{i\in S},(T^{p_i})_{i\notin S}\bigr).
 \tag{I.70}
\]
This is the coordinate cylinder with the coordinates in $S$
free. It has individual degree at most $\delta$ in each $Z_i$.
The $\delta+1$ elements
$T^{p_i},T^{q_{i,1}},\ldots,T^{q_{i,\delta}}$ are distinct
in $F(T)$. Consequently, if every curve from $\mathcal W_k$
kills $f$, tensor-product interpolation over $F(T)$ shows that
$f_S=0$ for every $|S|\le k$. Conversely, if all these cylinders
are zero, every such curve is zero. This equivalence uses only
formal roots and is valid in all characteristics.

\paragraph{2. Nullity one: the circuit argument.}
We prove that if one $A_i$ has nullity at most one, then a
nonzero $f$ has a nonzero one-coordinate cylinder. Suppose
instead that every such cylinder is zero. Put
\[
 B_j(T)=\prod_{\ell\ne i}F_{\ell j}(T^{p_\ell}).
\]
Every $B_j$ is nonzero. The identity $f_{\{i\}}=0$ says
$A_i(B_1,\ldots,B_w)^T=0$ over $F(T)$.
If $A_i$ has full column rank this is impossible. Otherwise
its kernel is one-dimensional and defined over $F$, so
\[
 B_j(T)=h(T)c_j\quad(1\le j\le w),
 \qquad h\ne0,\quad c_j\in F^\times.
 \tag{I.71}
\]
The nonvanishing of every $c_j$ follows from that of every
$B_j$. Thus the columns of $A_i$ form a circuit: their unique
relation, up to a scalar, has full support.

For each $j$, the ratio $B_j/B_1$ is constant. Apply the
divisor separation lemma I-33 to the $n-1$ remaining variables.
Its threshold is no larger than our $H$; empty products when
$n=1$ cause no problem. We obtain
\[
 \frac{F_{\ell j}(X)}{F_{\ell1}(X)}
       =d_{\ell j}X^{b_{\ell j}}
 \quad(\ell\ne i),
 \qquad d_{\ell j}\in F^\times,
 \quad b_{\ell j}\in\mathbb Z.
 \tag{I.72}
\]
The constants lie in $F$: the ratio has coefficients in $F$
and equals a constant times a monomial. Negative exponents
are allowed in this rational-function identity. Set
$b_{\ell1}=0$.

Fix $\ell\ne i$ and use the zero cylinder with $X_\ell=Z$.
Dividing by nonzero rational functions and using (I.71)--(I.72)
reduces it to the Laurent identity
\[
 \sum_{j=1}^w c_j F_{ij}(T^{p_i})
       (Z/T^{p_\ell})^{b_{\ell j}}=0.
 \tag{I.73}
\]
For every nonempty group $J_b=\{j:b_{\ell j}=b\}$, its
coefficient gives
$\sum_{j\in J_b}c_jF_{ij}(T^{p_i})=0$.
The substitution $X\mapsto T^{p_i}$ is injective on $F[X]$,
so $\sum_{j\in J_b}c_jF_{ij}(X)=0$.
The vector equal to $c_j$ on $J_b$ and zero elsewhere is
therefore in $\ker A_i$. Since that kernel is generated by
the full-support vector $c$, every nonempty $J_b$ must be
all of $[w]$. In particular all $b_{\ell j}$ equal
$b_{\ell1}=0$.

It follows that every off-$i$ factor is a constant multiple
of its first-column factor. Equation (I.71) then implies that
$f$ is the product of these common off-$i$ factors with the
zero polynomial $\sum_jc_jF_{ij}(X_i)$, up to a common
nonzero scalar. Thus $f=0$, a contradiction. This proves
the nullity-one assertion. The proof includes the case where
a proposed one-dimensional kernel has a zero coordinate:
then (I.71) cannot occur, and the $i$th cylinder is already
nonzero.

\paragraph{3. Three terms and the general local-nullity reduction.}
For three nonzero product terms, either some $A_i$ has rank
at least two, when the preceding argument applies, or every
$A_i$ has rank one. In the second case every factor list is
proportional, and $f$ is a scalar multiple of one separated
product. If nonzero, it survives the base curve. The cases
of one or two terms follow directly as well. Hence every
nonzero sum of at most three products has a nonzero cylinder
with at most one free coordinate.

For later use, we prove by induction that a sum of at most
$s\ge3$ separated products has a nonzero cylinder with at
most $s-2$ free coordinates (or all $n$ coordinates if fewer
are available). The base $s=3$ was just proved. If every
local factor matrix has rank one, the assertion again
reduces to one product. Otherwise choose a coordinate $i$
whose local rank is $r\ge2$, and choose $r$ of its factors
as a basis. Expanding the other factors in this basis gives
\[
 f=\sum_{a=1}^r F_{i,j_a}(X_i)Q_a(\mathbf X_{\ne i}),
 \tag{I.74}
\]
where each $Q_a$ is a sum of at most $s-r+1\le s-1$
separated products: the original pivot term and at most one
term for each nonpivot column. If every cylinder of $f$
with at most $s-2$ free coordinates were zero, independence
of the basis factors would imply that every $Q_a$ vanished
on cylinders with at most $s-3$ free coordinates. The
induction hypothesis would force every $Q_a=0$, contrary
to $f\ne0$. The prime threshold inherited by a subproblem
on fewer variables is still sufficient.

Now suppose a specified $A_i$ has nullity $c\ge2$.
Expansion (I.74) makes each $Q_a$ a sum of at most $c+1$
products. The result just proved detects each nonzero $Q_a$
with at most $c-1$ free coordinates. Adding $i$ gives at
most $c$ free coordinates for $f$. Nullity zero is handled
by freeing $i$ alone, and nullity one by paragraph~2. This
proves all the stated cylinder bounds and, by paragraph~1,
the explicit cocharacter and sampling bounds.

The same finite family jointly preserves the coefficient
span of the vector of product terms under the local-nullity
promise. More precisely, if $V\subseteq F^w$ is that
coefficient span and the sampling grid lies in a working
extension $E$, the sampled vectors span $V\otimes_F E$
over $E$. Indeed, every $E$-linear combination is still
such a sum over $E$ after scaling or deleting columns of
$A_i$; column deletion cannot increase its nullity, and
scalar extension preserves the original matrix rank.
If an $E$-linear functional annihilates every sampled vector,
the scalar hitting theorem applied over $E$ forces it to
annihilate $V\otimes_F E$. The reverse inclusion follows
because every evaluated vector is a combination of the
original coefficient vectors. In particular no original-field
functional nonzero on $V$ is lost. This is the exact
joint-span criterion I-14. We do not assert that one member
of the family is always a BIWA or preserves the entire span
individually.

\paragraph{4. An actual additive cancellation of three dense products.}
The multiplicative conclusion in I-33 does not extend to an
additive relation among three products. Let $n=6$, $\delta=12$,
and assign the following distinct prime weights to variables
$A,B,C,Q,R,S$, respectively:
\[
 (1663,523,139,127,1093,331).
\]
All exceed $H=120$ and are below the polynomial bound above.
The polynomial
\[
 \begin{aligned}
 F={}&(1-A)(1-S^2)-(1-R^2)(1-C)\\
     &{}-C(1-B)(1-Q^{12})
 \end{aligned}
\]
is nonzero: the monomial $AS^2$ has coefficient one.
Its three displayed terms are separated products, and no
pairwise ratio is a monomial in the original variables.
Nevertheless this prime curve kills it identically in every
characteristic. To verify this, put $a=1663$, $b=523$,
$c=139$. The remaining weights satisfy
$2\cdot331=b+c$, $2\cdot1093=a+b$, and
$12\cdot127=a-c$. The elementary identity
\[
 (1-U)(1-VW)-(1-UV)(1-W)=(1-V)(W-U)
\]
with $U=T^a$, $V=T^b$, $W=T^c$ proves the assertion.

The repaired family detects this example. Freeing $A$ and
holding the other variables on their prime curve gives
$(T^{1663}-A)(1-T^{662})\ne0$. This example refutes the
specific proposed extension of the divisor lemma; it is
not a lower bound for sparse PIT or for the repaired family.

\paragraph{5. Nullity two already defeats every single-coordinate repair.}
Let $n=8$, $\delta=1$, and take the first eight primes
greater than $H=14$:
\[
 (17,19,23,29,31,37,41,43).
\]
Set
\[
 F=(X_1X_8-X_2X_7)(X_3X_6-X_4X_5).
\]
Both binomials vanish on the base curve, because
$17+43=19+41=60$ and $23+37=29+31=60$.
On any cylinder freeing one coordinate, at least one of
the two disjoint binomials remains unchanged and zero.
Thus every one-coordinate cylinder is identically zero,
including all substitutions that change that coordinate
to any alternate prime power or any rational function.

The polynomial has tensor rank exactly four in the separated
product model. Expansion gives the upper bound four. For
the lower bound, partition the variables as
$\{1,8,3,6\}$ and $\{2,7,4,5\}$.
Each binomial has a rank-two coefficient matrix across its
corresponding partition, and the product on disjoint
variables gives their Kronecker product, of rank four.
Every separated product has coefficient-matrix rank one,
so at least four terms are necessary. This argument works
in every characteristic. In the displayed four-term
representation each local factor matrix has rank two and
nullity two.

Freeing, for example, $X_1$ and $X_3$ makes both binomials
nonzero, so the two-coordinate repair proved above removes
this precise obstruction. The next unresolved efficiency
question is to avoid a sampling exponent proportional to
unbounded local nullity. Polynomially many alternate primes
alone cannot do that while only one coordinate changes.
No conversion of this repair into a Hoza path-isolation
construction has been proved.\hfill$\square$

\paragraph{Derivation record.}
The full proof is a deduction in this track from the I-33
divisor lemma and elementary coefficient elimination.
The specific prime numbers, the three-term identity, and
the tensor-rank-four obstruction were checked directly.
No literature-priority claim is made.

\subsection{I-35. A smaller contraction invariant, its exact composition rule, and the rank-one lift obstruction}
\label{I-35}

\paragraph{Scope.}
This result concerns the rank-$r$ outer matrix products left after
factoring higher-rank resets. It separates three assertions:
preserving all constant scalar boundary contractions of one bridge,
composing bridge substitutions with independent parameters, and
synchronizing those parameters. The first is strictly weaker than
preserving coefficient span or generic matrix rank. The second has
an exact positive proof. The third does not follow, and explicit
counterexamples below identify the missing quantifier. No general
ROABP hitting set is claimed.

\paragraph{1. The exact rank-one-annihilator criterion.}
Let $M(\mathbf X)$ be an $r\times r$ polynomial matrix over a field
$K$, and let $U\subseteq\operatorname{Mat}_r(K)$ be the span of its
coefficient matrices. After a polynomial substitution $G$, write
$U_G$ for the coefficient-matrix span of $M(G)$; then $U_G\subseteq U$.
All statements can instead be made over a specified working extension.
For rational substitutions, clear a common nonzero polynomial
denominator first. The span is independent of this clearing choice:
a constant functional annihilates the cleared vector if and only if
it annihilates the original rational vector.

Use the pairing $\langle C,A\rangle=\operatorname{tr}(C^{\mathsf T}A)$,
and let $\Sigma_1$ be the set of matrices of rank at most one. Then
\[
 \begin{gathered}
 uM(\mathbf X)v\ne0\ \Longrightarrow\ uM(G)v\ne0
       \quad\hbox{for every constant row }u\hbox{ and column }v
 \\
 \Longleftrightarrow\qquad
 U_G^\perp\cap\Sigma_1=U^\perp\cap\Sigma_1.
 \end{gathered}
 \tag{I.75}
\]
Indeed, $uAv=\langle u^{\mathsf T}v^{\mathsf T},A\rangle$, and every
nonzero rank-one matrix has this form. A scalar contraction is the
zero polynomial exactly when its functional annihilates the entire
coefficient space. Since $U_G\subseteq U$, one inclusion of the two
intersections is automatic and the other is precisely nonvanishing
preservation. For preservation after every scalar extension, impose
this equality over an algebraic closure. Thus the target is equality
of base loci on the Segre variety of boundary pairs, rather than
equality of the entire linear spaces or matrix image cones.

\paragraph{2. An explicit optimal test space of dimension $2r-1$.}
Index rows and columns by $0,\ldots,r-1$. Let $H_k$ have entry one
when $i+j=k$ and zero otherwise, for $0\le k\le2r-2$. Their disjoint
nonempty supports make them linearly independent. Put
\[
 a(T)=(1,T,\ldots,T^{r-1})^{\mathsf T},\qquad
 H(T)=a(T)a(T)^{\mathsf T}=\sum_{k=0}^{2r-2}T^kH_k.
 \tag{I.76}
\]
For any nonzero constant $u,v$,
\[
 uH(T)v=
   \left(\sum_{i=0}^{r-1}u_iT^i\right)
   \left(\sum_{j=0}^{r-1}v_jT^j\right)\ne0.
 \tag{I.77}
\]
This is a product of two nonzero polynomials over a domain, in every
characteristic. Consequently the $(2r-1)$-dimensional Hankel space
$\mathcal H_r=\operatorname{span}\{H_k\}$ has no nonzero rank-one
annihilator. It satisfies (I.75) with $U=\operatorname{Mat}_r(K)$,
although $\mathcal H_r$ is a proper subspace when $r\ge2$.
Evaluating $H(T)$ at $2r-1$ distinct field values spans this space by
Vandermonde interpolation; an extension may supply the values.

The dimension $2r-1$ is optimal over an algebraically closed field.
This is known completely-entangled-subspace background, not a new
claim: Parthasarathy (2004), Propositions~1.3--1.4 and Theorem~1.5
(P-B16), gives this two-factor construction and optimal codimension
over $\mathbb C$. The algebraic-closure formulation follows from
the same Segre dimension argument; its checked proof is retained
in the separate background audit. The lower bound need not hold
when only rational boundary pairs over a nonclosed field are tested.

This is an actual matrix-bridge substitution. Take
$M(\mathbf X)=(X_{ij})$, whose coefficient span is all of
$\operatorname{Mat}_r(K)$, and substitute $X_{ij}=T^{i+j}$.
The generic matrix rank drops from $r$ to one, while (I.77) preserves
every nonzero constant boundary contraction. Moreover this generic
matrix is realized by an invertible width-$2r$ interior word:
\[
 \prod_{i,j}\bigl(I_{2r}+X_{ij}E_{i,r+j}\bigr)
     =\begin{pmatrix}I_r&M(\mathbf X)\\0&I_r\end{pmatrix}.
 \tag{I.78}
\]
Here one may relabel matrix-unit indices from one instead of zero.
All products of two off-diagonal generators vanish, so the identity
holds in every order. Selecting the upper-right block uses constant
full-rank rectangular endpoints. Thus generic rank preservation is
strictly stronger than scalar boundary preservation even for an
actual bridge with invertible interior factors.

More generally, after scalar extension to an algebraic closure,
every coefficient space $U$ has a subspace $W\subseteq U$ of
dimension at most $2r-1$ with the same rank-one annihilators.
To prove this, remove the base locus from the Segre
variety and map a boundary pair to its nonzero evaluation functional
in $\mathbb P(U^*)$. The closure $Y$ of this image has dimension at
most $2r-2$. If $\dim U\le2r-1$, take $W=U$. Otherwise set
$D=\dim U$ and $c=2r-1$. The Grassmannian of $(D-c)$-dimensional
vector subspaces of $U^*$ has dimension $c(D-c)$. Requiring such
a subspace to contain one fixed projective point imposes $c$
conditions: its fiber is a Grassmannian of dimension
$c(D-c-1)$. The incidence variety of subspaces meeting $Y$
therefore has dimension at most
$\dim Y+c(D-c-1)<c(D-c)$. Its image is closed because $Y$ is
projective, so a nonempty open set of subspaces misses $Y$.
The annihilator of any one of them defines the desired $W$.
This proves existence of a small contraction test space, not
an explicit coefficient-oblivious substitution realizing it for
an arbitrary bridge. The Hankel construction realizes it for the
full generic matrix example only. No field-of-definition or
effective-construction bound is asserted for the general existence
statement.

\paragraph{3. The valid composition statement and its independent-seed cost.}
Let $B_j(\mathbf X_j)$, $1\le j\le s$, be square matrices with
pairwise disjoint variable sets. Suppose each substitution
$\mathbf X_j=G_j(\mathbf S_j)$ preserves every constant scalar
boundary contraction of $B_j$, and does so after every field extension.
If the seed sets $\mathbf S_j$ are algebraically independent of one
another and of all remaining variables, then
\[
 uB_1(\mathbf X_1)\cdots B_s(\mathbf X_s)v\ne0
 \ \Longrightarrow\
 uB_1(G_1)\cdots B_s(G_s)v\ne0.
 \tag{I.79}
\]
Replace the matrices successively. At step $j$, take as coefficient
field the rational function field in all unprocessed variables and
all previously introduced seeds. The product on the left of $B_j$
is a constant row over this field, and the product on its right is
a constant column. The current scalar polynomial is nonzero, so
the hypothesis over that extension field preserves its nonzeroness
when the fresh seed set $\mathbf S_j$ is introduced. Induction proves
(I.79). This argument requires no equality of coefficient spaces.

If adjacent bridges share cut variables $\mathbf Z$, first work
over the field $K(\mathbf Z)$ and leave those variables unassigned.
Assume the bridge preservation hypothesis over this field and all
its extensions. The same proof then applies to substitutions of
their disjoint interior variables. A later specialization of $\mathbf Z$ must still preserve
the resulting nonzero numerator and any relevant denominators.
Nothing in (I.79) supplies that specialization automatically.

For a prescribed decomposition into $s$ blocks, suppose each map
is polynomial over $K$ in its seeds and the retained cut variables,
uses at most $d$ seeds, and has total coordinate degree at most $L$.
If $z$ shared cut variables remain free, then the resulting
original scalar polynomial of total degree at most $D$ has at most
$sd+z$ parameters and total degree at most $D\max\{1,L\}$.
A full grid therefore costs at most
\[
 \bigl(D\max\{1,L\}+1\bigr)^{sd+z}.
 \tag{I.80}
\]
This displays the cost rather than concealing it in a genericity
assertion. The argument presupposes the decomposition; it does not
produce a uniform unknown-order sampler or identify the reset
locations from a scalar black box. For maps with arbitrary
$K(\mathbf Z)$ coefficients instead, numerator and denominator
degrees require separate bounds; (I.80) is not asserted for them.

\paragraph{4. Why constant contraction preservation does not synchronize.}
Even preserving each entire local coefficient space does not justify
using the same seed in different blocks. For example, with
$u=e_1^{\mathsf T}$ and $v=e_2$,
\[
 u(I+XE_{12})(I-YE_{12})v=X-Y\ne0.
 \tag{I.81}
\]
Each substitution $X=T$ or $Y=T$ separately preserves the full
coefficient space of its invertible matrix factor. Simultaneously
making both substitutions kills the scalar contraction.

The stronger Hankel example makes the boundary quantifier explicit.
For $r=2$, the map $H(T)$ in (I.76) preserves every nonzero constant
boundary contraction, but
\[
 (-T,1)H(T)e_1=0.
 \tag{I.82}
\]
With an independent boundary parameter $S$, the corresponding
expression is $T-S\ne0$. Diagonal identification $S=T$ creates the
cancellation. Treating all constant boundary pairs over extensions
is not a license to identify a new seed with a preexisting element
of the coefficient field.

In fact an unrestricted same-seed rational-boundary condition would
be impossible for $r\ge2$. Given any matrix $A(\mathbf S)$ over
$K(\mathbf S)$, there are nonzero $u(\mathbf S),v(\mathbf S)$ with
$uAv=0$: choose a nonzero row and then a nonzero vector in the
kernel of its resulting row functional. For the generic matrix
$M(\mathbf X)$, however, $uM(\mathbf X)v$ is nonzero over
$K(\mathbf S)$ because at least one product $u_iv_j$ is nonzero.
Thus a viable synchronization theorem must use restrictions on the
boundary functions actually induced by neighboring bridges. Asking
for all rational functions of the same seed is too strong.

\paragraph{5. A rank-one multiplicative lift cannot retain an entry for $r\ge2$.}
There is also a precise obstruction to turning the higher-rank
outer product into scalar product closure by a multiplicative lift.
Let $r\ge2$, let $w\ge r$, and embed $r\times r$ matrices as the
upper-left corner $\iota(B)$ of $w\times w$ matrices. Put
$P=\iota(I_r)$. Suppose a multiplicative map
\[
 \rho:\operatorname{Mat}_w(K)\longrightarrow\operatorname{Mat}_N(E)
\]
is given over some extension $E$, and suppose $\rho(P)$ has rank
one. There cannot be a fixed $E$-linear readout $\Lambda$ satisfying
\[
 \Lambda(\rho(\iota(B)))=B_{11}
 \quad\hbox{for every }B\in\operatorname{GL}_r(K).
 \tag{I.83}
\]
No polynomiality or linearity of $\rho$ is needed for this obstruction;
multiplicativity and the fixed linear readout suffice.

Indeed $Q=\rho(P)$ is a rank-one idempotent. Since
$P\iota(B)P=\iota(B)$, every lifted corner matrix lies in
$Q\operatorname{Mat}_N(E)Q=EQ$ and has the form $\sigma(B)Q$.
The identity case of (I.83) gives $\Lambda(Q)=1$, so
$\sigma(B)=B_{11}$. Multiplicativity would then make $B\mapsto B_{11}$
a multiplicative function on $\operatorname{GL}_r(K)$. But for
\[
 B=I_r+E_{12},\qquad C=I_r+E_{21},
\]
one has $B_{11}=C_{11}=1$ and $(BC)_{11}=2\ne1$, in every
characteristic. This is a contradiction. Multiplying the proposed
readout by any integral power of $\det B$ does not evade it, because
$B,C,BC$ all have determinant one.

Exterior powers are consistent with this obstruction: at a rank-$r$
reset, the $r$th exterior power has rank one and retains maximal-minor
data, which acts by the determinant on an invertible $r$-dimensional
corner. It cannot also recover an arbitrary fixed entry there by
such a readout. This excludes the stated functorial scalarization,
not all nonfunctorial or coefficient-dependent constructions.

\paragraph{Remaining concrete gate.}
The relevant smaller target is the Segre-annihilator condition for
the actual contractions, together with a synchronization theorem
for the restricted boundary functions produced by adjacent reset
bridges. The independent-seed statement and its cost are proved
above. The Hankel bridge proves that this target can preserve scalar
information while destroying generic matrix rank. Neither fact
establishes an efficient synchronized map for general rank-$r$
outer words. That last assertion remains open in this report.
\hfill$\square$

\paragraph{Background and provenance.}
The optimal Hankel test-space statement is the known two-factor
case of Parthasarathy (2004), Propositions~1.3--1.4 and Theorem~1.5,
preprint pp.~4--6, already recorded in P-B16. The complete local
proofs were independently reopened and read for this audit.
The projective facts used here are the Segre embedding, Stacks
Project Authors (2026), Lemma~27.13.6, tag~01WD, and the projective
intersection dimension theorem, Lemma~33.34.3, tag~0B2R. Their
primary statements and displayed proofs were read. The Hankel
identity, actual unipotent bridge realization, composition rule,
synchronization counterexamples, and multiplicative-corner
obstruction are proved directly above. No historical priority
claim or general derandomization consequence is made.

\subsection{I-36. A pairwise dominant center preserves maximal minors; five pairs for singular width-three words}
\label{I-36}

\paragraph{Statement and scope.}
Fix distinct Lagrange nodes $a_1,\ldots,a_n$ in a working field
$E$, and let $\lambda_i(Y)$ be their degree-at-most-$(n-1)$
Lagrange polynomials. The coefficient field is a subfield of $E$.
For $n\ge1$, define
\[
 \alpha_i=t+z\lambda_i(y),\qquad
 \mathcal G_i^{(5)}=t+\sum_{j=1}^{5}z_j\lambda_i(y_j).
 \tag{I.84}
\]
All displayed parameters are algebraically independent.

The map $\boldsymbol\alpha$ preserves every nonzero maximal-size
minor of every interval in a univariate matrix word whose individual
factors have generic rank at most that size. More explicitly, if
for integers $1\le r\le w$, the $A_i(X_i)$ are $w\times w$
polynomial matrices of generic rank at
most $r$, with distinct variables at the different positions, then
every nonzero $r\times r$ minor of every interval product remains
nonzero after $\mathbf X=\boldsymbol\alpha$. Identically zero minors
remain zero. The map depends on neither the matrices nor their order.

Consequently $\mathcal G^{(5)}$ hits every nonzero width-at-most-three
ROABP admitting a representation in which every variable-dependent
matrix factor has generic rank at most two. Constant invertible
width-three factors may be absorbed into neighboring factors or
the boundaries. Arbitrarily many singular factors and arbitrary
drops in the ranks of longer intervals are allowed. The variable
order is unknown to the sampler. This improves I-30's six-pair bound
on this specified singular-factor subclass; it does not prove that
two centered pairs suffice, and it does not extend this improvement
to arbitrary variable-dependent invertible width-three factors.

\paragraph{1. Dominance on every coordinate pair.}
For $i\ne j$, the polynomial $\lambda_i(y)-\lambda_j(y)$ is nonzero:
its value at $y=a_i$ is one. Over $E(y)$, the linear substitution
\[
 \begin{pmatrix}X_i\\X_j\end{pmatrix}
 =\begin{pmatrix}1&\lambda_i(y)\\1&\lambda_j(y)\end{pmatrix}
   \begin{pmatrix}t\\z\end{pmatrix}
 \tag{I.85}
\]
is invertible. It therefore induces an injection
$E[X_i,X_j]\hookrightarrow E(y)[t,z]$. Every nonzero polynomial
depending on at most two specified coordinates stays nonzero.
The singleton case follows directly from the coefficient one on
$t$; it covers $n=1$. The same argument holds after any extension
of the coefficient field independent of the new parameters.

This is a formal polynomial assertion. It does not require a
single numerical parameter choice to work over a small base field;
the later interpolation step supplies the needed finite evaluations.

\paragraph{2. Rank-one words factor into adjacent bivariate contractions.}
Let $B_i(X_i)$ be nonzero polynomial matrices of generic rank one.
Choose any nonzero entry $b_i$ of $B_i$, let $p_i$ be its column,
and let $q_i$ be its row. The vanishing of the two-by-two minors gives
\[
 B_i=\frac{p_iq_i}{b_i}
 \quad\hbox{over the rational function field.}
\]
For every interval $[s,t]$ and constant boundary row $u$ and column
$v$ of the appropriate sizes, clearing the nonzero univariate
pivots gives the exact polynomial identity
\[
 \left(\prod_{i=s}^{t}b_i\right)
       uB_sB_{s+1}\cdots B_tv
 = (up_s)(q_tv)
       \prod_{i=s}^{t-1}(q_ip_{i+1}).
 \tag{I.86}
\]
The two endpoint factors are univariate. Each middle factor depends
on only the two distinct adjacent variables $X_i,X_{i+1}$.
This remains valid for a one-factor interval, with an empty middle
product. The variable sets of different factors on the right may
overlap; their nonzero product is still nonzero in a polynomial
domain.

If the original scalar contraction is nonzero, every factor on the
right of (I.86) is nonzero. Part~1 preserves all those factors and
all the pivots under the same substitution $\boldsymbol\alpha$.
The substituted scalar contraction is therefore nonzero. In
particular every nonzero entry of every rank-one interval word is
preserved. This proof uses scalar product closure after an exact
rank-one factorization, not an unsupported composition of arbitrary
constant-boundary tests with shared parameters.

\paragraph{3. Applying the factorization to maximal minors.}
For a matrix $A$, write $C_r(A)$ for its $r$th compound matrix.
Cauchy--Binet gives
\[
 C_r(A_s\cdots A_t)=C_r(A_s)\cdots C_r(A_t).
 \tag{I.87}
\]
If a factor has generic rank below $r$, its compound is zero and
every interval containing it has all $r$-minors identically zero.
Otherwise all factors in the interval have generic rank exactly
$r$, so each $C_r(A_i)$ has generic rank one. Part~2 applies to
this compound word and preserves each nonzero entry of (I.87),
which is precisely the corresponding $r$-minor of the original
interval. This proves the first assertion of the statement.

Compound matrices occur only in the proof. The sampler does not
construct them, recover their pivot entries, inspect the input
representation, or enumerate the unknown variable order. All
intervals and minors are preserved by the same formal map: the
argument establishes a nonzero image for each one individually,
and any finite product of these nonzero images is nonzero.

\paragraph{4. Completing the rank-one stage in width three.}
Remove constant invertible factors from the width-three word by
absorbing them into neighboring matrices or the constant boundary
vectors. This does not increase the generic rank of any remaining
factor. Unused variables cause no problem for (I.85). A representation
of smaller width may be padded with zero rows and columns to width
three, so its factors also have generic rank at most two.

Apply parts~1--3 with $r=2$. Every nonzero two-by-two interval
minor survives at the center $\boldsymbol\alpha$. There are no
generic rank-three intervals in the remaining word. The two
constant rank-one sentinels used in I-21 to encode the scalar
output introduce no rank-three interval or additional variable.

Now use the local rank-test construction proved in I-21 and
the endpoint absorption of I-30. In each selected nonzero
rank-one test, the interior intervals have generic ranks at
least two. Here their ranks are exactly two, and they have
already been preserved at $\boldsymbol\alpha$. Every nonzero
univariate polynomial endpoint vector is nonzero at its center:
each $\alpha_i$ is transcendental over the original coefficient
field, since its coefficient on $t$ is one. Constant nonzero
endpoints are covered directly.

Pass to the field of the first-stage parameters and translate
to independent residual variables. I-30 absorbs the endpoints
into unit factors, producing a word with matching generic and
constant interval ranks in $\{2,3\}$. The complete support
lemma I-29 with $(w,r)=(3,2)$ supplies a nonzero coefficient
whose monomial has support at most
\[
 2(3-2)+3-1=4,
 \tag{I.88}
\]
including both endpoint positions. Four fresh SV pairs preserve
every selected rank-one test, because a polynomial with such a
nonzero low-support coefficient cannot vanish under $\mathrm{SV}_4$.
The new pairs are independent of the first-stage parameter field.
The earlier two-by-two minors remain nonzero, as is seen by setting
the new amplitudes to zero. The local-rank induction of I-21 now
preserves the nonzero scalar output. Adding the first pair and
the four residual pairs gives exactly $\mathcal G^{(5)}$ in (I.84).

This last argument is the already checked rank-one completion
from I-30, with its rank-two-preserving center supplied by the
new rank-one-compound factorization. It does not assert that
the first-stage center alone preserves rank-one intervals.

\paragraph{5. Exact finite sampling and field requirements.}
For a total-degree-at-most-$D$ input, the substituted polynomial
has degree at most $D$ in $t$, at most $(n-1)D$ in each of the
five selectors, and total degree at most $D$ in the five
amplitudes. The tensor grid in $t$ and the selectors, together
with the Newton simplex in the amplitudes used in I-21, yields
at most
\[
 (D+1)\bigl((n-1)D+1\bigr)^5\binom{D+5}{5}
 \tag{I.89}
\]
evaluations. The working field must contain at least
$\max\{n,D+1,(n-1)D+1\}$ distinct elements. If it is a proper
extension of the coefficient field, black-box evaluation over
that extension is part of the evaluation model. All arguments,
including the simplex interpolation, hold in every characteristic.
The constant-polynomial case $D=0$ needs only one evaluation;
the displayed formula also equals one. No polynomial dependence
on arbitrary width is inferred from this fixed-width result.
\hfill$\square$

\paragraph{Provenance and the remaining gate.}
The rank-one pivot factorization and scalar product closure are
the established mechanism recorded in I-B06 and P-14. The new
deduction here is that pairwise dominance supplies an entire
maximal-minor preservation stage with one centered interpolation
pair, followed by the previously proved four-pair completion.
No historical priority claim is made. The centered two-pair
question for rank-two reset words remains open; the explicit
uncentered failure P-16 does not decide it.

\subsection{I-37. Two centered pairs for affine singular width-three words}
\label{I-37}

\paragraph{Statement and exact promise.}
Let $E$ be any field, let $a_1,\ldots,a_n\in E$ be distinct, and let
$\lambda_i(Y)$ be their Lagrange polynomials. A nonzero polynomial
$f\in E[X_1,\ldots,X_n]$ is assumed to admit a width-at-most-three
ROABP whose variable-dependent factors are affine-linear in their
own variables and have generic rank at most two. Constant factors
are allowed. The order is arbitrary and unknown to the map. Then
\[
 f\left(t+z_1\lambda_1(y_1)+z_2\lambda_1(y_2),\ldots,
        t+z_1\lambda_n(y_1)+z_2\lambda_n(y_2)\right)\ne0.
 \tag{I.90}
\]
This holds in every characteristic. It improves I-36's five-pair
bound on the stated affine subclass. It does not cover a mixture
containing variable-dependent generic-rank-three factors, or
replace the affine hypothesis by arbitrary individual degree.

We prove the slightly more convenient rational statement. For
arbitrary $c_1,\ldots,c_n\in E^\times$, set
\[
 R_{k,i}=t c_i+\sum_{h=1}^{k}\frac{z_h}{a_i-y_h}.
 \tag{I.91}
\]
Every input in the promised class satisfies $f(\mathbf R_2)\ne0$.
All parameters here are algebraically independent; rational
identities are understood in their rational function field.

\paragraph{1. Pairwise dominance and differential consequences.}
Write $\alpha_i=t c_i+z/(a_i-y)$. For $i\ne j$, the determinant
\[
 \frac{c_i}{a_j-y}-\frac{c_j}{a_i-y}
 =\frac{c_i(a_i-y)-c_j(a_j-y)}{(a_i-y)(a_j-y)}
 \tag{I.92}
\]
is nonzero. If its numerator were zero, comparison of its linear
and constant coefficients would give $c_i=c_j$ and
$c_i(a_i-a_j)=0$, a contradiction. The substitution in the two
coordinates is thus an invertible linear change of $(t,z)$ over
$E(y)$. It preserves every nonzero polynomial in at most two
specified coordinates, over any fixed extension of $E$ as well.

Suppose, toward a contradiction, that $f(\mathbf R_2)=0$. Setting
the second amplitude to zero gives $f(\boldsymbol\alpha)=0$.
The coefficient of a new amplitude $\epsilon$ in
$f(\alpha_i+\epsilon/(a_i-u):1\le i\le n)=0$ is
\[
 \sum_i\frac{f_i(\boldsymbol\alpha)}{a_i-u}=0,
 \qquad f_i=\partial f/\partial X_i.
\]
The distinct simple poles force $f_i(\boldsymbol\alpha)=0$ for
every $i$: multiply by $\prod_j(a_j-u)$ and set $u=a_i$.
Since $\boldsymbol\alpha+s\mathbf c$ is obtained by replacing
$t$ by $t+s$, these identities also give
\[
 \sum_j c_j f_{ij}(\boldsymbol\alpha)=0
 \quad\hbox{for every }i,
 \qquad H(\boldsymbol\alpha)\mathbf c=0,
 \tag{I.93}
\]
where $f_{ij}$ denotes the second partial derivative. This uses
only the coefficient of $s$, so it is valid in every characteristic.
For multilinear $f$, every diagonal entry $f_{ii}$ is zero.

\paragraph{2. Irreducible factors inherit the representation promise.}
We may extend the coefficient field to its algebraic closure.
This is a proof device, not a requirement that the final sampler
evaluate over the algebraic closure. A nonzero substitution over
an extension is a nonzero substitution over its original field.
The kernel of the rational substitution $\mathbf X=\mathbf R_2$
is prime, because the target is a field. Hence an assumed nonzero
annihilator has a nonconstant irreducible factor $g$ that is also
an annihilator.

The affine-factor representation makes $f$ multilinear. For a
product of nonzero polynomials, degree in each fixed variable is
additive. Thus the nonconstant irreducible factors of a multilinear
polynomial depend on disjoint variable sets. Specialize all
variables outside the set used by $g$ so that the product of the
other factors is a nonzero scalar. Such a specialization exists
over the infinite algebraically closed field. It realizes a
nonzero scalar multiple of $g$ as a specialization of the original
matrix word. Remove the now constant factors by absorbing their
products into a neighboring factor or a boundary. Multiplication
by constant matrices preserves affine dependence and cannot
increase generic rank. The surviving factors therefore still
have generic rank at most two. A width-two or width-one word may
be padded with zero coordinates to width three.

Consequently it suffices to treat an irreducible $f$ that depends
on every one of its remaining variables, with those variables
listed in their representation order. The corresponding nodes
remain distinct and the corresponding $c_i$ remain nonzero.
The proof does not enumerate or discover this order. If at most
two variables remain, part~1 already contradicts annihilation.
Assume henceforth that $n\ge3$.

\paragraph{3. Every strict internal interval has generic rank two.}
We use an elementary polynomial rank-one factorization. If a
nonzero matrix $B$ over a multivariate polynomial ring over a
field has rank one over its fraction field, then $B=pq$ for a
polynomial column $p$ and polynomial row $q$. Indeed, take a
nonzero polynomial column, divide its entries by their greatest
common divisor, and call the primitive resulting column $p$.
Every column of $B$ is a rational multiple $h p$. In a reduced
fraction for $h$, each irreducible denominator factor would have
to divide every nonzero entry of $p$. Primitivity excludes it,
so $h$ is polynomial. This proves the asserted factorization.
In particular, when $B$ is multilinear, the polynomial factors
$p$ and $q$ depend on disjoint variable sets: the maximum entry
degrees satisfy
$\deg_{X_i}B=\deg_{X_i}p+\deg_{X_i}q\le1$.

Write the irreducible scalar as
$f=u A_1(X_1)\cdots A_n(X_n)v$. For a strict internal interval
$2\le s\le t\le n-1$, the product $B=A_s\cdots A_t$ cannot
be zero, because then $f=0$. If it had rank one, the preceding
factorization would give
\[
 f=(uA_1\cdots A_{s-1}p)(qA_{t+1}\cdots A_nv).
 \tag{I.94}
\]
The first factor depends on $X_1$ and the second on $X_n$,
because $f$ does and each of these variables appears on only
its indicated side. Both factors are therefore nonconstant,
contradicting irreducibility. Since each original factor has
rank at most two, every strict internal interval has rank
exactly two.

The pairwise-dominant map of part~1 preserves every nonzero
two-by-two minor of every such interval. To see this with the
same quantifiers as I-36, take second compounds: each
$C_2(A_i)$ has generic rank one, and Cauchy--Binet turns an
interval minor into an entry of their product. The pivot-cleared
rank-one identity (I.86) writes that entry, times univariate
nonzero pivots, as a product of univariate endpoint polynomials
and adjacent bivariate contractions. Part~1 preserves every
factor. Thus every internal interval has rank two at
$\mathbf X=\boldsymbol\alpha$, simultaneously.

\paragraph{4. Compatible constants and nonzero active endpoints.}
Work over $K=\overline E(t,z,y)$, and translate to independent
variables $Y_i=X_i-\alpha_i$. Absorb the first and last factors
into the affine endpoint row $u(X_1)=uA_1(X_1)$ and column
$v(X_n)=A_n(X_n)v$. The interior consists of positions
$2,\ldots,n-1$.

The constant row $u(\alpha_1)A_2(\alpha_2)$ is nonzero. Indeed,
the polynomial row $u(X_1)A_2(X_2)$ is nonzero, since otherwise
$f$ would be zero, and one of its nonzero bivariate entries
survives by part~1. Similarly
$A_{n-1}(\alpha_{n-1})v(\alpha_n)\ne0$. Because every internal
constant interval has rank two, multiplication by the remaining
constant interior factors is injective on these two-dimensional
active spaces. In particular both endpoint projections to the
persistent active space are nonzero.

Here is an explicit justification for the compatible gauges.
View the interior constants $C_i=A_i(\alpha_i)$ as linear maps
from the state space at cut $i$ to the state space at cut $i-1$.
At each internal cut, the image of the next constant factor is
two-dimensional and is complementary to the kernel of the
preceding constant factor, because the product of those two
factors has rank two. Use that image as the active space and
that kernel as its one-dimensional complement. Choose an
active complement to the kernel at the final cut, and transport
its basis leftwards through the constants; the full interval
rank condition makes all transported active bases invertible.
At the initial cut choose any complementary line. In these
bases every interior constant is
\[
 P=\operatorname{diag}(1,1,0),\qquad
 A_i(\alpha_i+Y_i)=P+Y_iD_i.
 \tag{I.95}
\]
The nonzero active endpoint row $a$ and column $b$ satisfy
$ab=f(\boldsymbol\alpha)=0$. Choose an active basis with
$a=e_1^{\mathsf T}$ and $b=e_2$; this is possible by taking
the second active basis vector to be $b$ and choosing the first
to have pairing one with $a$. Use this same active basis change
at every cut, keeping all copies of $P$ fixed.

The passive endpoint components can also be removed while
leaving those constants fixed. At the initial cut use
$L=\left(\begin{smallmatrix}I_2&h\\0&1\end{smallmatrix}\right)$:
$L^{-1}P=P$ and the row changes from $(a,\rho)$ to
$(a,ah+\rho)$. Choose $ah=-\rho$. At the final cut use
$R=\left(\begin{smallmatrix}I_2&0\\k^{\mathsf T}&1\end{smallmatrix}\right)$:
$PR=P$ and the column becomes
$(b,\sigma-k^{\mathsf T}b)^{\mathsf T}$. Choose
$k^{\mathsf T}b=\sigma$. These independent endpoint changes
also work when there is only one interior factor. They preserve
all matrix ranks and affine dependence. We have obtained
\[
 f(\boldsymbol\alpha+\mathbf Y)
 =(e_1^{\mathsf T}+Y_1p)
   \prod_{i=2}^{n-1}(P+Y_iD_i)
   (e_2+Y_nq).
 \tag{I.96}
\]
Here $p$ is a row and $q$ a column. All gauges are constant with
respect to the residual variables $\mathbf Y$; their dependence
on the center parameters is allowed.

Since every affine interior pencil has determinant identically
zero, the coefficient of $Y_i$ in
$\det(P+Y_iD_i)$ gives $(D_i)_{33}=0$. The first-derivative
identities from part~1 give
\[
 p_2=0,\qquad q_1=0,\qquad (D_i)_{12}=0
 \quad(2\le i\le n-1).
 \tag{I.97}
\]
Thus each active two-by-two derivative block is lower triangular,
and the passive state has no self-loop in either $P$ or $D_i$.

\paragraph{5. A tridiagonal Hessian forces every output path to vanish.}
If $|i-j|>1$, at least one constant factor $P$ separates the
two differentiated positions in (I.96). It kills the passive
coordinate. The active derivative blocks are lower triangular,
the active part of $p$ is supported on state one, and the active
part of $q$ on state two. Therefore $f_{ij}(\boldsymbol\alpha)=0$
for $|i-j|>1$. The diagonal is zero by multilinearity.

Consequently $H(\boldsymbol\alpha)$ is symmetric tridiagonal
with zero diagonal. Equation (I.93), and $c_i\ne0$ for every
$i$, force all its entries to be zero. The first row gives
$H_{12}=0$; the second row then gives $H_{23}=0$; continuing
along the path proves $H_{i,i+1}=0$ for all $i$. This argument
uses no division by two.

The adjacent entries have particularly simple formulas:
\[
 \begin{aligned}
 H_{12}&=p_3(D_2)_{32},\\
 H_{i,i+1}&=(D_i)_{13}(D_{i+1})_{32}
                 &&(2\le i\le n-2),\\
 H_{n-1,n}&=(D_{n-1})_{13}q_3.
 \end{aligned}
 \tag{I.98}
\]
All alternative active contributions are zero by (I.97).

Expand (I.96) as a sum over state paths. An active path cannot
move directly from state one to state two, because every
$(1,2)$ entry is zero. The passive state three has no self-loop.
Hence any path that first reaches state two after starting in
state one must contain a transition $1\to3$ immediately
followed by $3\to2$. Its two entries have product equal to
one of the middle expressions in (I.98), times the two residual
variables, and the product is zero. If the path starts in the
passive endpoint component, its first possible move into state
two is covered by the first expression in (I.98); if it first
returns to state one, the preceding argument applies later.
Finally, an output through the passive component of the final
column is covered by the last expression in (I.98), unless
state two was already reached, in which case an earlier zero
transition has already occurred. A path staying in state one
cannot exit through the active final column because $q_1=0$.
Thus every contributing state path contains a zero factor.
The polynomial (I.96) is identically zero over $K$, contrary
to the nonzero input. This proves $f(\mathbf R_2)\ne0$.

\paragraph{6. Conversion to ordinary centered SV and exact sampling.}
Put $A(T)=\prod_i(T-a_i)$ and $d_i=A'(a_i)\ne0$. For the
diagonally rescaled polynomial
$F(W_1,\ldots,W_n)=f(W_1/d_1,\ldots,W_n/d_n)$, the affine
singular-factor promise is unchanged. Take $c_i=d_i$ in (I.91).
The Lagrange identity
\[
 d_i\left(t+\sum_{h=1}^{2}z_h\lambda_i(y_h)\right)
 =d_i t-\sum_{h=1}^{2}\frac{z_h A(y_h)}{a_i-y_h}
 \tag{I.99}
\]
shows that the ordinary centered map is the rational map just
proved, after the invertible change of amplitudes
$z_h\mapsto-z_hA(y_h)$ in the rational parameter field.
Therefore its polynomial substitution is nonzero, proving (I.90).

For total degree at most $D$ (here $D\le n$), that substitution
has degree at most $D$ in $t$, at most $(n-1)D$ in each selector,
and total degree at most $D$ in the two amplitudes. The Newton
simplex interpolation argument already proved in I-21 yields
a hitting set of at most
\[
 (D+1)\bigl((n-1)D+1\bigr)^2\binom{D+2}{2}.
 \tag{I.100}
\]
The working field must supply at least
$\max\{n,D+1,(n-1)D+1\}$ distinct elements. A proper working
extension uses the same extension-evaluation convention as
I-21 and I-36. The proof holds in all characteristics, and
$D=0$ requires just one evaluation. The selected order and
all coefficient-dependent factorizations occur only in the
proof, not in the sampler.\hfill$\square$

\paragraph{Provenance and the remaining obstruction.}
This deduction uses the pairwise dominance and rank-one compound
factorization proved in I-36. The first-order pole-extraction
observation was supplied independently in this run by the purity
agent while investigating the Hu--van Melkebeek--Morgan
vanishing-ideal framework; the full argument above proves the
needed arbitrary-slope and arbitrary-characteristic version
directly. No historical priority claim is made.

The new closing step is specific to affine singular pencils:
after rank preservation their passive constant entry and passive
linear coefficient both vanish. A variable-dependent invertible
factor can propagate the passive state, so nonadjacent Hessian
entries need not vanish and (I.93) need not force all adjacent
transition products to vanish. Higher-degree singular factors
also need not have a zero passive entry at every positive degree.
Thus the mixed rank-two/rank-three centered-two-pair gate remains
open, as does an unrestricted polynomial-in-width construction.

\subsection{I-38. Two centered pairs with isolated invertible affine factors}
\label{I-38}

\paragraph{Statement.}
The conclusion and sample bound of I-37 also hold for a width-at-most-three
affine univariate matrix word with the following larger promise:
after constant factors have been absorbed, no two consecutive
variable-dependent factors both have generic rank three. Thus
arbitrarily many generic-rank-three affine factors may occur,
provided a singular factor separates each two of them. The order
remains unknown, the nodes are any distinct working-field elements,
and the conclusion holds in every characteristic.

Equivalently, for any nonzero slopes $c_i$ and distinct nodes $a_i$,
the rational map $R_{2,i}=tc_i+z_1/(a_i-y_1)+z_2/(a_i-y_2)$
hits the stated class. The diagonal rescaling (I.99) then proves
ordinary common-center $t\mathbf1+\mathrm{SV}_2$ nonvanishing.
We prove the rational formulation.

\paragraph{1. The one-pole center dominates every three coordinates.}
Let $\alpha_i=tc_i+z/(a_i-y)$. Pairwise dominance is (I.92).
For distinct $i,j,k$, solve for $t,z$ from $X_i,X_j$ over $E(y)$.
The third coordinate has the form
\[
 X_k=A(y)X_i+B(y)X_j,
 \qquad A,B\in E(y).
 \tag{I.101}
\]
The two coefficients cannot both be constant. If they were,
comparison of the $z$ coefficients would give
$1/(a_k-y)=A/(a_i-y)+B/(a_j-y)$, which the distinct poles
exclude. Hence the right side of (I.101) is a nonconstant
rational function of $y$ over $E(X_i,X_j)$, and is transcendental
over that field. Thus the three coordinates are algebraically
independent. This proves dominance in every characteristic,
without a Jacobian argument. It also proves the claim over every
fixed extension of the coefficient field.

\paragraph{2. The representation promise survives irreducible reduction.}
Assume a nonzero annihilator and apply I-37's prime-kernel and
disjoint-variable factor argument. Specializing the variables
outside an irreducible factor makes the intervening matrix
factors constant. Absorb each constant block into the next
remaining factor, and absorb the final block into the terminal
column. All remaining factors are affine and their ranks cannot
increase.

The no-adjacent-rank-three promise also persists. If two adjacent
remaining factors both had rank three, both corresponding
original factors would have had rank three. The original promise
puts a singular original factor between them. It has now been
specialized into the intervening constant block, so that block
has rank at most two. Absorbing it into the later factor makes
that factor singular, a contradiction. This reasoning does not
require any specialized singular factor to keep its original
rank. A zero specialization cannot occur in the realized nonzero
scalar word.

We may therefore assume that the input is irreducible and depends
on every remaining variable. As in I-37, a strict internal
rank-one interval has a polynomial outer factorization and
splits the scalar into two nonconstant factors. Every strict
internal interval consequently has generic rank two or three.
If there are at most three remaining variables, part~1 already
hits the scalar. Assume $n\ge4$ below.

\paragraph{3. One pole preserves every internal interval rank.}
Each generic-rank-three factor is invertible at its univariate
center, because its nonzero determinant is univariate and the
center is transcendental. Its second compound is likewise
invertible there.

Consider a generic-rank-two interval. In its second compound
word, a singular original factor becomes a rank-one factor,
and every intervening rank-three factor becomes invertible.
There is at most one such intervening factor between consecutive
rank-one factors. Write a nonzero rank-one factor as $p_iq_i/b_i$
using a nonzero univariate pivot. The connecting scalar between
two successive rank-one factors is therefore either
\[
 q_i(X_i)p_j(X_j)
 \quad\hbox{or}\quad
 q_i(X_i)C_2(A_k(X_k))p_j(X_j).
 \tag{I.102}
\]
It depends on at most three distinct coordinates. Since the
original compound interval product is nonzero, every connecting
scalar is nonzero. Part~1 preserves all these scalars and all
univariate pivots. The remaining rank-one endpoint vectors are
nonzero at their centers, and any full-rank prefix or suffix is
invertible there. The compound interval product stays nonzero,
so the original interval still has rank two. A generic-rank-three
interval consists entirely of invertible factors and retains
rank three. Thus all strict internal interval ranks are preserved
simultaneously. The proof requires preservation of their ranks;
it does not assert preservation of every individual two-minor.

Absorb $A_1,A_n$ into the affine endpoint row and column. There
is an internal rank-two factor: otherwise the at-least-two
internal positions would be consecutive rank-three factors.
The first internal rank-two factor is at position two or three.
The endpoint row times the prefix through that factor is a
nonzero polynomial row in at most three variables, and remains
nonzero at the center by part~1. All subsequent constant
intervals have rank at least two, so propagation of its active
two-dimensional row space is injective. Hence the left endpoint
has a nonzero active projection for the full constant interior
word. The reversed argument at the last internal rank-two
factor gives a nonzero active projection of the right endpoint.

\paragraph{4. Compatible $I/P$ constants and endpoint shears.}
Over the field $K=\overline E(t,z,y)$, the constant interior word
has full-product rank two, and every subinterval rank two or
three. We give the compatible-space argument explicitly.
Use row-vector orientation. Choose a two-dimensional initial
space $U_0$ complementary to the left kernel of the full
constant product, and propagate it through the constants to
spaces $U_i$. Each propagation is an isomorphism on $U_i$.
Choose a complementary line $V_m$ at the final cut. For the
suffix product $T_i$ from cut $i$ to the end, put
\[
 V_i=\{v:vT_i\in V_m\}.
\]
If $T_i$ has rank two, its row image equals $U_m$, so
$V_i=\ker T_i$ is a line. If it has rank three, the inverse
image of $V_m$ is a line. In either case $U_i\cap V_i=0$,
and each constant factor sends the corresponding passive line
into the next one. Active bases transported through the word
make every active block $I_2$. The passive block is zero at
rank-two factors and a nonzero scalar at rank-three factors;
rescale the successive passive bases to make the latter scalar
one. Thus all interior factors have the form
\[
 A_i(\alpha_i+Y_i)=C_i+Y_iD_i,
 \qquad C_i\in\{P,I_3\},\quad P=\operatorname{diag}(I_2,0),
 \tag{I.103}
\]
with $C_i=P$ precisely at the rank-two positions. No consecutive
interior constants are both $I_3$.

The nonzero active endpoint vectors pair to
$f(\boldsymbol\alpha)=0$. Normalize them to $e_1^{\mathsf T}$
and $e_2$ by one common active change of basis. The endpoint
passive components must also be removed. Use the upper shear
$L$ from I-37 at every cut of the initial string of $I_3$
constants, up to the first $P$. Then
$L^{-1}I_3L=I_3$ along that string and $L^{-1}P=P$ at its end,
while the initial passive row component is killed. Similarly
use the lower shear $R$ from I-37 at every cut of the terminal
string of $I_3$ constants after the last $P$:
$PR=P$ and $R^{-1}I_3R=I_3$. This kills the passive column
component. If the first and last $P$ are the same factor,
$L^{-1}PR=P$ permits both shears simultaneously. Thus all
constants in (I.103) remain as stated, and
\[
 f(\boldsymbol\alpha+\mathbf Y)
 =(e_1^{\mathsf T}+Y_1p)
   \prod_{i=2}^{n-1}(C_i+Y_iD_i)
   (e_2+Y_nq).
 \tag{I.104}
\]

The pole extraction of I-37 gives $f_i(\boldsymbol\alpha)=0$
for every $i$. Therefore $p_2=q_1=(D_i)_{12}=0$. At every
singular factor, the coefficient of $Y_i$ in
$\det(P+Y_iD_i)=0$ also gives $(D_i)_{33}=0$. An invertible
factor may have a nonzero passive derivative entry; it is not
silently discarded.

\paragraph{5. Two tangent directions kill the bandwidth-two Hessian.}
Set $h_i=1/(a_i-y)$. Differentiate the identities
$f_i(\boldsymbol\alpha)=0$ with respect to $z$ and $y$.
Equivalently, extract their first coefficients under independent
parameter increments. Since $z\ne0$ in the rational function
field, this gives in every characteristic
\[
 H(\boldsymbol\alpha)\mathbf h=0,
 \qquad H(\boldsymbol\alpha)\mathbf h^{\circ2}=0,
 \tag{I.105}
\]
where $\mathbf h^{\circ2}=(h_1^2,\ldots,h_n^2)$. The diagonal
of $H$ is zero by multilinearity.

If $i<j$ and a singular interior position lies strictly between
them, the intervening constants contain $P$. Their product
is $P$, which removes the passive coordinate, while the active
derivative blocks are lower triangular. Hence $H_{ij}=0$.
If there is no such position, all intervening constants are
$I_3$, and $H_{ij}$ is exactly the product of the state-one-to-
passive entry at position $i$ and the passive-to-state-two
entry at position $j$, with the endpoint versions $p_3$ and
$q_3$ when appropriate. There are no alternative active
contributions because all $(1,2)$ entries vanish.

Since no two consecutive internal positions are invertible,
every gap of length at least three contains a singular
intermediate position. Thus $H_{ij}=0$ whenever $|i-j|>2$.
In the first row, at most two future entries remain. For two
indices $j\ne k$, the corresponding coefficient determinant
in (I.105) is
\[
 \det\begin{pmatrix}h_j&h_k\\h_j^2&h_k^2\end{pmatrix}
 =h_jh_k(h_k-h_j)\ne0.
 \tag{I.106}
\]
Both entries are therefore zero. If only one entry remains,
$h_j\ne0$ suffices. By symmetry the earlier entries in the
next row are now zero; apply the same argument to its at most
two future entries. Induction gives $H=0$.

Finally expand (I.104) into state paths, as in I-37. A path
cannot move directly from active state one to state two.
Before its first arrival at state two, take its last departure
from state one to the passive state, or its initial passive
endpoint entry if it has not yet visited state one. Every
intermediate passive stay must occur at an invertible factor:
a singular factor has both passive constant and passive
derivative entries zero. The product of the departure and
arrival entries is the corresponding Hessian entry $H_{ij}$,
and is zero. Intermediate passive stays only multiply this
zero product. The same argument covers an output through the
passive terminal component; an output through active state one
is excluded by $q_1=0$. Therefore every output path is zero,
contradicting the assumed nonzero scalar. This proves the
rational two-pair assertion and hence the ordinary centered
two-pair assertion.\hfill$\square$

\paragraph{Sampling, provenance and the remaining gate.}
The exact field and interpolation requirements are unchanged
from I-37: for total degree at most $D\le n$, the sample count
is (I.100), over a working field with at least
$\max\{n,D+1,(n-1)D+1\}$ distinct elements, or an extension
with the explicitly stated evaluation convention. No input
order or representation is reconstructed by the sampler.

The extra proof ingredients are the function-field proof of
triple dominance, the preserved adjacency promise under
irreducible specialization, the rank-one compound connectors
through one invertible factor, and the two tangent kernel
identities. Purity independently suggested the use of
$\mathbf h,\mathbf h^{\circ2}$ instead of slope-dependent
minors, avoiding an invalid Jacobian shortcut in characteristic
two. No historical priority claim is made.

Two or more consecutive invertible factors remain outside
this proof. Their compound connectors involve at least four
coordinates, so three-coordinate dominance alone does not
preserve their rank; even at a rank-preserved center, the
Hessian can have more than two future entries in a row. The
larger semiseparable Hessian structure and the consequences
of the assumed two-pair identity for possible one-pole rank
loss remain the specific next gate.

\subsection{I-39. Two exact failed repairs for unrestricted mixed affine words}
\label{I-39}

\paragraph{Scope.}
We prove two obstructions to the next proposed extensions of I-38.
First, the ordinary common-center one-pair map can lower a strict
internal interval from generic rank two to rank one even in an
irreducible multilinear width-three scalar word whose every generic
interval rank is two or three. Second, the symmetric semiseparable
Hessian pattern, all three tangent kernel identities, and the whole
additional-pole quadratic identity do not by themselves force the
Hessian to be zero. Neither construction is a counterexample to
the centered two-pair sampler. The first scalar is hit already
by one pair; the second is an abstract family of quadratic jets,
with a stated failure of integration to a fixed polynomial.
Both constructions are over characteristic zero.

\paragraph{1. A rational identity and its residue forms.}
We recall, and verify the needed parts of, the explicit construction
of P-16. Put $a=\sqrt{3/2}$ and, for an integer $d\ge2$, set
\[
 \begin{aligned}
 P_1(T)&=T^{2d}-1,&P_2(T)&=T^d-a,\\
 P_3(T)&=-2(T^d+a),&P_4(T)&=T^{2d}-2,\\
 R_1&=P_2/P_1,&R_2&=1/P_2,\\
 R_3&=1/P_3,&R_4&=P_3/P_4.
 \end{aligned}
 \tag{I.107}
\]
Work in a splitting field of these polynomials. Their four pole
sets are disjoint and simple, of sizes $2d,d,d,2d$ respectively;
none of the displayed numerators cancels a pole. For a pole
$b$ of $R_j$, let $\rho_b\ne0$ be its residue. Partial fractions
and decay $R_j(T)=O(T^{-d})$ give
\[
 R_j(T)=\sum_{b\in\mathcal A_j}\frac{\rho_b}{T-b},
 \qquad \sum_{b\in\mathcal A_j}\rho_b b^e=0
                    \quad(0\le e\le d-2).
 \tag{I.108}
\]
Moreover $P_2P_3=-P_1-P_4$, so
$R_1R_2+R_1R_4+R_3R_4=0$.
For the disjointly supported nonzero linear forms
$L_j=\sum_{b\in\mathcal A_j}\rho_b X_b$, define
\[
 Q=L_1L_2+L_1L_4+L_3L_4\ne0.
 \tag{I.109}
\]
Nonvanishing follows, for example, from any coefficient using
one variable from group one and one from group two. If $g$ has
degree at most $d-1$, division by $T-u$ and (I.108) give
\[
 L_j\left(\frac{g(b)}{b-u}:b\in\bigcup\mathcal A_j\right)
 =-g(u)R_j(u).
 \tag{I.110}
\]
Therefore $Q$ vanishes on all these numerator-at-most-$(d-1)$,
one-pole evaluations. These are the only P-16 facts used below.

\paragraph{2. Ordinary one-pair rank loss on an irreducible affine word.}
Take $d=2$, giving twelve core nodes. They come in opposite pairs,
and every $R_j$ is an even rational function. Add the five nodes
$0,3,-3,4,-4$, giving seventeen distinct nodes in total. Assign
the four nonzero new nodes to variables $s,u,v,t$, respectively;
the coordinate at zero is unused by the input polynomial.
Let
\[
 A(T)=\prod_{b\text{ among all seventeen nodes}}(T-b),
 \qquad d_b=A'(b),\qquad
 L_j^*=\sum_{b\in\mathcal A_j}\rho_b d_bX_b.
 \tag{I.111}
\]
The polynomial $A$ is odd and $A'$ is even. The residues of an
even rational function at opposite poles are negatives of one
another. It follows that $\sum_{b\in\mathcal A_j}\rho_b d_b=0$
for every group. Hence under the ordinary common-center map
$X_b=\tau+z\lambda_b(y)$ on these seventeen nodes,
\[
 L_j^*=zA(y)R_j(y),\qquad
 Q^*=L_1^*L_2^*+L_1^*L_4^*+L_3^*L_4^*=0.
 \tag{I.112}
\]
Indeed $d_b\lambda_b(y)=A(y)/(y-b)$, so this is an exact
polynomial identity after clearing the displayed rational notation.
All coefficients $\rho_bd_b$ are nonzero, and $Q^*$ is a
nonzero quadratic before substitution.

Put $A_0=L_1^*+L_3^*$, $B_0=L_2^*+L_4^*$ and $C_0=L_2^*L_3^*$.
The middle word consists of the affine unipotent factors
$I-\rho_bd_bX_bE_{21}$ in groups one and three, and
$I-\rho_bd_bX_bE_{32}$ in groups two and four, in group order.
Within each group their product is $I-L_j^*E$, since $E^2=0$.
Thus its full product is exactly
\[
 M=I-A_0E_{21}-B_0E_{32}+C_0E_{31}.
 \tag{I.113}
\]
Let $P_L=\operatorname{diag}(0,1,1)$ and
$P_R=\operatorname{diag}(1,1,0)$. Their corner is
\[
 P_LMP_R=
 \begin{pmatrix}0&0&0\\-A_0&1&0\\C_0&-B_0&0\end{pmatrix},
 \qquad
 \det\begin{pmatrix}-A_0&1\\C_0&-B_0\end{pmatrix}=Q^*.
 \tag{I.114}
\]
Use genuinely variable-dependent singular anchor factors
\[
 J_L(u)=(I+uE_{32})P_L,\qquad
 J_R(v)=P_R(I+vE_{21}).
\]
They have rank two at every value. The strict internal interval
$J_L(u)M J_R(v)$ has generic rank two because $Q^*\ne0$.
It has rank exactly one at the center (I.112): its indicated
two-by-two determinant is zero, while its $(2,2)$ entry remains
one. The outer unipotent multipliers in the anchors preserve
this rank calculation.

Now take the actual scalar word
\[
 F=e_3^{\mathsf T}(I+sE_{32})J_L(u)M J_R(v)
                      (I+tE_{21})e_1
   =(s+u-B_0)(t+v-A_0)-Q^*.
 \tag{I.115}
\]
Every variable is read once; every matrix factor is affine.
The two anchors are the only singular factors. An interval
containing no anchor is invertible, one containing exactly one
anchor has rank two, and one containing both has rank two
because (I.114) has nonzero determinant. Thus all generic
interval ranks are in $\{2,3\}$.

The scalar $F$ is irreducible. The invertible polynomial change
of variables $S=s+u-B_0$, $T=t+v-A_0$, keeping the remaining
variables fixed, converts it to $ST-Q^*$. As a polynomial in
$T$ over the UFD of the other variables, its coefficients $S$
and $-Q^*$ are coprime. It is primitive and linear, hence
irreducible over that UFD by Gauss's lemma. This also proves
irreducibility after extending to the algebraic closure.

Nevertheless $F$ is detected by the very center that loses
the internal rank. By (I.112), its image is the product of
the two factors in (I.115). Each has coefficient two on
$\tau$, while $A_0,B_0$ have no $\tau$ term. The image has
coefficient four on $\tau^2$ and is nonzero. This demonstrates
that irreducibility alone does not repair one-pole interval-rank
preservation, and that requiring preservation of every interval
rank is stronger than the desired scalar nonvanishing.

\paragraph{3. The longer-block Hessian conditions alone also fail.}
Take $d=4$ in part~1, with its twenty-four core nodes, and let
$H_0$ be the symmetric zero-diagonal Hessian of $Q$ in (I.109).
In the group order, define vectors $\alpha,\beta$ by putting
$\alpha_b=\rho_b$ in groups one and three and zero otherwise,
and $\beta_b=\rho_b$ in groups two and four and zero otherwise.
For $i<j$ the upper entries satisfy
$(H_0)_{ij}=\alpha_i\beta_j$. This is exactly the single-block
upper rank-one semiseparable pattern arising from an arbitrary
invertible stretch in the I-38 Hessian calculation.

Let $h_i(y)=1/(a_i-y)$, where $a_i$ now lists the twenty-four
core nodes, and set
\[
 D(y)=\operatorname{diag}((a_i-y)^2),\qquad
 H(y)=D(y)H_0D(y).
 \tag{I.116}
\]
It is nonzero, symmetric and zero-diagonal, and its upper
entries still factor as
$[\alpha_i(a_i-y)^2][\beta_j(a_j-y)^2]$.
Every $L_j$ kills polynomial evaluations of degrees at most
two by (I.108). Therefore $H_0$ annihilates each such evaluation
vector, and
\[
 H(y)\mathbf1=H(y)\mathbf h(y)
              =H(y)\mathbf h(y)^{\circ2}=0.
 \tag{I.117}
\]
The three vectors entering $H_0$ here are respectively
$((a_i-y)^2)_i$, $((a_i-y))_i$, and $\mathbf1$.

It also satisfies the complete extra-pole quadratic identity:
for an independent pole $u$,
\[
 \mathbf h(u)^{\mathsf T}H(y)\mathbf h(u)
 =2Q\left(\frac{(a_i-y)^2}{a_i-u}:1\le i\le24\right)=0.
 \tag{I.118}
\]
This is (I.110) with numerator $(T-y)^2$ of degree two.
Consequently the tangent kernels and the additional-pole
quadratic restriction do not suffice, even with the exact
semiseparable pattern, to force $H=0$.

This family is not itself the Hessian pullback of a fixed
polynomial at $\tau\mathbf1+z\mathbf h(y)$. If it were,
specializing $z=0$ would make that pullback independent of $y$,
whereas every nonzero entry of (I.116) is a nonconstant
degree-four polynomial in $y$. Thus the construction explicitly
exhibits the missing coordinate-ring and Hessian-integrability
constraints; it does not silently assert they have been met.
Multiplying or clearing denominators to enter the center's
coordinate ring would still require checking those constraints,
and the original width-three affine representation, anew.
\hfill$\square$

\paragraph{Provenance and next gate.}
The rational identity and residue forms are from the fully
checked P-16 construction; all extra matrix, parity, irreducibility
and Hessian deductions needed here are written out. No historical
priority claim is made. The remaining two-pair target must use
the assumed scalar annihilation, rather than irreducibility
alone, to control rank loss; or retain the integrability and
affine compound constraints beyond the displayed Hessian
identities. Neither obstruction settles that target or the
growing-width polynomial-size problem.


\subsection{I-40. A determinant-weight budget for translation-invariant matrix words}
\label{I-40}

\paragraph{Statement.}
Let $E$ be any field and let $0\ne f\in E[X_1,\ldots,X_n]$ be
homogeneous and multilinear of degree $D$, where $n\ge1$.
Assume the formal common-translation identity
\[
 f(X_1+s,\ldots,X_n+s)=f(X_1,\ldots,X_n).
 \tag{I.119}
\]
Suppose that the coefficient-cut ranks of $f$ in some order are
at most $w$. Put $m=n-2D$. Then $m\ge0$.
There is a minimal affine univariate matrix representation of
$f$ in that order, with bond sizes equal to its coefficient-cut
ranks, having the following property: every consecutive run of
generic-rank-$w$ factors in its consistent $w\times w$ padding
has length at most $m$. The same bound survives absorption of
constant factors. This assertion is valid in every characteristic;
the translation identity is a polynomial identity, not merely an
identity on the rational points of a finite field.

For $w\le3$ and $m\le1$, I-37 and I-38 consequently give
\[
 f\bigl(t\mathbf1+\mathrm{SV}_2(\mathbf y,\mathbf z)\bigr)\ne0,
 \qquad
 f\left(\left(tc_i+\frac{z_1}{a_i-y_1}
                         +\frac{z_2}{a_i-y_2}\right)_{i=1}^n\right)\ne0
 \tag{I.120}
\]
for arbitrary distinct nodes $a_i$ and arbitrary nonzero slopes
$c_i$ in the working field. No order is supplied to either map.
For this corollary use padding to width three even if the actual
width is smaller; if no rank-three core occurs, I-37 applies.

Together with P-17, the result also proves that balanced rational
two-pair evaluation hits every nonzero multilinear polynomial
of ROABP width at most three on at most nine variables, in every
characteristic and with unknown order. This last assertion is
about the balanced rational map. The proof does not identify it
with ordinary unweighted centered $\mathrm{SV}_2$ for every
nine-variable input.

\paragraph{1. The established highest-weight input and conventions.}
Write $V=E e_0\oplus E e_1$ and use the standard tensor action of
$G=\mathrm{SL}_2$. Let
\[
 u_+(s)=\begin{pmatrix}1&s\\0&1\end{pmatrix},
 \qquad h(t)=\operatorname{diag}(t,t^{-1}).
\]
We use the standard characteristic-free fact that a nonzero
$u_+$-fixed vector in a rational $G$-module, on which $h(t)$ acts
as $t^d$, has $d\ge0$. Jantzen (n.d.), Lecture~I, p.~2
(Frobenius reciprocity) and p.~3 (the explicit
$\mathrm{SL}_2$ induction computation), supply this exact input
in every characteristic; see I-B09. Upper- versus lower-Borel
conventions are changed by dualizing and conjugating by the Weyl
element. The split-torus weight decomposition is the established
integer-character grading of G-B02. These facts concern rational
group representations, and do not invoke complete reducibility
or Lie algebra eigenvalues. The convention check and an
elementary Bruhat verification are retained in the separate
\nolinkurl{I40_highest_weight_audit.md}.

Multihomogenize $f$ as the tensor
\[
 F=\sum_{S\subseteq[n]}c_S
       \bigotimes_{i=1}^n e_{\mathbf1_{i\in S}}
   \ \equiv\ \left(\prod_i U_i\right)
                f(V_1/U_1,\ldots,V_n/U_n).
 \tag{I.121}
\]
The upper unipotent fixes $e_0$ and sends $e_1$ to $e_1+s e_0$.
Thus (I.119) says exactly that $u_+(s)F=F$. Homogeneity gives
$h(t)F=t^{n-2D}F=t^mF$. The cited highest-weight fact gives
$m\ge0$. Scalar extension to an algebraic closure is harmless
here: it preserves polynomial identities and every coefficient
rank. The support spaces and matrix cores below may all be
defined over the original field $E$.

\paragraph{2. Both sides of a coefficient cut bound its determinant weight.}
Let $L_i\subseteq V^{\otimes i}$ and
$R_i\subseteq V^{\otimes(n-i)}$ be the left and right support
spaces of $F$ at the cut after position $i$. Concretely, $L_i$
is the image of contraction with all linear functionals on the
right tensor factors, and $R_i$ is defined symmetrically.
Their common dimension $r_i$ is the coefficient-cut rank;
$r_0=r_n=1$. Since $F$ is a semi-invariant of the upper Borel,
both support spaces are stable under that Borel. Indeed applying
a Borel element to a contraction is, up to its nonzero scalar
on $F$, another contraction after applying the inverse action
to the functional. Applying the inverse element gives equality
of the resulting spaces.

Let the determinant characters of the torus on these spaces be
$t^{\delta_i}$ and $t^{\delta_i'}$, respectively. They are integer
characters. The upper unipotent has determinant one on either
space: its representation matrix $G(s)$ has entries in $E[s]$
and inverse $G(-s)$, so its determinant is a unit in $E[s]$;
its value at zero is one. Consequently the determinant line of
$L_i$ is an upper-unipotent-fixed line in the rational
$G$-module $\bigwedge^{r_i}(V^{\otimes i})$. The highest-weight
fact gives $\delta_i\ge0$. The same argument on the right gives
$\delta_i'\ge0$.

In bases of $L_i$ and $R_i$, the tensor $F$ has a nonsingular
$r_i\times r_i$ coefficient matrix $C$. Torus semi-invariance
therefore gives
\[
 G_L(t) C G_R(t)^{\mathsf T}=t^m C.
\]
Taking determinants yields the exact budget
\[
 \delta_i+\delta_i'=r_i m,
 \qquad 0\le\delta_i\le r_i m.
 \tag{I.122}
\]
This is a statement in the character lattice $\mathbb Z$.
In positive characteristic, neither the inequality nor the
equality is reduced modulo the characteristic.

\paragraph{3. Minimal affine cores and their exact covariance.}
There is a natural inclusion
\[
 \iota_i:L_i\longrightarrow L_{i-1}\otimes V.
\]
To verify it, write $F$ as an element of
$L_{i-1}\otimes V\otimes V^{\otimes(n-i)}$ and contract its last
$n-i$ factors. These inclusions are Borel-equivariant. Choose
bases of the $L_i$, with the basis of $L_0=E$ equal to $1$ and
that of $L_n=EF$ equal to $F$. Write
\[
 \iota_i(b_{i,b})
   =\sum_a b_{i-1,a}\otimes
       \bigl(A_i^0[a,b]e_0+A_i^1[a,b]e_1\bigr),
 \qquad A_i(X_i)=A_i^0+X_iA_i^1.
 \tag{I.123}
\]
Repeatedly applying the inclusions reconstructs every tensor
coefficient, so the scalar matrix product
$A_1(X_1)\cdots A_n(X_n)$ equals $f$. The matrix shape is
$r_{i-1}\times r_i$, and these sizes are minimal because they
are the actual coefficient-cut ranks. This construction uses
only support-space inclusions and linear algebra, over any
field.

Let $G_i(s)$ be the upper-unipotent action matrix on $L_i$.
Equivariance of (I.123), comparing the coefficients of $e_0$
and $e_1$, gives
\[
 A_i(X_i)G_i(s)=G_{i-1}(s)A_i(X_i+s).
 \tag{I.124}
\]
In particular this formula includes both endpoints, with
$G_0=G_n=1$. It is an identity of matrices over $E[X_i,s]$;
there are no coefficient-dependent rational denominators.

Suppose $r_{i-1}=r_i=r$ and $A_i$ has generic rank $r$.
The determinant-one property of $G_{i-1}(s)$ and $G_i(s)$
implies
\[
 \det A_i(X_i+s)=\det A_i(X_i).
\]
Setting $X_i=0$ proves that this nonzero determinant is a
nonzero constant. Thus $A_i^0$ is invertible. For the torus,
comparison of the $e_0$ coefficient of (I.123) gives
\[
 A_i^0G_i(t)=t\,G_{i-1}(t)A_i^0.
\]
Taking determinants and cancelling $\det A_i^0\ne0$ proves
\[
 \delta_i=\delta_{i-1}+r.
 \tag{I.125}
\]
This is the key local restriction: a square full-rank core
increases the determinant character by its entire bond size.

\paragraph{4. The run bound and the width-three corollary.}
Pad each core consistently to $w\times w$, embedding each
$E^{r_i}$ as the first $r_i$ coordinates of $E^w$. Use the first
coordinate row and column as the scalar endpoints. This
preserves the product exactly. A padded factor can have generic
rank $w$ only if $r_{i-1}=r_i=w$ and its unpadded core is full
rank. On a run of $\ell$ such factors, (I.125) telescopes to
\[
 w\ell=\delta_{\mathrm{end}}-\delta_{\mathrm{start}}
           \le wm,
 \qquad \ell\le m,
 \tag{I.126}
\]
by (I.122) at the two bounding cuts. This proves the asserted
budget.

To absorb constant factors, multiply each intervening constant
block into the following variable-dependent factor, and absorb
the terminal block into the final column. Ranks cannot increase.
If two or more remaining full-rank factors form a consecutive
run, every constant factor removed between them must itself
have had full rank; otherwise its product would make the later
factor singular. Thus those remaining factors belong to one
full-rank run before absorption. The length of the remaining
run is no larger. Initial constant blocks and the final boundary
cannot create an extra adjacent variable factor. The same run
bound therefore holds after absorption.

For width at most three, use padding to three. If $m=0$, every
factor is singular, so I-37 applies. If $m=1$, no adjacent
variable-dependent factors have generic rank three after
absorption, so I-38 applies. This proves (I.120), with its exact
node and slope quantifiers. Its ordinary polynomial sampler
has the already proved size
\[
 (D+1)\bigl((n-1)D+1\bigr)^2\binom{D+2}{2}
 \tag{I.127}
\]
over a working field with at least
$\max\{n,D+1,(n-1)D+1\}$ elements. For $D=0$ one evaluation
suffices. A proper extension uses the explicit extension-input
convention of I-21 and I-37; the proof does not assume access to
extension evaluations when the black box only accepts base-field
points.

\paragraph{5. The nine-variable balanced two-pair theorem.}
Fix distinct nodes and let $\mathcal I_2$ denote the vanishing
ideal of balanced rational evaluation with numerator and
denominator degrees at most two, as in P-17. P-17 has already
proved the entire multilinear case on at most eight variables
and the homogeneous cubic case on at most ten variables, in
every order and characteristic. We use that stable result
without changing its statement.

Suppose that a nonzero multilinear $f\in\mathcal I_2$ on nine
variables has width at most three in some order. The established
homogeneous-degree decomposition of this ideal and the
lowest-component closure for ROABP width (P-B20, Theorem~9.13
and Proposition~8.3) give a nonzero homogeneous annihilator of
width at most three, of degree between three and six. The cubic
case of P-17 excludes degree three. Multilinear reversal
preserves both the balanced ideal and all coefficient-cut ranks,
so it excludes degree six and reduces degree five to degree
four. We may therefore suppose that $f$ is a homogeneous
quartic.

Balanced rational evaluation is closed under adding a constant.
Thus every coefficient of $f(X+s\mathbf1)$ in $s$ lies in
$\mathcal I_2$. Each coefficient of $s^j$ has degree $4-j$,
and this ideal has no nonzero multilinear component of degree
at most two. Hence
\[
 f(X+s\mathbf1)=f(X)+s g(X),
 \qquad g\in\mathcal I_2\text{ homogeneous cubic}.
 \tag{I.128}
\]
If $g\ne0$, it has the same cut-rank upper bound as $f$.
Indeed common translation is an invertible product of local
linear changes on the coefficient tensor, preserving cut ranks
over $E(s)$. With $q=s^{-1}$, the coefficient matrix of
$s^{-1}f(X+s\mathbf1)=qf(X)+g(X)$ is regular at $q=0$.
Every four-by-four minor vanishes before specialization and
therefore at $q=0$. Thus every coefficient-cut matrix of $g$
has rank at most three. P-17's nine-variable cubic theorem
forces $g=0$.

We have proved the translation identity (I.119) for this
quartic. Its weight is $m=9-2\cdot4=1$. Part~4 therefore gives
nonzero evaluation by the rational map with all slopes $c_i=1$.
That map is balanced of degrees at most two, contradicting
$f\in\mathcal I_2$. This proves the full nine-variable result.
All coefficient extractions and rank specializations in this
argument are valid in every characteristic.

\paragraph{6. An explicit finite balanced sampler.}
For any multilinear polynomial $f$ of total degree at most $D$,
write
\[
 R_i=t+\frac{z_1}{a_i-y_1}+\frac{z_2}{a_i-y_2},
 \qquad
 Q(\mathbf y)=\prod_{i=1}^n(a_i-y_1)(a_i-y_2).
 \tag{I.129}
\]
The rational map $\mathbf R$ has the same vanishing ideal as
balanced degree-two rational evaluation. One way to check the
scope is to pass to an algebraic closure: a generic quadratic
denominator has two distinct roots, and partial fractions give
exactly a constant plus two simple poles. The denominators with
distinct roots form a dense open subset in every characteristic.
Thus the parametrizations have the same formal polynomial
identities; no assertion about every rational function splitting
over the original finite field is needed.

For a nonzero pullback, $P=Qf(\mathbf R)$ is a nonzero polynomial.
Multilinearity is used here: expanding one monomial of $f$,
each index contributes either $(a_i-y_1)(a_i-y_2)$ or
\[
 t(a_i-y_1)(a_i-y_2)+z_1(a_i-y_2)+z_2(a_i-y_1).
\]
Consequently $P$ has degree at most $n$ in each selector $y_h$
and total degree at most $D$ in $(t,z_1,z_2)$. Choose $n+1$
distinct selector values, all outside $\{a_1,\ldots,a_n\}$,
and $D+1$ distinct amplitude values $b_0,\ldots,b_D$.
Sample the three amplitudes at the Newton simplex
\[
 (t,z_1,z_2)=(b_j,b_k,b_\ell),\qquad
               j,k,\ell\ge0,\quad j+k+\ell\le D.
\]
Univariate interpolation in the two selectors and the lower-set
Newton interpolation in the three amplitudes detect every
nonzero polynomial of these degrees. Every selected denominator
is nonzero. Thus the nine-variable theorem has a finite hitting
set of at most
\[
 (n+1)^2\binom{D+3}{3}
 \tag{I.130}
\]
points, for $n\le9$, over a working field with at least
$\max\{2n+1,D+1\}$ elements. The nodes and the sets of sample
values are fixed in advance and are independent of the hidden
order and all input coefficients. Extension evaluations, if
needed, obey the same explicitly stated black-box convention
as above.\hfill$\square$

\paragraph{Scope and the next gate.}
The determinant budget is not a polynomial-in-width PIT
construction. It controls runs of square full-width cores of a
homogeneous translation-invariant tensor; it does not constrain
every lower-rank transition or arbitrary nonhomogeneous words.
At weight two it permits two consecutive invertible cores, so
I-38 alone no longer closes the scalar question. The
translation-invariant ten-variable quartic is a concrete next
case, and the unrestricted mixed width-three problem and the
growing-width isolation transfer remain open here.

The highest-weight input is established representation theory.
The determinant-character budget, its core covariance
application, and the nine-variable deduction are proved above
as deductions in this research record, without a historical
priority claim. The standard dominance calculation and the
independent all-characteristic checks are recorded separately.


\subsection{I-41. Weight-two connectors and the balanced ten-variable theorem}
\label{I-41}

\paragraph{Statement and scope.}
Let $E$ be any field, and fix arbitrary distinct nodes
$a_1,\ldots,a_n$ in $E$. Every nonzero homogeneous multilinear
polynomial $f$ of width at most three in some variable order,
satisfying
\[
 f(X+s\mathbf1)=f(X),\qquad 0\le n-2\deg f\le2,
 \tag{I.131}
\]
has nonzero pullback under the balanced rational map
\[
 R_i=t+\frac{z_1}{a_i-y_1}+\frac{z_2}{a_i-y_2}.
 \tag{I.132}
\]
The order is arbitrary and unknown to the map, and the result
holds in every characteristic. Together with P-17 and I-40,
this proves the same balanced two-pair conclusion for every
nonzero multilinear width-at-most-three polynomial on at most
ten variables, without a homogeneity or translation-invariance
promise.

For total degree at most $D$, both assertions admit the finite
sample bound
\[
 (n+1)^2\binom{D+3}{3}
 \tag{I.133}
\]
over a working field of size at least
$\max\{2n+1,D+1\}$, with the extension-input convention of I-40.
The new weight-two argument uses the common rational center
$t\mathbf1$. It does not claim arbitrary nonzero center slopes
or the ordinary unweighted centered-$\mathrm{SV}_2$ conclusion
for this entire larger class.

\paragraph{1. Reduction to an irreducible tensor of weight two.}
Work over an algebraic closure as a proof device. Suppose a
nonzero input in (I.131) vanishes under (I.132). Since the
substitution target is a field, some nonconstant irreducible
factor $g$ also vanishes. Irreducible factors of a multilinear
polynomial use disjoint variable sets; specializing the other
sets to nonzero scalar factor values realizes $g$, up to scalar,
as a width-at-most-three polynomial in the induced order, as
proved in I-37.

These factors are homogeneous. Indeed the minimum and maximum
total degrees in a product add, and equality of those degrees
for a homogeneous product forces equality for every factor.
Each factor is also invariant under common translation of its
own variables. To see this, divide the translated product
identity by the original product over $E(s)$. The resulting
factors are rational functions in disjoint sets of variables,
so each ratio belongs to $E(s)^\times$. Comparing its highest
homogeneous part with the original homogeneous factor forces
that ratio to be one. I-40 therefore gives a nonnegative weight
for every factor. Their weights sum to $n_{\rm used}-2\deg f$,
which is at most the original $n-2\deg f$; unused coordinates
only increase the latter. Thus the chosen factor has weight
at most two.

Weights zero and one are already covered by I-40, as are all
inputs on at most nine variables. We may henceforth assume
that $f$ is irreducible, depends on every remaining variable,
has weight $m=n-2D=2$, and has $n\ge10$. Construct its actual
minimal affine cores by I-40's support inclusions, and pad
consistently to width three. We use those new minimal cores,
not a possibly nonminimal representation inherited from the
original factorization. Write the unpadded bond ranks as $r_i$.
Every interior cut has $r_i\ge2$: rank one would split the
scalar into a product in two nonempty disjoint variable sets.
Moreover
\[
 r_1=r_{n-1}=2,
 \tag{I.134}
\]
since a single multilinear variable supplies at most two
coefficient vectors. In particular the first and last internal
factors, $A_2$ and $A_{n-1}$, are singular after padding.

As in I-37, every strict internal interval has generic rank at
least two. A rank-one interval has a polynomial outer
factorization over the multivariate UFD and splits the scalar
into two nonconstant factors, because the first and last
variables occur on opposite sides. Each such interval therefore
has rank two or three. I-40 bounds every run of rank-three
factors by two. The only new rank-preservation issue is a pair
of consecutive rank-three cores between rank-two anchors.

\paragraph{2. A saturated run has invariant bounding support spaces.}
Suppose $A_{j+1},A_{j+2}$ are both rank three. They are a
maximal such run. The determinant characters of I-40 satisfy
\[
 \delta_j=0,\qquad \delta_{j+1}=3,
 \qquad\delta_{j+2}=6,
 \qquad\delta'_{j+2}=0.
 \tag{I.135}
\]
Indeed both increments are three and the allowed interval is
$0\le\delta_i\le3m=6$ at these cuts.

A highest-weight line of weight zero is fixed by the whole
$\mathrm{SL}_2$. For completeness, the Bruhat identity in the
I-40 audit shows that its opposite-unipotent orbit is a
polynomial of degree at most zero; it is therefore constant.
The two unipotent subgroups generate the group. Apply this to
the determinant lines in (I.135). A subspace is determined by
its nonzero decomposable exterior vector, equivalently by the
vectors whose exterior product with that vector is zero.
Thus $L_j$ and $R_{j+2}$ themselves are $\mathrm{SL}_2$-stable
three-dimensional spaces.

Taking all single-coordinate contractions propagates this
stability to $L_{j-1}$ and to $R_{j+3}$, respectively. These
are the actual support spaces: contracting $L_j$ by all
functionals on its final physical factor spans $L_{j-1}$,
and the analogous statement holds on the right. Every torus
weight in $L_j\subseteq V^{\otimes j}$ is congruent to $j$
modulo two. Since its determinant weight is zero and its
dimension is three, $j$ is even. Also $n$ is even because
$n-2D=2$. The adjacent support spaces lie in odd tensor powers.
Weyl pairing of their nonzero integer weights makes their
dimensions even; the width bound and nonzeroness therefore give
\[
 \dim L_{j-1}=\dim R_{j+3}=2.
 \tag{I.136}
\]
This also shows that both adjacent anchors exist away from the
scalar endpoints.

\paragraph{3. The canonical anchors, without a semisimplicity assumption.}
Let $W=L_{j-1}$. Its two torus weights are $q,-q$ for an odd
positive integer $q$. If $q>1$, all four weights of $W\otimes V$
are nonzero, namely $q+1,q-1,-q+1,-q-1$. Every
$\mathrm{SL}_2$-stable subspace of this tensor product then has
even dimension, again by Weyl pairing. This contradicts its
three-dimensional stable subspace $L_j$. Therefore $q=1$.
This elementary argument excludes nontrivial Frobenius twists
as well as every other possible larger character.

A two-dimensional rational $\mathrm{SL}_2$-module with weights
$1,-1$ is the standard module in every characteristic. One can
check the convention directly: torus conjugation forces the
upper and lower unipotent entries to be $b s$ and $c s$;
neither is zero, since the two unipotents generate the group.
The Gauss identity for their product forces $bc=1$, and a
basis rescaling gives the standard matrices. This is also the
degree-one case of the established I-B09 input.

Now identify $W$ with $V$. A stable three-plane in $V\otimes V$
must contain the one-dimensional weight-two space: otherwise
Weyl symmetry would remove weight minus two as well, leaving
only the two-dimensional weight-zero space. Applying the
opposite unipotent to $e_0\otimes e_0$ shows that the three-plane
contains, and therefore equals,
\[
 \Gamma^2 V=\operatorname{span}\{
 e_0e_0,\ e_0e_1+e_1e_0,\ e_1e_1\}.
 \tag{I.137}
\]
Here juxtaposition denotes tensor product. The three displayed
vectors are independent in characteristic two as well. We do
not identify this submodule with its dual or invoke a
semisimple decomposition of $V\otimes V$.

Using its displayed weight basis, the left anchor is exactly
\[
 C(u)=\begin{pmatrix}1&u&0\\0&1&u\end{pmatrix},
 \qquad u=X_j.
 \tag{I.138}
\]
The same argument on the right identifies $R_{j+2}$ with the
three-plane (I.137) in $V\otimes R_{j+3}$, where the latter
two-dimensional residual space is standard.

The two full-rank cores have nonzero constant determinants by
I-40. Make their constant matrices the identity by constant
changes of the two intermediate bases, starting from the chosen
left weight basis. The torus actions on those bases are shifted
successively by weights one and two. Their linear coefficient
matrices consequently have entries only one step above the
diagonal: their row weight minus column weight must equal two.
Thus their product has the form
\[
 M(x,y)=
 \begin{pmatrix}1&A&B\\0&1&C\\0&0&1\end{pmatrix},
 \tag{I.139}
\]
where $A,C$ are homogeneous linear forms in $x=X_{j+1}$ and
$y=X_{j+2}$, and $B$ is homogeneous quadratic (in fact a scalar
multiple of $xy$). All these changes of basis are over the
coefficient field and are independent of the evaluation-center
parameters.

The left weights at the end of the run are $(4,2,0)$, while
the right weights in $R_{j+2}$ are $(2,0,-2)$. The tensor has
weight two and full rank three across that cut. Its coefficient
pairing is therefore antidiagonal with three nonzero coefficients
$d_1,d_2,d_3$. Expand the canonical right three-plane (I.137),
and order the next left residual basis oppositely to the right
standard basis. The right anchor is
\[
 D(v)=\begin{pmatrix}d_1v&0\\d_2&d_2v\\0&d_3\end{pmatrix},
 \qquad v=X_{j+3},\qquad d_1d_2d_3\ne0.
 \tag{I.140}
\]
This uses the actual dual pairing. In particular it remains
valid when $\Gamma^2V$ is not self-dual in characteristic two.
The full-rank constant gauges identify torus actions only;
they do not assert an isomorphism of the neighboring
upper-unipotent modules.

\paragraph{4. The four-coordinate connector cannot vanish at one pole.}
The connecting matrix between the two-dimensional cuts is
$C(u)M(x,y)D(v)$. Its determinant is
\[
 \begin{split}
 \det(CMD)={}&d_1d_2v^2+d_1d_3v(C+u)\\
 &+d_2d_3\bigl(u^2+Au+AC-B\bigr).
 \end{split}
 \tag{I.141}
\]
This follows by multiplying the two-by-three and three-by-two
matrices; equivalently its $u^2$ coefficient is the bottom
two-row minor of $MD$. It is a homogeneous quadratic in the
four distinct coordinate variables, and its $u^2$ coefficient
is the constant $d_2d_3\ne0$.

Set $\alpha_i=t+z/(a_i-y_0)$, where $t,z,y_0$ are independent.
At $y_0=a_j$, only $u=\alpha_j$ has a pole among these four
coordinates. The term $d_2d_3u^2$ has a double pole with nonzero
leading coefficient $d_2d_3z^2$, and every other term in
(I.141) has pole order at most one there. Hence the connector
determinant survives this substitution in every characteristic.

We now verify all strict internal interval ranks at this same
one-pole center. Every individual full-rank core remains
invertible; its determinant is a nonzero constant. Every
rank-two core keeps rank two, because some nonzero univariate
minor is evaluated at a transcendental coordinate. In the
second compound word, these are respectively full-rank and
rank-one factors. A rank-two interval has a nonzero compound
product. Factoring each rank-one compound by a nonzero
univariate pivot writes it as nonzero endpoint vectors times
the scalar connectors between consecutive rank-one factors.

If no or one full-rank factor lies between two such factors,
the connector depends on at most three coordinates and survives
by I-38's three-coordinate dominance. If two lie between them,
they form exactly the saturated run just analyzed. At its
adjacent two-dimensional cuts the compound connector is the
determinant (I.141), up to nonzero univariate pivot factors,
so it too survives. Thus every compound interval product that
was nonzero stays nonzero. A rank-three interval consists
entirely of invertible factors and stays invertible. All
strict internal interval ranks are therefore preserved
simultaneously. This asserts preservation of ranks, not of
every individual minor.

\paragraph{5. Endpoints and compatible constant gauges.}
Absorb the first and last factors into the affine row and column.
By (I.134), the first and last internal factors are rank two.
The polynomial row $A_1(X_1)A_2(X_2)$ is nonzero, since otherwise
the full scalar would be zero. Pairwise dominance preserves a
nonzero entry at the center. The same holds for
$A_{n-1}(X_{n-1})A_n(X_n)$. Propagation through the remaining
rank-preserved constant intervals is injective on the persistent
two-dimensional active spaces, so both endpoint active
projections are nonzero. No four-coordinate endpoint assertion
is being assumed.

The full constant interior product has rank two, and every
subinterval has rank two or three. The compatible-space proof
of I-38, part~4, applies to this rank profile without any bound
on run length: propagate a two-dimensional complement of the
full-product row kernel, and take the suffix preimages of a
final complementary line. Each passive line is either the
rank-two suffix kernel or the inverse image under an invertible
suffix. In the transported bases the active blocks are
$I_2$, the singular passive blocks are zero, and the invertible
passive blocks can be scaled to one.

The endpoint active vectors pair to $f(\boldsymbol\alpha)=0$.
Normalize them to $e_1^{\mathsf T}$ and $e_2$, and remove the
passive endpoint components by the upper and lower shears of
I-37. They leave the first and last singular constants fixed;
if those coincided the same identities would permit both
shears. Over $K=\overline E(t,z,y_0)$, with independent residual
variables $Y_i=X_i-\alpha_i$, the resulting scalar is
\[
 f(\boldsymbol\alpha+\mathbf Y)=
 (e_1^{\mathsf T}+Y_1p)
 \prod_{i=2}^{n-1}(C_i+Y_iN_i)
 (e_2+Y_nq),
 \quad C_i\in\{P,I_3\},\quad P=\operatorname{diag}(1,1,0).
 \tag{I.142}
\]
The $I_3$ runs have length at most two, and $C_2=C_{n-1}=P$.
Pole extraction from the assumed two-pole identity, exactly
as in I-37, gives $f_i(\boldsymbol\alpha)=0$ for every $i$.
It follows that
\[
 p_2=q_1=(N_i)_{12}=0.
 \tag{I.143}
\]
At each singular factor the coefficient of $Y_i$ in
$\det(P+Y_iN_i)=0$ also gives $(N_i)_{33}=0$.
No such assertion is made at invertible factors.

\paragraph{6. Three tangent directions kill the bandwidth-three Hessian.}
Let $H=(f_{ij}(\boldsymbol\alpha))$, and put
$h_i=1/(a_i-y_0)$. Differentiate the already zero vector
$(f_i(\boldsymbol\alpha))$ with respect to $t,z,y_0$.
Equivalently take first coefficients under separate parameter
increments. Since $z\ne0$ in $K$, this gives
\[
 H\mathbf1=H\mathbf h=H\mathbf h^{\circ2}=0.
 \tag{I.144}
\]
All three identities hold in every characteristic. The diagonal
of $H$ is zero by multilinearity.

If a singular position lies strictly between $i$ and $j$, the
intervening constant product contains $P$. It removes the
passive state, while the active derivative blocks in (I.143)
are lower triangular, so $H_{ij}=0$. Otherwise all those
constants are identities, and $H_{ij}$ is exactly the product
of the state-one-to-passive derivative entry at $i$ and the
passive-to-state-two derivative entry at $j$. The endpoint
versions are $p_3$ and $q_3$. These statements follow directly
from (I.142); the other active contributions contain a zero
$(1,2)$ entry. Since no invertible run has length three,
\[
 H_{ij}=0\qquad\text{when }|i-j|>3.
 \tag{I.145}
\]

The first row of $H$ has at most three remaining entries.
For three distinct future indices $j,k,\ell$, the three
equations (I.144) have coefficient determinant
\[
 \det\begin{pmatrix}
 1&1&1\\h_j&h_k&h_\ell\\h_j^2&h_k^2&h_\ell^2
 \end{pmatrix}
 =(h_k-h_j)(h_\ell-h_j)(h_\ell-h_k)\ne0.
 \tag{I.146}
\]
The one- and two-entry cases use the corresponding initial
rows of this Vandermonde matrix. Thus the first row is zero.
By symmetry, the earlier entries of the next row are now zero,
leaving at most three future entries to which the same argument
applies. Induction proves $H=0$.

Finally expand (I.142) in state paths. Direct passage from
active state one to active state two is excluded by (I.143).
Before a path first reaches active state two, take its last
departure from active state one into the passive state; if it
starts passively, take its passive endpoint entry instead.
Any intervening passive stays must occur at invertible factors,
because a singular factor has both its passive constant and
linear entry zero. The product of the departure and arrival
entries is precisely the corresponding $H_{ij}=0$.
Additional passive stays only multiply this zero product.
The same reasoning applies to a final output through $q_3$;
an output through active state one is excluded by $q_1=0$.
Therefore every output path is zero, contradicting the
nonzero polynomial $f$. This proves (I.132) under (I.131).

\paragraph{7. All multilinear inputs on at most ten variables.}
I-40 already covers at most nine variables. For ten variables,
the established balanced-ideal degree interval is $3\le d\le7$.
The cubic theorem of P-17 excludes degree three, and reversal
excludes degree seven. Consider a homogeneous quartic
annihilator of width at most three. Its common translation has
the form $f+s g$ because the ideal has no components of degree
at most two. A nonzero $g$ would be a cubic annihilator with
the same width bound: as in I-40, it is the leading coefficient
of a product of invertible local translations, and the
coefficient-cut four-minors remain zero at the leading limit.
P-17 forces $g=0$. The quartic is translation-invariant of
weight $10-2\cdot4=2$, so the theorem just proved excludes it.
Reversal also excludes degree six.

For a homogeneous quintic annihilator, every positive-power
translation coefficient has degree at most four and remains
in the balanced ideal. If any were nonzero, the coefficient
of the highest occurring power would retain width at most
three by the same limit argument. The already excluded cubic
and quartic cases, and the absence of lower degrees, rule this
out. The quintic is therefore translation-invariant of weight
zero, and I-40 excludes it as well. This covers every allowed
homogeneous degree. Passing to a lowest nonzero homogeneous
component, which preserves the ROABP width in the same order,
proves the general multilinear ten-variable assertion.

The finite balanced sampler is exactly I-40, part~6: multiply
the pullback by $\prod_i(a_i-y_1)(a_i-y_2)$, interpolate its two
selector degrees at most $n$ away from every node, and its
three amplitude parameters of total degree at most $D$ on
the Newton simplex. This proves (I.133) and its field scope
for both classes.\hfill$\square$

\paragraph{Provenance and the next obstruction.}
The proof uses the established representation input I-B09 and
the fully proved I-37--I-40 normalization and rank arguments.
The purity agent independently checked the canonical
$\Gamma^2$ anchors, the actual dual pairing in characteristic
two, the connector determinant and the endpoint simplification.
Those checks do not assume that $\Gamma^2V$ is self-dual or
that constant full-rank gauges intertwine the entire Borel
action. The present result makes no historical priority claim.

At weight three, I-40 permits an invertible run of length three.
The three tangent directions in (I.144) do not by themselves
kill a bandwidth-four symmetric matrix. Shorter nonsaturated
runs also need not have both invariant bounding support spaces.
The general mixed affine two-pair gate, the unrestricted
eleven-variable case, and polynomial dependence on a growing
ROABP width remain separate questions.


\subsection{I-42. Reset Hessians, sharp block thresholds, and saturated runs}
\label{I-42}

\paragraph{Scope.}
The abstract Hessian examples below are not ROABP annihilators.
The positive polynomial consequence retains its explicit
saturated-run promise.

\paragraph{1. An exact reset-Hessian lemma.}
Let $K$ be a field, let $h_1,\ldots,h_n$ be distinct elements,
and write $v_i=(1,h_i,h_i^2)^{\mathsf T}$. Let $S\subseteq[n]$
be a set of reset positions. Suppose a symmetric zero-diagonal
matrix $H$ has the following form for $i<j$:
\[
 H_{ij}=\begin{cases}
 a_i b_j,&\text{if no element of $S$ lies strictly between $i,j$},\\
 0,&\text{otherwise},
 \end{cases}
 \qquad a_i,b_i\in K.
 \tag{I.147}
\]
Let $t(i)$ be the first reset strictly after $i$, or $n$ if
there is none. Assume
\[
 H\mathbf1=H\mathbf h=H\mathbf h^{\circ2}=0.
 \tag{I.148}
\]
If $t(i)-i\le5$ for every $i<n$, then $H=0$ in every
characteristic. If $\operatorname{char}K\ne2$, the weaker
bound $t(i)-i\le6$ still implies $H=0$.

\paragraph{Proof.}
Suppose $H\ne0$, and let $i$ be its first nonzero row.
All earlier rows and, by symmetry, earlier columns are zero.
Put $t=t(i)$. Some future entry in this row is nonzero,
so $a_i\ne0$. Let $k$ be the first index in $(i,t]$ with
$b_k\ne0$. The three row equations imply
\[
 \sum_{j=k}^t b_jv_j=0.
 \tag{I.149}
\]
Any at most three of the $v_j$ are linearly independent.
Consequently at least four of these $b_j$ are nonzero, and
$t-k\ge3$. Under the first bound $t-i\le5$, this implies
$k-i\le2$.

Since $k\le t-3$, it is not the terminal reset. Every future
entry in its row is still in the same block. Earlier columns
before $i$ are zero, so its row equation is
\[
 b_k\sum_{j=i}^{k-1}a_jv_j+
 a_k\sum_{j=k+1}^t b_jv_j=0.
\]
Use (I.149) and divide by $b_k\ne0$ to obtain
\[
 \sum_{j=i}^{k-1}a_jv_j-a_kv_k=0.
 \tag{I.150}
\]
There are at most three distinct vectors here, so their
independence forces $a_i=0$, a contradiction.

Now suppose $\operatorname{char}K\ne2$ and only $t-i\le6$
is assumed. The same argument applies unless $k=i+3$.
In that remaining case $t=i+6$, and all four coefficients
$b_k,b_{k+1},b_{k+2},b_{k+3}$ are nonzero. Equation (I.150)
is a nonzero relation among four Vandermonde columns; every
coefficient is nonzero, since any three are independent.
In particular $a_k\ne0$, and
\[
 \sum_{j=i}^{k-1}a_jv_j=a_kv_k.
\]
The row at $\ell=k+1<t$ has no contribution beyond this same
block. Substitute the preceding identity and (I.149) into its
equation. This gives
\[
 (2a_kb_\ell-a_\ell b_k)v_k-a_\ell b_\ell v_\ell=0.
 \tag{I.151}
\]
The two columns are independent and $b_\ell\ne0$, so first
$a_\ell=0$ and then $2a_kb_\ell=0$, a contradiction outside
characteristic two. This proves both bounds.

\paragraph{2. Sharpness for abstract Hessians.}
Both thresholds are sharp for matrices of the stated form,
provided the field contains the displayed number of distinct
nodes. In characteristic two, take seven distinct $h_i$.
Choose nonzero vectors $a$ and $b$, supported respectively on
$\{1,2,3,4\}$ and $\{4,5,6,7\}$, with
\[
 \sum_{i=1}^4a_iv_i=0,
 \qquad\sum_{i=4}^7b_iv_i=0.
 \tag{I.152}
\]
Each kernel is one-dimensional and all its displayed entries
are nonzero, by the three-column Vandermonde property. Define
$H_{ij}=a_ib_j$ for $i<j$, impose symmetry, and put $H_{ii}=0$.
The first three rows satisfy (I.148) by the second relation;
the last three do so by the first. The fourth row is
$-2a_4b_4v_4=0$. Thus $H\ne0$ has all three kernels. It is
one reset block of seven positions, with $S=\{1,7\}$, corresponding to five
intermediate invertible positions in the matrix-word setting.

In every characteristic, use eight distinct nodes and take
$a$ supported on $\{1,2,3,4\}$ and $b$ on
$\{5,6,7,8\}$, each annihilating its four-column Vandermonde
matrix. Then
\[
 H=ab^{\mathsf T}+ba^{\mathsf T}
 \tag{I.153}
\]
is a nonzero symmetric zero-diagonal matrix, with $S=\{1,8\}$, of the same
semiseparable form and satisfies (I.148). This proves the
outside-characteristic-two threshold is also exact. These
examples concern an abstract evaluated Hessian; no fixed
polynomial integrability or ROABP realization is asserted.

\paragraph{3. Consequence for an actual rank-preserved width-three word.}
Suppose a nonzero affine multilinear scalar word has at least
one internal rank-two factor and has, at the common
one-pole center $\alpha_i=t+z/(a_i-y)$, all the internal rank
and nonzero active-endpoint properties used in I-41. Assume
its nonzero generic internal interval ranks are two or three,
and that each invertible run has length at most four, or at
most five if the characteristic is not two. Then its nonzero
scalar cannot annihilate the balanced two-pole map.

Indeed, an assumed annihilation gives the same normalized
$I/P$ word, zero first derivatives and passive-zero entries
at singular factors as (I.142)--(I.143). If an endpoint is
followed or preceded by an invertible constant string, transport
its endpoint shear unchanged through that $I_3$ string to the
first or last $P$, exactly as in I-38; this leaves every
constant $I_3$ and $P$ intact. Its Hessian has exactly
the form (I.147), with resets at singular factors and the
endpoint versions of $a_i,b_i$. The distance to the next reset
or final endpoint is at most the invertible-run length plus
one. The three tangent identities give (I.148), so part~1
forces $H=0$. The full state-path argument of I-41 then kills
every possible passive excursion and every output path,
contradicting nonzero scalar output. Thus no additional
Hessian condition remains under these run bounds once the
specified ranks and endpoints have been preserved.

\paragraph{4. A positive saturated-run class beyond weight two.}
Here is one explicit structural application that supplies those
rank hypotheses rather than assumes them. Let $f$ be an
irreducible homogeneous multilinear polynomial, depending on
all its $n$ variables, of width at most three in some order,
and invariant under common translation. Let
$m=n-2\deg f$. Suppose either $m\in\{3,4\}$ in any
characteristic, or $m=5$ in characteristic different from two.
Construct the minimal affine cores from the actual coefficient
support spaces, as in I-40. Assume that every maximal run of
rank-three cores is either isolated or has length exactly $m$.
Then balanced two-pole evaluation hits $f$, in unknown order.
Its finite grid is the same
$(n+1)^2\binom{D+3}{3}$ grid as I-40, with the same field
and extension-input conditions.

To prove this, generic strict internal ranks are at least two
by irreducibility, and the first and last internal factors
are singular by minimality, exactly as in I-41. Isolated
invertible cores give connectors in at most three coordinates,
so I-38's dominance preserves them. Consider a length-$m$
run beginning after cut $j$. Every full-rank core increments
the determinant weight by three, and its entire allowed
budget is $3m$. Hence
\[
 \delta_j=0,\qquad\delta_{j+m}=3m,
 \qquad\delta'_{j+m}=0.
 \tag{I.154}
\]
Its two bounding support spaces are actual three-dimensional
$\mathrm{SL}_2$ submodules. Since $\delta_j=0$, the index
$j$ is even; since $n-m=2D$, the right bounding suffix also
has even length. The adjacent support spaces therefore have
dimension two and odd weights. The argument of I-41,
part~3, excludes all weights other than $\pm1$ and identifies
both bounding three-planes with the actual submodule
$\Gamma^2V\subset V\otimes V$.

Normalize the $m$ full-rank constant cores to the identity.
Each linear coefficient has nonzero entries only on the first
superdiagonal by torus covariance. Their product is still
\[
 M=\begin{pmatrix}1&A&B\\0&1&C\\0&0&1\end{pmatrix},
 \tag{I.155}
\]
where $A,C$ are linear and $B$ quadratic in all $m$ run
variables: a product of three strictly upper triangular
three-by-three matrices is zero. The final left weights are
$(m+2,m,m-2)$ and the right three-plane weights are
$(2,0,-2)$. Semi-invariance of weight $m$ gives the same
nonzero antidiagonal pairing and the same right anchor
as (I.140). Thus the complete connector still has determinant
(I.141), with nonzero constant coefficient $d_2d_3$ on the
square of the left-anchor variable.

Its unique double pole preserves this connector, independently
of how many coordinates occur in $A,B,C$. The other connectors
and univariate pivots survive as above. Therefore every strict
internal interval rank is preserved at the one-pole center.
The minimal endpoint ranks again reduce their active
nonvanishing to pairwise dominance. The maximum run is at
most four in the unrestricted-characteristic cases and at
most five otherwise. Part~3 now proves the asserted scalar
nonvanishing and the claimed finite sampler.

\paragraph{Exact remaining rank gate.}
At weight three, the only missing run length for this
structural theorem is two; at weight four it is two or three.
The reset-Hessian lemma shows that these cases need only a
rank-preservation repair, not an additional bandwidth argument.
The theorem does not claim that the saturated-run promise
survives arbitrary factorization or specialization. It is
stated for an irreducible input and its actual minimal cores.
Products of inputs already known to be hit also remain hit
by the same formal map, since its target is a field, but no
unproved structural promise for their factors is inferred.

The overlapping seven-node characteristic-two example and the
disjoint eight-node example are barriers to extending this
specific Hessian-kernel argument without further information.
They do not disprove a balanced two-pair sampler or a general
polynomial-in-width hitting set. No historical priority claim
is made.


\subsection{I-43. Grassmann lines repair weight-three connectors}
\label{I-43}

\paragraph{Statement.}
Let $E$ be any field, and let $a_1,\ldots,a_n$ be distinct
nodes in $E$. Every nonzero homogeneous multilinear polynomial
$f$, having ROABP width at most three in some variable order
and satisfying
\[
 f(X+s\mathbf1)=f(X),\qquad 0\le n-2\deg f\le3,
 \tag{I.156}
\]
has nonzero pullback under the balanced two-pole map
\[
 R_i=t+\frac{z_1}{a_i-y_1}+\frac{z_2}{a_i-y_2}.
 \tag{I.157}
\]
The order need not be known, and the assertion holds in every
characteristic. Together with the arbitrary-variable cubic
theorem P-19, this proves the same conclusion for every
nonzero multilinear width-at-most-three polynomial on at most
eleven variables, without a homogeneity or translation promise.

For total degree at most $D$, the finite hitting set in both
assertions has at most
\[
 (n+1)^2\binom{D+3}{3}
 \tag{I.158}
\]
points over a working field of size at least
$\max\{2n+1,D+1\}$, with the extension-input convention of
I-40. The new conclusion is for the balanced rational map
with common center $t\mathbf1$. It does not assert the same
result for arbitrary center slopes or for the ordinary
unweighted centered-$\mathrm{SV}_2$ map.

\paragraph{Background input I-B10.}
We use the established description of a projective line in a
Grassmannian: its three-planes all contain a fixed two-plane
and lie in a fixed four-plane. The exact exterior-algebra
input is Ghorpade, Patil, and Pillai (2009), Lemma~3, printed
p.~4 of the author/arXiv version: two independent decomposable
$\ell$-vectors have decomposable sum exactly when their
associated $\ell$-planes intersect in dimension $\ell-1$.
That statement and proof are over an arbitrary field. The
highest-weight-zero and highest-weight-one inputs are the
degree-zero and degree-one cases of I-B09, with the conventions
and elementary Bruhat audit already recorded in I-40.

\paragraph{1. The precise connector problem.}
Work over an algebraic closure for the proof. The factor and
minimal-core reduction of I-41 applies unchanged: homogeneous
multilinear factors use disjoint variable sets, inherit the
width by specializing complementary factors, and each is
translation-invariant. Their nonnegative weights add to the
weight on the used coordinates. Thus a counterexample to
(I.156) would have an irreducible factor of weight at most
three. I-41 covers weights at most two. We may assume that
$f$ is irreducible, depends on all its variables, and has
weight $m=n-2D=3$.

Construct its actual minimal affine cores from the coefficient
support inclusions, as in I-40, and pad consistently to width
three. Every proper bond rank is at least two, and the first
and last are exactly two. Every strict internal interval has
generic rank two or three; a rank-one polynomial outer
factorization would factor the scalar into two nonconstant
polynomials, as in I-37. In particular the first and last
internal factors are singular. I-40 bounds every run of
rank-three factors by three.

It suffices first to preserve the rank-two contractions
between consecutive singular factors at the common one-pole
center
\[
 \alpha_i=t+\frac{z}{a_i-y}.
 \tag{I.159}
\]
Here the two singular factors themselves are included in the
connector. Consecutive singular factors with no intervening
full-rank factor, or with one, involve at most three distinct
coordinates; I-38's all-characteristic triple dominance
preserves their nonzero compound matrices. A run of length
three is saturated and is covered by I-42, part~4.

Consider therefore two full-rank cores $A_{j+1},A_{j+2}$,
between singular anchors $A_j,A_{j+3}$. Let $\delta_j$ be the
determinant weight of the left support $L_j$. The two increments
of three and the total budget nine give
\[
 \delta_j\in\{0,1,2,3\},\qquad
 \delta'_{j+2}=3-\delta_j.
 \tag{I.160}
\]
Reversing the witnessing variable order exchanges left and
right support spaces and transposes this connector. Thus it
is enough to handle $\delta_j=0$ and $\delta_j=1$. This is
order reversal, not reciprocal multilinear reversal, and
does not change the generator.

\paragraph{2. A single invariant anchor is enough.}
Suppose $\delta_j=0$. The argument of I-41 identifies the
left support with the actual invariant submodule
$\Gamma^2V\subset V\otimes V$, its preceding support with
the standard two-dimensional module, and its anchor with
\[
 C(u)=\begin{pmatrix}1&u&0\\0&1&u\end{pmatrix},
 \qquad u=X_j.
 \tag{I.161}
\]
Normalize the two full-rank constant cores to the identity by
constant basis changes. Torus covariance gives
\[
 M(x_1,x_2)=
 \begin{pmatrix}1&A&B\\0&1&C\\0&0&1\end{pmatrix},
 \tag{I.162}
\]
where $A,C$ are linear in the two internal variables and
$B$ is quadratic. Write the arbitrary opposite affine anchor
as $D(v)$, with rows $D_1(v),D_2(v),D_3(v)$, where
$v=X_{j+3}$. Its generic matrix rank is two. Its three rows
are linearly independent as coefficient tensors in
$E[v]^{r_{j+3}}$: a constant relation between them would be
a relation between the three actual suffix coefficient
functions at the preceding minimal bond. The constant gauges
above preserve this independence.

The wedge of the two rows of $C(u)M D(v)$ is the row of its
second compound. Its coefficient on $u^2$ is exactly
\[
 D_2(v)\wedge D_3(v).
 \tag{I.163}
\]
If this polynomial vector is nonzero, substitute (I.159) and
consider the pole at $y=a_j$. Every other coordinate is
regular there. The coefficient of the double pole is a
nonzero scalar multiple of
$z^2[D_2(v_0)\wedge D_3(v_0)]$, where
$v_0=t+z/(a_{j+3}-a_j)$. This vector is nonzero because
$v_0$ is transcendental over the coefficient field. Hence
this connector retains rank two.

It remains to treat $D_2\wedge D_3=0$. A primitive rank-one
factorization over the univariate polynomial ring gives
\[
 D_2(v)=p_2(v)q(v),\qquad D_3(v)=p_3(v)q(v).
\]
Both rows are nonzero and linearly independent over $E$.
They are affine. If the primitive vector $q$ had positive
degree, the degree bound would force both $p_2,p_3$ to be
constants, a contradiction. Thus $q$ is a nonzero constant
row, $p_2,p_3$ are independent affine polynomials, and
$D_1(v)\wedge q$ is nonzero because $D$ has generic rank two.
Direct multiplication now gives the complete exterior row
\[
 \bigl[p_2(v)+(C+u)p_3(v)\bigr]
       \bigl[D_1(v)\wedge q\bigr].
 \tag{I.164}
\]
At $y=a_j$ its first factor has a nonzero simple-pole
coefficient proportional to $z p_3(v_0)$, whereas the second
is regular and nonzero. The terms $A,B$ cancel from this
exterior expression, so no other internal dependence can
cancel that pole. Again the connector retains rank two.
This proves the case $\delta_j=0$, without assuming that the
opposite anchor is invariant or self-dual.

\paragraph{3. Determinant weight one forces two equal positive weights.}
Suppose now that $\delta_j=1$. The nonzero Pl\"ucker vector
of $L_j$ is fixed by the upper unipotent and has torus weight
one. Its $\mathrm{SL}_2$ orbit spans the standard
two-dimensional module, by the highest-weight-one input.
Its projective orbit is therefore a line, all of whose
points are decomposable three-vectors. I-B10 gives fixed
spaces
\[
 K\subset L_j\subset T,\qquad \dim K=2,\quad \dim T=4.
 \tag{I.165}
\]
They are the intersection and span of all planes in that
orbit line, so both are $\mathrm{SL}_2$-stable.

We give the needed small-width consequence explicitly; it
does not assume semisimplicity in positive characteristic.
The parity of tensor weights and $\delta_j=1$ imply that
$j$ is odd. Thus the two torus weights of the stable space
$K\subset V^{\otimes j}$ are $q,-q$ for an odd positive
integer $q$, by Weyl symmetry. Take all contractions of its
last physical factor. Their span $W_K$ is an
$\mathrm{SL}_2$-stable subspace of the preceding actual
support $L_{j-1}$, and hence has dimension at most three.

Assume $q>1$, and choose weight vectors $k_+,k_-$ of $K$.
Their coefficient vectors on the last factor can have only
the four distinct weights
\[
 q-1,\ q+1,\ -q-1,\ -q+1.
 \tag{I.166}
\]
Weyl symmetry pairs the first with the last and the second
with the third. Nonzero coefficient vectors of distinct
weights are independent. Since $\dim W_K\le3$, exactly one
of those pairs can occur. Up to exchanging $e_0,e_1$, this
means
\[
 k_+=p\otimes e_0,\qquad k_-=r\otimes e_1,
 \tag{I.167}
\]
where $p,r$ are independent, or the same display with the
two final basis vectors exchanged.

Apply the unipotent that sends $e_0$ to $e_0+s e_1$ in
(I.167), and write its action on $p$ as $\rho(s)p$.
Since $K$ is stable, the resulting tensor
$\rho(s)p\otimes e_0+s\rho(s)p\otimes e_1$ must remain in
the span of the two displayed tensors. Its $e_0$ coefficient
forces $\rho(s)p$ into $E(s)p$; its $e_1$ coefficient then
forces the same nonzero vector into $E(s)r$, a contradiction.
For the exchanged case use the opposite unipotent. This
argument works over the rational function field in $s$ and
in every characteristic. Hence $q=1$.

The determinant character of the stable space $K$ is zero.
The quotient line $L_j/K$ therefore has weight one. The
stable two-dimensional quotient $T/K$ has weights $1,-1$,
because a rational $\mathrm{SL}_2$-module has trivial
determinant character. Consequently the weights on $L_j$ are
\[
 (1,1,-1).
 \tag{I.168}
\]
No splitting of either short exact sequence in (I.165) was
used, and this deduction excludes Frobenius-twisted larger
weights directly.

\paragraph{4. A unique internal pole detects the remaining connector.}
Choose a weight basis as in (I.168) and again normalize the
two full-rank constant cores to the identity. Torus covariance
allows their linear matrices only in the common
negative-to-positive weight step. They thus have the form
\[
 N_\ell=v_\ell e_3^{\mathsf T},\qquad
 (v_\ell)_3=0,\qquad N_\ell N_k=0.
 \tag{I.169}
\]
Their product is
$M=I+x_1N_1+x_2N_2$. Its second compound is affine in
$x_1,x_2$: the quadratic exterior term of a rank-one update
vanishes. Multiplying by the compounds of the two arbitrary
affine anchors gives
\[
 C_2\bigl(A_j(u)M(x_1,x_2)A_{j+3}(v)\bigr)
   =P_0(u,v)+x_1P_1(u,v)+x_2P_2(u,v).
 \tag{I.170}
\]
Here the $P_\ell$ are polynomial matrices and the displayed
matrix is not identically zero, because the original
connector has generic rank two.

If some $P_\ell$ with $\ell=1,2$ is nonzero, take the pole
at the node belonging to $x_\ell$. Every other summand is
regular. The coefficient of its simple pole is a nonzero
scalar multiple of $zP_\ell(u_0,v_0)$, where
\[
 u_0=t+\frac{z}{a_j-a_{j+\ell}},\qquad
 v_0=t+\frac{z}{a_{j+3}-a_{j+\ell}}.
 \tag{I.171}
\]
The two coefficients of $z$ are distinct, so $(u_0,v_0)$
is an invertible linear change of $(t,z)$. A nonzero
polynomial matrix $P_\ell$ therefore stays nonzero. If both
$P_1,P_2$ vanish, the nonzero matrix $P_0(u,v)$ survives by
pairwise dominance. This proves preservation in the
$\delta_j=1$ case and, by order reversal, in every case of
(I.160).

\paragraph{5. Simultaneous ranks, endpoints, and the scalar output.}
Every univariate singular-factor pivot is nonzero after
(I.159), since each $\alpha_i$ is transcendental. The
preceding arguments preserve the nonzero rank-two compound
of every connector between consecutive singular factors.
All these assertions hold as rational-function identities
in the same independent parameters $t,z,y$. They therefore
hold simultaneously; no independent seed is assigned to an
individual connector.

For clarity, use the same pivot-cleared rank factorizations
as I-36: a rank-two anchor is $U_iV_i$, with both factors
of full respective rank. A connector of rank two says that
the $2\times2$ contraction between its adjacent anchors is
invertible. Products of these invertible contractions
preserve every interval containing anchors; intervals made
only of full-rank cores remain invertible. Thus all strict
internal interval ranks are preserved. This deduction uses
the common evaluated pivots, not preservation of only
constant boundary functionals.

Minimality gives singular first and last internal factors.
Pairwise dominance preserves the nonzero endpoint/first-anchor
and last-anchor/endpoint contractions. The preserved internal
rank-two identifications propagate their nonzero active
projections, exactly as in I-41. Hence all the hypotheses of
I-42, part~3, are now proved at the same one-pole center.
The maximum invertible run is three, which is inside its
all-characteristic bound four. An assumed annihilation of
(I.157) gives the reset Hessian with its three tangent kernels;
I-42 kills that Hessian and then every output path. This
contradicts $f\ne0$. The factor reduction in part~1 proves
the full assertion (I.156).

\paragraph{6. All multilinear inputs on at most eleven variables.}
I-41 already handles at most ten variables. Suppose a
homogeneous width-three annihilator on eleven variables
exists. P-19 and the known minimum balanced degree leave
only degrees $4,5,6,7$. For degree four, every nonconstant
coefficient of common translation has degree at most three.
If any were nonzero, the coefficient of the highest power
of the translation parameter would retain width at most
three by the coefficient-cut leading-limit argument of I-40.
It would contradict P-19 or the absence of smaller degrees.
Thus the quartic is translation-invariant, has weight
$11-8=3$, and is excluded by the theorem just proved.
Multilinear reversal excludes degree seven.

For degree five, the same highest-translation-coefficient
argument would produce a width-three annihilator of degree
at most four, all of which have now been excluded. The
quintic is therefore translation-invariant of weight one,
and I-40 excludes it. Reversal also excludes degree six.
Passing to a lowest nonzero homogeneous component preserves
the witnessing width and proves the assertion for arbitrary
multilinear inputs on at most eleven variables.

Finally apply the denominator-safe interpolation in I-40:
after multiplication by
$\prod_i(a_i-y_1)(a_i-y_2)$, each selector degree is at most
$n$, and the total degree in the three amplitudes is at most
$D$. The selector grid outside all nodes and the amplitude
Newton simplex give exactly (I.158), with the stated field
and extension scope.\hfill$\square$

\paragraph{Provenance and remaining scope.}
The Grassmann-line fact is an established input, not a new
claim about Grassmannians. The new use here is its interaction
with the actual adjacent coefficient support of dimension at
most three and the resulting pole calculation for constrained
ROABP connectors. All endpoint, interval, and simultaneous
substitution assertions above concern the same parameter
field; none follows merely from preserving each local space
against constant rank-one functionals.

At weight four, determinant weights at the ends of an
unsaturated run can both exceed one. Those connectors are
not covered by the two classifications above. The unrestricted
balanced two-pair question, the next higher-degree finite-
variable cases, and the polynomial-in-growing-width isolation
gate remain separate. No historical priority claim is made.


\subsection{I-44. Torus-chain connectors and the next balanced degree windows}
\label{I-44}

\paragraph{Statement.}
Fix arbitrary distinct nodes $a_1,\ldots,a_n$ over a field
$E$, and use the balanced rational two-pole map of I-43.
The following assertions hold in an unknown witnessing
ROABP order.

First, let $f$ be irreducible, homogeneous, multilinear,
common-translation invariant, and of width at most three,
and remove its unused variables. Construct the actual minimal
affine cores from its coefficient support spaces. If every
run of rank-three cores has length at most five, the common
one-pole center $\alpha_i=t+z/(a_i-y)$ preserves all strict
internal interval ranks. This assertion holds in every
characteristic. Consequently balanced two-pole evaluation
hits $f$ whenever those runs have length at most four, or
at most five if $\operatorname{char}E\ne2$.

Second, without an irreducibility or run promise, every
nonzero homogeneous multilinear common-translation-invariant
width-at-most-three polynomial is hit when
\[
 0\le n-2\deg f\le4,
 \quad\hbox{or when}\quad
 0\le n-2\deg f\le5\ \hbox{and }\operatorname{char}E\ne2.
 \tag{I.172}
\]
Together with P-19, this proves the balanced two-pole
conclusion for all nonzero multilinear width-at-most-three
polynomials on at most twelve variables in every
characteristic, and on at most thirteen variables outside
characteristic two.

All the scalar conclusions have the same explicit finite
grid as I-40:
\[
 (n+1)^2\binom{D+3}{3},\qquad
 |E|\ge\max\{2n+1,D+1\},
 \tag{I.173}
\]
for total degree at most $D$, with evaluation over a stated
working extension if necessary. These are conclusions for
the balanced rational map with common slope one. No ordinary
centered-$\mathrm{SV}_2$ or arbitrary-slope promotion is made.

\paragraph{1. Torus spectra that need no anchor classification.}
Use I-40's actual minimal cores. Normalize the constant
matrices in any full-rank run to the identity. If the initial
three integer torus weights are $\lambda_1,\lambda_2,
\lambda_3$, the successive weight lists are obtained by
adding one to every entry. Each linear matrix $N_i$ has an
entry only from a column weight to a row weight two larger.

Unless the three weights are distinct and have the form
\[
 (q+2,q,q-2),
 \tag{I.174}
\]
all the allowed entries have a common one-dimensional image
or a common one-dimensional domain, or there is at most one
allowed entry. Indeed, with three distinct weights two
composable steps require exactly (I.174); without such a
chain there is at most one step. If a weight is repeated,
the only nonzero case has two levels of multiplicities one
and two. All products $N_iN_j$ then vanish, and their linear
span consists of matrices of rank at most one.

It follows that the run product is $I+\sum_i x_iN_i$ and
its second compound is affine in all run variables. After
multiplication by the two anchor compounds it has the form
\[
 P_0(u,v)+\sum_i x_iP_i(u,v).
 \tag{I.175}
\]
The original connector has rank two, so this polynomial
matrix is nonzero. The unique-internal-pole argument of
I-43, part~4, applies without a restriction on the number of
summands: at the node of $x_i$, the two anchor values are
independent linear forms in $t,z$. A nonzero $P_i$ therefore
has nonzero residue; if only $P_0$ remains, pair dominance
applies. This proves preservation for every spectrum other
than (I.174), in every characteristic. The differences of
weights here are integer characters of the torus, not
integers reduced modulo the characteristic.

\paragraph{2. All minimal singular anchors into a three-weight chain.}
Assume (I.174), in the displayed order. A preceding singular
anchor $A(u)$ has three columns and two or three rows. Both
its row and column coefficient tensors are independent over
$E$, by actual minimality, and its generic matrix rank is
two. Torus covariance permits preceding row weights only
from
\[
 q+3,\ q+1,\ q-1,\ q-3.
 \tag{I.176}
\]
Its columns have respectively the following possible
coefficient vectors in those row-weight spaces:
\[
 (u a,b,0,0)^{\mathsf T},\quad
 (0,u c,d,0)^{\mathsf T},\quad
 (0,0,u e,f)^{\mathsf T}.
 \tag{I.177}
\]

No preceding weight can have multiplicity two. At either
extreme of (I.176) all row coefficients use a single vector,
contrary to row independence. At weight $q+1$, multiplicity
two requires $b,c$ independent. A nonzero third column must
then use one of the two lower weights, and the resulting
three-column minor is nonzero, contrary to rank two.
The argument at weight $q-1$ is the same in reverse. Three
copies of a weight cannot be spanned by its at most two
coefficient vectors either.

Thus the preceding weights are distinct. The complete list
of possibilities, with rows ordered by decreasing weight,
is as follows. Every displayed scalar is nonzero, except
that in the middle two-row case $c,d$ may vanish individually
but not both.
\[
\begin{array}{c|c}
 \hbox{preceding weights}&A(u)\\ \hline
 (q+3,q-1)&
 \begin{pmatrix}a u&0&0\\0&d&e u\end{pmatrix}\\[3pt]
 (q+1,q-1)&
 \begin{pmatrix}b&c u&0\\0&d&e u\end{pmatrix}\\[3pt]
 (q+1,q-3)&
 \begin{pmatrix}b&c u&0\\0&0&f\end{pmatrix}\\[3pt]
 (q+3,q+1,q-1)&
 \begin{pmatrix}a u&0&0\\b&0&0\\0&d&e u\end{pmatrix}\\[3pt]
 (q+1,q-1,q-3)&
 \begin{pmatrix}b&c u&0\\0&0&e u\\0&0&f\end{pmatrix}
\end{array}
 \tag{I.178}
\]
To verify completeness, the other two-row subsets miss an
entire column. Among the four three-row subsets, the two
omitted subsets give a nonzero triangular determinant.
For the two retained three-row subsets, rank two forces
the middle diagonal entry to vanish; row and column
independence force all the entries shown as nonzero.
This proves the list directly from (I.177).

Over $E(u)$ the row spaces in (I.178) have one of three
forms:
\[
\begin{aligned}
 G_1(u)&=\begin{pmatrix}b&c u&0\\0&d&e u\end{pmatrix},
       &&bce\ne0,\quad d\in E,\\
 G_2(u)&=\begin{pmatrix}b&c u&0\\0&0&1\end{pmatrix},
       &&bc\ne0,\\
 B(u)&=\begin{pmatrix}1&0&0\\0&d&e u\end{pmatrix},
       &&de\ne0.
\end{aligned}
 \tag{I.179}
\]
The passage to these row spaces uses only full-rank
univariate factors and nonzero univariate denominators.
All remain valid after the one-pole substitution, since
$\alpha_i$ is transcendental over $E$.

\paragraph{3. Either of the first two anchor forms suffices.}
The normalized run product in the chain case is
\[
 M=\begin{pmatrix}1&A&B_0\\0&1&C\\0&0&1\end{pmatrix},
 \quad
 A=\sum_i p_ix_i,\quad C=\sum_i q_ix_i,\quad
 B_0=\sum_{i<j}p_iq_jx_ix_j.
 \tag{I.180}
\]
Here each individual core is
$I+x_i(p_iE_{12}+q_iE_{23})$; no quadratic term occurs
within a single core.

Let the opposite, unreduced minimal anchor have rows
$D_1(v),D_2(v),D_3(v)$. They are independent affine
coefficient tensors, and the matrix $D(v)$ has generic
rank two. If the left row space is $G_1$, the coefficient
of $u^2$ in the exterior row of $G_1MD$ is
$ce\,D_2\wedge D_3$. If this is nonzero, its unique double
pole preserves the connector. If $D_2\wedge D_3=0$,
I-43's primitive affine factorization gives
$D_2=p_2(v)q$, $D_3=p_3(v)q$, with $q$ constant and
$D_1\wedge q\ne0$. The full exterior row becomes
\[
 b\,[d p_2+(dC+e u)p_3]\,(D_1\wedge q).
 \tag{I.181}
\]
Its simple pole at the left-anchor node is nonzero.

For $G_2$ the exterior row is
\[
 b\,D_1\wedge D_3+(bA+c u)D_2\wedge D_3.
 \tag{I.182}
\]
If $D_2\wedge D_3\ne0$, use its unique simple pole in
$u$. Otherwise the expression is the nonzero univariate
vector $bD_1(v)\wedge D_3(v)$, which also survives.
These arguments do not bound the length of the middle run.
Order reversal supplies the same conclusion if the opposite
anchor has either of these two forms in the reversed word.

\paragraph{4. Two remaining anchor forms and a quadratic pole test.}
It remains to consider form $B$ at both ends. At the end
of a run of length $\ell$, the left weights are
$(q+\ell+2,q+\ell,q+\ell-2)$. Homogeneity pairs them
antidiagonally, with nonzero constants, with the three
distinct right weights. Apply (I.179) to the reverse anchor
and remove its full-rank univariate outside factor. After
constant column rescaling and ordering, the actual right
factor in the two-dimensional contraction is
\[
 \begin{pmatrix}\alpha v&0\\\beta&0\\0&1\end{pmatrix},
 \qquad \alpha\beta\ne0.
 \tag{I.183}
\]
This is the actual coefficient pairing; it does not assume
that a three-dimensional module is self-dual.

The two-by-two contraction is therefore
\[
 \begin{pmatrix}
  \alpha v+\beta A&B_0\\
  d\beta&dC+e u
 \end{pmatrix},
 \qquad
 F=(\alpha v+\beta A)(dC+e u)-\beta dB_0,
 \tag{I.184}
\]
where $\alpha,\beta,d,e$ are nonzero. In particular $F$ is
nonzero: its $uv$ coefficient is $\alpha e$.

We record the elementary quadratic restriction imposed by
one-pole annihilation. For any homogeneous quadratic in
distinct coordinate variables, a nonzero square coefficient
has a unique double pole at that coordinate's node.
Consequently an annihilating quadratic must be multilinear.
Write it as $\sum_{i<j}c_{ij}X_iX_j$, with symmetric
coefficients off the diagonal. The coefficient of $tz$ in
its one-pole pullback, and then the simple pole of its $z^2$
coefficient at $y=a_i$, give respectively
\[
 \sum_{j\ne i}c_{ij}=0,
 \qquad
 \sum_{j\ne i}\frac{c_{ij}}{a_j-a_i}=0.
 \tag{I.185}
\]
Two distinct possible neighbors give an invertible two-row
matrix in these equations. Thus every nonzero row of an
annihilating quadratic has at least three nonzero neighbors.
This derivation uses no factor of two and holds in every
characteristic.

Apply this to (I.184). The coefficient of $x_i^2$ is
$\beta d p_iq_i$. If any is nonzero, the connector already
survives. Otherwise define the disjoint sets
$P=\{i:p_i\ne0\}$ and $Q=\{i:q_i\ne0\}$.
In the resulting multilinear quadratic, a row indexed by
$i\in P$ has precisely its nonzero edge to $u$ and its
edges to $Q$ positions strictly before $i$. A row indexed
by $i\in Q$ has its edge to $v$ and its edges to $P$
positions strictly after $i$. This follows from
\[
 AC-B_0=\sum_{i>j}p_iq_jx_ix_j
 \quad\hbox{when all }p_iq_i=0.
 \tag{I.186}
\]

Suppose $\ell\le5$. Then at least one of $|P|,|Q|$ is
at most two. If either set is empty, $B_0=0$ and $F$ is
a product of two nonzero linear forms. Such a product
survives because distinct simple poles make the one-pole
map nonzero on every nonzero linear form.

If both sets are nonempty and $|Q|\le2$, (I.185) applied
to any $P$ row forces $|Q|=2$ and forces both $Q$ positions
to precede that row. Hence every $Q$ precedes every $P$.
This again makes $B_0=0$, so the same nonzero product
survives. If $|P|\le2$, apply the symmetric argument to
the $Q$ rows; again every $Q$ must precede every $P$.
Thus (I.184) cannot annihilate the one-pole map for
$\ell\le5$.

This proves that every constrained rank-two connector with
at most five intervening full-rank cores preserves rank.
All coefficient identities, univariate factors and residues
are in the same field $E(t,z,y)$.

\paragraph{5. Completion of the rank and scalar assertions.}
The irreducible/minimal reduction ensures that strict
internal interval ranks are two or three and that the first
and last internal factors are singular, exactly as in I-41.
The two cases of part~1 and parts~2--4 cover every torus
spectrum and every connector of length at most five.
Univariate pivots remain nonzero. Their evaluated rank
factorizations and the invertible two-by-two adjacent
contractions therefore preserve all strict interval ranks,
as proved in I-43, part~5. Pair dominance supplies the
nonzero active endpoint projections. All these assertions
hold simultaneously with the same one-pole parameters.
Small words with no such internal interval are directly
covered by pair or triple dominance.

I-42, part~3, now proves scalar nonvanishing when all runs
have length at most four, or at most five outside
characteristic two. This proves the first assertion.

For (I.172), factor a hypothetical annihilator and recompute
the chosen factor's actual minimal cores, as in I-41.
Each factor is homogeneous and translation-invariant, uses
a disjoint set of variables, and has nonnegative weight.
Its weight is bounded by the total weight. I-40 bounds its
full-rank runs by that weight, so the first assertion applies
to the required factors. The target of the evaluation map
is a field, so no product of these nonzero images can vanish.
This proves (I.172), without assuming that a run promise
for one nonminimal representation passes to its factors.

\paragraph{6. Twelve and thirteen variables.}
I-43 already handles at most eleven variables. On twelve
variables, P-19 and multilinear reversal leave only
homogeneous degrees four through eight. A quartic
annihilator is translation-invariant: a highest nonzero
positive translation coefficient would inherit width at
most three and have degree at most three, contradicting
P-19 and the known minimum degree. Its weight is four,
contradicting (I.172). Reversal excludes degree eight.
For degree five the same leading-translation argument
uses the already excluded lower degrees and produces
translation invariance of weight two; I-41 excludes it,
as well as degree seven by reversal. Degree six similarly
becomes translation-invariant of weight zero and is
excluded by I-40.

On thirteen variables, outside characteristic two, the
same induction starts with a quartic of weight five,
excluded by (I.172). A quintic becomes translation-invariant
of weight three and is excluded by I-43; a sextic has
weight one and is excluded by I-40. Reversal handles their
complementary degrees. In both cases a lowest nonzero
homogeneous component inherits the witnessing width,
so these homogeneous conclusions cover every multilinear
input. The interpolation proof of I-40 gives (I.173)
without a new degree estimate.\hfill$\square$

\paragraph{Remaining scope.}
The rank-preservation conclusion for five-core runs is
valid even in characteristic two. The scalar conclusion
there remains restricted to four-core runs because the
seven-node abstract Hessian of I-42 shows that its three
kernel equations alone are insufficient. In particular,
the weight-five characteristic-two case is reduced to its
possible saturated length-five runs, with their interval
ranks now preserved; no abstract Hessian is asserted to
integrate to a counterexample.

For unrestricted run length, the two remaining anchor forms
lead to the explicit quadratic (I.184). When both internal
types have at least three positions, the short-run argument
above no longer suffices. This is a constrained connector
problem, not a claim that general shared-parameter matrices
compose. The higher-degree and growing-width sampling gates
remain open. No historical priority claim or unperformed
numerical verification is made.

% END SOURCE work/isolation/findings.tex

% BEGIN SOURCE work/purity/findings.tex
% Purity/zero-status contribution, run CUE_GCT-20261121-purity.
% New proved deductions, not claims of priority over the literature.
\section{Purity, evaluation maps, and the precise derandomization gate}

\subsection{P-01: a quantitative Gram criterion}
\paragraph{Statement P-01.}
Let $\Omega$ be a finite probability space, let $\phi_1,\ldots,\phi_s:\Omega\to\mathbb C$, and suppose
\[
 \mathbb E|\phi_j|^2\ge a>0,\qquad
 |\mathbb E\phi_i\overline{\phi_j}|\le b\quad(i\ne j).
\]
Then every $c\in\mathbb C^s$ satisfies
\[
 \mathbb E\left|\sum_jc_j\phi_j\right|^2
 \ge \{a-(s-1)b\}\sum_j|c_j|^2.
\]
Consequently $a>(s-1)b$ implies that the coefficient-to-evaluation map $\mathbb C^s\to\mathbb C^\Omega$ is injective and that every nonzero combination is nonzero at a sampled point.
\paragraph{Proof.}
Expand the square. The diagonal sum is at least $a\sum_j|c_j|^2$, and the absolute value of the off-diagonal sum is at most $b\sum_{i\ne j}|c_i||c_j|$. Cauchy--Schwarz gives $(\sum_j|c_j|)^2\le s\sum_j|c_j|^2$, so the off-diagonal sum is at most $b(s-1)\sum_j|c_j|^2$. This proves the inequality. If a nonzero combination vanished at every sampled point, the left side would vanish while the right side would be positive, a contradiction. $\square$

For multilinear $f=\sum_{S\in\mathcal S}c_Sx_S$ on the sign cube and an $\varepsilon$-biased distribution (meaning $|\mathbb E x_T|\le\varepsilon$ for nonempty $T$), take $\phi_S=x_S$. Then $a=1$, $b=\varepsilon$, because $x_Sx_T=x_{S\triangle T}$. Thus $\varepsilon(|\mathcal S|-1)<1$ is a sufficient hitting criterion. The sharper coefficient-dependent criterion is
\[
 \mathbb E|f|^2\ge (1+\varepsilon)\|c\|_2^2-\varepsilon\|c\|_1^2>0.
\]
This criterion depends on coefficient energy or sparsity; ROABP width alone is not such a bound.

\subsection{P-02: an explicit purity-based hitting set in the sparse regime}
\paragraph{Statement P-02.}
For positive integers $n,s$, choose an odd prime $p>16n^2s^2$. Let $\chi$ be the quadratic character of $\mathbb F_p$, and put $b_i(t)=\chi(t+i)$ when $t+i\ne0$, and $b_i(t)=1$ otherwise, for $1\le i\le n$. The multiset
\[
 H_{n,s}=\{(b_1(t),\ldots,b_n(t)):t\in\mathbb F_p\}\subset\{\pm1\}^n
\]
is a hitting set for every nonzero $s$-sparse multilinear polynomial over any field of characteristic zero. Taking the first such prime gives cardinality $O(n^2s^2)$ and an elementary deterministic construction in time polynomial in $n,s$.
\paragraph{Proof.}
The quadratic-character construction and its bias bound are AGHP, \S4, Theorem~1, Construction~2, and Proposition~2 (P-B01). For completeness, the short form required here is as follows. For nonempty $T\subseteq[n]$, the squarefree polynomial $\prod_{i\in T}(t+i)$ has $|T|$ distinct roots since $p>n$. Weil's character-sum bound gives a sum of absolute value at most $(|T|-1)\sqrt p$. Replacing each zero character value by $1$ changes at most $|T|$ terms by at most $1$. Therefore
\[
 \left|\frac1p\sum_t\prod_{i\in T}b_i(t)\right|
 \le \varepsilon:=\frac{n-1}{\sqrt p}+\frac np.
\]
Our choice gives $\varepsilon(s-1)<1$, so P-01 proves the assertion over $\mathbb C$. For any set of at most $s$ monomials the evaluation matrix has integer entries; injectivity over $\mathbb C$ implies an integer maximal minor is nonzero, hence it remains nonzero over every characteristic-zero field. Bertrand's theorem supplies a prime less than twice the chosen threshold (enlarging an endpoint by one if necessary). Trial division and Euler's criterion already yield polynomial time in the numerical parameters $n,s$; no GRH is needed. $\square$

The word ``sparse'' is essential. This applies to sparse ROABPs, but an $n$-variate width-one ROABP can have $2^n$ nonzero coefficients.

\subsection{P-03: two small-width counterexamples and the attempted repair}
\paragraph{Statement P-03(a).}
Over a field of characteristic different from $2$, a subset $H\subseteq\{\pm1\}^n$ hits all multilinear width-one ROABPs if and only if $H=\{\pm1\}^n$.
\paragraph{Proof.}
If $a\notin H$, the nonzero polynomial
\[
 f_a(x)=\prod_{i=1}^n(1+a_ix_i)
\]
is a width-one ROABP. For $x\in\{\pm1\}^n$ it is nonzero exactly when $x=a$, where its value is $2^n$. Thus $H$ misses it. Conversely the sign-cube evaluation matrix for multilinear monomials is the tensor product of $\begin{pmatrix}1&1\\1&-1\end{pmatrix}$ and is invertible when $2$ is invertible. Hence the full cube hits all nonzero multilinear polynomials. $\square$

\paragraph{Statement P-03(b).}
Allowing a common affine transformation $x_i=u+vb_i(t)$ of every sign point does not repair the construction for general ROABPs. The polynomial
\[
 V(x_1,x_2,x_3)=(x_2-x_1)(x_3-x_1)(x_3-x_2)
\]
is a nonzero width-three ROABP of individual degree two and vanishes at all such points, even when $u,v$ range over the entire coefficient field.
\paragraph{Proof.}
Among three coordinates taking values in $\{u-v,u+v\}$, two coincide. Thus one factor vanishes. Nonzeroness follows from its displayed factorization in a polynomial ring over a field. The explicit ROABP is
\[
 [1\quad x_1\quad x_1^2]
 \begin{pmatrix}
 x_2&x_2^2&0\\-1&0&x_2^2\\0&-1&-x_2
 \end{pmatrix}
 \begin{bmatrix}x_3^2\\-x_3\\1\end{bmatrix}=V.
\]
Expansion verifies this identity in $\mathbb Z[x_1,x_2,x_3]$, so it holds in every characteristic. $\square$

These are refutations of particular proposed transfers, not refutations of polynomial-size ROABP hitting sets over a larger field or of weight generators with richer output alphabets.

\subsection{P-04: curve sampling after a surviving restriction}
\paragraph{Statement P-04.}
Let $C/\mathbb F_q$ be a smooth projective geometrically connected curve of genus $g$, let $r_i\in\mathbb F_q(C)$, and choose effective divisors $D_i$ containing their polar divisors. For multidegrees $d_i$, put $D=\sum_i d_iD_i$ and let $B$ be the union of supports of the $D_i$. If
\[
 q+1-2g\sqrt q-\#B(\mathbb F_q)>\deg D,
\]
then evaluations at $C(\mathbb F_q)\setminus B$ hit every polynomial $f$ of individual degrees at most $d_i$ whose pullback $f(r_1,\ldots,r_n)$ is nonzero in $\mathbb F_q(C)$.
\paragraph{Proof.}
Every monomial in the pullback has poles bounded by $D$, and the same is true of their sum. If the pullback is a nonzero constant it never vanishes. Otherwise its principal divisor has degree zero, so the degree of its zero divisor equals that of its pole divisor and is at most $\deg D$. In particular it vanishes at at most $\deg D$ rational points outside $B$. Weil's bound (P-B02) supplies at least $q+1-2g\sqrt q-\#B(\mathbb F_q)$ available rational points, which is larger. $\square$

This separates the two requirements: purity supplies enough rational points; it does not prove that the restriction to $C$ is nonzero. An explicit algorithm also needs an explicit model of $C$ and an algorithm for enumerating its rational points with the stated complexity. For $C=\mathbb P^1$ this enumeration is immediate, and the assertion reduces to univariate root counting.

\subsection{P-05: a low-width annihilator for a union of curves}
\paragraph{Statement P-05.}
Let $C_j/k$ ($1\le j\le M$) be smooth projective geometrically connected curves and let $r_{1j},r_{2j}\in k(C_j)$ have poles bounded by a common effective divisor $D_j$ of degree $\delta_j\ge1$. There is a nonzero polynomial $P\in k[X,Y]$, of degree at most $2\sum_j\delta_j$ in each variable, with
\[
 P(r_{1j},r_{2j})=0\quad\text{for every }j.
\]
Moreover $P$ has a two-variable ROABP of width at most $2\sum_j\delta_j+1$. Thus a union of such curve images does not hit that entire width/degree class.
\paragraph{Proof.}
For each $j$, put $b_j=2\delta_j$. The $(b_j+1)^2$ monomials $r_{1j}^a r_{2j}^b$ ($0\le a,b\le b_j$) lie in $L(2b_jD_j)$. The elementary divisor-space bound $\dim_k L(E)\le\deg E+1$ for effective $E$ gives
\[
 \dim L(2b_jD_j)\le2b_j\delta_j+1=4\delta_j^2+1<(2\delta_j+1)^2.
\]
For clarity, this divisor-space bound follows after scalar extension to an algebraic closure: imposing vanishing at any $\deg E+1$ distinct points away from $E$ is injective, since a nonzero section cannot have more zeros than poles. Base extension preserves dimension. Hence there is a nontrivial $k$-linear relation, defining a nonzero formal polynomial $P_j(X,Y)$ of bidegree at most $(b_j,b_j)$ that vanishes on $C_j$. Set $P=\prod_jP_j$. The polynomial ring is a domain, so $P\ne0$; its individual degrees are at most $B_0:=\sum_jb_j$. Finally write $P(X,Y)=\sum_{a=0}^{B_0}X^a Q_a(Y)$ and represent it by the row $[1,X,\ldots,X^{B_0}]$ times the column $(Q_0(Y),\ldots,Q_{B_0}(Y))^{\mathsf T}$. This is an ROABP of width $B_0+1$. $\square$

The quantitative conclusion is a lower bound on a specific curve-based construction. It does not exclude a family whose total pole degree is much larger than the target ROABP width or degree.

\subsection{P-06: what general Frobenius orthogonality proves}
\paragraph{Statement P-06.}
Let $U/\mathbb F_q$ be a smooth geometrically connected affine curve. Let $\mathcal F_1,\ldots,\mathcal F_s$ be geometrically irreducible lisse $\overline{\mathbb Q}_\ell$-sheaves, pure of weight zero for a fixed complex embedding, pairwise geometrically nonisomorphic. Set $t_i(x)=\operatorname{tr}(\operatorname{Frob}_x\mid\mathcal F_i)$, and
\[
 B=\max_{i,j}\dim H_c^1(U_{\overline{\mathbb F}_q},\mathcal F_i\otimes\mathcal F_j^\vee).
\]
If $q>s^2B^2$, the functions $t_1,\ldots,t_s$ on $U(\mathbb F_q)$ are linearly independent. More precisely, with $Q=\#U(\mathbb F_q)$,
\[
 \frac1Q\sum_x\left|\sum_i c_it_i(x)\right|^2
 \ge\frac{q-sB\sqrt q}{Q}\sum_i|c_i|^2.
\]
\paragraph{Proof.}
Purity of weight zero identifies the complex conjugate of a Frobenius trace with the trace of the dual representation. For $\mathcal G_{ij}=\mathcal F_i\otimes\mathcal F_j^\vee$, the trace formula gives
\[
 \sum_{x\in U(\mathbb F_q)}t_i(x)\overline{t_j(x)}
 =\sum_{a=0}^2(-1)^a\operatorname{tr}(\operatorname{Frob}_q\mid H_c^a(U_{\bar{\mathbb F}_q},\mathcal G_{ij})).
\]
Here $H_c^0=0$. By duality and geometric irreducibility, $H_c^2$ is zero for $i\ne j$ and is $\overline{\mathbb Q}_\ell(-1)$ for $i=j$: in the latter case its invariant scalar endomorphisms are fixed by arithmetic Frobenius. Thus the $H_c^2$ contribution is $q\delta_{ij}$. Deligne's weight theorem bounds the absolute values of the $H_c^1$ eigenvalues by $\sqrt q$, so the error in every Gram entry is at most $B\sqrt q$. If $E$ is the error matrix, $|c^*Ec|\le B\sqrt q(\sum|c_i|)^2\le sB\sqrt q\sum|c_i|^2$. Dividing by $Q$ gives the result; the displayed strict inequality also implies $Q>0$. $\square$

For an explicit conductor budget, let $\overline U$ have genus $g$ and let $b=\#(\overline U\setminus U)(\overline{\mathbb F}_q)$ count geometric boundary points. Write $r_i=\operatorname{rank}\mathcal F_i$ and $\sigma_{ij}=\sum_{x\in\overline U\setminus U}\operatorname{Swan}_x(\mathcal F_i\otimes\mathcal F_j^\vee)$. Then the Grothendieck--Ogg--Shafarevich formula gives the exact expression
\[
 \dim H_c^1(U_{\bar{\mathbb F}_q},\mathcal G_{ij})
 =\delta_{ij}+r_ir_j(2g-2+b)+\sigma_{ij}.
\]
Indeed $\chi_c(U,\mathcal G_{ij})=r_ir_j(2-2g-b)-\sigma_{ij}$, while the proof above gives $h_c^0=0$ and $h_c^2=\delta_{ij}$. This proves the displayed expression without discarding the conductor term. The Euler--Poincar\'e source is SGA~5, Expos\'e~X, Theorem~7.1 and Corollary~7.12, formula~(7.13), with the $\ell$-adic form also written in Katz--Sarnak, proof of Lemma~10.1.12, p.~296 (P-B14). In particular, if all tensor sheaves are tame and $r_i\le R$, then $B\le1+R^2\max(0,2g-2+b)$ is a sufficient bound. A claimed polynomial-size trace sampler must bound these geometric quantities polynomially as well as $s$.

The trace formula and weights are cited as P-B03; the passage from cohomological estimates to the explicit Gram inequality is the deduction here. The condition uses the dimensions of the relevant cohomology groups. Saying only ``equidistribution'' without controlling these quantities does not supply the claimed complexity. Tensor products used to model many monomials can increase both their number $s$ and $B$.

\subsection{Geometric dictionary for P-01--P-06}
For a monomial set $\mathcal S$, the coefficient space is the affine scheme $\operatorname{Spec}\mathbb Z[c_S:S\in\mathcal S]$. Evaluation on a finite set is a linear morphism of free modules. Injectivity over characteristic zero in P-01/P-02 is witnessed by a nonzero integral maximal minor, so it holds after base change to $\mathbb Z[1/\Delta]$ for some nonzero minor $\Delta$, but need not hold in every positive characteristic. P-03(a) is the idempotent/Lagrange basis of the finite etale sign-cube algebra over $\mathbb Z[1/2]$. P-03(b) is the alternating exterior product of three sections of the rank-three bundle generated by $1,X,X^2$; its vanishing is rank at most two.

In P-04 the map $k[x_1,\ldots,x_n]\to k(C)$ is pullback along a rational morphism to affine space, and the degree budget is a divisor controlling the corresponding section of $\mathcal O_C(D)$. P-05 concerns the nontrivial kernel of that pullback, detected by dimension of global sections. In P-06 the functions are images under the trace-of-Frobenius map on the Grothendieck group of constructible $\ell$-adic complexes. Tensor product becomes pointwise multiplication, duality becomes complex conjugation for weight zero, and compactly supported pushforward to $\operatorname{Spec}\mathbb F_q$ gives the trace sum. Purity bounds the nonconstant part of this pushforward. These assertions do not posit an unconstructed motive for an arbitrary ROABP. Exact trace/weight locators are SGA $4\frac12$, Rapport sur la formule des traces, Theorem~3.2; Deligne, Weil II, Theorem~3.3.1 and \S3.3.10; and Katz--Sarnak, \S\S9.0.13--9.0.16. The general equidistribution statement is Katz--Sarnak, Theorem~9.2.6 (compare Weil II, Theorem~3.5.3). Further SGA/EGA locators are supplied in the global dictionary.

\subsection{P-07: the next construction---bounded-support isolating weights}
\paragraph{Statement P-07.}
Let $S\subseteq\{0,\ldots,d\}^n$ have at most $s$ elements. Fix a prime
\[
 p>\max\{2,d,n\tbinom{s}{2}\}.
\]
For each $t\in\mathbb F_p$, let
\[
 w_i(t)=1+\operatorname{rep}_p(t^i)\in\{1,\ldots,p\},\qquad 1\le i\le n.
\]
Some member of this explicit $p$-element family assigns pairwise different integer weights $w(t)\cdot\alpha$ to the elements of $S$. In particular it isolates their minimum. Over any field containing $ndp+1$ distinct scalars, the substitutions $x_i=z^{w_i(t)}$ followed by evaluation at these scalars form a hitting set of size at most $p(ndp+1)$ for all $s$-sparse polynomials of individual degree at most $d$.
\paragraph{Proof.}
For distinct $\alpha,\beta\in S$, equality of their integer weights implies
\[
 0=\sum_i(\alpha_i-\beta_i)(1+t^i)\quad\text{in }\mathbb F_p.
\]
This polynomial in $t$ is nonzero: at least one nonconstant coefficient $\alpha_i-\beta_i$ is nonzero modulo $p$, because its absolute value is at most $d<p$. Its degree is at most $n$, so it has at most $n$ roots. The union of all pair-collision sets has size at most $n\binom{s}{2}<p$. Choosing $t$ outside that union proves the weight assertion. Under that substitution, every nonzero coefficient remains on a distinct power of $z$, so the resulting univariate polynomial is nonzero, of degree at most $ndp$. It cannot vanish at all $ndp+1$ distinct scalar values. The union of all evaluation points therefore hits every polynomial in the stated class. $\square$

This gives a polynomial-size family for bounded sparsity and allows arbitrary coefficients and cancellation; it does not require purity, because elementary root counting is already stronger for the polynomial collision predicates involved. Its geometric form is a family of one-parameter subgroups of $\mathbb G_m^n$, and the successful member lies in the interior of a normal cone exposing a unique vertex of the support's Newton polytope. For general width-$w$ ROABPs the support size can be exponential even at $w=1$, so substituting $s=\operatorname{poly}(w,n,d)$ here is unjustified. The remaining transfer must exploit coefficient-space structure, as in basis isolation, rather than enumerate support pairs.

\subsection{P-08: a precise conditional arithmetic encoding reduction}
\paragraph{Statement P-08 (conditional reduction).}
Let $\mathcal C_N$ be a polynomial class. Suppose the following data can be constructed in time polynomial in $N$: a common finite Galois extension $E_N/\mathbb Q$, its nonzero discriminant, and a map $p\mapsto h_N(p)$ from unramified rational primes to evaluation points. Suppose $[E_N:\mathbb Q]$ and $\log|\Delta_{E_N}|$ are polynomially bounded, that $h_N(p)$ is computable in polynomial time in $N+\log p$, and that for every nonzero $f\in\mathcal C_N$ some conjugacy class $C_f\subseteq\operatorname{Gal}(E_N/\mathbb Q)$ has
\[
 \operatorname{Frob}_p\in C_f\quad\Longrightarrow\quad f(h_N(p))\ne0
\]
for every unramified $p$. Under GRH for $\zeta_{E_N}$ these data yield a polynomial-size, polynomial-time hitting set for $\mathcal C_N$.
\paragraph{Proof.}
Lagarias--Odlyzko, Corollary~1.2 (P-B04), gives an unramified prime in every conjugacy class with
\[
 p\le c\,\log^2|\Delta_{E_N}|\,\log^4\log|\Delta_{E_N}|
\]
for an effective absolute constant, with the trivial extension treated separately. Increasing the bound to cover the finitely many small discriminants gives a polynomial bound $B(N)$. Enumerate all primes up to this bound by trial division, discard primes dividing the supplied discriminant, and evaluate $h_N$ on the rest. The set has polynomial size and is polynomial-time constructible; it contains a point from $C_f$ for every nonzero $f$, proving the claim. $\square$

If $E_N/\mathbb Q$ is abelian, factorization into Dirichlet $L$-functions makes this GRH assumption a consequence of the stipulated Alp\"oge theorem. For arbitrary $E_N$, it follows conditionally from Masullo's asserted Hecke GRH, by using the trivial Hecke character of $E_N$. The extension and evaluation map must be common to the polynomial class: an extension depending on the hidden polynomial would not construct a black-box hitting set. No construction satisfying these hypotheses for general ROABPs is supplied by the stated RH/GRH theorems or by the finite-field purity argument. The exact remaining gate is the explicit common encoding with polynomial discriminant, together with the Frobenius-to-nonvanishing implication. Deterministic quadratic nonresidues and modular square roots are the consequences assigned to Prompt~28 and are not rederived here.

\subsection{P-09: approximate positive cubature already costs the full tensor dimension}
\paragraph{Statement P-09 (adapted trace argument).}
Let $n\ge1$, $\delta\ge0$, $D=(\delta+1)^n$, and let $\xi_1,\ldots,\xi_M\in\mathbb C^n$ have positive real weights $w_j$. Suppose that constants $0<a\le b$ satisfy
\[
 a\|\operatorname{coeff} f\|_2^2
 \le \sum_{j=1}^M w_j|f(\xi_j)|^2
 \le b\|\operatorname{coeff} f\|_2^2
 \tag{P09.1}
\]
for every width-one product $f(x)=\prod_{i=1}^n p_i(x_i)$ with $\deg p_i\le\delta$. Then
\[
 M\ge \frac{a}{b}D.
 \tag{P09.2}
\]
In particular a relative error $0\le\varepsilon<1$ forces
$M\ge D(1-\varepsilon)/(1+\varepsilon)$. Nonnegative weights are allowed after zero weights are removed.
\paragraph{Proof.}
Put $v(z)=(1,z,\ldots,z^\delta)$ and
$e_j=\bigotimes_{i=1}^n v(\xi_{j,i})\in\mathbb C^D$.
Coefficient vectors are paired bilinearly with $e_j$ for evaluation.
Since $e_j\ne0$ and is decomposable, the vector
$c^{(j)}=\overline{e_j}/\|e_j\|_2$ is the coefficient vector of a width-one product, of coefficient norm one. Its value at $\xi_j$ is $\|e_j\|_2$. All summands in (P09.1) are nonnegative, so the upper bound applied to this product gives
\[
 w_j\|e_j\|_2^2\le b.
\]
Apply the lower bound to each of the $D$ monomials $x^\alpha$, $\alpha\in\{0,\ldots,\delta\}^n$, each itself a normalized width-one product, and sum. This gives
\[
 aD\le\sum_\alpha\sum_j w_j|\xi_j^\alpha|^2
 =\sum_j w_j\|e_j\|_2^2\le Mb,
\]
which proves the result. Equivalently, averaging over independently chosen Haar unit vectors in the $n$ coefficient factors gives covariance $I_D/D$ and the same trace bound. The finite basis sum avoids any integration assumption. $\square$

The proof also applies to arbitrary decomposable measurement vectors, not only polynomial evaluations. It is related to the rank-one measurement obstruction for ordinary $\ell_2$ restricted isometry: Cai and Zhang (2015, Lemma~2.1, p.~109; P-B15) prove an ambient-dimension obstruction for their Gaussian matrix model and use a different $\ell_1$ measurement condition instead. The weighted tensor statement above is proved here directly; no priority claim is made. Any proposed purity sampler that uniformly approximates coefficient energy for all bounded-width ROABPs includes width one and therefore meets this barrier. The barrier concerns the quantitative two-sided energy requirement; it does not exclude hitting sets.

\subsection{P-10: positivity alone, tensor-train kernels, and exact width-one sampling}
\paragraph{Statement P-10.}
Let $n,w\ge1$, $\delta\ge0$, and let $\mathcal T_w\subseteq\mathbb C^{(\delta+1)^n}$ be the cone of coefficient tensors whose consecutive prefix/suffix flattenings all have rank at most $w$, for a fixed variable order. Let $A_Hc=(f_c(\xi_j))_{j=1}^M$ and let all sampling weights be positive. The following are equivalent:
\[
\begin{gathered}
 H\text{ hits every nonzero tensor in }\mathcal T_w;\\
 \ker A_H\cap\mathcal T_w=\{0\};\qquad
 \inf_{c\in\mathcal T_w,\,\|c\|_2=1}
       \sum_jw_j|f_c(\xi_j)|^2>0.
\end{gathered}
 \tag{P10.1}
\]
The cone is exactly the coefficient class of width-at-most-$w$ ROABPs in that order. If $a$ and $b$ are the minimum and maximum sampled energies on its unit sphere, every such hitting set satisfies
\[
 \frac ba\ge\frac{(\delta+1)^n}{M}.
 \tag{P10.2}
\]
If $q_w=\dim\mathcal T_w$, every complex hitting set has at least $q_w$ points, and a nonempty Zariski-open set of $q_w$-tuples of evaluation points consists of hitting sets. Moreover
\[
 q_w\le n(\delta+1)w^2.
 \tag{P10.3}
\]
These last assertions establish existence, with integer choices possible, but do not provide an efficient deterministic construction or a coordinate-height bound.
The same equivalences, energy bound, and exact generic cardinality $q_w$ hold for unknown variable order: replace $\mathcal T_w$ by the union of its images under all variable permutations.

For width one there is an exact explicit answer over every field containing at least $n\delta+1$ elements: $n\delta+1$ distinct diagonal points $(t,\ldots,t)$ hit all nonzero products, and no set with at most $n\delta$ arbitrary evaluation points does so.
\paragraph{Proof.}
Flattening rank bounds are the vanishing of homogeneous minors, so $\mathcal T_w$ is closed. Rank factorization gives the standard tensor-train description. One way to see the converse explicitly is to let $U_i$ be the column space of the $i$th flattening. Then
$U_i\subseteq U_{i-1}\otimes\mathbb C^{\delta+1}$: contractions in the remaining factors produce the columns, and the preceding left support contains each such contraction. Choosing bases in these spaces writes their inclusions as cores of sizes at most $w\times(\delta+1)\times w$, with one-dimensional boundary cores. Their contraction is the original coefficient tensor. Conversely every core contraction factors each prefix flattening through its bond space, giving the rank bound. This is the exact rank-factorization result of Oseledets (2011, Theorem~2.1, pp.~2298--2299; P-B17), restated for polynomial coefficients. Zero-padding the cores gives a surjective polynomial parameterization with at most $n(\delta+1)w^2$ parameters, proving (P10.3).

The first two conditions in (P10.1) are identical by definition. With positive weights, a sampled energy vanishes exactly when all evaluations vanish. The intersection of the closed cone with the unit sphere is compact. Thus absence of a zero of the continuous energy function there is equivalent to a strictly positive minimum. Its maximum is finite as well. Since $\mathcal T_w$ contains all product tensors, P-09 gives (P10.2).

For the generic assertion let $X=\mathbb P(\mathcal T_w)$, of dimension $q_w-1$. Imposing $M<q_w$ homogeneous linear equations on its affine cone leaves dimension at least $q_w-M\ge1$, by the principal ideal dimension bound; over $\mathbb C$ the common zero set therefore contains a nonzero tensor. This proves the lower bound. Evaluation hyperplanes have no common base point on $X$, because a nonzero polynomial is nonzero at some point of $\mathbb C^n$. At each step choose an evaluation hyperplane containing no irreducible component of the remaining projective intersection. Such a choice exists: for each of finitely many components, choose one nonzero coefficient point on it; its polynomial excludes a proper hypersurface of possible evaluation points. Intersecting with the chosen hyperplane lowers every positive component's dimension by one and removes all zero-dimensional components. After $q_w$ steps the intersection is empty. Finally the incidence subset
\[
 \{([c],\xi_1,\ldots,\xi_{q_w}):
      [c]\in X,\ f_c(\xi_j)=0\text{ for all }j\}
 \subseteq X\times(\mathbb A^n)^{q_w}
\]
is closed and its projection is closed because $X$ is projective. The choices just made show that this projected bad set is proper. Its complement is the asserted nonempty open set. Integer tuples are Zariski dense in affine complex space, so some integer choice lies in that open set. The argument does not bound how far one must search.

For unknown order, the union $\bigcup_{\pi\in S_n}\pi\mathcal T_w$ is a finite union of closed cones of the same dimension $q_w$. It is therefore closed of dimension $q_w$, contains all product tensors, and has no common evaluation base point after projectivization. The compactness, dimension, and incidence arguments just given apply to this union unchanged. Although there are $n!$ possible orders, the existential number of generic evaluations remains $q_w$; no effective way to select them has been supplied.

For the explicit width-one upper bound, the restriction
$f(t,\ldots,t)=\prod_i p_i(t)$ is nonzero and has degree at most $n\delta$, hence cannot vanish at all $n\delta+1$ distinct scalars. For the lower bound, partition any collection of $M\le n\delta$ sample indices into $n$ sets $J_i$, each of size at most $\delta$, and set
\[
 f(x)=\prod_{i=1}^n\prod_{j\in J_i}(x_i-\xi_{j,i}).
\]
This is a nonzero product of monic univariate factors, has individual degree at most $\delta$, and vanishes at every listed point. Empty factors equal one. $\square$

The width-one construction and dimension threshold agree with the maximal completely entangled subspaces of Parthasarathy (2004, Propositions~1.3--1.4 and Theorem~1.5; P-B16). The kernel of the diagonal evaluation map is such a subspace of dimension $D-n\delta-1$. In this terminology ``completely entangled'' means precisely that the kernel contains no nonzero decomposable tensor; it is not a physical assumption on the polynomials. For larger tensor-train ranks, (P10.1) is the precise weaker gate. Proving positivity directly could avoid the exponential constant-distortion requirement, but generic existence and compactness do not construct that gate efficiently.

\subsection{P-11: an explicit number-field point is already an unconditional hitting-set encoding}
\paragraph{Statement P-11.}
Let $\mathcal C\subseteq\mathbb Q[x_1,\ldots,x_n]$ have individual degree at most $\delta$. Suppose a common algebraic number $\theta$, polynomials $A_i(T)\in\mathbb Z[T]$ of degree at most $a$, and an integer $b\ne0$ are supplied such that
\[
 f\ne0,\ f\in\mathcal C\quad\Longrightarrow\quad
 f\bigl(A_1(\theta)/b,\ldots,A_n(\theta)/b\bigr)\ne0.
 \tag{P11.1}
\]
Put $d=n\delta$. The completely explicit integer set
\[
 H=\{(uA_1(t),\ldots,uA_n(t)):
        0\le u\le d,\ 0\le t\le da\}
 \tag{P11.2}
\]
hits $\mathcal C$ and has size at most $(d+1)(da+1)$. If $b=1$, the single grid $t=0,\ldots,da$ with $u=1$ suffices. These conclusions need no GRH, no discriminant bound, no coefficient-height bound on $f$, and no algorithm for arithmetic in $\mathbb Q(\theta)$.

Conversely suppose an explicitly given integer hitting set $H_1,\ldots,H_M$, $M\ge1$, hits a class $\mathcal C\subseteq\mathbb Q[x]$ of total degree at most $d$. Polynomial-time interpolation and a prime $\ell>d(M-1)+1$ construct data as in (P11.1) with $\theta$ a primitive $\ell$th root of unity. The degree, coordinate degrees, coefficient bit lengths, and logarithmic discriminant are polynomial in $d,M$ and the bit lengths of the given points, choosing $\ell=O(d(M-1)+2)$. Thus the two explicit encodings are polynomially interconvertible over $\mathbb Q$.
\paragraph{Proof.}
For nonzero $f\in\mathcal C$, define
\[
 G_f(U,T)=f(UA_1(T),\ldots,UA_n(T)).
\]
It is nonzero, since its evaluation at $(U,T)=(1/b,\theta)$ is nonzero by hypothesis. Its separate degrees are at most $d$ in $U$ and $da$ in $T$. If it vanished on the entire rectangular grid in (P11.2), then, for each fixed $t$, univariate interpolation at the $d+1$ values of $u$ would force all its $U$-coefficients to vanish at $t$. Interpolation at the $da+1$ values of $t$ would then make every coefficient polynomial zero, a contradiction. When $b=1$, $G_f(1,T)$ is itself nonzero at $\theta$ and has degree at most $da$, proving the smaller bound. If $|\operatorname{coeff} A_i|\le2^K$, the coordinates in (P11.2) have bit length at most
\[
 K+\left\lceil\log_2(d+1)\right\rceil
  +\left\lceil\log_2(a+1)\right\rceil
  +a\left\lceil\log_2\max(1,da)\right\rceil+2,
\]
so polynomial input sizes give polynomial bit complexity as well as polynomial cardinality.

For the reverse direction interpolate rational polynomials $a_i(T)$ of degree at most $M-1$ with $a_i(j)=H_{j+1,i}$ for $0\le j<M$. A common denominator $b=(M-1)!$ clears their coefficients, because each Lagrange denominator is $(-1)^{M-1-j}j!(M-1-j)!$, which divides $b$. Hence $A_i=ba_i$ are integral and have polynomial coefficient bit lengths; explicitly, if the coordinates of $H$ have absolute values at most $2^K$, the Lagrange product expansion gives coefficient bit length $O(K+M\log(M+1))$. For nonzero $f\in\mathcal C$, the polynomial $f(a_1(T),\ldots,a_n(T))$ is nonzero because it is nonzero at at least one interpolation node. Its degree is at most $d(M-1)$. Choose a prime $\ell>d(M-1)+1$, and let $\theta=\zeta_\ell$. Its minimal polynomial over $\mathbb Q$ is $\Phi_\ell(T)=1+T+\cdots+T^{\ell-1}$: irreducibility follows from Eisenstein applied after $T\mapsto T+1$. Its degree exceeds $d(M-1)$, so it cannot divide the nonzero rational polynomial above. This proves (P11.1). Bertrand's bound and trial division choose $\ell$ of the stated size. The polynomial discriminant has absolute value $\ell^{\ell-2}$: differentiate $(T^\ell-1)/(T-1)$ at its roots and use $\prod_{a=1}^{\ell-1}(1-\zeta_\ell^a)=\Phi_\ell(1)=\ell$ in the product of the derivatives. The field discriminant divides this order discriminant. On every conjugate of $\theta$, $|a_i(\theta)|$ is bounded by the sum of the absolute values of its coefficients. Thus all claimed representation and height bounds are polynomial. $\square$

In particular, if a primitive-element presentation of a degree-$e$ number field is given with coordinate polynomials of degree below $e$, the remaining arithmetic gate is exactly the nonvanishing condition (P11.1), and P-11 uses $a\le e-1$. The argument does not assert that primitive-element coordinates of small degree and polynomial bit size can be obtained from an arbitrary implicit field description. It also does not extend the reverse direction to unrestricted complex coefficients: minimal-polynomial divisibility is a statement over $\mathbb Q$.

There are concrete restricted encodings. For width-one products over $\mathbb Q$, take $\theta=\zeta_\ell$ with $\ell-1>\delta$ and every coordinate equal to $\theta$; no nonzero factor of degree at most $\delta$ vanishes there. This fails already for the width-two polynomial $x_1-x_2$ when $n\ge2$. P-07 provides an explicit integer hitting set of polynomial size and polynomial coordinate bit length for bounded sparsity, so the reverse construction gives a polynomial-degree cyclotomic encoding for that larger restricted class as well. These constructions do not give a new route to general ROABP PIT: they show exactly why supplying a common effective nonvanishing encoding has already supplied essentially the missing hitting set. The strictly prime-dependent encoding in P-08 need not have the explicit polynomial coordinate representation assumed here and is not covered by this simplification.

\subsection{P-12: an explicit width-three algebraic point and trace-coordinate span map}
\paragraph{Statement P-12.}
Let $n,D\ge1$.  Let $\mathcal C_{3,n,D}$ consist of the rational-coefficient polynomials of total degree at most $D$ having a width-at-most-three ROABP in some variable order.  The order is not given.  Put
\[
 B=nD+1,\qquad N=B^{21},\qquad b=(n-1)!,\qquad
 a=(n-1)B^{20}+B^{19}.
\]
For the nodes $0,\ldots,n-1$, let $\lambda_i$ be the Lagrange polynomial for node $i-1$, and define the explicit integer polynomials
\[
 A_i(T)=bT+
    \sum_{j=1}^{10}T^{B^{2j-1}}\,
       b\lambda_i\bigl(T^{B^{2j}}\bigr),
       \qquad 1\le i\le n.
 \tag{P12.1}
\]
Let $\theta$ be any root of $T^N-2$ and set $\alpha_i=A_i(\theta)$.  Then
\[
 0\ne f\in\mathcal C_{3,n,D}\quad\Longrightarrow\quad
 f(\alpha_1,\ldots,\alpha_n)\ne0.
 \tag{P12.2}
\]
The common number field $E=\mathbb Q(\theta)$ has degree $N$, and
\[
 \log|\Delta_E|\le N\log N+(N-1)\log2.
 \tag{P12.3}
\]
Every $A_i$ has degree at most $a$, at most $10n+1$ nonzero terms, and integer coefficients of bit length $O(n\log(n+1))$.  These data are constructible in time polynomial in $n,D$; the polynomial exponents are deliberately not optimized.

There is also an explicit one-parameter zero bound, valid over every characteristic-zero coefficient field:
\[
 f\ne0\quad\Longrightarrow\quad
 0\ne f(A_1(T),\ldots,A_n(T)),\qquad
 \deg f(A(T))\le Da<N.
 \tag{P12.4}
\]
Consequently the integer points $A(t)$ for $0\le t\le Da$ form a hitting set of size at most $Da+1$, with polynomial coordinate bit lengths.  This is a concrete construction for width three, including arbitrarily many rank-two cuts.

For a constant-boundary row product
\[
 R(x)=uA^{\rm mat}_1(x_{\sigma(1)})\cdots
          A^{\rm mat}_n(x_{\sigma(n)})\in\mathbb Q[x]^{1\times3}
\]
whose entries have total degree at most $D$, let $V\subseteq\mathbb Q^3$ be its coefficient-row span.  Write
\[
 R(\alpha)=\sum_{k=0}^{N-1}\theta^k r_k.
\]
Then the explicit arithmetic evaluation preserves that span in the following exact sense:
\[
 V=\operatorname{span}_{\mathbb Q}\{r_k:0\le k<N\}
  =\operatorname{span}_{\mathbb Q}
       \{\operatorname{Tr}_{E/\mathbb Q}(\beta_k R(\alpha)):0\le k<N\},
 \tag{P12.5}
\]
where
\[
 \beta_0=1/N,\qquad \beta_k=\theta^{N-k}/(2N)\quad(1\le k<N).
\]
The trace is applied entrywise.  The assertion concerns rational coefficient spans after expansion in the field basis, not the span of a single row over $E$.

\paragraph{Proof.}
The complete width-three rank-preservation proof in I-21 supplies the generator
\[
 G_i(t,\mathbf z,\mathbf y)=t+\sum_{j=1}^{10}z_j\lambda_i(y_j)
\]
for every variable order.  This is the specialized construction independently derived and checked from Liu's unrefereed rank-profile mechanism (2026, \S3.1, Theorem~3.2/Lemma~3.3 and \S3.2, Lemmas~3.4/3.6; P-B18).  I-21, rather than an assumption about the entire preprint, is the mathematical input here.

For a nonzero $f$, first put $f_b(x)=f(bx_1,\ldots,bx_n)$.  The nonzero scalar $b$ preserves nonvanishing, degree, and the width-three representation.  Hence $f_b(G)$ is a nonzero polynomial in its twenty-one parameters.  Its degree in $t$ and in each amplitude $z_j$ is at most $D$, and its degree in each selector $y_j$ is at most $(n-1)D$.  All these individual degrees are strictly below $B$.

Order the parameters as $t,z_1,y_1,\ldots,z_{10},y_{10}$ and substitute $T^{B^k}$ for the parameter in position $k$, starting with position zero.  Distinct parameter monomials have distinct resulting exponents, by uniqueness of base-$B$ expansion.  The substitution is therefore injective on the bounded-degree coefficient space.  Its value on $f_b(G)$ is precisely $f(A_1(T),\ldots,A_n(T))$, proving nonvanishing in (P12.4).  The same argument works over any characteristic-zero coefficient field.

Formula (P12.1) gives degree at most $a$ for every coordinate, so the composition has degree at most $Da$.  Explicitly,
\[
 Da=B^{20}\bigl((n-1)D+D/B\bigr)
      <B^{20}(nD)<B^{21}=N.
\]
The polynomial $T^N-2$ is Eisenstein at two and is consequently the minimal polynomial of $\theta$ over $\mathbb Q$.  A nonzero rational polynomial of degree below $N$ cannot vanish at $\theta$, proving (P12.2).

The Lagrange denominator at node $i-1$ is
$(-1)^{n-i}(i-1)!(n-i)!$, which divides $(n-1)!$.  Thus $b\lambda_i$ is integral.  Expansion of its linear factors bounds every coefficient by $(n!)^2$, a sufficient uniform bound including $n=1$; there are at most $n$ coefficients.  Distinct summands in (P12.1) have disjoint exponent patterns, except for zero coefficients, and give at most $10n+1$ terms.  This proves the stated degree, sparsity, and coefficient-size estimates.  Evaluation at the integers through $Da$ has coordinate bit length
\[
 O\bigl(n\log(n+1)+\log(10n+1)+a\log(Da+2)\bigr),
\]
which is polynomial in $n,D$.  The integer hitting claim follows from the univariate zero bound.

The order $\mathbb Z[\theta]$ has discriminant of absolute value
$N^N2^{N-1}$, obtained by multiplying the derivative values
$N\theta_j^{N-1}$ over the $N$ roots of $T^N-2$.  The field discriminant divides the order discriminant by the square of the order index.  This proves (P12.3).  No discriminant estimate is needed for the nonvanishing proof itself.

For the row-span assertion, the coefficients $r_k$ belong to $V$, because substitution and reduction in a rational basis take rational linear combinations of the original coefficient rows.  In fact the degree bound below $N$ means that no reduction is needed.  If their span were a proper subspace of $V$, there would be a rational column $v$ annihilating every $r_k$ but not every coefficient row of $R$.  The scalar polynomial $Rv$ would then be a nonzero member of $\mathcal C_{3,n,D}$ with $(Rv)(\alpha)=0$, contradicting (P12.2).  This proves the first equality in (P12.5).

Finally $\operatorname{Tr}_{E/\mathbb Q}(\theta^j)=0$ when $j$ is not divisible by $N$, while the traces of $1$ and $\theta^N$ are $N$ and $2N$.  For $0\le j,k<N$ this gives
$\operatorname{Tr}(\beta_k\theta^j)=\delta_{kj}$ directly.  Thus the displayed $\beta_k$ are the trace-dual basis, and the traces in (P12.5) equal the coefficient rows $r_k$.  $\square$

The arithmetic trace here is the ordinary number-field trace.  This construction supplies the previously missing effective algebraic point for the width-three class through rank preservation.  It does not infer that point from Frobenius equidistribution or provide a bound polynomial in a variable width.  Evaluating a black box at $\alpha$ requires access over $E$; the integer curve points in (P12.4) give an alternative using rational or integer queries.  Quantitative positive coefficient-energy equivalence remains subject to P-09.

\subsection{P-13: a general arithmetic encoding of a centered interpolation generator}
\paragraph{Statement P-13.}
Let $n,D,K\geq1$. Let $\mathcal C$ be a class of rational
polynomials in $n$ variables of total degree at most $D$, closed
under a common nonzero rational scaling of all variables.
Use the Lagrange polynomials $\lambda_i$ for the integer nodes
$0,\ldots,n-1$. Assume the following centered generator hits
every nonzero member of $\mathcal C$:
\[
 G_i^{(K)}(t,\mathbf z,\mathbf y)
       =t+\sum_{j=1}^Kz_j\lambda_i(y_j).
 \tag{P13.1}
\]
Put
\[
 B=nD+1,\quad N_K=B^{2K+1},\quad b=(n-1)!,\quad
 a_K=(n-1)B^{2K}+B^{2K-1},
\]
and define the explicit integer polynomials
\[
 A_i^{(K)}(T)=bT+\sum_{j=1}^{K}
       T^{B^{2j-1}}\,b\lambda_i\bigl(T^{B^{2j}}\bigr).
 \tag{P13.2}
\]
If $\theta_K$ is a root of $T^{N_K}-2$ and
$\alpha_i^{(K)}=A_i^{(K)}(\theta_K)$, then
\[
 0\ne f\in\mathcal C
       \quad\Longrightarrow\quad f(\alpha^{(K)})\ne0.
 \tag{P13.3}
\]
The common field has degree $N_K$, and
\[
 \log|\Delta_{\mathbb Q(\theta_K)}|
       \leq N_K\log N_K+(N_K-1)\log2.
 \tag{P13.4}
\]
Every coordinate in (P13.2) has degree at most $a_K$, at most
$Kn+1$ terms, and integer coefficients of bit length
$O(n\log(n+1))$. There is also a nonzero one-parameter
restriction with an explicit degree bound:
\[
 0\ne f\in\mathcal C
   \quad\Longrightarrow\quad
 0\ne f(A^{(K)}(T)),\qquad
 \deg f(A^{(K)}(T))\leq Da_K<N_K.
 \tag{P13.5}
\]
Thus $A^{(K)}(t)$ for $0\leq t\leq Da_K$ is an explicit
integer hitting set of size $Da_K+1$. Its coordinate bit lengths
are at most
\[
 O\bigl(n\log(n+1)+\log(Kn+1)+a_K\log(Da_K+2)\bigr).
 \tag{P13.6}
\]
For fixed $K$ all these numerical budgets are polynomial in
$n,D$. The sparse coordinate formulas themselves have only
$Kn+1$ terms per coordinate and exponent bit length
$O(K\log(nD+1))$; this does not make arithmetic in the degree
$N_K$ field polynomial in a variable $K$.

More generally, let $R(x)\in\mathbb Q[x]^{1\times w}$ have
entries of degree at most $D$, and suppose every nonzero scalar
contraction $R(x)v$, for $v\in\mathbb Q^w$, belongs to
$\mathcal C$. If $V$ is its rational coefficient-row span and
\[
 R(\alpha^{(K)})=\sum_{j=0}^{N_K-1}\theta_K^j r_j,
\]
then
\[
 V=\operatorname{span}_{\mathbb Q}\{r_j:0\leq j<N_K\}
  =\operatorname{span}_{\mathbb Q}
    \{\operatorname{Tr}(\beta_jR(\alpha^{(K)})):0\leq j<N_K\},
 \tag{P13.7}
\]
where the trace is entrywise and
\[
 \beta_0=1/N_K,\qquad
 \beta_j=\theta_K^{N_K-j}/(2N_K)\quad(1\leq j<N_K).
\]

\paragraph{Proof.}
For $f\ne0$, the polynomial $f_b(x)=f(bx)$ is a nonzero member
of $\mathcal C$. By the hypothesis, $f_b(G^{(K)})$ is nonzero.
Its individual degrees are at most $D$ in $t$ and each $z_j$,
and at most $(n-1)D$ in each $y_j$, all strictly below $B$.
Send the $2K+1$ parameters, in the order
$t,z_1,y_1,\ldots,z_K,y_K$, to
$T,T^B,\ldots,T^{B^{2K}}$. Unique base-$B$ expansion makes
the exponents of distinct parameter monomials distinct.
The resulting nonzero polynomial is $f(A^{(K)}(T))$.
Moreover
\[
 Da_K=B^{2K}\bigl((n-1)D+D/B\bigr)
       <B^{2K}(nD)<B^{2K+1}=N_K.
\]
This includes $n=1$. Eisenstein's criterion at two shows that
$T^{N_K}-2$ is the minimal polynomial of $\theta_K$ over
$\mathbb Q$, proving (P13.3). The one-parameter nonvanishing
argument also holds over any characteristic-zero coefficient
field for which the same class-closure and generator hypotheses
hold. The algebraic-point assertion (P13.3) is specifically over
$\mathbb Q$.

The denominator of $\lambda_i$ is
$(-1)^{n-i}(i-1)!(n-i)!$, which divides $b$. Hence every
$b\lambda_i$ is integral. Its coefficients have absolute values
at most $(n!)^2$, and it has at most $n$ terms. The exponent
patterns in (P13.2) are distinct by the same base-$B$ argument,
giving the sparsity and coefficient-size bounds. The degree
bound and ordinary integer evaluation prove (P13.6). The
polynomial discriminant has absolute value
$N_K^{N_K}2^{N_K-1}$, and the field discriminant divides it by
the square of the order index; this proves (P13.4).

Every expanded row coefficient $r_j$ is a rational combination
of the original coefficient rows. In fact (P13.5) means no
reduction modulo the minimal polynomial is needed. If their
span were smaller than $V$, a rational annihilator column
would produce a nonzero scalar contraction in $\mathcal C$
vanishing at $\alpha^{(K)}$, contrary to (P13.3).
Finally $\operatorname{Tr}(\theta_K^s)$ vanishes when $s$ is not
divisible by $N_K$, while the traces of $1$ and $\theta_K^{N_K}$
are $N_K$ and $2N_K$. Consequently
$\operatorname{Tr}(\beta_j\theta_K^h)=\delta_{jh}$ for
$0\leq j,h<N_K$. This proves (P13.7). $\square$

\paragraph{Proved specializations.}
I-30 supplies the hypothesis for all rational width-at-most-three
ROABPs in one hidden variable order with $K=6$. Therefore
\[
 N_6=(nD+1)^{13},\qquad
 a_6=(n-1)(nD+1)^{12}+(nD+1)^{11}
\]
give a common algebraic point, an integer hitting curve, and exact
coefficient-row-span recovery for that class. This improves the
earlier valid twenty-one-parameter construction in P-12.

I-36 reduces the budget to $K=5$ when every variable-dependent
factor in a width-three representation is singular; constant
invertible factors may be absorbed. Hence this rational subclass
has $N_5=(nD+1)^{11}$. On the further affine-factor subclass,
I-37 proves $K=2$, so
\[
 N_2=(nD+1)^5,\qquad
 a_2=(n-1)(nD+1)^4+(nD+1)^3.
\]
These hypotheses are preserved by common nonzero variable scaling
and by changing the constant scalar output column. Consequently
both specializations give the common algebraic point, the explicit
integer hitting curve, and exact coefficient-row-span recovery in
(P13.3)--(P13.7). I-38 further extends the same $K=2$ conclusion to affine
representations in which no two consecutive variable-dependent
factors have rank three after constant factors are absorbed.
That adjacency promise is also preserved by common scaling and
constant output contractions. Two or more consecutive invertible
variable factors, and nonlinear singular factors, remain outside
these two-pair specializations.

For the promised width-$w$ class in I-29, whose internal interval
ranks belong to $\{r,w\}$, put $c=w-r$. Its proved pair budget
$K=2rc+w-1$ gives
\[
 N_K=(nD+1)^{4r(w-r)+2w-1}.
\]
Common variable scaling preserves these interval ranks, and
changing the constant scalar output column does not alter the
internal word. Thus both the scalar and coefficient-span
hypotheses above hold. This supplies an actual arithmetic
encoding for unboundedly many rank resets under that precise
promise. Its field degree and evaluated hitting-set size are
$(nD+1)^{O(w^2)}$, not polynomial in a growing width.

These traces are ordinary number-field traces. No RH,
equidistribution of Frobenius, or equivalence with positive coefficient
energy enters the proof. Neither the parameter count nor
the resulting field degree is asserted to be optimal.

\subsection{P-14: dense interpolation preserves coefficient rows across rank-one resets}
\paragraph{Statement P-14.}
Let $F$ be a field, let $n,w\geq1$, and consider an ROABP in any
one variable order
\[
 f(x)=uA_1(x_{\pi(1)})\cdots A_n(x_{\pi(n)})v,
 \qquad A_i\in F[X]^{w\times w},\quad
 u\in F^{1\times w},\quad v\in F^{w\times1}.
 \tag{P14.1}
\]
Assume that every factor has generic rank either $1$ or $w$ over
$F(X)$. No bound is placed on the number of rank-one factors.
Choose distinct $a_1,\ldots,a_n\in F$, their Lagrange polynomials
$\lambda_i$, and arbitrary $c_i\in F^*$. Put $K=w-1$ and
\[
 G_i(t,\mathbf y,\mathbf z)
       =c_it+\sum_{j=1}^{K}z_j\lambda_i(y_j).
 \tag{P14.2}
\]
Then $f\ne0$ implies $f(G)\ne0$. The generator does not depend
on the order $\pi$ or on which factors have rank one.

More precisely, for any fixed $u$ let
$R(x)=uA_1(x_{\pi(1)})\cdots A_n(x_{\pi(n)})$ and let
$V\subseteq F^{1\times w}$ be its coefficient-row span. Then
\[
 \operatorname{span}_F\{\text{coefficient rows of }R(G)\}=V.
 \tag{P14.3}
\]
This is equality of coefficient-row spaces. It does not assert
equality of the full matrix image cones.

If $\deg f\leq D$, where $D\geq0$, the generator gives a hitting
set of size at most
\[
 (D+1)\bigl((n-1)D+1\bigr)^{w-1}
             \binom{D+w-1}{w-1}.
 \tag{P14.4}
\]
It suffices to work over a field with at least
$\max\{n,D+1,(n-1)D+1\}$ elements; an extension may be used.
All characteristics are allowed. For $w=1$, the empty sum in
(P14.2) is the diagonal line $x_i=c_it$.

In particular every nonzero width-two ROABP in one hidden order,
of arbitrary individual degree, is hit by a common center plus
one interpolation pair. For $n\geq2$, the common center alone
does not hit all width-two ROABPs.

There is also a rational-function form of the result. For distinct
abscissas $a_i$, the evaluation map
\[
 x_i=\frac{g(a_i)}{h(a_i)},\qquad
 \deg g,\deg h\leq w-1,
 \tag{P14.5}
\]
hits the same scalar class wherever its denominators are defined.
The corresponding rational coefficient vectors preserve the
coefficient-row span after expansion in any common polynomial
denominator. The rational form is not needed to construct the
polynomial hitting set (P14.4).

\paragraph{1. A bridge with invertible interior.}
We first prove an elementary local lemma. Let
\[
 b(X)=L(X_0)B_1(X_1)\cdots B_q(X_q)Q(X_{q+1})
 \tag{P14.6}
\]
be a nonzero scalar polynomial, with separate variables in the
displayed factors. Either endpoint may instead be constant.
Assume that $L,Q$ are nonzero polynomial row and column vectors,
respectively, and every $B_i$ is a generically invertible
$w\times w$ polynomial matrix. At the center $X_j=c_jt$, over
$F(t)$, every interior constant is invertible and both endpoint
constants are nonzero. Here each $c_j$ is nonzero and $t$ is
transcendental; substitution into a nonzero univariate polynomial
therefore remains nonzero.

Write $X_j=c_jt+Y_j$. If $L_0=L(c_0t)$, choose an index $a$ with
$(L_0)_a\ne0$. Then
\[
 L(c_0t+Y)=L_0\left(I+
       \frac{e_a}{(L_0)_a}\bigl(L(c_0t+Y)-L_0\bigr)\right).
 \tag{P14.7}
\]
The matrix in parentheses has constant term $I$. Likewise, if
$Q_0=Q(c_{q+1}t)$ and $(Q_0)_b\ne0$, then
\[
 Q(c_{q+1}t+Y)=\left(I+
       \bigl(Q(c_{q+1}t+Y)-Q_0\bigr)
                      \frac{e_b^{\mathsf T}}{(Q_0)_b}\right)Q_0.
 \tag{P14.8}
\]
Thus endpoints can be absorbed into two additional univariate
matrix factors with invertible constants, without increasing
the width or repeating any variable inside this bridge.

For completeness, a product with all constant factors invertible
can be normalized so that every constant is $I$. If its factors
are $M_i(Y_i)$, put $C_i=M_i(0)$, $P_i=C_1\cdots C_i$, $P_0=I$,
and
\[
 N_i(Y_i)=P_{i-1}M_i(Y_i)P_i^{-1}.
\]
Then $N_i(0)=I$ and
$N_1\cdots N_s=M_1\cdots M_sP_s^{-1}$; absorb $P_s$ into
the final constant column.

Let $W_i$ be the coefficient-row span of
$L_0N_1\cdots N_i$ over $F(t)$. Because $N_i(0)=I$,
$W_{i-1}\subseteq W_i$. Let $g_i$ count the strict inclusions
up to position $i$. Induction shows that $W_i$ is spanned by
coefficient rows whose monomials use at most $g_i$ of the
variables. If there is no growth, use the already available
spanning rows with exponent zero in $Y_i$. If there is growth,
multiply the preceding spanning rows by coefficient matrices
of $N_i$; this uses at most one additional variable. Since
$\dim W_0=1$ and $\dim W_i\leq w$, one has $g_i\leq w-1$.
Consequently a nonzero scalar contraction of the final row
has a nonzero coefficient supported on at most $w-1$ variables.

The standard interpolation map with $w-1$ pairs hits any
polynomial with such a coefficient. Indeed, set its selectors
to the abscissas of that support and set any unused amplitudes
to zero. The map then keeps the chosen coordinates as independent
amplitudes and sets every other residual coordinate to zero.
The resulting coordinate restriction is nonzero. This proves
that $b(G)\ne0$, allowing the variables absent from the bridge
to retain their original indices among all $n$ Lagrange nodes.

\paragraph{2. Exact splitting at every rank-one factor.}
Return to (P14.1), and first suppose there is at least one rank-one
factor. Denote its positions by $i_1<\cdots<i_s$ in the program
order. For each such factor choose a nonzero entry $a_j(X)$,
in row $p_j$ and column $q_j$, and put
\[
 U_j(X)=A_{i_j}(X)_{*,q_j},\qquad
 V_j(X)=A_{i_j}(X)_{p_j,*}.
\]
All two-by-two minors vanish, so the following is an identity
over the univariate rational function field:
\[
 A_{i_j}(X)=\frac{U_j(X)V_j(X)}{a_j(X)}.
 \tag{P14.9}
\]
Every $U_j,V_j$ is a nonzero polynomial vector, and $a_j\ne0$.
Writing $B^{(j)}$ for the product of the full-rank factors
strictly between consecutive rank-one positions, and including
the two exterior products, gives the polynomial identity
\[
 \left(\prod_{j=1}^{s}a_j(x_{\pi(i_j)})\right) f(x)
   =\bigl(uB^{(0)}U_1\bigr)
      \left(\prod_{j=1}^{s-1}V_jB^{(j)}U_{j+1}\right)
       \bigl(V_sB^{(s)}v\bigr).
 \tag{P14.10}
\]
The displayed factors on the right are scalar polynomials.
Since $f\ne0$ and the coefficient ring is a domain, each one
is nonzero. Each is a bridge of the form (P14.6), including
the case of an empty interior.

The variable at a rank-one reset appears in two adjacent
bridges, through $U_j$ and $V_j$. This causes no independence
assumption: the bridge lemma proves that each factor remains
a nonzero polynomial under the \emph{same} substitution $G$.
Their product is therefore nonzero in
$F[t,\mathbf y,\mathbf z]$. Moreover every pivot denominator
$a_j(G_{\pi(i_j)})$ is nonzero because setting all amplitudes
to zero gives $a_j(c_{\pi(i_j)}t)\ne0$. Substituting into the
polynomial identity (P14.10) now implies $f(G)\ne0$ directly.
There is no evaluation at a pole and no need to choose pivot
entries algorithmically for the black-box sampler.

If there are no rank-one factors, (P14.1) itself is one bridge
with constant endpoints and generically invertible interior.
The local lemma applies. When $w=1$ every nonzero factor is
already a scalar univariate polynomial; their product remains
nonzero on $x_i=c_it$, also covering the overlap of the two
rank alternatives in this case.

\paragraph{3. Coefficient rows, finite sampling, and rational form.}
Every coefficient row of $R(G)$ belongs to $V$, because the
substitution has coefficients in $F$. If their span were a
proper subspace of $V$, a constant column $v\in F^w$ would
annihilate all these rows while not annihilating $V$. Then
$R(x)v\ne0$ but $R(G)v=0$, contradicting the scalar result.
This proves (P14.3). Over an extension containing the nodes,
the same assertion is $V\otimes_F E$ equality over $E$.

For finite sampling, $f(G)$ has degree at most $D$ in $t$,
degree at most $(n-1)D$ in each selector $y_j$, and total degree
at most $D$ in all amplitudes $z_j$ together. Use a grid of
$D+1$ distinct center values, a grid of $(n-1)D+1$ values for
each selector, and the standard total-degree interpolation
simplex of size $\binom{D+K}{K}$ for the amplitudes. One
description of the last set is
$(b_{e_1},\ldots,b_{e_K})$ with $e_i\geq0$ and
$\sum e_i\leq D$, where $b_0,\ldots,b_D$ are distinct.
The triangular Newton basis proves its unisolvence in any
characteristic. This proves (P14.4).

For the rational form, put $A(T)=\prod_i(T-a_i)$ and
$c_i=1/A'(a_i)$. Then
\[
 \frac{\lambda_i(y)}{c_i}=\frac{A(y)}{y-a_i}.
 \tag{P14.11}
\]
Apply (P14.2) to $f(X_1/c_1,\ldots,X_n/c_n)$, whose factor
ranks are unchanged. Its argument in the original coordinates
is $t+\sum_j b_j/(y_j-a_i)$ after the dominant parameter
change $b_j=z_jA(y_j)$. These values are evaluations of a
rational function with numerator and denominator degrees at
most $K$. Conversely a generic rational function of these
degrees has a separable denominator and such a partial-fraction
expansion over an algebraic closure. Hence the two maps have
the same image closure after the stated rescaling. This
proves (P14.5); scalar contractions give its row-span version
in exactly the same way. Multiplication by a common nonzero
denominator does not alter the constant linear annihilators
of a vector of rational functions.

Finally a nonzero width-two program cannot contain an
identically zero matrix factor, and every nonzero two-by-two
factor has generic rank one or two. Thus $K=1$ applies to
all width-two programs. For $n\geq2$, the polynomial
$x_1/c_1-x_2/c_2$ has a width-two representation by the two
unipotent factors $I+(x_1/c_1)E_{12}$ and
$I-(x_2/c_2)E_{12}$, with endpoints $e_1^{\mathsf T},e_2$.
It vanishes on the common center, establishing the asserted
failure for $K=0$. $\square$

\paragraph{Scope and provenance.}
The support argument for invertible interiors and the endpoint
absorption are written out above; they are the elementary
mechanisms used in I-17 and I-30. The deduction here uses
rank-one factorization and simultaneous product nonvanishing
to replace a full matrix-image requirement by scalar bridge
conditions. It improves I-29's $r=1$ pair budget from $3w-3$
to $w-1$ under the stated factor-rank hypothesis. For a
nonzero scalar program that hypothesis also forces every
interval rank into $\{1,w\}$: an interval containing a rank-one
factor has rank at most one and cannot be identically zero,
because that would make the whole scalar output zero.

The rank-one splitting is the dimension-independent form of
the established width-two reduction in Agrawal, Gurjar,
Korwar, and Saxena (2015), author-hosted version,
\S5.4, Lemma~34, printed pp.~24--25. Their Lemma~35,
printed p.~25, gives product closure by interpolation, and
Theorem~37, printed p.~26, proves polynomial-time black-box
PIT for width two. Those complete local arguments were read.
Here endpoint absorption and coefficient-row growth give the
stated common centered generator directly, so its product
closure requires no additional interpolation of a finite set.

The rational-function interpretation uses the primary
definitions and interpolation relationship in Hu, van
Melkebeek, and Morgan (2024), Definition~1.2 and
Proposition~2.2, printed pp.~5 and 16. Their Theorem~1.9,
printed p.~10, has an additional low-degree-monomial
hypothesis; it is not being used to justify the arbitrary-degree
statement here. No literature-priority claim is made for
this particular generator or the bridge formulation.
For growing $w$, (P14.4) is $(nD+1)^{O(w)}$, not polynomial
in $w$. Arbitrary intermediate factor ranks remain outside
this proof. No purity, GRH, or approximate energy equivalence
is required.


\subsection{P-15: a width-three obstruction to one centered interpolation pair}
\paragraph{Statement P-15.}
Work in characteristic zero. Put
\[
 F(x_1,x_2,x_3,x_4)
   =2x_1x_2-(x_1+x_2)(x_3+x_4)+2x_3x_4.
 \tag{P15.1}
\]
This nonzero rational polynomial has minimum ROABP width exactly
three, even when the variable order may be chosen freely. It has
a representation in which every matrix factor, and hence every
nonempty interval matrix product, is invertible.

For the four Lagrange nodes
\[
 (a_1,a_2,a_3,a_4)=(1,-1,\mathrm i,-\mathrm i),
 \qquad \mathrm i^2=-1,
 \tag{P15.2}
\]
the standard centered one-pair generator fails:
\[
 F\bigl(t+z\lambda_1(y),\ldots,t+z\lambda_4(y)\bigr)=0.
 \tag{P15.3}
\]
There are also balanced rational-function evaluation failures.
For either ordered set of abscissas
\[
 (1,-1,\mathrm i,-\mathrm i),\qquad (0,1,2,2/3),
 \tag{P15.4}
\]
one has
\[
 F\left(\frac{\alpha a_1+\beta}{\gamma a_1+\delta},
        \ldots,
        \frac{\alpha a_4+\beta}{\gamma a_4+\delta}\right)=0
 \tag{P15.5}
\]
as a rational identity. Thus balanced $\mathrm{RFE}_1^1$ does not
hit this width-three class. For standard integer abscissas
$a_i=i-1$, $n\geq7$, the same failure is witnessed by
\[
 F_{\rm int}
   =2x_1x_4-(x_1+x_4)(x_3+x_7)+2x_3x_7,
 \tag{P15.6}
\]
which uses the harmonic ordered subset $(0,3,2,6)$.

These are distinct coordinate conventions. Equation (P15.3)
concerns the ordinary center $t\mathbf1$ and the four nodes
(P15.2). Equation (P15.6) concerns balanced rational evaluation
at the full standard integer grid. It does not by itself give
a counterexample to the ordinary centered one-pair Lagrange
generator at that integer grid.

\paragraph{Proof: the exact width and rank profile.}
Let $E_{ab}$ be the three-by-three matrix unit. Define
\[
 \begin{aligned}
 M_j(x_j)&=I+u_jx_jE_{12}+v_jx_jE_{23},\\
 (u_1,u_2,u_3,u_4)&=(1,1,-2,0),\\
 (v_1,v_2,v_3,v_4)&=(0,2,-1,-1).
 \end{aligned}
 \tag{P15.7}
\]
Every $M_j$ is upper triangular with diagonal entries one,
so every interval product is invertible. In the $(1,3)$ entry
of the product, a path must take an $E_{12}$ step in an earlier
factor and an $E_{23}$ step in a later factor. Therefore
\[
 e_1^{\mathsf T}M_1M_2M_3M_4e_3
    =\sum_{i<j}u_iv_jx_ix_j=F.
 \tag{P15.8}
\]
This proves width at most three.

For the middle cut $\{1,2\}\mid\{3,4\}$, the left coefficient
rows include the three independent right polynomials
$2x_3x_4$, $-(x_3+x_4)$, and $2$. Their distinct degrees prove
coefficient-matrix rank three, and the fourth possible row is
a repetition. For either other balanced partition, the
degree-one cross block is, up to row and column permutations,
\[
 \begin{pmatrix}2&-1\\-1&2\end{pmatrix},
\]
of determinant three, while the degree-zero and degree-two
blocks both have rank one. Thus those cuts have rank four.
Every order has one of these three balanced partitions at its
middle cut, so every order needs width at least three. The
displayed order has width exactly three.

\paragraph{Proof: ordinary centered interpolation.}
For the nodes (P15.2), Lagrange interpolation gives
\[
 \lambda_a(y)=\frac14
    \left(1+\frac{y}{a}+\frac{y^2}{a^2}+\frac{y^3}{a^3}\right).
 \tag{P15.9}
\]
The polynomial $F$ is invariant under common translation and
homogeneous of degree two. More explicitly, if
$S=x_1+x_2-x_3-x_4$, $U=x_1-x_2$, and $V=x_3-x_4$, then
\[
 2F=S^2-U^2-V^2.
 \tag{P15.10}
\]
At the four interpolants,
\[
 S=y^2,\qquad U=\frac{y^3+y}{2},\qquad
 V=\frac{\mathrm i(y^3-y)}2.
\]
Consequently $S^2=U^2+V^2$, so $F(\lambda(y))=0$.
Translation invariance and homogeneity now give (P15.3).
This is a polynomial identity over $\mathbb Q(\mathrm i)$;
the annihilated polynomial itself has rational coefficients.

\paragraph{Proof: balanced rational evaluation.}
The alternative expression
\[
 F=(x_1-x_3)(x_2-x_4)+(x_1-x_4)(x_2-x_3)
 \tag{P15.11}
\]
makes the rational identity immediate. For
$r(T)=(\alpha T+\beta)/(\gamma T+\delta)$,
\[
 r(a)-r(b)=
 \frac{(\alpha\delta-\beta\gamma)(a-b)}
      {(\gamma a+\delta)(\gamma b+\delta)}.
 \tag{P15.12}
\]
Every ordered quadruple in (P15.4), and also $(0,3,2,6)$,
satisfies the harmonic relation
\[
 (a_1-a_3)(a_2-a_4)+(a_1-a_4)(a_2-a_3)=0.
 \tag{P15.13}
\]
Substitution of (P15.12) into (P15.11) proves (P15.5).
The identity holds on the allowed open set where denominators
are nonzero; singular or constant rational maps cause no
exception. The subset $(0,3,2,6)$ gives (P15.6), with the
other factors in a width-three program taken to be identities.
$\square$

\paragraph{What this closes and what it leaves open.}
The obstruction already occurs with a full-rank matrix word,
so singular rank resets are not needed for failure of one pair.
For these specified node choices, P-14 supplies two pairs for
the full-rank width-three subclass, and (P15.3) proves that one
does not suffice. This disproves a node-uniform one-pair
strengthening of P-14 to width three.

It does not refute balanced $\mathrm{RFE}_{w-1}^{w-1}$ for
general width $w$, nor prove that two pairs suffice for all
width-three programs containing rank-two resets. The broader
width-three construction in I-30 remains valid with six pairs.
The balanced rational ideal and its degree-two generators are
the primary-source framework of Hu, van Melkebeek, and Morgan
(2024), Theorem~1.3, printed p.~6; the identities above are
independently written out and do not assume a general width
theorem from that source. No literature-priority claim is made.


\subsection{P-16. An explicit width-three obstruction to uncentered two-pair interpolation}
\label{P-16}

\paragraph{Statement and scope.}
Work in characteristic zero. There are eighteen explicitly specified,
pairwise distinct algebraic Lagrange nodes and a nonzero multilinear
polynomial $H$ of degree $16$, of minimum ROABP width exactly three,
such that
\[
 H\bigl(\operatorname{SV}_2(\mathbf y,\mathbf z)\bigr)=0.
 \tag{P16.1}
\]
Its displayed width-three representation has every univariate factor
generically invertible. In contrast, P-14 proves
\[
 H\bigl(t\mathbf1+\operatorname{SV}_2(\mathbf y,\mathbf z)\bigr)\ne0.
 \tag{P16.2}
\]
Thus this is an obstruction to the uncentered map, not to balanced
$\operatorname{RFE}_{2}^{2}$ or to centered two-pair interpolation.
The latter two maps on the full width-three class with rank-two resets
remain separate questions.

We use $\operatorname{RFE}_{k}^{\ell}$ for evaluations of rational
functions whose numerator degree is at most $k$ and denominator degree
at most $\ell$. In fact the construction below gives, for every
$k\ge1$, a quadratic width-three annihilator of
$\operatorname{RFE}_{k}^{1}$ on $6(k+1)$ nodes, and its multilinear
reversal annihilates $\operatorname{RFE}_{1}^{k}$. Only $k=2$ is needed
for (P16.1).

\paragraph{Explicit disjoint pole groups.}
Put $d=k+1$, choose positive real radicals
\[
 \alpha=2^{1/(2d)},\qquad \beta=(3/2)^{1/(2d)},
 \qquad a=\beta^d=\sqrt{3/2},
\]
and let $\zeta$ be a primitive $2d$th root of unity. Work over the
number field
\[
 E=\mathbb Q(\zeta,\alpha,\beta),\qquad
 [E:\mathbb Q]\le(2d)^2\varphi(2d).
 \tag{P16.3}
\]
This bound follows by adjoining the two radicals, each of degree at
most $2d$, to the cyclotomic field. Define
\[
 \begin{array}{ll}
 P_1(T)=T^{2d}-1,&P_2(T)=T^d-a,\\
 P_3(T)=-2(T^d+a),&P_4(T)=T^{2d}-2.
 \end{array}
\]
Let $\mathcal A_i$ be the roots of $P_i$. Explicitly these four
ordered groups can be listed as
\[
 \begin{split}
 \mathcal A_1&=(\zeta^j)_{0\le j<2d},\qquad
 \mathcal A_2=(\beta\zeta^{2j})_{0\le j<d},\\
 \mathcal A_3&=(\beta\zeta^{2j+1})_{0\le j<d},\qquad
 \mathcal A_4=(\alpha\zeta^j)_{0\le j<2d}.
 \end{split}
 \tag{P16.4}
\]
All roots are simple. The first, middle, and last groups have distinct
absolute values $1,\beta,\alpha$, and the two middle root sets are
separated by their $d$th powers $a$ and $-a$. Thus all $6d$ nodes
are distinct. For $k=2$, $d=3$, their field has degree at most $72$.

Set
\[
 R_1=\frac{P_2}{P_1},\qquad R_2=\frac1{P_2},\qquad
 R_3=\frac1{P_3},\qquad R_4=\frac{P_3}{P_4}.
 \tag{P16.5}
\]
Each $R_i$ has precisely the simple poles in $\mathcal A_i$, with
nonzero residues, and satisfies $R_i(T)=O(T^{-d})$ at infinity.
Moreover $P_2P_3=-P_1-P_4$, and hence
\[
 R_1R_2+R_1R_4+R_3R_4
 =\frac{P_4+P_2P_3+P_1}{P_1P_4}=0.
 \tag{P16.6}
\]
For $b\in\mathcal A_i$, define the explicit nonzero residue
$\rho_b=\lim_{T\to b}(T-b)R_i(T)$. Introduce one variable $X_b$ for
each node and put
\[
 L_i(\mathbf X)=\sum_{b\in\mathcal A_i}\rho_bX_b,
 \qquad
 Q=L_1L_2+L_1L_4+L_3L_4.
 \tag{P16.7}
\]
The four variable sets are disjoint, so $Q$ is nonzero, multilinear,
and homogeneous of degree two.

\paragraph{The unbalanced rational-evaluation identity.}
The partial fraction identity
\[
 R_i(T)=\sum_{b\in\mathcal A_i}\frac{\rho_b}{T-b}
\]
and its expansion at infinity give
\[
 \sum_{b\in\mathcal A_i}\rho_b b^j=0
       \quad(0\le j\le d-2=k-1).
 \tag{P16.8}
\]
A generic rational function with numerator degree at most $k$ and
linear denominator has the form
\[
 r(T)=p(T)+\frac{c}{T-u},\qquad \deg p\le k-1.
\]
On the open set on which no evaluation denominator vanishes,
(P16.8) implies
\[
 L_i\bigl((r(b))_b\bigr)=-cR_i(u).
\]
Equation (P16.6) proves that $Q$ vanishes on these evaluations.
Clearing the finitely many evaluation denominators gives a polynomial
identity in the numerator and denominator coefficients. Its vanishing
on the dense open set of genuine linear denominators therefore extends
to all defined evaluations in $\operatorname{RFE}_{k}^{1}$, including
constant denominators. No assertion that a degree-$k$ polynomial part
is killed term by term is needed for this last specialization.

\paragraph{The exact width and its reversal.}
Read the variables in the four group orders in (P16.4). For each node,
use
\[
 M_b(X_b)=
 \begin{cases}
 I_3+\rho_bX_bE_{12},&b\in\mathcal A_1\cup\mathcal A_3,\\
 I_3+\rho_bX_bE_{23},&b\in\mathcal A_2\cup\mathcal A_4.
 \end{cases}
 \tag{P16.9}
\]
The $(1,3)$ entry of their product is exactly $Q$: the possible
ordered $E_{12}E_{23}$ pairs come from groups $(1,2),(1,4),(3,4)$,
and products of three such matrix units are zero. Thus $Q$ has a
width-three unit-factor representation. Its minimum width in any
order cannot be two: it is a nonzero annihilator of balanced
$\operatorname{RFE}_{1}^{1}$, whereas P-14 hits every nonzero
width-two ROABP with that family. Alternatively, its rank in the
specified order is already three at a cut strictly inside group two:
the independent left polynomials $1,L_1,L_1L_{2,\mathrm{left}}$
are paired with independent right polynomials
$L_3L_4,L_{2,\mathrm{right}}+L_4,1$, of three different degrees.

Let $N=6d$, and reverse the coefficient basis in every variable:
\[
 Q^{\mathrm{rev}}(\mathbf X)=
       \left(\prod_bX_b\right)Q\bigl((X_b^{-1})_b\bigr).
 \tag{P16.10}
\]
It is a nonzero multilinear polynomial of degree $N-2$. For every
variable partition, reversal merely permutes the rows and columns
of the coefficient matrix. Hence it preserves every coefficient-cut
rank and the minimum width in every order. Replacing (P16.9) by
$X_bI_3+\rho_bE_{12}$ or $X_bI_3+\rho_bE_{23}$ gives its explicit
width-three word; each factor has determinant $X_b^3$.

If $r$ has numerator degree at most one and denominator degree at
most $k$, then $1/r$ has numerator degree at most $k$ and denominator
degree at most one. On the dense open set where all relevant values
are nonzero and defined, (P16.10) therefore vanishes on $(r(b))_b$.
Clearing denominators extends the identity to every defined
$\operatorname{RFE}_{1}^{k}$ evaluation.

\paragraph{The ordinary SV convention and the repair by centering.}
Specialize now to $k=2$, $d=3$, $N=18$. Let
\[
 A(T)=\prod_b(T-b),\qquad D_b=A'(b)\ne0,
 \qquad \lambda_b(Y)=\frac{A(Y)}{(Y-b)D_b},
 \qquad H(\mathbf X)=Q^{\mathrm{rev}}\bigl((D_bX_b)_b\bigr).
 \tag{P16.11}
\]
For the ordinary two-pair generator
$\operatorname{SV}_{2,b}=z_1\lambda_b(y_1)+z_2\lambda_b(y_2)$, one has
\[
 D_b\operatorname{SV}_{2,b}
   =-\sum_{j=1}^{2}\frac{z_jA(y_j)}{b-y_j}.
 \tag{P16.12}
\]
As a function of the node $b$, the right-hand side is a rational
function with numerator degree at most one and denominator degree
at most two. The proved $\operatorname{RFE}_{1}^{2}$ identity gives
(P16.1) on a dense open parameter set and thus identically. This
also agrees with the fixed diagonal equivalence in Hu et al. (2024),
Proposition~2.2, pp.~16--17 (P-B20).

The displayed factors for $H$ are
$D_bX_bI_3+\rho_bE_{12}$ or $D_bX_bI_3+\rho_bE_{23}$, with determinant
$D_b^3X_b^3$. All are generically invertible, and diagonal variable
scaling preserves all coefficient-cut ranks. Thus no variable order permits width two for $H$, while the displayed
order attains width three. P-14, with $w=3$ and center slopes
all one, applies directly to this word and proves (P16.2).

This construction does not invoke a monotonicity assertion for the
width of factors. Its width claims follow from explicit matrix words
and exact coefficient-basis permutations. It also does not establish
failure of the centered map on rank-two reset words; that is the
remaining balanced question.
\hfill$\square$



\subsection{P-17. Balanced two-pair rank bounds from coefficient cuts and $\mathrm{SL}_2$ symmetry}
\label{P-17}

\paragraph{Statement and scope.}
Fix distinct nodes $a_1,\ldots,a_n$ in a field $E$.
Let $\mathcal I_2$ be the vanishing ideal of evaluations of rational
functions with numerator and denominator degrees at most two.
All assertions below hold in every characteristic, with field
extensions used as proof devices when necessary.

Every nonzero multilinear polynomial in $\mathcal I_2$ on at most
eight variables has ROABP width at least four in every variable
order. The same conclusion holds for every nonzero homogeneous
multilinear cubic in $\mathcal I_2$ on at most ten variables,
and in any number of variables for the decomposable cubic class
specified below. In any number of variables, a nonzero multilinear
width-three polynomial whose multihomogenization is
$\mathrm{SL}_2$-invariant is also hit by balanced two-pair
evaluation, indeed by I-37's ordinary centered two-pair map.

These are statements about the balanced rational generator.
They do not identify it with the ordinary unweighted centered SV
map outside the separately stated invariant subclass. The
proof here leaves a common-translation-invariant nine-variable
quartic. The subsequent I-40 closes that residual, and I-41
proves the full multilinear conclusion through ten variables.
The unrestricted case beyond ten variables remains open.

\paragraph{Known ideal input and the decomposable class.}
Put
\[
 M=\begin{pmatrix}1&\cdots&1\\a_1&\cdots&a_n\\a_1^2&\cdots&a_n^2\end{pmatrix},
 \qquad K=\ker M\subseteq E^n.
 \tag{P17.1}
\]
For a three-element set $S$, use a consistent ordered Vandermonde
factor $\Delta_S=\det M_S$. Its sign convention is immaterial as
long as it is used consistently. Hu, van Melkebeek, and Morgan
(2024), Definition~9.5, p.~55, and Theorem~9.13, pp.~61--62
(P-B20), give the following exact input. A homogeneous multilinear
cubic
\[
 f=\sum_{|S|=3}c_S X_S
\]
lies in $\mathcal I_2$ if and only if its decoded alternating tensor
\[
 W=\sum_{|S|=3}\frac{c_S}{\Delta_S}e_S
       \quad\hbox{belongs to }\quad\bigwedge^3 K.
 \tag{P17.2}
\]
The same theorem says that the multilinear homogeneous degrees in
$\mathcal I_2$ lie in the interval $3\le d\le n-3$. We use these
published characterizations, rather than reprove them.

The decomposable class consists of the nonzero tensors
$W=p\wedge q\wedge v$ in (P17.2). If $Z$ is the $3\times n$
matrix with these three rows, the corresponding polynomial is,
by Cauchy--Binet,
\[
 f=\det\bigl(M\operatorname{diag}(X_1,\ldots,X_n)Z^{\mathsf T}\bigr),
 \qquad MZ^{\mathsf T}=0.
 \tag{P17.3}
\]
Every three columns of $M$ are independent. In particular, every
nonzero vector of $K$ has at least four nonzero coordinates.

\paragraph{1. A three-block cut obstruction.}
Let $f\in\mathcal I_2$ be a nonzero homogeneous multilinear cubic.
For any variable partition $U\sqcup V$, its coefficient matrix is
block diagonal after grouping row and column monomials by degree.
The block $C_d$ has left degree $d$ and right degree $3-d$.
We claim that three consecutive nonzero blocks $C_0,C_1,C_2$, or
$C_1,C_2,C_3$, force total coefficient rank at least four.

Suppose instead that $C_0,C_1,C_2$ are nonzero and the total rank
is at most three. Then $C_3=0$ and each of those three blocks has
rank one. Consequently
\[
 f=A(V)+L(U)B(V)+Q(U)N(V),
 \tag{P17.4}
\]
where $L,N$ are nonzero linear forms, $B,Q$ are nonzero homogeneous
quadratics, and $Q$ is multilinear. The pole-extraction identity
proved in I-37 gives
$\partial_i f\in\mathcal I_1$, where $\mathcal I_1$ denotes the
balanced one-pair ideal. The ideal $\mathcal I_1$ is prime and
contains no nonzero linear form: it is the kernel of a rational
parameter substitution, and the distinct simple poles
$1/(a_i-y)$ are linearly independent. Thus it contains no
nonzero product of two linear forms.

For every constant direction $b$ on $U$ with $b(L)=0$,
\[
 b(f)=b(Q)N\in\mathcal I_1,
\]
so $b(Q)=0$. Write $L=\ell^{\mathsf T}X_U$ and let $T$ be the
symmetric zero-diagonal coefficient matrix of the multilinear
quadratic $Q$, so that $b(Q)=(Tb)^{\mathsf T}X_U$ in every
characteristic. We have $\ker\ell^{\mathsf T}\subseteq\ker T$.
Therefore $T=u\ell^{\mathsf T}$ for a column $u$; symmetry implies
$u=c\ell$. Its zero diagonal then implies $c\ell_i^2=0$ for every
$i$, hence $c=0$, since $\ell\ne0$. Thus $T=0$ and $Q=0$, a
contradiction. This avoids division by two. Exchanging $U,V$
proves the claim for $C_1,C_2,C_3$.

\paragraph{2. Every decomposable cubic has a cut of rank at least four.}
Take (P17.3) and fix any variable order. Let $z_i\in E^3$ denote
column $i$ of $Z$. Choose the first prefix $U$ whose columns span
a two-dimensional space. Immediately before its last column,
the prefix span has dimension one. The remaining suffix $V$
has column rank at least two. Otherwise a nonzero linear
functional on $E^3$ annihilating the previous prefix and the
suffix would yield a nonzero row combination of $Z$ supported
only on that last prefix coordinate. This contradicts the
minimum support four of nonzero vectors in $K$.

If the suffix has rank three, there are nonzero monomial
coefficients of left degrees zero, one, and two. Indeed, choose
respectively three independent suffix columns, one nonzero
prefix column and two completing suffix columns, or two
independent prefix columns and one completing suffix column.
The Vandermonde factors in (P17.2) do not vanish. Part~1 therefore
forces coefficient rank at least four.

It remains to consider the case in which both projections of
the row space of $Z$ have rank two. Its kernels under these two
projections are distinct lines. Choose a row basis $p,q,v$ in
which $p$ is supported on $U$ and $q$ on $V$. Both $p$ and $q$
belong to $K$ and have at least four nonzero coordinates.
Moreover $v_U$ is not proportional to $p_U$, and $v_V$ is not
proportional to $q_V$, since the two projection ranks are two.

For $i\in U$ and $j,k\in V$, the left-degree-one coefficients
are, up to one consistent sign,
\[
 p_i(q_jv_k-q_kv_j)(a_j-a_k)(a_i-a_j)(a_i-a_k).
 \tag{P17.5}
\]
Evaluation of a polynomial of degree at most two at the nonzero
coordinates of $p$, followed by multiplication by those $p_i$,
is injective, because their nodes include at least four distinct
points. It follows that
\[
 \operatorname{rank}C_1
 =\dim\operatorname{span}\{
      (T-a_j)(T-a_k):q_jv_k-q_kv_j\ne0\}.
 \tag{P17.6}
\]
There are at least two distinct pairs with nonzero coefficients
in this set. If only one pair were nonzero, the decomposable
nonzero two-vector $q_V\wedge v_V$ would be supported on that
pair, so its two-dimensional row space, including $q_V$, would
be supported on only two coordinates. This contradicts the
support bound for $q$. Two distinct monic quadratics with
these distinct two-element root sets are linearly independent.
Thus $\operatorname{rank}C_1\ge2$.

The symmetric argument, using the nonzero coordinates of $q$
and the independent pair $p_U,v_U$, gives
$\operatorname{rank}C_2\ge2$. The degree blocks have disjoint
row and column indices, so the total coefficient rank is at
least four. This proves the decomposable assertion in every
order and for every $n$.

\paragraph{3. At most seven variables.}
For $n\le5$, the degree interval supplied by Theorem~9.13 is
empty. For $n=6$, the only possible degree is three and
$\dim K=3$, so every nonzero three-vector is decomposable.
For $n=7$, the possible degrees are three and four and
$\dim K=4$; every nonzero three-vector in a four-dimensional
space is decomposable. Thus part~2 covers every cubic.

Multilinear reversal
\[
 f^{\mathrm{rev}}(X)=\Bigl(\prod_iX_i\Bigr)f(X_1^{-1},\ldots,X_n^{-1})
 \tag{P17.7}
\]
preserves the balanced ideal, since reciprocal rational functions
still have numerator and denominator degrees at most two.
It preserves each coefficient-cut rank by complementing the
row and column subset indices, as already used in P-16. For
$n=7$, it turns a homogeneous quartic into a cubic, covering
that case as well.

Finally the ideal decomposes by homogeneous degree. The lowest
nonzero homogeneous component of a polynomial has ROABP width
no greater than the original in the same order (Hu et al., 2024,
Proposition~8.3, p.~42). Hence a nonhomogeneous multilinear
annihilator of width at most three would have a nonzero cubic
or quartic component of width at most three, contradicting the
preceding cases. This proves the full seven-variable assertion.

\paragraph{4. The additional eight-variable cubic case.}
Let $n=8$ and suppose that a cubic annihilator has width at most
three in a fixed order. Take the first three variables as $U$
and the remaining five as $V$. The restriction map
$K\to E^V$ is an isomorphism: its kernel would be a vector of
$K$ supported on three coordinates, and both spaces have
dimension five. Its third exterior power is also an
isomorphism. Therefore $C_0\ne0$ for every nonzero $W$.

We first show that $C_2=0$. If $C_1=0$, fix $i\in U,k\in V$.
The quadratic $\partial_i f\in\mathcal I_1$ has both zero row
sum and zero simple-pole residue at node $k$. Since $C_1=0$,
these two equations involve only the two possible indices
$j\in U\setminus\{i\}$:
\[
 \sum_j c_{ijk}=0,\qquad
 \sum_j\frac{c_{ijk}}{a_k-a_j}=0.
 \tag{P17.8}
\]
Their two distinct reciprocal weights make the coefficient
matrix invertible. Thus all these $c_{ijk}$ vanish and $C_2=0$.
If instead $C_1\ne0$ and $C_2\ne0$, part~1 contradicts the
rank bound because $C_0\ne0$. This proves $C_2=0$ in all cases.
The zero row-sum identity for a left pair then gives $C_3=0$,
since the sole possible third left index is its only remaining
term. The row-sum and residue identities used here follow
by substituting $t+z/(a_i-y)$ into the relevant quadratic
and taking the coefficient of $tz$ and the simple-pole
residues of the coefficient of $z^2$; no division by two is
used.

Put $K_0=K\cap E^V$, which has dimension two, and choose a
complement $J$ in $K$. The projection $J\to E^U$ is an
isomorphism. In the decomposition of $\bigwedge^3(K_0\oplus J)$,
the highest left-coordinate degree of $\bigwedge^3J$ is three,
and that of $K_0\wedge\bigwedge^2J$ is two; these leading
projections are injective. The absence of all coefficients
with at least two left variables therefore forces
\[
 W\in\bigwedge^2K_0\wedge J.
 \tag{P17.9}
\]
Because $\dim K_0=2$, every nonzero tensor in (P17.9) is a
wedge of a basis of $K_0$ with one additional vector. Thus
$W$ is decomposable. Part~2 now contradicts the width bound.
This proves the eight-variable cubic assertion.

\paragraph{5. The nine- and ten-variable cubics from two end anchors.}
The argument at the start of part~4 holds for any number of
variables. If $U$ is any three-element set and $V$ its complement,
the projection $K\to E^V$ is an isomorphism: it is injective by
the three-column independence of $M$, and both spaces have
dimension $n-3$. Thus the block $C_0$ is nonzero for every
nonzero cubic tensor $W$.

If this cut has total coefficient rank at most three, the same
two-equation argument (P17.8) and part~1 force $C_2=C_3=0$.
For $K_U=\ker(K\to E^U)$, choose a complement projecting
isomorphically to $E^U$. The leading-projection argument of
part~4 then gives
\[
 W\in\bigwedge^2K_U\wedge K.
 \tag{P17.12}
\]
This includes its possible component in $\bigwedge^3K_U$;
no decomposability is claimed for arbitrary $n$.

Now let $n=9$ and fix a putative width-three order. Apply this
anchor to its first three and its last three variables. The
corresponding kernels $K_L,K_R$ each have dimension three.
Their intersection is zero, since a vector in it is supported
on the middle three variables and belongs to $K$. Therefore
\[
 K=K_L\oplus K_R.
\]
The first anchor requires at least two factors from $K_L$ in
the third exterior power, and the second requires at least
two from $K_R$. These are disjoint summands in
\[
 \bigwedge^3(K_L\oplus K_R)
 =\bigoplus_{j=0}^3
      \bigwedge^jK_L\otimes\bigwedge^{3-j}K_R.
\]
Their intersection is zero, contradicting $W\ne0$.
This proves the nine-variable cubic assertion in every order.

For $n=10$, the two end kernels have dimension four and their
intersection $P$ is the one-dimensional kernel supported on the
middle four coordinates. Write
\[
 K=P\oplus A\oplus B,\qquad
 K_R=P\oplus A,\quad K_L=P\oplus B,
 \tag{P17.13}
\]
where $A$ and $B$ each have dimension three. The first-coordinate
projection maps $A$ isomorphically to the first three coordinates,
and the last-coordinate projection does the same for $B$ and the
last three. In the exterior-power decomposition associated with
(P17.13), the first anchor requires at least two factors from
$P\oplus B$, and the last anchor at least two from $P\oplus A$.
Since $\dim P=1$, their intersection is exactly
$P\wedge A\wedge B$. Choose $0\ne p\in P$ and bases
$A_1,A_2,A_3$, $B_1,B_2,B_3$ whose respective outer projections
are the coordinate bases. Then
\[
 W=p\wedge\sum_{i,j=1}^3 Z_{ij}A_i\wedge B_j,
 \qquad Z\ne0.
 \tag{P17.14}
\]
Every one of the four middle coordinates of $p$ is nonzero,
by the minimum support four of vectors in $K$.

Take the cut after the first four variables. Fix $Z_{ij}\ne0$.
For each middle coordinate $k$, the coefficient of the triple
consisting of first coordinate $i$, middle coordinate $k$, and
last coordinate $j$ is, up to its ordered sign,
$p_kZ_{ij}\Delta_{\{i,k,j\}}\ne0$. Using the first and second
middle coordinates gives respectively a nonzero entry in $C_2$
and in $C_1$. If $C_0\ne0$, part~1 contradicts the cut-rank
bound. If $C_0=0$, the projection from $K$ to the last six
coordinates has a one-dimensional kernel $Eq$, consisting of
the vectors supported on the first four coordinates. Its
induced map on $K/Eq$ is injective, so the kernel of its third
exterior power is exactly $q\wedge\bigwedge^2K$. Consequently
$W$ is divisible by both $p$ and $q$. The two lines are distinct,
because their supports intersect in only one coordinate and
each nonzero vector of $K$ has support at least four. Thus
\[
 W\in\left(p\wedge\bigwedge^2K\right)
       \cap\left(q\wedge\bigwedge^2K\right)
   =p\wedge q\wedge K.
\]
The nonzero tensor $W$ is therefore decomposable, contradicting
part~2. This proves the ten-variable cubic assertion in every
order and characteristic.

\paragraph{6. An invariant multilinear width-three tensor has a singular affine representation.}
Work over an algebraic closure; polynomial identities and
coefficient-cut ranks are unaffected by this field extension.
We use the characteristic-free weight decomposition for the split
torus $\mathbb G_m$: its rational representations are
$\mathbb Z$-graded vector spaces. The exact primary reference is
SGA~3, Expos\'e~I, Proposition~4.7.3 (G-B02), whose proof is the
orthogonal-projector decomposition of its comodule map. This
established torus fact is used here, not reproved.

For a multilinear polynomial $f$ in an even number $n$ of variables,
write its multihomogenization as
\[
 \widetilde f((U_i,V_i)_{i=1}^n)
   =\Bigl(\prod_i U_i\Bigr)f(V_1/U_1,\ldots,V_n/U_n).
 \tag{P17.10}
\]
Suppose that this tensor is invariant under the simultaneous
$\mathrm{SL}_2$ action on all pairs, and that its coefficient-cut
ranks in some order are at most three. For a cut after $i$
variables, its left Schmidt space is the image of contraction
with the dual of the right tensor factors. Invariance makes
this image an $\mathrm{SL}_2$ submodule of $(E^2)^{\otimes i}$,
and its dimension is exactly the coefficient-cut rank $r_i$.
Restrict this Schmidt space to the diagonal torus
$\operatorname{diag}(t,t^{-1})$. The ambient tensor power has
integer character weights $i,i-2,\ldots,-i$. For odd $i$, all
these weights are odd and none is zero. The Weyl element
\[
 w=\begin{pmatrix}0&-1\\1&0\end{pmatrix}\in\mathrm{SL}_2
\]
conjugates the torus parameter $t$ to $t^{-1}$, and hence gives
an isomorphism between the weight-$r$ and weight-$-r$ subspaces.
The nonzero integer weights pair without fixed points, so the
Schmidt-space dimension is even. These are torus characters,
not eigenvalues reduced modulo the field characteristic; the
argument also holds in characteristic two. Since $r_i\le3$
and $f\ne0$, we obtain $r_i=2$ at every odd cut.

Use the standard minimal tensor-train factorization with bond
sizes $r_i$ (P-B17, Theorem~2.1), with $r_0=r_n=1$.
Each local core is affine in its one variable and has matrix
shape $r_{i-1}\times r_i$. One of its two adjacent cuts is odd,
so one side has dimension at most two. Padding these cores
consistently to $3\times3$ matrices gives an affine width-three
word whose every variable-dependent factor has rank at most
two. I-37 therefore hits $f$ with two centered pairs, both in
its ordinary and its rational formulations. This proves the
stated invariant subclass for every even $n$.

\paragraph{7. The eight-variable quartic in every characteristic.}
Suppose that a homogeneous multilinear quartic $f\in\mathcal I_2$
on eight variables has width at most three. Balanced rational
evaluation is closed under adding a constant, so every coefficient
of $f(X+t\mathbf1)$, as a polynomial in $t$, lies in
$\mathcal I_2$. Its coefficient of $t^j$ has homogeneous degree
$4-j$. The known minimum degree three therefore gives
\[
 f(X+t\mathbf1)=f(X)+t g(X),\qquad
 g\in\mathcal I_2\text{ homogeneous of degree three}.
 \tag{P17.11}
\]
If $g\ne0$, it has width at most three in the same order.
Indeed translating each univariate factor preserves the cut
ranks over $E(t)$; multiply the coefficient matrices by $t^{-1}$
and specialize $t^{-1}$ to zero. All minors of size four still
vanish, so the limiting coefficient matrices of $g$ have rank
at most three. Part~4 excludes such a nonzero $g$. Hence
$f(X+t\mathbf1)=f(X)$.

The reversal $f^{\mathrm{rev}}$ is also a homogeneous quartic
annihilator of width at most three. Applying the same argument
to it proves its common-translation invariance. On the
multihomogenization (P17.10), common translation is one
unipotent subgroup of $\mathrm{SL}_2$, while reversal conjugates
it to the opposite unipotent subgroup. Those two subgroups
generate $\mathrm{SL}_2$. Thus $\widetilde f$ is invariant,
and part~6 contradicts its membership in $\mathcal I_2$.

The only possible homogeneous degrees on eight variables are
three, four, and five. Part~4 excludes degree three, reversal
excludes degree five, and the preceding argument excludes
degree four. Passing to a lowest nonzero homogeneous component
as in part~3 proves the full eight-variable conclusion.\hfill$\square$

\paragraph{The remaining balanced-ideal gate.}
The degree-zero-block anchor holds uniformly for any $n$.
The extra four-coordinate cut resolves the one-dimensional
intersection of the two end kernels at $n=10$. At $n=11$ their
intersection has dimension two, so that argument no longer
forces the same decomposable tensor. The unrestricted cubic
case from eleven variables onward remains open here.

On nine variables, the possible homogeneous degrees are three
through six. Part~5 and reversal exclude degrees three and six.
A putative width-three annihilator therefore reduces, after
homogeneous-component extraction and possibly reversal, to a
quartic. The translation-leading argument of part~7, now using
the nine-variable cubic bound, makes that quartic invariant under
common translation. Its multihomogenization has nonzero torus
weight, so the invariant-tensor argument does not automatically
apply. This common-translation-invariant nine-variable quartic
is the smallest case left by the proof.

The nine-variable residual of this proof is subsequently closed
by I-40, and I-41 proves the full multilinear balanced two-pair
theorem through ten variables. The unrestricted case beyond ten
variables remains a separate gate. The proof does not infer
cut-rank bounds from a fixed linear change of coordinates to a
Hankel ideal; such a change need not preserve ROABP coefficient
cuts. The all-characteristic odd-weight argument and the
ten-variable cubic extension were supplied independently by
the isolation agent during the proof review,
replacing the author's initial characteristic-zero
complete-reducibility argument. No historical priority claim
is made.



\subsection{P-18. An arbitrary-length trilinear class and a two-boundary cubic gate}
\label{P-18}

\paragraph{Statement.}
Fix distinct nodes $a_1,\ldots,a_n$ over any field $E$.
Suppose that a nonzero multilinear polynomial $f$ has width at
most three in some variable order and is separately homogeneous
of degree one in each block of a partition of all its variables
into three disjoint nonempty consecutive blocks of that same
order. The blocks and order need not be supplied to
the evaluator. Then balanced rational two-pair evaluation
\[
 R_i=t+\frac{z_1}{a_i-y_1}+\frac{z_2}{a_i-y_2}
 \tag{P18.1}
\]
hits $f$. This gives a uniform finite hitting set of at most
\[
 20(n+1)^2
 \tag{P18.2}
\]
points over a working field of size at least
$\max\{2n+1,4\}$, with the explicitly stated extension-input
convention of I-40. The assertion holds in every characteristic.

There is also a precise reduction of the unrestricted cubic
question. If a homogeneous multilinear cubic of width at most
three annihilates (P18.1), choose such an annihilator with the
fewest essential variables and a witnessing order. Its support
triples have their first, second, and third positions in three
consecutive intervals that overlap in at most one variable at
each interface. At least one overlap must occur. Thus at most
two shared boundary variables are the remaining obstruction
in this reduced cubic problem; this statement does not resolve
them or the unrestricted cubic theorem.

\paragraph{Inputs and notation.}
Write $\mathcal I_k$ for the formal vanishing ideal of balanced
rational functions with numerator and denominator degrees at
most $k$. These ideals are prime kernels of rational parameter
substitutions. We use the established HMM degree and exterior
characterizations (P-B20, Definition~9.5, p.~55, and
Theorem~9.13, pp.~61--62), the pole extraction proved in I-37,
and the three-consecutive-degree-block lemma of P-17, part~1.
In particular $\mathcal I_2$ has no nonzero multilinear component
of degree less than three, $\mathcal I_1$ contains no nonzero
linear form, and every partial derivative of a cubic in
$\mathcal I_2$ belongs to $\mathcal I_1$.
The argument below uses those inputs rather than reproving
rational-function ideal membership.

\paragraph{1. A balanced cubic that is trilinear in disjoint blocks cannot have width three.}
Suppose for contradiction that $f\in\mathcal I_2$ satisfies
the stated disjoint-block promise, with blocks $U,V,W$ in
that order. Remove unused variables; this preserves both the
width bound and formal generator membership at the retained
nodes. A nonzero cubic in $\mathcal I_2$ is invariant under
common translation: every positive-power coefficient of
$f(X+s\mathbf1)$ still belongs to $\mathcal I_2$ and has
multilinear degree at most two, so it is zero.

The coefficient of $s$ is the sum of the derivatives in the
three blocks. Those three summands have distinct block degrees
$(0,1,1)$, $(1,0,1)$, and $(1,1,0)$, respectively. Each
summand is therefore zero. Since the degree in each entire
block is one, this also gives separate common-translation
invariance within each block. This uses no division by the
characteristic.

Let $r$ be the coefficient rank across $U\mid(V,W)$. Write
\[
 f=\sum_{i\in U}X_i B_i(V,W),\qquad
 r=\dim\operatorname{span}\{B_i:i\in U\}.
 \tag{P18.3}
\]
Separate translation in $U$ gives $\sum_i B_i=0$. Let $i_*$
be the last variable of $U$ in the witnessing order. It is
essential, so $B_{i_*}\ne0$, and the earlier $B_i$ span the
full $r$-dimensional space because their sum is $-B_{i_*}$.
At the cut just before $X_{i_*}$, the rows indexed by earlier
single variables span these $r$ dimensions. The constant-prefix
row is $X_{i_*}B_{i_*}$, independent of them by its degree in
$X_{i_*}$. Hence this cut has rank at least $r+1$, and
\[
 r\le2.
 \tag{P18.4}
\]
If $U$ had only one essential variable, the same sum relation
would already make $f=0$. Applying the identical argument in
the reversed order gives rank at most two across
$(U,V)\mid W$.

Neither rank can be one. A rank-one extreme block cut would
factor $f$ into a nonzero linear form and a nonzero quadratic
in disjoint variables. Primality of $\mathcal I_2$ would put
one factor in that ideal, contrary to its minimum multilinear
degree. Thus both ranks are exactly two. Choose independent
linear-form vectors $L(U)$ and $N(W)$, each of length two.
The two coefficient-space factorizations give
\[
 f=L(U)^{\mathsf T}M(V)N(W),\qquad
 M(V)=\sum_{j\in V}X_j M_j,\quad M_j\in E^{2\times2}.
 \tag{P18.5}
\]
All middle variables are essential, so $M_j\ne0$ for each.
Pole extraction gives
$\partial_j f=L^{\mathsf T}M_jN\in\mathcal I_1$.
If $M_j$ had rank one, this nonzero quadratic would be a
product of two nonzero linear forms, which a prime ideal
without nonzero linear forms cannot contain. Therefore every
$M_j$ is invertible. Separate translation in $V$ gives
$\sum_jM_j=0$, so there are at least two middle variables.

Cut between two consecutive middle variables, leaving a
nonempty middle prefix and suffix. The part using a prefix
variable has left degree two and right degree one, and its
coefficient rank is two: one invertible $M_j$ in the prefix,
together with the independent $L,N$, already attains that
rank. The part using a suffix variable has left degree one
and right degree two and likewise has rank two. Their degree
blocks have disjoint row and column indices. The total cut
rank is at least four, a contradiction. This proves the
nonvanishing assertion for all $n$ and every characteristic.

The denominator-safe interpolation argument of I-40, part~6,
now applies at total degree $D=3$. Each selector degree is at
most $n$ after multiplication by
$\prod_i(a_i-y_1)(a_i-y_2)$, and the three-amplitude total
degree is at most three. Its grid size is
$(n+1)^2\binom{6}{3}=20(n+1)^2$, with the field size stated
in (P18.2). The grid does not require the witnessing blocks
or variable order.

\paragraph{2. A smallest unrestricted counterexample has at most two shared boundary variables.}
Suppose that a cubic counterexample exists, choose one with
the fewest essential variables, and fix a width-three order.
For any proper prefix cut $A\sqcup B$, let $C_d$ be its
coefficient block of left degree $d$, $0\le d\le3$.
We first note a direct consequence of the established exterior
characterization. Decode coefficients by
\[
 W=\sum_{|S|=3}\frac{c_S}{\Delta_S}e_S,
 \qquad
 W\in\bigwedge^3\ker[\mathbf1;\mathbf a;\mathbf a^{\circ2}].
 \tag{P18.6}
\]
The three contraction equations connect only adjacent left
exterior degrees: an output two-vector of left degree $e$
receives a left contraction from degree $e+1$ or a right
contraction from degree $e$. Therefore a zero degree block
separates the kernel equations on its two sides. In particular,
if $C_1=0$, the degree-zero part alone still satisfies all three
contractions, so the polynomial represented by $C_0$ belongs
to $\mathcal I_2$. If $C_2=0$, the analogous statement holds
for $C_3$. This is the adjacent-block separation mechanism
of HMM Proposition~8.6, printed pp.~44--45, applied here to
the balanced cubic's stated contraction equations.

A nonzero $C_0$ or $C_3$ so separated is obtained by setting
one side of the variables to zero. Its width cannot increase,
and it has fewer essential variables, contradicting the
choice of the counterexample. Consequently $C_0\ne0$ forces
$C_1\ne0$. P-17's three-block lemma then forces $C_2=0$;
minimality forces $C_3=0$. Symmetrically $C_3\ne0$ allows
only $C_2,C_3$. If both extremes vanish, only $C_1,C_2$ remain.
For every prefix, all occupied degree blocks are thus contained
in two adjacent positions.

Equivalently, any two support triples have prefix counts that
differ by at most one. Define
\[
 \begin{aligned}
 \alpha&=\max\{\text{first position of a support triple}\},\\
 \beta&=\min\{\text{second position of a support triple}\},\\
 \gamma&=\max\{\text{second position of a support triple}\},\\
 \delta&=\min\{\text{third position of a support triple}\}.
 \end{aligned}
 \tag{P18.7}
\]
Then $\alpha\le\beta$: otherwise the prefix ending at $\beta$
contains two variables of one support triple and none of a
triple whose first position is $\alpha$. Similarly
$\gamma\le\delta$, since the prefix ending at $\delta$
would otherwise contain three variables of one triple and
at most one of a triple whose second position is $\gamma$.
Every monomial therefore uses one first position in
$[1,\alpha]$, one second position in $[\beta,\gamma]$, and
one third position in $[\delta,n]$.

No variable can lie in a gap between these intervals, since
all variables are essential. The intervals are consecutive,
and the only possible intersections are the singleton
$\{\alpha\}$ if $\alpha=\beta$ and the singleton
$\{\gamma\}$ if $\gamma=\delta$. If neither intersection
occurs, the polynomial is trilinear in three disjoint
consecutive blocks and part~1 excludes it. Thus any smallest
counterexample has one or two shared boundary variables.
This proves the reduction.\hfill$\square$

\paragraph{Scope and the unresolved boundary cases.}
The trilinear promise requires three consecutive blocks in
one witnessing width-three order. It does not assert that an
arbitrary linear change of coordinates preserves such an order
or the coefficient-cut ranks. The bound is a concrete
arbitrary-length restricted sampler, not polynomial-in-growing-width
PIT. P-19 subsequently proves the full cubic two-pair theorem,
using the interval reduction established here.

Setting a shared boundary variable to zero cannot simply be
used to remove the overlap: it need not preserve the balanced
ideal. Imposing a zero at one node lowers the rational numerator
budget after the corresponding diagonal rescaling. The coupled
quadratic partials at the one or two boundary variables are
the remaining gate at this stage of the argument; P-19
subsequently closes both boundary cases. The separate audit
\nolinkurl{cubic_ordered_gate.md} also records a direct
all-characteristic cubic prolongation consequence of the HMM
kernels; it is not needed for part~1. No historical priority
claim is made.



\subsection{P-19. Balanced two-pair evaluation hits every multilinear width-three cubic}
\label{P-19}

\paragraph{Statement.}
Fix pairwise distinct nodes $a_1,\ldots,a_n$ over a field $E$.
Every nonzero homogeneous multilinear cubic having ROABP width
at most three in some variable order is hit by the balanced
two-pair map
\[
 R_i=t+\frac{z_1}{a_i-y_1}+\frac{z_2}{a_i-y_2}.
 \tag{P19.1}
\]
The variable order need not be known. The statement holds in
every characteristic and for every number of variables.
Consequently the denominator-safe finite grid of I-40 gives
a uniform hitting set of at most
\[
 20(n+1)^2
 \tag{P19.2}
\]
points over a working field of size at least
$\max\{2n+1,4\}$, with the extension-input convention stated
there. By multilinear reversal, the corresponding formal
nonvanishing assertion also holds for homogeneous degree
$n-3$.

\paragraph{Inputs.}
We use the prime balanced one-pair ideal $\mathcal I_1$, its
absence of nonzero linear forms, and pole extraction
$\partial_i\mathcal I_2\subseteq\mathcal I_1$ for multilinear
cubics, as specified in P-17--P-18. We also use the proved
minimal-counterexample interval reduction in P-18, part~2,
and the established cubic exterior characterization of
Hu, van Melkebeek, and Morgan (2024), Definition~9.5, p.~55,
and Theorem~9.13, pp.~61--62 (P-B20). No new characterization
of the full vanishing ideal is assumed.

\begin{proof}
Suppose a counterexample exists. Among counterexamples on
subsets of its variables, choose one with the fewest essential
variables, remove unused variables, and relabel a witnessing
width-three order as $1,\ldots,N$. The corresponding nodes
remain distinct. Denote this nonzero cubic by $f$.

P-18, part~2, supplies consecutive position intervals
\[
 A=[1,\alpha],\qquad B=[\beta,\gamma],\qquad
 C=[\delta,N],\qquad
 \alpha\le\beta\le\gamma\le\delta,
 \tag{P19.3}
\]
such that every support triple has its first, second, and third
positions in $A,B,C$, respectively. There are no inessential
positions. The intersections $A\cap B$ and $B\cap C$ have
at most one position each. This use of P-18 retains its
minimal-counterexample hypothesis; no arbitrary coordinate
specialization is asserted to preserve the balanced ideal.

We first show that at most one position of $B$ lies outside
$A\cup C$. Suppose instead that there are two such positions
$j<k$. Each is essential and can occur only as the second
position of a support triple. Therefore
\[
 Q_j=\partial_j f
   =\sum_{i\in A,\,\ell\in C}q^{(j)}_{i\ell}X_iX_\ell,
 \qquad
 Q_k=\partial_k f
   =\sum_{i\in A,\,\ell\in C}q^{(k)}_{i\ell}X_iX_\ell
 \tag{P19.4}
\]
are nonzero bilinear quadratics. Here every $i\in A$ precedes
$j$, and every $\ell\in C$ follows $k$, so the two extreme
variable sets are disjoint.

Both quadratics belong to $\mathcal I_1$ by pole extraction.
Their displayed bilinear coefficient matrices each have rank
at least two. Indeed, rank one would express the corresponding
nonzero quadratic as a product of two nonzero linear forms.
Primality of $\mathcal I_1$ would put one of those forms in
$\mathcal I_1$, contrary to the absence of linear equations.
This argument applies over $E$ and uses no division by the
characteristic.

Take the coefficient cut immediately after position $j$.
In its left-degree-two block, restrict rows to $X_iX_j$ with
$i\in A$ and columns to $X_\ell$ with $\ell\in C$.
The resulting submatrix is the coefficient matrix of $Q_j$,
so this degree block has rank at least two. In its
left-degree-one block, restrict rows to $X_i$ with $i\in A$
and columns to $X_kX_\ell$ with $\ell\in C$. This gives the
coefficient matrix of $Q_k$, so this block also has rank at
least two. Since $f$ is homogeneous, these degree blocks
have disjoint row and column index sets. The full coefficient
cut therefore has rank at least four, contradicting the
witnessing width bound.

We have proved
\[
 |B\setminus(A\cup C)|\le1.
 \tag{P19.5}
\]
Together with the two possible singleton intersections in
(P19.3), this implies $|B|\le3$.

Now put
\[
 \mathscr K=\ker
 \begin{pmatrix}
 1&\cdots&1\\ a_1&\cdots&a_N\\ a_1^2&\cdots&a_N^2
 \end{pmatrix},
 \qquad
 W=\sum_{|S|=3}\frac{c_S}{\Delta_S}e_S
       \in\bigwedge^3\mathscr K,
 \tag{P19.6}
\]
where $c_S$ is the coefficient of $X_S$ in $f$ and
$\Delta_S$ is the consistently ordered three-node Vandermonde
determinant. The cited characterization gives the membership
in (P19.6), and $W\ne0$ because all $\Delta_S$ are nonzero.

Let $\pi:E^N\to E^{[N]\setminus B}$ be coordinate projection.
Its restriction to $\mathscr K$ is injective. A vector in its
kernel would be supported on at most the three coordinates
of $B$, whereas any at most three columns of the displayed
Vandermonde matrix are independent. The exterior cube of an
injective linear map is injective as well. But every support
triple of $f$, and hence of $W$, meets $B$ by (P19.3).
It follows that
\[
 (\bigwedge^3\pi)(W)=0,
 \tag{P19.7}
\]
contradicting this injectivity and $W\ne0$.

Thus no cubic counterexample exists. The interpolation
argument of I-40 at total degree three gives (P19.2), with
the stated selector and amplitude field sizes. Multilinear
reversal preserves both every coefficient-cut rank and the
balanced rational family, as established in P-17, so it gives
the degree-$n-3$ assertion. This completes the proof.
\end{proof}

\paragraph{Scope and remaining higher-degree problem.}
This proves the full cubic statement, including both shared
boundary cases left by P-18. It does not prove the full
multilinear two-pair theorem in arbitrarily many variables,
nor does it identify balanced two-pair evaluation with an
unweighted ordinary centered SV map. For a remaining
multilinear width-three annihilator, the lowest nonzero
homogeneous degree must be at least four, and the highest
must be at most $n-4$. For the first claim use the known
inheritance of width by the lowest homogeneous component
(HMM Proposition~8.3, p.~42), together with the absence of
degrees below three and the theorem just proved; apply the
same argument after multilinear reversal for the second.
The higher-degree and growing-width sampling gates remain
separate. No historical priority claim is made.

% END SOURCE work/purity/findings.tex

% BEGIN SOURCE work/purity/zero_status.tex
\section{What the stipulated RH and GRH inputs actually discharge}

\paragraph{Input convention.}
Write A for the supplied Alp\"oge theorem: RH for $\zeta$ and GRH for all Dirichlet $L$-functions. The relevant exported declarations are \nolinkurl{riemannHypothesis}, \nolinkurl{LFunction_eq_zero_re}, and \nolinkurl{LFunction_eq_zero_re_of_isPrimitive} in \nolinkurl{ALPOGE_Main.lean}; their predicates are in \nolinkurl{ALPOGE_Statements.lean}. We use A as stipulated and do not claim to have rebuilt its imported archive. Write M for Masullo's claimed GRH for Hecke $L$-functions over every number field, cited conditionally as the prompt's \S1.2. The supplied index does not contain a numbered \S1.2 statement or the full proof, so no theorem number is invented for M. Theorem~13.3.1.1 (PDF p.~522 in the index) is a trace-formula claim; Conjecture~14.1.2.1 (PDF p.~529) explicitly concerns automorphic $\mathrm{GL}_n$ GRH and simplicity. Neither is silently substituted for M or for a proof of general automorphic GRH.

\subsection{P-Z01: Montgomery's supported pair correlation}
For an even Schwartz function $f$ with $\operatorname{supp}\widehat f\Subset(-1,1)$, set $W(u)=4/(4+u^2)$ and $N_0(T)=T\log T/(2\pi)$. Montgomery's theorem and Fourier identity (1973, \S1, pp.~181--182, equations~(1)--(4); P-B05), combined with A, give
\[
 \lim_{T\to\infty}\frac1{N_0(T)}
 \sum_{0<\gamma,\gamma'\le T}
 f\left(\frac{\log T}{2\pi}(\gamma-\gamma')\right)W(\gamma-\gamma')
 = f(0)+\int_{\mathbb R}f(u)\left[1-\left(\frac{\sin\pi u}{\pi u}\right)^2\right]du.
\]
The sum is over ordered zero indices with multiplicity. Subtracting the identical-index diagonal gives the sine-kernel two-point density. To check the normalization directly, the right side equals
$\widehat f(0)+\int_{-1}^1|u|\widehat f(u)\,du$ under the support restriction, exactly the contribution of Montgomery's delta term at frequency zero and its $|u|$ term. This is now unconditional under the input convention. The unrestricted-support assertion, the complete nearest-neighbour spacing law, and all-height GUE statistics require additional estimates. The missing analytic input is control of the off-diagonal prime correlations occurring beyond the supported range; critical-line location does not provide that estimate in the cited theorem.

\subsection{P-Z02: the Hughes--Rudnick hypothesis correction}
Put $x_{\chi,j}=\gamma_{\chi,j}\log q/(2\pi)$ and $W_f(\chi)=\sum_jf(x_{\chi,j})$. Hughes--Rudnick (2003), Introduction, expressly permits complex $\gamma_{\chi,j}$ in $\rho=1/2+i\gamma_{\chi,j}$ (P-B06). Its analytic mean, variance, and restricted moments are therefore already unconditional before A. For admissible $f$ (Definition~2.1), averaged over the $q-2$ nonprincipal characters modulo a prime $q$, Theorems~3.1, 4.1, and 5.1 respectively give
\[
 \mathbb E W_f\longrightarrow\int f,\qquad
 \operatorname{Var}(W_f)\longrightarrow\int_{-1}^1|u|\widehat f(u)^2\,du,
\]
with supports $[-2,2]$ and $[-1,1]$, and Gaussian centred moments through order $m$ when $\operatorname{supp}\widehat f\subseteq[-\alpha,\alpha]$ and $m\alpha<2$. A makes these ordinates real. It also discharges the explicit GRH in Theorem~8.1, Corollary~8.2 ($\liminf_q\min_{\chi\ne\chi_0}x_{\chi,1}\le1/4$), and Theorem~8.3 (its positive-proportion bound for $\beta\ge0.633$). Here $x_{\chi,1}$ is the first nonnegative scaled ordinate. A finite collection of Gaussian moments is not a Gaussian limit theorem at microscopic scale: the same paper calls the phenomenon mock-Gaussian, and the full random-matrix moments are not all Gaussian.

\subsection{Dependency table}
\begingroup\small
\begin{longtable}{p{0.25\textwidth}p{0.33\textwidth}p{0.34\textwidth}}
\hline
Statement and locator & Exact RH/GRH input or other gate & Status under A and M\\\hline
\endfirsthead
\hline
Statement and locator & Exact RH/GRH input or other gate & Status under A and M\\\hline
\endhead
Montgomery (1973), \S1 theorem and (3)--(4), supported pair correlation & RH for $\zeta$; unchanged $(-1,1)$ Fourier support & Discharged by A (P-Z01).\\
Rudnick--Sarnak (1996), Theorem~1.2, $m=1$ & RH for the fixed primitive Dirichlet $L$-function (or $\zeta$); support $\sum|\xi_j|<2$ & Discharged by A. Theorem~1.1 already proves a smoothed analytic version with complex ordinates.\\
Hughes--Rudnick (2003), Theorems~3.1,4.1,5.1 & No GRH for the analytic statistic; A makes all ordinates real & Already unconditional formulas; real-zero interpretation now supplied by A (P-Z02).\\
Hughes--Rudnick (2003), \S8 & Dirichlet GRH & Discharged by A; original support and constants retained.\\
Gao--Zhao (2020), Theorems~1.1 and~1.3 & GRH for the quadratic/quartic Hecke families over $\mathbb Q(i)$ and auxiliary Hecke characters in Lemma~2.8 & Conditional on M; symplectic density for quadratic family with support $(-2,2)$, unitary for quartic with $(-20/19,20/19)$. A alone does not cover these general Hecke characters.\\
Iwaniec--Luo--Sarnak (2000), Theorems~1.1--1.3 & Standing RH for target cuspidal $\mathrm{GL}_2$ functions, plus Dirichlet RH for extended support; see p.~56, and pp.~88--89 for auxiliary symmetric-square estimates & Beyond A and M for the general family: missing RH for $L(s,f)$. Extra symmetric-square/Rankin--Selberg hypotheses used in individual estimates can be avoided by the averaged replacements explicitly noted in \S4.\\
Iwaniec--Luo--Sarnak (2000), Theorems~1.4--1.5 & RH for target $L(s,\operatorname{sym}^2f)$ ($\mathrm{GL}_3$); auxiliary hypotheses as specified in \S4 & Beyond A and M; missing target symmetric-square RH.\\
Iwaniec--Luo--Sarnak (2000), Theorem~1.9 & Weighted, parity-unsplit analytic one-level statistic, support $(-2,2)$ & Already unconditional; target zero ordinates may be complex.\\
Iwaniec--Luo--Sarnak (2000), Theorem~1.10 & Hypothesis S, a prime exponential-sum estimate & Beyond A and M; the precise missing input is Hypothesis S, not merely another GRH label.\\
Katz--Sarnak, Theorem~9.2.6; Weil II, Theorem~3.5.3 & Purity, the actual geometric/arithmetic monodromy hypotheses, and the stated field-extension limit & Already unconditional in that scope. Growing rank/genus/conductor requires separate uniformity.\\
Sawin (2020), Hypothesis~1.3, Theorem~1.4 (label \nolinkurl{main}), Remark~1.9(1),(5) & $H(n,r_+,r_-,n-1)$ follows from duality and gives fixed-parameter $O(q^{-1/2})$; stronger degree-aspect conclusions require vanishing and stronger uniform Betti bounds & Basic $q$-aspect comparison already unconditional. General stronger cohomology and uniformity gates remain; the unconditional unbalanced case is recorded below.\\
Keating--Snaith / CFKRS general number-field moment predictions & Arithmetic moment asymptotics and off-diagonal terms; CFKRS labels its zeta assertion a conjecture & A and M do not discharge the moment recipe. The finite-$N$ CUE integrals are separate exact theorems.\\
Berry--Keating $H=xp$ and BGS universality heuristics & A self-adjoint spectral realization with the required boundary/domain conditions and appropriate chaotic universality theorem & Beyond A and M; RH alone is not that construction.\\\hline
\end{longtable}
\endgroup

Sawin also proves an unconditional unbalanced case: Corollary~1.8 (repeated at source label \nolinkurl{combined-result}) treats $r_-=1$ and $r_+=r$, with $n>2r+1$, the same stated characteristic exclusions, and error
\[
 O\!\left(\sqrt q\,\bigl(q^{-(p-2r)/(2pr)}(r+2)^{r+1}\bigr)^n n^{r+1}\right),
 \qquad p=\operatorname{char}\mathbb F_q.
\]
For fixed $r,p,q$ satisfying $q^{(p-2r)/(2pr)}>(r+2)^{r+1}$, this error decays exponentially in $n$ up to the displayed polynomial factor. Thus some degree-aspect cases are already proved; the unresolved general moment claim is broader than this restricted theorem.

For Rudnick--Sarnak in degree $m>1$, Theorem~1.2 also needs the target RH and, for $m>3$ in the stated source, their coefficient Hypothesis~(1.7). M covers a Hecke case only after all the relevant identification and other hypotheses have been verified; it is not a blanket automorphic GRH theorem. For ILS, ``Hecke cusp form'' refers to a $\mathrm{GL}_2$ automorphic object; it must not be confused with a one-dimensional Hecke character over a number field.

\paragraph{Quantum-chaos source discipline.}
The connection recalled from memory is the symmetry-class heuristic: chaotic quantum systems without an antiunitary time-reversal symmetry suggest unitary statistics, while spinless time-reversal invariance with $T^2=+1$ suggests orthogonal statistics (the $T^2=-1$ symmetry class is symplectic). The original BGS paper (Bohigas, Giannoni, \& Schmit, 1984, pp.~1--4) concerns chaotic quantum spectra; it does not construct the zeta-zero operator. The bibliographic locator and Berry--Keating's explicit $H_{\rm cl}=XP$ speculation were checked against the primary publication records. Berry and Keating (1999), \emph{The Riemann zeros and eigenvalue asymptotics}, SIAM Review 41, 236--266, and their \emph{$H=xp$ and the Riemann zeros} contribution in \emph{Supersymmetry and Trace Formulae: Chaos and Disorder}, pp.~355--367, are the relevant sources. Detailed internal section/equation locators for the latter were not inspected here, so they are not fabricated.

There is a distinction from bare operator existence. Once real zero ordinates $\gamma_j$ are given, the diagonal operator $He_j=\gamma_je_j$ on $\ell^2$, with domain $\{(a_j):\sum_j\gamma_j^2|a_j|^2<\infty\}$, is self-adjoint. Indeed testing the adjoint against the coordinate vectors forces the same diagonal action and domain. This merely packages the supplied spectrum. It does not supply a specified differential or pseudodifferential Hamiltonian, arithmetic trace formula, boundary conditions, or a proof of chaotic local statistics. Those are the substantive missing constructions in the Berry--Keating/Hilbert--P\'olya direction.

% END SOURCE work/purity/zero_status.tex

\clearpage\part{The Grothendieck restatement}
\section{A separate Grothendieck dictionary: objects, operations, and proof transport}
\label{sec:grothendieck-dictionary}

This part restates the checked deductions in the requested geometric language. A locator identifies the framework in which a proof lives; it is not a claim that Grothendieck stated a modern complexity result. The original manuscript images are not in the attachment. Cote and scan references below therefore refer to the supplied transcriptions, and the mathematical corrections do not certify an unseen handwritten reading.

\subsection{Exact framework locators and scope}

\paragraph{D-S3: group schemes and weights.}
For an affine group scheme, its representations are comodules over its Hopf algebra. The relevant reference is SGA~3, Expos\'e~I, Proposition~4.7.2. For a diagonalizable group $D(M)$, such comodules are $M$-graded modules, by Proposition~4.7.3; the following remark gives the corresponding graded-algebra description of affine schemes with action. Proposition~4.7.4 gives equivariant splittings when the underlying sequence splits. These statements, inspected in the original Expos\'e~I PDF, give the integral meaning of every torus weight and cocharacter below. For stabilizers, SGA~3, Expos\'e~VI$_B$, Proposition~4.1 gives upper semicontinuity of group-fiber dimension and explicitly cites EGA~IV$_3$, Theorem~13.1.3. The latter locator is corroborated through that primary SGA text; its entire EGA proof has not been reread here.

\paragraph{D-S6: $K_0$, exterior powers, Adams operations, and filtration.}
SGA~6, Expos\'e~V, Definition~2.1 and Example~2.2(b), construct the pre-$\lambda$ structure of locally projective modules over a ringed topos; Lemma~2.2.1 uses a filtration on exterior powers of an exact sequence. Definition~2.4 and equation~(2.4.2) specify the universal product and composition identities. Expos\'e~VI, section~3.2 gives the special-$\lambda$-ring assertion for the $K'$ of a quasi-compact scheme. Expos\'e~V, Proposition~3.2 defines $\gamma$ from $\lambda$; section~3.10, equation~(3.10.1), defines the augmentation filtration. Section~7.1, equation~(7.1.1), and Proposition~7.5 give Adams operations and their ring identities. The electronic transcriptions of the cited Expos\'e~V passages were inspected; the exact VI~3.2 statement was corroborated in the primary-text transcription during independent review. We do not describe an integral character identity as an integral semisimple decomposition.

For split reductive group schemes over a principal ideal domain, in particular over $\mathbb Z$, a useful complementary reference is Serre (1968), Theorem~5, section~3.7, and the explicit $\mathrm{GL}_n$ character ring in section~3.8. The base-change statements in that theorem concern the representation Grothendieck group, not preservation of a characteristic-zero basis of simple objects in every characteristic. Nor do they automatically answer the separate finite-symmetric-group lattice question recorded in the supplied study of cote~133.


\paragraph{D-Tor: derived restriction and its obstruction groups.}
SGA~6, Expos\'e~III, Definition~1.5 defines Tor independence of two
schemes over a base by vanishing of their positive Tor sheaves; in
the affine case it is exactly the assertion that derived tensor
product agrees with ordinary tensor product. Notation~1.6 specifies
the relative derived tensor product. Expos\'e~I, Proposition~5.1
and Definition~5.2 give the flat-amplitude criterion, and paragraph~5.4
identifies it with ordinary finite flat-resolution length for a
module. These exact passages were read in the electronic French
primary-text transcription. III~1.5 refers to EGA~III, section~6.5,
for local hyper-Tor; that reference is corroborated here through
SGA~6, without a new whole-section reading of EGA~III.

\paragraph{D-Proper: closed images of projective incidence schemes.}
EGA~II, Definition~5.4.1, printed p.~100, defines properness
using separated finite type and universal closedness.
Theorem~5.5.3(i), printed p.~104, states that a projective
morphism is proper. The corresponding French primary-text
passages, including the displayed statements and the proof
following the latter, were read through the indexed original
PDF text. The complete PDF did not load in the web reader.
These locators justify the closed incidence images below;
they do not supply a computational bound for their equations.

\paragraph{D-FGA: images, quotients, and projective parameter spaces.}
For a homogeneous orbit map over $\mathbb Z$, take its schematic image by the kernel of the coordinate homomorphism, as in G-04. Degreewise kernels and images are given by finite matrices and their minors. EGA~II, sections~2.4--2.5, in particular Proposition~2.5.2, gives the projective scheme and associated sheaf of a graded module. The electronic EGA translation was inspected for this locator. FGA, Bourbaki Expos\'e~221, Theorem~3.1 gives the Quot parameter space, Theorem~3.2 its fixed-Hilbert-polynomial projective form, and Proposition~3.8 its Grassmannian realization in the required degree range. The original FGA scan was inspected. These results place the families in parameter spaces; they supply no polynomial bound for finding the relevant minors or their nonvanishing loci.

\paragraph{D-ET: traces, weights, and their arithmetic size.}
SGA~4$\frac12$, \emph{Rapport sur la formule des traces}, Theorems~3.1--3.2 and Proposition~3.3, gives the determinant expression for the $L$-function, the trace formula, and the logarithmic-derivative identity. SGA~5, Expos\'e~X, Theorem~7.1 and Corollary~7.12, formula~(7.13), gives the Euler--Poincar\'e conductor accounting. Deligne's Weil~II, Theorem~3.3.1 and section~3.3.10, supplies the compact-support weight estimate used in P-06; Theorem~3.5.3 gives the equidistribution setting. The Weil~II locators were checked against Katz and Sarnak (1999), sections~9.0.13--9.0.16 and Theorem~9.2.6. The original large Weil~II PDF was not successfully opened here. This is a statement-and-dependency check, not an independent verification of the complete proof of Weil~II.

\subsection{Rectangular characters and orbit-closure multiplicities}

\paragraph{R-01. Tensor pullback of a Schur operation.}
For locally free $E,F$ of ranks $a,n$, the tensor morphism $\mathrm{GL}(E)\times\mathrm{GL}(F)\to\mathrm{GL}(E\otimes F)$ is defined over $\mathbb Z$. The class $s_{(d^n)}[E\otimes F]$ restricts, after passage to characteristic zero, to the sum of $[S_\mu E\otimes S_\nu F]$ with coefficients $g((d^n),\mu,\nu)$. The Haar projector is the complex realization of extraction of that character component. Its finite CUE dimension is a rank of a Schur construction; its normalized limiting constant need not be an integer rank. This is precisely the proof of R-01 transported through D-S3 and D-S6. Cote~133, scan~028, lines~6--17, provides the tensor-power and symmetric-operation language, while scan~019 distinguishes rank from the full character class.

\paragraph{R-02. Two different meanings of positivity.}
The tableau expansion is an integral character identity in the universal special $\lambda$-ring, with the Weyl alternating factor extracting a highest weight. Over $\mathbb C$, the invariant summand of $S_\lambda(E\otimes F)\otimes\det(E)^{-d}\otimes\det(F)^{-d}$ is an actual vector space. A Hermitian metric identifies its projector trace with a sum of squared absolute values. The former is an integral identity; the latter uses complex conjugation and compact integration. D-S6 and cote~133, scan~028, lines~19--24 (the direct-summand question), are the exact conceptual distinction. No cancellation-free combinatorial basis is inferred from the analytic squares.

\paragraph{R-03. Central quotient, disconnected stabilizer, and rank.}
The tensor morphism has the central subgroup $t\mapsto(tI,t^{-1}I)$ in its kernel, a statement of group schemes over $\mathbb Z$. A compact block-Levi branching rule is therefore a different restriction functor. The transpose component of the determinant stabilizer is additional group-scheme data beyond the connected tensor subgroup. The $S_4$ counterexamples compare distinct invariant functors and survive as finite character computations. D-S3 places these groups and actions; cote~133, scan~019, lines~45--56, is the corresponding warning that augmentation loses equivariant information. The negative Weingarten coefficient is a complex integration fact, not a negative rank of a locally free sheaf.

\paragraph{R-04. Plethysm as an operation on representation $K_0$.}
The cycle-index formula expresses $\sigma^d(\sigma^n[E])$ using the Adams operations $\psi^j$. Its displayed rational coefficients are interpreted in $K_0\otimes\mathbb Q$ or the universal torsion-free character ring; the resulting symmetric-power operation is integral. D-S6, V~7.1 and V~7.5, makes their algebraic status precise. They are ring operations; $g\mapsto g^j$ need not be a group homomorphism, and $\psi^j$ need not preserve the cone of honest representations. The $U(2)$ mixed-moment counterexample therefore translates to the statement that separate character distributions do not determine the joint operation. The wreath-subgroup invariants give the characteristic-zero tensor-functor realization. Cote~133, scan~028, lines~12--17 and~26--35, places tensor powers and induction products, with each $\rho_i$ in its own homogeneous degree.

\paragraph{R-05. Image modules, not merely stabilizer invariants.}
The intrinsic object is the image of the degree-$d$ orbit-coefficient map. Its dual is the span of the orbit of $f^{\otimes d}$. Over $\mathbb C$, complete reducibility identifies the multiplicity in each Schur summand with the rank of its multiplicity matrix. The Haar covariance and its reduced Gram matrices have exactly those images; their principal exterior coefficients are squared Pl\"ucker norms. This transports the entire rank proof to determinantal rank strata and Grassmannians in D-FGA. The Hermitian form is an auxiliary real structure and is not asserted to define a motive. Cote~133, scan~028's question about a direct summand is answered here only after passing to characteristic zero.

\paragraph{R-06. Integral matrices and special fibers.}
The polynomial orbit map is defined over $\mathbb Z$ after the specified symmetrizer scalings. Its finite coefficient matrix is an integral map of free modules. The Gaussian covariance is another complex Gram form on its image, while the integral character-sum matrix computes the required Schur-isotypic rank after scalar extension. The Schur projectors require the stated factorial denominators. G-04 specifies the open localization on which the ranks and kernels specialize correctly. D-FGA records the associated rank and Quot strata; D-S6 records the tensor and Schur operations. This translation preserves the exponential matrix sizes, so changing language does not remove the computational gate.

\paragraph{R-07. Binary restriction and highest weights.}
The determinant-pencil substitution is a morphism of coefficient schemes, defined after clearing the displayed binomial denominators. The sections $Q_j$ are highest-weight polynomial functions, and independence is proved by their restrictions to the chosen diagonal-pencil chart. Over $\mathbb C$ every binary form is a product of linear forms, so the binary coefficient subscheme lies in the determinant orbit closure. Highest-weight sections of length at most two therefore cannot belong to its ideal. This is the same proof expressed as ideal restriction in D-FGA and weights in D-S3. Cote~133, scan~028 supplies the polynomial-functor viewpoint; it does not supply a new multiplicity formula by rank alone.

\paragraph{R-08. The quotient class in equivariant $K_0$.}
For the polynomial ring $S=\operatorname{Sym}(W)$ and orbit quotient $A=S/I$, an equivariant free resolution yields $[A]=\sum_i(-1)^i[S\otimes T_i]$ in the appropriate graded Grothendieck group. Symmetric powers produce the Molien denominator and the alternating resolution produces the numerator. Exactness, not pointwise sign, proves the formula. D-S6, V~2.2(b), and D-FGA place the additivity and coherent quotient. Cote~133, scan~019, lines~9--20, contains the exterior/symmetric generating-series relation. Ignoring the numerator forgets the defining ideal and is not an operation allowed by $K_0$ additivity.

\paragraph{R-09. A four-factor invariant and a finite-rank cutoff.}
The squared rectangular multiplicity sum is the determinant-isotypic component of $\operatorname{Sym}(E_1\otimes E_2\otimes E_3\otimes E_4)$ in degree $nd$, obtained by composing two tensor restriction functors. The $U(N)$ cutoff means that Schur functors with more than $N$ rows vanish on a rank-$N$ bundle. The stable ratios denominator is therefore an identity of symmetric operations on four tensor factors. D-S3 and D-S6 carry the proof without an arithmetic purity hypothesis. Cote~133, scan~028, is the corresponding tensor-power locus. Nonnegative character norms are exact, but no effective algorithm for the remaining coefficient extraction has been added by the translation.

\paragraph{R-10. The quadratic case and the component group.}
Over $\mathbb C$, the $2\times2$ determinant is a nondegenerate quadratic form in four variables. Its orbit closure is the whole quadratic-form coefficient space. The connected group is the tensor image of $\mathrm{SL}_2\times\mathrm{SL}_2$; the full quadratic stabilizer includes the other orthogonal component. Accordingly the all-odd four-row part present in the ordinary rectangular Kronecker coefficient is removed from the symmetric coefficient and the coordinate ring. D-S3 records the group and component distinction, while D-S6 identifies the even-part Schur decomposition. Cote~133, scan~019's augmentation does not retain this distinction.

\paragraph{R-11. Low-degree surfaces and restricted ideals.}
The linear and quadratic determinant constructions are coefficient-scheme morphisms. The cubic statement uses a dense open set of smooth cubic surfaces and Beauville's determinantal representation theorem; closedness extends the containment. Block padding is a morphism over the coefficient ring, and the highest-weight ideal inclusion is the same D-S3/D-FGA argument as G-08. It excludes four-row multiplicity obstructions for padded forms of degree at most three without asserting a global orbit containment for arbitrary many-variable cubics. Cote~133, scan~028 provides the Schur-functor interpretation of the resulting coordinate modules.

\paragraph{R-12. A Jacobian certificate and its characteristic-zero consequence.}
The permanent-pencil coefficient map is a morphism of affine schemes over $\mathbb Z$. A maximal Jacobian minor nonzero modulo the stated prime is a nonzero integer and therefore nonzero over $\mathbb Q$. In characteristic zero the differential-rank criterion gives dominance, so the four-variable projection of the permanent orbit closure fills the coefficient space in the certified sizes. D-FGA then translates the hypersurface quotient sequence of G-07 into the exact four-row multiplicity difference. The finite-field calculation is a certificate for an integer minor; it is not an appeal to equidistribution or an assertion of semisimplicity modulo that prime. The tensor-power language is cote~133, scan~028; the rank/character distinction remains scan~019.


\paragraph{R-13. A finite family of integral rank certificates and a torus quotient.}
For each stated size the Jacobian is a map of free modules on the affine pencil-parameter scheme. Its displayed minor defines a principal open rank locus over $\mathbb Z[1/101]$ and over $\mathbb Q$. The independent permanent identities check that these are the derivatives of the same morphism, rather than arbitrary matrices. For larger sizes the row-column torus acts in the fibers with the computed one-dimensional stabilizer. D-S3, VI$_B$~4.1 and the cited EGA dimension result, together with the elementary fiber-dimension calculation, give the upper bound. Highest-weight ideal restriction in D-FGA then yields the stated finite-range multiplicities. This extends the rank and Schur-functor dictionary of cote~133, scan~028; it does not turn a modular certificate into a characteristic-free semisimple decomposition.

\paragraph{R-14. The projective multiplication image and Hilbert growth.}
The map sending a line and a degree-$m$ form to their padded product is a morphism from $\mathbb P^3\times\mathbb P(\mathcal P_{m,4})$. Properness makes its image closed, and the unique linear factor of a general product makes the morphism generically injective. The resulting affine cone has dimension $c_m+3$. Its homogeneous ideal and Hilbert function are the D-FGA objects of EGA~II~2.4--2.5 and FGA Expos\'e~221, Theorems~3.1--3.2. The dimension comparison forces a larger Hilbert function in large degree, and D-S3 identifies the corresponding highest-weight multiplicities. The same proof therefore transports exactly; no effective degree threshold or new purity statement is inserted by this reformulation.

\paragraph{R-15. A normal direction to a family of pencil images.}
The cyclic pencil is a point of the affine parameter scheme, and $W_0$ is the image of its tangent map. Restriction to the explicitly defined lines $\Lambda_i\subset\mathbb P^3$ and the independent interval sections computes a lower bound for that image. The residue calculation proves that the section $T$ gives a nonzero class in the quotient by $W_0$. The two-chord family makes this class lie in the tangent image away from the special parameter, while an augmented determinant witnesses persistence of independence on a principal open. These are explicit maps of coefficient modules and open rank loci in D-FGA, defined over an integer coefficient ring after choosing the stated integral parameters. The generic characteristic-zero dimension and Hilbert argument then gives the four-row inequality for every equal size. D-S3 and D-S6 identify its weight and Schur class. Cote~133, scan~028 supplies the tensor-module setting; scan~019 explains why this is a multiplicity assertion rather than a numerical CUE dimension alone.

\subsection{Isolation as a torus action and a problem of coefficient modules}

\paragraph{I-01. Independent variable coordinates.}
The vertex-variable path polynomial is a section of the structure sheaf of an affine toric coordinate space. Its path monomials are distinct characters of its split torus. A Hoza weight is a cocharacter, and uniqueness of the minimum produces a one-dimensional initial character with a nonzero coefficient. D-S3 identifies the action and grading; I-10 supplies its flat hypersurface family over a field. Cote~133, scan~019, is only a finite-set analogy here, not a claim that path existence already computes a representation multiplicity.

\paragraph{I-02. Descent through a monomial quotient.}
The monomial map is a map of affine toric schemes whose character-lattice map is the integer matrix $A$. A cocharacter is induced by an integral weight on the full variable torus exactly when its degree vector lies in the image of $A^{\mathsf T}$. Homogeneity of the binomial kernel gives the grading on the image algebra; extending it to the full variable torus additionally requires the integral image-lattice condition in I-02. For example $z\mapsto x^2$ has zero kernel, but weight one on $z$ is not induced by an integer weight on $x$. The diamond is an explicit failure even of kernel homogeneity. This is D-S3, I~4.7.3, applied to the quotient algebra. Cote~133, scan~019 supplies no rule permitting independent path gradings to be identified with the quotient grading.

\paragraph{I-03. A complete coefficient flag.}
Filter the coefficient module by the span of monomials below each cocharacter weight. Basis isolation says that at each growth degree precisely one indexed coefficient survives, and that these surviving degrees are distinct. The dual formulation quantifies over every nonzero quotient functional. This is a statement about filtered locally free modules over a field and their graded pieces, not just the presence of one initial path. D-S3 gives the associated grading, and D-FGA places the open minor conditions on the relevant flag and Grassmann schemes. Cote~133, scan~028's direct-summand question motivates the distinction between a filtration and a chosen splitting.

\paragraph{I-04. The Pl\"ucker weight polytope.}
Take the coefficient row plane in a Grassmannian and its nonzero Pl\"ucker coordinates. Their torus weights are the sums of the corresponding monomial exponents. The vertex-fiber assertion is the symmetric-exchange proof stated in I-04; the probabilistic estimate counts walls crossed in one coordinate. The resulting open cones are normal cones refined by basis-weight inequalities. D-S3 identifies their cocharacters and D-FGA, Expos\'e~221, Proposition~3.8 places the Grassmannian. No finite-field purity theorem chooses an integer point in all these cones automatically.

\paragraph{I-05. Real Gram positivity and algebraic minors.}
With two independent dual coefficient matrices, Cauchy--Binet is an integral identity of sections of a determinant line. On the real locus where one is the conjugate transpose of the other, it becomes a sum of nonnegative squared Pl\"ucker coordinates. The counterexample shows that this determinant line does not retain the whole coefficient-image flag after colliding weights. D-FGA supplies the rank-stratum language, and D-S6 supplies the determinant as top exterior power. Cote~133, scans~019 and~028, respectively mark the distinction between numerical positivity and actual tensor-module information.

\paragraph{I-06. Secant rank of an evaluation functional.}
An evaluation at a point is a rank-one decomposable functional on the bounded-degree polynomial space. The moment matrix of the proposed nonlinear exponent functional has tensor-product rank. Expressing it as a fixed sum of evaluations would place it in a secant locus of bounded rank; the computed determinant excludes that locus below the stated exponential size. This is an algebraic rank statement over $K(t)$, hence belongs to the determinantal geometry in D-FGA. D-S3 also explains why a nonlinear exponent function is not a character-lattice homomorphism. No motive is needed for this counterexample.

\paragraph{I-07. Taking invariants changes the inserted data.}
The invariant functor on complex representations is realized by the Reynolds projector. An independent twirl sends each state insertion to its invariant scalar component. The two-insertion calculation resolves the commutant into the identity and transposition operations and exhibits a negative coefficient. Thus invariants are extracted in the representation category of D-S3/D-S6, with the full tensor operation retained. Cote~133, scan~028, describes why the symmetric-group action on tensor powers matters; scan~019 warns against replacing that action by dimension.

\paragraph{I-08. A finite quotient morphism, conditional on a hitting set.}
The trace-generating functions are regular invariant functions on the affine scheme of matrix tuples. A hitting set separating their coefficient polynomials gives a separating family in the quotient coordinate algebra; the cited invariant-theory theorem makes that algebra finite over the generated subalgebra. This is a finite morphism of affine schemes, not automatically a choice of algebraically independent parameters of minimal cardinality. D-FGA gives the general morphism and parameter-space setting. The unresolved explicitness condition is exactly the one retained in I-08 and GCT~V; the dictionary does not discharge it.

\paragraph{I-09. A tensor contraction computing a Hermitian norm.}
The matrix product computes a regular coefficient map. Tensoring with its conjugate and contracting equal monomial indices gives its squared norm; the torus integral extracts the zero character in that tensor product. D-S3 explains the character extraction. Over an integral algebraic base one must use independent dual variables, not treat complex conjugation as a regular morphism. The metric realization lives over $\mathbb R$ or $\mathbb C$. It is the same coefficient module as before, and access to its presenting matrices remains a white-box assumption.

\paragraph{I-10. Rees modules and specialization.}
The cocharacter gives an increasing filtration on coefficient spans. Its Rees module has the associated graded as its special fiber and the original module on the punctured line. The explicit hypersurface homogenization in I-10 is flat over a field by the torsion-free PID argument in its proof; general ideal degenerations require the entire weight-homogenized ideal. D-S3 identifies the grading and D-FGA identifies the flat family and its Hilbert/Quot points. The associated graded preserves total filtered rank, whereas adding colliding coefficients in a substituted polynomial can reduce their span, as I-05 proves. These are different maps.

\paragraph{I-11. Exterior powers do not remove equations.}
The exterior dimension and graph-of-linear-map count describe the sizes of coordinate and Grassmannian models. The two-term zero determinant is already an equation of the rank-deficient locus. Assigning different path weights to its cancelling terms changes that equation. D-S6 gives exterior powers and D-FGA gives the Grassmannian, but neither turns a signed determinant expansion into unsigned reachability. Cote~133, scan~028's direct-summand margin remains a relevant distinction between a functorial module and its termwise expansion.

\paragraph{I-12. Support is the nonvanishing locus of a regular function.}
A specified coefficient of an integer ROABP is a regular polynomial in its edge parameters. Its nonvanishing open subset is determined by a signed path sum, whereas ordinary reachability detects whether a parameterizing path set is empty. The GapL reduction is a computation of that regular function; the conditional complexity-class inclusion follows only under the specified logarithmic-space reduction. D-FGA locates its zero scheme. Cote~133, scan~019 supplies the nearby conceptual distinction between the cardinality of a finite set and a signed class; it states no complexity theorem.

\paragraph{I-13. Evaluation quadrature and the Segre variety.}
Coefficient tensors of width-one products form the affine cone over a Segre variety. Equality of the quadratic forms on all these tensors, with their complex conjugates, forces the moment-matrix identity proved in I-13. Its rank then bounds the number of evaluation vectors. This restates the entire proof as an assertion about a Segre cone, its dual evaluation vectors, and determinantal rank loci in D-FGA. The Hermitian positivity is not an integral $\lambda$-ring assertion. The much smaller hitting set in the same finding concerns only kernel avoidance, not equality of quadratic forms.


\paragraph{I-14. A Grassmannian open condition for several substitutions.}
The row plane of the coefficient matrix is a point of $\operatorname{Gr}(r,F^s)$. A fixed family of cocharacters gives the integral linear map $P_{\mathcal W}^{\mathsf T}$. The desired locus is the open set where its restriction to the tautological rank-$r$ bundle is injective. Its complementary equations are the displayed linear forms in Pl\"ucker coordinates. D-FGA, Expos\'e~221, Proposition~3.8 is the Grassmannian framework, and D-S3 identifies the cocharacters. The example proves that this open set depends on coordinate values and can be larger than the union of the individual basis-isolation loci.

\paragraph{I-15. A finite character-lattice obstruction.}
The bounded seed decoder produces a subspace of cocharacters of the split torus. Pigeonhole counting gives two distinct characters whose difference annihilates every cocharacter in that subspace. Their binomial is therefore a nonzero regular function in the kernel of every pullback. This is precisely the character/cocharacter pairing of D-S3, I~4.7.3, with the finite integer bounds retained. The extension to bounded-degree Boolean decoders is a dimension bound for their cocharacter span. It does not forbid more general maps whose spans are larger.

\paragraph{I-16. A finite coefficient-state scheme and base-field points.}
Over $\mathbb F_q$ the coefficient row space has $q^w$ rational points, and each letter acts on it by a linear morphism. The state automaton records these morphisms exactly. The separate evaluation lower bound counts the rational points avoiding products of prescribed coordinate divisors. Both constructions live in affine schemes over the finite field; D-S3 gives the coefficient representation, while D-FGA distinguishes the finite set of rational points from the whole scheme or its extension-field points. Purity does not change the explicit counting inequality.

\paragraph{I-17. A trivialization on an invertibility open set.}
The nonzero determinant polynomials define an open set on which the constant transition maps are bundle isomorphisms. The displayed gauge change trivializes those constant maps. The nested coefficient row subspaces can then grow at most $w-1$ times, and their chosen sections have the stated support. This is a concrete filtered-module argument in D-S3 and a rank-stratum argument in D-FGA. The known AGKS invertible-factor result is cited separately; the proof establishes the particular permutation-invariant sampler stated here.

\paragraph{I-18. A unit minor over a discrete valuation ring.}
Over $F[[\varepsilon]]$, a maximal minor defines an open rank locus precisely when it is a unit. A nonzero nonunit meets the generic fiber and misses the special fiber. The explicit matrix family separates these two cases, and the ideal of sampled scalar values shows why saturation in the fraction field does not recover an original evaluation. This is the elementary local-module form of the determinantal loci in D-FGA. The distinction from a flat Rees degeneration is essential: flatness of an ambient family does not imply surjectivity of this particular compressed module map.

\paragraph{I-19. Rank-factorization charts and the coefficient field.}
On the principal open where a chosen minor is invertible, a rank-$r$ bundle map factors through a rank-$r$ bundle. The adjugate formula clears its denominator and is an integral polynomial identity. Passing to the function field of the cut-variable scheme makes the intervening blocks disjoint; it does not remove the cut variables from the original base. Multilinearity proves coefficient-span preservation under independent local substitutions. D-FGA supplies these rank charts and D-S3 the tensor contractions. The unresolved descent from the cut-variable function field is explicitly retained.

\paragraph{I-20. Bounded losses in a sequence of coefficient subspaces.}
The product of a nonzero coefficient and the selected pivot minors defines a principal open on the cut-coordinate affine scheme. The finite grid meets this open by the degree bound. After specializing there, the coefficient subspaces have total dimension loss at most $\kappa$; hence their total growth is at most $w-1+\kappa$. This is an exact rank and filtration argument on the same D-FGA charts as I-19, followed by the D-S3 coordinate grading. It yields the stated sampler with its parameter dependence, without asserting that arbitrary singular sequences have bounded total loss.

\subsection{Purity, arithmetic sampling, and the limits of numerical shadows}

\paragraph{P-01. A Gram form on a finite space of sections.}
The coefficient support spans a vector space of functions. Their pairwise correlations are entries of a Gram matrix; its spectral lower bound is the elementary estimate in P-01. In geometric applications these functions may be sections or Frobenius traces, but no sheaf is manufactured by naming the Gram matrix. D-ET applies only after an actual trace-function realization exists. The distinction from rank augmentation in cote~133, scan~019, is that all cross terms have been retained and bounded.

\paragraph{P-02. A Kummer-character construction.}
The signs arise from the rank-one quadratic Kummer sheaf pulled back by $t\mapsto t+i$ on the line with its zeros removed. Products over distinct indices remain governed by the explicitly squarefree polynomial, so the curve character estimate controls the off-diagonal trace sums. The altered values at the excluded points contribute the stated elementary error. D-ET, together with the exact Alon--Goldreich--H\aa stad--Peralta locator, carries the proof. This yields the stated sparse multilinear hitting set; it does not create a polynomial bound on arbitrary coefficient supports.

\paragraph{P-03. The kernel of restriction.}
The sign-cube counterexample is a nonzero section whose restriction vanishes on the selected finite subscheme; the common-affine repair is annihilated by the explicit Vandermonde section. These are regular-function identities and survive independently of improved point-counting estimates. They belong to the affine-scheme language of D-FGA. Cote~133, scan~019 is an analogy for information lost by passing to a smaller numerical image, not a source of these new counterexamples.

\paragraph{P-04. Pullback to a curve and its zero divisor.}
For the curve morphism specified by the rational functions $r_i$, the section $f(r_1,\ldots,r_n)$ has polar divisor bounded by the displayed $D$. If it is nonzero, its zeros have total degree at most $\deg D$. The trace formula and curve RH then provide enough rational points. D-ET controls the count; the pullback's nonvanishing is a separate scheme-theoretic condition. This proof uses curves and sheaves actually specified in the statement, so a motivic interpretation is available for their cohomology, not for an unspecified ROABP.

\paragraph{P-05. A curve family can lie in a small polynomial kernel.}
Dimension comparison in the spaces $H^0(C,\mathcal O_C(aD))$ produces an algebraic relation between the two chosen functions. Multiplying the relations for the finite curve family gives its annihilating section and the explicit bivariate ROABP. D-FGA places the image and its ideal; D-ET cannot remove a function from that ideal by increasing the field size. The relation is proved directly with its degree and width bounds and is not a lower bound against every geometric sampler.

\paragraph{P-06. Purity with a complete conductor budget.}
The relevant class is $[\mathcal F_i\otimes\mathcal F_j^\vee]$ in the Grothendieck group of constructible Weil sheaves. The compactly supported trace gives its Gram entry. Geometric irreducibility and inequivalence determine $H_c^2$, affine vanishing removes $H_c^0$, and Weil~II bounds $H_c^1$. SGA~5 supplies the displayed Swan-conductor expression for its dimension. This is exactly D-ET, with every rank and conductor retained. A claim of polynomial explicitness additionally requires that these data and the enumeration of the points be polynomially bounded. Cote~133, scan~028 supplies a tensor-category analogy, but not the missing sheaf realization.

\paragraph{P-07. Sparse cocharacters over the integer lattice.}
The weight family is an explicit finite set in $\operatorname{Hom}(\mathbb G_m,\mathbb G_m^n)=\mathbb Z^n$. The collision polynomial over $\mathbb F_p$ is nonzero by its degree and coefficient bounds, so a seed outside the union of its root sets makes the support characters distinct. D-S3 translates the construction exactly; only the elementary finite-field root bound is used. Thus the result is independent of an additional purity or RH hypothesis.

\paragraph{P-08. Arithmetic covers with a specified nonvanishing encoding.}
The common number field determines a finite \emph{\'etale} cover of $\operatorname{Spec}\mathbb Z$ after inverting its discriminant. Frobenius conjugacy classes are the arithmetic data on that cover. Effective Chebotarev produces a small unramified prime only after the stated common map from primes to evaluations and its nonvanishing property are supplied. D-ET provides the parallel finite-field trace perspective; it is not a proof of that arithmetic encoding. The abelian case uses the stipulated Dirichlet GRH, while the general stated Hecke extension remains conditional on Masullo. No claim about arbitrary automorphic $\mathrm{GL}_n$ GRH is made.


\paragraph{P-09. Positive Hermitian forms on a Segre cone.}
The width-one coefficient tensors are decomposable sections, forming a Segre cone. Applying the upper sampled-energy bound to the conjugate evaluation tensor bounds each weighted evaluation norm, while applying the lower bound to the monomial sections bounds their total trace. This is the complete rank-one tensor proof of the cardinality inequality. D-S6 identifies tensor products and D-FGA the Segre coefficient scheme. The inequalities themselves use a complex Hermitian realization and are not inequalities in integral $K_0$.

\paragraph{P-10. Projective incidence and avoidance of a determinantal cone.}
The tensor-train class is the closed cone defined by the stated flattening minors. Its projectivization and evaluation hyperplane sections give a projective incidence correspondence. Proper projection makes the bad tuples a closed set; successive hyperplane sections show this set is proper at the stated cardinality. For unknown order one uses the finite union of the permuted cones. This is D-FGA's projective and determinantal geometry with D-S3's coefficient grading. Compactness of the complex unit sphere gives a positive Hermitian lower bound only after kernel avoidance is known. Neither the projective argument nor the $K_0$ class bounds the integer heights of an explicit good tuple.

\paragraph{P-11. Pullback along an explicit arithmetic parameterization.}
The presentation $\alpha_i=A_i(\theta)/b$ specifies a morphism from an affine parameter scheme to the variable space after inverting $b$. Pulling back a polynomial to $G_f(U,T)$ and using the stated degree grid gives the unconditional forward construction. Conversely Lagrange interpolation and a cyclotomic point give a finite \emph{\'etale} arithmetic realization after inverting the discriminant; minimal-polynomial degree prevents a nonzero rational pullback from vanishing there. D-FGA describes the morphism and its kernel, and D-ET places the finite cover. The proof does not require a Frobenius estimate: the missing hypothesis is the explicit common nonvanishing map itself, rather than a further trace bound.

\paragraph{P-Z01 and P-Z02. Realizations of zero statistics.}
These are analytic consequences of the precisely stated RH predicates and supported explicit-formula theorems. Their expression by unitary integrals does not make the zeta zeros the eigenvalues of a constructed geometric Frobenius. D-ET explains the function-field comparison only in the families with an established sheaf and monodromy group. There is no cote~133 scan containing a classical-RH-to-GCT implication; scans~019 and~028 supply representation operations, not that missing implication.

\subsection{The integral controls G-01--G-12}

\paragraph{G-01. Burnside classes, representation classes, and rank.}
The Burnside-ring counterexample is an equivariant finite-set computation. Linearization sends a $G$-set to a permutation module, but the two-subset operation does not become exterior power: the $C_2$ sign in the proof witnesses this failure. In D-S6 the correct relation is $x^2=2\lambda^2(x)+\psi^2(x)$, and only augmentation replaces $\psi^2(x)$ by the rank. This is the precise mathematical correction to cote~133, scan~019, line~53 and the supplied study's equivariant reading.

\paragraph{G-02. The universal terminal variable.}
Specializing one factor to the tensor unit in the universal product identity forces the top elementary symmetric function to remain present. SGA~6, V~2.4, equation~(2.4.2), explicitly includes that terminal variable. The proof in G-02 is its necessity in the universal character ring, so the correction to scan~006, lines~25--26, is mathematical and independent of paleography.

\paragraph{G-03. Equal finite filtration data, unequal representations.}
The two sums of line characters are honest locally free representations of $\mathbb G_{m,\mathbb Z}$, whose difference is the $(D+1)$st power of an augmentation-zero element. D-S6, V~3.10, places the difference in the required filtration step. The associated-bundle construction transports this identity to any scheme carrying the line bundle; it does not force the bundles to remain distinct on every base. Thus finite Chern or $\gamma$ data cannot recover all the character coefficients needed in scans~019 and~028.

\paragraph{G-04. The integral image and its good rank locus.}
The schematic orbit closure is defined over $\mathbb Z$ by a kernel, and its image algebra is torsion-free. In a fixed degree a nonzero maximal minor gives an explicit principal open on which the coefficient matrix is equivalent to a split rank matrix. This is a direct proof of base change for the displayed module, with D-FGA providing the image, Grassmannian and Quot framework. At the excluded primes the flat integral model's fiber may differ from the orbit closure of the reduced point. This is why integral $K_0$ identities from scan~028 do not imply unchanged characteristic-zero multiplicities in all fibers.

\paragraph{G-05. Plane curves and the three-row sector.}
The nonempty complement of the theta divisor in $\operatorname{Pic}^{g-1}(C)$ supplies Beauville's line bundle, hence a linear determinantal equation. Density of smooth curves gives the scheme-theoretic containment needed for every ternary coefficient section. Highest-weight sections of length at most three belong to that coefficient subalgebra, so its determinant ideal has none. D-FGA records the ideal restriction and parameter spaces; D-S3 records the weights. The numerical positive count is not extracted from Chern classes of the curve alone.

\paragraph{G-06. Stabilizer dimensions and four-row equations.}
The polynomial matrix-pencil map is invariant under the explicitly stated subgroup scheme $H$ of $\mathrm{GL}_n^2$. Its generic stabilizer is the scalar subgroup by the concrete matrix calculation. D-S3, VI$_B$~4.1 and its cited EGA theorem justify the dimension behavior, and the resulting proper image has a nonzero ideal. D-FGA and highest-weight restriction give an equation of exactly four rows, using G-05 to exclude shorter weights. This supplies no nonvanishing on a particular padded form until that extra statement is proved.

\paragraph{G-07. A Cartier divisor and a determinant twist.}
The reduced determinantal-quartic locus is a prime hypersurface. Its equation spans the character $\det^D$, so its ideal is the line-character twist of the ambient graded algebra shifted by $D$. The exact sequence in G-07 is the equivariant $K_0$ statement of a principal Cartier divisor. D-S6 gives additivity, D-FGA gives the projective hypersurface and quotient, and D-S3 identifies the character. It produces the shifted plethysm subtraction exactly; the corrected reduced degree is a literature issue, not a new transcription of cote~133.

\paragraph{G-08. A universal morphism of coefficient schemes.}
The two monomial trees and their cofactor determinant give a matrix over the universal integer coefficient ring. Homogenization and block padding therefore commute with any base change. Over $\mathbb C$ the containment of all low-variable restrictions yields an inclusion of highest-weight ideal spaces, hence the multiplicity inequality. This is D-FGA applied to an explicit scheme morphism, with the weights read through D-S3 and the symmetric powers through D-S6. Its complexity bound is part of the construction, not a consequence of representability alone.

\paragraph{G-09. A local padding coordinate and smaller universal matrices.}
Once the nonzero padding form is one coordinate, dehomogenization reduces the number of affine variables by one. The shared exponent chains are again defined over the universal coefficient ring; their exact vertex count gives the stated finite size. Changing an arbitrary padding form to that coordinate is local on its nonzero coefficient charts. D-FGA carries these local morphisms and their closure statement, and the D-S3 highest-weight argument yields the improved row exclusion. The known asymptotic uniform-representation theorem is cited separately; neither the chart construction nor its GCT consequence requires RH.


\subsection{Further rectangular, rank-profile, and arithmetic consequences}

\paragraph{R-16. Several normal directions on one deformation.}
The cyclic pencil lies on the coefficient affine scheme; disjoint row-pair modifications define one morphism from the affine line into that parameter scheme. The exact permanent identities produce polynomial sections of its differential-image sheaf away from zero. Their special-fiber classes are separated by characters of the finite diagonalizable group of $m$th roots of unity after adjoining $\zeta$ and inverting $m$. Dual Jacobi--Trudi is the Schur operation in D-S6, equivalently the exact cote~133 symmetric/exterior dictionary. A nonzero maximal minor makes the normal directions independent on a principal open. D-FGA and D-S3 therefore carry the dimension and padding proof; the characteristic-zero conclusion does not assert uniform good reduction at every prime.

\paragraph{G-10. Projective degree controls the first forced Hilbert gap.}
The pencil map induces a rational morphism of projective coefficient schemes. Pullback of general hyperplanes and the degree of their intersections bound the image degree, including components of the base locus. Chardin's bound then compares the graded dimensions with a polynomial subalgebra supplied by linear Noether normalization. These are D-FGA's projective, graded-module and Hilbert frameworks, with associated sheaves as in EGA~II, Proposition~2.5.2. The explicit bound is geometric effectivity; no small-size representation of every graded module is inferred.

\paragraph{G-11. Rectangular classes in a reductive quotient.}
Restriction from $\mathrm{GL}_4$ to $\mathrm{SL}_4$ sends precisely the determinant characters among polynomial irreducibles to the tensor unit. In D-S3 and D-S6, the degree-$d$ invariant module is therefore the rectangular multiplicity space. Explicit moment-map witnesses establish a nonempty stable open locus over $\mathbb C$; its quotient has dimension reduced by $15$, by the fiber-dimension statement already located in D-S3. The invariant rings are quotients because the group is linearly reductive in characteristic zero. The explicit Veronese is a standard graded coordinate algebra, so D-FGA applies again to its Proj, degree and Hilbert function. The analytic compact-group witness proves the stated characteristic-zero fact; it is not asserted to define a positive integral motive or to make reductive invariants exact over $\mathbb Z$.

\paragraph{G-12. A cocharacter degeneration and its invariant ideal.}
The homomorphism $t\mapsto\operatorname{diag}(t^3,t^{-1},t^{-1},t^{-1})$ is an integral cocharacter of $\mathrm{SL}_{4,\mathbb Z}$, as in D-S3. Its weight filtration extends the padded form over the affine line whenever $3(n-m)\geq m$. Strict inequality gives the zero special fiber; equality gives a ternary residual form with its padding factor. The ternary determinantal containment of G-05 then implies containment of invariant ideals in characteristic zero. In D-FGA this is the specialization of a morphism and the inclusion of its ideal of functions; in D-S6 it selects the determinant-character summands. It excludes precisely the specified rectangular sector.

\paragraph{I-21. Rank strata over local Artin bases.}
Every interval coefficient product is a map of free modules. The proof works over finite local Artin rings and requires a unit maximal minor, so the rank-two image and its complement split on the corresponding determinantal chart. Projection growth and its last-zero reset control the necessary coefficient supports. The two selector stages pull these finitely many rank conditions back to one nonempty open parameter locus. D-FGA supplies the determinantal and Grassmannian language, while D-S3 translates the coefficient supports to character gradings. This gives the explicit unknown-order width-three morphism; it does not assume that an arbitrary rank stratum over a nonreduced base splits without a unit minor.

\paragraph{P-12. A finite flat arithmetic evaluation and its trace dual.}
The algebra $\mathbb Z[T]/(T^N-2)$ is finite free of rank $N$. After inverting its discriminant it is finite \emph{\'etale}; over $\mathbb Q$ it is a field by Eisenstein. The integral coordinate polynomials $A_i$ define the evaluation morphism. Their bounded-degree pullback does not annihilate any nonzero member of the specified ROABP class, by I-21 and the explicit monomial encoding. The dual basis for the trace pairing extracts the transformed rational coefficient rows, which span the original coefficient-row space. This is an actual span-preserving arithmetic realization, not an injectivity assertion for differences of arbitrary members of the nonlinear width-three class. D-FGA describes the affine finite morphism and its kernel, D-S3 the coefficient lattice, and D-ET the finite \emph{\'etale} cover. Frobenius distribution and additional purity bounds are unnecessary for this construction.


\paragraph{R-17. Identical tangent-image modules and distinct isotypic sectors.}
At the cyclic pencil, the determinant and permanent differential maps have the same image up to signs of their columns. Their rank equality, together with the group-scheme fiber bound, gives the exact determinant-image dimension. D-FGA and D-S3 transport this computation. The two finite padded pairs then compare Hilbert polynomials, while the cocharacter specialization of G-12 removes the determinant-character summands. Their remaining positive differences are necessarily nonrectangular Schur summands in D-S6.

\paragraph{R-18. Two compatible normal families on the same scheme curve.}
Both local permanent identities hold on the same one-parameter deformation. The character decomposition of the special-fiber quotient now has two independent sections per permitted mode. The four-column coefficient determinant is the Schur operation $s_{(4^b)}$, interpreted in D-S6; its root-of-unity specialization is checked by a Weyl alternant. D-FGA turns their joint nonzero minor into an open rank condition, giving the improved padding range. The simultaneous nature of the deformation is necessary: ranks obtained at unrelated pencils are not added.

\paragraph{G-13. A weighted filtration inside an ordinary projective embedding.}
Generators of the reductive invariant quotient have specified positive degrees. Its cumulative graded module embeds in the ordinary polynomial-degree filtration of their affine coordinate embedding. Homogenizing that embedding gives an integral projective closure, to which the D-FGA Hilbert and degree comparison applies. Algebraically independent homogeneous sections provide the lower filtered dimension. This reduces the finite first-gap bound without asserting equality of the two filtrations or generation by a small Veronese degree.

\paragraph{I-22. Ordered coefficient modules and the radical quotient.}
The actual ROABP coefficients occupy an ordered product subspace, which need not equal the associative algebra generated by all layer coefficients. Passing to the semisimple quotient loses radical-valued sections even when their specified boundary functional is nonzero. In D-FGA these are different linear maps and different kernels; in representation $K_0$ the composition factors alone do not recover the off-diagonal extension data. The explicit triangular examples show the missing information before any trace or purity estimate can be applied.

\paragraph{I-23. Cocharacters preserving a Pl\"ucker subscheme.}
The Grassmann cone is cut out by the Pl\"ucker ideal. A diagonal cocharacter preserves that subscheme precisely when its weights satisfy the proved exchange equalities, which identify the integral torus weights induced from the original rank-$w$ bundle, up to a common character. D-S3 supplies the cocharacter lattice and D-FGA the Grassmannian. The binary classification and code-distance sum then quantify the cost of imposing this exact subscheme-preservation condition on every Hoza seed.

\paragraph{I-24. Constant-rank maps with repeated loss of remembered directions.}
The nested projections and their univariate perturbations define maps of free modules over an affine parameter scheme. Every interval lies on its specified unit-minor rank chart, including over the finite Artin quotients used in I-21. Nevertheless the selected boundary section is one monomial of exponential support. D-FGA's rank stratification records this phenomenon exactly; the ranks alone do not bound the length of a coefficient-support filtration.

\paragraph{I-25. Nested projective images on a Grassmannian.}
Invertible factors act regularly on the Grassmannian of rank-$r$ quotients or subspaces. The irreducible prefix image closures are nested because the constant factor is the identity; every strict inclusion raises dimension. D-FGA's Grassmannian chart gives dimension $r(w-r)$, and its projective geometry identifies scalar vanishing with containment in a hyperplane. Removing equal-image steps therefore gives the stated support existence bound. The invertible-factor hypothesis is what allows this proof to omit the reset behavior in I-24.

\paragraph{R-19. Compatible row and column image sections.}
Three-row motifs retain both supported row and supported column factorizations in one parameter family. Their four normal sections are separated by a six-column character calculation, whose Toeplitz determinant is the Schur operation $s_{(6^b)}$. D-S6 gives that operation, D-S3 the diagonalizable character splitting, and D-FGA the common nonzero determinantal minor. The resulting improved range is a statement about one actual pencil deformation; no addition of ranks from incompatible families is used.


\paragraph{G-14. Truncated coefficient jets and Schur-weight exclusion.}
The restriction modulo the square of the padding coordinate defines a linear
jet projection between affine parameter schemes. At a cyclic ternary pencil,
the differential is surjective: its two blocks are the cofactor-span map and
its multiplication by linear forms. The nonzero differential minor gives a
dominant morphism in characteristic zero (D-FGA). Every weight in the Schur
functor of shape $\lambda$ uses any fixed basis direction at most
$\lambda_1$ times, by its exterior-column construction in D-S6. After padding,
the low-weight part of an equation therefore factors through the first jet.
Its vanishing on the dominant determinant-jet image forces vanishing on all
padded forms. This is containment of the corresponding isotypic ideal
submodules, followed by exact passage to the characteristic-zero quotient;
it is not an assertion that the full padded varieties are contained.

\paragraph{I-26. Determinantal equations of an interpolation image.}
The maps $U$ and $V$ define bundle maps with $UV^T=0$. The determinant of
$U\operatorname{diag}(X)V^T$ is a nonzero section by a maximal-minor chart,
but its pullback to the centered interpolation parameter scheme vanishes
because that pulled-back bundle map factors through a rank-$k$ bundle.
The common center acts trivially. All exterior minors form a finite free
module on which rank-one additive translations act by commuting unipotent
matrices; the proof uses polynomial identities and remains valid in the
stated characteristics. This is the D-S3 additive group action and the
D-FGA determinantal rank locus. The surviving column-matroid bases give the
separate cocharacter nonvanishing certificate of D-S3. No purity assertion
turns the annihilated interpolation pullback into a nonzero section.


\paragraph{G-15. Harmonic covariants make the third-jet morphism dominant.}
The $\mathrm{SL}_2$ representation-ring decomposition of the endomorphism
bundle and the symmetric algebra of its adjoint bundle gives the covariants
$T_l$ in D-S3/D-S6. Their scalar coefficients in the adjugate are determined
by a finite trace pairing and its Jacobi matrix. The strictly positive
squared norm in the proof is an ordinary Hermitian linear-algebra fact,
not a weight-purity assertion. Linear multiples of these cofactor sections
surject onto the quotient by the fourth power of the padding-coordinate
ideal. Thus the differential minor gives dominance of the corresponding
third-jet morphism in D-FGA, while dimension excludes the next full jet.
The D-S6 exterior-column weight bound transfers dominance to containment
of the specified isotypic ideal submodules. This closes the precise
higher-jet gate stated conditionally in G-14.

\paragraph{I-27. A Boolean face in the cocharacter lattice.}
The Pl\"ucker ideal of the D-FGA Grassmann cone supplies integral linear
equations in the diagonal cocharacter lattice of D-S3. Under the actual
base-six code embedding, bounded integer deficits cannot carry between
digits. Its intersection with that lattice condition is consequently a
coordinate face of the seed cube. The induced original-bundle characters
then agree on two untested basis directions. The displayed rank-two
unipotent representation detects this equality as exact vanishing of a
nonzero section. No statement about arbitrary decoders follows.

\paragraph{I-28. A hyperplane dimension drop pays for the final parameter.}
The left endpoint defines a curve through its nonzero constant point of
the Grassmannian. Including it in the nested projective-image chain uses
the same dimension budget as the internal bundle automorphisms. If the
constant right functional vanishes on the final image, that image lies
in a proper hyperplane section and loses at least one dimension; this
pays for retaining the right endpoint parameter. D-FGA therefore gives
the sharper $r(w-r)$ support existence statement with both endpoints
included. Its width-three specialization has eight SV pairs, and the
finite flat trace construction of P-12 uses the corresponding seventeen
parameters when sharpened by this result.


\paragraph{R-20. Second-order tangent sections and the exact eventual rectangle range.}
The permanent coefficient morphism has a family of differentials over
the one-parameter motif scheme. Row and column relations cancel its
constant and first-order normal terms. Dividing their polynomial
sections by the square of the parameter gives regular limiting sections
of the ambient coefficient bundle. The finite certificate proves a
nonzero minor after one residue-field specialization; the explicit
degree bounds and cyclotomic argument lift that nonvanishing to the
required characteristic-zero family. D-FGA records the resulting image
dimension, and the D-S3 Fourier characters separate its translated
motifs. The invariant quotient and Veronese construction of G-11 then
give the specified rectangular Schur classes. This transports the
actual proof, including its finite threshold, rather than treating a
dimension inequality as an unspecified positivity assertion.

\paragraph{R-21. Integral image ranks and fixed-rank character extraction.}
The orbit-coefficient homomorphism of free modules has an integral Gram
matrix. Its positive-rank characteristic coefficient is a nonzero
integer, giving the displayed analytic spectral separation. A bilinear
coefficient kernel presents the same image ranks without enumerating
the full tensor basis. The diagonalizable group of D-S3 splits that
kernel into weight blocks; the D-S6 Weyl-character identity extracts
each Schur multiplicity from their ranks. D-FGA identifies the ranks
with the degreewise coordinate-image modules. The algorithm is
polynomial only when the bundle rank is fixed and degree is measured
numerically. The metric estimate is ordinary linear algebra, not an
unproved arithmetic slope transfer.

\paragraph{I-29. A persistent active subbundle and one passive rank.}
On the preserved rank-profile chart, constant transition maps split
into an identity on a rank-$r$ active subbundle and invertible or zero
maps on a rank-$(w-r)$ complement. At the last zero map, the vanishing
higher minors give a row graph over the active summand. The stable
projection argument bounds the support filtration of its coefficient
image by $r(w-r)+w-1$. The Grassmannian dimension argument for rank
preservation adds $r(w-r)$ interpolation pairs. These are explicit
rank-stratum and filtration calculations in D-FGA and D-S3, valid with
arbitrarily many resets; the polynomial pair budget retains its
separate sampling exponent.

\paragraph{I-30. Endpoint sections extend to local automorphisms.}
A nonzero constant endpoint section has a split coordinate functional.
The displayed rank-one update extends the section to a bundle
automorphism on the principal open where its determinant is nonzero.
The determinant is a unit at the center. Thus all generic and constant
interval ranks are preserved after the endpoints are absorbed in the
word. In D-FGA this is an explicit trivialization on the required rank
chart, and in D-S3 it preserves the original coordinate support
filtration. The two width-three rank stages consequently use two and
four pairs, respectively.

\paragraph{I-31. Irreducible matrix cones and deletion of stable blocks.}
The closure of scalar multiples of a polynomial matrix image is an
irreducible affine cone in D-FGA. The rank-$s$ chart identifies the
image of its last block with the graph of a map over the first $s$
columns. Equality of two consecutive projected cones allows deletion:
an invertible right action gives a closed subcone of the same dimension,
which is equal to the irreducible cone. This geometric equality
preserves every homogeneous minor and composes under multiplication
on disjoint parameter schemes. Rank descent and split rectangular
endpoints yield the stated recurrence in the original rank-$w$ bundle.
No exponential compound bundle is introduced, but the recursive block
factor remains and prevents a uniform polynomial sampling conclusion.

\paragraph{P-13. A finite flat algebra encoding any centered generator.}
After clearing Lagrange denominators, distinct base-$B$ exponents embed
the bounded parameter monomials in a single polynomial variable. The
Eisenstein quotient by $T^{B^{2K+1}}-2$ is finite free over
$\mathbb Z$; its generic fiber is a separable field and the displayed
trace-dual basis recovers every coefficient coordinate. The
annihilator argument identifies the resulting image span with the
original coefficient-image module. D-FGA gives the finite algebra and
module interpretation, and D-S3 gives the tensor-coordinate transport.
This is the ordinary algebra trace, distinct from the Frobenius traces
of D-ET, and does not acquire a purity or RH hypothesis by analogy.


\paragraph{G-16. An auxiliary dual line selects a nonrectangular Schur class.}
Form the Segre--Veronese cone of the degree-$n$ coefficient scheme with
the $a$th Veronese of the dual rank-four bundle. Its degree-$d$
coordinate module is exactly the original degree-$d$ module tensored
with $\operatorname{Sym}^{ad}V$. The D-S3 invariant functor and the
D-S6 Schur decomposition select precisely
$(ad+md/4,md/4,md/4,md/4)$. The auxiliary dual line cancels the padding
weight of the explicit Fermat point, so the enhanced cone has a stable
point for every positive padding exponent. D-FGA and the reductive
quotient locators in G-11 give the quotient dimensions; the separate
action and parameter degrees give the effective Hilbert bounds.
Thus the same dimension gap determines an actual nonrectangular
highest weight and a degree progression. The proof does not infer a
highest-weight label from a total dimension inequality alone.


\paragraph{G-17. Partial flags classify the stretched directions.}
For a partition direction $\nu\vdash ns$, subtract its last row to
obtain the polynomial highest weight $\eta$. The orbit of its extremal
line is the partial flag scheme embedded by the exterior-power Schur
construction. Its degree-$d$ coordinate module is the Cartan component
$S_{d\eta}V$, as follows directly by dualizing the orbit span of the
extremal tensor power. Pairing this cone with the $s$th Veronese of the
coefficient scheme makes the D-S3 invariant functor extract exactly
the multiplicity at $d\nu$, in original coordinate degree $sd$.
Positive Fermat coefficients balance the moment map precisely in the
strict region $\nu_4>0$, $\nu_1>(n-m)s$. D-FGA gives the quotient
dimension comparison and its effective Hilbert bounds. The two
excluded faces follow from the earlier ideal containments, so this
is an exact stretched-direction criterion, with its dimension
hypothesis retained. D-S6 records every exterior, symmetric and Schur
operation used in the construction.

\paragraph{I-32. A common-center rank chart with exponential deletion cost.}
Each singular factor is lifted to two univariate bundle maps whose
product depends only on the difference of their coordinates. On the
diagonal center those differences vanish, and every partial interval
is a product of nested projections between invertible boundary maps.
Thus all interval ranks are already correct on that D-FGA chart.
Nevertheless a homogeneous matrix-coordinate section is the product
of all paired differences. A coordinate subspace meets its nonvanishing
locus only if it retains one coordinate from every pair. This gives
the exact exponential threshold for preserving the image cone, while
a dense D-S3 interpolation map keeps the same section nonzero. The
obstruction concerns this coordinate-deletion criterion only.

\paragraph{I-33. Divisors on the one-dimensional torus and compatible weights.}
The prime-weight substitutions are integral cocharacters of the
split torus in D-S3. Pulling back the rational ratio of two products
gives divisors on the multiplicative line. Distinct prime root orbits
can overlap in at most one point, while the inseparable case is
controlled by integer divisor multiplicities. This forces any lost
ratio to be a character, which the coordinate variants separate.
The subset-sum example identifies the exact failure of closure under
products for the raw family. Interpolating its finite evaluation set
by one affine curve repairs scalar product closure because the
coordinate ring of the line is a domain. D-FGA places this curve and
the coefficient-image restriction; no purity theorem is needed.


\paragraph{G-18. Removing scalar orbit periods and comparing every large degree.}
The auxiliary Schur cone and its D-S6 Cartan coordinate modules remain
those of G-17. The smooth Fermat factor has finite projective stabilizer;
varying its padding line in a common stable open avoids every nontrivial
fixed locus. Unique factorization then gives a free projective orbit.
In the D-S3 central group scheme $\mu_4\subset\mathrm{SL}_4$, the
character exponent $ns-|\eta|=4\nu_4$ is trivial. Consequently two
distinct nonzero scalar multiples of the chosen tensor have distinct
closed orbits. Exact reductive invariants separate them, forcing the
degrees of common nonzero invariant sections to generate $\mathbb Z$.
The explicit numerical-semigroup argument gives sections in every
degree at least $\beta^2$. Multiplication by these sections is injective
in both quotient coordinate domains. D-FGA's standard Veronese and
Hilbert-function comparison can therefore be transported to every
integer degree, with the stated effective thresholds and common
asymptotic exponent. The rectangular central congruence is already
encoded by $ns=4\nu_4$. All stability, exact invariant-functor, and
quotient assertions here are in characteristic zero; an integral
model alone is not an assertion of the same behavior at every prime.

\paragraph{I-34. Local relation modules and coordinate cylinders.}
At one factor position, take the kernel of the univariate coefficient
map from the free module on product terms. A rank-one kernel with full
support is a circuit. The D-S3 prime cocharacters and the divisor
argument of I-33 force every canceled multiplicative ratio off this
position to be a character. A free coordinate separates its character
exponents, and the circuit property excludes every proper relation.
Higher local nullity is reduced by coefficient elimination. The result
gives the explicit bound on the number of free coordinate directions;
interpolation replaces the cylinders by integral cocharacters. In
D-FGA terms, the finite family preserves the coefficient-image module
after extension to the working field. The rank-four flattening example
defines a genuine nonzero tensor in the kernel of every one-coordinate
cylinder and identifies the precise failure of that smaller family.

\paragraph{P-14. Rank-one reset maps and scalar bridge sections.}
On a rank-one bundle-map chart, a nonzero pivot writes the map as a
column times a row divided by that pivot. Clearing all pivots splits
the scalar contraction into bridge sections. Invertible interior maps
and absorbed endpoint maps have identity constants over $F(t)$; the
ascending coefficient-row subspaces have at most $w-1$ strict growth
steps. One common interpolation morphism keeps every nonzero bridge
section nonzero, and their product remains nonzero in its coordinate
domain even when adjacent bridges share a reset variable. Constant
linear annihilators give exact coefficient-row preservation. This is
the D-FGA image-module criterion and the D-S3 dense substitution; no
equality of matrix image cones or purity statement is required. The
balanced rational-evaluation form follows from the stated partial
fractions and the same scalar-annihilator argument.


\paragraph{R-22. Cyclic characters and one open differential-rank condition.}
The permanent coefficient morphism is an integral polynomial map of
affine schemes. R-22 works over a cyclotomic extension in characteristic
zero. The D-S3 diagonal torus restricts to $\mu_m$, whose characters
split the differential coefficient space. A single staircase parameter
space supplies both local row and column cofactor relations. Their
highest components define two subspaces with explicitly computed
coprime-factor intersection. The D-S6 Schur determinant is the nonzero
minor of coefficient restriction; its principal geometric specialization
is a Vandermonde determinant with distinct cyclotomic exponents.
The product of finitely many nonzero parameter polynomials defines one
nonempty open subset on which all bidegrees are simultaneously valid.
An integer grid contains a point of that open after the root of unity
is fixed; translated motifs use the prescribed cyclotomic twists.
The final D-FGA differential minor is nonzero in one common deformation,
so its image dimension is at least the stated bound. The proof can be
spread over the localization where these particular denominators and
minors remain nonzero; it does not assert the same rank in every residue
characteristic. Its quotient and highest-weight consequences use the
characteristic-zero dictionaries of G-10 and G-17--G-18.

\paragraph{The G-17 quantitative corollary.}
The rank of a symmetric power and the six-factor Schur-rank polynomial
bound the ambient bundle rank in the D-S6 construction. Together with
the polynomial action degree, these give the binary-length estimate
for the existing first successful quotient degree. This is an effective
bound on the degrees of sections in the D-FGA graded algebra; it is
separate from computing the section or deciding its multiplicity.


\paragraph{R-23. A single graded part of a derived intersection.}
Let $Z=\operatorname{Spec}(R/J_m)$ be the permanent Jacobian scheme,
and let $L=\operatorname{Spec}S$ be a general four-plane in
$\operatorname{Spec}R$. Their D-Tor derived fiber product has
coordinate complex $(R/J_m)\otimes_R^{\mathbf L}S$.
Its first homology in internal degree $m$ is precisely the dimension
defect $\epsilon_m$. The known full-variable linear resolution gives
the preceding map, and the cyclic coefficient calculation proves
that its $2m-2$ old syzygies remain independent while no new constant
relation appears. The degree-$m$ differential then gives
$c_m=4m^2-2m+2-\epsilon_m$. This is an equality of dimensions of
explicit cohomology groups, not an inference from an alternating
$K_0$ class. The low-size positive defects prove failure of the
unqualified restriction argument. The depth bound rules out a full
regular-sequence justification but leaves the single graded
vanishing question open. D-FGA places the associated generic rank
stratum, and D-S6 distinguishes its actual homology from its
Euler characteristic.

\paragraph{P-15. A nonzero section killed by a dense interpolation morphism.}
The displayed quadratic is a rational coefficient tensor with
minimum tensor-train width three: its three balanced flattenings
have ranks $3,4,4$. The upper-triangular bundle-map factorization
keeps every interval on the invertibility open set. Nevertheless
its pullback under the specified centered Lagrange morphism is
zero. The identities in the three linear coordinates $S,U,V$
prove this directly; the harmonic cross-ratio identity supplies
the separately stated rational-evaluation examples. D-FGA records
the nonzero section and its zero pullback, while D-S3 records the
coordinate choices. The example is over $\mathbb Q(\mathrm i)$
for the fourth-root Lagrange nodes, and uses rational abscissas
for the other convention. No scalar change of coordinates is
silently omitted in comparing these two maps.


\paragraph{R-24. The first zero of the graded derived-intersection defect.}
The specified integer Jacobian minor is a section of an exterior
power of the permanent coefficient-map differential. Its nonzero
residue modulo $101$ proves that its principal open locus has
nonempty characteristic-zero fiber. On that locus the image
dimension is $1124$, the exact D-S3 torus quotient bound.
R-23 then identifies $\epsilon_{17}=0$ as vanishing of the
indicated degree of the D-Tor first homology. This is one
explicit generic-fiber result; no complete Tor independence
or assertion in every residue characteristic follows. The
padded coefficient image and resulting highest weights use
the D-FGA and D-S6 dictionaries already proved in G-17--G-18.

\paragraph{R-25. A single family, canceled tangent sections, and induction.}
The bordered matrix is a morphism from an affine parameter scheme
times $\mathbb A^1_\varepsilon$ to the space of forms. Two
coordinated parameter variations cancel their order-one terms
exactly; division by $\varepsilon^3$ on the punctured parameter
line gives fixed Hessian sections in the same differential image.
Their restrictions to the hyperplane $\ell=0$ are the stated
ternary coefficient columns. The intersection of two nonempty
rank opens in the irreducible parameter space permits one matrix
with both required properties. A single minor with nonzero
constant term proves the dimension increment. This transports
the proof to D-FGA rank loci and tangent maps; it does not infer
rank from an alternating $K_0$ class. The finite integer minor
certifies the antecedent at $m=17$ and hence $c_{18}=1262$.
The uniform nonemptiness of those rank opens for all larger
sizes remains the specific geometric question.

\paragraph{I-35. Base loci on a Segre scheme and the seed diagonal.}
The two boundary vectors define a point of
$\mathbb P^{r-1}\times\mathbb P^{r-1}$, embedded as rank-one
functionals. Equality of their annihilator intersections is
equality of base loci of two linear systems on this Segre
scheme. The incidence dimension argument and D-Proper give
an existential subspace with the same base locus and at most
$2r-1$ generators. It need not be realized by an explicit
substitution on the original variables. For the generic matrix,
the Hankel curve realizes such a linear system, although the
matrix value has rank one. This is the known optimal test-space
dimension cited in P-B16, with the new bridge deduction stated
separately.

Independent seed sets give a product of parameter schemes.
At each replacement the other parameters belong to the
coefficient function field, so the extension-field contraction
hypothesis applies. Identifying those parameters afterward is
pullback along a diagonal; a nonzero section can become zero
under that pullback, as the explicit examples prove. Thus
base change of the coefficient field is distinct from seed
identification. Finally the multiplicative lift sends the
rank-$r$ idempotent corner into a one-dimensional endomorphism
corner; its fixed entry readout would have to be a multiplicative
character, contradicted by the two unipotent matrices. D-S3
places these actions, D-S6 places their exterior powers, and
D-FGA places their coefficient-image modules. No new purity
or representation-rank conclusion is inserted by this language.

\paragraph{G-19. A discriminant line and a specified Schur submodule.}
Over the integer coefficient scheme, the homogeneous form and
its partial derivatives define a closed incidence scheme with
projective point coordinates. Its projection is closed by
D-Proper. On the characteristic-zero fiber its reduced
discriminant equation is a section of coefficient degree
$D=k(m-1)^{k-1}$ and variable-change determinant weight
$\delta=m(m-1)^{k-1}$. Benoist's exact degree theorem and
primitive integer construction supply this line; D-S3 records
its character. The highest-weight inclusion into the plethysm
generates one $S_{(\delta^k)}V$ over $\mathbb C$, in the D-S6
Schur-functor language. Its restriction map to one orbit
coordinate ring is zero and to the other injective. This is
a statement about a specified submodule of the coefficient
algebra, not the rank of its entire isotypic multiplicity space.

The Haar identity is its compact complex realization; Gaussian
orthogonality is a positive norm of its integer pullback
coefficients. Neither norm is asserted to arise from a pure
arithmetic sheaf or motive. The explicit cubic Macaulay minor
defines an open localization where the graded Jacobian map
is onto, so the projective gradient-zero scheme is empty.
Its nonzero residue certifies the characteristic-zero smooth
fiber and hence the displayed discriminant nonvanishing.
Positive padding factors the section and forces a singular
intersection, so the discriminant pullback becomes zero.
This same geometric argument proves the stated limitation.
The row-minimality at size three is the D-S3 highest-weight
subalgebra argument applied to the dense D-FGA four-variable
determinant image. Cote~133, scan~028, supplies the Schur
operation viewpoint; scan~019's rank-versus-character warning
remains essential for the specified-copy distinction.


\paragraph{P-16. Reciprocal coordinates, residue functionals, and a necessary center.}
The four disjoint pole divisors lie on the projective line over
the explicitly bounded number field. Their residue functionals
annihilate polynomials of the stated degrees because of the
orders of vanishing at infinity. The displayed rational identity
then makes one nonzero width-three tensor section vanish on
the unbalanced evaluation image. Passing to the coordinate
torus and applying inversion exchanges numerator and denominator
degree. Multiplication by the full coordinate monomial extends
this pullback to the multilinear coefficient tensor; on each
rank-two variable factor it exchanges the two basis vectors.
Every coefficient flattening is consequently changed only by
row and column permutation. This is a direct equality of ranks
of bundle maps in D-FGA, with D-S3 supplying the torus inversion
and diagonal node rescaling; no rank-monotonicity assumption
about polynomial factors is made.

The simple-root discriminant of the node polynomial makes all
$A'(b)$ invertible in the working field, so the Lagrange and
unbalanced rational maps have the stated exact diagonal
identification. Clearing denominators extends the identity
from the dense torus chart to the whole defined evaluation
family. Adding the common center is a different affine
parameter morphism. Its nonvanishing follows from P-14 because
all displayed transition determinants are nonzero polynomials.
Thus the same section is zero on one image and nonzero on the
centered image. D-S6 places the coefficient tensor operations;
it does not identify these two pullback maps merely because
their target bundle ranks agree. The construction uses the
specified algebraic nodes, and supplies no counterexample
for an arbitrary preassigned node set.


\paragraph{I-36. Pairwise dominant pullback and rank-one exterior maps.}
The centered one-pair morphism is dominant on every coordinate
pair: its determinant over the selector function field is the
difference of two Lagrange basis sections and is nonzero.
Hence its pullback is injective on each such bivariate
coordinate algebra. On a rank-one bundle-map chart, a pivot
writes the map as a column and a row divided by a nonzero
univariate section. Clearing the pivots turns an interval
contraction into a product of univariate endpoint sections
and adjacent bivariate sections. Each survives the same
pullback, and their product survives in the parameter domain.

For factors of rank at most $r$, the $r$th exterior map is
zero or rank one. Functoriality of exterior powers, equivalently
Cauchy--Binet, identifies its interval entries with the
original maximal minors. The first paragraph therefore
preserves each nonzero such section. D-S6, V~2.1--2.2.1,
places the exterior-power construction, D-S3 places its
coordinate action, and D-FGA places the resulting determinantal
rank opens. In width three with singular factors this supplies
the rank-two-preserving center. The four residual interpolation
pairs from I-30 preserve the remaining local rank-one tests,
giving five pairs in total. This proof uses actual nonzero
minors throughout; no numerical rank or purity argument is
substituted for the exact pullback assertion. Arbitrary
variable-dependent invertible width-three factors remain
outside this five-pair refinement.


\paragraph{R-26. Characterwise determinantal opens and a linear derived defect.}
The Jacobi parameters form an integral affine scheme; restricting
to nonzero diagonal and edge parameters gives an integral torus
chart. The finite cyclic action on the ternary variables decomposes
the coefficient bundle over characteristic zero into its character
subbundles. In each ordinary character, eight specified normal
sections have a nonzero determinant on a nonempty open chart.
The seventeen finite cores certify those opens by reduction modulo
$101$; their denominators are invertible at the stated certificates.
These are exact algebraic determinants, not approximate ranks.

The cut argument takes place over a Laurent-series field: divide
the second weight component by its path product, replace each cut
edge by a parameter times a chosen nonzero constant, and specialize
the resulting regular normalized matrix at parameter zero. Equal
outer cut orders allow its path-product ratio to have any specified
nonzero limit. Inserting isolated vertices multiplies the indicated
component sections by the common coordinate $x$, an injective map
of the coefficient modules. The two insertion rules therefore
transport each certified determinantal open to all required sizes.
The exceptional character sections are independent by the
coprimeness of consecutive continuants. Irreducibility of the
parameter scheme makes the finite intersection of all these opens
nonempty. This is the scheme-theoretic content of the construction;
it does not assert that one literal zero-edge pencil realizes all
the normalized certificates simultaneously.

Apply the R-25 normal-direction injection and the D-Tor dictionary
of R-23. The dimension of the degree-$m$ first derived tensor
homology is at most $6(m-18)$ for $m\ge18$. Thus the exact remaining
failure of the restricted linear strand is bounded linearly.
D-S3 places the character action, D-FGA places the coefficient
rank opens and orbit-image dimensions, and D-Tor (SGA~6, III~1.5--1.6
and I~5.1--5.4) places the derived restriction. The consequence is a
characteristic-zero bound; the use of modular certificates alone
does not give the whole statement in every positive characteristic.
Combining these dimension bounds with the flag quotients of
G-17/G-18 gives the exact rectangular padding range already for
$m\ge14$, as detailed in R-26. No new positivity conjecture is
substituted for those proved quotient comparisons.

\paragraph{I-37. Rank-two active bundles and the second coefficient form.}
The rational centered two-pole family defines a morphism from an
open parameter scheme to the coordinate affine space. Its one-pole
subfamily is dominant on each coordinate pair. An irreducible
multilinear section can be considered separately because its
irreducible factors use disjoint coordinate sets; specializing the
other coordinates preserves the affine singular factorization.
On the nonzero minor charts, every strict internal product has
rank two. A rank-one internal product would have a polynomial
column--row factorization over the coordinate UFD, and would split
the irreducible scalar into two nonconstant factors.

Exterior squares and the pairwise dominant pullback preserve these
rank-two charts, by I-36. Compatible bases split each intermediate
rank-three bundle into a persistent active rank-two summand and
a passive line, with constant transition $\operatorname{diag}(1,1,0)$.
The affine singular condition makes the linear passive self-entry
zero. If the original section vanished on the two-pole image,
simple-pole residues would annihilate all first coefficient forms
at the one-pole center. Translation along its nonzero common
slope vector annihilates the contraction of the second coefficient
form with that vector. In the compatible bases this form has only
adjacent off-diagonal entries. Induction along the ordered path
then makes each of them zero in every characteristic. Every
possible active-to-passive-to-active output transition consequently
has a zero product, forcing the whole section to be zero.

This restates the complete contradiction in bundle-map language.
The second coefficient form is extracted formally from multilinear
polynomials, so division by two is never used. D-S6, V~2.1--2.2.1,
places the exterior-square maps; D-S3 places the changes of frame;
D-FGA places the rank strata and nonvanishing charts. The final
diagonal node rescaling identifies the rational family with
ordinary centered Lagrange interpolation. Its five parameter
coordinates yield the exact grid count in I-37. Invertible
variable-dependent transitions can propagate the passive line,
so the tridiagonal argument does not extend to that unresolved
mixed-rank case.


\paragraph{I-38. Triple-dominant pullback and isolated invertible transitions.}
The one-pole parameter morphism is dominant on every three
coordinate factors. Indeed, after solving for the center and
amplitude using two coordinates, the third is a nonconstant
rational function of the selector over their function field:
three distinct simple poles forbid both interpolation coefficients
from being constant. This argument establishes dominance even
when a characteristic-two Jacobian test is inseparable.

For the promised alternating rank pattern, exterior-square maps
of the rank-two factors have rank one, and there is at most one
invertible exterior-square map between two of them. Each scalar
connector is therefore a section depending on at most three
coordinates. Triple dominance preserves every required nonzero
connector. Compatible active rank-two summands and passive lines
then put the constant maps in the forms $I_3$ and
$\operatorname{diag}(I_2,0)$. The rank pattern survives the
irreducible multilinear reduction: a specialized singular factor
between two surviving invertible ones is absorbed into the latter
and makes it singular.

The second coefficient form at the one-pole center annihilates
both vectors $h_i=(a_i-y)^{-1}$ and $h_i^2$. Its only possible
entries are within two positions of the diagonal. The determinant
$h_jh_k(h_k-h_j)$ is a nonzero section, in every characteristic,
so row induction annihilates the whole form. Every passive
excursion of an output path has one of its zero entries as a
factor. This transports the contradiction of I-38 without
asserting preservation of each individual minor, and without
discarding passive propagation at invertible factors.
D-S6, V~2.1--2.2.1, supplies exterior powers; D-S3 supplies the
bundle changes of frame; D-FGA supplies the rank-open and
function-field formulation. Longer consecutive invertible runs
require connectors on four or more coordinates and remain outside
this triple-dominance proof.


\paragraph{R-27. Formal endpoint jets and exact derived restriction.}
The endpoint deformation gives a matrix of coefficient sections
over the formal parameter ring $\mathbb C[[\varepsilon]]$.
Its first and second normal maps are the successive leading
coefficient maps after eliminating the constant-rank image.
The lower triangular coefficient matrices represent its
reductions modulo powers of $\varepsilon$. Their ranks therefore
count the unit, first-order and second-order pivots exactly.
The boundary projection retains complete character subbundles.
The support of the two successive normal maps lies wholly
inside those subbundles, so the projected quotients agree
with the full quotients; discarded character sections cannot
create an additional correction.

The fixed size-seventeen certificate proves nonempty opens
for the three coefficient ranks. Central insertion transports
its normalized sections by the injective monomial maps
$f\mapsto xf$, $f\mapsto yf$, $f\mapsto zf$ on disjoint
weight windows. The cut is interpreted in a Laurent-series
chart after canceling the long-path edge product. The exact
endpoint identities give a separate upper bound of four
on the first normal rank. This upper bound is required:
lower rank bounds alone would not control the second normal
map under specialization. Intersecting the rank opens and
this identity chart gives all six missing normal directions
at a single actual pencil. R-25 then makes the restriction
complex exact in the specified degree in every size at least
seventeen. D-S3 places the character action and changes of
frame, D-FGA places the determinantal opens, and D-Tor
(SGA~6, III~1.5--1.6 and I~5.1--5.4) places this precise
derived-tensor vanishing. No regular-sequence assertion in
all degrees is made.

\paragraph{G-20. Exact comparison of quotient Hilbert dimensions.}
Insert R-27's exact image dimension into the enhanced
partial-flag quotients of G-17/G-18. Their dimension difference
is the integer $C_m-2n^2+1$. The two standard-Veronese
Hilbert polynomials have positive leading coefficients;
their original-degree ratio therefore has exactly that
growth exponent. The exponent is never zero for the integer
pairs in the theorem. This removes the former unresolved
dimension transition, without determining every finite
graded piece. D-FGA and EGA~II, 2.5.2, place the projective
coordinate rings; the invariant-quotient and stabilizer
locators are the same as those already recorded for G-17/G-18.
The first-row condition for rectangular weights converts
to $n<4m/3$. The size-thirteen example proves sharpness of
the threshold fourteen for the uniform eventual statement,+not absence of an individual earlier-degree equation.

\paragraph{G-21. Additional-variable tangent summands.}
The cofactor sections form an independent rank-$m^2$
subbundle on a nonempty quaternary parameter open. Intersect
it with the exact differential-rank open of R-27. Extending
the coefficient bundle by unused variables adds one copy
of that cofactor subbundle in each new variable multidegree.
These summands have zero intersections with one another
and with the old tangent image. The resulting rank is
$rm^2-2m+2$, equal to the row--column orbit-parameter bound.
The same degree-$m$ tensor complex from D-Tor has kernel
of dimension $2m-2$, exhausted by the old linear relations.
Independence of the cofactors handles degree $m-1$; lower
degrees have zero source. This proves the exact graded
homology statement, including its stated limits. For $r>m^2$
the coefficient specialization is an algebra morphism into
a larger polynomial ring, not a quotient by a negative
number of equations.

\paragraph{I-39. Failure of rank-stratum preservation and nonintegrable jets.}
The ordinary centered one-pair morphism sends the explicitly
given generic rank-two interval to the rank-one stratum,
even though the scalar coordinate section is irreducible
and stays nonzero. The two singular anchors and the unipotent
interior give an exact determinantal equation for this
specialization. Thus preservation of every internal rank
stratum is a stronger invariant than scalar nonvanishing.
D-FGA supplies the determinantal-stratum language, while
D-S3 places the affine unipotent transition factors.

The separate Hessian family is a nonzero symmetric section
with the asserted tangent kernels and rational-curve
quadratic restriction. It cannot be the second coefficient
pullback of a fixed polynomial along the whole one-pole
parameter morphism: restriction to zero amplitude would
force independence of the selector, contrary to its explicit
entries. This is an exact failure of the proposed abstract
jet criterion. It is not a counterexample to the two-pair
scalar nonvanishing problem; the remaining construction
must retain the coordinate-ring and integrability conditions.


\paragraph{R-28 and G-22. Added tangent bundles and Fourier-frame stability.}
The determinant differential has an independent cofactor subbundle
of rank $n^2$ on a nonempty quaternary parameter open. Each added
variable contributes a copy in a separate multidegree. The resulting
rank $(r-2)n^2+2$ equals the parameter dimension minus the dimension
of the generically finite-stabilizer left--right special-linear
orbit. The permanent padding map is generically injective after
projectivization: the irreducible nonlinear factor and padding
line are recovered uniquely. D-S3 identifies the group actions,
and D-FGA and EGA~IV$_3$, 13.1.3, identify these rank opens and
fiber-dimension calculations.

In G-22, the partial Fourier matrix is a frame in a rank-$r$
locally free representation, after scalar extension to $\mathbb C$.
Its rank-one Hermitian moments sum to the required diagonal
weight. The ordered tensor and the extremal flag line have
zero total traceless moment and finite projective stabilizer.
Multiplication of the linear sections and the outer symmetric
power preserve the minimum weight for every cocharacter, since
the corresponding initial sections have nonzero products in a
symmetric algebra. The Hilbert--Mumford criterion therefore
transports stability to the multiplied section. This is a
cocharacter and initial-term dictionary in D-S3; it does not
assert equality of the two Hermitian moment maps. The Schur
and exterior tensor operations belong to D-S6, V~2.1--2.2.1,
V~3.2 and V~3.10, with the precise representation-ring convention
already fixed above. The square roots give a complex existence
witness, not an integral point or an effective arithmetic sampler.

The cyclic product-plus-power family is an irreducible parameter
scheme. Nonzero coefficient minors impose the finite projective
Lie-stabilizer condition. Irreducibility of a form is open because
the reducible locus is a finite union of proper projective
multiplication images. An explicit Eisenstein section and the
coefficient calculation of its stabilizer show these opens are
nonempty. Their intersection with the stable open remains nonempty.
Varying the padding line avoids the fixed loci of the finite
automorphism group and leaves only the central group scheme
$\mu_r$ upstairs. Its character on the enhanced representation
is $r\nu_r$ and is trivial. Thus distinct nonzero scalar multiples
have distinct closed orbits. Invariant separation makes the
nonzero evaluation degrees have gcd one; the numerical-semigroup
argument gives every degree from $\beta^2$. Proper projective
images use D-Proper, and the character and stabilizer statements
use D-S3. All these orbit assertions are over characteristic zero.

\paragraph{G-23. Exact growth of enhanced invariant coordinate algebras.}
The Cartan cone attached to the flag weight has degree-$j$
coordinate module $S_{j\eta}V$. The diagonal Segre--Veronese
algebra is therefore
$\mathbb C[Y]_{sj}\otimes S_{j\eta}V$; taking invariants
extracts the single label $j\nu$. The stable quotient dimensions
are $u=(r-2)n^2+r_\eta-r^2+3$ and
$v=rm^2-2m+r+r_\eta-r^2+2$. This is a dimension computation
on two projective invariant schemes, not an equality of their
coordinate modules in $K_0$.

The explicit standard Veronese has $E=rT\beta!$.
Its positive leading Hilbert coefficients, combined with
nonzero sections in all sufficiently large original degrees,
give the ratio $\kappa j^{v-u}(1+O(j^{-1}))$ through every
integer stretch. The character-ring reduction cannot remove
the difference $v-u$; its nonvanishing is a separate elementary
congruence argument. D-FGA and EGA~II, 2.5.2, place the graded
coordinate algebras and projectivizations. The coefficient-degree
bound on the action and the parameter-degree bound on the orbit
map remain distinct. The construction proves the eventual sign
in the stated strict cone with $4\le r\le m$; boundary faces,
more rows, and effective positive extraction at individual
finite degrees remain outside this conclusion.

\paragraph{P-17. Exterior decoders, coefficient-rank strata, and odd weights.}
On the integral node scheme obtained by inverting all node
differences, the three-row Vandermonde matrix is a split
surjection and its kernel $K$ is locally free of rank $n-3$.
The cubic coefficient decoder is a diagonal isomorphism with
unit Vandermonde factors. Its balanced-ideal part is
$\bigwedge^3 K$, by the explicitly cited ideal theorem.
The coefficient-cut rank bounds are determinantal closed
conditions. The three-block contradiction, the support-at-least-four
property, and the leading exterior projections are valid in
every residue characteristic on this node scheme. The two end
kernel subbundles have intersection ranks zero and one at
$n=9$ and $n=10$, yielding respectively the incompatible wedge
degrees and the forced common exterior divisor in the proof.

For an invariant multilinear tensor, each Schmidt support is
an $\mathrm{SL}_2$ subrepresentation on its fixed-rank chart.
At an odd cut its split-torus characters are nonzero odd
integers. The Weyl element pairs opposite character spaces, so their total rank is even, including over characteristic two.
With rank at most three this forces rank two at every odd cut,
and each affine core factors through a rank-two bundle.
The two-pair theorem I-37 then applies. Translation and reversal
are the two unipotent operations on the multihomogenization;
the closed coefficient-rank degeneration to the leading
translation term gives the eight-variable closure. The exact
torus reference is SGA~3, I~4.7.3; exterior and symmetric
operations use D-S6, and the rank-stratum formulation uses
D-FGA. No division by two or semisimplicity in positive
characteristic is used.

\paragraph{I-40. Determinant characters of Borel-stable support bundles.}
For the homogeneous common-translation-invariant tensor, the
two support bundles at a cut are Borel-stable, although they
need not be $\mathrm{SL}_2$-stable. Their determinant lines
are upper-unipotent fixed lines inside exterior powers of
full rational $\mathrm{SL}_2$ representations. The standard
dominant-character fact cited from Jantzen, Lecture~I,
pp.~2--3, gives nonnegative integer weights $\delta_i,\delta_i'$.
The nonsingular coefficient pairing of total weight $m=n-2D$
gives $\delta_i+\delta_i'=r_i m$. This is an equality in the
character lattice of the split torus, not a congruence in
the residue field.

The support inclusions $L_i\to L_{i-1}\otimes V$ give
equivariant affine core maps. Their unipotent covariance
makes every nonzero square determinant constant in the local
variable. Its constant term is then invertible, and its
torus covariance gives the exact increment
$\delta_i-\delta_{i-1}=r_i$ across a full-rank square core.
Telescoping the determinant-line character proves the
full-width run bound $n-2D$. Constant absorption preserves
it because a removed singular map cannot leave the next
full-width map invertible. At width three, weights zero and
one put the word in I-37 and I-38 respectively. The cubic
closure of P-17 then makes a nine-variable quartic annihilator
translation invariant of weight one and gives the contradiction.

This restatement uses the determinant exterior operation in
D-S6, the split torus and Borel group schemes in D-S3, and
the coefficient-rank charts in D-FGA. The rational evaluation
map is regular after inverting the node-pole differences.
Its generic two-simple-pole chart is dense in the balanced
degree-two rational family even in characteristic two.
Clearing the squarefree coordinate denominator and using the
explicit selector and Newton-simplex grids gives the finite
sample count in I-40, with its working-field and extension-input
conditions retained.


\paragraph{G-24. Boundary cocharacters and restriction of invariant algebras.}
At $n=ra$ and $m=(r-1)a$, the cocharacter with weights
$(r-1,-1,\ldots,-1)$ sends a padded form to
$x_0^a q(0,x_1,\ldots,x_{r-1})$. Every invariant section
has the same value on this limit. Restriction therefore
injects the padded invariant coordinate algebra, with its
original grading, into the $(r-1)$-variable permanent
invariant algebra. Injectivity follows from the limit on
the dense padded image, not from an assumed equality of
the two algebras. This is a cocharacter-specialization
statement in D-S3 and a morphism of graded coordinate
algebras in D-FGA.

For $r\ge5$ the lower algebra has the exact stable quotient
dimension supplied by G-21. For $r=4$ it is bounded degree
by degree by the full ternary invariant algebra, whose map
onto the lower permanent algebra is a surjection by
exactness of invariants. R-28 supplies a separate free
stable point in each comparison ring by adding one linear
factor; the integer $a\ge2$ makes both auxiliary directions
strict. Thus both degree semigroups have gcd one. Their
dimensions satisfy $u>v$, and the explicit common generator
bound and Veronese make the Hilbert comparison strict at
every sufficiently large original degree. Above the
boundary the same cocharacter sends the form to zero.
Together with G-23 this proves the exact eventual rectangular
range through $m$ rows, with the stated finite-degree and
nonrectangular exclusions from its scope. EGA~II, 2.5.2,
places the projectivizations, and the stable quotient and
image-degree locators are those already recorded for G-23.


\paragraph{R-29 and G-25. Two-product families and enlarged stable quotient schemes.}
Work over $\mathbb C$ with a rank-$r$ vector space split as two
summands of ranks at most $m-1$. The ordered linear-factor tensors,
their symmetric products, and the auxiliary Cartan flag cone are
formed by the tensor and Schur operations of D-S6. Minimum
cocharacter weights are preserved under multiplication in the
symmetric algebra. Kempf--Ness and unitary transitivity on the
complex partial flag variety give a product section with any
prescribed positive rational moment spectrum in each summand.
This analytic normalization is an existence argument on the
complex fiber; it does not assert an integral section or an
effective arithmetic change of basis.

The orthogonal sum of the two product sections has the required
moment spectrum. The Lie-stabilizer equations split by bidegree:
mixed blocks vanish by independence of directional derivatives,
and diagonal blocks are scalar. The affine scheme parametrizing
two products and a padding line contains the stable open, the
irreducible-form open, and the finite-projective-stabilizer open.
Their nonempty intersection and avoidance of finitely many fixed
line loci give a common stable section with only the central
projective stabilizer. Its central character is $r\nu_r$ and
is trivial on $\mu_r$. These are group-scheme and character-lattice
statements in D-S3 on the characteristic-zero fiber; proper
projective multiplication images use D-Proper.

The degree-$j$ Cartan coefficient algebra still extracts precisely
the label $j\nu$. Hence the projective invariant schemes have
the same dimensions $u,v$ as G-23 for $r\le2m-2$. G-25 checks
the only changed dimension lower bounds and carries across the
same standard Veronese, positive Hilbert leading terms and
gcd-one degree semigroup. The exact difference is unchanged.
For the new positive family $r=n=m+1$ it equals one, and the
label is $(3j,(2j)^{m-1},j)$ in degree $2j$.
D-FGA and EGA~II, 2.5.2, place these graded algebras and their
projectivizations. This comparison of actual invariant Hilbert
functions is stronger than exhibiting one equation copy, but
does not provide a general positive formula for individual
finite-degree multiplicities.

\paragraph{I-41. Saturated weights, divided squares, and pole ranks.}
On the fixed-rank coefficient charts of I-40, a width-three run
of two invertible cores exhausts the determinant-character budget
at weight two. The bounding determinant lines have character zero
and are fixed by the whole split $\mathrm{SL}_2$. Their decomposable
exterior tensors determine invariant rank-three support bundles.
Single-coordinate contractions give neighboring invariant
rank-two bundles. Torus character parity and the Weyl involution
force weights $1,-1$: larger opposite weights would leave no
odd-rank invariant subbundle in their tensor product with the
standard bundle. This uses the integral character lattice of
SGA~3, I~4.7.3, rather than reduction of weights modulo the
residue characteristic.

The invariant rank-three subbundle in $V\otimes V$ is explicitly
spanned by $e_0e_0$, $e_0e_1+e_1e_0$, and $e_1e_1$.
It is the divided-square tensor submodule in every characteristic.
The proof retains its actual dual coefficient pairing; it neither
identifies it with its dual in characteristic two nor assumes
semisimplicity. Determinant and tensor operations use D-S6,
and the support charts use D-FGA. The character shifts make
the two middle affine maps upper unitriangular in compatible
weight bases. Their two-dimensional connector determinant has
the nonzero square coefficient $d_2d_3$.

Over the node scheme with all node differences inverted, the
one-pole rational evaluation is regular after also inverting
the node-pole differences. The unique double pole of that square
term proves nonvanishing of each saturated connector. Compound
rank-one factorization then preserves every required interval
rank. On the resulting compatible charts the second-coefficient
matrix has bandwidth three and annihilates the three sections
$1,h,h^2$. Their three-column Vandermonde determinants are units,
so the row induction and passive-path argument force the
contradiction in every characteristic. Translation and reversal
are root-group operations; their leading limits stay in the
closed determinantal width locus. They reduce the ten-variable
degrees to the weight-two and weight-zero cases. The balanced
finite interpolation grid is the same regular-open construction
as in I-40, with all field-size and extension conventions retained.


\paragraph{R-30 and the G-25 endpoint extension. Compact forms and auxiliary flags.}
At a block of rank $m$, the nonconstant torus character defines
a nontrivial partial flag. Its stabilizer intersects the Fourier
basis torus only in scalars: the corresponding circulant matrix
cannot preserve a proper consecutive coordinate subspace unless
every nonconstant cyclic coefficient vanishes. This produces a
stable ordered tensor even though its polynomial product alone
has a positive-dimensional complex stabilizer. After normalization,
the product has only a finite compact projective stabilizer.
The tilted padding line subtracts a positive rank-one Hermitian
operator from the desired moment matrix, with two explicitly
chosen rational positive eigenvalues. Both product blocks then
have nonconstant spectrum.

The full ordered padded tensor is the point at which the moment
vanishes. Losev's real-form stabilizer calculation applies there
and makes its complex stabilizer finite. The two-product parameter
scheme also has a nonempty open with finite padded-form stabilizer.
A generic section of the nontrivial partial flag orbit avoids
every fixed locus of that finite projective group. The resulting
enhanced point has no scalar period. This is a characteristic-zero
stable-quotient construction using D-S3 group stabilizers and
cocharacters, D-S6 tensor/Schur operations and D-FGA flag and
coordinate schemes; the Hermitian normalization is not transported
as an integral arithmetic algorithm. The exact G-25 quotient
comparison consequently extends to $r\le2m$ with unchanged
explicit constants.

\paragraph{G-26. A Pieri face of the bidegree pullback.}
The padding parameter morphism pulls a degree-$d$ coordinate
section back to bidegree $(ad,d)$ on the line and the permanent
cone. Its injectivity comes from dense image. In the representation
Grothendieck group its Schur components are governed by
$[\operatorname{Sym}^{ad}V^*]\,[S_\alpha V^*]$ and the
Pieri identity. The direct-sum multiplicity comparison is taken
over $\mathbb C$, where these representations are semisimple;
no integral direct-sum decomposition is inferred from a $K_0$
identity. When $\lambda_1=ad$, the horizontal strip must use
every column, so the unique surviving label is the tail of
$\lambda$. The zero-last-character coefficient subalgebra
identifies that tail multiplicity in one fewer variable.
These tensor and Schur identities belong to D-S6, with the exact
Pieri source locator retained in G-B20.

Two different enhanced invariant schemes now provide the
comparison: the determinant scheme in $r$ variables and the
tail permanent scheme in $r-1$ variables, or its full ternary
ambient bound at $r=4$. Their separate stable points have trivial
scalar period. Deleting the first flag part removes a nonnegative
number of strict weight comparisons, and the exact quotient
dimensions satisfy $u>v$ throughout the equality face with
$4\le r\le m$. The standard Veronese and the two separate
degree semigroups give strict comparison in every sufficiently
large original degree. D-FGA and EGA~II, 2.5.2, locate these
graded coordinate algebras and projectivizations. The argument
uses an injection and an upper bound, not a surjective lifting
of lower-dimensional sections to padded forms.

\paragraph{P-18. Three block characters and adjacent contraction degrees.}
The blockwise homogeneous coefficient tensor is graded by
three independent split tori. The three summands of its common
translation derivative have different block characters, hence
each is zero in every residue characteristic. At an extreme
block cut, the last-variable coefficient and the constant-prefix
coefficient occupy different degree blocks; their ranks add.
This forces two-dimensional extreme support bundles. The
prime rational-evaluation kernel makes every nonzero middle
coefficient map between those bundles invertible, and an
interior cut has two disjoint rank-two blocks. It cannot lie
in the width-three determinantal locus.

For a smallest unrestricted cubic, the exterior decoder of P-17
lives in $\bigwedge^3K$ over the node scheme. Decomposition by
left exterior degree makes the three contraction operators
connect only adjacent degree blocks. A zero block therefore
separates a smaller kernel element. Minimality and the P-17
rank lemma leave only adjacent occupied degrees at every
prefix; the ordered support is consequently trilinear on
three intervals with at most two shared boundary variables.
The degree-three interpolation grid has $20(n+1)^2$ points.
This dictionary uses the torus characters in D-S3, exterior
operations in D-S6 and determinantal charts in D-FGA. It neither
changes the variable order by a general linear coordinate map
nor asserts that setting a shared boundary variable to zero
preserves the balanced ideal.

\paragraph{I-42. Rank-one reset blocks and the special fiber at two.}
On the compatible rank charts of I-41, the mixed second
coefficients factor as $a_i b_j$ inside each reset block and
vanish across a strict internal reset. These are polynomial
rank-one equations; they retain more information than a
bandwidth condition alone. The three tangent sections
$1,h,h^2$ give a split rank-three evaluation map on any three
distinct nodes. The first nonzero row and the first future
nonzero coefficient reduce a hypothetical relation to at most
three Vandermonde columns when the reset distance is at most
five. At distance six the remaining relation yields exactly
the factor $2a_kb_\ell$. This also vanishes only trivially on
the open base where two is invertible.

The characteristic-two fiber is genuinely different: two
four-column kernels overlapping in one node give the sharp
seven-node abstract counterexample. Disjoint four-column
kernels give the eight-node example in every characteristic.
Neither is asserted to be the jet of an actual annihilating
ROABP. Use mixed coefficients, rather than division by two,
for this second-order statement. The fixed-rank and kernel
schemes belong to D-FGA, while the split characters and
divided-square anchors of saturated runs use D-S3 and D-S6
as in I-40/I-41. A saturated run exhausts the determinant
character budget, so the same anchor square coefficient
preserves its one-pole rank regardless of its length. The
sharp reset bounds then prove the stated irreducible classes
at weights three and four, and weight five away from two;
shorter nonsaturated runs retain their explicit rank gate.


\paragraph{P-19. A three-coordinate transversal and an injective exterior projection.}
Over the distinct-node base, the rank-three Vandermonde map
defines a kernel bundle $K$. Every coordinate subbundle on at
most three nodes meets $K$ trivially; projection away from such
nodes is consequently injective on $K$ and on its third exterior
power. The P-18 ordered-support reduction and two disjoint
rank-two coefficient blocks show that a smallest width-three
cubic counterexample would have all its monomial supports
meeting one set of at most three coordinates. Its decoded
section of $\bigwedge^3 K$ would be killed by that injective
projection. This contradiction works on every field fiber of
the distinct-node base. The proof uses determinantal coefficient
charts from D-FGA, split-torus degree blocks from D-S3, and the
exterior operations of D-S6. It closes both boundary cases of
P-18, with $20(n+1)^2$ evaluation points, and leaves degrees
four through $n-4$ as the remaining balanced width-three gate.

\paragraph{R-31. A residual scalar character on a common parameter scheme.}
On the open parameter scheme where the $2m+1$ product and
padding factors are independent, a compensated transposition
fixes the polynomial and its auxiliary covector. Its determinant
character is $-1$. After lifting to the special linear group,
the enhanced tensor has character $(-1)^b$. Density forces
every odd graded invariant to vanish when $b$ is odd. The
coordinate ring and its Cartan extraction are the D-FGA and
D-S6 constructions of G-22; the stabilizer character belongs
to D-S3. Over $\mathbb C$, closed-orbit separation then rules
out a nonzero closed point with trivial scalar period on this
common enhanced subcone. It does not impose the same character
on the larger permanent or determinant schemes. The obstruction
therefore narrows the next stable-point construction rather
than supplying an equation or multiplicity comparison for the
full cones.


\paragraph{R-32 and G-27. Unpadded stable quotients and a one-sided degree semigroup.}
The same two-product parameter morphism, without a padding
line, meets every prescribed positive moment direction when
each factor block has rank at most $d-1$. Its enhanced tensor
is stable. A nonzero flag weight allows a generic flag section
to remove the remaining finite projective stabilizer. The
central character then vanishes on the special linear center.
For a rectangular label the flag is a point, so only stability
and the quotient dimension are retained. This is the D-S3
stabilizer and cocharacter construction with the D-S6 Schur
line and D-FGA enhanced coordinate scheme; no integral
moment-map assertion is made.

G-27 uses two such schemes, in ranks $r$ and $r-1$. The upper
comparison needs no degree-period-one property on the lower
scheme. If $g$ is the gcd of its homogeneous generator
degrees, every sufficiently large multiple of $g$ is present.
Whenever its degree-$j$ piece is nonzero, a nonzero multiplier
places that piece inside a suitable degree of the standard
Veronese. The determinant algebra, used for a lower bound,
does have period one. Thus the projective Hilbert comparison
works at every original degree without identifying the two
semigroups. These are the graded-domain and Proj constructions
of D-FGA and EGA~II, 2.5.2. Together with the Pieri tail
dictionary of G-26, they classify every fixed nonrectangular
direction through $2m-1$ rows. The finite degree thresholds
are effective bounds, not polynomial-size algorithms.

\paragraph{I-43. A Pl\"ucker line constrains actual adjacent support bundles.}
A determinant-weight-one line in the third exterior power
generates a standard rank-two $\mathrm{SL}_2$ module. Its
projectivization lies in the Grassmannian, whose line
description gives an invariant rank-two subspace inside an
invariant rank-four space. The exact known Grassmann input
is I-B10, Lemma~3, p.~4. The actual preceding coefficient
support has rank at most three. Its split-torus characters
and the physical unipotent action exclude a larger pair
of opposite weights in that invariant two-plane, leaving
weights $(1,1,-1)$ on the rank-three support.

No splitting of these module extensions is needed in a
positive-characteristic fiber. The normalized core updates
have a common covector, so their second exterior map is
affine in the internal variables. A single node pole detects
each nonzero coefficient while the two anchor evaluations
remain independent. The determinant-weight-zero case uses
one divided-square anchor and an arbitrary opposite affine
anchor. These identities hold on the actual coefficient
rank charts over the same rational parameter field. The
preserved two-dimensional contractions compose, including
the endpoint projections, before applying the I-42 reset
argument. D-S3 locates the group scheme and characters,
D-S6 the exterior/divided operations, and D-FGA the
Grassmann and determinantal charts. The result is valid on
every field fiber with distinct nodes: all weights at most
three and, using P-19, every multilinear input on at most
eleven variables are hit by the balanced two-pair sampler.


\paragraph{R-33 and G-28. A generic Pl\"ucker chart and a quotient upper bound.}
The maximal differential rank of the permanent coefficient
morphism identifies its generic tangent kernel with the
row-column diagonal gauge directions. On the Grassmann
open where every Pl\"ucker coordinate of the coefficient
matrix span is nonzero, any diagonal operator preserving
that span has equal sums on every fixed-size coordinate
subset. Exchanging one index forces all individual weights
equal. The zero gauge character makes the scalar zero.
Thus the generic projective stabilizer has zero Lie algebra
and is finite over $\mathbb C$.

These are the determinantal and Grassmann opens of D-FGA
and the torus character calculation of D-S3. Stabilizer
openness retains the precise SGA~3 VI$_B$, Proposition~4.1
and cited EGA~IV$_3$, Theorem~13.1.3 framework. The generic
fiber dimension equality is separately cited from the
Stacks Project, Lemma~33.20.4(2), tag~0B2L (R-B18), whose
complete primary proof was read. A fiber of the invariant
quotient contains the orbit through its generic point.
This gives only an upper bound on quotient dimension;
it asserts neither stability nor equality.

G-28 uses exactly that upper bound on the lower permanent
algebra. Its actual dimension may be smaller, or the ring
may contain only constants. The projective image degree and
Hilbert upper estimates decrease with that dimension, so
the G-27 comparison still holds. This closes all fixed
nonrectangular directions through $2m$ rows and extends
the boundary-only comparison to
$r\le\min\{2n-2,m^2\}$. The Schur tail, Proj and degree
semigroup constructions retain the D-S6/D-FGA interpretation
already given; generic freeness is not substituted for
the stable-point hypothesis needed for the upper determinant
algebra's period one.

\paragraph{I-44. Torus chains, anchor charts, and a residue-support condition.}
The normalized cores move between integer torus characters
differing by two. Unless the rank-three support is a chain
of three successive such characters, their updates have
a common image or covector, and the second exterior map
is affine. In the chain case, the actual preceding support
allows four weights. Minimal coefficient independence and
rank two leave precisely the five anchor charts displayed
in I-44. These statements use the D-S3 character grading
and D-FGA rank charts and hold on every characteristic fiber.

Two anchor types have a unique nonzero pole coefficient.
The third, paired with the actual reverse anchor, reduces
the second exterior map to an explicit quadratic. A square
coefficient gives a double pole. In the multilinear case,+the coefficient of $tz$ and the simple residues impose two
independent equations on each row. Every nonzero row would
therefore need at least three neighbors. With at most five
internal positions, one of the two coefficient types has
at most two entries, forcing an order that makes the
quadratic factor into nonzero linear forms. The evaluation
kernel cannot contain that product.

The resulting interval and endpoint assertions use a single
parameter field before invoking the reset-Hessian theorem.
The exterior operations belong to D-S6, without division
by two. The characteristic-two distinction occurs only in
the subsequent I-42 scalar Hessian step: four-core runs
are covered on every fiber, and five-core runs where two
is invertible. Factor weights then give the complete balanced
theorem through twelve variables in every characteristic
and thirteen away from two, with the stated finite grid.

\subsection{What the manuscript record does and does not establish}
The supplied cote~133 study also reports decisions about the name on scan~008, a characteristic-two pre-$\lambda$ issue on scan~006, and a crossed-out Stiefel--Whitney assertion on scan~009. Those reports are retained as reports of earlier readers, with their own scope. No new visual decision about the name or the deleted text is made in this run. The five attached LaTeX transcriptions contain no original raster pages. The separate transcription table records the forced algebraic readings and the distinction between a textual ambiguity, a mathematical correction, and a claim whose ink cannot be checked here.

\clearpage\part{Unchecked targets and unresolved gates}
\section{Unchecked targets and the current smaller gates}

This part contains no asserted proof of the unresolved targets. A conditional reduction in the checked part means that the implication was proved; it does not mean its antecedent was established. The live question register and dated ledger give the current states, including additional deductions made after this text was first assembled.

\paragraph{U-01: general coefficient extraction and determinant-boundary ranks.}
The exact rectangular Haar transforms, ratios kernel, positive finite-group projectors and Gram-rank identities are established in R-01--R-09. The remaining general gate is a uniformly efficient positive/combinatorial realization of the multiplicity spaces and effective control of the orbit-map ranks or quotient numerator. R-27 closes the quaternary image-dimension and graded-defect gate at every size. G-20 therefore determines the eventual sign of every fixed strict four-row direction, with explicit bounds and no integer dimension transition left undecided. G-21 gives maximal permanent-image parameter dimension in every ambient variable count at least four for m at least seventeen. G-22, R-28/R-29 and G-23/G-25 now determine the eventual sign in every strict direction with $4\le r\le2m$ for $m\ge17$, with both occurrences at all sufficiently large integer stretches and no integer equality case. More than $2m$ rows, nonrectangular boundary faces beyond $r\le\min\{2n-2,m^2\}$, individual finite degrees, changing directions, and the superlinear-padding GCT regime remain unresolved. The exact eventual rectangular range starts uniformly at size fourteen; this sharp eventual threshold is not a theorem excluding every earlier finite-degree gap. G-19 supplies one unpadded discriminant copy but no total multiplicity comparison and no positive-padding separator. R-21 remains an exact fixed-ambient computation with an exponent growing with the ambient dimension. The general nonstandard quantum-group conjectures and algorithmic requirements retain their separately stated hypotheses.

\paragraph{U-02: a polynomial ROABP hitting set.}
I-30 closes the width-at-most-three, one-unknown-order case with six interpolation pairs, including arbitrarily many rank-two cuts. I-36 improves this to five pairs for the class with singular variable-dependent factors, through one-pair preservation of all nonzero maximal interval minors. I-29 proves a polynomial pair budget when there is one allowed singular rank. I-31 proves a bound controlled by the number of allowed rank levels and polynomial-size sampling for every fixed width. Its recurrence is still $2^{O(w\log w)}$ for all ranks, and sampling uses the pair count as an exponent. I-32 proves that polynomial exact coordinate-cone compression at every preserved common center is impossible. I-34 extends the prime-cocharacter repair to separated products with bounded local nullity and proves its exact one-coordinate limitation. P-14 replaces full image-cone preservation by scalar bridges for factors of rank one or full rank, reducing that case to $w-1$ pairs. General intermediate-rank matrix cancellation remains outside these divisor and product arguments. P-15 refutes a node-uniform one-pair generator already at width three with invertible factors. P-16 refutes the uncentered two-pair map on eighteen specified algebraic nodes, even for generically invertible factors, while P-14 proves that centering repairs that example. I-37 closes centered two-pair isolation for affine singular-factor width-three words. I-38 additionally permits isolated invertible affine factors. I-40 supplies a determinant-character run bound for homogeneous translation-invariant tensors and closes weights zero and one through I-37/I-38. Together with P-17 it proves the balanced two-pair theorem through nine variables, including an explicit grid. I-41 closes the weight-two class and the complete balanced problem through ten variables by a nonzero connector square coefficient and a bandwidth-three argument. I-43 closes weight three and the complete balanced theorem through eleven variables. I-44 further closes the balanced theorem through twelve variables in every characteristic and thirteen away from two. The next unrestricted small cases are thirteen variables in characteristic two and fourteen otherwise. Longer runs in unrestricted words, nonlinear singular factors, and the growing-width question remain open. The remaining target is a hitting set whose size and construction time are polynomial jointly in length, degree and growing width. Explicit joint coefficient-span preservation in I-14 remains a sufficient gate; I-03--I-04 give the stronger cone-hitting alternative. The direct Hoza transfer must respect the monomial quotient and coefficient cancellation. The independent-variable theorem, rank-profile specialization, and arithmetic encoding do not yet supply the required uniform width dependence. I-35 gives the strictly weaker rank-one-annihilator criterion and proves composition with independent seeds; the remaining bridge gate is an efficient common substitution preserving the actual restricted seed-dependent boundary functions. Arbitrary rational-boundary preservation and the stated rank-one multiplicative scalarization are refuted. The attempts register records the exact constructions tested and their failure points.

\paragraph{U-03: matrix invariants and GCT~V.}
Once U-02 is solved with the stated field and construction bounds, I-08 gives the polynomial-size separating family and finite matrix-quotient morphism. The implication is complete; the polynomial hitting-set antecedent is not. General explicit Noether normalization and general PIT have stronger hypotheses than this matrix-specific consequence. No circuit lower bound or $P$ versus $NP$ separation is asserted here.

\paragraph{U-04: a geometric or arithmetic encoding with polynomial budgets.}
The sparse and actual-sheaf constructions are proved. For arbitrary ROABPs the remaining task is to construct a sampler whose evaluation kernel avoids the bounded-width tensor variety, without requiring near-isometry of the entire coefficient norm. P-10 makes that kernel criterion exact. P-11 removes a detour: a sufficiently explicit common algebraic point is itself a hitting-set encoding, independently of GRH. P-13 constructs that point and an integer curve for every class hit by the stated centered interpolation map, with degree $(nD+1)^{2K+1}$. This gives polynomial arithmetic budgets for width three and for each fixed width through I-31. The uniform polynomial-in-width arithmetic construction remains open. A Frobenius construction would have to supply new nonvanishing structure; applying Chebotarev after assuming it does not solve the original problem.

\paragraph{U-05: unrestricted number-field random-matrix laws.}
The dependency table gives the supported statements now discharged under the stipulated inputs. Full pair correlation, spacing laws, the general arithmetic moment recipe, arbitrary automorphic target RH, and a zeta-zero quantum Hamiltonian retain their explicitly named analytic or spectral prerequisites. No literature theorem is strengthened beyond its stated support or uniformity range by changing an RH label.

\paragraph{Next work and stopping discipline.}
The attempts files record the constructions actually tried and their failure locations. The next substantive work is on the remaining joint-compression and special-fiber gates and on the restricted projected-orbit ideals. Packaging and repeated exposition are not counted as movement. Checkpoint and review times are actual elapsed times recorded in the ledger; no scheduled review is claimed before it has occurred.


\subsection{Updated gates after the generic-stabilizer and chain-anchor arguments}
R-27 closes the uniform quaternary dimension and specified Tor gate.
R-29/R-30 and G-25 give the strict-direction sign through
$2m$ rows. G-26 classifies every fixed direction through
$m$ rows. R-32/R-33 and G-27/G-28 classify every fixed
nonrectangular direction through $2m$ rows, with a further
boundary-only range $r\le\min\{2n-2,m^2\}$.
Rectangular directions beyond $m$ rows, strict directions
past $2m$ rows, changing directions and efficient finite-degree
extraction remain open. R-31 excludes the unchanged two-product
no-scaling strategy in specified $2m+1$-row directions;
it imposes no vanishing on the full cones.

P-19 proves the cubic theorem in every variable count and
characteristic. I-44 preserves all one-pole interval ranks
through five full-rank cores by its complete torus-chain
anchor calculation. The I-42 scalar step then gives weights
at most four in every characteristic and at most five away
from two. The complete balanced two-pair theorem holds
through twelve variables, and thirteen away from two.
The characteristic-two weight-five gate consists of possible
saturated five-core runs whose ranks are now preserved but
whose reset Hessian needs an additional constraint.
At longer lengths the remaining two-anchor quadratic is
explicit; both internal coefficient types can have at least
three positions. Nonlinear singular factors, general
shared-seed synchronization and uniform polynomial dependence
on growing width remain open.

\appendix
\clearpage\part{References, exact locators, and audit scope}
\section{Rectangular bridge: sources and exact
locators}\label{rectangular-bridge-sources-and-exact-locators}

All new R statements are deductions made and checked in this run; no
claim of priority over the literature is made. The older numerical
formulas are explicitly credited. Attached LaTeX was read directly, not
inferred from the prompt summary.

\subsection{R-B01 --- Bump--Gamburd rectangular product
identity}\label{r-b01--bumpgamburd-rectangular-product-identity}

Bump, D., \& Gamburd, A. (2006). On the averages of characteristic
polynomials from classical groups. \emph{Communications in Mathematical
Physics, 265}, 227--274. \url{https://doi.org/10.1007/s00220-006-1503-1}

Exact locator: §3.1, Proposition 4, equation source label
\nolinkurl{e:produ}; theorem label \nolinkurl{cpaverage}. Attached source
\nolinkurl{LIT_Bump_Gamburd_averages_of_characteristic_polynomials/acp20.tex},
lines 994--1039. General parameters are \nolinkurl{N,L,K}; the polynomial
is \texttt{s\_(N\^{}L)} in \nolinkurl{K+L} variables. This supports R-01.

\subsection{R-B02 --- Bump--Gamburd swaps and
ratios}\label{r-b02--bumpgamburd-swaps-and-ratios}

Same work. Exact locators: §2.3, Lemma 1 (\nolinkurl{lem:min1}), §3.1,
Proposition 4 (\nolinkurl{cpaverage}), and §4 (\nolinkurl{comf}) for
rectangular Weyl coset/Laplace sums. §3.3, Theorem 3 (\nolinkurl{t:1}),
equation \nolinkurl{unitaryratio} for inverse characteristic-polynomial
ratios, with \nolinkurl{N≥Q,R} and the denominator parameters of modulus
less than one. The specialization \nolinkurl{K=L=0,Q=R=n²} used in R-09 is
within those hypotheses when \nolinkurl{N≥n²}. Attached source lines
1539--1563 and 1650--1735.

\subsection{R-B03 --- Finite CUE moments, hook content, Schur
orthogonality}\label{r-b03--finite-cue-moments-hook-content-schur-orthogonality}

Same work. §3.1, Corollary 1 (\nolinkurl{ratioskeatingsnaith}), equation
\nolinkurl{ksformula}; §2.2, equations \nolinkurl{ratiosschurorthog} and
\nolinkurl{e_hhl}. Finite moment equals the rectangular representation
dimension; the leading asymptotic coefficient need not be an integer.

\subsection{R-B04 --- Arithmetic CFKRS status and swap
formula}\label{r-b04--arithmetic-cfkrs-status-and-swap-formula}

Conrey, J. B., Farmer, D. W., Keating, J. P., Rubinstein, M. O., \&
Snaith, N. C. (2005). Integral moments of L-functions. \emph{Proceedings
of the London Mathematical Society, 91}(1), 33--104.
\url{https://doi.org/10.1112/S0024611504015175}

Version matching attachment: arXiv:math/0206018v3 (2005),
\url{https://arxiv.org/pdf/math/0206018} . Exact locators: Conjecture
1.5.1 (\nolinkurl{thm:zeta2kconjecture}), Theorem 1.5.2
(\nolinkurl{thm:UN2k}), §2.1 recipe, §2.2 equation \nolinkurl{conj:M}, Lemma
2.5.1 (\nolinkurl{thm:concisesumunitary}). The author\textquotesingle s
older PDF at Waterloo is v1 and uses different section numbers; do not
cite its §3.2 as if it matched the attached §2.5. The statements are
checked as conjecture versus theorem, not assumed proved from RH.

\subsection{R-B05 --- Schur--Weyl and exact
projectors}\label{r-b05--schurweyl-and-exact-projectors}

Collins, B., \& Śniady, P. (2006). Integration with respect to the Haar
measure on unitary, orthogonal and symplectic group.
\emph{Communications in Mathematical Physics, 264}(3), 773--795.
\url{https://doi.org/10.1007/s00220-006-1554-3}

Primary source: \url{https://arxiv.org/pdf/math-ph/0402073} . Exact
locators: Theorem 2.1 (Schur--Weyl duality), equation (3) and the
central idempotents preceding equation (4), Proposition 2.2 (Haar
conditional expectation as orthogonal projection).

\subsection{R-B06 --- Weingarten calculus with small-dimension
correction}\label{r-b06--weingarten-calculus-with-small-dimension-correction}

Same work: Proposition 2.3, equation (9), includes the necessary
\nolinkurl{ell(lambda)≤dimension} cutoff; Corollary 2.4, equation (11), is
the entry integration formula. The example after Proposition 2.5 gives
the negative transposition value \texttt{−1/{[}N(N²−1){]}}. No
unrestricted stable Weingarten inverse is substituted at singular small
dimensions.

\subsection{R-B07 --- Rains theorem}\label{r-b07--rains-theorem}

Rains, E. M. (2003). Images of eigenvalue distributions under power
maps. \emph{Probability Theory and Related Fields, 125}(4), 522--538.
\url{https://doi.org/10.1007/s00440-002-0250-2}

Primary source: \url{https://arxiv.org/pdf/math/0008079} . Exact
locators: Theorem 1.3, equation (8), and Corollary 1.4. They describe
one power\textquotesingle s eigenvalue law. They do not assert
independence of different powers of the same matrix.

\subsection{R-B08 --- GCT multiplicity criterion and rectangular
obstruction
limit}\label{r-b08--gct-multiplicity-criterion-and-rectangular-obstruction-limit}

Ikenmeyer, C., \& Panova, G. (2017). Rectangular Kronecker coefficients
and plethysms in geometric complexity theory. \emph{Advances in
Mathematics, 319}, 40--66.
\url{https://doi.org/10.1016/j.aim.2017.08.024}

Primary source: \url{https://arxiv.org/pdf/1512.03798} (v2); exact
locators: Theorem 1.4 (\nolinkurl{thm:strategy}), equations (1.3) and
(1.5), Theorem 1.6 (\nolinkurl{thm:together}), and Theorem 2.1
(\nolinkurl{thm:manivel_rect}). Theorem 1.6 requires
\texttt{n\textgreater{}3m\^{}4}. Source hard-coded cross-references
occasionally refer to earlier numbering, so numerical locators were
confirmed against the rendered primary PDF.

\subsection{R-B09 --- Orbit versus closure and no occurrence
obstructions}\label{r-b09--orbit-versus-closure-and-no-occurrence-obstructions}

Bürgisser, P., Ikenmeyer, C., \& Panova, G. (2019). No occurrence
obstructions in geometric complexity theory. \emph{Journal of the
American Mathematical Society, 32}(1), 163--193.
\url{https://doi.org/10.1090/jams/908}

Primary source matching attachment:
\url{https://arxiv.org/pdf/1604.06431v3} . Exact locators: Theorem 1.4
(\nolinkurl{th:main}), \texttt{n≥m\^{}25}; §1(b), paragraphs defining the
GCT coefficients and bounding them by the symmetric rectangular
Kronecker coefficient of the orbit, then by the ordinary rectangular
coefficient. The old v1 theorem number 1.5 must not be silently used for
the supplied v3 source.

\subsection{R-B10 --- Known stable numerical
positivity}\label{r-b10--known-stable-numerical-positivity}

Manivel, L. (2011). On rectangular Kronecker coefficients. \emph{Journal
of Algebraic Combinatorics, 33}, 153--162.
\url{https://arxiv.org/abs/0907.3351}

Primary PDF: \url{https://arxiv.org/pdf/0907.3351} . Exact locators:
Theorem 1(2); §3.3, Corollary 2(1), which explicitly states that for
\nolinkurl{d≥q}, \texttt{g((nd−q,q),(d\^{}n),(d\^{}n))} is the number of
partitions of \nolinkurl{q} into parts between 2 and n. §3.4 gives a
birational flag-quotient interpretation. R-07 credits this numerical
identity and proves explicit centered-polynomial witnesses and their
orbit-closure survival, rather than claiming a new count.

\subsection{R-B11 --- GCT IV: proved two-row rule, its exact rectangular
specialization, and
limits}\label{r-b11--gct-iv-proved-two-row-rule-its-exact-rectangular-specialization-and-limits}

Blasiak, J., Mulmuley, K. D., \& Sohoni, M. (2015). \emph{Geometric
complexity theory IV: Nonstandard quantum group for the Kronecker
problem}. Memoirs of the American Mathematical Society, 235(1109).
\url{https://doi.org/10.1090/memo/1109}

The supplied file is the later three-author
\nolinkurl{Apr13KroneckerGCT4.tex}, not the original two-author March 2007
version cited by the older GCT VII/VIII manuscripts. Matching author
PDF: \url{https://websites.umich.edu/~jblasiak/Apr13KroneckerGCT4.pdf} .
Page and theorem numbers below are for that Apr13 manuscript.

\begin{itemize}
\tightlist
\item
  Exact attached file:
  \nolinkurl{sources/LIT_Mulmuley_GCT_IV_nonstandard_quantum_group_Kronecker/Apr13KroneckerGCT4.tex}.
\item
  §12.1, theorem source label
  \texttt{t\ nonstandard\ schur-weyl\ duality}, source lines 4236--4290:
  nonstandard Schur--Weyl duality proved for general dimensions. This
  established structural theorem should be distinguished from the later
  two-row limitation on explicit canonical bases.
\item
  Theorem 15.21, especially (vii)--(viii), manuscript p.102, label
  \texttt{t\ main\ canonical\ basis}, source lines 5963--5978:
  highest-weight counting rule for two-row Kronecker coefficients and
  its symmetric/exterior refinement. The two-row restriction means both
  \emph{inner} GL partitions have length at most two; the outer
  partition may have four rows.
\item
  Theorem 18.2, manuscript pp.124--125, label
  \texttt{t\ consequence\ for\ Kronecker}; equation (163), label
  \texttt{e\ g\ =\ invariant-free\ g}: constructive positive two-row
  rule from five types of invariant-free components, each zero or one.
\item
  Corollary 18.4, manuscript p.127, label
  \texttt{c\ kronecker\ generating\ function\ sym\ ext}, source lines
  7736 onward; introduction equivalent equation
  \texttt{e\ intro\ Kronecker\ generating\ function}, lines 1435--1455:
  symmetric/exterior generating functions. Taking the constant
  coefficient gives the all-even/four-odd-part formula recorded as R-10.
  The numerical formula is known; our work specializes and compares it
  with the closure ring.
\item
  Example 17.8, manuscript p.122, label \texttt{ex\ not\ positive},
  source lines 7326--7351: a divided-power matrix coefficient has mixed
  Laurent signs, namely \texttt{q\^{}10-q\^{}4-q\^{}-4+q\^{}-10}. The
  source first states that no global sign readjustment makes every
  specified Chevalley matrix coefficient nonnegative. The displayed
  algebraic identity was separately recomputed in
  \nolinkurl{check_n2_formulas.py}.
\item
  Conjecture 19.1, §19.1, label \texttt{cj\ canonical\ basis\ X\^{}r},
  lines 7920--7967: proposed compatible basis for the full tensor power.
  §19.2, source lines 8191--8209: naive extension can have too small a
  dimension; it has not supplied the general positive rule. Intro §1.1,
  source lines 893--903, explicitly describes the unresolved cases.
\item
  R-10 gives a direct principal-minor highest-weight proof of the
  classical identity
  \texttt{Sym\^{}d(Sym\^{}2\ C\^{}s)=direct\_sum\ S\_(2mu)\ C\^{}s}, so
  no unverified book exercise number is used for that identity.
\end{itemize}

\subsection{R-B12 --- GCT VII: proved structure versus conjectural
plethysm
reciprocity}\label{r-b12--gct-vii-proved-structure-versus-conjectural-plethysm-reciprocity}

Mulmuley, K. D. (2008a). \emph{Geometric complexity theory VII:
Nonstandard quantum group for the plethysm problem} (revised extended
abstract; arXiv:0709.0749v2; original Technical Report TR-2007-14,
2007). \url{https://arxiv.org/abs/0709.0749v2}

\begin{itemize}
\tightlist
\item
  Exact attached file:
  \nolinkurl{sources/LIT_Mulmuley_GCT_VII_nonstandard_quantum_group_plethysm/main.tex}.
\item
  Introduction, source lines 212--229, Theorem 1.2, label
  \nolinkurl{tnonstd}: existence of the nonstandard quantization, complete
  reducibility, and quantum Peter--Weyl theorem.
\item
  Introduction, lines 257--285, Conjecture 1.3, label \nolinkurl{cduality}:
  nonstandard duality. Lines 290--305, Conjecture 1.4, label
  \nolinkurl{creciprocity}: nonstandard reciprocity. §6, Conjecture 6.1,
  label \nolinkurl{crecirefined}: refined specialization statement.
\item
  Introduction, lines 309--335: conditional positive decomposition into
  \texttt{m\_lambda\^{}alpha\ n\_pi\^{}alpha} and the two subproblems
  P1/P2. Lines 336--366 explicitly state that algorithm correctness and
  the required canonical-basis properties are not proved in this source.
\item
  Abstract and introduction distinguish this proposed positive plethysm
  formula from already established existence and complete reducibility.
  Latest attached arXiv version is v2, September 1, 2008; the first
  version is September 5, 2007.
\end{itemize}

\subsection{R-B13 --- GCT VIII: precise conjectural positivity and
finite-field RH
language}\label{r-b13--gct-viii-precise-conjectural-positivity-and-finite-field-rh-language}

Mulmuley, K. D. (2008b). \emph{Geometric complexity theory VIII: On
canonical bases for the nonstandard quantum groups} (revised extended
abstract; arXiv:0709.0751v2; original Technical Report TR-2007-15,
2007). \url{https://arxiv.org/abs/0709.0751v2}

\begin{itemize}
\tightlist
\item
  Exact attached file:
  \nolinkurl{sources/LIT_Mulmuley_GCT_VIII_canonical_bases_nonstandard_quantum_groups/main.tex}.
  Primary HTML for this version:
  \url{https://arxiv.org/html/0709.0751v2} .
\item
  Abstract, lines 144--159, and Introduction, lines 195--218: algorithms
  are explicitly called conjecturally correct; required properties are
  also conjectural.
\item
  §2.2, Conjecture 2.2, label \nolinkurl{clocalcrystalb}, lines 634--655:
  local crystal basis existence.
\item
  §2.3, Conjecture 2.5, label \nolinkurl{cuniquemindegleft}, lines
  947--952: unique minimum degree complexity required for globalization.
\item
  §2.4.1, Conjecture 2.6, label \nolinkurl{ccelldecompleft}, lines
  1041--1046: cellular decomposition compatible with the restricted
  standard quantum group.
\item
  §2.4.2, Conjecture 2.7, label \nolinkurl{cposkroneckerleft}, lines
  1073--1090: Kronecker positivity, with the precisely stated sign and
  \texttt{(q-q\^{}-1)\^{}a} factors allowed for structure constants. The
  following paragraph defines positive coefficients to mean nonnegative
  \textbf{rational} coefficients.
\item
  §2.4.3, Conjectures 2.8--2.9, labels \nolinkurl{cnonstdposleft0} and
  \nolinkurl{cnonstdposleft1}, lines 1151--1214: weaker general saturation
  and nonstandard almost-positivity claims; not an established
  nonnegative integer tableau count.
\item
  §3.3, Conjectures 3.1--3.2, labels \nolinkurl{cunimindegright},
  \nolinkurl{cunimindegright2}; §3.4.1, Conjecture 3.3, label
  \nolinkurl{ccelldecompright}; §3.5, Conjecture 3.4, label
  \nolinkurl{cposkroneckerright}; §3.5.1, Conjectures 3.5--3.6, labels
  \nolinkurl{cnonstdposright}, \nolinkurl{cnonstdposright1}; §3.5.2,
  Conjecture 3.7, label \nolinkurl{coppquasicell}.
\item
  §6, label \nolinkurl{scomplexity}, lines 2810--2854: efficient
  combinatorial labels and crystal operations needed for a \#P rule are
  additional proposed properties, with details deferred to a full
  version.
\item
  Introduction, lines 238--247, and §2.4.2, lines 1118--1126:
  perverse-sheaf theory and the finite-field Riemann-hypothesis
  machinery are invoked as a prospective extension needed for nonflat,
  noncommutative varieties. These passages supply no implication from
  the classical analytic RH or Dirichlet GRH to the stated algebraic
  conjectures.
\item
  Exact theorem/conjecture numbers were checked against the shared TeX
  theorem counter; principal positivity and complexity locators were
  also matched to the primary HTML. No full independent proof of these
  background quantum-group constructions is claimed.
\end{itemize}

\subsection{R-B14 --- Cubic surfaces and the small-degree four-row
exclusion}\label{r-b14--cubic-surfaces-and-the-small-degree-four-row-exclusion}

Beauville, A. (2000). Determinantal hypersurfaces. \emph{Michigan
Mathematical Journal, 48}, 39--64.
\url{https://math.univ-cotedazur.fr/u/beauvill/pubs/det.pdf}

\begin{itemize}
\tightlist
\item
  Corollary 6.4, printed pp.52--53 (PDF pages 14--15): over an
  \textbf{algebraically closed} field, every \textbf{smooth cubic
  surface} admits a \nolinkurl{3×3} linear determinantal equation. Its
  additional count of 72 representations is not needed for R-11.
\item
  Corollary 6.5, printed p.53, states the extra
  twisted-cubic/birationality/six-lines conditions over a field that
  need not be algebraically closed. R-11 works over C and does not
  assert the unrestricted statement over other fields.
\item
  §1.9, printed p.42, cross-references the plane-curve and cubic-surface
  statements and explains the dimensional boundary for generic
  hypersurfaces.
\item
  R-11 uses Corollary 6.4 only for smooth quaternary cubics, then proves
  density, block-diagonal padding, linear-substitution closure, and
  highest-weight ideal containment. It does not claim that every
  singular cubic itself is a linear determinant; membership in the
  closure is sufficient.
\end{itemize}

\subsection{R-B15 --- Dual Jacobi--Trudi identity used in the
simultaneous Fourier
proof}\label{r-b15--dual-jacobitrudi-identity-used-in-the-simultaneous-fourier-proof}

Sawin, W. (2020). A representation theory approach to integral moments
of L-functions over function fields. \emph{Algebra \& Number Theory,
14}(4), 867--906. \url{https://doi.org/10.2140/ant.2020.14.867}

\begin{itemize}
\tightlist
\item
  Exact attached file:
  \nolinkurl{sources/LIT_representation_theory_approach_to_moments_function_fields/main.tex}.
\item
  Lemma 2.3, source label \nolinkurl{representation-generating-identity};
  equation source label \nolinkurl{Jacobi-Trudi}, lines343--345. In
  arXiv:1810.01303v2 this is equation(2.4), printed p.12. The proof
  explicitly identifies the determinant of exterior powers with the
  conjugate-partition Schur character, citing Fulton--Harris Formula A6.
\item
  R-16 uses only this unconditional classical identity, with three equal
  column lengths j, to identify the 3x3 Toeplitz determinant with
  \texttt{s\_(3\^{}j)}. The new nonvanishing specialization is then
  proved directly from the Weyl alternant and the distinct exponents
  in0..k\textless m.
\item
  No cohomology-vanishing hypothesis or number-field/function-field
  transfer is invoked by R-16. The arithmetic parts of
  Sawin\textquotesingle s paper are audited separately by the
  root/purity agents.
\item
  Matching primary preprint: \url{https://arxiv.org/pdf/1810.01303v2} .
  Published metadata also checked against the journal\textquotesingle s
  primary PDF endpoint and author metadata.
\end{itemize}

\subsection{Cross-reference used in R-21:
G-B10}\label{cross-reference-used-in-r-21-g-b10}

Alpöge, L. (2026, November20). \emph{Remark} {[}Supplied note and
stipulated analytic input{]}, printedp4, step13. Local PDF:
\nolinkurl{sources/ALPOGE_remark.pdf}; extraction
\nolinkurl{work/root/ALPOGE_remark.txt}, line192. The step reads ``Bost
slope inequality''; the R21 scope audit concerns this step, without
re-auditing stipulated RH/Dirichlet-GRH. The complete bibliographic
source is registered by root as G-B10, not duplicated under a new R-B
identifier.

\subsection{R-22 source and derivation
map}\label{r-22-source-and-derivation-map}

R-22 is a proved in-run deduction, with no priority claim. Its new local
identities, finite parameter-grid selection, and uniform rank estimate
are proved completely in r22\_fragment.tex.

\begin{itemize}
\tightlist
\item
  R-15/R-17 provide the cyclic base differential dimension 2m\^{}2+2;
  R-16 gives its exact two-column character quotient. These internal
  dependencies are stated explicitly in R22 paragraph3.
\item
  R-B15 (Sawin2020, Lemma2.3, equation2.4/source label Jacobi-Trudi) is
  used only for the unconditional dual Jacobi--Trudi identity in R22
  equation R22:schur. Nonvanishing of the geometric specialization is
  proved directly by the Weyl alternant, rather than imported from an
  arithmetic theorem.
\item
  R-13 supplies the row--column torus upper bound 4m\^{}2-2m+2; R-14
  supplies the padding dimension c\_m+3; R-17 gives exact determinant
  dimension 2n\^{}2+2.
\item
  G-10, G-11/G-12, and G-17 supply the stated effective degree,
  rectangular boundary, and stretched-label consequences after R22
  proves their dimension hypothesis.
\item
  The primitive root is fixed before the parameter grid is selected. The
  motif may depend on m and q. No statement about a universal
  coefficient choice or a polynomial-time construction is taken from a
  source or inferred from finite tests.
\end{itemize}

\subsection{R-B16 --- ELSW: the actual linear strand of
subpermanents}\label{r-b16--elsw-the-actual-linear-strand-of-subpermanents}

Efremenko, K., Landsberg, J. M., Schenck, H., \& Weyman, J. (2018). On
minimal free resolutions of sub-permanents and other ideals arising in
complexity theory. Journal of Algebra, 503, 8--20.
\url{https://doi.org/10.1016/j.jalgebra.2018.01.021}

\begin{itemize}
\tightlist
\item
  Primary author PDF: \url{https://people.tamu.edu/~jml/ELSWmfr2-29.pdf}
  . This is the17-page author version; theorem/page locators below refer
  to it.
\item
  Theorem1.1, printedp2, gives the complete linear-strand dimensions.
  Section4.1, Example4.1, printedp11, gives the first linear syzygy
  dimension 2κ binom(n,κ+1)\^{}2 and the explicit row/column
  Laplace-difference basis.
\item
  Section4.1, Theorem4.3 and Lemma4.4, printedpp11--13, state and prove
  the full linear-strand description. For R23 only the first-syzygy case
  κ=m−1 is used; The standalone audit r23\_syzygy\_audit.tex checks that
  special case by monomial comparison; the manuscript cites the known
  result rather than reproving it.
\item
  The group description at the start of section4, printedp11, also
  identifies the permanent stabilizer, but R23 does not need to assume
  it.
\item
  Publication metadata verified against the first
  author\textquotesingle s institutional research record:
  \url{https://cris.bgu.ac.il/en/publications/on-minimal-free-resolutions-of-sub-permanents-and-other-ideals-ar-2/}
  .
\end{itemize}

\subsection{R-B17 --- BCMV: a codimension bound that blocks the
regular-sequence
shortcut}\label{r-b17--bcmv-a-codimension-bound-that-blocks-the-regular-sequence-shortcut}

Boralevi, A., Carlini, E., Michałek, M., \& Ventura, E. (2024). On the
codimension of permanental varieties (arXiv:2402.17839v2). arXiv.
\url{https://arxiv.org/pdf/2402.17839v2}

\begin{itemize}
\tightlist
\item
  Matching primary PDF retrieved at
  \url{https://arxiv.org/pdf/2402.17839} ; first page identifies
  v2,4December2024.
\item
  Theorem4.14, printedp14, records von zur Gathen\textquotesingle s
  lower bound5 for the codimension of the permanent singular locus for
  every m\textgreater=3, with upper bound2m.
\item
  Theorem4.23, printedp17, proves the improved lower bound6 for
  m\textgreater=6. R23 only needs the lower bound5 to show
  depth(R/J\_m)\textless m\^{}2−4.
\item
  The first derivative ideal defines this singular locus; its height is
  the codimension of the underlying variety even if the ideal is not
  radical. Neither cited codimension result asserts preservation of
  linear syzygies under a generic four-variable restriction.
\end{itemize}

\subsection{R-24: exact finite certificate at
size17}\label{r-24-exact-finite-certificate-at-size17}

This is an in-run finite computation and deduction, with no external
priority claim. The complete pencil, points, selected indices, and raw
matrix are in \nolinkurl{permanent_maximal/permanent_range_n17.json} and
\nolinkurl{.u8}; the compact omission list is also stated in the
manuscript. \nolinkurl{check_permanent_maximal.cpp} uses the R13 subset
evaluator, while the separately executed Gray-code Ryser evaluator
checked every entry. Root independently verified the selected minor by
modular Bareiss and the 32-dimensional scaling kernel in
\nolinkurl{work/root/check_r24_certificate.py}; this is distinct from the
author all-entry recomputation. The only theoretical dependencies are
the already recorded R13 upper bound, R23 defect identity, and
R14/G17/G18 consequences.

\subsection{R-25: exact bordered induction and size18
certificate}\label{r-25-exact-bordered-induction-and-size18-certificate}

An in-run conditional reduction and finite corollary, without a priority
claim. The permanent block identity, its parameter derivatives, and
irreducible-open intersection are proved in full. The ternary size17
data, evaluation matrix, and selected minor are in
\nolinkurl{border_ternary/border_dense_m17.json} and \nolinkurl{.u16};
\nolinkurl{check_border_ternary.cpp} constructs the truncated subset
recurrence. \nolinkurl{check_border_ryser.cpp} separately checks all
entries by Gray-code Ryser and prefix/suffix products. Root checks the
specified minor by fraction-free Bareiss in
\nolinkurl{work/root/check_border_m17.py}. No statement about untested m
follows from these finite calculations.

\subsection{R-26: uniform Jacobi-continuant
induction}\label{r-26-uniform-jacobi-continuant-induction}

This is an in-run deduction, with no external priority claim.
Dependencies are the already proved R25 exact border cancellation, R23
Tor identity, R25 base size18, and R14/G17/G18 consequences. The
manuscript supplies the exact cofactor and Hessian forms, the normalized
pair matrix, all17 core row lists/residues, and both insertion proofs.
\nolinkurl{check_jacobi_cores.py} and
\nolinkurl{border_ternary/jacobi_uniform_cores.json} supply full finite
matrices; author integer Bareiss checks were performed separately. Root
reconstructed the cores by enumerating path matchings, not by the
continuant recurrence, and checked their determinants by modular
fraction-free Bareiss in \nolinkurl{work/root/check_r26_cores.py}. The
theorem is not a numerical extrapolation and does not assert exact Tor
vanishing.

\subsection{R-27: uniform endpoint repair and exact four-variable
dimension}\label{r-27-uniform-endpoint-repair-and-exact-four-variable-dimension}

This is a proved in-run deduction, with no external priority claim. R26
supplies the continuant/path formulas and exact generic monomial-base
ranks; R25 supplies the single-family border induction; R24 supplies the
first maximal-size base; R13 supplies the dominant sizes and torus upper
bound; R23 identifies the degree-specific Tor defect. The manuscript
proves the general-endpoint K11 identity, the three first-order lifts,
the finite-order rank lemma, complete-character support, and the
whole-epsilon insertion in full. The finite certificate is fully
specified by integer alpha/delta vectors, endpoint forms, row/column
ordering, recurrence, and deterministic pivot rule in
\nolinkurl{r27_fragment.tex}, so its proof is not a JSON-only assertion.

The author builders are \nolinkurl{check_jacobi_repair.py},
\nolinkurl{check_jacobi_boundary_repair.py}, and
\nolinkurl{check_boundary_toeplitz.py}; all three finite arrays and
selected minors are preserved under \nolinkurl{border_ternary}. Root
independently reconstructed every finite coefficient by direct
path-matching enumeration and epsilon-polynomial multiplication in
\nolinkurl{work/root/check_r27_boundary_jets.py}, obtaining the
same33852 entries and ranks84/172/262. The separate author
\nolinkurl{check_r27_one_insertion.py} verifies only the single m17-to18
row embedding; uniformity rests on the written interval/weight argument,
not on this bounded computation. Root\textquotesingle s complete proof
reading and the author whole-proof reread both passed; other received
reviews are recorded by their exact scope.

\subsection{R-28: higher-variable dimensions and common scalar-period
elimination}\label{r-28-higher-variable-dimensions-and-common-scalar-period-elimination}

This is an in-run deduction, without a priority claim. R17/R23 provide
the four-variable determinant rank and cyclic cofactor independence, G21
the permanent dimension, G22 the stable Fourier-frame product point, and
G18 the elementary numerical-semigroup bound. The new rank-r determinant
extension, irreducible finite-stabilizer cyclic witness, common-open
intersection, generic padding-line freeness including exponent1, and
center calculation are proved in \nolinkurl{r28_fragment.tex}.

The cited external background is already G-B15 rather than a duplicate
R-B identifier: Mukai, S. (2003). \emph{An introduction to invariants
and moduli} (W. M. Oxbury, Trans.; Cambridge Studies in Advanced
Mathematics, Vol.81). Cambridge University Press. Theorem5.3 with
Corollary5.4, printedpp161--162, supplies closed-orbit separation;
Proposition5.15, printedp166, supplies openness of stability. Primary
book: \url{https://www.math.ens.psl.eu/~benoist/refs/Mukai.pdf} .
Reopened here as , targeted statements /. The statements/hypotheses were
read; no independent whole-book verification is claimed.

The common-free-point part holds for3\textless=r\textless=m. Its
determinant and padded-image dimension formulas retain their separately
stated r\textgreater=4 hypotheses. No boundary assertion at nu1=(n−m)s
is imported from the four-row theorem.

\subsection{R-29: prescribed moments and two-block common
witnesses}\label{r-29-prescribed-moments-and-two-block-common-witnesses}

This is an in-run deduction, without a priority claim. The full
prescribed-product-moment lemma, two-block Lie-stabilizer calculation,
padding weight identity, and enlarged-family intersection are proved in
\nolinkurl{r29_fragment.tex}. G22 supplies the previously checked
Fourier-frame and Hilbert--Mumford framework; R28 supplies the central
no-scaling and numerical-semigroup conclusion. The exact rank-r
coordinate-ring extraction is the same Cartan calculation in G22 and
does not rely on the old r\textless=m witness range.

The external background is G-B14, not a duplicate R-B identifier: Losev,
I. V. (2006). \emph{The Kempf--Ness theorem and invariant theory}
{[}Preprint{]}. arXiv:math/0605756.
\url{https://arxiv.org/pdf/math/0605756} . Theorem4, printedp1, states
both directions of closed orbit iff intersection with the moment-zero
set; its proof is on printedp2. The primary PDF was reopened as
specifically to check the reverse direction needed to normalize the
polynomial product. The original Kempf--Ness1979 source is cited by
Losev but was not independently opened here. Mukai\textquotesingle s
Hilbert--Mumford, stability-openness and separation inputs retain the
exact G-B15 locators already recorded for G22/R28.

The change of coordinates is an existence consequence of a closed orbit,
not a claimed polynomial-time normalization. The proof carefully
distinguishes the additive moment on ordered tensor factors from the
directly calculated moment of a sum in orthogonal polynomial blocks. It
never uses moment additivity for multiplication of forms.

\subsection{R-30: endpoint products, compact stabilizers, and tilted
padding}\label{r-30-endpoint-products-compact-stabilizers-and-tilted-padding}

An in-run algebraic deduction with no priority claim. R29 supplies the
tensor-induced metric convention and reverse Kempf--Ness normalization.
The new proper-interval circulant argument, rational eigenvalue tilt,
finite padded-form witnesses, and auxiliary-flag elimination are proved
in \nolinkurl{r30_fragment.tex}. The product in exactly m variables
retains a complex torus stabilizer; only its unitary projective
stabilizer is made finite.

The external stabilizer input is already G-B14: Losev (2006), \emph{The
Kempf--Ness theorem and invariant theory}, proof of Theorem1, printedp2,
equations(3)--(4) and the sentence immediately following. At a
moment-zero vector, that passage explicitly concludes that the compact
stabilizer Lie algebra is a real form of the complex stabilizer Lie
algebra. Primary PDF \url{https://arxiv.org/pdf/math/0605756} was
reopened as ; the statement is applied only to the FULL balanced ordered
padded tensor. No assumption that an unbalanced factor has a reductive
stabilizer in the fixed unitary position is made.

The checked alternative finite-symmetrization route is preserved in
\nolinkurl{r30_finite_symmetrization_audit.md}; it is not needed for the
manuscript proof and does not on its own give finite stabilizers. Root,
purity and isolation each completed a whole-proof review; exact hashes
and the later scope-only clarification are recorded in the check events
and \nolinkurl{r30_author_checks.md}.

\subsection{R-31: a proved scalar-period obstruction for the next-row
route}\label{r-31-a-proved-scalar-period-obstruction-for-the-next-row-route}

This is an in-run deduction without a priority claim. Its only matrix
calculation is written explicitly in \nolinkurl{r31_fragment.tex}: a
compensated transposition fixes the padded two-product form and generic
auxiliary covector, has determinant minus one, and after SL
normalization acts on the enhanced tensor by(-1)\^{}b. The exact
invariant extraction is the previously proved G22 Cartan calculation and
requires no stability hypothesis.

The external input is already G-B15: Mukai (2003), Theorem5.3 with
Corollary5.4, printedpp161--162, closed-orbit separation. Since every
odd homogeneous invariant vanishes on the common subcone, separation
places any nonzero closed orbit and its negative in the same orbit. The
result concerns only the actual two-product subcone, not either full
permanent/determinant orbit closure. The separate rank-deficient-moment
failure and unproved richer-pencil routes are recorded in
\nolinkurl{r31_next_row_attempts.md}. No finite computation or unseen
source theorem is used.

\subsection{R-32: stable unpadded comparison
algebras}\label{r-32-stable-unpadded-comparison-algebras}

This is an in-run deduction without a priority claim. It applies the
already proved R29 prescribed-moment and two-block argument with no
padding, and therefore has no first-row padding inequality. The known
external input remains G-B14 (Losev2006 Theorem4) for moment-zero
closedness; R30 supplies the faithful nontrivial partial-flag action
argument and R28 supplies closed-orbit separation plus the semigroup
bound. The exact Cartan extraction and stable quotient dimension are
G22. All new hypothesis checks and the rectangular/nonrectangular
distinction are proved in \nolinkurl{r32_fragment.tex}. Rectangular
directions receive a stable algebra of the claimed dimension but are not
asserted to have invariant-degree gcd one.

\subsection{R-33: exact tangent kernels control generic projective
stabilizers}\label{r-33-exact-tangent-kernels-control-generic-projective-stabilizers}

This is an in-run deduction without a priority claim. G21 supplies the
already proved generic coefficient-map rank; the new proof identifies
its entire kernel with the visibly independent row-column diagonal gauge
on a dense open, then uses every-nonzero Pluecker coordinate to rule out
projective variable Lie symmetries. The final quotient inequality uses
the exact generic-fiber equality in the Stacks Project, Lemma33.20.4(2),
tag0B2L (R-B18 below). Dictionary entry D-S3\textquotesingle s EGA/SGA3
statements concern fiber-dimension upper semicontinuity and stabilizer
openness; they are not used as a substitute for that equality. No new
full EGA proof reading is claimed. The proof avoids any
rational-invariant-field identification and never treats generic finite
stabilizer as stability. Exact reviewed and final source-only hashes are
recorded in \nolinkurl{r33_author_checks.md}.

\subsection{R-B18 --- Exact generic fiber
dimension}\label{r-b18--exact-generic-fiber-dimension}

The Stacks Project Authors. (2026). \emph{The Stacks Project},
Lemma33.20.4(2), tag0B2L. Retrieved November21,2026, from
\url{https://stacks.math.columbia.edu/tag/0B2L}

\begin{itemize}
\tightlist
\item
  Exact statement: for irreducible X and Z=closure f(X), the fiber
  dimension equals dimX-dimZ on a dense open. The displayed proof
  computes the generic fiber dimension as the difference of
  transcendence degrees.
\item
  The earlier D-S3/EGA/SGA source attribution was corrected before final
  compilation; the mathematical proof is unchanged.
\end{itemize}
\section{Isolation sources and exact
locators}\label{isolation-sources-and-exact-locators}

All source paths below refer to the supplied archive extracted once to
\nolinkurl{sources/}. Existing results are not claimed as new. The dates in
bibliographic entries follow publication/version dates, not regenerated
PDF title-page timestamps.

\subsection{I-B01 --- new graph input}\label{i-b01--new-graph-input}

Hoza, W. M. (2026). \emph{NL = UL} (ECCC TR26-323). Supplied file
\nolinkurl{ECCC_TR26-323_Hoza_NL_equals_UL.pdf}.

\begin{itemize}
\tightlist
\item
  Definition 1.2, printed p. 2: internal-vertex path weights; uniqueness
  for all reachable vertex pairs.
\item
  Theorem 1.3, printed pp. 2--3: one seed per fixed simple graph; seed
  length O(log(WN)); logarithmic-space computability; polynomial integer
  magnitudes.
\item
  Formula (1), printed p. 3: rho\_x(a,i)=sum\_j 6\^{}(nu\_2(i)+ell-j)
  C(a)\_j x\_j.
\item
  Theorem 2.4, printed p. 5: failure probability at most 2\^{}(-d)
  W\^{}4 N\^{}2.
\item
  Used as the supplied theorem. This track did not claim a formal or
  full independent verification of the theorem.
\end{itemize}

\subsection{I-B02 --- basis isolation}\label{i-b02--basis-isolation}

Agrawal, M., Gurjar, R., Korwar, A., \& Saxena, N. (2014).
\emph{Hitting-sets for ROABP and sum of set-multilinear circuits}
{[}Preprint{]}. arXiv:1406.7535v1. \url{https://arxiv.org/abs/1406.7535}

Supplied source:
\nolinkurl{LIT_AGKS_hitting_sets_ROABP_basis_isolation/deltaDistanceSICOMP.tex}.

\begin{itemize}
\tightlist
\item
  Theorem 1: source label \nolinkurl{thm:ROABPhs}; source lines 241--246;
  PDF p. 2.
\item
  Definition 5: basis isolating weight assignment; source lines
  902--919; PDF p. 7.
\item
  Lemma 6: source label \nolinkurl{lem:basisIsolation}, lines 927--988; PDF
  pp. 7--8.
\item
  Observation 7: \nolinkurl{obs:smallerWeight}, lines 1005--1028; PDF p. 8.
\item
  Lemma 8: \nolinkurl{lem:weightFunction}, lines 1040--1055 and proof
  thereafter; PDF pp. 9--11.
\item
  Additional I-17 context: Theorem 3 (\nolinkurl{thm:invROABPHS}), printed
  pp. 5 and 19; Lemma 27 (\nolinkurl{lem:bSConc}), printed p. 22, source
  lines 2354--2400, invertible-constant-term matrix concentration.
  Primary PDF refs , , , .
\item
  Primary PDF inspected: \url{https://arxiv.org/pdf/1406.7535} .
\end{itemize}

\subsection{I-B03 --- matrix invariants}\label{i-b03--matrix-invariants}

Forbes, M. A., \& Shpilka, A. (2013). \emph{Explicit Noether
normalization for simultaneous conjugation via polynomial identity
testing} {[}Preprint{]}. arXiv:1303.0084v2.
\url{https://arxiv.org/abs/1303.0084}

Supplied source:
\nolinkurl{LIT_Forbes_Shpilka_explicit_Noether_normalization_matrix_invariants/pit-nn.tex}.

\begin{itemize}
\tightlist
\item
  Theorem 1.3 (\nolinkurl{generatinginvariants}): trace generators of words
  of length at most m² in characteristic zero.
\item
  Theorem 1.13 (\texttt{thm:integral\ subring}) and Corollary 1.14:
  quasipolynomial separating invariants/integral subring.
\item
  Theorem 1.15 (\texttt{thm:orbit\ intersection}): deterministic NC
  orbit-closure intersection already follows from white-box PIT.
\item
  Construction 3.1 (\nolinkurl{roabpinv}), source lines 510--516: M(x),
  trace-generating polynomial.
\item
  Lemma 3.2 (\nolinkurl{diffofinv}), lines 520--526: difference has width
  2m².
\item
  Lemma 3.4 (\nolinkurl{roabpcoeff}), lines 540--567; Corollary 3.6
  (\nolinkurl{reducetoroabp}), lines 581--587.
\item
  Construction 3.7 (\nolinkurl{const:hittingset}), lines 591--595.
\item
  Theorem 3.8 (\nolinkurl{thm:sepinv}), lines 599--613:
  m²\textbar H\textbar{} homogeneous separating invariants.
\item
  Theorem 3.9: specialization of Derksen--Kemper Theorem 2.3.12,
  homogeneous separating invariants give integrality over algebraically
  closed characteristic zero.
\item
  Corollary 3.10: integral-subring consequence for the same set.
\item
  Primary PDF inspected: \url{https://arxiv.org/pdf/1303.0084} .
\end{itemize}

\subsection{I-B04 --- revised GCT V
numbering}\label{i-b04--revised-gct-v-numbering}

Mulmuley, K. D. (2017). Geometric complexity theory V: Efficient
algorithms for Noether normalization. \emph{Journal of the American
Mathematical Society, 30}, 225--309.
\url{https://doi.org/10.1090/jams/864}

Supplied source:
\nolinkurl{LIT_Mulmuley_GCT_V_Noether_normalization/main.tex},
matching arXiv:1209.5993v5 (26 May 2016).

\begin{itemize}
\tightlist
\item
  Theorem 5.14, source label \nolinkurl{tequisdit}, lines 3094--3110; PDF
  pp. 45--46: strengthened general PIT versus strict e.s.o.p. for
  universal-circuit orbit closures. This is stronger than ordinary ROABP
  PIT and has an infinitesimal/border requirement.
\item
  Section 7.5 and Lemma 7.9 (\nolinkurl{lforbes}): replace symbolic
  determinant testing by ROABP testing for matrix invariants.
\item
  Theorem 7.1 (\nolinkurl{tmatrixstrongform}), Theorem 7.2
  (\nolinkurl{tintroorbitclosureNew}) and Theorem 7.7
  (\nolinkurl{tsepderksen}): exact scope of matrix normalization and
  already-known orbit-closure algorithms.
\item
  Primary PDF inspected: \url{https://arxiv.org/pdf/1209.5993} ; the
  cited theorem numbering was checked in this version.
\item
  A reference in the 2013 Forbes--Shpilka paper to an earlier GCT V
  Theorem 1.6 must NOT be copied as the current theorem number. In the
  supplied revised source Theorem 1.6 is a different theorem about
  constant-dimensional groups.
\end{itemize}

\subsection{Independent source check for root finding
G-07}\label{independent-source-check-for-root-finding-g-07}

Leal, M., Lozano Huerta, C., \& Vite, M. (2024). The Noether--Lefschetz
locus of surfaces in P³ formed by determinantal surfaces.
\emph{Mathematische Nachrichten, 297}(12), 4671--4688.
\url{https://doi.org/10.1002/mana.202400132}

Primary preprint: \url{https://arxiv.org/pdf/2303.09028} , version 3 (1
November 2024), inspected in this independent review.

\begin{itemize}
\tightlist
\item
  Definition 2, PDF p. 2: closed determinantal locus det(a,b).
\item
  Proposition 1.1, PDF p. 4: irreducibility and nonempty smooth locus.
\item
  Theorem 2, PDF pp. 3 and 15: prime divisor F1 degree 320112.
\item
  Table 2, PDF p. 13: F1 has a=(5,5,5,5), b=(6,6,6,6), hence linear
  entries.
\item
  Section 4.1, PDF p. 15 immediately before Theorem 2:
  \nolinkurl{P_(20,2)=(1/2)D_(23,14)}.
\end{itemize}

\subsection{I-B05 --- arithmetic support and logarithmic-space
counting}\label{i-b05--arithmetic-support-and-logarithmic-space-counting}

Allender, E., \& Ogihara, M. (1996). Relationships among PL, \#L, and
the determinant. \emph{RAIRO--Theoretical Informatics and Applications,
30}(1), 1--21. \url{https://doi.org/10.1051/ita/1996300100011}

Primary scan inspected:
\url{https://www.numdam.org/item/ITA_1996__30_1_1_0.pdf}

\begin{itemize}
\tightlist
\item
  Section 2, printed p. 4: halting logarithmic-space machines and their
  accepting-minus-rejecting gap.
\item
  Theorem 1, printed p. 2: GapL and logarithmic-space functional
  reductions to integer determinant; following paragraph explicitly
  passes through iterated multiplication of \{0,1,-1\} matrices.
\item
  Section 5, Definition 2, printed p. 11: C\_=L defined by gap zero.
\item
  Theorem 14, printed p. 12: singular integer matrices are
  C\_=L-complete.
\item
  I-12 proves its particular ROABP formulation directly. It does not
  assert that coC\_=L is different from NL or that arbitrary
  polynomial-time algebraic support manipulation is impossible.
\end{itemize}

\subsection{Independent source check for
P-08}\label{independent-source-check-for-p-08}

Lagarias, J. C., \& Odlyzko, A. M. (1977). Effective versions of the
Chebotarev density theorem. In A. Fröhlich (Ed.), \emph{Algebraic number
fields} (pp. 409--464). Academic Press.

Primary-text TeXromancers transcription inspected:
\url{https://aareyanmanzoor.github.io/assets/articles/lagarias-odlyzko.pdf}

\begin{itemize}
\tightlist
\item
  Theorem 1.1 and Corollary 1.2, transcription p. 4; equation (1.4): GRH
  for the Dedekind zeta function of L alone, an effective absolute
  constant, and an unramified prime ideal for every conjugacy class with
  norm bounded by c₂(log Δ\_L)²(log log Δ\_L)⁴.
\item
  The parenthetical immediately following (1.4) handles L=Q by p=(2).
\item
  The source\textquotesingle s full Chebotarev proof was not
  independently checked. P-08\textquotesingle s entire reduction proof
  was independently read and checked; see
  \nolinkurl{independent_review_purity.md}.
\end{itemize}

\subsection{I-B06 --- known rank-one cut and product
closure}\label{i-b06--known-rank-one-cut-and-product-closure}

Agrawal et al. (2014), the same primary supplied source and
arXiv:1406.7535v1 as I-B02.

\begin{itemize}
\tightlist
\item
  Section 5.3, Lemma 32, printed p. 25, source label
  \nolinkurl{lem:width2}: rank-one factorization of singular 2×2 factors
  after multiplication by a nonzero entry.
\item
  Lemma 33, printed p. 26, source label \nolinkurl{lem:lagrange}: Lagrange
  interpolation through a common hitting set preserves each scalar
  factor and its product. Its degree parameter is total degree, and
  disjoint variables are not required.
\item
  Theorem 34, printed p. 26, source label \nolinkurl{thm:ROABP22HS}:
  unknown-order width-two sparse-block result; univariate factors have
  sparsity at most delta+1 and maximum monomial support one, yielding
  polynomial size with the necessary field-size/extension convention.
\end{itemize}

\subsection{I-B07 --- one hidden order, known order, and every
order}\label{i-b07--one-hidden-order-known-order-and-every-order}

Gurjar, R., Korwar, A., \& Saxena, N. (2017). Identity testing for
constant-width, and any-order, read-once oblivious arithmetic branching
programs. \emph{Theory of Computing, 13}(2), 1--21.
\url{https://doi.org/10.4086/toc.2017.v013a002}

\begin{itemize}
\tightlist
\item
  Primary publisher PDF:
  \url{https://theoryofcomputing.org/articles/v013a002/v013a002.pdf}
\item
  Theorem 3.6 and Corollary 3.7, printed p. 11: constant width gives
  polynomial size when the variable order is known, in characteristic
  zero or the stated sufficiently large characteristic.
\item
  Definition in §2.2.2, printed p. 6, and Theorem 4.9: the every-order
  class requires a width bound in each order, unlike a single hidden
  order.
\end{itemize}

Guo, Z., \& Gurjar, R. (2020). Improved explicit hitting-sets for
ROABPs. In \emph{APPROX/RANDOM 2020}, LIPIcs 176, Article 4, 4:1--4:16.
\url{https://doi.org/10.4230/LIPIcs.APPROX/RANDOM.2020.4}

\begin{itemize}
\tightlist
\item
  Primary PDF:
  \url{https://drops.dagstuhl.de/storage/00lipics/lipics-vol176-approx-random2020/LIPIcs.APPROX-RANDOM.2020.4/LIPIcs.APPROX-RANDOM.2020.4.pdf}
\item
  Theorem 1, printed p. 4:3: unknown-order size polynomial in
  \texttt{d*(nr)\^{}(log\ n/max(1,log\ log\ n-log\ log\ r))} in
  characteristic zero or greater than d; an explicitly weaker
  all-characteristic bound is also stated.
\item
  Theorem 3, printed p. 4:4, treats the distinct every-order class.
\end{itemize}

Hu, I., van Melkebeek, D., \& Morgan, A. (2024). Polynomial identity
testing via evaluation of rational functions. \emph{Theory of Computing,
20}(1), 1--70. \url{https://doi.org/10.4086/toc.2024.v020a001}

\begin{itemize}
\tightlist
\item
  Primary PDF:
  \url{https://theoryofcomputing.org/articles/v020a001/v020a001.pdf}
\item
  Theorem 1.9, printed p. 10: the SV-type width assertion includes a
  nonzero monomial of degree at most l+1. That extra promise must be
  retained.
\end{itemize}

Shpilka, A. (2026, November 9). \emph{Polynomial-size hitting sets for
multilinear any-order read-once algebraic branching programs and for
set-multilinear depth-3 circuits with constant-size variable sets}. ECCC
TR26-314. \url{https://eccc.weizmann.ac.il/report/2026/314/}

\begin{itemize}
\tightlist
\item
  Theorem 1.1, printed p. 3: multilinear every-order ROABPs; Remark 1.3
  gives the border extension. The theorem permits arbitrary
  characteristic and a sufficiently large field/extension.
\end{itemize}

Liu, Y. (2026, November 19). \emph{Black-box identity testing for affine
orbits of read-once oblivious algebraic branching programs} (Version 1,
Hexagon 2611.00634v1) {[}Unrefereed preprint{]}.
\url{https://hexagonmath.org/2611.00634v1}

\begin{itemize}
\tightlist
\item
  Primary PDF read: \url{https://hexagonmath.org/pdf/2611.00634v1}
\item
  Section 3.1, Theorem 3.2 and Lemma 3.3, printed pp. 5--6: truncated
  coefficient spaces and the rank-profile support lemma.
\item
  Section 3.2, Lemma 3.4, printed p. 6, and Lemma 3.6, printed p. 7:
  lowest homogeneous components and local rank tests.
\item
  Proposition 3.1, printed p. 4, and §3.4, printed pp. 8--10: generator
  implementation inspected for context. The source\textquotesingle s
  fixed-width polynomial bound is of the form
  \texttt{(nD+1)\^{}(2\^{}\{O(w\^{}2)\})}, not
  \texttt{(nD+1)*2\^{}\{O(w\^{}2)\}}.
\item
  I-21 proves a narrower original-coordinate width-three specialization
  independently in full, with ten interpolation pairs and a common
  scalar center. It does not rely on the Hasse-jet lemma. I-31 later
  proves a separate original-coordinate result for arbitrary width with
  width-dependent exponential pair count; the source\textquotesingle s
  broader affine-orbit and multivariate-layer assertions are not
  certified.
\item
  Host metadata inspected directly (lines 4, 8, 16--18): primarily
  AI-generated text; human understanding no parts; contributed by Yilin
  Liu; CC BY 4.0; OpenAI models used for formulation, proof
  development/review, text, and finite checks. This is source
  provenance, not a mathematical rejection or validation.
\item
  Hexagon has its own identifier namespace. Official
  \url{https://arxiv.org/abs/2611.00634} is an unrelated robotics
  article; it must not be cited for this paper.
\end{itemize}

Singh, S., \& Singh, V. (2026, February 13). \emph{An algebraic rigidity
framework for order-oblivious deterministic black-box PIT of ROABPs}
(arXiv:2602.13449v1) {[}Preprint{]}.
\url{https://arxiv.org/abs/2602.13449}

\begin{itemize}
\tightlist
\item
  Official arXiv record and primary PDF
  \url{https://arxiv.org/pdf/2602.13449} were located .
\item
  The full proof-scope audit I-22 refutes two foundational reductions;
  it does not disprove the headline existence statement. Definition 2.1
  (p. 4) and section 3.1 (p. 6) conflate a scalar output with its
  generated coefficient algebra. Lemma 3.1 (pp. 6--7), Remark 3.1.4 (p.
  9), and Lemma 3.2 (p. 10) incorrectly use radical nilpotence to
  discard nonzero products lying in the radical. Theorem 7.4 (p. 38)
  relies on these steps. Exact counterexamples and full proofs appear in
  I-22.
\end{itemize}

\subsection{I-26 --- determinantal SV vanishing
ideals}\label{i-26--determinantal-sv-vanishing-ideals}

Hu, I., van Melkebeek, D., \& Morgan, A. (2024). Polynomial identity
testing via evaluation of rational functions. \emph{Theory of Computing,
20}(1), 1--70. \url{https://doi.org/10.4086/toc.2024.v020a001}

\begin{itemize}
\tightlist
\item
  Primary PDF:
  \url{https://theoryofcomputing.org/articles/v020a001/v020a001.pdf}
\item
  Theorem 1.3, printed p. 6, and Corollary 1.5, printed p. 7:
  determinantal multilinear degree-(k+1) annihilators of SV\_k on 2k+1
  variables. Theorems 1.9--1.10, printed p. 10, retain the stated
  low-degree-monomial hypothesis for their ROABP assertions. Proposition
  2.2, printed p. 16, gives the relation to rational-function
  evaluation.
\end{itemize}

Fournier, H., \& Korwar, A. (2018, March 16). \emph{Limitations of the
Shpilka--Volkovich generator} {[}WACT presentation slides{]}.
\url{https://wact.imj-prg.fr/pdfs/slides-korwar.pdf}

\subsection{I-B08 --- known logarithmic SV bound for the every-order
class}\label{i-b08--known-logarithmic-sv-bound-for-the-every-order-class}

Forbes, M. A., Saptharishi, R., \& Shpilka, A. (2014). Hitting sets for
multilinear read-once algebraic branching programs, in any order. In
\emph{Proceedings of the 46th Annual ACM Symposium on Theory of
Computing} (pp. 867--875). \url{https://doi.org/10.1145/2591796.2591816}

\begin{itemize}
\tightlist
\item
  Primary full version: \emph{Pseudorandomness for multilinear read-once
  algebraic branching programs, in any order},
  \url{https://arxiv.org/pdf/1309.5668} . Definition 1.1, printed p. 3,
  defines its ``commutative'' class by a width bound in every variable
  order, not by a single hidden order. Pairwise commuting coefficient
  matrices imply membership but are a sufficient representation
  condition, not the definition.
\item
  Theorem 4.1, printed pp. 21--22, and Corollary 4.3, printed p. 22: SV
  with \texttt{1+2\ floor(2\ log\_2\ w)} pairs hits
  individual-degree-less-than-d width-w every-order polynomials, with
  stated field condition \texttt{\textbar{}F\textbar{}\textgreater{}nd}
  and quasipolynomial evaluation size. The proof\textquotesingle s local
  factorization use of the every-order condition and additive SV
  composition were read. Subsidiary rank-condenser proofs were not all
  rechecked.
\item
  Primary locators: , , . The later primary exposition
  Forbes--Shpilka--Volk (2018),
  \url{https://theoryofcomputing.org/articles/v014a018/v014a018.pdf} ,
  was used to locate the original result, not as a replacement for it.
\item
  I-B08 makes explicit that the commuting unipotent I-26 family already
  has a known logarithmic sufficient pair bound. The open question in
  I-26 is for the full one-hidden-order rank-test class.
\end{itemize}

\subsection{I-31 --- full matrix cone
specialization}\label{i-31--full-matrix-cone-specialization}

Liu (2026), the same explicitly unrefereed primary source as I-B07.

\begin{itemize}
\tightlist
\item
  I-31 gives a direct full proof with the looser count2ws-1, avoiding
  the source\textquotesingle s first-block function-field factorization,
  and adds rectangular endpoint absorption in the ground GL\_w space. It
  proves the original-coordinate, one-hidden-order result for all
  widths, with2\^{}\{O(w log w)\} pairs in general; the field,
  interpolation and fixed-width versus variable-width distinctions are
  explicit.
\item
  Root and purity independently checked the complete standalone proof.
  This does not certify the source\textquotesingle s affine-orbit,
  bounded-overlap, or bivariate refinements.
\end{itemize}

\subsection{I-32 --- paired reset
construction}\label{i-32--paired-reset-construction}

This is a deduction authored in the purity track from I-24, with
complete independent proof readings by root and isolation. The exact
retained-coordinate threshold and all partial-interval ranks are proved
in the report. No external priority claim or new source theorem is
invoked.

\subsection{I-33 --- established PIT scope and the explicit prime
construction}\label{i-33--established-pit-scope-and-the-explicit-prime-construction}

Bisht, P., Gupta, N., \& Volkovich, I. (2023). Towards identity testing
for sums of products of read-once and multilinear bounded-read formulae.
In \emph{43rd IARCS Annual Conference on Foundations of Software
Technology and Theoretical Computer Science (FSTTCS 2023)}, LIPIcs 284,
Article 9, 9:1--9:23. \url{https://doi.org/10.4230/LIPIcs.FSTTCS.2023.9}

\begin{itemize}
\tightlist
\item
  Primary publisher PDF:
  \url{https://drops.dagstuhl.de/storage/00lipics/lipics-vol284-fsttcs2023/LIPIcs.FSTTCS.2023.9/LIPIcs.FSTTCS.2023.9.pdf}
\item
  Theorem 1, printed p.9:4: polynomial-time blackbox PIT for the sum of
  two products of read-once formulas. The immediately preceding
  paragraph states \textbar F\textbar\textgreater n*d or access to an
  adequate extension.
\item
  Section1.3.1, printed pp.9:5--9:7, gives the proof overview;
  Theorem48, printed p.9:20, gives the generator G\_(n,4), followed by
  the proof of Theorem1. The statements and proof overview were
  inspected; all subsidiary proofs have not been independently
  certified.
\item
  Primary retrieval locators from the original reads: and . The
  publisher PDF was reopened at the recorded I33 integration event.
  Official ECCC provenance also exists at
  \url{https://eccc.weizmann.ac.il/report/2023/109/} .
\item
  I-33 proves its specific divisor-orbit prime family directly and
  claims no historical priority. Polynomial PIT for its width-two
  subclass is already supplied by I-B06; higher-degree univariate
  factors also split into linear read-once factors after scalar
  extension. The scalar product-closure repair is explicitly the
  AGKS2014 Lemma33 construction, with its source locators in I-B06.
\end{itemize}

\subsection{I-34 --- prime repair under a local nullity
promise}\label{i-34--prime-repair-under-a-local-nullity-promise}

The complete proof is a deduction from the I-33 divisor lemma and
elementary coefficient elimination. Constant-term-count PIT is
background, supplied for example by I-B08; no new general complexity
claim is made for that established case. The arbitrary-term-count
statement retains its local-nullity promise. The three-product additive
identity and the tensor-rank-four example are proved directly and
corroborated by check\_I34.py. Root and purity each independently read
the complete proof. No additional external theorem or historical
priority claim is invoked.

\subsection{I-35 --- contraction tests and exact composition
scope}\label{i-35--contraction-tests-and-exact-composition-scope}

Parthasarathy, K. R. (2004). On the maximal dimension of a completely
entangled subspace for finite level quantum systems. \emph{Proceedings
of the Indian Academy of Sciences---Mathematical Sciences, 114}(4),
365--374. \url{https://doi.org/10.1007/BF02829441}

\begin{itemize}
\tightlist
\item
  Primary full version: \url{https://arxiv.org/pdf/quant-ph/0405077} .
  Propositions 1.3--1.4 and Theorem 1.5, preprint pp.4--6, were reopened
  and their complete local proofs read .
\item
  The 2r−1-dimensional two-factor Hankel test space and its optimality
  over the complex numbers are known background, already recorded as
  P-B16. Its reproof is retained only in
  I35\_known\_background\_audit.md; it is not presented as a new result.
  The same Segre argument gives the explicitly stated algebraically
  closed field formulation.
\end{itemize}

The Stacks Project Authors. (2026). \emph{The Stacks Project}.

\begin{itemize}
\tightlist
\item
  Lemma 27.13.6, tag01WD (Segre embedding):
  \url{https://stacks.math.columbia.edu/tag/01WD} . Primary statement
  and displayed proof read .
\item
  Lemma 33.34.3, tag0B2R (projective intersection dimension):
  \url{https://stacks.math.columbia.edu/tag/0B2R} . Primary statement
  and displayed proof read .
\item
  I-35 proves the rank-one-annihilator criterion, the generic unipotent
  bridge realization, extension-field composition with fresh seeds and
  an explicit polynomial-map degree budget, both synchronization
  failures, and the rank-one multiplicative-corner obstruction. The
  generic small-subspace existence argument asserts neither an effective
  construction nor a field-of-definition bound. Root and purity each
  independently read the full repaired proof.
\end{itemize}

\subsection{I-36 --- pairwise dominance and singular-factor
completion}\label{i-36--pairwise-dominance-and-singular-factor-completion}

This deduction invokes no new external theorem. The rank-one pivot
identity and scalar product closure are the established mechanism
recorded in I-B06 (AGKS2014 Lemmas32--33) and the fully proved P-14. The
full-width-three rank-one completion is the already audited
I-21/I-29/I-30 argument. I-36 proves directly that t+SV1 is dominant on
every pair, applies the pivot-cleared identity to rank-one compound
words, and thereby supplies one whole maximal-minor preservation stage.
All-characteristic scalar extension and finite-grid requirements are
stated explicitly. Root and purity independently read and rederived the
full proof. No computation certificate or historical priority claim is
made.

\subsection{I-37 --- affine singular width-three words and two centered
pairs}\label{i-37--affine-singular-width-three-words-and-two-centered-pairs}

This is a fully proved deduction from I-36 and elementary polynomial
algebra, with no additional external theorem accepted as an unproved
input. The purity agent supplied the first-order pole-extraction idea
while studying Hu, van Melkebeek, and Morgan (2024), \emph{Polynomial
identity testing via evaluation of rational functions}, whose primary
published source and exact RFE/SV locators are already recorded under
I-26 and P-B20. I-37 directly proves the arbitrary-nonzero-slope and
every-characteristic version required here. It also supplies complete
proofs of irreducible multilinear factor inheritance, the
multivariate-UFD rank-one factorization, compatible endpoint gauges, the
tridiagonal Hessian induction and every boundary case of the final
state-path contradiction. These deductions are not attributed as a
theorem of the source paper, and no historical priority claim is made.

The rank-one compound/pair-dominance input is exactly the already
audited I-36. The result is for representations by affine singular
factors, not every polynomial that happens to be multilinear with an
unrelated higher-degree singular representation. Root and purity
independently read and rederived the entire standalone proof; both
passed it without a requested mathematical change. The final polynomial
sampler has five parameters and the exact grid count
(D+1)((n-1)D+1)\^{}2 binom(D+2,2), with the stated working-field
extension convention. No numerical certificate is required or claimed
for this algebraic proof.

\subsection{I-38 --- isolated invertible affine
factors}\label{i-38--isolated-invertible-affine-factors}

This is a proved extension of I-37, not a new external result being
accepted. Every new step is written in full: algebraic triple dominance
without a Jacobian, inheritance of the no-adjacent-rank-three promise,
exact compound connector scope, rank (rather than all-minor)
preservation, endpoint projections, compatible I/P constants with
transported endpoint shears, and the two tangent-kernel elimination.
Purity independently suggested Hh=Hh² rather than a slope-dependent
determinant; the latter can be degenerate in characteristic two. The
proof and the five-parameter polynomial sampler hold in every
characteristic. Root and purity each read and rederived the entire
standalone proof and passed it without requesting a repair. No
historical priority claim or numerical verification claim is made.

\subsection{I-39 --- concrete failures of the next mixed-word
repairs}\label{i-39--concrete-failures-of-the-next-mixed-word-repairs}

The rational identity and residue forms are the explicit P-16
construction, independently read in full earlier in this run. I-39
rederives every needed identity rather than relying on an unstated RFE
theorem. The additional ordinary-node parity scaling, affine matrix
corner and variable anchors, absolute irreducibility, all-interval rank
assertions, and abstract Hessian identities are proved in full. Purity
independently checked the abstract Hessian idea and then read the
complete written proof; root also read and rederived the entire proof.
Both passed it without repair. The all-characteristic assertions of
I-37/I-38 are unchanged; I-39\textquotesingle s two obstructions are
explicitly over characteristic zero. No historical priority claim is
made. The actual symbolic check\_I39\_matrix\_identities.py execution
corroborated the finite displayed matrix and rational identities,
without being described as a general proof certificate.

\subsection{I-B09 --- characteristic-free highest-weight
lines}\label{i-b09--characteristic-free-highest-weight-lines}

Jantzen, J. C. (n.d.). \emph{Representations of algebraic groups and
their Lie algebras, Lecture I}. Primary author lecture:
\url{https://www.math.uni-bielefeld.de/~baumeist/sommerschule/Jantzen1.pdf}
. The exact inputs are printed p.2 (Frobenius reciprocity and both
adjunction maps) and p.3 (for the lower Borel, induction of character d
is Sym\^{}d(V*) for d\textgreater=0 and zero for d\textless0). These
statements and their immediate computations were independently read in
full by isolation, purity and root in this run. No characteristic-zero
simplicity or complete-reducibility result from the later discussion is
used. The upper-Borel eigenline convention follows by dualizing and Weyl
conjugation; a complete sign check and direct Bruhat audit are in
I40\_highest\_weight\_audit.md. This is established background, not a
claimed new representation-theoretic theorem.

\subsection{I-40 --- determinant-character budget and the balanced
nine-variable
closure}\label{i-40--determinant-character-budget-and-the-balanced-nine-variable-closure}

The new deduction is proved in full from I-B09, the already audited
split-torus grading G-B02, I-37/I-38, and the stable P-17 theorem. Both
Borel support spaces provide nonnegative determinant weights, whose sum
is the bond rank times n-2D. The actual minimal-core inclusion yields
the exact matrix covariance and the full-rank increment, rather than
assuming an equivariant representation. Root and purity independently
read and rederived the complete proof and the separate convention audit,
passing without repair. The balanced-ideal degree and lowest-component
inputs retain their exact Hu--van Melkebeek--Morgan
Definition9.5/Theorem9.13/Proposition8.3 locators recorded under P-B20
and in the independent P17 review. The full nine-variable conclusion is
distinguished from the ordinary centered-SV result for the
weight-zero/one translation-invariant subclass. The denominator-safe
finite grid and all extension-field requirements are proved directly. No
historical priority or numerical certificate is claimed.

\subsection{I-41 --- saturated weight-two connectors and the balanced
ten-variable
closure}\label{i-41--saturated-weight-two-connectors-and-the-balanced-ten-variable-closure}

This is a fully proved deduction from I-40\textquotesingle s determinant
budget, I-B09\textquotesingle s established characteristic-free group
representation input, the I-37/I-38 rank and normalization arguments,
and stable P-17. The saturation argument makes only the
determinant-weight-zero bounding support spaces SL2-stable. The adjacent
standard-module conclusion is proved by integer weight pairing,
explicitly excluding Frobenius twists without semisimplicity. The
characteristic-two three-plane is the actual submodule Gamma\^{}2(V)
inside V tensor V; the proof never identifies it with its dual. The
right anchor uses the actual antidiagonal coefficient pairing. Purity
independently checked that matrix calculation and the endpoint
simplification during development, then root and purity independently
read and rederived the entire proof. Both passed with only the same TeX
dim typo, repaired before integration. The resulting five-parameter map
is balanced rational evaluation with common slope one; no ordinary
centered-SV or arbitrary-slope extension is claimed at weight two. No
new external theorem, historical priority, or numerical certificate is
asserted.

\subsection{I-42 --- reset-Hessian thresholds and saturated-run
subclasses}\label{i-42--reset-hessian-thresholds-and-saturated-run-subclasses}

The reset-Hessian argument and both sharp abstract examples are proved
directly by Vandermonde column independence. They invoke no new external
theorem. Root and purity independently rederived the complete draft,
including the characteristic-two overlap and the characteristic-not-two
row(k+1) argument. The saturated-run application uses precisely the
established I-B09 input and I40/I41 canonical support-space arguments.
Its positive promise is stated for an irreducible homogeneous
translation-invariant input and its actual minimal cores; factor
inheritance of that additional run-length promise is not assumed. The
abstract seven- and eight-node Hessians carry no fixed-polynomial
integrability or ROABP-annihilator claim. Root requested explicit
transport of endpoint shears through initial/terminal identity strings,
which is now included from the already checked I38 proof. No computation
certificate or historical priority is claimed.

\subsection{I-B10 --- lines in a Grassmannian over an arbitrary
field}\label{i-b10--lines-in-a-grassmannian-over-an-arbitrary-field}

Ghorpade, S. R., Patil, A. R., \& Pillai, H. K. (2009). Decomposable
subspaces, linear sections of Grassmann varieties, and higher weights of
Grassmann codes. \emph{Finite Fields and Their Applications, 15}(1),
54--68. \url{https://doi.org/10.1016/j.ffa.2008.08.001} . Primary
author/arXiv version:
\url{https://www.math.iitb.ac.in/~srg/preprints/DecomposableSubspaces.pdf}
; arXiv0710.5161v1, \url{https://arxiv.org/abs/0710.5161} . The exact
input is Lemma3, printedp4, and its full proof: two independent
decomposable exterior vectors have decomposable sum precisely when their
associated planes intersect in codimension one. Its ambient field is
arbitrary (specified in section2, printedp3). Isolation, root, and
purity each independently opened this primary PDF and read the complete
lemma/proof in this run. The arXiv primary record confirms the journal
metadata and DOI.

Internal retrieval record: isolation (author PDF, lines159--221), (arXiv
bibliographic record); root . These identifiers are audit locators, not
manuscript bibliography text.

\subsection{I-43 --- weight-three connectors and balanced
eleven-variable
closure}\label{i-43--weight-three-connectors-and-balanced-eleven-variable-closure}

This is a fully proved deduction from the established I-B09/I-B10 inputs
and the already independently checked I40--I42 and P19 results. The
determinant-weight-one Grassmann orbit line supplies invariant
two/four-planes; actual adjacent coefficient support of dimension at
most3 excludes larger torus weights by a physical-unipotent calculation
valid in every characteristic. That constraint makes the normalized
middle compound affine, and a unique internal pole tests its coefficient
using independent anchor values. The determinant-weight-zero case needs
only one canonical anchor; the opposite affine rows are handled
directly, without invariant or self-dual assumptions. Root and purity
independently read and rederived the entire exact standalone proof,
including shared-parameter pivot/endpoint composition, the reset-Hessian
application and the n\textless=11 degree induction, and passed it
without repair. The arbitrary-variable cubic result P19 was
independently reread in full by isolation. No numerical certificate,
ordinary-SV promotion, or historical priority claim is made.

\subsection{I-44 --- exact torus-chain anchors and the next balanced
windows}\label{i-44--exact-torus-chain-anchors-and-the-next-balanced-windows}

This deduction uses no new external theorem. It applies the already
proved exact integer-character covariance of I40 to classify every
possible three-weight step pattern, then enumerates all minimal singular
affine anchors into the unique chain pattern. The coefficient-row and
coefficient-column independence are those of the actual minimal tensor
representation. Both independent reviewers explicitly re-enumerated the
six two-row and four three-row subsets, the excluded multiplicities, the
surviving five matrix patterns, and the actual antidiagonal right
pairing. The two quadratic pole equations are derived directly and use
no division by2. Root and purity each read and rederived the complete
exact standalone proof, passing without repair, including the
common-parameter rank/endpoints, factor reduction and finite degree
windows. I42 supplies only its already audited scalar reset thresholds,
so the distinction between all-characteristic rank preservation for
five-core runs and the characteristic-not-two scalar conclusion remains
explicit. No numerical certificate, new background theorem, ordinary-SV
promotion or historical priority claim is made.
\section{Purity and zero-statistics source
ledger}\label{purity-and-zero-statistics-source-ledger}

Run: CUE\_GCT-20261121-purity. Researcher: GPT-6 Astra Pro
(\nolinkurl{/root/purity_zero_status}); configured reasoning effort is
not exposed. Sources below are background inputs or exact
statement/hypothesis checks, not claimed independent verifications of
entire literature proofs. Deductions P-01--P-11 appear in
\nolinkurl{findings.tex}; their identifiers designate development in this
run and do not assert historical priority. Status deductions P-Z01/P-Z02
appear in \nolinkurl{zero_status.tex}.

\subsection{Background IDs and exact
locators}\label{background-ids-and-exact-locators}

\subsubsection{P-B01 --- Quadratic-character small
bias}\label{p-b01--quadratic-character-small-bias}

Alon, N., Goldreich, O., Håstad, J., \& Peralta, R. (1992). Simple
constructions of almost k-wise independent random variables.
\emph{Random Structures \& Algorithms, 3}(3), 289--304.
\url{https://doi.org/10.1002/rsa.3240030308}

Primary author copy: \url{https://johanhastad.se/almostindep.pdf}
(19-page copy dated September 25, 1997). Exact locator: §4, Theorem 1
(Weil character bound), Construction 2, Proposition 2 and its proof,
printed pp. 9--10 / PDF pages 9--10. P-02 uses only the monic squarefree
case, and changes the arbitrary zero-character convention to +1; the
same n/p error applies. Publication metadata checked against publisher
and author page
\url{https://www.wisdom.weizmann.ac.il/~oded/p_aghp.html}.

\subsubsection{P-B02 --- Curve point
bound}\label{p-b02--curve-point-bound}

Deligne, P. (1974). La conjecture de Weil. I. \emph{Publications
Mathématiques de l'IHÉS, 43}, 273--307.
\url{https://doi.org/10.1007/BF02684373}

Primary scan:
\url{https://www.numdam.org/item/PMIHES_1974__43__273_0.pdf}. Exact
locator: formula (1.5.1), p. 275, and Theorem (1.6), p. 276. Applying
the trace formula to a smooth projective geometrically connected curve,
with H⁰ of dimension 1, H¹ of dimension 2g, and H²=Q\_l(−1), yields
\textbar\#C(F\_q)−(q+1)\textbar≤2g√q. This recovers the classical curve
theorem; no historical-priority claim is made for Deligne in dimension
one. The original theorem statement and trace formula were read, not the
entire proof of Weil I.

\subsubsection{P-B03 --- Trace formula and
weights}\label{p-b03--trace-formula-and-weights}

Deligne, P. (1977). \emph{Cohomologie étale (SGA 4½)} (Lecture Notes in
Mathematics, Vol. 569). Springer. Exact locator: Rapport sur la formule
des traces, Theorem 3.2 (trace formula); Theorem 3.1 and Proposition 3.3
give the determinant/L-function formulation. French primary-text
transcription inspected:
\url{https://grothendiecksga.com/read/sga4.5/fr/II_3.html}.

Deligne, P. (1980). La conjecture de Weil. II. \emph{Publications
Mathématiques de l'IHÉS, 52}, 137--252.
\url{https://doi.org/10.1007/BF02684780}. Exact locators: Theorem 3.3.1
and §3.3.10 for the compact-support weight bound; Theorem 3.5.3 for
equidistribution. The large original PDF was not successfully opened
through this tool; these exact locators were verified in the
author-hosted Katz--Sarnak text below, §9.0.15, which states the needed
result and explicitly cites both 3.3.1 and 3.3.10. No full-proof check
of Weil II is claimed.

Katz, N. M., \& Sarnak, P. (1999). \emph{Random matrices, Frobenius
eigenvalues, and monodromy} (American Mathematical Society Colloquium
Publications, Vol. 45). American Mathematical Society. Author-hosted
book: \url{https://web.math.princeton.edu/~nmk/RMFEM.pdf}. Exact
locators: §§9.0.13--9.0.16, printed pp. 269--270 (duality, affine
vanishing, trace formula, weights, geometric coinvariants); Theorem
9.2.6 (equidistribution under the stated monodromy hypotheses);
§§9.1.5--9.1.7 for smooth proper cohomology and purity. These
sections/theorem were read; no claim of reading the entire book.

\subsubsection{P-B04 --- Effective Chebotarev under Dedekind
GRH}\label{p-b04--effective-chebotarev-under-dedekind-grh}

Lagarias, J. C., \& Odlyzko, A. M. (1977). Effective versions of the
Chebotarev density theorem. In A. Fröhlich (Ed.), \emph{Algebraic number
fields} (pp. 409--464). Academic Press.

Original paper in a later TeXromancers transcription:
\url{https://aareyanmanzoor.github.io/assets/articles/lagarias-odlyzko.pdf}.
Exact locator: Theorem 1.1, Corollary 1.2, formula (1.4), transcription
p. 4. Corollary 1.2 assumes GRH for the Dedekind zeta function of L and
provides an unramified prime in each conjugacy class with norm at most
c₂(log Δ\_L)²(log log Δ\_L)⁴, with an effective absolute constant. L=Q
is explicitly handled by prime 2. Bibliographic identity checked at
author publication list
\url{https://websites.umich.edu/lagarias/PUBLICATIONS/zeta.html}. P-08
is a new conditional reduction using this background bound; it is not a
claim that a suitable common arithmetic encoding exists for ROABPs.

\subsubsection{P-B05 --- Montgomery supported pair
correlation}\label{p-b05--montgomery-supported-pair-correlation}

Montgomery, H. L. (1973). The pair correlation of zeros of the zeta
function. In \emph{Analytic number theory} (Proceedings of Symposia in
Pure Mathematics, Vol. 24, pp. 181--193). American Mathematical Society.

Author-hosted scan:
\url{https://public.websites.umich.edu/hlm/paircor1.pdf}. Exact locator:
§1, theorem and equations (1)--(4), pp. 181--182. RH is assumed at the
opening. The weight is 4/(4+u²); normalization is T log T/(2π); the
theorem gives the form factor uniformly on compact subintervals of
\textbar α\textbar\textless1. P-Z01 is the Fourier consequence for
compactly supported transforms in (−1,1), preserving the original weight
and ordered-index diagonal convention. The full-support conjecture is
distinguished in equations (11)--(12), pp. 183--184. The needed theorem,
normalization, and Fourier conversion were checked, not the complete
13-page proof.

\subsubsection{P-B06 --- Hughes--Rudnick}\label{p-b06--hughesrudnick}

Hughes, C. P., \& Rudnick, Z. (2003). Linear statistics of low-lying
zeros of L-functions. \emph{The Quarterly Journal of Mathematics,
54}(3), 309--333. \url{https://doi.org/10.1093/qmath/hag021}

Primary supplied source:
\nolinkurl{sources/LIT_Hughes_Rudnick_low-lying_zeros_linear_statistics/main.tex}.
Exact locators: Introduction (complex ordinates allowed, no GRH assumed,
source lines 94--150); Definition 2.1 (admissible f, line 340); Theorem
3.1, label \nolinkurl{thm:expectation} (line 383, support {[}−2,2{]});
Theorem 4.1, label \nolinkurl{thm:Wf_variance} (line 505, support
{[}−1,1{]}); Theorem 5.1, label \texttt{thm-L:Gaussian\ moments} (line
621, m\textless2/α); Theorem 8.1, label \nolinkurl{thm:inf_less_1/4R}
(line 1349); Corollary 8.2, label \nolinkurl{thm:inf_less_1/4} (line
1454); Theorem 8.3, label \nolinkurl{thm:pos_prop_small_zeros} (line
1484, β≥0.633). Definition x\_\{χ,j\}=γ\_\{χ,j\} log q/(2π); x\_\{χ,1\}
is the first nonnegative ordinate. Mean/variance/restricted moments are
already unconditional analytic statements; §8 explicitly assumes
Dirichlet GRH. Publication year verified as 2003, not 2004, using the
primary OUP record and author publication list.

\subsubsection{P-B07 --- Rudnick--Sarnak}\label{p-b07--rudnicksarnak}

Rudnick, Z., \& Sarnak, P. (1996). Zeros of principal L-functions and
random matrix theory. \emph{Duke Mathematical Journal, 81}(2), 269--322.

Primary paper copy:
\url{https://empslocal.ex.ac.uk/people/staff/mrwatkin/zeta/rudnick-sarnak.pdf}.
Exact locators: Theorems 1.1--1.2, coefficient Hypothesis (1.7), and
support condition Σ\textbar ξ\_j\textbar\textless2/m. Theorem 1.1
permits complex ordinates in its smoothed analytic statistic; Theorem
1.2 adds target RH for the real-zero/unsmoothed form. The coefficient
condition is required in the stated treatment for degree m\textgreater3.
The m=1 Dirichlet/zeta instances are the instances discharged by
stipulated Alpöge. Statement/hypothesis audit, not full-paper proof
verification.

\subsubsection{P-B08 ---
Iwaniec--Luo--Sarnak}\label{p-b08--iwaniecluosarnak}

Iwaniec, H., Luo, W., \& Sarnak, P. (2000). Low lying zeros of families
of L-functions. \emph{Publications Mathématiques de l'IHÉS, 91},
55--131. \url{https://doi.org/10.1007/BF02698741}

Primary scan: \url{https://www.numdam.org/item/10.1007/BF02698741.pdf}.
Exact locators: §1 standing RH convention p. 56; Theorems 1.1--1.3 for
cuspidal GL₂ target families; Theorems 1.4--1.5 for symmetric-square GL₃
target families; Theorem 1.9 and Theorem 1.10/Hypothesis S (both p. 66);
§4, pp. 88--89, on auxiliary symmetric-square/Rankin--Selberg individual
prime sums and their averaged replacements. The proof does not justify
declaring all the auxiliary individual RH assumptions irreducibly
necessary: the authors explicitly remove some by averaging in the
claimed ranges. Conversely, the target cuspidal and symmetric-square RH
is not Dirichlet or one-dimensional Hecke-character GRH. Statement and
dependency audit only.

\subsubsection{P-B09 --- Conditional Hecke-family
example}\label{p-b09--conditional-hecke-family-example}

Gao, P., \& Zhao, L. (2020). One-level density of low-lying zeros of
quadratic and quartic Hecke L-functions. \emph{Canadian Journal of
Mathematics, 72}(2), 427--454.
\url{https://doi.org/10.4153/S0008414X1900021X}

Primary preprint: \url{https://arxiv.org/pdf/1708.01701} (v3, March 22,
2019). Exact locators: Theorem 1.1, p. 3, quadratic Q(i) family with
support (−2,2), symplectic one-level density; Theorem 1.3, p. 3, quartic
family with support (−20/19,20/19), unitary density; Lemma 2.8, p. 9,
auxiliary nonprincipal Hecke-character prime sum under GRH. These
dependencies are covered conditionally by the user-stated Masullo Hecke
GRH, not by Dirichlet GRH alone.

\subsubsection{P-B10 --- Function-field moment
limitations}\label{p-b10--function-field-moment-limitations}

Sawin, W. (2020). A representation theory approach to integral moments
of L-functions over function fields. \emph{Algebra \& Number Theory,
14}(4), 867--906. \url{https://doi.org/10.2140/ant.2020.14.867}

Primary supplied source:
\nolinkurl{sources/LIT_representation_theory_approach_to_moments_function_fields/main.tex};
arXiv \url{https://arxiv.org/abs/1810.01303}. Exact locators: Hypothesis
1.3, label \nolinkurl{hypo}, line 77; Theorem 1.4, label \nolinkurl{main},
line 85; Remark 1.9(1),(5), lines 204--224. The theorem assumes
H(n,r,tilde r,w), n\textgreater2max(r,tilde r)+1, and characteristic ≠2
if n=4 or 5. Its error includes q\^{}((w−n)/2) C\_\{r,tilde r\}\^{}n
n\^{}(r+tilde r). Remark 1.9(5) obtains w=n−1 from duality, hence
O(q\^{}(−1/2)) with other parameters fixed. Remark 1.9(1) retains
stronger cohomology and uniform Betti-number gates for general stronger
degree-aspect/ratio/correlation conclusions. The paper also proves an
unconditional unbalanced case in Corollary 1.8 (source line 156,
repeated later at label \nolinkurl{combined-result}, line 941), with r\_-=1
and error
sqrt(q)\emph{(q\^{}\{−(p−2r)/(2pr)\}}(r+2)\^{}\{r+1\})\^{}n*n\^{}\{r+1\},
under the stated n/characteristic conditions. This restricted
degree-aspect result must not be omitted. Publisher PDF metadata:
\url{https://msp.org/ant/2020/14-4/ant-v14-n4-p03-s.pdf}; web refs , .
The original source is supplied, so no current universal claim about all
later function-field moment literature is inferred from its 2020
discussion.

\subsubsection{P-B11 --- CFKRS and CUE
distinction}\label{p-b11--cfkrs-and-cue-distinction}

Conrey, J. B., Farmer, D. W., Keating, J. P., Rubinstein, M. O., \&
Snaith, N. C. (2005). Integral moments of L-functions. \emph{Proceedings
of the London Mathematical Society, 91}(1), 33--104.
\url{https://doi.org/10.1112/S0024611504015175}

Primary supplied source:
\nolinkurl{sources/LIT_CFKRS_integral_moments_of_L-functions/integralmoments2v.tex}.
Exact locators: §1.5, \nolinkurl{sec:mainconjectures}, from source line
883; Conjecture 1.5.1, \nolinkurl{thm:zeta2kconjecture}, line 919; exact
finite-N CUE Theorem 1.5.2, \nolinkurl{thm:UN2k}, line 968. The recipe
itself is §2.1. The rectangular/CUE agent owns detailed numbering and
the exact matrix-integral derivation. The arithmetic conjecture is not
obtained merely by discharging RH/GRH in unrelated supported
zero-statistics theorems.

\subsubsection{P-B12 --- Quantum-chaos sources; memory
qualification}\label{p-b12--quantum-chaos-sources-memory-qualification}

Bohigas, O., Giannoni, M. J., \& Schmit, C. (1984). Characterization of
chaotic quantum spectra and universality of level fluctuation laws.
\emph{Physical Review Letters, 52}, 1--4.
\url{https://doi.org/10.1103/PhysRevLett.52.1}

Berry, M. V., \& Keating, J. P. (1999). The Riemann zeros and eigenvalue
asymptotics. \emph{SIAM Review, 41}(2), 236--266.
\url{https://doi.org/10.1137/S0036144598347497}

Berry, M. V., \& Keating, J. P. (1999). H=xp and the Riemann zeros. In
J. P. Keating, D. E. Khmelnitskii, \& I. V. Lerner (Eds.),
\emph{Supersymmetry and trace formulae: Chaos and disorder} (pp.
355--367). Plenum Press. Author-institution bibliographic record:
\url{https://research-information.bris.ac.uk/en/publications/ih-xpi-and-the-riemann-zeros/}.

The symmetry-class connection was recalled from memory and is explicitly
marked that way in the paper. Publication locators and the SIAM
abstract\textquotesingle s H\_cl=XP speculation were verified against
primary publication/author records (; APS DOI above). No uninspected
chapter or equation number is claimed. The stipulated statements do not
construct a specified differential or pseudodifferential Hamiltonian
with the desired arithmetic spectral interpretation. The paper
explicitly distinguishes this from the elementary abstract diagonal
self-adjoint operator obtained by inserting the already-given real zero
ordinates into its definition.

\subsubsection{P-B13 --- User-stipulated inputs and limits of the
attachments}\label{p-b13--user-stipulated-inputs-and-limits-of-the-attachments}

Alpöge: \nolinkurl{sources/ALPOGE_Main.lean} and
\nolinkurl{sources/ALPOGE_Statements.lean}. Read the actual exported
declarations \nolinkurl{riemannHypothesis},
\nolinkurl{riemannZeta_eq_zero_re},
\nolinkurl{LFunction_ne_zero_of_half_lt_re},
\nolinkurl{LFunction_eq_zero_re}, and
\nolinkurl{LFunction_eq_zero_re_of_isPrimitive}, and their statement
predicates. These are used as stipulated theorem input; imports were not
rebuilt, and this ledger is not a truth audit.

Masullo: \nolinkurl{sources/MASULLO_STATEMENTS_INDEX.md} only. The prompt
attributes Hecke GRH to §1.2; this section\textquotesingle s numbered
statement/full proof is not present in the supplied index. Theorem
13.3.1.1, indexed PDF p. 522, is a trace-formula claim; its following
malformed row contains editorial placeholders. Conjecture 14.1.2.1, p.
529, concerns general automorphic GL\_n GRH/GSH and is explicitly a
conjecture. Proposition 14.1.2.3, p. 531, assumes Conjectures 14.1.2.1
and 14.1.2.2. These entries do not validate a general automorphic
theorem. All uses of M remain conditional on the user-stated Hecke
claim, with no invented theorem number or proof check.

\subsubsection{P-B14 --- Explicit conductor
budget}\label{p-b14--explicit-conductor-budget}

Grothendieck, A. (1977). Formule d'Euler--Poincaré en cohomologie étale
(I. Bucur, rédaction). In \emph{Cohomologie l-adique et fonctions L (SGA
5)} (Lecture Notes in Mathematics, Vol. 589), Exposé X. Springer. Exact
locators: Theorem 7.1, Corollary 7.12, formula (7.13). Primary-text
English transcription inspected:
\url{https://grothendiecksga.com/read/sga5/en/X-7.html} ; its French
counterpart did not load in this tool. The l-adic formula used in P-06
is explicitly stated in the author-hosted Katz--Sarnak book, proof of
Lemma 10.1.12, printed p. 296 . That formula plus h\_c\^{}0=0 and
h\_c\^{}2=δ\_ij gives the exact conductor expression; the whole SGA
proof is not claimed independently verified.

\subsubsection{P-B15 --- Rank-one measurements and energy
distortion}\label{p-b15--rank-one-measurements-and-energy-distortion}

Cai, T. T., \& Zhang, A. (2015). ROP: Matrix recovery via rank-one
projections. \emph{The Annals of Statistics, 43}(1), 102--138.
\url{https://doi.org/10.1214/14-AOS1267}

Primary author-hosted article:
\url{https://tony-cai.com/assets/pdf/papers/ROP.pdf}. Exact locators:
Definition 2.1, Theorems 2.1--2.2, p. 107, for the \ensuremath{\ell}₁ measurement
condition; §2.2, pp. 108--110; Lemma 2.1, p. 109, for failure of
constant-distortion \ensuremath{\ell}₂ restricted isometry with subambient sampling in
the Gaussian rank-one matrix model. The lemma\textquotesingle s proof is
assigned by the paper to supplementary material, which was not read.
This is a relevant predecessor, not an attribution of the exact weighted
tensor statement P-09 to the paper. P-09 has its own complete
deterministic trace proof.

\subsubsection{P-B16 --- Width-one kernels as completely entangled
subspaces}\label{p-b16--width-one-kernels-as-completely-entangled-subspaces}

Parthasarathy, K. R. (2004). On the maximal dimension of a completely
entangled subspace for finite level quantum systems. \emph{Proceedings
of the Indian Academy of Sciences---Mathematical Sciences, 114}(4),
365--374. \url{https://doi.org/10.1007/BF02829441}

Primary author preprint: \url{https://arxiv.org/pdf/quant-ph/0405077}.
Exact locators: Proposition 1.3, Proposition 1.4, and Theorem 1.5,
preprint pp. 4--6. Read both complete proposition proofs and the theorem
deduction. The Vandermonde product-vector construction is the diagonal
polynomial-evaluation construction in P-10; the optimal codimension is
Σ\_i(d\_i−1)+1. The elementary arbitrary-node annihilator and the
tensor-train kernel/conditioning formulation are written out in P-10.
Publication metadata corroborated by
\url{https://repository.ias.ac.in/36155/}. The later sections on
perfectly entangled subspaces were not reviewed.

\subsubsection{P-B17 --- Tensor-train ranks and closed coefficient
classes}\label{p-b17--tensor-train-ranks-and-closed-coefficient-classes}

Oseledets, I. V. (2011). Tensor-train decomposition. \emph{SIAM Journal
on Scientific Computing, 33}(5), 2295--2317.
\url{https://doi.org/10.1137/090752286}

Primary published article copy:
\url{https://users.math.msu.edu/users/iwenmark/Teaching/CMSE890/TENSOR_oseledets2011.pdf}.
Exact locators: §2, definition and equation (2.1); Theorem 2.1 and
complete rank-factorization proof, pp. 2298--2299. These were read.
Corollary 2.4, p. 2300, also records closedness/existence for bounded TT
ranks; its proof was read. P-10 uses only exact rank factorization,
homogeneous minor equations, compactness, and the elementary parameter
upper bound. No numerical TT-SVD algorithm was implemented or certified.

\subsection{Scope of citation checks}\label{scope-of-citation-checks}

The bibliography distinguishes primary text, author-hosted copies,
publisher metadata, later transcriptions, supplied source files, and
memory-based explanations. Numerical theorem labels derived from
supplied LaTeX use the source\textquotesingle s shared theorem counters;
integration should preserve the source label as a second locator. This
ledger does not certify source completeness, detect unseen subsequent
literature, or claim formal verification.

\subsubsection{P-B18 --- Source of the rank-profile mechanism used in
the fully proved width-three
specialization}\label{p-b18--source-of-the-rank-profile-mechanism-used-in-the-fully-proved-width-three-specialization}

Liu, Y. (2026, November 19). \emph{Black-box identity testing for affine
orbits of read-once oblivious algebraic branching programs} (Version 1,
Hexagon 2611.00634v1) {[}Unrefereed preprint{]}.
\url{https://hexagonmath.org/2611.00634v1}

Primary PDF: \url{https://hexagonmath.org/pdf/2611.00634v1} . Read
Theorem 3.2 and Lemmas 3.3--3.6 completely, printed pp. 5--7, and the
interpolation-pair construction in §3.4, pp. 9--10. Independently
derived the profile-\{0,2,3\} support-four specialization and reviewed
the complete standalone I-21 proof. P-12 uses that proved
specialization, not an acceptance of the entire preprint. No full
general-width, affine-orbit, or later-section certification is claimed.
The host explicitly labels this primarily AI-generated text and marks no
parts human-understood; version1 is CC BY4.0. Its provenance describes
OpenAI-assisted formulation, proof development/review, drafting, and
checks. This is a Hexagon identifier, not an arXiv identifier. The host
policy says dissemination is not a correctness/novelty judgment:
\url{https://hexagonmath.org/about} .

\subsubsection{P-13 --- General centered-generator arithmetic
encoding}\label{p-13--general-centered-generator-arithmetic-encoding}

The final general encoding lemma uses a stated centered-SV\_K hitting
hypothesis and elementary Kronecker substitution, Eisenstein
irreducibility, the order-discriminant formula, and the trace-dual basis
proof already written in P-12. Its proved specializations are I-30 (K=6
for arbitrary width three) and I-29 (K=2r(w-r)+w-1 under the two-rank
interval promise). I-28, I-29, and I-30 were independently read in full;
detailed reviews and actual events are in work/isolation. No additional
external theorem or literature-priority claim is introduced. The
underlying mechanism provenance remains P-B18, limited to the statements
actually rederived and checked. The earlier eight-pair draft was
superseded before integration by this single general lemma.

P-B18 later source-reading supplement (2026-11-21T02:47:16Z): the
complete image-cone subsection 3.3, printed pp. 7--8, was reopened and
read in the primary PDF, The independent full I-31 proof review is
recorded in work/isolation/independent\_review\_purity\_I31.md. That
original-coordinate specialization is checked in full; the
source\textquotesingle s broader affine-orbit theorem and unrelated
later sections remain outside this review.

\subsubsection{P-B19 --- Established rank-one splitting and width-two
PIT}\label{p-b19--established-rank-one-splitting-and-width-two-pit}

Agrawal, M., Gurjar, R., Korwar, A., \& Saxena, N. (2015). Hitting-sets
for ROABP and sum of set-multilinear circuits. \emph{SIAM Journal on
Computing, 44}(3), 669--697. \url{https://doi.org/10.1137/140975103}

Primary author-hosted version:
\url{https://www.cse.iitk.ac.in/users/nitin/papers/roabp-2014.pdf} .
Exact locators in this 29-page author copy: §5.4, Lemma 34, printed pp.
24--25, gives the width-two rank-one-factor splitting; Lemma 35, p. 25,
gives interpolation closure of a hitting set under products; Lemma 36
and Theorem 37, p. 26, conclude width-two black-box PIT. The complete
Lemma 34 and Lemma 35 proofs and Theorem 37 deduction/context were read
here. P-14 adapts that established splitting in arbitrary dimension
under the rank-\{1,w\} promise, with the fully written
endpoint/row-growth argument and one common polynomial generator. It
makes no literature-priority claim. The author-copy numbering may differ
from another preprint version. Local reading copies are
Hu\_vanMelkebeek\_Morgan\_2024.pdf/txt and
Agrawal\_Gurjar\_Korwar\_Saxena\_2015.pdf/txt in this work directory.

\subsubsection{P-B20 --- Exact rational-evaluation ideal and the scope
of its width
theorem}\label{p-b20--exact-rational-evaluation-ideal-and-the-scope-of-its-width-theorem}

Hu, I., van Melkebeek, D., \& Morgan, A. (2024). Polynomial identity
testing via evaluation of rational functions. \emph{Theory of Computing,
20}(1), 1--70. \url{https://doi.org/10.4086/toc.2024.v020a001}

Primary published PDF:
\url{https://theoryofcomputing.org/articles/v020a001/v020a001.pdf} .
Exact locators: Definition 1.2, printed p. 5, defines RFE; Proposition
2.2, statement p. 16 and complete proof p. 17, identifies ordinary SV
with RFE\_(ell)\^{}(ell−1) after diagonal rescaling. The centered
balanced form RFE\_k\^{}k in P-14 is derived explicitly by adding the
common constant and using partial fractions. Theorem 1.8, pp. 8--9, is
an exact multilinear membership test. Theorem 1.9, p. 10, requires width
less than ell/3+1 \textbf{and a monomial of degree at most ell+1}; it is
not an unrestricted arbitrary-degree ROABP theorem. Proposition 8.3, p.
42, preserves the width of the lowest nonzero homogeneous component; its
proof was read. The §8.2 proof setup, pp. 43--45, uses the
degree-(ell+1) adjacent-block feature. Theorem 9.13, p. 61 with proof
through p. 62, gives the exterior-contraction test for every homogeneous
multilinear degree. Its complete local proof was read using the earlier
image/kernel and boundary lemmas as stated inputs; this was not a
full-paper proof audit. For balanced k=ell and degree d, the condition
uses d−k contractions by polynomial evaluation weights of degree at most
d−1. No unrestricted width conclusion is inferred from that
characterization.

\subsubsection{P-14 --- Scope of the new bridge
deduction}\label{p-14--scope-of-the-new-bridge-deduction}

P-14 gives a completely written dense sampler and exact
coefficient-row-span preservation for arbitrary numbers of original
rank-one/full-rank factors. The simultaneous product argument tolerates
a reset variable appearing in two adjacent bridge polynomials. All pivot
denominators are cleared in a polynomial identity and are nonzero at the
transcendental common center. The class has a pair budget w−1 and
sampling cost (nD+1)\^{}\{O(w)\}; this is a restricted growing-width
construction and an improvement over the r=1 budget in I-29, not
polynomial-in-width PIT. General intermediate factor ranks remain
outside the proof. Root and isolation agents independently read and
checked the full proof; the exact event records are in check\_events.tsv
and the isolation review file.

\subsubsection{P-15 --- Exact width-three one-pair
counterexample}\label{p-15--exact-width-three-one-pair-counterexample}

The polynomial (2x\_1x\_2-(x\_1+x\_2)(x\_3+x\_4)+2x\_3x\_4) and its full
proof are written in P15\_width\_three\_one\_pair\_obstruction.tex. It
is an elementary four-point harmonic relation in the balanced RFE ideal
of P-B20; its explicit width-three unipotent representation, exact cut
ranks, and ordinary centered-SV identity at fourth-root-of-unity nodes
are proved directly. The standard integer-grid statement concerns
balanced RFE on the subset (0,3,2,6), not an unproved identification
with the ordinary centered interpolation generator at that grid. Root
and isolation independently checked the whole proof. check\_P15.py and
P15\_symbolic\_checks.json provide exact symbolic corroboration of the
matrix identity, centered substitution, cut ranks, and harmonic
constants. These finite checks are not a general-width theorem or a
novelty certification.

\subsubsection{P-B21 --- Generic shifts and their characteristic
restriction}\label{p-b21--generic-shifts-and-their-characteristic-restriction}

Forbes, M. A., Ghosh, S., \& Saxena, N. (2018). Towards blackbox
identity testing of log-variate circuits. In \emph{45th International
Colloquium on Automata, Languages, and Programming (ICALP 2018)}
(LIPIcs, Vol. 107, Article 21, pp. 21:1--21:16).
\url{https://doi.org/10.4230/LIPIcs.ICALP.2018.21}

Primary author copy:
\url{https://www.cse.iitk.ac.in/users/nitin/papers/icalp18-log-var-hsg.pdf}
. Theorem 2, printed p. 21:4, requires characteristic zero or greater
than the polynomial degree and says that a basis-isolating-weight shift
gives a cone-closed coefficient basis. The paragraph on printed p. 21:5
explicitly draws the formal/random independent-shift consequence. Lemma
11, p. 21:8, gives the support bound at most log2 of the
coefficient-space dimension. These statements and the short Lemma 11
proof were inspected; the complete Theorem 2 proof was not re-audited in
this lookup. No all-characteristic version or explicit efficient shift
is inferred.

\subsubsection{P-B22 --- Factor width is not monotone, even for
irreducible
factors}\label{p-b22--factor-width-is-not-monotone-even-for-irreducible-factors}

Andrews, R., Armand, J., Dwivedi, P., Hansen, M. R. D., Limaye, N.,
Srinivasan, S., \& Tavenas, S. (2026). On closure properties of
read-once oblivious algebraic branching programs. In \emph{17th
Innovations in Theoretical Computer Science Conference (ITCS 2026)}
(LIPIcs, Vol. 362, Article 9, pp. 9:1--9:21).
\url{https://doi.org/10.4230/LIPIcs.ITCS.2026.9}

Official publisher record:
\url{https://drops.dagstuhl.de/entities/document/10.4230/LIPIcs.ITCS.2026.9}
. Exact statement: Theorem 1, printed p. 9:3; repeated with its local
proof on p. 9:11. It gives, over every field, width w=2\^{}\{O(n)\}
polynomials having an irreducible factor requiring width
w\^{}\{Omega(log d)\} in every order. The publisher metadata and theorem
statement were verified; the local proof/preceding matching argument
were inspected without claiming a complete paper audit. This reinforces
the scope restriction in the RFE-ideal investigation: passing to an
irreducible factor does not preserve an available width bound.

\subsubsection{P-16 --- Centering is necessary for the proved two-pair
full-rank
class}\label{p-16--centering-is-necessary-for-the-proved-two-pair-full-rank-class}

P16\_uncentered\_two\_pair\_obstruction.tex gives the complete explicit
rational-function construction. Its general unbalanced identity has
6(k+1) algebraic nodes in a field of degree at most 4(k+1)\^{}2
phi(2k+2); the k=2 specialization has eighteen nodes and field degree at
most72. Multilinear reversal and the exact diagonal
D\_b=A\textquotesingle(b) turn it into an ordinary uncentered SV2
annihilator of degree16 and minimum width3, with explicit invertible
factors. P14 supplies the centered repair. The only external
identification used is the established ordinary SV/RFE diagonal
equivalence P-B20, also written out with signs in the proof. Root and
isolation independently read the whole proof and requested only a
closing-display-delimiter repair, which was made. No literature-priority
claim or balanced-generator counterexample is asserted.

\subsubsection{P-B23 --- Known secant/prolongation
background}\label{p-b23--known-secantprolongation-background}

Sidman, J., \& Sullivant, S. (2009). Prolongations and computational
algebra. \emph{Canadian Journal of Mathematics, 61}(4), 930--949.
\url{https://doi.org/10.4153/CJM-2009-047-5}

Primary author preprint: \url{https://arxiv.org/pdf/math/0611696} (v3,
April3,2008). Definition1.1, p1, defines prolongation in characteristic
zero. Theorems4.1--4.2 and Proposition4.5, pp13--14, identify the first
possible secant-degree equations and give the symbolic-power
containment. The complete local proofs were read, with earlier
Corollary2.10 used as a stated input; no full-paper audit is claimed.
The specialized balanced-RFE derivative/cylinder argument is kept in
balanced\_rfe\_prolongation\_attack.md as an audit, not a reproof of
known background in the manuscript. I37 separately proves exactly the
elementary arbitrary-characteristic pole identity it needs.

\subsubsection{P-13 --- Later proved singular-factor
specializations}\label{p-13--later-proved-singular-factor-specializations}

I36 supplies K5 for rational width-three words with singular
variable-dependent factors, giving field degree(nD+1)\^{}11. I37
supplies K2 for the further affine singular-factor class, giving
degree(nD+1)\^{}5. Common scaling and changes to a constant output
column preserve both promises, so the general P13 trace-span statement
also applies. I37 was independently whole-read here; its review is
work/isolation/independent\_review\_purity\_I37.md. The arbitrary
width-three degree remains exponent13, and unrestricted mixed-factor or
nonlinear two-pair claims are not added.

I38 extends the same P13 exponent5 specialization to affine width-three
words with no two consecutive variable-dependent rank-three factors
after constants are absorbed. The full I38 proof was independently read
and rederived; see work/isolation/independent\_review\_purity\_I38.md.
Its triple-dominance argument and two tangent directions are algebraic
in every characteristic. Longer invertible runs remain outside that
proof.

\subsubsection{P-B24 --- Characteristic-zero SL2 representation
input}\label{p-b24--characteristic-zero-sl2-representation-input}

Etingof, P. (2020). \emph{Lecture11: Representations of Lie groups and
Lie algebras}. In \emph{18.745 Lie Groups and Lie Algebras I}, Fall2020.
MIT OpenCourseWare.
\url{https://ocw.mit.edu/courses/18-745-lie-groups-and-lie-algebras-i-fall-2020/mit18_745_f20_lec11.pdf}

Theorem 11.16, parts (iii)--(iv), printed pp. 66--67 (PDF pp. 5--6),
classifies finite-dimensional irreducibles as symmetric powers of the
standard two-dimensional representation and states complete
reducibility. The preceding polynomial realization gives their dimension
r+1. The complete local proofs were read; no complete-course proof audit
is claimed. The official course page identifies Pavel Etingof and
Fall2020. This source was used in the initial characteristic-zero P17
draft. The final proof replaces it with the characteristic-free
split-torus grading theorem in SGA3, ExposéI, Proposition4.7.3 (G-B02),
so no complete-reducibility input remains in the final claim.

\subsubsection{P-17 --- Exact external inputs for the balanced two-pair
cut
argument}\label{p-17--exact-external-inputs-for-the-balanced-two-pair-cut-argument}

The cubic decoder and the full multilinear degree interval use P-B20,
Definition9.5 p55 and Theorem9.13 pp61--62; the lowest homogeneous
component uses Proposition8.3 p42. The minimal tensor-train core
factorization is P-B17, Theorem2.1 pp2298--2299. The all-characteristic
odd-cut parity argument uses only the split-torus grading theorem from
SGA3, ExposéI, Proposition4.7.3 (G-B02), typeset p22/PDFpage22, original
margin p36. Primary scan:
\url{https://webusers.imj-prg.fr/~patrick.polo/SGA3/origExp1.pdf} . Its
complete local proof was read: the comodule coefficients are orthogonal
projectors whose images give the grading. No full SGA volume audit is
claimed. The resulting odd-character pairing and ROABP application are
proved in the fragment; the published representation theorem itself is
not reproved. The final scope is all characteristics, every multilinear
polynomial on at most eight variables, every homogeneous cubic on at
most ten variables, and the stated decomposable/invariant subclasses at
arbitrary length. No general n-variable width-three theorem or
historical-priority certification is asserted.

\subsubsection{Subsequent I40 refinement of
P17}\label{subsequent-i40-refinement-of-p17}

The separate I40 result closes the n9 quartic left by P17 using an
all-characteristic Borel determinant-weight budget and the previously
checked I38 sampler. The complete stable proof and separate
highest-weight convention audit were independently read here; see
work/isolation/independent\_review\_purity\_I40.md (actual
event05:37:04UTC). I independently read Jantzen\textquotesingle s
primary LectureI, printed pp2--3, for Frobenius reciprocity and the
explicit lower-Borel SL2 induction computation. Source URL:
\url{https://www.math.uni-bielefeld.de/~baumeist/sommerschule/Jantzen1.pdf}
. The final I40 manuscript cites this established representation input
and retains the elementary Bruhat calculation only in an audit. P17
itself was not modified after its exact-version reviews.

\subsubsection{P-18 --- Ordered trilinear sampler and exact remaining
cubic support
gate}\label{p-18--ordered-trilinear-sampler-and-exact-remaining-cubic-support-gate}

P18 uses the already inspected HMM cubic decoder, Definition 9.5 p. 55
and Theorem 9.13 pp. 61--62 (P-B20), and the independently proved P17
three-block lemma. The adjacent-block truncation mechanism is HMM
Proposition 8.6, statement p. 44 and complete local proof p. 45; that
text was reopened and read while deriving the balanced cubic
specialization. The fragment spells out only its precise application to
the three balanced contraction equations. I37 supplies pole extraction
and I40 supplies the denominator-safe interpolation grid. The
all-characteristic arbitrary-length disjoint-trilinear conclusion and
the minimal-counterexample reduction to at most two shared boundary
variables are completely proved in P18. Root and isolation independently
read the final exact fragment, SHA256
87282923d4e8b02925d01cd6ca00ac3ed2cb57d7fa97634d8055ac9031714147. The
separate cubic prolongation calculation remains in
cubic\_ordered\_gate.md as an audit and is not a hidden input to the
positive theorem. No unrestricted cubic or growing-width claim is made.

\subsubsection{P-19 --- Every multilinear width-three cubic is hit by
balanced two-pair
evaluation}\label{p-19--every-multilinear-width-three-cubic-is-hit-by-balanced-two-pair-evaluation}

P19\_all\_cubic\_balanced.tex proves the arbitrary-length and
every-characteristic cubic conclusion from the exact P18 interval
reduction, prime one-pair membership of first partials, and
P-B20\textquotesingle s cubic exterior decoder (Definition 9.5, p. 55;
Theorem 9.13, pp. 61--62). The two literal coefficient submatrices and
the at-most-three-coordinate Vandermonde/exterior projection argument
are written out completely. I40 supplies the finite grid of size
20(n+1)\^{}2. Root and isolation independently read the complete exact
fragment, SHA256
4bb5f2f9588840bf41986b657ba3493ebd92cda0f01bea09f52ba896b1535cee,
including the full P18 dependency, and passed it without repair. No
additional external theorem or historical-priority assertion is used.
HMM Proposition 8.3, p. 42, and multilinear reversal give the explicit
remaining degree window {[}4,n-4{]}; arbitrary higher-degree or
growing-width PIT is not asserted.
\section{Integral geometry and archive source
register}\label{integral-geometry-and-archive-source-register}

These are framework and source-dependency checks. They do not claim full
independent verification of every cited literature proof. New finite
constructions and transfers are proved under their finding IDs in the
manuscript.

\subsection{G-B01 --- Grothendieck\textquotesingle s supplied cote 133
transcriptions}\label{g-b01--grothendiecks-supplied-cote-133-transcriptions}

Grothendieck, A. (n.d.). \emph{Cote 133, scans 006, 008, 009, 019, and
028} {[}Supplied LaTeX transcriptions{]}.

Exact files: \nolinkurl{ARCHIVE_133_scan006.tex},
\nolinkurl{ARCHIVE_133_scan008.tex}, \nolinkurl{ARCHIVE_133_scan009.tex},
\nolinkurl{ARCHIVE_133_scan019.tex}, \nolinkurl{ARCHIVE_133_scan028.tex}.
All five transcriptions were read. No original image is supplied. See
TRANSCRIPTION\_NOTES.tsv for file-line locators and separate
mathematical and paleographic confidence. The accompanying
\nolinkurl{SA17_STUDY_cote133_lambda.md} is an earlier research record,
not an independently reverified image reading.

\subsection{G-B02 --- SGA 3}\label{g-b02--sga-3}

Demazure, M., \& Grothendieck, A. (Eds.). (1970). \emph{Schémas en
groupes (SGA 3)} (Lecture Notes in Mathematics, Vols. 151--153).
Springer.

Exact locators: Exposé I, Propositions 4.7.2--4.7.4 (comodules,
diagonalizable groups and gradings, splitting), Theorem 5.3.3
(diagonalizable-group cohomology over an affine base). Original Exposé I
inspected at
\url{https://webusers.imj-prg.fr/~patrick.polo/SGA3/origExp1.pdf} .
Exposé VI\_B, Proposition 4.1 (fiber dimension for group schemes),
electronic primary-text transcription
\url{https://grothendiecksga.com/read/sga3/fr/VIB-4.html} . It
explicitly cites EGA IV\_3, Theorem 13.1.3. The entire EGA proof was not
reread here.

\subsection{G-B03 --- SGA 6}\label{g-b03--sga-6}

Berthelot, P., Grothendieck, A., \& Illusie, L. (Eds.). (1971).
\emph{Théorie des intersections et théorème de Riemann--Roch (SGA 6)}
(Lecture Notes in Mathematics, Vol. 225). Springer.
\url{https://doi.org/10.1007/BFb0066283}

Exact locators: Exposé V, Definition 2.1, Example 2.2(b), Lemma 2.2.1,
Definition 2.4 and equation (2.4.2), Proposition 3.2, §3.10 equation
(3.10.1), §7.1 equation (7.1.1), Proposition 7.5; Exposé VI, §3.2
(special-lambda assertion for quasi-compact schemes). The V passages
were read in \url{https://grothendiecksga.com/read/sga6/fr/V-2.html} ,
\url{https://grothendiecksga.com/read/sga6/fr/V-3.html} and
\url{https://grothendiecksga.com/read/sga6/fr/V-7.html} . The VI 3.2
locator was located/corroborated, not claimed as a complete proof
reading. A separate 24-page PDF found during retrieval was an incomplete
front-matter excerpt and was not treated as the full volume. For R23,
additionally read Exposé III, Definition1.5 and Notation1.6 at
\url{https://grothendiecksga.com/read/sga6/fr/III-1.html} , and Exposé
I, Proposition5.1, Definition5.2 and paragraph5.4 at
\url{https://grothendiecksga.com/read/sga6/fr/I-5.html} . These are
electronic French primary-text transcriptions. III1.5 explicitly refers
to EGAIII6.5; that reference is corroborated through SGA6 only.

\subsection{G-B04 --- Integral representation
rings}\label{g-b04--integral-representation-rings}

Serre, J.-P. (1968). Groupe de Grothendieck des schémas en groupes
réductifs déployés. \emph{Publications Mathématiques de
l\textquotesingle IHÉS, 34}, 37--52.
\url{https://doi.org/10.1007/BF02684589}

Exact locators: Theorem 5, §3.7; §3.8 for GL\_n. Author-hosted original
scan inspected:
\url{https://www.college-de-france.fr/media/jean-pierre-serre/UPL8576975072934521334_Serre_Grpes_de_Gr.pdf}
. The theorem concerns Grothendieck groups for split reductive group
schemes over a principal ideal domain; it is not a claim that
characteristic-zero simple decompositions remain semisimple in positive
characteristic, or a solution of the distinct finite-S\_n lattice
question.

\subsection{G-B05 --- Hilbert and Quot
schemes}\label{g-b05--hilbert-and-quot-schemes}

Grothendieck, A. (1961). Techniques de construction et théorèmes
d\textquotesingle existence en géométrie algébrique IV: Les schémas de
Hilbert. \emph{Séminaire Bourbaki, 6}, Exposé 221, 249--276.
\url{https://www.numdam.org/article/SB_1960-1961__6__249_0.pdf}

Exact locators: Theorem 3.1 (Quot representability), Theorem 3.2 (fixed
Hilbert polynomial), Proposition 3.8 (Grassmannian immersion). Relevant
original pages inspected. These are parameter-space theorems, not
efficient algorithms for the present rank computations.

\subsection{G-B06 --- EGA graded modules and
dimension}\label{g-b06--ega-graded-modules-and-dimension}

Grothendieck, A., \& Dieudonné, J. (1961). Éléments de géométrie
algébrique II: Étude globale élémentaire de quelques classes de
morphismes. \emph{Publications Mathématiques de l\textquotesingle IHÉS,
8}, 5--222.

Exact locator: §§2.4--2.5, Proposition 2.5.2 for the sheaf associated to
a graded module (2.5.1 is the setup; Definition 2.5.3 refers back to
Proposition 2.5.2). Electronic translation inspected at
\url{https://jcreinhold.github.io/ega/ii/02-02-homogeneous-prime-spectra.html}
. The original large PDF did not load in the web tool. For G19 and I35,
the indexed French original passages for Definition5.4.1, printedp100,
and Theorem5.5.3(i), printedp104, including the latter\textquotesingle s
displayed proof, were read at
\url{https://geotoppos.org/divers/ega2.pdf} . They give universal
closedness and projective implies proper. A direct temporary download
failed with HTTP403; no complete original-PDF reading is claimed. EGA
IV\_3, Theorem 13.1.3 is used through its explicit citation and
application in SGA 3 VI\_B Proposition 4.1, with this check scope
recorded.

\subsection{G-B07 --- Determinantal plane curves and cubic
surfaces}\label{g-b07--determinantal-plane-curves-and-cubic-surfaces}

Beauville, A. (2000). Determinantal hypersurfaces. \emph{Michigan
Mathematical Journal, 48}, 39--64.
\url{https://doi.org/10.1307/mmj/1030132717}

Author PDF: \url{https://math.univ-cotedazur.fr/u/beauvill/pubs/det.pdf}
. Exact locators: Proposition 3.1(a), §3.2, Proposition 3.3 and Remark
3.4 for plane curves and their line bundles; Corollary 6.4 for smooth
cubic surfaces; Corollary 1.12 for the relevant linear resolution. These
statements were inspected. G05 supplies the required
density/highest-weight deduction; R11 supplies the low-degree padding
deduction.

\subsection{G-B08 --- The reduced determinantal-quartic
divisor}\label{g-b08--the-reduced-determinantal-quartic-divisor}

Leal, M., Lozano Huerta, C., \& Vite, M. (2024). Noether--Lefschetz loci
of determinantal surfaces. \emph{Mathematische Nachrichten, 297}(12),
4671--4688. \url{https://doi.org/10.1002/mana.202400132}

Primary version inspected: arXiv:2303.09028v3,
\url{https://arxiv.org/pdf/2303.09028} . Exact locators: Theorem 2;
Table 2, component F\_1 with a=(5,5,5,5), b=(6,6,6,6), degree 320112;
§4.1 immediately preceding Theorem 2,
\nolinkurl{P_(20,2)=(1/2)D_(23,14)}. Original table and theorem were
visually checked. The attached Landsberg survey\textquotesingle s
section zhanglee/theorem zlthm gives 640224; the reduced divisor degree
required for a principal equation is 320112. G07 records this
cycle-multiplicity correction, separately from manuscript transcription.

\subsection{G-B09 --- Known optimal order for uniform determinantal
representations}\label{g-b09--known-optimal-order-for-uniform-determinantal-representations}

Boralevi, A., van Doornmalen, J., Draisma, J., Hochstenbach, M. E., \&
Plestenjak, B. (2017). Uniform determinantal representations. \emph{SIAM
Journal on Applied Algebra and Geometry, 1}(1), 415--441.
\url{https://doi.org/10.1137/16M1085656}

Primary author copy:
\url{https://hochsten.win.tue.nl/pdf/siag-unidetrep.pdf} . Exact
locators: Theorem 1.3, printed p.417, and its construction in §4,
printed pp.425--427; Proposition 4.5; Remark 4.8. The asymptotic order
d\^{}(v/2) for v fixed affine variables is a known theorem. G09 gives a
displayed finite shared-chain bound and the low-row multiplicity
consequence, without claiming the asymptotic exponent as new. Both
independent reviewers checked the stated scope and the
construction\textquotesingle s coordinate-chart qualification.

\subsection{G-B10 --- Supplied analytic
inputs}\label{g-b10--supplied-analytic-inputs}

Alpöge, L. (2026, November 20). \emph{Remark} {[}User-stipulated theorem
and supplied note{]}. \url{https://alpo.ge/remark.pdf}

Exact locators: steps 2--5 p.3; steps 7,8,13,17 p.4; exported theorem
names and predicates in ALPOGE\_Main.lean and ALPOGE\_Statements.lean.
The source files and four-page note were read; the imported Lean archive
was not supplied or rebuilt.

Masullo, A. M. (2026, November 20). \emph{The music of the primes}
{[}Draft; supplied statement index only{]}.

The general Hecke-GRH claim is the prompt\textquotesingle s §1.2
attribution; its numbered statement is absent from the supplied index.
Theorem13.3.1.1 p.522 and Conjecture14.1.2.1 p.529 were inspected in the
index. Every dependency on the general Hecke claim is conditional.

\subsection{G-B11 --- All-degree Hilbert upper
bound}\label{g-b11--all-degree-hilbert-upper-bound}

Chardin, M. (1989). Une majoration de la fonction de Hilbert et ses
conséquences pour l'interpolation algébrique. \emph{Bulletin de la
Société Mathématique de France, 117}(3), 305--318.
\url{https://doi.org/10.24033/bsmf.2124}

Primary PDF:
\url{https://www.numdam.org/article/BSMF_1989__117_3_305_0.pdf} . Exact
locator: central unnumbered theorem on printed p.306, with its
definition of EGR and projective dimension D; proof continues on
pp.307--310. The displayed all-degree bound was read and visually
checked. G10/G11 apply it to reduced irreducible characteristic-zero
projective varieties, without a Cohen--Macaulay hypothesis.

\subsection{G-B12 --- Degree of an algebraic
image}\label{g-b12--degree-of-an-algebraic-image}

Sombra, M. (1997). Bounds for the Hilbert function of polynomial ideals
and for the degrees in the Nullstellensatz. \emph{Journal of Pure and
Applied Algebra, 117--118}, 565--599.
\url{https://doi.org/10.1016/S0022-4049(97)00028-5}

Primary author preprint:
\url{https://www.maia.ub.edu/~sombra/publications/hilbert/hilbert5.pdf}
. Exact locators in that preprint: §1, p.5, projective Bézout bounds on
the sum of degrees of components; Proposition2.10, pp.11--12, degree
bound for an everywhere regular polynomial map. G10 explicitly proves
the rational-map version it needs, accounting for the base locus. No
rational-map statement is attributed verbatim to Proposition2.10.

\subsection{G-B13 --- Degree bounds for invariant
generators}\label{g-b13--degree-bounds-for-invariant-generators}

Derksen, H. (2001). Polynomial bounds for rings of invariants.
\emph{Proceedings of the American Mathematical Society, 129}(4),
955--963. \url{https://doi.org/10.1090/S0002-9939-00-05698-7}

Primary author copy:
\url{https://bpb-us-e1.wpmucdn.com/sites.northeastern.edu/dist/7/497/files/2020/12/A.I.a.7.pdf}
. Exact locators: Theorem1.1, Proposition1.2, and Example2.1, printed
p.956. Example2.1 gives the symmetric-power generator-degree estimate
used for quaternary forms. The Veronese generation bound and ensuing
orbit-closure comparison in G11 are proved separately. For G16, the
general finite-kernel Proposition1.2 and Theorem1.1 were reread on p956;
the action on Sym\^{}n(V) tensor Sym\^{}a(V*) has coefficient degree
n+3a, giving the stated enhanced-representation bound. This is distinct
from its determinant parametrization degree n+a.

\subsection{G-B14 --- Closed orbits from the moment
map}\label{g-b14--closed-orbits-from-the-moment-map}

Losev, I. V. (2006). \emph{The Kempf--Ness theorem and invariant theory}
{[}Preprint{]}. arXiv:math/0605756.
\url{https://arxiv.org/abs/math/0605756}

Exact locator: Theorem4, p.1, with proof on p.2. The paper states that
the general results are established background and gives proofs. The
original source it cites is Kempf, G., \& Ness, L. (1979), ``The length
of vectors in representation spaces,'' \emph{Lecture Notes in
Mathematics, 732}, 233--243. The locator actually inspected and used
here is Losev Theorem4; no separate inspection of the 1979 original is
claimed.

\subsection{G-B15 --- Stable affine quotient
dimensions}\label{g-b15--stable-affine-quotient-dimensions}

Mukai, S. (2003). \emph{An introduction to invariants and moduli} (W. M.
Oxbury, Trans.; Cambridge Studies in Advanced Mathematics, Vol.81).
Cambridge University Press.

Primary book copy inspected:
\url{https://www.math.ens.psl.eu/~benoist/refs/Mukai.pdf} . Exact
locators: Theorem5.9 (affine quotient surjectivity), Proposition5.10 and
its proof (invariant closed subsets and exactness), printed pp.164--165;
Definition5.12 (closed orbit and finite stabilizer), p.165;
Proposition5.15 (open stable locus) and Theorem5.16 (separation), p.166;
Corollary5.17 (geometric stable quotient), p.167. G11 supplies explicit
stable points before applying the fiber-dimension deduction. For G18,
Theorem5.3 and its complete two-step proof, printed pp.161--162,
Corollary5.4, p.162, and Proposition5.15, p.166, were reopened and read.
They give separation of distinct closed orbits by invariants and
openness of the stable locus; the padding-line and effective semigroup
arguments are proved in G18 itself. For G19, Definition5.20 and
Corollary5.21, printedp170, and the character discussion in
Proposition5.19, printedp168, were reopened and checked. For G22,
Theorem7.4, printedp212, its linear ray-type formulation equation(7.3),
p213, and its displayed proof pp218--219 were freshly read. This
supplies the stable one-parameter-subgroup minimum-weight criterion; the
Fourier frame and multiplication step are proved in G22.

\subsection{G-B16 --- Existing quadratic border-complexity
bound}\label{g-b16--existing-quadratic-border-complexity-bound}

Landsberg, J. M., Manivel, L., \& Ressayre, N. (2013). Hypersurfaces
with degenerate duals and the Geometric Complexity Theory Program.
\emph{Commentarii Mathematici Helvetici, 88}(2), 469--484.
\url{https://doi.org/10.4171/CMH/292}

Primary author version inspected:
\url{https://math.univ-lyon1.fr/~ressayre/PDFs/lmr.pdf} , dated
April2010. Its exact locator is Theorem1.0.1, p.1, with the
border-complexity definition immediately preceding it. The published
bibliographic metadata was corroborated by the publisher record. This
source supplies established prior-result context: the present
linear-padding multiplicity constructions do not improve that quadratic
lower bound. No latest/best lower-bound assertion or acceptance of
unrelated newer preprints is made.

\subsection{G-B17 --- The sl2 tensor and harmonic
decompositions}\label{g-b17--the-sl2-tensor-and-harmonic-decompositions}

Fulton, W., \& Harris, J. (1991). \emph{Representation theory: A first
course} (Graduate Texts in Mathematics, Vol. 129). Springer.
\url{https://doi.org/10.1007/978-1-4612-0979-9}

Exact locators: section11.2, Exercise11.11, printed p.151
(Clebsch--Gordan); Exercise11.14, printed p.153 (symmetric powers of the
three-dimensional adjoint representation). The exact exercise statements
were inspected in indexed passages from the primary book scans hosted by
university mathematics departments. The publisher record lists the
first-edition print publication in1991; its current electronic record
also displays later reprint/digital dates. G15 supplies the specialized
harmonic covariants and finite Jacobi calculation itself. Full-PDF web
retrieval was too large; the attempted local full-book retrieval was
canceled. No complete-book reading is claimed.

Primary scan locators:
\url{https://websites.math.leidenuniv.nl/edixhoven/teaching/2019-2020/CIMPA/Fulton-Harris.pdf}
and
\url{https://www-fourier.univ-grenoble-alpes.fr/~panchish/ETE\%20LAMA\%202018-AP/lecturesZETAS2018/Fulton\%20Harris_Representation\%20theory\%20first\%20course.pdf}
. Publisher metadata:
\url{https://link.springer.com/book/10.1007/978-1-4612-0979-9} .

\subsection{G-B18 --- Singular linear determinant
sections}\label{g-b18--singular-linear-determinant-sections}

Alper, J., Bogart, T., \& Velasco, M. (2017). A lower bound for the
determinantal complexity of a hypersurface. \emph{Foundations of
Computational Mathematics, 17}(3), 829--836.
\url{https://doi.org/10.1007/s10208-015-9300-x}

Primary author version:
\url{https://sites.math.washington.edu/~jarod/papers/determinantal-complexity.pdf}
. Exact locator used: Theorem1.7, printedp3, and its complete proof,
printedp5. It states the characteristic-zero codimension bound for every
linear determinant restriction and the generic equality; in at least
five variables the corresponding projective hypersurface is always
singular. G19 uses this known theorem, together with the permanent
singular-locus bounds already cited as R-B17. The associated
smooth-section and specified-module deduction is proved in G19. The
published metadata and DOI were corroborated by the author-institution
publication record. No latest or strongest determinantal-complexity
claim is made.

\subsection{G-B19 --- Degree and integral normalization of the
discriminant}\label{g-b19--degree-and-integral-normalization-of-the-discriminant}

Benoist, O. (2012). Degrés d'homogénéité de l'ensemble des intersections
complètes singulières. \emph{Annales de l'Institut Fourier, 62}(3),
1189--1214. \url{https://doi.org/10.5802/aif.2720}

Primary author PDF:
\url{https://www.math.ens.psl.eu/~benoist/articles/degreic.pdf} . Exact
locators: Theorem1.3, equations(1.4)--(1.6), printedpp1190--1191, for
coefficient and variable-change degrees; Theorem1.7, p1191, for the
primitive integer equation and its characteristic-dependent reductions;
Example1.8, p1192, for the hypersurface specialization. The coefficient
degree is k(m−1)\^{}(k−1), and the determinant character exponent is
m(m−1)\^{}(k−1). The second displayed formula in Example1.8 omits the
factor d1 in that exponent; the primary equations(1.4),(1.6) and scalar
dilation give the stated correct exponent. The example was independently
checked visually. This is a source-formula correction, not a cote133
transcription claim. G19 only uses the primitive integer normalization
and characteristic-zero hypersurface properties, without asserting a
reduced discriminant in every finite characteristic. Bibliographic
metadata was checked against the primary NUMDAM publisher archive:
\url{https://archive.numdam.org/articles/10.5802/aif.2720/} .

\subsection{G-B20 --- Pieri on the first-row padding
face}\label{g-b20--pieri-on-the-first-row-padding-face}

Bump, D., \& Gamburd, A. (2006). On the averages of characteristic
polynomials from classical groups. \emph{Communications in Mathematical
Physics, 265}, 227--274. \url{https://doi.org/10.1007/s00220-006-1503-1}

Exact inspected locator: supplied source acp20.tex, section2.1, equation
label e:pieri, lines404--422, together with the horizontal-strip
definition immediately before it. The source cites Macdonald
equationI(5.16),p73 and Bump Theorem42.4; those two books were not
independently reopened for this result. Root and isolation read the
supplied primary statement. G26 uses the known Pieri rule, then proves
its unique-tail consequence and the actual-coordinate-ring injection. No
generic ambient-plethysm equality is asserted.

\subsubsection{G-B14 additional stabilizer locator for
R30}\label{g-b14-additional-stabilizer-locator-for-r30}

Losev2006, proof of Theorem1, printedp2, equations(3)--(4) and the
following sentence: at a moment-zero vector, the compact stabilizer Lie
algebra is a real form of the complex stabilizer Lie algebra. Root
freshly read the complete Theorem4 proof and this Theorem1 proof on the
primary PDF. The real-form conclusion is applied only at the balanced
full tensor inR30, never at an unbalanced block product.

\end{document}
